% Gerado por tools/tei2latex.py a partir de revisado.xml. Regenerar, não editar.

\pg{1} INTRODUCÇÃO

1. \textbf{Philosophia.} --- \emph{Philosophia}, que no sentido etymologico significa \emph{amor da sabedoria}, pode definir-se: \emph{a sciencia que trata das causas supremas dos entes, descobertas pela luz natural da razão.} --- A \emph{Philosophia} --- \textbf{a)} é \emph{sciencia}; porquanto \emph{sciencia}, em geral, é uma serie de conhecimentos certos, coordenados e relativos a um determinado objecto, e tal é a \emph{Philosophia}; --- \textbf{b)} que trata das causas \emph{supremas} dos entes; pois esta sciencia não se occupa das causas mais \emph{obvias} ou \emph{proximas} dos entes, mas sobe ás mais \emph{elevadas} ou \emph{supremas}; --- \textbf{c)} descobertas pela \emph{luz natural da razão}; visto que a \emph{Philosophia} se eleva ao conhecimento das causas \emph{supremas}, não com o auxilio da luz \emph{sobrenatural} da Fé, mas sim da luz propria da razão \footnote{\emph{Causa}, no sentido mais lato, é tudo o que explica de algum modo o ente, e que por isso nos pode levar ao seu conhecimento. --- Em cada ente encontramos muitas e diversas causas, que o constituem, ou de que depende, ou que de algum modo o explicam. Estas causas, em cada genero, estão por tal forma relacionadas e ordenadas, que uma depende de outra e todas de uma primeira, á qual é a razão sufficiente de todas e a todas dá luz e firmeza. Ora a nossa intelligencia não se satisfaz com a superficial observação dos factos, nem com a consideração das causas \emph{proximas ou infimas}, mas eleva-se e progride sempre e não descança, se não chega ás causas \emph{ultimas e supremas}, se não descobre o \emph{ultimo porque} da coisa, de modo que não tenha outras perguntas a fazer. Esta é a \emph{Philosophia}.\par As causas \emph{ultimas} são \emph{intrinsecas} ou \emph{extrinsecas}, conforme se referem á \emph{essencia} ou á \emph{existencia} do ente. --- Tanto as \emph{intrinsecas}, como as \emph{extrinsecas}, podem ser ultimas \emph{relativamente} (num certo genero), ou \emph{absolutamente}. São \emph{intrinsecas e relativamente} ultimas as leis do pensamento na ordem dos actos da intelligencia e as leis moraes na ordem dos actos da vontade; --- são causas \emph{intrinsecas e absolutamente} ultimas as essencias dos entes. A causa \emph{extrinseca e absolutamente} ultima de todos os entes finitos é Deus.\par A \emph{Philosophia} chama-se commummente \emph{sciencia}, e não \emph{sabedoria}. A razão é porque \emph{sciencia} significa, em geral, o conhecimento das causas dos entes, mas a \emph{sabedoria (sapientia)} denota um certo \emph{sabor} ou \emph{gosto}, que experimenta o nosso espirito, quando possue um \emph{perfeitissimo} conhecimento de \pgn{2}uma coisa. Por isso a \emph{sabedoria} é um \emph{perfeitissimo} conhecimento das coisas, pelo qual as proprias coisas são percebidas em todas as suas causas, mesmo nas mais remotas, e em todos os seus principios, mesmo nos mais elevados. Ora como a razão humana, entregue a si mesma, dificilmente pode alcançar esse \emph{perfeitissimo} conhecimento, com muita exactidão e propriedade attribuimos á \emph{Philosophia} a dignidade de \emph{sciencia} e não a de \emph{sabedoria}. --- Poderia objectar-se, que a \emph{Philosophia} é \emph{sabedoria}, porque a \emph{sabedoria} consiste no conhecimento da causa \emph{suprema} e \emph{altissima}, que é Deus, e a \emph{Philosophia} trata de Deus. Respondemos que a \emph{Philosophia} seria simplesmente \emph{sabedoria}, se julgasse das coisas creadas pelo criterio das coisas divinas; porque a \emph{sabedoria} é, propriamente, o conhecimento das coisas pela causa \emph{altissima}. Mas não é assim. Nós não julgamos das coisas creadas pelo criterio das coisas divinas, mas, pelo contrario, julgamos das coisas divinas pelo criterio das coisas creadas; isto é, não conhecemos as creaturas por meio de Deus, mas conhecemos Deus por meio das creaturas, e por isso o nosso conhecimento é \emph{sciencia}, e não \emph{sabedoria}. Tal é a doutrina de S. Thomaz (\emph{Sum. Theol.}, IIa-IIae, q. 9, a. 2 ad 3: q. 45, a. 1). Por isso o nome de \emph{sabedoria} convem propriamente á Theologia revelada e, sobretudo, áquelle dom sobrenatural, que tem esse nome e que consiste em julgar de todas as coisas em conformidade das relações, que ellas têm com Deus. --- Todavia, ás vezes, \emph{sciencia} e \emph{sabedoria} tomam-se no mesmo sentido.}.

\pg{2} 2. \textbf{Objecto da Philosophia.} --- \emph{O objecto}, pois, da Philosophia é constituido pelas \emph{causas supremas} dos entes, emquanto são conhecidas pela luz natural da razão. --- Por isso a Philosophia --- é \emph{uma} só sciencia, porque, embora considere muitos e diversos entes, quer na ordem \emph{subjectiva} quer na \emph{objectiva}, considera-os, comtudo, sob \emph{um} só aspecto, isto é, nas respectivas \emph{causas supremas}, --- e distingue-se não só das sciencias \emph{naturaes}, que se occupam das causas \emph{proximas} dos entes, mas tambem da Theologia \emph{revelada}, que se eleva á Causa \emph{altissima} de todos os entes pela luz \emph{sobrenatural} da Fé \footnote{Em toda e qualquer sciencia deve distinguir-se um duplice objecto: \emph{material} e \emph{formal}. O \emph{material é o ente}, que a sciencia attinge: assim o \emph{corpo humano} é o objecto \emph{material} da medicina; --- o \emph{formal} é a razão, o \emph{aspecto} sob o qual e pelo qual a sciencia attinge o ente; assim a \emph{sanidade} é o objecto \emph{formal} da medicina, que não se occupa do nosso corpo senão emquanto deve ser sanado. Ora o objecto \emph{formal} é propria e verdadeiramente o objecto da sciencia. --- E é este objecto, e não o \emph{material}, que dá ás sciencias a \emph{unidade} e a \emph{especie}; --- a \emph{unidade}, porque a sciencia é e diz-se \emph{una}, quando, embora trate de coisas materialmente varias e diversas, comtudo considera-as sob \emph{um só} aspecto, ou razão; assim a \emph{zoologia} é \emph{uma só} sciencia, porque trata de varios seres sob a unica razão de \emph{animalidade}; --- a \emph{especie}, porque uma sciencia differe de outra, quando, embora considerem o mesmo objecto \emph{material}, comtudo o consideram sob differente aspecto; assim a \emph{anatomia} distingue-se da \emph{physiologia}, porque, embora se occupem do mesmo \emph{corpo humano}, todavia a \emph{anatomia} considera a \emph{estructura} \pgn{3}d'elle, e a \emph{physiologia}, as \emph{funcções}. --- Feita esta distincção, facilmente percebemos --- que o objecto \emph{material} da Philosophia é constituido pelos varios \emph{entes}, que são o \emph{mundo}, o \emph{homem}, Deus, os \emph{actos da intelligencia} e da \emph{vontade}, --- e que o objecto \emph{formal} consiste nas \emph{causas supremas} d'esses varios entes, emquanto são conhecidas pela \emph{luz natural} da razão. Por isso a Philosophia é \emph{uma só} sciencia, porque considera tudo sob o mesmo unico aspecto, e é \emph{especificamente} diversa de todas as outras sciencias, que ou consideram apenas as causas \emph{proximas} dos entes, como as sciencias \emph{naturaes}, ou consideram as causas \emph{supremas}, emquanto são conhecidas pela luz \emph{sobrenatural} da Fé, como a Theologia revelada.\par É, pois, manifesto o erro --- de \emph{Laromiguière} e de \emph{Jouffroy}, quando dizem que a Philosophia não tem um objecto \emph{proprio} e \emph{determinado}, --- e dos \emph{positivistas}, quando affirmam que a Philosophia não é uma sciencia \emph{especial}, mas é uma \emph{classificação} de sciencias. As causas \emph{supremas}, descobertas pela luz \emph{natural} da razão, são o objecto \emph{formal} talmente proprio da Philosophia, que não compete a nenhuma outra sciencia. Da \emph{unidade} e da \emph{especie} d'este objecto deriva a \emph{unidade} e a \emph{especie} d'esta sciencia. --- Os adversarios partem do falso principio que a \emph{unidade} e a \emph{especie} de uma sciencia depende da \emph{unidade} e da \emph{especie} do objecto \emph{material}, e não do objecto \emph{formal}. S. Thomaz diz: «\emph{Illa scientia est una, quae est unius generis subjecti (objecti) formaliter sumpti... Materialis diversitas scibilium non diversificat scientiam, sed formalis}» (\emph{Tot. log. summa}, tr. S, c. 14).}.

\pg{3} 3. \textbf{Divisão da Philosophia.} --- A Philosophia \textbf{divide-se em} tres partes, que são a \emph{Logica}, a \emph{Metaphysica} e a \emph{Moral}. Na verdade, as causas \emph{supremas} dos entes, que esta sciencia \textbf{investiga}, podem referir-se --- ou aos entes \emph{objectivos}, ou \emph{reaes}, que não dependem de nós no seu ser, --- ou aos entes \emph{subjectivos}, que dependem de nós no seu ser, enquanto derivam da nossa \emph{intelligencia}, e são os entes \emph{logicos}, ou derivam da nossa \emph{vontade}, e são os entes \emph{moraes}. Dos entes \emph{objectivos}, considerados \textbf{nas suas} \emph{causas supremas}, occupa-se a \emph{Metaphysica}; dos entes \emph{logicos}, a \emph{Logica}; dos entes \emph{moraes}, a \emph{Moral}. --- A \emph{Metaphysica} subdivide-se em quatro partes, que são a \emph{Ontologia}, a \emph{Cosmologia}, a \emph{Anthropologia} e a \emph{Theodicea}, conforme se occupa do \emph{ente}, considerado \emph{em geral}, ou trata do \emph{mundo}, do \emph{homem}, de \textbf{Deus} \footnote{Diz S. Thomaz: «O officio do homem \emph{sabio} é \emph{ordenar} (isto é, dispor as coisas de modo que uma se refira a outra); porque a sabedoria é a principal perfeição da razão, e é proprio d'esta faculdade conhecer a ordem... Ora a ordem refere-se á razão por \emph{quatro} modos. Por quanto, ha uma ordem, que a razão \emph{não estabelece}, mas unicamente \emph{considera}; tal é a ordem das \emph{coisas naturaes}. Ha uma outra ordem, que a razão, reflectindo, estabelece nos \emph{proprios actos}, quando coordena as suas ideas, e os signaes das ideas, que são as palavras. Ha uma terceira ordem, que a razão, reflectindo, estabelece nas coisas externas, que produz, p. e., numa casa. A ordem, que a razão apenas considera nas coisas, mas não produz, é objecto proprio da\pgn{4}Philosophia \emph{natural} (á qual se reduz a \emph{Metaphysica}). A ordem, que a razão estabelece nos actos proprios, pertence á Philosophia \emph{racional} (que é a \emph{Logica}). A ordem, que a razão estabelece nos actos da vontade, pertence á Philosophia \emph{moral} (que é a \emph{Ethica}). A ordem, finalmente, que a razão estabelece nas coisas externas, que produz, pertence ás \emph{artes mechanicas} (d'ahi a Philosophia das \emph{artes}, chamada \emph{Esthetica}» (\emph{In X Eth. ad Nic.}, I, I. I). D'este modo a Philosophia abrange \emph{quatro} partes, que são: a \emph{Logica} (sciencia da \emph{verdade subjectiva}), a \emph{Metaphysica} (sciencia da \emph{verdade objectiva}), a \emph{Esthetica} (sciencia do \emph{bello}), a \emph{Moral} ou \emph{Ethica} (sciencia do \emph{bem}). Como a Philosophia não trata \emph{ex-professo} das \emph{artes mechanicas}, mas limita-se a expor os principios do \emph{bello}, por isso a \emph{Esthetica} pode reduzir-se á \emph{Metaphysica}, e assim teremos tres grandes partes irreductiveis, que são a \emph{Logica}, a \emph{Metaphysica} e a \emph{Moral}. --- Note-se que o homem, embora seja um só ente, pode, comtudo, considerar-se sob um triplice aspecto, e assim ser objecto das tres partes da Philosophia. Por quanto. --- elle é um ser \emph{real}, e, como tal, é objecto da \emph{Metaphysica}, que o considera na sua natureza, --- é um ser \emph{racional}, isto é, que raciocina, e os actos da sua razão formam o objecto da \emph{Logica}, que os regula para a consecução da verdade, --- é um ser \emph{moral}, e os actos da sua vontade são o objecto da \emph{Moral}, que os \textbf{ordena} para a consecução do fim conveniente.\par A divisão da Philosophia em tres partes: \emph{Logica}, \emph{Metaphysica} e \emph{Moral}, introduzida primeiramente por \emph{Platão}, foi admittida pelos mais celebres escriptores. Cicero diz: «\emph{Fuit jam a Platone accepta philosophandi ratio triplex}: una de \emph{vita} et \emph{moribus}, altera de \emph{natura} et \emph{rebus occultis}, tertia de \emph{disserendo}» (\emph{Acad.}, l. 1, n. 6). O mesmo affirma \emph{S. Agostinho} (\emph{De Civ. Dei}, l. XI, c. 26). S. Thomaz escreve: «\emph{Sunt tres Philosophiae partes, scilicet moralis, logica et naturalis}» (\emph{In c. II Matth.}). --- Muitos auctores dividem a Philosophia em \emph{theoretica} e \emph{pratica}. A \emph{theoretica} é constituida pela \emph{Metaphysica}; a \emph{pratica} abrange a \emph{Logica}, a \emph{Moral} e a \emph{Esthetica}. Tambem esta divisão é excellente. Mas não podemos admittir a divisão em \emph{Logica}, \emph{Dynamilogia}, \emph{Anthropologia}, \emph{Ontologia}, \emph{Cosmologia}, e \emph{Theologia natural}, e muito menos a divisão em \emph{Psychologia}, \emph{Logica}, \emph{Metaphysica} e \emph{Moral},}.

%%>>> t4

\tese{4}{Disposição das partes da Philosophia.}
A \emph{Logica} deve prepor-se ás outras partes da Philosophia; pois ensina as regras, que a nossa intelligencia deve seguir, em toda e qualquer sciencia, para alcançar, provar e defender a verdade. --- A \emph{Metaphysica} deve proceder a \emph{Moral}; pois é pela consideração do mundo, do homem e de \textbf{Deus}, que chegamos a conhecer o principio, o fim e a norma dos nossos actos \emph{moraes}. --- Na \emph{Metaphysica} estuda-se --- primeiramente a \emph{Ontologia}, que, tratando do \emph{ente em geral}, subministra os principios fundamentaes de todas as sciencias, --- depois a \emph{Cosmologia}, porque a primeira realidade, que se offerece á nossa consideração, é o mundo exterior, --- em seguida a \emph{Anthropologia}, não só porque o homem é a synthese de todos os entes visiveis, mas tambem porque elle reflecte sobre si mesmo depois de ter contemplado o mundo, \pg{4} \pg{5}  --- finalmente a \emph{Theodicea}, porque é pela consideração do mundo, e sobretudo do homem, que nos elevamos ao conhecimento da existencia e das infinitas perfeições de \textbf{Deus} \footnote{A ordem, que seguimos na disposição das partes da Philosophia, pode chamar-se e é \emph{doutrinal ou logica}, porque corresponde á natural evolução e ao progresso dos nossos conhecimentos. Pomos em primeiro logar a \emph{Logica}; porque, antes de raciocinar, o homem deve apprehender as leis e o modo do raciocinio. Por isso a \emph{Logica} deve prepor-se não só ás \emph{outras} partes da Philosophia, mas tambem a todas as outras sciencias, porque em todas se emprega o raciocinio. Diz S. Thomaz: « \emph{Prius philosophus debet addiscere logicam, quam alias scientias, quia logica tradit communem modum procedendi in omnibus aliis scientiis} » (\emph{In II Met.}, c. V. l. V). --- Outros escriptores, seguindo a ordem \emph{ontologica}, principiam pela \textbf{\emph{Ontologia}}, porque esta parte da Metaphysica trata das primeiras e universalissimas ideas e propriedades de todas as coisas, e considera os primeiros e supremos principios de todas as sciencias. Mas como poderemos nós raciocinar ácerca do ente e dos seus attributos, se não conhecermos o modo de raciocinar legitimamente? --- Outros, como \emph{Gratry}, seguindo a ordem da \textbf{\emph{dignidade}}, sustentam que o estudo da Philosophia deve principiar pela \textbf{\emph{Theodicea}}, porque, dizem, Deus é o Ser absoluto e necessario, de quem todas as \textbf{coisas} recebem a existencia e a intelligibilidade. Mas estes escriptores confundem a \textbf{ordem} da realidade com a ordem do conhecimento. Deus é o primeiro Ente na ordem da realidade, porque tudo deriva d' Elle; mas não é o primeiro Ente na ordem do conhecimento, porque é pelas creaturas que subimos ao Creador. --- Outros, e são os \emph{Cartesianos}, seguindo a ordem \emph{psychologica}, começam pela \emph{Psychologia} (que é uma parte da \emph{Anthropologia}); porque, dizem, sendo as faculdades da alma os meios, de que, como de instrumentos, homem se serve para alcançar a verdade, é necessario conhecer primeiramente essas faculdades e por isso a essencia da alma, de que derivam. Mas a razão adduzida não é convincente. Porquanto, do facto de as faculdades da alma serem os instrumentos, de que nos servimos para a consecução da verdade, não se segue que, para o recto uso d'essas faculdades, seja necessario conhecel-as na sua natureza e no seu principio; porque, para isso, basta que se admitta a sua existencia e o seu recto e legitimo modo de operar, assim como, para o recto uso de um instrumento, não é preciso que se \emph{conheça} a \emph{qualidade} da materia, de que foi feito, mas basta que se conheça a \emph{sua virtude ou} operação natural, e o modo por que se emprega.}.

%%>>> t5

\tese{5}{Relações entre a Philosophia e as sciencias naturaes.}

a) \emph{A Philosophia é mais excellente que as sciencias naturaes}. --- A excellencia de uma sciencia mede-se pela excellencia do \emph{objecto}, de que trata, e do \emph{fim}, que attinge. Ora o objecto da Philosophia, emquanto \emph{theoretica} ou \emph{especulativa}, é mais excellente que o objecto das sciencias naturaes; pois a Philosophia considera o mundo, o homem, Deus, e considera estes seres nobilissimos nas causas ou razões \emph{supremas} e \emph{remotas}, ao passo que as sciencias \pg{6} naturaes ou não se occupam d'estes objectos, ou se se occupam d'alguns d'elles (do mundo, do homem), limitam-se a descobrir as causas ou razões mais \emph{obvias} e \emph{proximas}: --- como tambem o \emph{fim}, que a Philosophia, emquanto \emph{pratica}, attinge e que é a consecução do fim ultimo, isto é, do Bem infinito na ordem natural, é superior ao \emph{fim}, que as sciencias naturaes attingem e que é a consecução de um bem material ou temporal. Logo a Philosophia é mais excellente que as sciencias naturaes \footnote{Cf. \emph{Summ. Th.}, p. I, q. 1, a. 5. --- Podemos estabelecer a nossa these sobre outra base. O grau de perfeição ou excellencia de uma sciencia depende do seu grau de abstracção. Uma sciencia é tanto mais perfeita, quanto mais é abstracta. Ora nenhuma sciencia é mais abstracta que a Philosophia: e por isso esta excede, em excellencia, todas as outras sciencias. Diz \emph{Taine}: «Os tres quartos dos homens consideram como especulações ociosas as grandes syntheses do pensamento: e erram, porque a formação d'essas syntheses é o escopo da vida de um povo ou de uma edade, e é por ellas que a humanidade chega ao seu pleno desenvolvimento. Se um habitante dos outros planetas descesse a perguntar-nos onde se encontra a nossa especie, seria necessario mostrar-lhe as cinco ou seis ideas, que possuimos acerca do espirito e do mundo: estas dar-lhe-hiam o conceito exacto dos nossos conhecimentos »  (\emph{Le positivisme anglais}, p. 11-12).}.

b) \emph{A Philosophia é necessaria para a comprehensão das sciencias naturaes.} --- As causas dos entes estão por tal modo coordenadas entre si, que as mais \emph{obvias} ou \emph{proximas} dependem das \emph{supremas} ou \emph{remotas}. Sendo assim, é claro que as sciencias naturaes, que têm por objecto as mais \emph{obvias} ou \emph{proximas} causas, somente na Philosophia, que tem por objecto as causas \emph{supremas} ou \emph{remotas}, podem encontrar a ultima razão das suas incognitas, o termo ultimo das suas investigações; porque as razões \emph{proximas} dependem das \emph{supremas}, e d'estas recebem intelligibilidade e certeza. Logo a Philosophia é necessaria para a comprehensão das sciencias naturaes \footnote{A Philosophia, tendo por objecto proprio e especifico as causas \emph{supremas} das coisas, descobre e ensina os principios primeiros e universaes, que devem servir de fundamento a todas as sciencias. Por isso ella, além de ser o vinculo, que liga e reune numa \emph{synthese} grandiosa e admiravel todas as sciencias humanas, é absolutamente necessaria para a \emph{comprehensão}, isto é, para o perfeito conhecimento das outras sciencias, sobretudo da \emph{Mathematica}, da \emph{Physica}, do \emph{Direito positivo}. --- É necessaria para a comprehensão da \emph{Mathematica}. Na verdade, a \emph{Mathematica} procede sempre pela \emph{abstracção} e pelo \emph{raciocinio}. Mas é a Philosophia, que, além de dar á \emph{Mathematica} os \emph{primeiros principios}, que são a base de toda a sciencia, forma o espirito para o \emph{raciocinio} e para a \emph{abstracção}. --- É necessaria para a\pgn{7}comprehensão da \emph{Physica}. Por quanto, a \emph{Physica}, por se basear na \emph{observação} e na \emph{inducção}, e estudar as causas \emph{proximas} dos phenomenos sensiveis, não pode prescindir da Philosophia, que determina as condições e as regras da \emph{observação} e da \emph{inducção}, e se eleva á Causa Primeira e absolutamente Suprema. --- É necessaria para a comprehensão do \emph{Direito positivo}. Na verdade, fundando-se o Direito \emph{positivo} no \emph{natural}, é claro que ninguem poderá estudar proficientemente aquelle, sem conhecer este. Ora o Direito \emph{natural} é o objecto da Philosophia \emph{moral}.}.

%%>>> t6

\pg{7} \tese{6}{Relações entre a Philosophia e a Theologia revelada.}

a) \emph{A Philosophia é uma sciencia independente da Theologia revelada.} --- Uma sciencia é independente de outra, quando não recebe d'esta --- nem os \emph{principios}, em que se funda, --- nem as \emph{conclusões}, que deduz, --- nem os \emph{meios}, de que se serve para deduzir dos principios as conclusões. Ora a Philosophia não recebe da Theologia revelada --- nem os seus \emph{principios}, porque estes são conhecidos pela luz natural da razão, --- nem as suas \emph{conclusões}, porque estas são deduzidas d'esses principios, --- nem os seus \emph{meios}, que são o \emph{syllogismo} e a \emph{inducção}, porque estes são formados á luz dos mesmos principios racionaes. Logo a Philosophia é independente da Theologia revelada \footnote{A Philosophia pertence ao grupo das sciencias \emph{racionaes}, isto é, d'aquellas sciencias, que a razão forma com as suas forças naturaes, e por isso é uma sciencia autonoma, independente de toda e qualquer sciencia, mesmo da Theologia revelada. De facto, a Philosophia existia já, quando appareceu a divina luz do Christianismo, e não pode dizer-se que, ao appa-recer d'essa luz, ella perdesse a propria independencia; visto que N. S. Jesus Christo veio a este mundo, não para destruir a natureza, mas para a aperfeiçoar e ennobrecer. E todos sabem que, quando, no meado do seculo XIX, dois escriptores francezes, \emph{de Bonald} e \emph{Lamennais}, pretenderam obrigar a razão humana a pedir os primeiros principios e os primeiros motivos de certeza á Revelação sobrenatural, a Egreja, não só não acceitou, mas reprovou e condemnou essa pretenção, mais generosa do que sabia. A ordem natural não pode fundar-se na ordem sobrenatural; cada uma tem as suas bases. --- Mas d'ahi não se segue que a Philosophia deva prescindir da Theologia revelada, ou da Fé. Esta permitte á razão do philosopho seguir os proprios principios e as suas livres tendencias, sem que seja dirigida, de um modo \emph{positivo} e \emph{directo}, pela Revelação, comtanto que evite de contradizer as verdades divinamente manifestadas e de invadir uma \emph{esphera} superior. É uma subordinação \emph{indirecta}, que não se opõe á \emph{autonomia}. Diz o Concilio Vaticano: «A Fé não prohibe ás sciencias de servirem-se, \emph{cada uma na sua esphera}, dos seus \emph{principios proprios} e do seu \emph{methodo particular}; mas, embora reconheça sempre esta justa liberdade, vela com cuidado para que não aceitem erros, mettendo-se em opposição com a doutrina divina, nem invadam e perturbem, sahindo dos seus limites, a \emph{esphera} e as verdades da Revelação» (\emph{Const. Dei Filius}, c. II).}.

\textbf{b)} A \emph{Philosophia é inferior}, \emph{em dignidade}, á \emph{Theologia revelada}. --- É inferior quanto ao \emph{objecto}; porque a Philosophia, embora, pela luz natural da razão, se eleve ás causas supremas e até á Causa absolutamente primeira, comtudo não descobre senão poucas verdades, relativas a Deus, e de um modo muito imperfeito; ao passo que a Theologia revelada, guiada pela luz sobrenatural da Fé, descobre muitas outras-verdades, relativas a Deus, que transcendem a capacidade de toda a intelligencia creada. --- É inferior quanto á \emph{certeza}; porque a Philosophia é guiada, no conhecimento dos principios e na deducção das conclusões, pela luz natural da razão, que pode enganar-se e muitas vezes se engana; ao passo que a Theologia revelada é guiada --- no conhecimento dos principios, que são os artigos da Fé, pela luz da sciencia divina, que não se engana e não pode enganar, --- e na deducção das conclusões, pelo menos das que se relacionam com esses artigos, pelo magisterio infallivel da Egreja. --- É inferior quanto ao \emph{fim}; porque a Philosophia dispõe o homem para o fim ultimo \emph{natural}, que consiste na contemplação de Deus pelas creaturas; ao passo que a Theologia revelada ordena o homem para o fim ultimo \emph{sobrenatural}, que consiste na visão \emph{intuitiva} de Deus. Logo a Philosophia é inferior, em dignidade, á Theologia revelada \footnote{S. Thomaz (\emph{Sum}. \emph{Th}. P. I, q. 1, a. 5) demonstra que a Theologia revelada, especulativa e pratica, é superior, em dignidade, a todas as outras sciencias, especulativas e praticas, que se guiam pela luz natural da razão. --- Adverte o S. Doutor que as duvidas ou incertezas, que se levantam no nosso espirito acerca de um ou outro artigo de Fé, nascem, não da incerteza \emph{objectiva} dos mesmos artigos, mas da fraqueza \emph{subjectiva} da intelligencia creada (\emph{Ib.,ad} 1).}.

%%>>> t7

\pg{8} \tese{7}{Utilidade da Philosophia.}
A Philosophia é utilissima para o \emph{individuo} e para a \emph{sociedade}.

\textbf{a)} \emph{É utilissima para o individuo}. --- Porquanto, a Philosophia aperfeiçoa as mais nobres faculdades, de que o homem é dotado e que são a \emph{intelligencia} e a \emph{vontade}. --- Aperfeiçoa a \emph{intelligencia}, elevando-a ao conhecimento das causas supremas dos seres creados, e até ao conhecimento do proprio Deus. --- Aperfeiçoa a \emph{vontade}, guiando-a, pela pratica das virtudes moraes, á consecução da felicidade eterna, ainda que na ordem natural. Logo a Philosophia é utilissima para o \emph{individuo} \footnote{Que a Philosophia aperfeiçoe a \emph{intelligencia}, é uma verdade, que\pgn{9}todos admittem facilmente. Mas, nem todos pensam em que a Philosophia tende tambem a aperfeiçoar a \emph{vontade}, o que é ainda mais importante. A rectidão da vontade é immensamente mais preciosa, mais util, mais necessaria que a illustração da intelligencia. Toda a perfeição da nossa vontade, mesmo na ordem natural, resume-se e consiste no amor do summo e infinito Bem. Ora este amor desperta-se e augmenta á vista das perfeições creadas, emquanto estas são um reflexo das infinitas perfeições de Deus. Se são amaveis as creaturas, infinitamente mais amavel é o Creador.}.

\textbf{b)} \emph{E' utilissima para a sociedade}. --- Por quanto, a sociedade, sendo uma reunião de individuos para o alcance de um bem commum, exige o concurso efficaz e proporcionado de todos os seus membros. Tal concurso dos membros da sociedade consiste na somma das forças, especialmente moraes, com que elles conspiram para o fim e que derivam da applicação dos principios da verdade e da moralidade. Ora os principios da verdade e da moralidade, de cuja applicação derivam as forças moraes, são ensinados e inculcados pela Philosophia. Logo a Philosophia é utilissima para a \emph{sociedade} \footnote{Os costumes correspondem ás ideas. Se isto se verifica nos individuos, muito mais se verifica na sociedade; porquanto, se os individuos podem estar e estão, mais ou menos, em contradicção comsigo mesmos, as multidões seguem sempre o impulso dos seus principios. Negae Deus, o livre arbitrio, a propriedade; e vereis a que pode chegar e chega uma sociedade. D'ahi a insipiencia dos homens politicos, que prohibem a liberdade da acção, e depois permittem a liberdade do pensamento. É a Philosophia que dirige o mundo. O conflicto dos interesses e das paixões nasce, muitas vezes, de uma idea, que um escriptor lançou no seio da sociedade e que parecia dever ficar sempre no campo da abstracção. --- A Philosophia é sempre a luz da historia. Se quizermos conhecer a civilisação de um povo ou de uma epocha, devemos examinar a maneira por que esse povo ou essa epocha comprehendeu o direito, a auctoridade, a liberdade, etc. É na Philosophia de Socrates, de Platão e de Aristoteles que deve procurar-se a intelligencia da historia grega nos seculos \emph{V} e VI antes de Christo, como nas theorias d' Epicuro se encontra a explicação da corrupção e da decadencia, que se seguiram. As glorias do sec. XIII coincidem com a grande epocha da Philosophia scholastica. Foi o espiritualismo christão de Descartes, de Bossuet, de Malebranche, de Leibnitz e de tantos outros, que fez o seculo de Luiz XIV. É impossivel ter uma idea exacta da revolução franceza e dos principios, que a inspiraram, se não se conhecer a Philosophia do seculo XVIII. A theoria que o \emph{capital é um roubo} não tem sido e não é hoje, por ventura, a causa da revolução social, que ameaça cobrir o mundo de sangue e de ruinas? --- Podemos, portanto, concluir que a historia do mundo é a historia das suas ideas.}.

%%>>> t8

\pg{9} \tese{8}{Disposições para o estudo da Philosophia.}
As principaes são tres: \emph{oração humilde} e \emph{fervorosa}, \emph{grande pureza de}

\emph{coração}, \emph{ardente amor da verdade}. --- É necessaria a \emph{oração}; porque, sendo Deus o \emph{Senhor das sciencias}, \emph{o Pae das luzes}, a Elle deve recorrer quem deseja a sapiencia. --- É necessaria a \emph{pureza de coração}; porque um coração, que se arrasta no lodo das mais vis paixões, não pode elevar-se á contemplação de coisas nobilissimas, como são as de que trata a Philosophia. --- É necessario \emph{o amor da verdade}; porque só este amor provoca o desejo de saber, conforta as forças e vence as difficuldades, subordina as ideas ás coisas e não as coisas ás ideas, e, se detesta o erro, ama sinceramente o errante.

%%>>> t9

\pg{10} \tese{9}{Fontes da Philosophia.}
Sendo a Philosophia o conhecimento natural dos entes pelas causas supremas, é claro, que as suas legitimas fontes devem ser os livros dos escriptores, que trataram, com verdade e sapiencia, os grandes problemas, relativos ao mundo, ao homem, a \textbf{Deus}. Sobre todos os escriptores levanta-se, como o sol entre os outros astros, o grande S. Thomaz de Aquino, a quem, pela sublimidade do talento e pela innocencia da vida, a Egreja deu o nome de \emph{Doutor angelico}. Tudo o que de verdade se encontra nos escriptos dos Philosophos pagãos, dos Padres e Doutores da Egreja, e de todos os sabios, que o precederam, S. Thomaz não só conheceu perfeitamente, mas augmentou, completou e ordenou com tal firmeza de principios, com tal severidade de raciocinio, com tal perspicuidade e propriedade de linguagem, que a sua doutrina não só resolve e esclarece os mais arduos problemas, mas se adapta maravilhosamente ás necessidades de todos os tempos, e tem a força de refutar todos os erros, passados, presentes e futuros. Levado por esses motivos, o Santo Padre Leão XIII declarou S. Thomaz \emph{Patrono de todas as escholas catholicas}\footnote{Alguns dizem que a Philosophia de S. Thomaz envelheceu, que não corresponde ás necessidades dos nossos tempos, nem conserva a força e a efficacia necessaria para confutar os adversarios da verdade. Mas erram. S. Thomaz não deduz as suas conclusões senão dos principios e da essencia das coisas; e, como os principios são immutaveis e as essencias necessarias, tambem as conclusões são immutaveis e necessarias. Por isso a doutrina do Angelico satisfaz ás necessidades e exigencias da mente humana em todo e qualquer tempo, e tem a efficacia de confutar todas as difficuldades dos adversarios, em todos os tempos. A Philosophia scholastica, de que S. Thomaz é o principal representante, foi chamada justamente a \emph{Philosophia do senso commum}; e, por isso, deve ser seguida por todos os que \pgn{11}ainda não renunciaram a esse \emph{bom senso}. --- Nem se diga que a Philosophia de S. Thomaz é contraria ás descobertas e ao progresso das \emph{sciencias naturaes}. S. Thomaz não só não é contrario ás \emph{sciencias naturaes}, mas, ensinando que a nossa intelligencia não pode elevar-se ao conhecimento das coisas immateriaes senão pelo conhecimento das coisas materiaes, ensinou, ao mesmo tempo, com a \emph{palavra} e com \emph{o exemplo}, que o philosopho deve investigar diligentemente os segredos da natureza e occupar-se do estudo das sciencias physicas. Os melhores naturalistas confessam que entre os principios da Escholastica e as descobertas \emph{certas} das sciencias naturaes não só não existe verdadeira opposição, mas, pelo contrario, reina a mais perfeita harmonia: o que teremos occasião de mostrar no decurso d'estes elementos. E se houvesse, num caso determinado, uma verdadeira e averiguada opposição (o que não é impossivel, pois as sciencias naturaes não tinham, no seculo XIII, o desenvolvimento, que têm hoje), o philosopho, amigo unicamente da verdade, não deveria ter duvida em rejeitar uma conclusão, embora defendida pelos mais afamados mestres, ou pelo proprio S. Thomaz. É esta a norma inculcada pelo mesmo S. P. Leão XIII na celebre Encyclica \emph{Aeterni Patris} (4 agosto 1879).}.

LOGICA

%%>>> t10

\pg{11} \tese{10}{Logica.}
\emph{Logica} (do grego \emph{logos} --- \emph{discurso}, \emph{razão}) é a sciencia \emph{que dirige}, \emph{por meio de leis}, \emph{as operações da nossa razão}, \emph{para que ordenada e facilmente alcancemos} a \emph{verdade}.

Todos os homens têm uma certa aptidão natural para

dirigir e coordenar as suas operações intellectuaes. É a \emph{Logica natural}. --- Mas essa aptidão não era sufficiente. Foi necessario que, por meio da reflexão e da analyse, se formasse um systema completo de leis, reguladoras do pensamento; de modo que chegassemos á acquisição da verdade por um caminho facil e direito. Este systema de leis é a \emph{Logica artificial}, ou \emph{scientifica}, de que nos occupamos \footnote{O homem nasce com um ardente desejo de conhecer a verdade. Juntamente com esse desejo o sapientissimo Creador deu-nos uma aptidão proporcionada, um criterio natural, pelo qual nos servissemos, com rectidão, das faculdades cognitivas, e descobrissemos, anteriormente a todo e qualquer magisterio ou ensino, algumas leis ou regras fundamentaes, que nos ajudassem a conseguir a verdade. Estas regras, assim descobertas, constituem a \emph{Logica natural}. A ella devemos o conhecimento de muitas verdades, que são dotadas de evidencia immediata, e de muitas conclusões, que espontaneamente derivam d'essas verdades. --- Mas a Logica \emph{natural}, embora muito util, não basta. Por quanto, a verdade está, por vezes, tão escondida aos nossos olhos e tão cercada de difficuldades, que não é facil percebel-a. Por \pgn{12}isso, podendo aperfeiçoar-se a natureza, pela arte, foi necessario aperfeiçoar essas regras da Logica \emph{natural}, completando-as e coordenando-as de modo, que a nossa intelligencia podesse --- produzir os seus actos com facilidade, sem confusão e sem erro, --- fazer bem as demonstrações, embora complexas, --- deduzir recta e legitimamente as conclusões, remota e confusamente contidas nos principios, --- defender a verdadeira doutrina das objecções, apresentadas sob as apparencias da verdade. --- Nem este trabalho de aperfeiçoamento é difficil para o homem. Dotado de reflexão, pode elle estudar e analysar os actos da propria intelligencia e descobrir a maneira de os produzir com facilidade, com ordem e sem erros. O conjuncto d'essas regras, que foram descobertas pela reflexão e completam as da Logica \emph{natural}, constitue a Logica \emph{artificial}, ou \emph{scientifica}. D'onde se vê que a Logica \emph{artificial} se distingue da \emph{natural}, como o perfeito se distingue do imperfeito. --- S. Thomaz define assim a Logica \emph{artificial}: « \emph{Scientia}, \emph{quae est directiva ipsius actus rationis}, \emph{per quam scilicet homo in ipso actu rationis ordinate et faciliter et sine errore procedit} » (\emph{Post. anal., lib. I, l. 1}).}.

11. \textbf{Objecto da Logica}. --- O \emph{objecto} da \emph{Logica} é constituido pelas nossas \emph{operações intellectuaes}, \emph{emquanto devem ser dispostas} e \emph{ordenadas para a consecução da verdade}. --- D'ahi a differença entre a \emph{Logica} e as demais sciencias; pois nenhuma d'estas estuda o pensamento para o \emph{dirigir} para o seu fim \footnote{O \emph{objecto} da Logica, como o de qualquer outra sciencia, divide-se em \emph{material} e \emph{formal}. O \emph{material}, que é \emph{tudo} o que a sciencia attinge, é constituido pelas \emph{nossas operações intellectuaes}; o \emph{formal}, que é \emph{o aspecto}, sob \emph{o} qual a sciencia attinge ou considera o objecto \emph{material}, é a \emph{ordem}, que a razão deve estabelecer nas suas operações para alcançar a verdade. --- Em todas as sciencias, é indispensavel a \emph{ordem} entre os varios entes, de que ellas tratam; mas nenhuma d'ellas, exceptuada a \emph{Logica}, tem por \emph{objecto} a propria \emph{ordem}. Por isso a Logica poderia definir-se: a sciencia da \emph{ordem} nas operações intellectuaes. Diz S. Thomaz: « \emph{Ordo}, quem ratio considerando \emph{facit} in proprio actu, pertinet ad \emph{rationalem} philosophiam » (\emph{In I Eth.}, l. 4).\par Por isso a \emph{Logica} não só differe de todas as demais sciencias racionaes, mas tambem das outras partes da Philosophia, com as quaes parece identificar-se e que são a \emph{Dynamilogia} (tratado das \emph{faculdades}), a \emph{Ideologia} (tratado das \emph{ideas}), e a \emph{Ontologia}. --- Differe da \emph{Dynamilogia}, pois esta considera as \emph{faculdades}, emquanto são \emph{realidades derivadas} da essencia da alma ou do composto: mas a \emph{Logica} considera essas faculdades, emquanto são \emph{meios para} alcançar a verdade. --- Differe da \emph{Ideologia}: porque esta trata das \emph{ideas}, emquanto são \emph{actos} ou \emph{productos vitaes}, determinando a sua \emph{natureza}, \emph{propriedades} e \emph{origem}: ao passo que a \emph{Logica} falla das ideas, emquanto \emph{devem ser ordenadas} para a consecução da verdade. --- Differe da \emph{Ontologia}; porque esta attinge a verdade emquanto \emph{objectiva}, isto é, emquanto existe na \emph{realidade}, ao passo que a \emph{Logica} considera a \emph{verdade} emquanto \emph{subjectiva}, isto é, emquanto é \emph{termo} dos nossos conhecimentos.}.

12. A \textbf{Logica} é arte e \textbf{sciencia}.

a) É \emph{arte.} --- Por quanto, \emph{arte} é o conjuncto de regras ou \pg{13}  preceitos, que ensinam a executar. com perfeição, alguma coisa. Ora a Logica é um conjuncto de regras ou preceitos, que ensi- -nam a fazer, com perfeição, a \emph{definição}, a \emph{divisão}, o \emph{raciocinio}, etc. Logo a Logica é \emph{arte} \footnote{A Logica não só é \emph{arte}, mas é a \emph{arte das artes}. porque dirige o acto da razão, do qual derivam todas as artes. Diz S. Thomaz «A razão ' pode dirigir não só os actos das outras faculdades, : mas até os seus proprios actos. Por quanto, á faculdade intellectual tem o- poder. de reflectir sobre si mesma; poisque a inlelligencia entende-se a si mesma e a3 razão pode raciocinar ácerca do seu acto. E por isso, assim como, pelo (acto de x razão raciocinar ácerca do acto das mãos, foi inventada a arte do pedreiro e do exrpinteiro, pela qual o homem pode exercer, com (acilidude e com ordem, os actos d'estas artes, assim tambem, pelo mesmo motivo, é necessaria uma arte, directiva do proprio acto da razão, pela qual o homem, nos seus raciocinios, proceda com \emph{ordem}, com \emph{facilidade} e sem \emph{erro}. Esta arte é a \emph{Logica}, isto é, a sciencia \emph{racional}. A qual é \emph{racional}, não só porque é conforme a razão (o que é commum a todas as artes;,, mas tambem porque se refere ao acto da razão, como a sua propria materia. Per isso pode chamar-se « \emph{arte das} artes, porque dirige o acto da razão, do qual derivam 7 todas ás artes» (I \emph{Post}. \emph{Analyt}., l. 1). ! (2) A Logica não só é sciencia, mas é a sciencia \emph{das} sciencias, porque dirige a propria razão, da qual derivam Lodas as sciencias. E é uma sciencia. \emph{especial}, isto é, distineta de Lodas as outras sceiencias. Por quanto, as sciencias distinguem-se pelos seus objectos \emph{formaes}. Ora o objecto \emph{formal} da Logica, que, como dissemos, é x \emph{ordem} nos actos intellectuaes, distingue-se do-objecto- \emph{formal} das-outras seteneizs, (que se oceupamr de uma secte de verdades, relativas a uma coisa, mus nio se occupam do modo de conseguir a verdade. --- Por isso a Logica não é o instrumento \emph{proximo} e \emph{proprio} de cada sciencia, mas é o instrumento \emph{remoto} e conmmum a todas as sciencias: porque presereve as leis do raciocinio, que se hão de observar em toda c qualquer sciencia. --- Concluimos que az Logica é, ao mesmo tempo, \emph{arte} e \emph{sciencia}, mas sob differentes aspectos. X arte, porque dá as normas para a recta formação do \emph{raciocinio}, etc. ; --- é \emph{sciencia}, porque dá a razão das normas, que prescreve, Logo a Logica é uma \emph{sciencia pratica}. (Cf. \emph{Sunmm}. \emph{Th}. I*-1 5, q. 57, a. 3 ad 3). --- De tudo o que deixamos dito resulta. : \emph{exvellencia}. da Logica. Todavia esta é inferior, em dignidade, ás outras partes \emph{da Philosophia}\pgn{14}. Com \emph{effeito}, o que se procura por amor de outra coisa é inferior ao que se procura por si mesmo; pois \emph{o} meio é inferior ao fim. Ora a instrucção, que se procura no estudo da. Logica, é um meio para o conhecimento perfeito das outras partes da Philosophia, e, em geral, de todas as sciencias.}. m

b) É \emph{sciencia}. ---- Com effeito, é proprio da \emph{sciencia} deduzir dos principios intrinsecos, que constituem a essencia de um objecto, as propriedades ou os attributos do mesmo objecto, e, se se trata de dar regras, deduzil-as dos proprios principios ' das coisas, a que essas regras se referem; poisque a sciencia é o conbecimento das coisas pelas suas causas. Ora a Logica deduz dos principios, que constituem a natureza do \emph{raciocinio}, da \emph{definição}, da \emph{divisão}, ele., as regras para a receta formação do mesmo \emph{raciocinio}, etc. Logo a Logica é \emph{sciencia} ().

%%>>> t13

\pg{12} \pg{14} \tese{13}{Necessidade da Logica.}
O estudo da \emph{Logica} é summamente necessario. Porquanto, a nossa intelligencia nunca poderá alcançar um \emph{profundo} e \emph{completo} conhecimento de qual- ' quer sciencia, se não souber ordenar as. suas operações, exLremar a verdade do erro, formar raciocinios, julgar da legitimidade dos argumentos. Ora tudo isto constitue o objecto da \emph{Logica}. Logo o estudo da \emph{Logica} é summamente necessário (\textasciicircum{}).

%%>>> t14

\tese{14}{Divisio da Logica.}
As operações intellectuaes, que constituem o objecto da Logica, podem considerar-se, --- ou quanto ii \emph{ordem}, que deve existir entre ellas, para a consecução da verdade, --- ou quanto á \emph{realidade objectiva}, que ellas exprimem. D'ahi a divisão da. Logica em \emph{formal} e \emph{material}, --- A Logica \emph{format} occupa-se da \emph{ordem}, com que devem ser dispostas, e das \emph{leis}, ás quaes devem estar sujeitas as operações inlelleeluaes, que são a \emph{idem}, \emph{o juizo} e o \emph{raciocinio}. --- A Logica \emph{material} occupa-se da \emph{realidade objectiva}, que as nossas operações exprimem e que é a \emph{verdade}; e por isso occupa-se --- da” \emph{natureza} da verdade, --- dos meios, que empregamos para a acquisição da verdade (\emph{faculdades}) --- do \emph{caminho}, que percorremos (\emph{methodo}) e do \emph{resultado}, que alcançamos (\emph{sciencia}). --- A Logica, pois, trata-se em \emph{duas} secções. A primeira secção é dedicada \& Logica \emph{formal}; a \emph{segunda}, \emph{á material} 0).\footnote{A maior parte. senão á totalidade dos erros, que se encontram nos livros dos modernos escriptores, sobretudo naturalistas, derivam da ignoranceia da \emph{Logica}; Apresentam elles, ás vezes, hypotheses. arbitrarias, jnducções incompletas, factos não averigzuadós; outras vezes, de factos certos e averiguados deduzem consequencius falsas, E qual é o motivo verdadeiro d' estes desvarios ? É a ignoraúcia, se não é o desprezo da \emph{Logica}. Sem o estudo accurado d' esta sciencia, o homem, se pode chegar ao conliecimento das verdades \emph{comuns}, que se percebem immediatamente, pela força do hom senso, não pode alcançar as outras verdades senão misturadas com muitos erros e depois de muitas difficuldades. (Ctr, Cicero, \emph{Orat}. IV, 16).}\footnote{A Logica \emph{forma}! chama-se \emph{formal}, porque a \emph{ordem}, de que ella se occupa, é a \emph{forma} das nossas operações. Diz-se tambem --- \emph{c«nalgica}, porque \emph{analysa} o pensamento e tambem porque se funda em principios \emph{analyticos}, --- \emph{dialectica}, porque ensina a raciocinar, --- \emph{subjectiva}, porque trata dos actos do sujeito pensante, --- menor, porque trata da \emph{forma} das nossas operações intellectuaes, que é inferior á sua \emph{materia} (\emph{a} verdade), e porque\pgn{15}trata de questões menos difficeis c menos elevadas. --- A \textbf{Logica} \emph{material} chama-se assim, porque « verdade é a \emph{materia}, ou o conteudo dos pensamentos. Diz-se tambem --- \emph{critica}, porque trata dos \emph{criterios}, isto é, dos meios uteis para alcançar, com certeza, a verdade, --- \emph{objectica}, porque se occupa da \emph{verdade}, que é o \emph{objecto}, ou. a \emph{materio} do pensamento, --- \emph{maior}, porque estuda a \emph{materia} do pensamento, que é a \emph{verdade} e que é superior á sua \emph{forma}, c porque trata das questões mais difficeis e elevadas. --- Esta divisão da Logica em \emph{formal} e \emph{material} lunda-se na theoria ácerca da \emph{materia} é da \emph{forma}, que Aristoteles e os escholasticos expazeram em confonnidade com o \emph{semso} commum dos homens. --- \emph{Kant} adoptou as palavras, mas alterou o sentido d'ell:s, Dividiu tambem elle « Logica em \emph{formal} e \emph{material}; mas i Logica \emph{formal}, para elle, é a que se occupa de umas formas \emph{stójectóvas} ou \emph{ideas} do pensamento, inteiramente independentes dz realidade e anicriores : Lodo e qualquer dado da experiencia. Para os homens de bom senso, a Logica \emph{formal}, tratando da \emph{forma} do pensamento, € por isso dos varios modos de combinar as \emph{ideas}, os \emph{juizos}, OS raciocinios, se prescinde da \emph{realidade objectiva} do pensamento, não a \emph{nega}, nem \emph{falseia}, como fazem os \emph{Kantístas}. --- Auctor \emph{e} pace da Logica \emph{scientifica} foi \emph{Aristoteles}, no seu. Organum, --- obra perfeita, a qual, segundo a \textbf{confissão} do proprio \emph{Kant}, nada foi depois acerescentado. Na edade média, os \emph{escholasticos}, e sobretudo S. Thomaz, expuzeram-na e commentaramena. A \emph{Logica}, que teve \emph{Aristoteles} por fundador, é a \emph{formal}.\par (\emph{f}) Todas as operações di nossa mente, para alcançar a verdade, resumem-se nus tres indicadas. Na verdade, quando percebemos alguma coisa, ou emitlimos o nosso juízo, ou não. Não ha outra hypothese. Se não emitlimos o nosso juizo, é mero \emph{apprehenszo}, \emph{é idea}. Se emittimos um juizo. este --- ou emille-xe \emph{inunediatamente}, pela simples comparação do predicado com \emph{o} sujeito, e fazemos um \emph{juizo}, --- ou deduz-se de um juizo precedente, € lemos \emph{o raciocinio}. Qutrus operações, que alguns escriptores enumeram, são apenas \emph{modalidades} de uma ou outra das tres mencionadas, e reduzem-se a estis. (Cfr. S. Thom., in 1 \emph{Periher}., 1. 1).\par A causa, pela qual a intelligencia humana, para aleangar a verdade, precisa de executar varias e diversas operações, é a sua limitação. Se ella,}

» Operações da intelligencia. --- A intelligencia humana não apprepende sempre a verdade de um modo intuitivo e simul-

laneo, mas porgraus eom momentos successivos, executando \emph{Zres} diversas Operações, que são a \emph{idea}, O \emph{juizo} e O \emph{raciocinio}. Estas operações estão por tal fórma dispostas, que a \emph{deze} se refere ao \emph{juizo}, e este ao \emph{raciocinio}; pois o \emph{juizo} é composto de \emph{ideas}, e o \emph{raciocinio} de \emph{juízos}. A expressão verbal da \emph{ides} é o \emph{termo}; do \emph{juizo}, é a \emph{proposição}: do \emph{raciocinio}, é a \emph{argumentação} \footnote{\pgn{16}percebendo um objecto. podesse comprehender, com uma simples intuição, tudo o que. a esse objecto se refere, não haveria necessidade de formar \emph{juizos}; --- como tambem, se num principio conhecido podesse distinctamente conhecer todas as conclusões, virtual e implicitamente contidas nesse principio, não haveria logar para \emph{raciocinios}. Mas a. intelligencia humana não é dotada de trnla perfeição (\emph{Sum}. \emph{Theol}, p. 1, q. 98, aa. 3, 4).}.

\emph{TE} - \emph{oA AMA IUe} nem Yat 16 LOGICA --- SECCAO PRIMEIRA.

16. Divisão \textbf{da Logica} formal --- Dividimos a Logica \emph{formal} em \emph{tres} capitulos: no \emph{primerio}, tratamos da \emph{idea} e do \emph{termo}; no \emph{segundo}, do \emph{juizo} e da \emph{proposição}: no \emph{terceiro}, do \emph{raciocinio} e da \emph{argumentação}.

%%>>> cap-1

\chapter*{CAPITULO PRIMEIRO\\[0.3em]\large IDEA E TERMO}
\addcontentsline{toc}{chapter}{CAPITULO PRIMEIRO --- IDEA E TERMO}

\pg{15} \pg{16} SussAnt0: --- dia, sua definição e unalyse. ---- Divisão geral da idea. ---

Divisio da idea universal-reflexa. --- Reduccüo das idens a eulezorias,

--- “Termo, suas especies e propriedades. --- Definição e divisio.

Amrigo ]

\emph{Idea}, \emph{sua definigao e analyse}

17. idea. --- \emph{Idea é a simples representação intellectual de uma coisa}; assim, percebendo o \emph{homem}, temos uma \emph{idea}. --- A \emph{idea} --- é \emph{representação}, porque é uma Imagem, que exprime e reproduz a coisa; --- é representação \emph{intellectual}, porque é uma imagem gravada na \emph{intelligencie}, e. não na \emph{phantasia}, ou numa outra faculdade sensitiva; --- é \emph{sümples} representação, porque a coisa, reproduzida por ella, não se attribue nem se nega propriedade alguma \footnote{A reproducção ou representação de uma coisa.na \textbf{nossa} intelDgeneia Lem o \emph{nome} de \emph{idea}, porque nessa. representação a inlelligenceia vé a propria coisa, --- A idea chama-se tamhem --- \emph{verbo mental}, porque a nossa intelligencia, informada pela idea, falla e diz a si mesma a coisa pereehida, --- \emph{conceito}, porque é uma coisa, que é gerada pela propria intelligencia e que começa a viver a vida intellectual do sujeito cognoscente, --- \emph{fermo mental}, porque todos os actos intellectuues se resolvem nas ideas, como” nos seus ultimos elementos, --- \emph{percepção} ou \emph{cpprelheusão}, porque a intelligencia forma á idea de uma coisa, quando a \emph{percebe} ou \emph{apprehende}. --- Notamos que. rigorosamente fallando, a \emph{percepção} É O0 acto \emph{subjectivo} do \emph{cognoscente}, \emph{ao} passo que a \emph{idem é o producto} d' esse acto: ou, por outras\pgn{17}palavras, á \emph{percepção} é o conbecimento considerado \emph{subjectivamente}, e a. \emph{idea}. € \emph{o} proprio conhecimento considerado \emph{objectivanente}, Todavia, no uso commum, a \emph{percepção} e a \emph{idea} empresam-se para significar uma e a mesma coisa.\par Sendo uma representação \emph{ínfellectiral}, \emph{a idea} não deve confundir-se com ü \emph{imagem sensivel}, que se forma na \emph{phantusia} Juntamente com a percepção e que ncompanba todas ás operações intelleclaaes ; porque essa \emph{imagen sensivel} represenka uma. coisa \emph{material}, como existe na reslidade, ao passo que a \emph{deo} represente a coisa de um modo \emph{espiritual} ou \emph{Gmnaterial}.}.

%%>>> t18

\pg{17} \tese{18}{Objecto da idea.}
O objecto, que a \textbf{idea} representa, é a \emph{essencia}, ou \emph{natureza} do ser percebido. Por quanto, a \emph{coisa}, que é representada ou reproduzida na intelligencia pela idea, não é \emph{Lodo} o ser percebido, mas é apenas o que \textbf{esse} ser tem de mais intimo e que o consliLue numa determinada especie, e não noutra: assim a coisa, que a idea de \emph{homem} representa, não é tudo o que o individuo humano possue, mas é só o que o torna individuo d'esta especie, que é a \emph{humanidade}. Ora a Coisa, que um ser tem de mais intimo e que o constitue numa certa especie, é a \emph{essencia}, ou \emph{natureza}. Logo o objecto, que a idea representa, é a \emph{essencia}, ou \emph{natureza}, do ser percebido \footnote{l'odos os seres finitos são compostos de muitos e diversos elementos; alguns dos quaes são \emph{commans} a outros seres da mesma especie ou genero, e outros são \emph{proprios} de cada ser. Assim um certo homem, D. ex. \emph{Pedro}, \emph{é animal}, \emph{é racional}, \emph{é branco}, \emph{é alio}, \emph{é sabio}, ele. Alguns d' estes elementos, como \emph{animal}, \emph{racional}, são commis a todos os homens; outros, como \emph{franco}, \emph{alto}, ele,, são proprios d' este homem, que se chama \emph{Pedro}, e não de outro homem. Os elementos comunais constituem a essencia, ou a \emph{natureza} de Pedro, porque a \emph{essencia} do homem é constituida pela. \emph{animalidade} e pela \emph{racionalidade}; ao passo que os outros estão fora da natureza, . Ora a iden de \emph{lomem} representa a \emph{essencia}, ou a natureza, d' \textbf{esse} ser, c não as propriedades individuaes, que o fazem distinguir dos outros homens. Na verdade, a. \emph{idea} é a forma de uma laculdade \emph{zmematerial}, como é a nossa intelligencia. Mas uma faculdade \emph{immaterial} só pode: receber um: coisa \emph{Tmanaferiol}, como é a \emph{essencia} dos seres. Ax imagens representativas das propriedades \emph{iudioidiaes} dos seres só podem ser recebidas nas faculdades d: ordem inferior, que é a \emph{sensitiva}. : : Advertimos que a \emph{idea} pode ser tomada segundo o seu objecto \emph{material} e segundo o seu objecto \emph{formal}. É tomada segundo o seu objecto \emph{material}, quando se toma no conjuncto de \emph{todos} os elementos, que i constituem; assim, quando digo: \emph{Pedro é homem}, a idea de \emph{homen} é tomada no \emph{conjuncto} de \emph{todos} os elementos, que x constituem e que são a \emph{animalidade} e x \emph{racionalidade}, X tomada segundo o seu objecto \emph{formal}, quando se toma na \emph{parte} dos seus elementos, que denota a dilferenç \emph{especifico}; assim quando digo: \emph{o} homem e \emph{o Anjo coulecem} o \emph{Creador}, a ideo de \emph{Lomem} é tomuda \emph{apegas} "um dos seus elementos --- a \emph{racionalidade}, pela qual o homem convem com o Anjo e \emph{que} é a \emph{cuusa adequada} d esse conhecimento.\par Yol.[ 3}.

19 Caracteres da idea. --- Os principaes caracteres da \emph{idea} são dois: a \emph{universalidade} e a \emph{necessidade}.

a) A \emph{idea é universal}. --- Por quanto, a idea representa a \emph{essencia}, ou \emph{natureza} dos seres. Ora tal essencia, ou natureza, pelo facto de determinar os seres a uma certa especie, é \emph{comcm} a \emph{todos} os seres, que pertencem a mesma “especie, .e a \emph{todos} pode attribuir-se no \emph{mesmo sentido}, e por isso é uma coisa \emph{universal} ; tal é a \emph{humanidade}. Logo a idea, representando a \emph{essencia}, ou \emph{natureza} dos seres, é \emph{universal} \footnote{A idea universal, como veremos, é formada pela abstracção, emquanto a intelligencia considera nos seres de uma especie o que é commum a \emph{todos}, sem se importar do que é \emph{proprio} de cada um. E só a este caracter de \emph{universalidade}, de que são dotadas as' ideas, que devemos o poder de formar \emph{juizos} e \emph{raciocinios}. Assim, se posso dizer com razão que \emph{Pedro é homem}, \emph{Paulo é homen}, só posso dizel-o emquanto a idea de \emph{homem} representa o que ha de \emph{comuna} nos individuos humanos, prescindindo do que é \emph{proprio} de \emph{Pedro}, de \emph{Paulo}, etc. : porque se a idea de \emph{homem} representasse o que é proprio de \emph{Pedro}, não poderia dizer-se que Paulo é \emph{homem}.}..

b) A \emph{idea é necessaria}. --- 'A idea representa a \emph{essencia}, ou a \emph{natureza} dos seres. Ora a. essencia dos seres é \emph{iminulavel}, e por isso \emph{necessaria}; porque, se lhe se tirar ou accrescentar um só dos elementos, que a constituem, muda immediatamente; assim se aO \emph{homem}, que é o \emph{animal racional}, se tirar o \emph{animal}, já não lemos o \emph{homem}, mas o \emph{anjo}. Sendo necessaria a \emph{essencia} dos entes, tambem \emph{necessaria} é a idea, que a representa \footnote{Como se vé, a \emph{necessidade}, ou \&\emph{nanutabitidade}, da \emph{idea}, é uma consequencia da \emph{necessidade}, ou. \emph{imanutabilidade} da \emph{essencia} das coisas. --- Da \emph{necessidade}, ou \emph{immutabilidade} da \emph{essencia} trataremos. e com o devido desenvolvimento, mais larde, na \emph{Ontologia}.}.

%%>>> t20

\pg{18} \tese{20}{A idea é signal natural.}
\emph{Signal} é uma coisa, que nos leva ao conhecimento de outra. A \emph{natural}, ou \emph{convencional}, conforme a relação entre a coisa significativa e a coisa significada depende da \emph{natureza} (assim o. fumo é signal \emph{natural} do fogo), ou depende do \emph{accordo} dos homens (assim, o ramo da oliveira é signal \emph{convencional} da paz) --- Ora a idea é \emph{signal}, porque apresenta um objecto ao nosso espirito; e é signal natural, porque a relação entre ella e o objecto depende da \emph{natureza}, independentemente de toda a convenção bumana \footnote{A divisão do \emph{signal} em \emph{natural} e \emph{convencional} funda-se no seu nexo com a \emph{coisa sígsificada}. --- Mas, se attendermos ao seu nexo com a \emph{faculdade cogmitiva}, \emph{o signal} divide-se tambem em \emph{formal} e \emph{objectivo}. \emph{É format},\pgn{19}quando é \emph{imagem} e \emph{semelhança} de uma coisa e por isso leva \emph{immediatamente}, antes d'elle ser percebido em si mesmo, ao conhecimento da coisa Signifieuda : assim a imagem de uma. flor, recebida na faculdade visual, leva \emph{immediatamente} esta faculdade a percepção da flor. E \emph{object}: vo, quando \emph{primeiramente} é conbecido em si mesmo, e só em seguida leva á percepção do objecto representado: assim o fumo é signal \emph{objectivo} do fogo, Ora a idea é signal \emph{format}: porque leva \emph{immediatamente} á percepção do objecto, e só depois é percebida pela \emph{reflexão}. --- Mas d' isto Occupar-nos-hemos mais tarde.}.

%%>>> t21

\pg{19} \tese{21}{Comprehensão e extensão da idea.}
A \emph{comprehensao} da \emph{idea} é a somma ou o \emph{conjuncto das notas}, ou \emph{dos elementos}, que «\& \emph{constituem}; assim a \emph{comprehensão} da idea de \emph{homem} é formada pelos dois elementos, ou notas: \emph{animalidade} e \emph{racionalidade}. \emph{A extensão} é a somma ou o \emph{numero dos individuos}, \emph{que a idea abrange}; assim a \emph{extensivo} da idea de \emph{homem} é constituida pela collecção de todos os individuos humanos. --- A \emph{comprehensão} e a \emph{extensão} da \emph{idea} estão na razão \emph{inversa}; porque, quanto, maior é a \emph{comprehenszo}, tanto menor é a \emph{extensão}, --- e, viceversa, quanto maior é a \emph{extensão}, tanto menor é a \emph{comprehensao}; assim a idea de \emph{Homem} tem mais \emph{comprehensio} que a idea de \emph{animal}, porque o homem possue todos os graus de perfeição, de que é dotado o animal, e, alem d'isso, possue a razão, de que o animal é destituido, --- mas, pelo contrario, a idea de \emph{animal} tem mais \emph{extensão} que a idea de \emph{homem}, porque a idea. de \emph{animal} abrange todos os homens, como a idea de \emph{homem}, e, alem d'isso, extende-se a muitos outros seres vivos, aos quaes não se extende a idea de \emph{homem} \footnote{A \emph{comprehensao} da \emph{idea} não pode augmeptar nem diminuir, sem que 1 idea deixe de ser o que é: porque a \emph{comprehensao} denota a \emph{essencia}, e esta é incapaz de augmento e de diminuição. A \emph{extensão} da \emph{idea}, porém. pode augmentar ou diminuir, sem que por isso augmente on diminua a sux \emph{comprelensio} ; pois a \emph{extensão} é uma coisa \emph{extrinseca} e \emph{accidental a idea}.}.

Anr:GO 1]

\emph{Divisao geral da idea}

%%>>> t22

\tese{22}{Divisio geral da idea.}
A \emph{ideu} pode considerar-se --- ou quanto ao \emph{objecto}, que ella representa, --- ou quanto ao \emph{modo}, por que o representa. D'abi uma duplice divisão,

%%>>> t23

\tese{23}{Divisão da idea quanto ao objecto representado.}
Quanto ao \emph{objecto}, que representa, a \emph{idea} divide-se --- \emph{a}) em \emph{real} e \emph{logica}, --- \emph{b}) \emph{positiva} e \emph{negativa}, --- \emph{c}) \emph{absoluta} e \emph{relativa}, --- \pg{20}  \emph{d}) \emph{simples} e \emph{composta}, --- \emph{e}) \emph{abstracta} e \emph{concreta}, --- \emph{f}) \emph{intuitiva} e \emph{abstractiva}, --- \emph{q}) \emph{univoca} e \emph{analoga}. --- \emph{h}) \emph{singular}, \emph{uni}. \emph{versal}, e \emph{particular}, --- \emph{i}) \emph{universal-directa} e \emph{universal-reflexa}.

a) Idea real e logica. --- Idea \emph{real} é a que representa um objecto \emph{real}, isto é, existente fóra do nosso pensamento; tal é a idea de \emph{homem}, considerado como é em si mesmo. --- Idea \emph{logica} é a que representa um ente \emph{logico}, isto é, só existente na intelliigencia; ta] é a idea de \emph{homem}, considerado como \emph{especie}, a que estão subordinados os individuos humanos \footnote{Um e o mesmo ser pode considerar-se \emph{em} \$\emph{i}, ou na \emph{éntelligencia}. ---- Considerado \emph{em si}, o ser possue propriedades, de que a natureza a dotou. independentemente do nosso pensamento. A idea, que representa este ser, é \emph{real}; tal é a idea de \emph{homem}, considerado \emph{em si} mesmo, isto é, nos seus elementos ou notas de \emph{ser sensitivo} e \emph{racional}. --- Considerado na \emph{intelligencia}, : o ser possue ou pode possuir attributos, que não possue na realidade c por isso lhe foram dados, não pela natureza, mas pela abstracção e reflexão do nosso pensamento. A idea, que representa este ser, é \emph{logica}; tal é a idea de \emph{homem}, considerado: como \emph{especie}; porque na realidade não existe o \emph{Jo mem-especie}, mas existe unicamente o \emph{homem-individuo}. .}.

b) Idea positiva e negativa. --- A idea \emph{posilica} é a que exprime uma \emph{entidade} ou. \emph{perfeição}; taes são as ideas de \emph{luz}, de \emph{virtude}. --- \textasciicircum{}A idea \emph{negative} é a que denota a \emph{ausencia} ou \emph{privação} de \emph{entidade} ou \emph{perfeição}: (aos são as ideas de \emph{trevas}, de \emph{peccado}, que indicam a \emph{ausencia} da luz, ou da virtude \footnote{A idea \emph{negativa} não representa a propria \emph{zwsencia} de entidade, porque a. ausencia é \emph{nada} e do: nada. não ha imagem, mas significa a \emph{opposta entidade emquanto negada}, Assim a ideu de \emph{trevas} representa --- não a ausencia da \emph{luz}, mas a negação da luz, isto é, a propria luz \emph{emquamio negada}. --- Para que. pois, uma idea possa chamar-se \emph{positca} ou negativa, é necessario attender a \emph{coisa significada}, e não ao seu \emph{termo} orat, Se a idea exprime uma. ausencia de entidade, é sempre \emph{negativa}, embora o termo oral seja \emph{positivo}, tal é a idea de \emph{morte}; como tambem. .se \& idea. exprime uma \emph{entidade}, é sempre \emph{positiva}, embora o termo oral \emph{seja negativo}, tal é a idea de \emph{imanort«l}. Por isso alguns termos são \emph{positivos} na expressão e \emph{negativos} na \emph{significação} (ex. \emph{cegueira}, \emph{morte}, \emph{peccado}) e outros são \emph{positivos} na. \emph{significação} e \emph{negativos} na \emph{expressão} (ex. \emph{immortalidade}, \emph{inmensidade}).}.

€) Idea absoluta e relativa. --- A idea \emph{absoluta} representa um objecto, que se não refere a outro; como a idea de \emph{planta}. ---.A \emph{relativa} representa um objecto, que se refere a outro; como a idea de \emph{filho}, que se refere á de \emph{pae}.

d) Idea simples e composta. --- A idea é \emph{simples}, quando é constituida por um \emph{unico} elemento; tal a idea de \emph{ente}. --- E \pg{21}  \emph{composta}, quando abrange \emph{varios} elementos; como é a idea de \emph{homem}, composto de alma e de \emph{corpo} \footnote{Pode dizer-se tambem que a. idea --- é \emph{siptples}, quando exprime uma essencia \emph{incomplexa}, isto é, quando exprime uma nica essencia (porque a essencia, ainda que constituida por varios elementos, é por si \emph{simples} ou \emph{indivisivel}, de modo que não lhe se pode tirar um só elemento, sem que ella fique destruída), tal é a idea de \emph{homem}; --- é \emph{composta}, quando exprime uma essencia \emph{complexa}, isto é, um complexo de essencias, tal é a idea de \emph{sabio}, que significa o \emph{homem} e a \emph{sabedoria}.}. . \textbar{}

€) Idea abstracta e concreta. --- Idea \emph{abstracta} é a que representa uma essencia ou propriedade, como que separada do sujeito e subsistente por si; tal é a idea de \emph{sabedoria}. --- \emph{Concreta} é. a que representa um sujeito com a sua fórma ou propriedade, como existe nas coisas; tal é a idea de \emph{sabio} \footnote{Ha Leis especies de \emph{concreto}: \emph{metaphiysico}, \emph{physico} e \emph{logico}. É \emph{metaplajsico}, quando a lorma ou propriedade é \emph{essencia}! ao sujeito : tal é homem, porque a \emph{Ammanidade} compete \emph{essencialmente} ao individuo da especie humana. É \emph{physico}, quando a propriedade, ainda que \emph{intrinseca} ao sujeito, € comtudo \emph{accidental}, emquanto não lhe compete necessariamente ; tal é \emph{sabio}, porque a \emph{sabedoria} não compete \emph{necessariamente} ao homem, aliás todo \emph{o} homem seria sabio. É \emph{logico}, quando a propriedade é \emph{accidental} e \emph{extrinseca} no sujeito: tal é \emph{Zewvado}, porque o \emph{louvor} está propriamente no conveilo e nas palavras dos que louvam, e não no sujeito, que é seu termo.\par G) A idea \emph{intuitiva}, se chega a representar ou exprimir a coisa tanto, quanto, esta. €- cognoscivel, de modo que \emph{o} conhecimento \emph{eguale} e como: que \emph{exqotte toda} a grandeza do objecto conhecido, chama-se \emph{comprehensiva}. Os bemaventurados, no ceu, têm um conhecimento \emph{intuitivo} de Deus, mas não \emph{comprehensivo}; Só uma intelligencia. infinito pode \emph{comprehender} um ser \emph{ivfinito}, --- Façamos uma advertencia, Para a iden de uma coisa ser \emph{iniuifiva}, não basta que a coisa seja conhecida \emph{immediatamente}, \emph{por si mesina}, mas e \emph{preciso} que se una a faculdade cognitiva ou pela sua essencia, ou por uma \emph{especie} ou imagem propria. Assim, o \emph{Zodo} e a \emph{parte} são coisas conhecidas \emph{inmediatamente}, \emph{por si mesmas}, sem que, para serem percebidas, seja preciso recorrer \emph{n} outras coisas, em que estejam contidas, e comtudo essas ideas não são \emph{intuitivas}, mas \emph{absiractioas}, porque se formam pela \emph{abstracção}.}.

f) Hdea intuitiva e abstractiva. --- A idea \emph{intuitiva} é a que representa uma coisa como é em si mesma, emquanto esta se une immediatamente á intelligencia ou pela sua essencia ou por uma imagem propria; tal é a idea de \emph{Pedro}, que está na minha presença. --- A idea \emph{abstractiva} é a que representa uma coisa, não como é em si mesma, mas como se encontra na imagem de outra coisa, que de algum modo a manifesta; tal é a idea, que temos agora de Deus pela vista das creaturas (7)

\$) \textbf{Idea univoca} e \textbf{analoga}. --- A idea \emph{univoca} representa um objecto, que se toma Sempre \emph{numa} e \emph{mesma} significação; tal é a idea de \emph{substancia}, que applicada a planta, ao sol, etc., significa sempre a mesma coisa. --- A idea \emph{analoga} representa um objecto, que se atiribue aos inferiores numa signilicação \emph{nem totalmente identica}, \emph{nem tolalimente diversa}; tal é a idea de \emph{sanidade}, que se attribue ào \emph{animal}, ao \emph{alimento}, \emph{a côr} \footnote{\pgn{23}S. '"homaz, a proporção ou a semelhança entre duas proporções: « Aequalitas proportionum vocalur \emph{proporcionalitas} » (Yn V \emph{Met}., lect. V).\par A idea \emph{analoga} segundo a \emph{proporcionalidade} pode ser analoga segundo a proporcionalidade \emph{propria} e segundo a proporceionalidade \emph{émpropria}, ou \emph{metaphorica}. --- É analoga segundo a proporcionalidade \emph{propria}, quando a propriedade, representada pela idea, se attribue aos dois termos analogados, porque num e noutro se encontra como \emph{forma intrinseca}, embora se não encontre em ambos \emph{do mese modo}. Ássim a idea de \emph{sapiencia}, quando se attribue a Deus e ao bomem, é analoga segundo a \emph{proporctonalidade}, porque a relação da \emph{sapiencia} com a intelligencia humana é \emph{semelhante} á relação da mesma sapiencia com a intelligencia divina, visto que numa e noutra inteligencia a sapiencie significa \emph{o} conhecimento das causas supremas; --- e é nnaloga segundo a proporcionalidade \emph{propria}. porque a \emph{sapiencia} é uma forma \emph{zn£rinseca} tanto da intelligencia divina quanto da bumana, embora na divina seja infinitamente mais perfeita do que na humana: --- É analoga segundo a proporcionalidade \emph{improprie}, ou \emph{metaphorica}, quando a propriedade, representada pela ides. se attribue aos dois termos analogados, porque num se encontra como forma \emph{in£rinseca} e noutro como semelhança d'essa forma. Assim, a idea de \emph{riso}, quando se attribue 20 homem e ao jardim, --- é analoga segundo a \emph{proporcionalidade}, porque a relação do viso com a face do homem é semelhante a relação da amenidade com o jardim: --- mas é analoga segundo a proporcionalidade \emph{impropria}, ou \emph{metaphorica}, porque o \emph{riso} compete: \emph{propria} e \emph{intrinsecamente} ao homem, ao passo que ao jardim só compete por uma certa \emph{sestelhanca}, que este tem com o proprio riso do homem. --- Esta analogia de proporcionalidade \emph{Impropria}, embora tenha o seu fundamento nas coisas, comtudo depende principalmente do arbitrio humano. --- D'onde se segue:\par a) Na analogia de \emph{proporcio}, ou de \emph{attribuiczo}, --- 1º) n forma é \emph{intrinseca} ao analogado \emph{principal}, mas é \emph{extrinseca}. e Só convem por \emph{denominação} (\emph{quanto} ao nome) aos outros analogados; assim a \emph{sanidade é intrinseca} ao animal, mas é \emph{extrimsecm} a cór, ao alimento, ete., que d' ella só tem o \emph{nome}, não a \emph{realidade}; --- 9") a forma do analogado \emph{principal} deve entrar na definição dos outros analogados, que recebem do principal o nome; assim não podemos entender o que é ;\& sanidade no alimento, se não entendermos primeiramente o que é a sanidade no animal; --- 3º) se nos analogados, \emph{principal} e secundarios, não ha unidade de \emph{forma}, porque 2 forma, enconlrando-se apemas no analorado principal, é numericamente una, ba comtudo \emph{unidade} de \emph{ordem}, emquanto os analogados secundarios se referem a uma e a mesma forma do anatogado principal, da qual tiram o seu nome; --- 4º) os \emph{analogados} não têm um e o mesmo conceito \emph{objectivo}, porque não têm uma \textbar{} a mesma forma, e o conceito é proporcionado \& forma.\par b) Na analogia de \emph{proporcionatidade propria}, --- 1º) a forma é \emph{intrinseca tanto} ao \emph{anaüle;nio principal} Como nos \emph{outros}, “\emph{ainda} que o não seja \emph{do}\par \emph{nwsmo medo}: e por isso us ideas \emph{analogas segundo} esta analogia differem\pgn{24}das analogas segundo a \emph{proporção} e a \emph{adtriluiécao}, c tambem das \emph{neivocas}; --- 8) o abalogado \emph{principal} não deve entrar na definição dos outros, porque estes têm a sua forma \emph{intrinseca}, de que liram o proprio nome: -- 39) todos estes analogados convém num só conceito objectivo, porque todos são intrinsecamente constituidos pela mesma forma; mas esta unidade é imperfeita, ' porque a forma não se encontra em Lodos do \emph{mesmo modo}.\par €) Na analogia. de \emph{proporcionatidade metaphorica} verifica-se o que se verifica na analogia de \emph{proporção} e de \emph{at£ribuicGo}; porque, p. ex., o \emph{viso} não é \emph{intrinseco} no jardim, como o é no homem, e por isso não tem e não pode ter um e o mesmo conceito objectivo, \textbar{}}.

( A idea \emph{analoga} merece especial attenção. \emph{Analogia} significa \emph{proporção}. \emph{velação}, \emph{semelhança}. Tdea \emph{enaloga}, como dizemos no texto, significa uma propriedade ou perfeição, a qual se altribue a entes, que são \emph{essenciclmente} diversos. mas que têm entre si uma certa \emph{Proporção} ou semelhança. ssim, quando se diz são o \emph{animal} e são o \emph{alimento}, a. \emph{sanidade} atiribue-se a0 animal e a0 alimerto numa significação, que não é --- nem totalmente \emph{identica}. porque a \emph{sanidade} ateibuida ao «\emph{alimento} significa uma coisa, e altribuida ao \emph{animal} significa outra, --- nem tot:lmente \emph{diversa}, porque a \emph{sanidade} encontra-se no \emph{animal} como no seu \emph{sujeito} e no \emph{alimento} como na sua propria cansa.

Ha uma triplice especie de ideas \emph{onalogas}; uma segundo a \emph{proporção}, outra segundo a \emph{af£ribwiczo}, e a terceira segundo a \emph{proporcionalidadqe}. --- Alguns auctores identificam a analogia de \emph{proporção} com a de \emph{atiribuição}.

3) A idea \emph{analoga}, segundo : \emph{proporção} representa uma propriedade, a. qual se attribue a dois.entes, --- a um. emquanto se encontra nelle \emph{principal} e \emph{intrinsecamente} (\emph{principal analogado}), --- a. outro, emquanto, .ainda que se não encontre nelle intrinsecamente, comtudo tem .uma certa relação com \emph{o} primeiro (\emph{analogado secundario}); tal é a idea de \emph{sanidade}, attribuida ao \emph{animal} e ao \emph{alimento}, porque o \emph{animal} diz-se \emph{são} por si, emquanto a. \emph{sanidare lhe} convem \emph{intrinsecamente}, a0 passo que o \emph{alimento} se diz são só em ordem ao animal, isto é, só emquanto é causa de sanidade n0 animal.

b) A idea \emph{analoga} segundo a \emph{attribuição} representa uma. propriedade, a qual se attribue a dois entes, --- não emquanto a um convem \emph{intrinsecamente} e ao outro em ordem ao primeiro, --- mas emquanto um e outro se referem \emph{a} um terceiro, ao qual essa propriedade convem \emph{principal} e \emph{intrinsecamente}; tal é a idea de \emph{sanidade}, attribuida ao \emph{alimento} e a \emph{cór}, porque o \emph{alimento} e a \emph{cór} não se dizem \emph{sãos} senão em relação \& um terceiro, --- ao animal, em que a sanidade se encontra principal e essencialmente, visto que o \emph{alimento} E causa da sanidade do animal. e a côr é \emph{signal} d' essa sanidade, --- Por isso, na. analogia de \emph{attribuíção} ha. pelo menos, \emph{tres} termos, 80 passo que nu analogia de \emph{proporção} ha só \emph{dois}. .

€) A idea \emph{analoga} segundo a \emph{proporcionatidade} representa. uma propriedade, que se altribue a dois entes, não porque um se refere a outro, mas Porque \emph{a reacao} d'essa propriedade com um ente é \emph{semelhante} ou \emph{proporcionada á relação} da mesma propriedade com outro ente: tal é a idea de \emph{rei}. que se \emph{attribue} ao homem e no leão, porque a relação do leão com os outros animaes é semelhante àá relação do soberano com os seus subditos, isto €. o leão está para os outros animaes como Oo soberano está para os seus subditos. ---- Como se vé, a analogia de \emph{preporcionalidade} encerra duas proporções e \emph{quatro} termos (embora nem sempre explicilos), e chama-se de

h) Idea \textbf{singular}, universal e \textbf{particular}. --- A idea é \emph{singular}, quando representa uma \emph{essencia}, \emph{propria} de um só \emph{individuo determinado}, de modo que só a este \textbf{convem} ou pode \emph{pr oporcionalidade}, por \textbf{isso} mesmo que a \emph{proporcionalidade} é, como diz \pg{24}  ser xitribuida; tal é a idea de \emph{Pedro}, --- É \emph{universal}, quando representa uma \emph{essencia}, comunwm a \textbf{muitos} entes, de modo que a todos e a cada um d'elles pode ser attribuida; tal é a idea de \emph{humanidade}, que é commun a todos os homens, e que se attribue a todos e a cada um d'elles. Por isso a idea \emph{universal} toma-se em toda a sua \emph{extensão} --- A \emph{particular} é a propria idea \emph{universal}, nias tomada. numa \emph{parie} da sua extensão; como é a idea de \emph{alguns homens} \footnote{Toda a idea representa uma \emph{essencia}. Mas a essencia, que a idea \emph{sinfWlar} representa, está cercada de certas propriedades, que no seu conjuncto só podem convir a um \emph{individuo determinado}. Vslas propriedades chamam-se tambem \emph{notas} ou \emph{propriedades individuaes}, e reduzem-se a \emph{sete}, contidas nos seguintes versos latinos:\par \emph{Forma}, \emph{figura}, \emph{locus}, \emph{tem pus}, \emph{stérps}, \emph{patria}, \emph{nomen}:\par Haec ea sunt septem, quae non babet unus cet alter,\par A essencia, porém, que n iden \emph{nmiieersal} exprime, considera-se em si ou nos seus elementos constitutivos. sem relação anos caracteres, que a tornam concreta nos individuos, é d' este modo represenLa se-nos como uma coisa, á qual'por tal modo existe ou pode existir em muitos Seres, que nelles se multiplique e a enda. um d\textasciicircum{} elles se attribue, Assim (insistimos no mesmo exemplo) a \emph{hwmanidace}, abstrahida das cireumstancias ou notas que a tornam concreta em \emph{Pedro}, \emph{Paulo}, ectc., --- é uma. coisa, que se encontra em muitos individuos, pois a humanidade existe em todos os homens, --- multiplica-se \emph{numericamente} nelles, porque !u, tantos homens, quantos são os que têm a humanidade, --- e attribue-se a. cada um d'elles, pois de cadi individuo de natureza. humans podemos dizer que é \emph{komen}.\par A idea \emph{universal} dilfere --- \emph{a}) da \emph{transcendental}, --- 0) du \emph{commun}, --- €) da \emph{collectiva}, ---- 4) e da \emph{particular}. T\par a) A ideo \emph{uamiversal} differe da \emph{transcendental}, X idea é \emph{Ére uscendenfat}, quando representa uma coisa, que node attribuit-se a lodos os eittes, existentes e possiveis: tal é a idea de \emph{ente}, de \emph{verdadeiro}, dc \emph{bom}, ele. Diz-se \emph{transcendental}, porque \emph{transcende}, ou \emph{excede}, lodos. os \emph{generos} e todas as \emph{especies}. Ora ha uma grande dilferenca. entre a idea \emph{universal} e a \emph{Drenscendental}. Por quanto. a iden \emph{nnuicersal} não representa senão o que é proprio de uma \emph{deterininada} categoria de entes: ao passo que a \emph{transcendental} applica-se a \emph{todos} os seres, existentes e possiveis, Alem d'isso, a idea \emph{universal}\pgn{25}representa uma coisa, que se attribue a todos os sujeitos, aos quaes se extende, numa significação wnivcora ou \emph{édenfic}« (assim \emph{homem}, applicado a \emph{Pedro}, \emph{a Paulo}, etc.. significa sempre a mesma coisa); ao passo que a \emph{irünscendental} representa uma coisa, que se atiribue aos seus sujeitos numa significação :\emph{analoga} (assim \emph{ente}, applicado ao \emph{homem}, a planta, a \emph{pedra}, não significa a mesma coisit).\par \emph{b}) A idea \emph{universal} dilfere da commam, Na verdade, a essencia, indila pela idea \emph{universal}, multiplira-se nos sujeitos, .a que se attribue, de modo que a essencia de um individuo, embora convenha \emph{especificamente} com a essencia de outro individuo, todavia dilfere \emph{weemericaimente}; assim são lantag as humanidades, quantos os homens. Mas a essencia, denotada pela idea \emph{commun}, embora se attribua a varios suppostos, todavia não exprime exigencia de \emph{meltiplicação} numerica nos mesmos suppostos, a que se attribue. Assim a Essencia Divina é commum ás Tres Pessoas, mas não pode chamar-se \emph{universal}, porque não pode estar sujeita á multiplicação \emph{unierica}, pois ha um só Deus, (Cf. S, Th. in ] \emph{Sen}, Dist. 19. q. 4, a. 29 ad 2. ---\par €) A idea \emph{universal}! diflere da \emph{collectiva}, (exercito, rebanho). Porquanto, a \emph{nnieersal} pode ser attribuida a \emph{todos} e a \emph{ca/a} um dos sujeitos; mas n \emph{collectiva} convem só a \emph{collecção inteira} dos sujeitos. Por exemplo, a. idea du \emph{planta}, que é \emph{universal}, altribue-se a \emph{cada} um. dos vegetaes, assim dizemos «a figueira é planta, o pinheiro é planta, a flor é planta», so passo que a idea de \emph{exercito}, por ser \emph{collectica}, não convem a. enda soldado, porque não podemos dizer «este soldado é o exercito» (\emph{Sm}. \emph{Th}.. p. 1, q. 31, a. 1 ad 2).\par \emph{d}) A idea \emph{suniversal} dilfere da. \emph{particular}. Por quanto, como dissemos, i idea \emph{sudversal} bLoma-se em \emph{toda} a sun \emph{exfewsão}, mas n \emph{particular} só numa \emph{parte} da sua \emph{extensão}; assim quando dizemos « \emph{alguns homens}» concebemas a essencia humana existente, não em todos os homens, mas só em alguns homens \emph{fudetermáinados}. --- Como se vê, a idea \emph{particular} approxima-se da. \emph{singular}, e ás vezes chega a identificar-se com esta, como neste exemplo : \emph{Certo homen}: escreveu \emph{os Lusiadas}.\par Portanto, é facil reconhecer as notas, que carnaelerisam e distinguem a ddea \emph{singular} é a \emph{universal}. --- A \emph{séngpilar} € \emph{éucommaunicaocl}. porque convem exclusivamente a um certo e deleriminado individuo: a \emph{universal é comnutricavel}, porque pode attribuir-xe e attribue-se a. todos e a cada um dos individuos, dotados da essencia, que essa idea exprime. --- A \emph{singular é concreta}, porque está envolvida nas propriedades \emph{zndieiditaes}, a0 passo que i \emph{miilversol é abstracta}, pois considera-se como que separada das propriedades \emph{zud£oédiaces}, que possue na realidade. --- A \emph{singular} é \emph{compleae}, pois que se compõe \emph{da essencia} e dos \emph{caracteres individaaes}: \emph{a wnisersal}. 6 \emph{iucomyles}, porque só exprime a \emph{essencia} do ente. --- A \emph{singular-tem} menos \emph{extensão} que a \emph{universal}, porque a \emph{singular} applica-se a um só \emph{individuo}, ao passo que a \emph{anicersal} appliea-se a \emph{meaifos}, --- mas tem mais \emph{compreensão} que a \emph{nnicersal}, porque a \emph{singular} não só representa a \emph{essencia}, como a \emph{iusicersal}, mas tambem representa ás \emph{aetas} ou \emph{propriedades}, que a \emph{iiioersal} não representa.}. \textbar{}

i) \textbf{Idea universal-directa} e \textbf{uaiversal-reflexa}. --- A idea \emph{universal-directa} representa uma essencia, considerada nos seus \emph{elementos constitutivos}, \emph{e sem relação aos individuos}, \emph{nos}

\emph{quaes se encontra ou pode encontrar-se}; tal é a idea de \emph{homem},

só considerada nas notas de «\emph{nimalidade} e de \emph{racionalidade}. --- “A idea \emph{universal-reflexa} representa a. \textbf{mesma} \emph{essencia}, \emph{mas considerada} em \emph{relaczo} aos \emph{individuos}. ou aos \emph{inferiores}, \emph{nos quaes pode existir} e \emph{aos quaes pode attribuir-se}; tal é a idea de \emph{homem}, considerada como \emph{especie}. --- Como se vé, a idea \emph{universal-refleza} é a propiria \emph{universal-directa}, a qual se accrescenta a relação aos inferiores O). E

%%>>> t24

\pg{22} \pg{23} \pg{25} \pg{26} \tese{24}{Divisio da idea quanto ao modo de representar o objecto.}
Em relação ao \emph{modo}, por gue representa o objecto, ü \emph{idea} divide-se --- \emph{a}) em \emph{clara} e \emph{obscura}, --- \emph{b}) \emph{distincta} e \emph{confusa}, --- \emph{oc}) \emph{completa} e \emph{incompleta} --- \emph{d}) \emph{adequada} e \emph{inadequada}.

4) Idea clara e obscura. --- A idea é \emph{clara}, quando o objecto, por ella representado, se exirema de qualquer outro. --- É \emph{obscura}, quando o objecto não se \textbf{extrema} dos outros. --- A idea \emph{obscura} pode tornar-se \emph{clara}.

; Db) Idea distineta e \textbf{confusa}. --- A idea é \emph{distincta}, quando não só extrema o seu objecto de qualquer outro, mas tambem indica os \emph{elementos essenciaes} do mesmo objecto; assim a idea de \emph{homem é distincta}, quando representa a \emph{animalidade} e a \emph{racionalidade}, que são os elementos constitutivos do homem. --- É \emph{confusa}, quando não exprime os \emph{elementos essenciaes} do objecto, ainda que o extreme de qualquer outro. aa\footnote{A idea \emph{universal-directa} representa um ente \emph{real}; porque os \emph{constitutivos} de uma essencia se acham.verdadeiramente na mesma essencia, independentemente da nossa intellizencia. A \emph{untiversal-reflexa}, --- emquanto contem a. \emph{universal directa}, is!o é, uma. essencia, é ente \emph{real}, --- mas, emquanto diz \emph{relação} nos inferiores.e se considera como uma \emph{forma} commun a varios inferiores, é ente \emph{logico}, isto é, é um producto da nossa \emph{intelligencia e} não existe nº reslidade, embora tenha nesta o seu fundamento: assim 0 \emph{homem}, considerado como \emph{especie}, não existe realmente, Mas d'isto mais \emph{tarde}, e com o devido desenvolvimento. ---, A idea \emph{universal-directa} é chamada \emph{directa}, porque se concebe por uma consideração \emph{directa} da nossa \emph{intellizencia}. Diz-se tambem --- \emph{metaphysica}. ou \emph{objectiva}, porque exprime a essencia considerada em \emph{si}, isto é. independentemente do nosso pensamento. --- \emph{fundamental}, porque, não designando \emph{explicita} e \emph{positivamente} os individuos, é o \emph{fundamento} da. \emph{unicersal-reflexão}. --- A \emph{mueiversal-reflexa} diz-se \emph{reflexa porque} se concebe por uma consideração \emph{reflexa} da inlellizencia. Chama-se \emph{Lambem} --- \emph{logica}, porque a \emph{relação} de uma essencia nos individuos, \emph{em que} se encontra ou posst encontrar-se, é obra inteiramente dependente \emph{da intelligencia}, --- \emph{formal}, porque exprime exactamente a indole da idea \emph{iueersul}, \emph{que} consiste em uma coisa, ser \emph{counmsar} a mulus.}

c) Idea \textbf{completa} e incompleta. --- A idea é \emph{completa}, quando abrange, não só os elementos \emph{essenciaes} do objecto, mas ainda as \emph{mais simples notas}, em que elles podem resolver-se; \textbf{assim} \textless{} a idea de \emph{homem} é \emph{completa}, quando exprime tambem as notas de \emph{ente}, \emph{subsistente}, \emph{corporeo}, \emph{vivente}, \emph{sensitivo}, \emph{intelligente}, \emph{livre}, --- incluídas na \emph{enimalidade} e na \emph{racionalidade}. --- E \emph{incompleta}, quando, embora exprima os elementos essenciaes, de que O \textbf{ente} é constituido, todavia não indica as mais \emph{simples notas}, em que esses elementos se resolvem.

d) Idea adequada e inadequada. --- Idea \emph{adequada} é a que apresenta i intelligencia todo o objecto em si mesmo e em Lodas as suas relações. por forma que entre elle e a idea existe uma completa \emph{equação}. Esta idea, superior ás forças da nossa inteligencia, é só propria de Deus. --- A \emph{inadequada} É a que não chega a tanta perfeição (').

AmwriGo 111 2 \textbar{} ae

\emph{Divisao da idea universal-reflexa}

Divisão da idea universal-reflexa. --- A idea \emph{universal reflexa} é, como dissemos, a que representa a \emph{essencia em relação} aos \emph{individuos}, e exprime o \emph{modo}, \emph{por que uma coisa pode existir em muitos sujeitos} e \emph{ser-lhes} «\emph{tlribuida na mesma significação}. --- Ora esta idea divide-se em \emph{especie}, \emph{genero}, \emph{differença}, \emph{propriedade} e \emph{accidente}. Estas ideas chamam-se \emph{predicaveis} \footnote{Para outros, 2 idea é \emph{adequada}, se do seu objecto representa O que é bastante para um cero e determinado fim: é \emph{inadeguada}, no caso contrario --- Todavia esta interpretação não está muito conforme com a propria \emph{elymologia} da palavra «\emph{adequada} ».}.\footnote{A idea \emph{universal} pode ser \emph{univoca} ou \emph{analoga}. --- A «\emph{eniversal-anafoga} é a que se attribue a diversos seres em significação \emph{analoga}; assim são \emph{analogas} ás ideas \emph{transcendentaes}, --- \emph{ente}, \emph{uno}, \emph{verdadeiro}, bom, --- \emph{porque a entidade}, \emph{a unidade}, \emph{a verdade}, 0 \emph{bem}, ainda que sejam possuidos por todos os seres, todavia não o são do mesino medo. --- A «\emph{nzcersal-unicoer E} a que se attribue a; diversos seres \emph{em} significação \emph{uivoca}; assim \emph{são muúcocas ás cinco ideas universaes-reflezas}. enumeradas no texto, pois altribuem-se aos sujeitos na mesma significação.\par Ás cinco ideas \emph{universaes-univoçcas} eram chamadas pelos antigos « \emph{praetica reis} » (\emph{praedivabilia}), pois exprimem, como dissemos, os «\emph{diversos nodos por gue iat coisa pode existir em muitos sujeitos e ser-Hies atribuida} (\emph{praedicarii}, e pelos modernos são chamadas « \emph{cafegoremas} v.}

%%>>> t26

\pg{27} \pg{28} \tese{26}{Especie.}
\emph{Especie} é uma idea \emph{universal}, que repre-

senta a essencia \emph{completa} de muitos \emph{individuos}, \emph{só mümericamente} distinctos entre si; tal é a idea de \emph{humanidade}, que, comprehendendó a \emph{animalidade} e a \emph{racionalidade}, exprime a essencia \emph{completa} de cada individuo humano. --- Para que, pois, uma idea (\emph{o conteudo}) possa attribuir-se como \emph{especie}, é necessario --- que denote um coisa \emph{commun} a varios entes, --- que essa coisa se refira a entes .só \emph{numericamente} distinctos entre si, --- € que exprima a essencia \emph{completa} dos mesmos entes \footnote{Um exemplo mostrayà claramente. o modo por que a nossa intelligencia forma a idea de \emph{especie}. Vendo m vitos individuos humanos, podemos, por meio da \emph{absfracczo}. prescindir das notas ou caracteres \emph{vndéviduaes}, \emph{proprios} de cada um, e só considerar os elementos \emph{constitutivos} da \emph{essencia}, \emph{communs} a todos os que sio dotados da MESA essencia. O conjunceto \emph{de todos} OS elementos \emph{commans} constitue a \emph{especie}. Assim a idea de \emph{homem}, considerada em relação aos individuos humanos, é \emph{especie}; porque --- denota \emph{zona}. Coisa cominum a varios entes, pois a \emph{humanidade} é uma coisa só e todavia se encontra multiplicada em todos os individuos humanos, --- diz respeito a entes só \emph{monericamente} distinctos, pois os individuos humanos \emph{differem}, não \emph{essencial}, mas \emph{numericamente}, --- exprime a. essencia \emph{completa} d'esses individuos, porque não se encontra nelles elemento substancial ou commum, que não se ache contido na iden de \emph{homem}.}.

%%>>> t27

\tese{27}{Genero.}
\emph{Genero} é uma idea. \emph{universal}, que representa o elemento essencial, \emph{commuum} a \emph{especies diversas}: assim \emph{animal é genero}, porque exprime o elemento \emph{comium} a essencia do homem e a do animal irracional. --- Para que, pois, uma idea possa attribuir-se como \emph{genero}, é necessario --- que denote \emph{uma} coisa commum a vários entes, --- que essa coisa se vefirá a entes \emph{especificamente} distinctos entre si, --- e qué exprima a essencia \emph{incompleta} dos mesmos entes, e, propriamente, o elemento \emph{determinavel} da essencia (\emph{O}.\footnote{Se compararmos um homem, não com outro homem, mas, por ex. com um cavallo, veremos que um e outro convém num elemento \emph{essencicl}, mas differem noutro. Convem emquanto são ambos \emph{animaes}: pois dizemos que o homem é \emph{animal}, o envallo é \emph{antinal}, Differem emquanto o homem é animal \emph{racional}, e o cavallo é animal \emph{irracional}. Esse elemento \emph{essencial}, \emph{que} é commum a diversas \emph{especies}, é 0 \emph{genero}, Assim \emph{anbnal é} "\emph{enero}, porque --- denota \emph{uma} coisa \emph{conamur} a varios entes, pois a \emph{unimalidade}, que é uma só coisa, é commum ao homem e ao \emph{bruto}, --- diz respeito \& entes, \emph{que} diferem na \emph{especie}, pois o homem e o bruto são entes \emph{especificamente diversos}, --- exprime a essencia \emph{incompleta}, porque o homem e o bruto, \emph{além} d'esse elemento essencial, que é commum a ambos e que é significado \emph{pela} idea de \emph{animal}, possuem outro elemento essencial, que é proprio de\pgn{29}cada um e de que deriva a sua diversidade especifica, --- e por isso exprime \emph{o} elemento \emph{determinavel} da essencia, porque, por ex.. \emph{anemia}, podendo ser \emph{racional}, ou \emph{irracional}, não é por si determinado a uma od outra especie, e por isso é \emph{determinavel}, e \emph{deve} ser \emph{determinado}. : :\par O \emph{genero} diflere da \emph{especie}; porque a \emph{especie} exprime a essencia \emph{completa}, e existe em sujeitos, SÓ \emph{nmenericamente} distinctos entre si; \&o passo que o \emph{genero} exprime uma \emph{parte} da essencia e existe em sujeitos, que \emph{especificamente} dilferem um do outro. Por isso a \emph{comprehensaão} do \emph{genero} é menor que ES da \emph{especie}, embora a \emph{e.ctensão} d'aquelle seja major. que a dº esta, pois o \emph{genero} comprehende varias \emph{especies}. . . \emph{EE}}

%%>>> t28

\pg{29} \tese{28}{Differenca.}
\emph{Differença} é.uma idea universal, representativa do elemento essencial, que é \emph{proprio} de \textbf{cada} especie, ---- istb é, que unido ao \emph{genero} constitue a \emph{especie}; tal é racional, que unido ao genero \emph{animal}. constitue a especie humana. Por isso, para que uma idea possa altribuir-se como \emph{difference}, deve --- representar uma coisa commum a varios entes, \emph{mumericauente} distinctos, --- exprimir a essencia \emph{incompleta}, e, propriamente, o elemento \emph{determinante} da essencia (\textasciicircum{}).\footnote{Assim, quando dizemos: o \emph{Jromem é animal} racional, 2 idea \emph{racional} \& uma \emph{differenca}, porque --- representa uma coisa que se encontra em £\emph{odos} os individuos humanos, só \emph{munericamente} distinctos, porque \emph{todo} o homem é \emph{racional}, --- exprime a essencia \emph{incompleta}, porque o homen, além de \emph{racional}, é tunbem \emph{animal}, --- € exprime o elemento \emph{determinante} da essencia, porque \emph{determina} O genero \emph{animal} e o limita a uma especie (humana) e não a outra (belluina).\par As seguintes notas, relativas á \emph{diferença}, são de grande importancia.\par 2) A \emph{differença} distingue-se da \emph{especie} e do \emph{genero}. --- Distingue-se da \emph{especie}; pois a \emph{especie} exprime a essencia \emph{completa}, e a \emph{differenga}- exprime so uma \emph{parte} da essencia, --- Distingue-se do \emph{genero}; pois .este exprime i pirle essencial comuniun a especies diversas, e, para exprimir isso, assume a forma de \emph{substantivo}, (por ex. «animal»), ao passo que a \emph{differenga}, denota a parte essencial \emph{propria} de cada, especie, e, para exprimir j850, assume a forma de \emph{adjectivo} (por ex: «\emph{racional} »). A razão é que a \emph{differença}, embora em si seja mais nobre que o genero (assim \emph{racional} é mais nobre que \emph{animal}), todavia, para o nosso modo de entender, é inferior ao \emph{genero}; pois este considera-se como- sujeito. d aquella. .\par b) A \emph{differenga} determina o \emph{genero} para este constituir uma certa \emph{especie}, \emph{e} não outra, Porquanto, O \emph{genero}, exprimindo o elemento comm 2 varias especies, e sendo por isso indifferente para constituir uma ou outra das especies subordinadas, deve ser limitado a uma certa especie e não a outra, Esta limitação é feita pela \emph{differença}. --- Y'este modo resulta a \emph{especie}, composta de \emph{genero} e de \emph{differenca}. Estes dois elementos não se distinguem \emph{physico} ou \emph{realmente}, mas só \emph{logica} ou \emph{mentalmente}; pois a \emph{especie} não é uma terceira realidade composta de duas autras realidades, mas é uma terceira idea composta de duas ideas incompletas.\par €) As \emph{differengas devem de algu nodo} ser contidas sob o \emph{genero}; aliás\pgn{30}£09 i " Lun}

%%>>> t29

\pg{30} \tese{29}{Divisão do genero, da especie e da differenca.}
Cada

uma d'estas tres ideas \emph{universaes} --- \emph{genero}, \emph{especie} e \emph{difference}

--- divide-se em \emph{suprema}, \emph{influe} ou \emph{proxima}, e \emph{intermedia}.

a) Genero \emph{supremo} é o que acima de si não tem outro ge-

nero e que por isso não pode ser \emph{especie}, --- \emph{Infimo}, ou \emph{proximo},

é o que abaixo de si não tem outro genero, mas só especies.

--- \emph{Inleriuedio} é o que tem; acima e abaixo de si, outro genero,

e, ao mesmo tempo mas sob. diverso aspecto, é \emph{genero} e especie.

b) Especie \emph{suprema} é a que está immediatamente abaixo do D genero \emph{supremo}... --- \emph{Infima} é a que está abaixo do genero \emph{in}-

: \emph{fimo}. --- \emph{Inlermedia} é a que está abaixo do genero \emph{intermedio}.

€) Differenca \emph{suprema} é a que, unida ao genero \emph{supremo},

constitue as especies \emph{supremas}, --- \emph{Infina} é a que, unida ao genero Él \emph{infino}, constitue as especies \emph{infimas}. --- --- \emph{Jntermedia} é a que,

\emph{b}\textasciicircum{} unida ao genero \emph{intermedio}, forma as especies \emph{intermedias} \footnote{\pgn{31}altribuir-se no £\emph{odo}; assim. se nesta proposição --- o \emph{homem é racional}, \emph{a differen}. "\& « racional » significasse só uma \emph{parte} da essencia humana, teriamos que o hometm é o ser dotado \emph{exclusivamente} de \emph{razão}: 0 que é falso.\par Todavia, o \emph{genero} e a \emph{difference}, se denotam a essencia \emph{completa} dos entes, não a denotam \emph{explicitamente}, Porquanto o \emph{genero}, \emph{explicitamente}, só denota a \emph{parie} essencial, comuni a varias \emph{especies}; e. \emph{implicitamente}, devota a outra \emph{parte} essencial, \emph{propria} de cada \emph{especie}. Como tambem a \emph{differença}, \emph{explicitamente}, só denota a \emph{parte} essencial, \emph{propria} de cada \emph{especie};.. e, \emph{implicitamente}, exprime a \emph{parte} essencial, commun a varias especies,\par Em resumo: --- \emph{especie} significa a essencia \emph{completa}, e exprime-a \emph{explicitamente}; --- O \emph{genero} significa a essencia \emph{completa}, mas \emph{explicitamente} só denota a. parte essencial, que é \emph{comiam} a varias especies e que por Isso é indeterminada e determinavel; --- a \emph{differença} tambem significa a essencia \emph{completa}, mas \emph{explicitamente} só denota a parte essencial, que é \emph{propria} de cada \emph{especie} é que por isso é determinada e determinante.\par A \emph{especie} é \emph{universal-reflexa} por tudo o que ella significa e denota; o: \emph{genero} e a \emph{differença} são \emph{unicersaes-vefiexas} só pelo que denotam \emph{eaplicitamente} (Cf. S. Th., \emph{Op}. 30, eo. 4 e 5 --- \emph{Op}. 48, tr. T. c. 4).}.

não \$e lhe poderiam unir. Dizemos --- de \emph{algum modo}; porque as \emph{differencas}

não estão contidas \emph{em acto} sob o \emph{genero} (aliás deveriamos altríibuir a este,

ao mesmo tempo e sob o mesmo aspecto, predicndos contradietorios), mas

estão contidas em \emph{potencia}, em quanto o \emph{fenero}, que é por si indifferente,

\emph{pode} ser determinado por ellas. !

d) A \emph{difference}, por exprimir uma qualidade \emph{essencial}. que, unida ao

\emph{genero}, constitue a \emph{expecie}, não pode deixar de ser \emph{positiva}, isUQ G. deve

exprimir uma \emph{perfeição} ou \emph{realidade}. E Se, ás vezes, é \emph{negativa}, isto é, se

indica a \emph{ausencia} de perfeição ou realidade, não devemos \textbf{pensar} que a

\emph{differença}: seja mera \emph{negação} (pois a \emph{negação} é \emph{nada} e o \emph{nada} não pode ser DT \emph{elemento essencial} das coisas), mas devemos entender que a qualidade, a qual " \emph{E} determina O \emph{genero} a uma certa especie, nos é desconhecida, e que pode ser N dotada de toda e. qualquer essencia, menos d aquella que é afastada do EL \emph{genero} pela differenca \emph{negativa}. Assim, quando dizemos que o cão é \emph{animal} - \emph{irracional}, esta diiferença « \emph{irracional}» não denota que a qualidade \emph{essencial}, D que unida ao genero « \emph{animal}» constitue o cão. seja uma \emph{negação} ou \emph{axsen}- \emph{Po}. \emph{cia}. de realidade, mas significa que essa qualidade nos é desconhecida, e que \emph{i} = pode ser tudo, menos \emph{racional} (Cf, S. Th., \emph{De Ente et Essentia}, c. 6). \emph{B} (1) Adverlimos que a especie, o \emph{genero e a differença}, quando se attrif buem aos entes, que lhes estão subordinados, se tomam no sentido, que uU estas ideas \emph{wuniversaes-reflexas} denotam \emph{explicitamente}, e não no sentido, \textbar{} n que ellus exprimem \emph{implicitamente}. Sejamos mais claros. 5 Não só a \emph{especie}, mas tambem o \emph{genero} e a \emph{differenca} significam a \emph{especie EN completa}. Na verdade, o \emph{genero} e a \emph{differença} attribuem-se, por \emph{identidade}, V aos entes, que lhes estão subordinados, pois dizemos: o \emph{homem É animal}, \emph{o E homem é racional}, Ora estes \emph{dois} wnicersaes, se não signilicassem a essencia \emph{e completa}, mas só um dos elementos essenciaes, não poderiam attribuir-se, \emph{o} por \emph{identidade}, aos seres, que lhes estão \emph{sujeitos}, \emph{pois a parte} não pode --- ] [ \emph{O Ed}

%%>>> t30

\pg{31} \tese{30}{Propriedade.}
\emph{Propriedade} é uma idea universal, representativa de uma qualidade, que \emph{necessariamente} emana da essencia \emph{completa}; tal é no homem a \emph{faculdade de fallar}. --- Para que, pois, uma qualidade possa dizer-se, \emph{propriedade}, é necessario --- que derive da essencia \emph{completa}, --- que derive necessariamente, --- e que por isso convenha a \emph{todos} os entes dotados da mesma \emph{essencia}, e só a elles, e convenha \emph{sempre} \footnote{A \emph{propriedade} não constitue a essencia, mas deriva da essencia \emph{constituída}, c deriva \emph{necessariamente}; de modo que. constituida a essencia, resulta a \emph{propriedade}. --- Por isso a \emph{propriedade} --- convem a \emph{todos} Os que têm essa essencia, porque, derivando \emph{necessariamente} da essencia, deve encontrar-se onde se encontra a essencia, --- convem \emph{exclusivamente} aos que têm essa essencia, porque só os que têm uma essencia podem participar do que necessariamente deriva d' ella, --- e convem sempre, porque 0 ente não pode perder a propriedade, sem perder a essencia. --- Deve, porém, advertir-se que uma qualidade, para ser verdadeira propriedade, deve derivar da essencia completa, isto é, do conjuncto de \emph{todos} os elementos, que constituem a essencia; porque, se derivasse \emph{de} «\emph{im} só elemento, não poderia convir \emph{exclusivamente} á uma especie, mas conviria tambem a outra especie, que é constituida por aquelle elemento; assim \emph{dormir}, com relação ao homem, não é uma \emph{propriedade} no sentido rigoroso da palavra, porque deriva da essencia \emph{incompleta} do homem, isto é, deriva do homem emquanto \emph{animal}, e por isso convem tambem ao bruto.\par A \emph{propriedade} distingue-se da \emph{especie}, do \emph{genero} e da \emph{differenca} ; porque estes Lres unicersaães constituem, no todo ou em parie, a essencia; mas a \emph{propriedade} emana, ainda que \emph{necessariamente}, \emph{da}. essencia \emph{eoustitida}.}.

%%>>> t31

\tese{31}{Accidente.}
\emph{Accidente} é uma idéa \emph{universal}, represendcc \pg{32}  LOGick --- SECQAO PRBEIRA i

tativa de uma qualidade, que \emph{indifferentemente} pode dar-se ou deixar de dar-se num sujeito; tal é a \emph{bondade} no liomem ().

223. As ideas, universaes-reflexas são cinco. --- As ideas \emph{universaes-reftexas} denotam os diversos modos, por que uma coisa pode existir em muitos sujeitos e ser-lhes attribuida. Ora os diversos modos, por que uma coisa pode existir em muitos sujeitos e ser-llies attribuida, são \emph{cinco}. Por quanto, o que.existe em muitos sujeitos --- ou \emph{pertence} u essencia d'elles, --- ou \emph{adhere} a essencia Já \emph{constituida}. --- Se pertence á \emph{essencia}, --- ou constitue a essencia \emph{completa}, e temos o \emph{especie}, --- ou apenas constitue \emph{uma parte} da essencia, e temos o \emph{genero} ou a \emph{differenca}, segundo essa parte. da essencia é \emph{indeterminada} e \emph{determinavel} ou é \emph{determinada} e \emph{determinante}. --- Se \emph{adhere} a essencia jà \emph{constituída}, --- ou adhere \emph{necessariamente}, e temos a \emph{propriedade}\footnote{O \emph{accidente} distingue-se da \emph{propriedade}; pois esta emana \emph{necessariamente} da essencia constituida, e aquelle pode adherir ou deixar de adherir a mesma essencia. --- Note-se que uma qualidade, para ser \emph{eecidente}, não é necessario que possa e deva separar-se \emph{realmente} do sujeito, mas hasta que possa separar-se \emph{mentalmente}. Assim a «loura no cysne é \emph{accidente}; visto que, \emph{embora} ella. não possa seperür-se \emph{reelanente} do eysne, todavia --- \emph{podemos} pensar que \emph{o} eysne perea a alvura, sem perder a sua essencia.\par Das noções expost resultam os seguintes corolarios :\par 3) O £6 \emph{genero} e uw \emph{especie} altribuem-se \emph{em concreto} aos entes, que lhes estão subordinados, quando estes entes são \emph{substancias}. isto é, seres subsistentes: por si; assim dizemos: o homem é \emph{animal}, e não podemos dizer: o homem é «\emph{qninalidade}. Porquanto, nesses juizos allirmamos a existencia da identidade entre o sujeito e o predicado. Ora essa identidade não existiria, se a \emph{especie} e O \emph{genero} se altribuissem em \emph{abstracto a substancia}; pois a \emph{substaneim}, por ser uma coisa. \emph{completa}, não é identica a uma essencia \emph{abstracta}, 3. qual, excluido o sujeito a que adbira, é uma .coisa \emph{incompleta}.\par 1) O \emph{genero} e a \emph{especie} attribuem-se em \emph{abstracto} aos inferiores, quando estes entes são \emph{accidentes}, isto é, são entes que precisam de outro ente, a que adhiram: assim dizemos: a alvura é \emph{côr}, e não podemos dizer; x alvura é \emph{colorida}. Na verdade, o \emph{genero} e a \emph{especie}, quando se altribuem ao \emph{accidente}, devem exprimir o que é essencial ao \emph{accidente}, e não \emph{o} que lhe é \emph{extranho}; porque esses dois \emph{universaes-reflexos} são predicados \emph{essenciaes}. Mas, 86 O \emph{genero} e a \emph{especie} se aliribuissem em \emph{concreto}, exprimiriam a \emph{essencia} do accidente, e tambem o \emph{sujeito}, em que a \emph{essencia} existe e que, poc ser um ente \emph{substancial}, \emph{é extranho} a idea de \emph{accidente}.\par 6) A \emph{differenca}, \emph{a propriedade}, o \emph{accidente} aliribuem-se aos sujeitos \emph{em concreto}; assim dizemos: Pedro é \emph{racional}, \emph{honesto}, etc., pois esses tres universaes assumem a forma de «\emph{adjectivo}, e todo \emph{o} «\emph{adjectivo é concreto}.} \pg{33} , --- ou adhere \emph{contingentemente}, e lemos o \emph{accidente}. \textbf{Logo}

as ideas \emph{universaes-reflexas} são \emph{cinco} (*).

%%>>> art-1

\section*{Artigo IV.}
\addcontentsline{toc}{section}{Artigo IV.}

\emph{Reducgao das ideas á categ gorias}

%%>>> t33

\tese{33}{Categorias.}
As nossas ideas, consideradas \emph{objectivamente}, são tantas e tão diversas, que, para evitar a confusão, devem reduzir-se a certas classes universaes. Estas \textbf{classes} \emph{x} chamam-se \emph{categorias}. --- \emph{Celegorias}, polis, são Os \emph{generos supre}-

.) mos, a que se reduzem e sob os quaes se subordinam todas as : = ideas, consideradas \emph{objectivamente}, e que representam todos os predicados, que podem altribuir-se \& um sujeito (\textasciicircum{}). , :

%%>>> t34

\tese{34}{Requisitos das categorias.}
As \emph{categorias} devem ter "5 7\emph{s} os seguintes requisitos: --- «) a sua significação deve ser clara c e determinada; porque, se for confusa e incerta, não. podem. servir para ordenar ideas: --- a) uma não deve incluir outra,” aliás teremos uma repetição inutil: --- 0\emph{c}) devem ser abstractas \textbar{} --- e universaes, de maneira que possam representar os \emph{generos} d \emph{supremos} das nossas ideas.

%%>>> t35

\tese{35}{As categorias das ideas são dez.}
Cada \emph{idea} repre- ---. senta um: \emph{objecto}. Serão, pois, tantas as categorias das \emph{ideas}, ; L. quantas forem as dos \emph{objectos}. As categorias dos \emph{objectos} são \emph{des},

(1) Devemos explicar estas duas affirmagóes --- que toda a idea repre- : --- senta uma \emph{essencia} e que a idea de \emph{propriedade} e de \emph{accidente} não constitue . . a essencia, mas uma coisa que \emph{adhere} a essencia já \emph{constituida}. Na ver- \emph{Ps} dade. toda a idea representa uma essencia. Esta essencia encontra-se, não . \emph{e} Só na substancia, mas tambem, enrbora secundarirmente, no accidente; pois E aos tambem o \emph{accidente} é uma coisa determinada. e vão é outra, e uma coisa é o que é pela sua \emph{essencia}. O \emph{accidente}, pois, tem á sua essencia e todavia não constitue a \emph{essencia} da \emph{substancia}. \emph{A substancia} tem a Sua essencia, constituida por elementos, que por isso se chamam substanciaes, e tem propriedades ou qualidades, cada uma das quaes é dotada de essencia propria, ------. a qual não \emph{constitue}, mus \emph{utodifice} 2. propria substancia.

(8) Os antigos davam ás \emph{categorias} o nome de \emph{praedicamenta}, porque estas, «, --- como dizemos, representam os \emph{predicados}, \emph{que} se \emph{attribuem} ou podem attri- : buir-se a um sujeito; assim como \emph{dava o} nome \emph{de praedicabiig} ás cinco D ideas \emph{universaes-reflexas}, de que nos occupámos no artigo antecedente. --- Às id \emph{categorias} dillerem das cinco ideas \emph{universaes-refiexas}, pois estas exprimem Os diversos modos, por que uma coisa pode ser ativibuida a outra; ao passo que ás \emph{categorias} denotam as proprias \emph{coisas}, \emph{que} são ou podem ser aliribuidas a um determinado sujeito por um d'aquelles diversos \emph{modos}.

Vol. Y \& \pg{34}  34

a saber: \emph{Substancia}, \emph{quantidade}, \emph{qualidade}, \emph{relação}, \emph{acção}, \emph{paixão}, \emph{localisação}, \emph{quandocação}, \emph{estado} e \emph{habito}: pois não ha objecto, que não se reduza a alguma d'estas. Logo as categorias das \emph{ideas} são \emph{dez}; e não ha idea, que não possa comprehender-se em alguma d'estas classes. --- A enumeração é completa e satisfaz aos requisitos indicados no numero antecedente \footnote{Os escholasticos exprimiam as \emph{dez} calegorijas com estes versos:\par Arbor sec servos ardore --- refrigerat ustos; (\emph{Substancia}) (\emph{quantidade}) (\emph{relação}) (\emph{qualidade}) (acção) (paixão)\par Ruri eras stabo; Sed tunicatus etc.\par (\emph{locatisagdo}) \textbar{} (qwamdocacao) (\emph{estado}) (\emph{habito})\par As categorias não podem ser mais do que as «\emph{ez} indicadas no texto. Na verdade, as \emph{categorias}, representando os diversos \emph{predicados}, que podem , attrimir-se a um sujeito, devem ser tantas, quantas são as perguntas, que podem fazer-se ácerca de um individuo. Ora essas perguntas não são mais nem; menos de \emph{dez}. Assim, ácerca de \emph{Pedro}, podemos pergunlar: O que é? é homem (\emph{substancia}). É srande? é pequeno (\emph{quantidade}), ---- De quem é filho? é filho de Paulo. (\emph{relação}). --- Como é? é bom (\emph{qualidade}). --- O que está fazendo? escreve (acção) --- O que solíre? soffre sêde (\emph{paíxão}). --- Aonde? em Coimbra (\emph{localisacao}). --- Quando? hoje (\emph{quandocação}). --- Como está? sentado (\emph{estado}). --- O que veste? a batina (\emph{habito}) --- Acabadas estas perguntas, já não ha mais que investigar. Logo us \emph{categorias} são \emph{dez}.\par Foi Aristoteles quem primeiro reduziu todas as ideas a \emph{dez} categorias, que são as enumeradas no texto. --- Es; classificação é a mais completa e racional, que se conhece. Na verdade, todo o \emph{predicado} --- ou exprime uma coisa, que existe em si, --- ou exprime uma coisa, que existe noutra, como no seu \emph{sujeito}, que a sustenta. Se a coisa existe \emph{em si}, é substancia; se existe num \emph{sujeito}, é accidente, \emph{Substancia} e \emph{accidente} sã 0, pois, as classes supremas, ás quaes.se reduzem todos os predicados; porque Ludo \emph{o} que existe \emph{em si} está contido sob a categoria da substancia, e tudo o que existe num \emph{sujeito} está contido sob a categoria do \emph{accidente}. --- O \emph{accidente} suhdivide-se em nove classes. Por quanto, o que existe noutro ente, como no seu \emph{sujeito}, refere-se a este por \emph{nove} modos; poisque, --- ou é a extensão das parles materiaes do mesmo sujeito, e é \emph{quantidade}, --- ou é uma perfeição, que determina o sujeito, e é \emph{qualidade}, --- ou é uma ordem mutus entre o sujeito e outro individuo, e é relação, --- ou é o exercicio da energia, de que o sujeito é dotado, e é \emph{acção}, --- ou é uma alteração, que o sujeito recebe pela acção de uma causa externa, e é \emph{paixão}, --- ou é uma determinação do sujeito quanto ao \emph{logar}, que occupa, e é \emph{localisação}. --- ou é a determinção do sujeito quanto á \emph{duração}, c é \emph{quandocação}, --- ou é uma disposição do sujeito no Jogar, e é \emph{estado}, --- ou é uma denominação, que se dá ao sujeito por causa do vestido, e é \emph{habito}. --- Aristoteles, expondo a sua doutrina no livro das \emph{Categorias}, diz que não entendeu dar uma Gnumeração nem\pgn{35}absolutamente \emph{exclusiva}. nem absolutamente \emph{perfeita}, mas que só desejou apresentar uma, que lhe pareceu a mais npta para uma coordenação \emph{logica} das \emph{ideas} e dos \emph{objectos}.\par Note-se a diflerencs entre o \emph{accidente universal}, que é uma das cinco ideas universaes-rellexas, e o \emph{accidente categorico}. de que nos oceupamos. O accidente \emph{wnibversal} significa. uma. qualidade \emph{contingente}, que pode \emph{indiffe}. . \emph{ventemente} dar-se ou deixar de dar-se no sujeito, O accidente \emph{catenorico} exprime um ente, \emph{que existe} num \emph{sujeito}, mas abstrahe do modo por que existe, não exprimindo se existe \emph{necessaria}, Qu \emph{contingentemente}. Por isso, entre a substancia e o accidente \emph{universal} pode collocar-se a. \emph{propriedade}, que reside \emph{necessariamente} na substancia; mas entre 2 Substancia.e o accidente \emph{caíegorico} não ha meio fermo, e a \emph{propriedade}, pelo facto de existir no sujeito, reduz-se ao accidente \emph{categorico} (\emph{Sum}. \emph{Th}; p. I. q. 77, a. 1 ad 9).\par Alguns philosophos apresentaram diversos systemas de \emph{categorias}.\par As categorias de Victor \emph{Cousin} são duas: \emph{Knie infinito} e \emph{ente finito}. --- As de \emph{Schleiermacher} são duas: \emph{sujeito} e \emph{predicado}. --- As de \emph{Lotce} são tres: Teues, \emph{atiributivas} e \emph{relativas}. --- Às de \emph{Sigwart} são quatro: \emph{coisa}, \emph{propriedade}, \emph{actividade} e \emph{relação}. --- As de \emph{Rosmini} são tres: \emph{ontologicas}, \emph{dialecticas} e \emph{ideologicas}. --- As de \emph{Spinoza} e \emph{Descartes} são tres: \emph{substancia}, \emph{airibulo} e \emph{modo}. --- As de \emph{Locke} tambem são tres: \emph{sabstancia}, \emph{modo} e relação. --- As de \emph{Leibnitz} e de \emph{Wolff} são seis: \emph{substancia}, \emph{quantidade}, \emph{qualidade}, \emph{paixão}, acção e relação. --- \emph{Renouvier} estabeleceu á \emph{relação} como base das categorias, --- Para \emph{Krause}, a. categoria fundamental é o ser, considerado nos tres differêntes aspectos de \emph{these}, \emph{anfithese} e \emph{Synthese}.\par As mais celebres categorias, que os modernos oppõem ás de. Aristoteles, são as de \emph{Kant}. \textbf{Este} philosopho reduziu todas as ideas \emph{puras} (meramente subjectivas) a \emph{doze} categorias, dispostas em quatro classes : \emph{Unidade}, \emph{Plurolidade} ou Multidão, \emph{Totalidade} (Quantidade), -- \emph{Realidade} ou \emph{Affirinacao}, \emph{Negação}, \emph{Limitação} (Qualidade), --- \emph{Substancia} e \emph{Modo} ou \emph{Accidente}, \emph{Causa} € \emph{Xffeito}. Acção e \emph{Paixão} (Relação), --- \emph{Possibilidade} e \emph{Impossibilidade}, \emph{Existencia}, e \emph{Não-existencia}, \emph{Necessidade} e \emph{Contingencia} Modalidade).\par De todas estas categorias diremos em geral --- \emph{a}) que algumas são de menos para abranger todo o objecto, ou são de mais e por isso se incluem umas nas outras: --- 0) que outras são falsus. porque não correspondem a realidade das coisas, ou se fundam em falsos principios.\par A classificação de Kaot não pode acceitar-se pelas seguintes razões. --- \emph{a}) Algumas das suas categorias são iranscendentes, vagas e indeterminadas, como são as dispostas sob a \emph{Qualidade} e \emph{a Modalódede}, e por isso não podem servir Dara a coordenação das ideas. --- \emph{b}) Às categorias das tres primeiras classes, segundo Kant, são superiores ás das 4.º, quando é proprio o contrario, porque não esiste \emph{realilade}, nem \emph{pisralédede}, \emph{non causa} e \emph{effeito}, que se não conceba como \emph{possivel} ou \emph{existente}. --- c) Algumas d'ellas\pgn{36}estão contidas em outras: assim a \emph{umidade} e a \emph{pluralidade} estão contidas na \emph{realidade}; \emph{a mnio-existencia} reduz-se 3 \emph{possibilidade} ou a \emph{dmpuossibilidade}. --- \emph{d}) Outras estão mal dispostas : assim a \emph{substancia}, sendo uma coisx \emph{absoluta}, não pode estar subordinada á \emph{relação}; \emph{á quantidade}, por depender da \emph{substancia}, não pode ser anterior á esta: a «ope € 3. \emph{paixão}, a cansa € o \emph{effeito}, por não serem relações, não podem estar subordinadas á \emph{relação}; o «\emph{ccidenie}, em geral, não pode tambem collocar-se sob a \emph{relação}, porque ha apenas um genero de \emph{accidente}, que se chama relação. --- \emph{e}) As categorias de Kant são meramente \emph{swhjecficis}, sem \emph{relagão} alguma a realidade; quando é certo que ás categorias são \emph{objectivas}, porque representam os predicados reaes, que se podem \emph{attribuir} a um \emph{sujeito}. --- \emph{f}) Finalmente, \emph{o} systema. de Kant é falso na:sua propria base. As categorias, Como as entendeu sempre o senso commum, são ideas, ou conceitos objectivos e simples, sob os quaes se coordenam todos os predicados reaes. Ora us categorias, como as entende Kant, não se referem á \emph{primeira} operação da intelligencia, mas a \emph{segunda}; não são \emph{conceitos}, são \emph{juízos}, ou, melhor, são \emph{formas de juizo}, e representam os differentes \emph{modos}, por que um predicado se attribue a um sujeito; e essus \emph{formas}, esses 1\emph{nodos} são inteiramente \emph{subjectivos}, aos quaes nenhuma realidade corresponde.}.

%%>>> t36

\tese{36}{Modo de considerar as categorias.}
As categorias podem considerar-se num duplice estado --- \emph{em si} e na \emph{intelligencio}. Consideradas \emph{em si}, são entes reaes, \emph{universaes-dárectos}, \pg{35}  e possuem attributos, de que a natureza as dotou. Consideradas na \emph{intelligencia}, são entes \emph{logicos}, \emph{universaes-reffexos}, eiém attributos, que lhes não convém na realidade, mas. foram. formados pela abstracção. --- Em quanto entes \emph{veaes}, as \emph{categorias} são objecto da \emph{Até(ephysice}; pois esta trata d'aquellas, para lhes \pg{36}  conhecer a essencia, e as \emph{propriedades}, que emanàm da essencia. Em quanto entes \emph{logicos}, constituem o objecto da \emph{Logica}; pois esta considera as \emph{crafegorias} como \emph{generos supremos}, para sob elles coordenar, como generos \emph{intermedios} e especies, lodos os predicados, que se podem, attribuir a um sujeito \footnote{As categorias, emquanto constituem o objecto da \emph{Metaphysica}, SãO noções reages e \emph{objectivas}, e podem ser atiribuídas aos objectos, que existem fóra da intelligencia; ao passo que, emquanto são objecto da \emph{Logica}, São noções \emph{logicas} e \emph{subjectivas}, \emph{e} não podem attribuir-se aos objectos existentes na realidade; porque, lomadas neste segundo sentido, as categorias são ideas universaes-reflexas, isto é, se consideram como \emph{generos} e \emph{especies}, que hão de ser coordenadas e dispostas conforme a diversidade da sua extensão. Por isso, se a categoria \emph{metapluysica} é constituida exclusivamente pelo \emph{genero supremo}, considerado nos seus atiribulos \emph{reaes} e \emph{objectioos}, a categoria \emph{logica} é toda a serie dos predicados essenciaes, coniprehendendo o genero \emph{Supremo}, \emph{OS} generos e as especies \emph{intermedias}, \emph{Me} a especie \emph{infini}.}.

%%>>> t37

\tese{37}{Definição das categorias.}
Reservando vara a \emph{Metaphysica} uma conveniente explicação das \emph{categorias}, aqui damos apenas a \emph{definição} de cada uma d'ellas: que é quanto basta para a coorderiação das ideas. \textbar{} P categ. - \emph{Substancia} é a coisa, que para existir não carece de

sujeito, a que adhira; assim o \emph{homem é substancia}. "IF \textbf{eateg}. \textbar{} - \emph{Quantidadeé} a propriedade, pela qual a substancia

é divisivel em partes; tal é a \emph{grandeza} na pedra. UF categ. - \emph{Qualidade} é a perfeição, que modifica a substan-

cia, sem alterar-lhe a essencia ; "aléa virtude

e a sciencia no \textbf{homem}. - IV* categ. - \emph{Relação} é o nexo, que liga uma substancia \&

, outra ; assim \emph{filho} indica uma relação. V* categ. \textbar{} - Acção É O exercicio da força, de que a substancia

é dotada; tal é escrever. VE categ. - \emph{Paixão} é a alteração, que a substancia \textbf{soffre} pela

acção do agente; tal é \emph{sentir frio}.

\emph{Localisação} é a determinação da substancia quanto VII* categ.

ao \emph{logar}; por ex., Pedro está em \emph{Roma}. VIP categ. - Quandocação .é a determinação da substancia

quanto ao \emph{tempo}; por ex., \textbf{Pedro} \emph{hontem} escreveu, IXº categ. - \emph{Estado é a-disposição} das partes quantitativas da

substancia no logar; por ex. André \textbf{está} \emph{sentado}. X* eateg. - \emph{Habito} é a denominação que se dá a substancia

por causa do vestido ou decoração; ex. o ma-.

gistrado £\emph{ogado} \footnote{Na \emph{Ontologia} trataremos das \emph{categorias} com o devido desenvolviinento. Mas a semples.explicacá\&o mostra que as \emph{categorias} SÃO, em si, noções reaes, porque denotam entes \textbf{reaes}.--- Por isso, não devemos pensar que aS \emph{categorias logicas} sejam coisas inteiramente \emph{ideaes}, 4s quaes nenhuma realidade corresponde. como erradamene imaginou \emph{Kant}. Às categorias, na \emph{Logica}, \emph{exprimem} a coordenação \emph{logica} dos predicados \emph{essenciaes} das coisas reges; e por isso alguma coisa, nellas, é \emph{ideal}, e algum coisa é \emph{real}. \textbar{} \emph{ideal} a qualidade de \emph{genero}, ou de \emph{especie}, que se concebe por uma. alstraeqüo: mas o gue esse \emph{genero}, Ou essa \emph{especie} exprime, é muito \emph{real ou objectivo}. Por quanto, o \emph{unécersal reflexo} suppõe e contem o \emph{directo}, e por isso, se tem attributos \emph{ideaes}. emquanto é \emph{reflexo}, tem attributos \emph{remes} Ou \emph{nhjectivos}, emquanto couvtem o \emph{directo}. Esta é a razão por que a definição das categorias exprime um coisa real ou objectiva.}.

%%>>> t38

\pg{37} \tese{38}{Ordenação da primeira categoria.}
Não basta determinar o numero das categorias; é mister estabelecer a ordem em cada uma d'ellas. --- Limitar-nos-hemos a apresentar a ordenação da \emph{primeira}, que é a \emph{substancia}.

\emph{Substancia} é o genero supremo. Divide-se em \emph{corporea} e \emph{incorporea} (ditferencas supremas) --- Substancia \emph{corporea} é o \emph{corpo} (especie suprema).

\emph{Corpo} (1.º genero inter medio) pode ser \emph{animado} e \emph{inanimado} (1.º differengas intermedias) \emph{Corpo animado} constitue o vivente (1\textasciicircum{} especie intermedia).

\emph{Vivente} (2." \textbf{genero} intermedio) divide-se em \emph{sensitivo} e \emph{insensitivo} (2. differenças intermedias). \emph{Vivente sensitivo} é o \emph{animal} (2. especie intermedia).

Animal (genero infimo) é \emph{racional} ou \emph{irracional} (ditferencas infimas). \emph{Anánal racional} é o \emph{homem} (especie infima).

\emph{Homem} é só \emph{especie}, c abaixo de si só contém \emph{individuos}.

Eis o quadro: \emph{E}

\textbf{Substancia}

(gen. \emph{supr}.) \textbar{} :

corporea. (incorporea)

(\emph{diff}. supr.)

\textbf{Corpo}

(\emph{esp}. \emph{supr}. e 1.º \emph{gen}. \emph{énterm}.)

animado : (inanimado) (\emph{L}5 \emph{diff}. \emph{énterm}.)

\textbf{Vivente}

(\emph{I}\textasciicircum{} \emph{esp}. \emph{iuferm}. \emph{e} 2.º \emph{gen}. \emph{feiern}.)

sensitivo (insensitivo) (2.º \emph{diff}. \emph{interim}.) ; \textbar{}

\textbar{} \textbf{Animal}

] (2-* esp. \emph{énierm}. \emph{e gen}. \emph{infimo})

racional (irracional)

(\emph{if}. \emph{inf}) C : MEME

\textbf{Homem}

(\emph{esp}. \emph{inf}.)

Pedro, Paulo, André \footnote{Façamos algumas rellexões acerca da ordenação da categoria de \emph{substancia}, --- O genero \emph{substancia} contém a especie \emph{corpo}, DÃO na sua \emph{comprehensáo}. mas na sua \emph{extensão}; porque a \emph{especie} tem uma comprehens\&o maior que o \emph{genero} (x \emph{especie} contém o \emph{genero}, e, alem d'esse, a \emph{differenca}). --- Estando as especies superiores contidas nas inferiores. a estas competem as propriedades d'aquellas: assim. ao homem compelem as propriedades de \emph{animal}, de \emph{vivente}, de \emph{corpo}, e, por consequencia, de \emph{substancia}, --- As \emph{differencas} dividem-se em \emph{consiitutivas} e \emph{divisivas}. \emph{Constitntivas} são as que constituem o \emph{genero}, e competem a todas as especies, que estão contidas debaixo d'aguelle genero. \emph{Divisivas} são As que \emph{dividem} o \emph{genero}. c por isso não são \emph{todas} participadas por todas ás especies, mas uma differença é participada por uma especie, e outra por outra: Assim \emph{animado} e \emph{sensiliuo} são dillerenças \emph{constitutivas} do \emph{animal}, é por isso são communs ao homem € 30 bruto; ao passo gue \emph{racional} e \emph{irracional}, que são differencas \emph{divisivas} (porque dividem \emph{o animal} não são communs ao homem e ao bruto, mas uma (\emph{racional}) é propria do homem, a outra (\emph{irracional}) é propria do bruto.\par Na ordenação d'es!a categoria só consideramos as differenças \emph{positivas}, As differencas \emph{negativas}, unidas ao respectivo genero, dão logar a outras classificações. A substancia \emph{incorporea} é o \emph{espirito}, que se subdivide em innumeraveis especies. O corpo \emph{inanimardo}, é e \emph{mineral}, de que se occupa\pgn{39}a Mineralogia. O vivente \emph{insensitivo} É a \emph{planta}, de que trata a Botanica. O animal \emph{irracional} é 0 \emph{bruto}, que forma o objecto da Zoologia.\par Agora uma advertencia. Os seres, que possuem a essencia completa do genero \emph{supremo} e dos \emph{infermedios}, como são as \emph{especies} e oS \emph{individuos}, \emph{directamente} e \emph{por} si se collocam na relativa categoria. Os seres, porém, que só \emph{incompletamente} participam do genero \emph{supremo e} dos énfermedios, como são os \emph{elementos} Ou as \emph{modificações} da especie ou do individuo, só \emph{indirectamente} se reduzem ás categorias, e se collocam sob aquella, em que se encontra a especie ou o individuo, de que são elementos ou modificações.}.

%%>>> art-2

\section*{Artigo V.\\ y}
\addcontentsline{toc}{section}{Artigo V. --- y}

\pg{38} \pg{39} \emph{Termo}, \emph{suas especies e propriedades}

%%>>> t39

\tese{39}{Termo.}
\emph{Termo}, ou \emph{palavra}, ou \emph{vocabulo}, é a \emph{expressão verbal} da \emph{idem} --- O homem, sendo naturalmente \emph{social}, deve manifestar aos outros as proprias ideas. Esta manifestação extrinseca só pode effectuar-se por meios \emph{sensiveis}; pois é proprio do homem, composto de alma espiritual e de corpo, elevar-se ao immaterial pelo sensivel, Entre todos os meios o mais efficaz e universal é o \emph{termo} \footnote{Entendemos fallar, como se vé, do termo \emph{oral}, ou \emph{vocat}, que exprime sensivelmente \emph{a idea} (chamada \emph{termo mental}) --- Diz-se \emph{termo}, porque é nelle que se resolve e acaba, como no seu ultimo elemento; a \emph{proposição}.\par \emph{A idea} e \emph{o termo} differem entre si pelas seguintes razões ; --- 1º) a \emph{idea exprime zutmediatamente} o \emph{objecto}: O \emph{termo} exprime \emph{smmediatamente a idea C mediatamente} o \emph{objecto}; --- 2º) a \emph{idea é anterior} ao \emph{fermo}, pois este só se emprega. e \emph{fórma} para exprimir aquella, e sem \emph{ideas} não haveria \emph{termos}; --- 3\textasciicircum{}) pela \emph{ideas} fallamos commoseo, indicando o objecto percebido: pelo \emph{terno} fallamos com. os outros; --- 4º) a \emph{idea} é signal \emph{natural}, a0 passo que o \emph{fernto} é Signal \emph{concencional}, como vamos demonstrar.\par Todavia.o nexo entre a \emph{ideo} e o \emph{fermo} é tão estreito, que não podem Separar-ge. Não podemos conceber uma \emph{idea}, sem que ao mesmo tempo se forme \emph{na} nossa phantasia o \emph{fermo}, que a representa e exprime; como \emph{tambem o fermo} é a propria \emph{idea} coucretisada e tornada sensivel.}.

%%>>> t40

\tese{40}{O termo é signal convencional.}
O \emph{termo}, --- a) é \emph{signal}, porque nos manifesta a idea, concebida pela intelligencia, --- \emph{b}) mas é signal \emph{convencional}, porque, embora nos seja natural a manifestação das nossas ideas por meios sensiveis, comtudo foi o livre arbitrio dos homens que ligou a um vocabulo determinado a significação de ama certa idea. --- No \emph{fermo}, como na \emph{idea}, considera-se a \emph{comprehensão} e a \emph{extensão} (\textasciicircum{}).\footnote{O \emph{fermo}, alem de ser signal \emph{convencional}, é tambem signal \emph{objectivo}: \emph{pois}. antes \emph{de levar} a \emph{intelligencia} ao conhecimento da \emph{idea} e do \emph{objecto}, O \emph{termo} deve ser primeiramente conhecido em si mesmo.}

%%>>> t41

\pg{40} \tese{41}{Especies do termo.}
O \emph{termo}, sendo a \emph{expressão verbal da idea}, divide-se em tantas especies, quantas são as especies da idea, considerada em relação ao \emph{objecto} representado \footnote{Ao £\emph{ermo} não podem convir as divisões, que são proprias da \emph{idea}, considerada em relação no \emph{modo} por que representa o objecto; pois o \emph{fermo} limita-se a exprimir o \emph{objecto}, existente na intelligencia, mas prescinde do \emph{modo} por que o mesmo objecto é percebido.}.

Para evitarmos repetições inuteis, consideraremos a divisão do \emph{fermo} --- \emph{aj} em \emph{positivo} e \emph{negativo} --- \emph{b}) em «\emph{absoluto} e \emph{connotativo} --- c) em \emph{abstracta} e \emph{concreto} --- \emph{d}) em \emph{singular}, \emph{universal} e \emph{particular} --- \emph{e}) em \emph{wnivoco}, \emph{equivoco} e \emph{analogo}. .

a) Termo positivo e negativo. --- Termo \emph{positivo} 6.0 que exprime una \emph{realidade}, ou \emph{perfeició}: camo o termo Deus, \emph{caridade}. --- \emph{Negativo} é o que denota a «\emph{wsencia} de \emph{realidade}, ou de \emph{perfeição} ; como o termo \emph{nada}, \emph{trevas} \footnote{Para conhecermos se um termo é \emph{positivo} ou \emph{negativo}, devemos considerar a \emph{idea}, que elle exprime. Se a. idea é \emph{positiva}. isto é, se representa uma \emph{perfeição} Ou \emph{realidade}, o termo é \emph{positivo}; como tambem se a idea é \emph{negativa}, isto é, se denota «\emph{wsencia} de \emph{realidade} ou \emph{perfeição}, o termo é \emph{negativo}. Assim são \emph{positivos} os termos: \emph{immenso}, \emph{immortal}, \emph{infinito} ete.: e são \emph{negativos}. os termos: \emph{morte}, \emph{peccado}, \emph{vicio}, ete. --- O termo \emph{negativo} distingue-se do \emph{privativo}. Porquanto. o \emph{negativo} denota apenas a \emph{ausencia} de realidade, sem indicar se \emph{o} sujeito é capaz ou incapaz d' essa realidade; ao passo que o \emph{privativo} significa a. \emph{ausencia} de uma perfeição nam = sujeito, que a devia possuir, tal é o termo \emph{cego}, com relação ao homem.}.

b) Termo absoluto e connotativo. --- Termo \emph{absoluto} é o : que significa uma coisa, gue é \emph{substancia}, ou que é considerada

o \emph{d} : como \emph{substancia}; tal é \emph{planta}, \emph{bondade}. O \emph{absoluto} chama-se tambem \emph{denotativo}. --- \emph{Connotativo} é o que exprime uma \emph{propriedade} ou \emph{qualidade}, que \emph{existe} e \emph{emquanto existe} numa substancia; tal é bom, \emph{sabio} \footnote{Come se vê, o termo \emph{connofativo} exprime, ao mesmo tempo, duas coisas: uma \emph{qualidade} e um \emph{sujeito}, on ama. substancia, que a possue. O termo \emph{absohdo} denota \emph{directamente} o \emph{sujeito} com as suas qualidades; o \emph{connotativo} significa. \emph{directamente a qualidade}, existente no sujeito. --- Por isso a. \emph{denotaczo} 6 \& propriedade de termo \emph{denofatico} ou \emph{absoluto}, emquanto este indica um \emph{sujeito}, ou uma \emph{substancia}: como a \emph{connmolação} é a propriedade do termo \emph{connotativo}, emquanto elle indica uma \emph{gualidede} ou um \emph{attributo} de um sujeito ou de uma substancia. --- O termo \emph{absoluto}, ou \emph{denoninativo}, costuma chamar-se tambem «\emph{substantivo}, e O \emph{connotativo} diz-se tambem \emph{adjectivo}. Todavia o \emph{adjectivo} tem mais extensão na Logica do que na Grammatica: porque muitas vezes os lermos, que ua Grammatica são \emph{substantizos}, na Logica são \emph{adjectivos}, \emph{ou cosnolatiros}: taes são Os levmos \emph{uceidico}, \emph{pintor}, musico, etc.. que exprimem uma \emph{qualidade} adherente a nm \emph{sujeito},}.

c) Termo \textbf{abstracto} e \textbf{concreto}. --- Termo \emph{abstracto} é o que exprime uma essencia ou propriedade, considerada em si mesma, sem relação ao sujeito, que a possue; tal é o termo \emph{humanidade}, \emph{vida}, sabedoria, etc. --- \emph{Concreto} é o que denota um sujeito com a sua essencia ou propriedade, camo existe na ordem real; tal é o termo \emph{homem}, vivente, \emph{sabio}, etc. (\textasciicircum{}).

d) Termo \textbf{singular}, universal e particular. --- O termo \emph{singular} é o que exprime uma essencia, \emph{propria} de um \emph{individuo determinado}, de modo que só a este individuo pode ser attribuida; tal é o termo \emph{Pedro}. --- \emph{Universal} é o que exprime ama \emph{essencia}. \emph{comuni} a muitos entes, do modo que a todos e a cada um d' estes pode ser attribuida : tal é o termo \emph{homem} --- \emph{Particular} é o proprio. termo :\emph{niversaó}, mas limitado na sua \emph{extensão}; pov ex, \emph{dlguue home}, \footnote{Todo o termo \emph{cannotatévo é concreío}; mus nem todo o .\emph{concreto É connotativo}; assim O termo \emph{Tromenm}. 6 \emph{concreto}, mas não é \emph{connobativo}. --- O \emph{concreto} pode ser, indifferentemente. \emph{adjectivo} € \emph{substantivo}; assim os termos \emph{sabio} e homem são ambos \emph{concretos}.}.

e) Termo univoco, equivoco e analogo. --- Termo :\emph{univoco} é o que se attribue a varios entes, mas sempre na \emph{m significação}: tal é o termo \emph{komen}, quando se altribue a \emph{Pedro}, \emph{Paulo}, \emph{André}, ele. --- \emph{Equivoco} É o que se attribue a. varios . entes, mas em \emph{diversa significação}; tal é o termo \emph{leão}, quando so attribue a uma \emph{constellação} e a um certo \emph{animal}. --- \emph{Ánalogo} é o que se attribue a varios entes puma significação, que não é inteiramente \emph{identica}, nem inteiramente \emph{diversa}; tal é o termo \emph{são}, quando se aíliibue ao \emph{animal} e ao ar \footnote{O termo \emph{universal}, como idea \emph{snicersal}, distingue-se do termo \emph{transcendental}. do comun, do \emph{collectico}, \emph{e} do \emph{particular}. .}.\footnote{O termo pode \emph{ser equivoco}, mas não a \emph{idea}. Por quanto o \emph{fermo}, semdo um signal \emph{consenciemal}, pode emprejrr-se livremente para denolar entes diversos; mas a idea, por ser um signal \emph{matural}. não pode representar senão und unico objecto. ---- Para se conhecer se um termo é \emph{iarivoro}, OU. \emph{equivoco}, é preciso eonusiderur os sujeitos, aos quaes se attribue. Se os sujeitos convém entre si não só na denominação, mas lambem nua essencia, 0 termo\par \emph{ssencia}, ainda que convenham na é wnicoco; se os sujeitos não convém ut\par in. o termo é \emph{equicoco}. --- É preciso evitar o uso dos termos \emph{equédenontins}\par is o meio, que nos foi dado por Deus para a manifestação dos cocos: \textbar{} nossos pensamentos, converlerese-ha num instrumento de confusão e de ruina social. Quando, pois, se empresa um termo, que tem differentes signilicacões, é necessario determinar o seu sentido,}

%%>>> t42

\pg{41} \pg{42} \tese{42}{Propriedades do termo.}
As principaes \emph{propriedades} do \emph{termo} são \emph{seis}: --- \emph{supposição}, \emph{appellação}, \emph{estado}, \emph{distracção}, \emph{ampliação} e \emph{restricção}.

a) Supposição. --- \emph{Supposicato} é a \emph{accepcüo} ou \emph{posição} do \emph{termo em logar da coisa}, \emph{que elle exprime}. --- Nào podendo iniroduzir no discurso as coisas, de que se trata, servimo-nos dos signaes, que façam as \textbf{vezes} d'ellas. Estes Signaes são os \emph{termos}, --- São tantas as especies da \emph{supposição}, quantos são os modos, por que um termo pode substituir uma coisa ().

b) Appellação. --- \emph{Appellação} é a \emph{applicação} de um \emph{termo} «\emph{a outro}: assim, quando digo: Deus é \emph{santo}, applico o termo \emph{santo} . ao termo Deus. --- O termo, que se applica, diz-se \emph{appellante}; o termo, a que o outro se appliea, chama-se \emph{appellado}. Nas \emph{proposições}, o \emph{predicado} é o termo \emph{appellante}, e o \emph{sujeito} é o termo \emph{appellado} \footnote{Para que \emph{um} termo (\emph{appellante}) possa appliear-se a outro \emph{fappellado}), \emph{aguelle} deve referir-se de alium modo a este, Ora um termo pode referir-se á outro, \emph{ou} por causa da \emph{forma}, que se encontra neste. on por causa do \emph{sujeito}, que possue tal forma. Se um termo é \emph{appellado} por causa da sua.\pgn{43}\emph{forma}, á appellação é \emph{formal}: se é \emph{appetlado por} causa do \emph{sujeito}, \emph{a appellação} € \emph{material}. Assim quando digo: o \emph{medico} cura \emph{os doentes}, a. appellação é \emph{formal}, porque o medico cura em quanto possue a scicheia da medicina. pela qual se chama e é medico: mas. quando digo: \emph{o mer/ro} come, a. appellação € \emph{material}, porque elle come. não emquanto é medico, nas emquanto é \emph{homem}, que, aliás, pode possuir e possue a sciencia da medicina. --- Para distingüirmos se a appellação é \emph{formal} ou \emph{material}, devemos attender ás seguintes regras: --- 1.) O termo \emph{appellante} emprega-se sempre \emph{formalmente}, € não \emph{materialImente}. Assim, quando digo: \emph{o homem é bom}, o termo appellante \emph{bom} iloma-se \emph{formabnenie}, isto é, por causa da \emph{ferma}, que é a \emph{bondade e} que reside no sujeito; porque, se se tomasse \emph{materialmente}, isto é, "por causa do \emph{sujeito}. o sentido da proposição seria: o \emph{homem é homem}. --- 9.) Quando \emph{o} termo \emph{appellado} é \emph{concreto}, toma-se \emph{materialmente}. e não \emph{formalmente}. Assim, quando dizemos: Deus morreu \emph{por nós}, o Vermo \emph{anpellado}. que é Deus, toma-se \emph{materialmente}, isto é, por causa do \emph{sujeito}, e não \emph{formaliente}, isto €, não por causa da \emph{forma}, que é a Natureza divina; porque não entendemos que morren \emph{a} natureza divina, mas que morreu o Sujeito, ou a Pessoa, que possue a natureza divina; e morreu quanto á sua natureza bumana.}. : \textbar{} :

(1:4 supposição divide-se primeiramente em propria e \emph{impropria}, A supposição \emph{prapria} é a accepeto do termo em Jogar da coisa. que o mesmo Lermo \emph{exprime} na sua significação \emph{propria}, como neste exemplo: \emph{Coimbra} é \emph{uma}. \emph{cidade} de Portugal. \emph{Inzpropria} é a accepeüo do termo em logar da coisa, que o mesmo termo denota na sua significação \emph{Zumpropria}, ou \emph{metaphorica}, como: \emph{Coimbra} é O \emph{cerebro} de Portugal. --- A \emph{supposiczo propria} pode ser \emph{real} e \emph{logica}. \emph{A real} é a accepeüo do termo em logar da coisa, considerada quanto aos seus attributos reces, como: o \emph{Romem é} £\emph{ntelligente}. A \emph{logica} é a. accepção do termo em logar da coisa, considerada quanto aos seus attributos \emph{logicos}, como : o \emph{homem é especie}. --- A supposição \emph{real} pode ser \emph{abso}- «\emph{ta} e \emph{pessoal}. \emph{A absoluta} é a posição do termo em lograr da essencia, considerada sem relação aos individuos, como: \emph{o homem é animal racional}. A \emph{pessoal} é a posição!do termo em logar-da essencia, emquanto esta existe nos individuos, como: \emph{o homem} morre, (porque não morre a essencia humana, mas sim o individuo. humano). --- A SUpposição \emph{pessoal} pode ser \emph{singular}, \emph{Tenéeersal} ou \emph{tlüstribuliva}, \emph{particular}, e \emph{collecti va}; conforme o termo representa \emph{jun sh} individuo, ou \emph{todos} € \emph{cada} um dos individuos contidos sob a sua exlensio, ou apenas \emph{alguns} individuos, ou a \emph{collecção} dos individuos. --- A \emph{supposicio particular é determinada}, ou. \emph{confusa}, conforme o termo representa \emph{alguns} dos seus inferiores \emph{determinados} (ex. \emph{a/guns Apostolos} escreveram o \emph{Evangelho}) ou \emph{éndeterminados} (ex. \emph{alguns soldados} são necessarios para a \emph{vietorin} - A sopposicio \emph{mrirergal é completa} on \emph{incompleta}, conforme o \emph{termo representa} --- £\emph{odos} e \emph{cada um} dos \emph{individuos} (ex. \emph{lodo} o \emph{animal} é sensitivo), --- \emph{ou todas} e \emph{eua} uma das \emph{especies}, contidas sob o \emph{genero}, significado \emph{pelo temo} (\emph{ex}. \emph{todo o animal} entrou na area de Noé).

c) Estado. --- \emph{Estado} é a \emph{accepcüo do termo no tempo}, que \emph{o verbo exprime}; assim quando digo: \emph{o surdo não} ouve, o termo \emph{surdo} toma-se no \emph{tempo presente}, significado pelo verbo \emph{não} ouve, e o sentido é: \emph{quem} presentemente é \emph{surdo}, Jão Ouve.

d) Distracção. --- \emph{Distracção} é a \emph{aceepção} do \emph{termo num tempo diverso d'aquelle}, que \emph{o verbo exprime}; assim quando, digo: \emph{os cegos} vêem, o termo \emph{cegos} não deve tomar-se no tempo \emph{presente}, como indica o verbo \emph{vêem}, mas.no tempo \emph{passado}, de modo que o sentido e: \emph{os que foram cegos}, \emph{hoje} véem. --- A \emph{distracção} oppóoe-se \emph{ao estado}.

€) Ampliação. --- \emph{dmpliação} é a \emph{extensão do termo de uma significação} menor « \emph{uma outra maior}; quando dizemos: \emph{bemaveníurados os mansos}, \emph{porque possuirão} a \emph{terra}, o termo \emph{meansos}, que pelo verbo \emph{possuirão} só deveria referir-se aos \emph{futuros}, extende-se a todos os mansos, \emph{que foram}, são e serão.

f) Restricção. --- \emph{Restricção} é a \emph{coarctação do termo de uma significação maior a uma outra} menor; assim nas palavras de Jesus aos Apostolos: '\emph{Pregai o Evangelho} a \emph{toda a creatura}, o termo «\emph{toda} \& \emph{creatura} » limita-se exclusivamente ás creataras \emph{humanas}. --- A \emph{vestriccito} oppõe-se á \emph{ampliação}.

ÁRrTIGO VI

\emph{Definição} e \emph{divisão}

%%>>> t43

\pg{43} \pg{44} \tese{43}{Definição, --- Definição é a explicação de uma coisa ou \textbar{} de um termo.}
Assim, quando dizemos: «o homem é \emph{animal racional} », damos a \emph{definição} do homem. --- Serve para conhecermos a \emph{comprehensão} das \emph{ideas} e dos \emph{termos} \footnote{A \emph{definição} serve para fixar-o estado da questão, ou para dar uma noção exacta das coisas, e assim evitar inuteijs controversias.\par A \emph{definição} é 0 \emph{principio} e \emph{O termo} da sciencia. Purtimos de uma definic\&o incompleta ou provisoria, para chegarmos-a uma definição completa, e \emph{perfeitamente} conhecida. É o que acontece com as ideas, que de obscuras e confusas se tornam claras e distinctas. A \emph{definição}, pois. é é \emph{principio} para os que apprebendem. mas é \emph{o fermo} para os que já appreheuderam.\par 2) Relativamente á definição \emph{momiwol} fazemos as seguintes obsoervações; --- 1?) só devem ser definidos os termos \emph{obsenros}. \emph{incertos} c \emph{equicocos}; --- 2*j \emph{a} definição \emph{eomaurum} deve ser seguida sempre: aliás não seriam possiveis nem as sciencias, hem as soc ;o--- 3?) a definição \emph{particular} só deve ser \emph{apresentada em} caso de necess\par (\$1 Acerca \emph{da} definição \emph{real}, notamos: --- l\textasciicircum{}) A definição \emph{essencial} deve, quando seja \emph{possivel}, preferir-se a \emph{deseriptica} ; pois aquella explica a \emph{not}}.

%%>>> t44

\tese{44}{Especies da definição.}
A \emph{definição é nominal} e \emph{real},

a) A definição \emph{nominal} é a que explica o \emph{fermo}; por ex, a Philosophia é o amor da sabedoria. --- Divide-se em \emph{etymologica}, \emph{con} e \emph{particular}, \emph{A etymologica} explica o \emph{termo} pela \emph{etymologia}, como: a Theologia é o discurso acerca de Deus. A cowmui explica o termo, declarando o sentido em que os homens geralmente o tomam, como: \textbf{Deus} é o ser necessario. A \emph{particular} explica o \emph{fermo}; indicando o sentido, que lhe se applica, como: a natureza é a collecção dos seres creados \footnote{\pgn{45}ponatur» In VII Met, l 1)}

b) A definição \emph{real} é a que explica a \emph{coisa} significada pelo \emph{fermo}; por ex., a Philosophia é a sciencia das causas supremas dos entes, descobertas pela luz natural da razão. --- Divide-se em \emph{essencial} e \emph{descriptiva}. \emph{A essencial}, de todas a mais nobre, explica a coisa pelos seus elementos \emph{essenciaes}; e subdivide-se em \emph{logica} e \emph{physica}, conforme esses elementos --- são \emph{logicos} (como: o homem é animal racional), --- ou são \emph{physicos} (como: o homem.é uma substancia, composta de alma racional e de corpo organico) A \emph{descriptiva} explica a coisa pelos elementos, que se referem á mesma coisa, mas que não pertencem á sua essencia, como são as propriedades, as \emph{notas individuaes}, etc. (7).

%%>>> t45

\pg{45} \tese{45}{Requisitos da definição.}
A \emph{definição} \textbf{deve}:

L. Ser composta do \emph{genero} proximo e da \emph{differenca especifica}; "porque d'este modo comprehenderá todos os elementos, (que constituem a essencia da coisa definida \footnote{Deve seguir-se esta rexra, quando se deseja a definição \emph{logico}, \emph{que} é a mais excellente, --- Todo o ente creado é composto de dois elementos\pgn{46}essenciaes, --- um comunwm e \emph{indeterminado}, outro \emph{proprio} e \emph{determinado}. Pelo elemento \emph{cosumua} e \emph{indeterminado}, \emph{o} ente é semelhante aos outros: pelo elemento \emph{proprio} e \emph{determinado}, distingue-se dos outros e fórma uma classe especial. O elemento comuna e \emph{indeterminado} é o \emph{genero}; o \emph{proprio} e \emph{determinado} é a \emph{differenga}. \emph{especifica}. Logo a \emph{definição}, para explicar a \emph{essencia} de um ente, deve apresentar o \emph{genero} e a \emph{differenca}, ---- Mas note-. Sé que o \emph{genero} deve ser \emph{proximo}; pois o remoto não contém todas as notas essenciaes da coisa definida. Assim. é falsa esta definição: o homem é o \emph{vivente racional}; pois que o \emph{vivente}, sendo genero \emph{remoto} de \emph{homem}, não inclue a nota de \emph{animalidade}, pela qual o homem é semelhante aos brutos, nem o extrema do Anjo. que tambem é \emph{vivente racional}. --- Se não podermos determinar o \emph{gênero proximo} e a \emph{differença especifica}. ou porque o ente é \emph{si»mtplicissimo} c não consta de varios elementos. ou porque não temos um conhecimento \emph{distincto} das coisas, devemos então designar duas propriedades, das quaes x primeira exprima a semelhança, que a coisa definida tem com as outras, ea segunda denole o que é exclusivamente proprio d' ella. Assim, não podendo dar uma \emph{rigorosa} definição de Deus, pois Elle é simplicissimo e transcende todos os \emph{generos} e todas as \emph{differencas}, dizemos que Elle é o \emph{Ente dotado de perfeição infinita},\par (\emph{D} Isto é evidente ; porque a delinicüo serve para declarar uma coisa obscura ou desconhecida. Por isso. na definição \emph{nominal}, devemos empregar os termos mais conhecidos, e, na. definição \emph{real}, devemos recorrer a conceitos, anteriormente explicados. --- Contradizem a este regra --- 1º) os que definem a coisa pela propria coisa. como: \emph{homen} é o \emph{ser} que \emph{possue a humanidade}; --- 2º) os que em pregam termos \emph{metaphoricos}, \emph{indeterminado}\$, \emph{equivocos} ou \emph{insolitos}, como: \emph{os olhos são as estrellas do homem}; --- 3") os que, sem \emph{necessidade}, se servem de termos \emph{negativos}, que indicam o que a coisa não é e não o que ella é, como: \emph{o homem é não-pedra}. Dizemos --- sem \emph{necessidade}; porque, ás vezes, ou pela limitação da nossa iatelligeneia, ou pela propria natureza das coisas, devemos recorrer a termos \emph{negativos}; assim dizemos que \textsc{Deus} é \emph{immenso}, que O \emph{cego} é o homem que \emph{não} vi.}. \emph{o} : \emph{reza} das coisas, e esta pára nº superficie. --- 2º) É absurdo exigir a. dofinigüo \emph{essencial} de todas as coisas; pois, se tudo precisasse de \emph{definição}, uma coisa devia ser definida por outra, e esta por noutra, até ao infinito: o que é absurdo, porque, então, nenhuma coisa podia ser definida. Por isso todas as definições devem” resolver-se em certas noções elementares, que, por serem evidentes, não precisam de definição. Estas noções são: \emph{ente}, \emph{unidade}, \emph{verdade}, etc. "Podavia d'estas noções podemos fazer uma descripção ou declaração por meio de noções, que sejam, \emph{relativamente} X nós, mais conhecidas. --- Alem d'isso, as mais elementares noções, por serem simplicissimas e constarem de uma unica \emph{rota}, não podem estar sujeitas a uma verdadeira definição, pois esta resolve a coisa nos seus varios elementos.

A definição \emph{deseriptiva}, que explica a coisa por elementos extra-essenciaes, divide-se em \emph{propria}, \emph{accidental}, causal e \emph{genetica} --- \emph{Propria} é a que explica a coisa pelas \emph{propriedades}, que derivam da essencia da mesma coisa; como: o homem é \emph{capas de rir} e de \emph{apprehender} « \emph{sciencia}. --- 'decidental é a que explica a coisa pelas suas qualidades: \emph{accidentaes}, que, no. sem \emph{conjuncto}, SÓ podem convir á mesma coisa; assim Virgilio (\emph{Aen}., 111) descreve Poliphemo: \emph{Monstrwumn horrendun}, \emph{informe}, \emph{ingens}, cui lumen \emph{ademplunt}, irinca. nmanmwn \emph{pinus regit}, \emph{et vestigia}. \emph{firmat}. --- Causal é a que explica a coisa pelas causas \emph{exivinsecas}, que São --- 3 \emph{efficiente} (como: o \emph{homem foi creado} por \textsc{Deus}) --- a \emph{exemplar} (como: o \emph{Romem foi creado} d, \emph{imagem} de \textsc{Deus}), --- á \emph{final} (como: \emph{o homen}. \emph{foi cre«do} para que dime u \textsc{Deus}) --- A \emph{genetica} É \& que explica a . €oisa pela. sua \emph{origeit}, como: « \emph{eclipse do sol deriva da interposição da} lut \emph{entre o sol} e a \emph{terra}. --- A definição \emph{descriptica} só num sentido \emph{laío} e \emph{improprio} se Chama \emph{definição} ; porque, embora faça distinguir uma coisa das outras, não explica o que é a coisa. --- Todavia, esta definição emprega-se frequentes vezes nas sciencias, que não penetram no interno das coisas, para lhes descobrir os elementos essenciaes.

Devemos advertir que de muitas coisas pode dar-se uma definição absoluta, e de muitas outras uma definição \emph{relativa}. Dá-se uma definição \emph{absoluta} das coisas, que, pela sua natureza, não se referem a outras; e esta delinição indica exclusivamente os elementos, que constituem a essencia da propria coisa ; assim a \emph{substancia} é susceptivel de uma definição \emph{absoluta}, Dá-se uma definição \emph{relativa} das coisas, que se referem essencinhnente 3 outras. e. esta definição indica tambem a coisa, á qual uma outra se relere; assim a definição do \emph{accidente} exprime a \emph{substancia}, a qual-o accidente se refere. Diz S. Thomaz: «\emph{Substantia}, quae habet quidditatem (essentiam) \emph{absolute}, non dependet in su\& quidditate ex alio. ita quod oporteat exteriorem essentimu in ejus definitione poni. \emph{Aceidentia} vero nom habent esse nisi per hoc quod iusunL in subjecto, et ideo eorum quidditas est dependens \& subjecto; et proptec hoc oportel quod subjectum accidenlis in delinilione

Il. Ser \emph{reciproca}, de modo que a \emph{definição} e a \emph{coisa defi}-

\emph{nida} possam converter-se uma na outra. Esta condição deriva

da primeira. É, pois, exacta a definição: «o homem é animal -

racional ».. pois podemos. dizer: « o animal racional é homem»,

II. Ser \emph{clara}; aliás não consegue o seu fim \footnote{Sendo a \emph{definiczo} como que a directriz de toda a investimação e diseussão. deve ser tal que.possa, i cada instante e com facilidade, ser citada: \emph{o} que vão poderia. fazer-se, se não fosse \emph{breve}. Dor isso, quando se defne o \emph{nome}, empresuém-se somente os termos necessarios e sufficientes para que se conheça a sua significação: e. quando se define a coisa, indiquem-se somente squellas notas, que explicam a. mesma coisa e a distinguem das outras: assim é defeituosa esta definição: o \emph{homem é animal racional} e \emph{sagas}, porque a. nota \emph{sagas} está contida na nota \emph{racional}.\par Aristoteles (II \emph{Post}, \emph{Anal}.) ensina doia methodos para procurar e encon-}.

IV. Ser \emph{breve}, para'poder conservar-se na memoria \footnote{\pgn{47}trar os elementos essenciaes, de que se compõe a definição, especialmente \emph{logica}; um chama-se \emph{descendente}, outro \emph{ascendente}. --- O methodo \emph{descendente} purie do que é mais universal para descer ao que é menos universal. Por quanto, toma o \emph{genero}, ao qual está sujeita a coisa, que se ha de definir, e a esse genero successiviumente accrescenta tantas diiferencas especificas, quantas são necessarias pura determinar e limitar o genero, de modo que: esta possa converter-se na coisa definida, Então temos a definição; porque aquelle genero com as dilfereuças, á excepção da ultima, constitue o genero \emph{proximo}, € a ultima diffevença é a \emph{difference especifica}. Assim, devendo definir \emph{o homen}, tomamos a \emph{substancia}, como \emph{genero}, sob o qual o homem está contido; a este genero de \emph{substancia} accrescentamos as diflerengas \emph{corporeo}, \emph{amimada}, \emph{sensitiva}, \emph{racional}, e temos a definição \emph{lorica} do \emph{homem}; porque o, genero \emph{substancia} com as tres primeiras differencas é o \emph{genero proximo}, e a ultima differenca \emph{racional} é a \emph{differenca especifica}. --- O methodo \emph{ascendente}, fazendo o caminho inverso, parte do que é menos universal para subir ao que é mais universal. Porquanto. começa por comparar a coisa, que se ha de definir, com as coisas, que mais lhe se approximam. \textbf{Esta} . comparação mostra que a coisa tem um elemento commum com as outras coisas, pelo qual Ibes é semelhante, e tem outro elemento \emph{proprio}, pelo qual diifere das outras. O elemento comm é O \emph{genero}: o elemento \emph{proprio} é a \emph{differenga}. Assim, querendo enipregar este methodo para definir o \emph{kromem}, comparamos o homem com o bruto. D'esta comparação resulta que c homem convem com o bruto emquanto é «\emph{animal}, mas difflere do bruto emquanto é \emph{racional}. Por isso \emph{animal} é o \emph{genero} na definição de \emph{homen}, e \emph{racional} é a \emph{differenca}..}.

%%>>> t46

\pg{46} \pg{47} \tese{46}{Divisão.}
\emph{Divisão} é a \emph{decomposição de} um \emph{todo} nas suas \emph{partes}, ou nos seus \emph{elementos}. --- A \emph{divisão} leva-nos ao conhecimento da \emph{extensão} das \emph{ideas} e dos \emph{termos} \footnote{Pela \emph{definição}, que extrema o objecto de tudo o que lhe é extranho, reconhecemos \emph{a unidade} do mesmo objecto; pela \emph{divisão}, que distingue as partes ou elementos do objecto, reconhecemos a sua \emph{pluralidade} Ou \emph{complexidade}, Pela \emph{definição}, 0 objecto torna-se \emph{claro}; pela \emph{divisão}, \emph{distincto}. \emph{A divisão}, pois, é o complemento da \emph{definição}. --- A \emph{utilidade} da divisão é inconteslavel. A nossa intelligencia é tão limitada, que não pode, com um só olhar, comprehender as coisas, sohretudo se são complexas, e, para as comprehender, \emph{deve} esmiuçabas e resolvel-as nas suas partes.}.

%%>>> t47

\tese{47}{Especies .da divisão.}
A \emph{divisão é logica} e \emph{real}.

a) A divisão \emph{logica} é a decomposição de um \emph{todo logico}, ou \emph{universal}, nas diversas partes, que o constituem; tal é a divisão do \emph{genero} nas diversas \emph{especies}, de que é composto.

b) A divisão \emph{real} é a decomposição de um \emph{todo real}, ou \emph{objectivo}, nas \textbf{suas} diversas partes; tal é a divisão de um organismo. --- A divisão \emph{real é physica} ou \emph{metaphasica}, segundo o \emph{Lodo real} é composto de partes distinetas --- ou \emph{realmente} (tal é \pg{48}  o homem composto de \emph{corpo} e de ala), --- ou só \emph{logicamente} (tal é o homem composto de \emph{animalidad}2 e \emph{racionclidade}) \footnote{A divisão \emph{physica} pode ser --- \emph{essencial}, ou \emph{integral}, oa \emph{potencial}, ou \emph{accidental}, segundo o \emph{todo physico} € composto de partes --- \emph{essenciaes} (tal é o homem composto de \emph{alum} e \emph{corpo}), --- ou \emph{integrantes} (tal é um \emph{orgasismo}), --- cu \emph{polenciaes} (tal \& a nossa alma, que é dotada de muitas e diversas \emph{potencias} Ou \emph{faculdades}), --- ou \emph{accidentaes} (tal é o Lodo \emph{artificial}).}.

%%>>> t48

\tese{48}{Regras da divisio.}
A \emph{divisão} deve: \emph{u}

l. Ser breve; de contrario produzirá confusão.

Il. Ser \emph{adequada}; de modo que a collecção das partes enumeradas seja egual ao todo» São, pois, defeituosas --- tanto as divisões em que, se omitte alguma parie do objecto dividido, ---. como as que enumeram uma parie, que lhe não pertence.

IH. Ser \emph{ordenada}; de modo que o todo primeiramente se resolva nas partes mais proximas, e depois nas mais remotas.

1V. Ser \emph{opposta}; de modo que as partes se excluam mutuamente: aliás unia parle será repelida sem utilidade \footnote{A" \emph{divisão} reduz-se a \emph{classificação}, frequentemente empregada pelos naturalistas, --- Diremos da sua \emph{definição}, \emph{especies}, \emph{graus} é \emph{requisitos}.\par 3) Classificação. --- \emph{Classificação} É a disposição ordenada de muitos individuos em classes. --- Todos os seres; como já notamos, se relacionam por alguma semelhança, A nossa intellizencia, considerando o caracter, em que muitos individuos são semelhantes, toma-o como base da \emph{classificação}.\par b) Especies de classificação. --- A classificação é \emph{artificial} ou \emph{natural}. \emph{É artificial}, quando o caracter, que se toma por base, consiste numa. propriedade \emph{extrinseca} aos individuos. É \emph{materal}, quando 0 caracter, que se estabelece como fundamento, é uma propriedade \emph{essencial} aos individuos,\par c) Varios graus de classificação. --- Os naturalistas distinguem varios graus nas suas classificações. O mais alto grau é o reino (genero supremo): o \emph{reino} divide-se em \emph{typos} (1.º genero intermedio): o \emph{typo} em \emph{classes} (2.º genero intermedio): a \emph{classe} em \emph{ordens} (3.º genero intermedio); a \emph{ordem} em . \emph{tribus} (4.º genero intermedio); a £\emph{ribw} em \emph{familias} (5.º genero intermedio); ua \emph{familia} em \emph{generos} (genero infimo); O \emph{genero} em \emph{especies}; a \emph{especie} em \emph{variedades}; \emph{a variedade} em \emph{individuos}. --- Note-se a. dilferenga entre especie e \emph{variedade}. \emph{A especie} é a collecção de individuos, que participam da mesma natureza e geranr outros individuos semelhaotes. A \emph{variedade} é a collecção de individuos, que, por causa de caracteres particulares accidentaes, se distinguem dos outros individuos da mesma especie. A \emph{variedade}, que. por meio da geração, transmitte 0s proprios caracteres, chama-se raça.\par d) Requisitos da classificação, --- A classificação, por ser uma divisão, deve salislazer a todas ás regras, que demos para esta, ------- Além disso, é neces-. sario que o caracter, que se estabelece como base da classificação, --- a) seja \emph{unico}, porque d' outro modo não poderemos conhecer as mutuas relações dos seres; --- \emph{b}) seja \emph{permanente}, porque, se fôr transitorio, acabará a classifienção com o caracter; --- \emph{c}) seja \emph{importante}, aliás será inutil.}.

%%>>> cap-2

\chapter*{CAPITULO SEGUNDO}
\addcontentsline{toc}{chapter}{CAPITULO SEGUNDO}

\pg{49} ES \textbf{JUIZO} E \textbf{PROPOSICAO} : \textsc{SUMMARIO}: 7--- Juizo \& proposição, sua divisão, --- Divisão da propo- “ gição simples. -- Divisão da proposição composta. --- fQualidades das

proposições.

: ArTIGO \textbar{}

\emph{Juizo e proposição}, sua \emph{divisao}

%%>>> t49

\tese{49}{Juizo.}
\emph{Juizo} é a \emph{união} ou a \emph{separação de duas ideas}, e \emph{por isso de dois objectos}, \emph{por meio da affirmação} ou \emph{da negação} ;: como, «o homem é livre» --- «o homem não é escravo». --- O \emph{juizo}, pois, é posterior á \emph{idea} (!*).

%%>>> t50

\tese{50}{Elementos do juizo.}
Os elementos do \emph{j:?zo} são tres: as \emph{duas} ideas, e a \emph{relação}, que ellas têm entre si. --- A \emph{idea}, de que alguma coisa se aflirma ou nega, diz-se \emph{sujeito}; a \emph{idea}, que se afhrma ou nega do sujeito, chama-se \emph{predicado}; a \emph{relação}, que une o' predicado ao sujeito, denomina-se \emph{copula}. --- O \emph{sujeito} e o \emph{predicado} são a \emph{materia} do juizo: a \emph{copula} é a \emph{forma} \footnote{O \emph{juico} é necessariamente um acto da \emph{intelligencia}; porque só esta faculdade pode conhecer a relação (que é uma coisa \emph{abstracía}) entre dois objectos e por isso entre duas ideas. --- Os escriptores sensualistas, como \emph{Tracy}, dizem que o \emph{juico} é um acto da \emph{sensibilidade}; algons \emph{cartesianos} - sustentam que elle é um acto da \emph{vontade}. Uns e outros erram; porque a intelligencia, e só ella, pode conhecer a \emph{relação} entre duas ideas, e unil-as ou separal-as por meio da \emph{affirmação} ou da \emph{negação}, 20 passo que a vontade e a sensibilidade são duas faculdades cegas e não julgam. É verdade que, ás vezes, é a vontade que determina a intelligencia a pronunciar o juizo (e isto verifica-se, quando a intelligencia não vé a relação entre o sujeito e o predicado); mas o juizo é essencialmente um acto da intelligencia. --- Lembremo-nos sempre de que o juizo é a affirmac\&o da conveniencia ou da desconveniencia entre duas ideas, consideradas não subjectivanmente, como meras formas ideaes (comoTerradamente opina \emph{Kant}), mas objectivzmente, emquanto representam as proprias coisas externas, Quem ignora ou duvida se duas ideas representam, ou não, uma coisa real, não pode dizer se essas ideas convém, ou não, entre-si. O juizo, relativo i conveniencia ou descorm- .veniencia entre duas ideas, funda-se uecessariamenie na conveniencia ou desconveniencia objectiva das coisas. que as ideas representam.}.\footnote{Para a formação do \emph{juizo} São necessarias varias e diversas operações da intelligencia. Primeiramente percebemos \emph{o sijeito} e o \emph{predicado}; em seguida comparamos estes dois elementos; finalmente pronunciamos o \emph{juizo}, allirmando que elles convém, ou não, entre si,\par O \emph{juízo}, embora considerado nos seus \emph{elementos}, \emph{seja} uma coisa multipla * composta, todavia, considerado \emph{n}\& sua \emph{essencia}, é uma coisa simplicissima; pois consiste ni \emph{affirmaczo} da conveniencia ou da desconveniencia entre o sujeito e o predicado, e esta «/\emph{firinação} € uma coisa simplicissima.\par Vol, I. 5}

%%>>> t51

\pg{50} \tese{51}{Divisão do juizo.}
O \emph{juizo} divide-se --- \emph{a}) em \emph{affirmativo} e \emph{negativo}, --- \emph{b}) \emph{verdadeiro} e \emph{falso}, --- \emph{c}) \emph{directo} e \emph{reflexo}, --- \emph{d}) \emph{analytico} e \emph{synthelico}, --- \emph{e}) \emph{inmediato} e \emph{mediato}.

a) \textbf{Juizo} affirmativo e \textbf{negativo}. --- É \emph{afirmativo}, se affirma a relação de conveniencia entre o sujeito e o predicado; ex. «a agua do Monáego é limpida ». --- É \emph{negativo}, se nega aquella relação de conveniecia; ex. «a agua do Tibre nao é limpida » ()

b) \textbf{Juizo verdadeiro} e \textbf{falso}. --- É \emph{verdadeiro}, quando ao 'sujeito convem o predicado que lhe attribuimos, ou não convem o que lhe negamos; ex. «o amor de patria é nobre» --- «o mundo não é eterno». --- É \emph{fdlso}, quando ao sujeito não convem o predicado que lhe attribuimos, ou convem o que lhe negamos; ex. « o egoismo é a norma da vida» --- «o animal não é sensitivo ». .

€) Juizo \textbf{directo} e reflexo. --- É \emph{directo}, quando a intellizencia affirma a conveniencia, ou a desconveniencia, do predicado com o sujeito, e, implicitamente, a propria conformidade com a realidade objectiva: assim quando digo: \emph{Pedro é homem}, affirmo que a \emph{Pedro} compete a \emph{humanidade}, e, ao mesmo tempo, que tal é o meu modo de ver., --- É \emph{reflexo}, quando, dobrando-se sobre si mesma, a intelligencia considera o seu proprio juízo, e reconhece a sua \emph{conformidade} com au realidade objectiva; como quando digo: \emph{estou certo de que Pedro é homem} \footnote{O \emph{juizo}, considerado na “sua \emph{materia}, div ide-se em \emph{affirmativo} e \emph{negativo}, emquanto o \emph{predicado} convem ou não Stonvem ao \emph{sujeito}. Mas, se \emph{o} consideramos emquanto é aclo da intellizencia, o \emph{juizo} € sempre \emph{a/firmativo}; pois sempre \emph{se afirma} que o \emph{pretécado} convem, ou não, ao \emph{sujeito}.}.

d) Juizo analytico e synuthetico. --- É \emph{analytico}, quando o , predicado está contido na propria essencia do sujeito, emquanto --- ou a constitue, --- ou d' ella deriva necessariamente: ex. o\footnote{D' onde resulta --- que o \emph{objecto} do juizo \emph{directo} é o proprio \emph{sujeito}, ' a que se attribue um \emph{predicado}, e o \emph{objecto} do juizo \emph{veflexo} é a propria. \emph{conformidade} da inLlelligencia com a realidade, --- e que o juizo \emph{reflexo} iuclue e suppõe necessariamente o \emph{directo}. --- Como tambem se vé o erro dos que dizem que a \emph{intelligencia}, para conhecer a sux \emph{conformidade} com a realidade objectiva, deve necessariamente \emph{reffectir} sobre os seus actos. A intelligencia, como dizemos no texto, conhece essa \emph{conformidade}, pelo facto mesmo de attribuir a um determinado sujeito um determinado predicado. Se a não conhecesse, não pronunciarià o seu juizo. --- Isto, porém, não obsta a que a intellivencia examine, pela rellexão, o seu juizo primitivo e se convença, cada vez mais, da verdade do seu acto.}

%%>>> cap-3

\chapter*{CAPITULO SEGUNDO}
\addcontentsline{toc}{chapter}{CAPITULO SEGUNDO}

\pg{51} \emph{homem} é racional. --- É \emph{synihetico}, quando o predicado convem ao sujeito, não porque está contido na essencia d'elle, ou deriva da essencia, mas porque assim o attesta a experiencia, ou o testemunho dos outros bomens; ex. \emph{Pedro é saóio} \footnote{No juizo analytico, o predicado convem ao sujeito \emph{essencialmente}, \textbar{} e por isso \emph{necessariamente}, e onde se encontra o sujeito, deve encontrar-se o predicado; 20 passo que, no juizo \emph{syníhketico}, o predicado convem ao sujeito \emph{accidentalmente}, e por isso \emph{conbingeniemenio}, e pode encontrar-se o sujeito sem esse predicado. --- O juizo \emph{analytico}, em que a relação entre o sujeito € o predicado é conhecida pela simples \emph{analyse} do mesmo sujeito, chame-se tambem --- \emph{a}/ « \emph{priori}, porque faz conhecer o objecto pelo seu \emph{principio},que é o \emph{conceito} du. sua vatureza ou essencia, --- 6] \emph{necessario}, porque o predicado convem \emph{necessariamente} ao sujeito, --- \emph{c}] \emph{metaplysico}, porque a relação entre o sujeito e o predicado se funda na propria essencia das coisas, --- \emph{d}] \emph{absoluto}, porque não admitte excepções, --- \emph{ej racional}, porque só a razão pode perceber que o predicado esta nu essencia do sujeito. --- O juizo \emph{synthetico}, em que a attribuição do predicado ao sujeito se funda na experiencia ou no testemunho, denomina-se Lumbem --- \emph{aj a-posteriori}, porque se funda na experiencia, --- 5] \emph{contingente}, em quanto o predicado só \emph{accidental} Ou \emph{contingenieménie} convem ao sujeito, --- \emph{cj physico}, porque o nexo entre o predicado e o sujeito baséa-se num facto physico, --- \emph{dj hpothetico}, porque pode ter excepções, --- \emph{e}] \emph{empirico}, porque a. relação entre \emph{o} predicado € o sujeito conhece-se pela experiencia ou pelo testemunho dos outros.\par \emph{Kant}, na sua \emph{Critica da razão pura}, admitte uma terceira classe de juizos, que não são absolutamente \emph{analyticos}, vem absolutamente \emph{syntheticos}, e que elle chama \emph{syntheticos-a-priori}.\par É evidente a falsidade d' esta opinião. Um predicado é, ou não é, da essencia do sujeito. Não ha meio termo, Se é da essencia \emph{do sujeito}, o juizo é \emph{analytico}, ou \emph{a-priori}; se o não é, \emph{o juizo é syntAtetico}. Dizer que um juizo é \emph{synthketico a-priori} € afirmar que um predicado convem e não convem essencial ou necessariamente no \emph{sujeito}: o que repu;ma ao senso commum. --- Mais tarde voltaremos \emph{a} este assumpto.}.-

e) \textbf{Juizo} immediato e mediato, --- É \emph{immediato}, se a relação entre o predicado e o sujeito se conhece \emph{intuttivaanente} ; ex. o mundo existe. --- É \emph{mediato}, se aquella relação só se conhece pelo discurso; ex. a nossa alma é immortal ().

%%>>> t59

\tese{59}{Proposição.}
\emph{Proposição} é a \emph{expressão verbal do juizo}. --- Como o \emph{fermo} é a expressão \emph{verbal} da \emph{idea}, assim a \emph{proposição} é a expressão \emph{verbal} do \emph{juizo}. Quando dizemos : “o ceu é \emph{bellissimo}, enunciamos uma \emph{proposição}, :

53. \textbf{Elementos} da proposição. --- A \emph{proposição} tem os mesmos elementos do \emph{juízo}, a saber: o \emph{sujeito}, o \emph{predicado} e a\footnote{O juizo \emph{inmmediato} diz-se tambem \emph{primitivo}, porque não \emph{suppõe} outro juizo; e o \emph{mediato} chama-se tambem: \emph{derivado}, \emph{pois} deriva de \emph{outro} juizo,\par Dl 4\emph{a} da cmas im \emph{ea} ete mr m e d} \pg{52}  \emph{copúla}. ---\emph{JQ sujeito} exprime-se, em geral, pelo \emph{nome substantivo};

o \emph{predicado}, pelo, nome \emph{adjectivo}: a \emph{copula}, pelo \emph{verbo substantivo} «\emph{ser}», posto no presente do \emph{indicativo}. --- O \emph{verbo adjectivo} contém o \emph{predicado} e \emph{a.copula}, como: Pedro \emph{estuda} \footnote{Advertimos que os \emph{termos} da proposição devem ser tomados \emph{formalmente}, isto é, segundo o sentido, e não \emph{materialmente}, isto é, não segundo o numero das palavras; por isso'muitos termos \emph{meateriaes} podem formar um só termo \emph{formal}. Assim nesta proposição: a \emph{alia lunnana}, \emph{espiritual e immortal é creada por DEUS}, ns palavras \textless{}a \emph{aline fnmana}, \emph{espiritual e} : \emph{émmorial}» formam um só termo, que é o \emph{sujeito}.\par Digamos alguma coisa acerca do \emph{nom}: e do \emph{verbo}.\par O nome é um termo, que denota \emph{pessoas}, ou \emph{coisas}, ou alguma \emph{propriedade}, que ás pessoas ou as coisas possuem, sem denotar o \emph{tempo} (embora baja termos, que signifiquem o £\emph{empo}, como se fosse uma coisa; p. ex. \emph{dia}, \emph{hora}, ete.). --- Divide-se em \emph{substantivo}, e \emph{adjectivo}, O \emph{substantivo} significa uma. coisa, que é \emph{substancia}, ou. que é considerada como \emph{substancia}. O \emph{adjectivo} exprime uma coisa, que é accidente, ou que é considerada como accidente. Por isso, O \emph{sujeito} da proposição, que indica uma substancia, Ou uma coisa considerada como substancia, deve ser significado pelo nome substantivo; e o \emph{predicado}, que indica um accidente, ou uma coisa considerada como accidente. deve ser significado pelo nome \emph{aurjectivo}.\par O verbo é um termo, que indica o \emph{estado}, a acção ou a \emph{paixão} de pessoas ou coisas, com a indicação tambem do \emph{tempo}, em que as coisas acontecem. --- Divide-se em \emph{substantivo} e \emph{adjectivo}. \emph{O substantivo é o} que propriamente significa o acto de \emph{ser}, em que se fundam todos os outros actos, e é um só: \emph{sou}, \emph{es}, é. O \emph{adjectivo} é o que, alem do acto de ser, exprime tambem uma acção, ou uma paixão, ou ima. modificação ou variedade do ser, como são os verbos : amo, quero, etc.\par O verbo \emph{substantivo} costuma chamar-se \emph{prünitivo}, porque todos os outros verbos são modificações ou determihações d' este verbo; assim quando digo : \emph{Pedro} ama, o verbo ama exprime a acção de Pedro, fundada sobre o seu acto de \emph{ser}, e resolve-se na seguinte proposição: \emph{Perro é amante}. --- Alem d' isso, o verbo \emph{substantivo} pode tomar-se num sentido \emph{absotutív} e NUM sentido \emph{relativo}. Toma-se num sentido \emph{absoluto}, quando significa a \emph{existencia} actual, e então não é acompanhado \emph{pelo} nome \emph{adjectivo}, como : Deis é. Toma-se num sentido \emph{relativo}, quando signilica o «\emph{cto de ser} e compõe \emph{o} predicado com o sujeito, como: \emph{o homem é racional}. Tomado absolutamente, o verbo \emph{ser} é sempre \emph{real}. "Tomado \emph{relativamente}, \emph{o} verbo ser é \emph{logico} ou \emph{real}, conforme a vatureza dos termos, que une ou separa. Se os termos são ambos \emph{logicos}, ou se um é \emph{real} e outro é \emph{logico}, \emph{o} verbo é exprime um eunle \emph{logico}. como: a \emph{chimera é inpossivel}, \emph{Fedro é surdo} (a negação ou a privação não adhere ao sujeito), --- e, se os termos são ambos \emph{reaes}, o verbo é tambem é \emph{real}, e denota o ser real e absoluto, em que se funda a união do predicado com o sujeito, como: Deus é \emph{misericordioso}.}.

54. \textbf{Divisão} da proposição. --- A \emph{proposição} divide-se em \emph{simples} e \emph{composta}. --- É \emph{simples}, quando contém \emph{um só sujeito} \pg{53}  e um só \emph{predicado}, como: «a primavera é agradavel » Lf \emph{composta}, quando, \emph{explicita} ou \emph{implicitamente}, contém. \emph{mais que}

uma \emph{proposição simples}; como: «a musica e a' "poesià são bellas artes». --- “Tanto a proposição \emph{simples} como a \emph{composta} subdividem-se em varias especies \footnote{Alguns confundiram, com a proposição \emph{composta} a. complexa. Proposição complexa é a que contém uma ou mais proposições \emph{incidentes}, que determinam a significação do sujeito ou do predicado. Mas é infundada tal opinião. \& proposição \emph{complexe}, embora apparentemente composta, é \emph{simples} na realidade, visto que a \emph{incidente} não forma um termo distincto.}. ! \textbar{}

ÁrTIGO Il

\emph{Divisão da proposição simples}

55. Fundamento d'esta divisão. --- Na proposição \emph{simples} podemos considerar tres coisas: --- a \emph{materia}, de que se trata, --- a \emph{quantidade}, ou \emph{extensão} dos lermos, --- e a \emph{qualidade} do nexo entre o sujeito e o predicado.

Em relação é \emph{materia}, a proposição \emph{simples} costuma dividir-se em \emph{analytica} e \emph{synihetica}; --- em relação á \emph{quantidade}, em \emph{universal}, \emph{particular} e \emph{singular}; --- em relação á \emph{qualidade}, em \emph{affirmativa} e \emph{negativa}. \emph{DU}. \emph{Quem}

56. Proposição analytica e synthetica. --- \emph{A} proposição é

\emph{analylica}, quando exprime um juizo \emph{analytico}; ex. «o homem é racional». --- É \emph{synthetica}, quando traduz um juizo \emph{synthetico}; ex. « Pedro é sabio ».

57. Proposição \textbf{universal} particular e singular. --- É «\emph{mniversal}, quando o-sujeito é \emph{universal}; ex. «toda a creatura é finita ». É \emph{particular}, quando o sujeito é \emph{particular}; ex. «alguns homens sio bons». --- É \emph{singular}, quando o sujeito é \emph{singular}, ex. « Camões é o principe dos poetas portuguezes \footnote{Na proposição \emph{simguicr}, o \emph{sujeito} é tomado em \emph{toda} a. sua \emph{extensão} ; por isso x proposição \emph{singular} equivale é universal, e reduz-se a esta. --- A \emph{wuiversalidaxde} de uma proposição pode ser \emph{metaphysica}, \emph{physica} € \emph{moral}. \emph{É metaphysica}, quando a relação entre o sujeito e o predicado se funda \emph{na propria} essencia das coisas, e por isso não admitte excepção alguma ; como: «a alma humana é \emph{espiritual} ». É \emph{physica}, quando a relação entre o \emph{predicado} e o sujeito se funda nas \emph{leis physicas}, e por isso pode ter excepções ; como: «\emph{uenhiwin}. \emph{morto vesuscita} ». Ij \emph{moral}, quando a relação entre o sujeito c o predicado se funda nas \emph{leis moraes}, e por isso admitte excepções: como: «\emph{os} pres amam seus \emph{filhos} ». --- Quando não é evidente a extensão do sujeito, a proposição reduz-se á \emph{universal}, ou \& \emph{particular}, OU a \emph{singular}, conforme a naiurza da materia, de que se trata.}.

58. Proposição affirmativa e \textbf{negativa}. --- \emph{Affirmativa} é a que exprime a \emph{conveniencia} do predicado com o sujeito; ex. « O B. Nano Alvares foi um grande capitão ». ---- \emph{Negativa} é a que exprime a \emph{descon teniencia} do predicado com o sujeito; ex. «O orgulho não é bom conselheiro » ().

59. Comprehensio e extensão do predicado nas proposi-

ções affirmativa e negativa, --- As seguintes observações são da maxima importancia: \textbar{}

3) Na proposição \emph{afirmativa}, o \emph{predicado} toma-se em \emph{toda} a sua \emph{comprehensão}; assim na proposição: «Pedro é \emph{homem}» affimamos que a Pedro convém \emph{todos} os elementos, que constituem o predicado « homem » \footnote{Se ao sujeito fallusse uma só das notas. que constituem o \emph{predicado}, não poderia \emph{afirmar-se} que este convem aquelle,}. \textbar{}

b) Na proposição \emph{affirmative}, o \emph{predicado} é tomado só em \emph{parte} da sua \emph{extensio}, e por isso é \emph{particular}; assim quando dizemos: «as flores são plantas », não se entende que \emph{todas} mas sómente \emph{clgunas} plantas são flores. --- “Todavia, o predicado, quando exprime a mesma coisa que o sujeito, tema-se em toda a sua \emph{extensão}, como neste exemplo: «o homem é animal racional », pois que. todo animal racional é homem (O). \textbar{}

€) Na proposição \emph{negativa}, o \emph{predicado} é tomado só em \emph{parte} da sua \emph{comprehensão}: assim na proposição « o animal não é planta », só negamos ao sujeito «\emph{alguns} e \emph{não todos} os elementos do predicado, pois a vida \emph{vegetativa} é \emph{comu} ao animal e a planta \footnote{Da proposição \emph{affirinativa} e \emph{negativa} differe a proposição \emph{infinita}, --- a que tem por. predicado um termo infinito, como: \emph{o homem é não-planta}. Esta proposição não é \emph{a/firmatica}, nem \emph{negativa}; porque não nltribue ao sujeito o' predicado, nem o nega; Mas, embora não seja \emph{afirmativa} nem \emph{negativa}, comtudo participa de uma e de outra. Participa da \emph{affirmativa}, emquanto indica que pode convir ao sujeito um predicado indeterminado, á excepção do que é significado pelo termo infinito: participa da \emph{negatiuu}, emquanto afasta do aujeito um predicado determinado.} - \textbar{}\footnote{Na proposição «\emph{ffir}: \emph{mativa}, dizemos que o sujeito está contido sob a \emph{extensão} do predicado, mas não dizemos que só \emph{elle} está, contido; porque o predicado pode conter outros sujeitos. O predicado (em maior «\emph{extensio} que o sujJeito, embora Leuba menor \emph{comprehensão}.\par (9 Para se negar \emph{um predicado} a um sujeito. nio é necessario que o sujeito seja destituido de \emph{todos} e \emph{cada uma} das notas. que constituem o predicado: basi que lhe \emph{falte uni} só, pois a ausencia de mie só nota destroe a \emph{identidade} entre o sujeito \emph{e} o peediesdo.}

d) Na proposição \emph{negativa}, o \emph{predicado} é \textbf{tomado} em \emph{toda} a extensão, e por \textbf{isso} é \emph{universal}; \textbf{assim} neste \textbf{exemplo}: «o animal não é planta » o predicado toma-se universalmente, pois \emph{nenhuma} planta é animal \footnote{Um predicado nio poderia \textbf{afastar-se} de um sujeito, se um só dos entes, contidos sob a \emph{extensão} do predicado, podesse affirmar-se do sujeito.\par A proposição \emph{simples} costuma dividir-se tambem em absoluta e \emph{modal}. \emph{É} «\emph{absoluta}, quando indica se o predicado convem, ou não, ao sujeito, como: o ceu é sereno, É \emph{modal}, quando não só indica se o predicado convem, ou não, ao sujeito, mas tambem exprime o mudo. por que o predicado convem, ou repugna, ao sujeito, como: é \emph{possivel que} eu \emph{hoje parta}. --- Digamos alguma coisa acerca da proposição modal. O \emph{modo} é uma determinação da. coisa \emph{feito} pelo \emph{adjectivo}. O \emph{adjectivo} é duplice: --- do nome, como \emph{branco}, \emph{preto}, ele., --- e do verbo, taes são os \emph{adverbios}, como \emph{brevemente}. velozmente, ete., que se chamam adjectivos do verbo, porque estão junto do verbo e nelle se firmam. D'ahi o duplice modo: \emph{nominal} e \emph{adverbial}, O modo \emph{nominal} é a determinação da coisa, feita pelo adjectivo do. nome. como : curso veloz; * --- o modo \emph{adverbial} é a determinação da coisa, feita pelo adjectivo do \emph{verbo}, isto é, pelo \emph{adcerbio} como: \emph{corre velozmenie}, --- O modo pode determinar --- ou o \emph{sujeito} da proposição, como: \emph{o homem} branco corre, --- ou O \emph{predicado}, como: o \emph{homem corre} velozimente, --- ou a \emph{relação} entre o predicado c o sujeito. como: é \emph{possivel} que \emph{o homem corra}. --- \emph{A} proposição \emph{modal} é propriamente a que declara o \emph{modo por} que O predicado se refere ao sujeito. A proposição, em que o \emph{nodo} determina o sujeito ou o predicado, é \emph{absoluta}.\par \emph{Os modos}, que exprimem a relação do predicado com o sujeito, são seis: \emph{verdadeiro e falso}, \emph{necessario} e \emph{impossivel}, \emph{possivel} e \emph{contingente}. --- Os modos « \emph{serdadeiro} e \emph{falso} » não alteram a relação do predicado com o sujeito, e não acerescentam nada ao sentido da proposição \emph{absoluta}: ussim estas proposições : é \emph{verdade} que \emph{Pedro} corre, é \emph{falso que Paulo} escreve, equivalem a cestas: \emph{Pedro corre}, \emph{Faulo} não escreve. Mas os outros quatro modos alteram a relação do predicado com o sujeito, e a proposição : « \emph{Pedro} corre » dilfere da seguinte: é \emph{possécel que Pedro corra}. São, pois. quatro as especies da proposição \emph{modal}: \emph{modal necessaria}, \emph{modal-contingente}, \emph{modal-impossivel} e \emph{modal-possivel}. Como : \emph{modal contingente} é equivalente A \emph{modal-possivel}, porque dizer : « é \emph{possivel que Pedro} escreva » é o mesmo que dizer: «é \emph{contingente que Pedro} escreva », por isso só consideramos tres especies de \emph{proposição nodal}: \emph{a modal-necessaria}, a \emph{modal-impossivel e a modal-possivel}.\par A proposição \emph{moda}? tem, como a \emph{ebsofuta}, \emph{a quantitade} e a \emph{qualidade}. \emph{Quanto é qientitade}, a proposição \emph{modal}? \emph{e universal ou particular}, segundo a diversa. indole. do \emph{modo}, Se o \emph{inodo} exprime \emph{necessidade}, \emph{a} proposição é\par \emph{porque o que}. convem \emph{necessaricinente a um sujeito}, \emph{iii cersad-a fiemativa convenrxihe} por cusa \emph{da essencia}, e por isso deve attribuir-se a todos OS \emph{individuos}, \emph{que} tem a mesma essencia: assim a proposição «é \emph{necessario que Lodo} 0 \emph{liomem seja racional} » corresponde á esta «\emph{todo o homem} é \emph{ractóael} 5», Se \emph{o modo exprime éoipossibilidide}, a proposição é \emph{wezeersal-tegativa}, \emph{porque} o \emph{que não pode} convir a um sujeito, nem pode convir a algum\pgn{56}individuo, que se acha contido na \emph{extensão} d' esse sujeito:.assim a proposição «é \emph{impossivel que o homem seja bruto}» --- corresponde á esta --- « nenhum \emph{homem é bruto}». Se 0 \emph{modo} exprime \emph{possibilidade}, \& proposição é \emph{parteular}, porque o que convem \emph{possivelmente}, ou \emph{contingentemente}, convem a. alguns \emph{individuos}; assim está \emph{proposição} --- « 6 \emph{possivel que o} homem \emph{seja douto}» --- corresponde a oulri --- « \emph{algum homem é donio}». --- Quanto á \emph{qua}- \emph{Tidade}, \emph{a modal é affirmativa} ou \emph{negativa}, conforme o \emph{modo} se \emph{afirma} ou se \emph{mega}; assim a proposição --- «não é \emph{possivel que o homem seja bom}» --- é \emph{negativa}.}. . :

\textbf{ÁrTIGO} III

\emph{Divisão da proposição composta}

%%>>> t60

\pg{54} \pg{55} \pg{56} \tese{60}{Especies da proposição! composta.}
As principaes especies da proposição \emph{composta} são \emph{cinco}: \emph{copulativa}, \emph{causal}, \emph{condicional}, \emph{disjunctiva} e \emph{a dversativa}. \textbar{}

%%>>> t61

\tese{61}{Proposição copnlativa.}
A proposição é \emph{copulativa}, quando tem \emph{diversos} sujeitos; ou \emph{diversos} predicados, unidos pelas conjunceoóes \emph{copulativas}; ex. «a virtude e a sciencia fazem o homem grande». --- Será \emph{verdadeira}, se o fôr \emph{cada} uma das proposições. \emph{simples}, em que póde resolver-se.

%%>>> t62

\tese{62}{Proposição causal.}
A proposição é \emph{causal}, quando dá a razão por que um predicado convem, ou não, a um determinado sujeito; ex. « a alma humana é immortal, por ser espirito », --- Para que esta proposição seja \emph{verdadeira}, é necessario que o seja tambem a \emph{razão} adduzida \footnote{A razão adduzida pela \emph{causal} será verdadeira, se existir nm nexo entre a mesma razão e a proposição \emph{simples}, isto é, se na proposição \emph{simples} o predicado convier ao sujeito em virtude da razão adduzida ; aliás será falsa, embora cada uma das partes seja por si verdadeira. Por isso é falsu esta proposição causal: « 3 alma humana é espiritual. porque existe », pois que a existencia da alma. não é a razão da sua espiritualidade, embora seja verdade que a nossa alma é espiritual e que existe.} -

; 63. Proposição condicional. --- A proposição é \emph{condicional}, quando resulta de varias proposições simples, unidas pelas particulas \emph{condiciomaes}; ex. « se Pedro é homem, é racional ». --- Compõe-se de duas partes: da \emph{condição}. que se chama \emph{antecedente}, e do \emph{condicionado}, que se denomina \emph{consequente}. --- É \emph{verdadeira}, quando entre a \emph{condição} e o \emph{condicionado} ha um nexo tal que, admittida a \emph{condição}, é forçoso admittir tambem o \emph{condicionado}, embora, consideradas separadamente, ambas as partes sejam falsas; por ex. «se a pedra sente, tem alma sensitiva \pg{57} ». É \emph{falsa}, quando não existe aquelle nexo, embora as duas partes, consideradas separadamente, sejam verdadeiras ; ex. «se a \textbf{planta} vegeta, o sol é brilhante ».

%%>>> t64

\tese{64}{Proposição disjunetiva.}
A proposição é \emph{disjunctiva}, quando tem diversos sujeitos, ou diversos predicados, unidos entre si pelas conjunegoes \emph{disjunctivas}, ex. «o mundo é infinito ou : finito ». --- É \emph{verdadeira}, quando um dos predicados convem ao sujeito, ou quando o predicado convem a um dos sujeitos (').

%%>>> t65

\tese{65}{Proposição adversativa.}
A proposição é \emph{adversaliva}, quando encerra varias proposições \emph{simples}, unidas pelas conjuncoes \emph{adversalivas}; ex. « Luiz é sabio, mas não é bom ». --- Será \emph{verdadeira}, se o forem todas as proposições \emph{simples} \footnote{A proposição \emph{disjuctioa}, pois, deve ter os dois seguintes \emph{requisitos} : --- 1º) as partes erzumeradas devem ser \emph{completas}: por isso é [alsa a se-\par \emph{p} guinie: « esta côr é branca ou verde »; --- 2º.) as partes não podem .ser.lo-. das verdadeiras nem todas falsas; mas devem ser taes' que a \emph{affirimacio} de uma d'ellas importe a \emph{exclusão} das outras; e a \emph{exclusão} de todas, menos uma, importe a \emph{afirmação} da que não foi excluída. Assim é verdadeira a seguinte: «este ser vivo ou é \emph{planta}, ou \emph{animal}, ou \emph{homem} »; pois, se alfirmamos que é homem, com isso negamos que seja planta ou animal, c. sc negarmos que é planta e animal, já afirmamos que é homem, --- A \emph{disjunctiva} resolve-se na \emph{condicional}. ;\par (\$) Entre as partes da proposição \emph{adversativa} deve exisUr \emph{uma certa} opposição, embora pouco sensivel, ou pelo menos uma diversidade: \emph{aliás n adversativa} reduzir-se-ba a \emph{copulativa}.\par 'Pertecem tambem ás proposições \emph{compostas a relaliee} e a \emph{exponicel}, --- A \emph{relativa} é a que resulta de varias proposições simples, unidas pelas particulas \emph{relativas} --- \emph{onde} e \emph{ahi}, \emph{tal} e \emph{qual}, \emph{então} C \emph{quando}, etc.. como: \emph{tal é o povo}, \emph{qual é o sacerdote}. Será verdadeira esta proposição, se existir realmente o nexo eatre as proposições simples. de modo que uma seja causa e determinação da outra: ilids, será falsa. Assim é falsa esta proposição \emph{então} será \emph{feliz o} homem, \emph{quando possuir as riquezas}. --- A \emph{exponivel} é \emph{n} que, embora seja \emph{simples} na apparencia, emquanto resulta de um só sujeito e dc um só predicado, comtudo é realmente \emph{composta} em força de algumir parüceula; unida ao sujeito ou ao predicado; tar é a proposição --- só' Deus é \emph{santo}, --- que se resolve nas duas seguintes --- Deus é \emph{santo} c \emph{todo outro} ser \emph{não} é santo, Será \emph{verdadeira}, se for verdadeira a. parlicula, que modifica a proposição: alias, será falsa. Assim é falsa a seguinte --- só \emph{o homen vive}. Chama-se \emph{expouicel}; porque \emph{deve ser exposta} ou \emph{explicada}, \emph{para que} se conheça a sua significação total.}.

%%>>> art-3

\section*{Artigo IV.}
\addcontentsline{toc}{section}{Artigo IV.}

\pg{58} \emph{Qualidades das proposições}

%%>>> t66

\tese{66}{Numero das qualidades.}
As \emph{qualidades} das proposições são tres: \emph{opposição}, conversão e \emph{equivalencia}, (\textasciicircum{}).

%%>>> t67

\tese{67}{Opposição.}
\emph{Opposiczo} é a \emph{repugnancia de duas proposições}, que \emph{têm o} mesmo \emph{sujeito} e \emph{o} mesmo \emph{predicado}, --- A \emph{repugnancia} está em ser uma \emph{affirmativa}, e a outra \emph{negativa}: ex. «a neve é branca» --- «a neve não é branca » (\textasciicircum{}).

%%>>> t68

\tese{68}{Divisão das proposições oppostas.}
Com relação á sua \emph{extensão}, as proposições \emph{opposías} podem ser ou ambas \emph{universaes}, ou "ambas \emph{particulares}, ou uma universal e outra \emph{particular}. --- Na primeira hypothese, \textbf{chamam-se} \emph{conírarias}, ex., « todos os homens são justos » --- « \textbf{todos os} homens não são justos » ; na segunda, dizem-se \emph{sub-contrarias}, ex. « alguns homens são justos » --- « alguns homens não são justos » ; na terceira, são \emph{contradictorias}, ex. «todos os homens são justos »

---:« alguns homens não são justos » \footnote{Esta ires \emph{qualidades} convém ás proposições, --- não emquanto estas se consideram em sé \emph{nesmas} (pois neste sentido são \emph{afirmativas} ou \emph{negativas}, universaes ou \emph{particulares}, ele.), --- mas em quanto se consideram na relação, que umas têm com as outras.}.

Eis o schema das proposições \emph{oppostas}, advertindo que a lettra A designa a proposição \emph{universal-affirmativa}, E a \emph{wniversalnegativa}, la \emph{particular-affirmativa}, O a \emph{particular-negativa}. Todos os homens CONTRARIAS Todos os homens

são \emph{justos e} " i e não são justos

“\emph{Oo} 7 " D»

"\$9 ES

\textasciicircum{}4. aS

2 9

\emph{SUO}.

"a \emph{Q}

; : * \textasciicircum{}4 Alguns homens Oo ; \emph{Ss} Alguns homens

são justos não são justos

SUB-CONTRARIAS\footnote{Duas proposições serão, pois. \emph{oppostas} --- 1º.) se tiverem o mesmo \emph{sujeito} € O mesmo \emph{predicado}, --- 9\textasciicircum{}.) se uma \emph{aflirmar} 0 que a. outra \emph{nega}. --- Se não se verifiearem estas duas condições, não poderá haver opposição.}\footnote{\textsc{Duas} proposições \emph{singulares} são \emph{contrasdictorias}. quando differem \emph{na qualidade}; ex. « Antonio é subio », « Antonio não é sabio ».\par Alguns incluem no numero das proposições \emph{oppostas} as \emph{subalternas}, que\pgn{59}convém entre si na \emph{qualidade}, emquanto são \emph{ambas affirmativas} ou \emph{anas negativas}, mas differem na \emph{quantidade}, emquanto uma é \emph{universal} e outra € \emph{particular}, como as seguintes: \emph{todos os homens são justos} --- \emph{alguns homens são justos}. --- Mas as \emph{subalternas} não são propriamente \emph{oppostas}; pois a \emph{opposição} dá-se quando uma proposição \emph{affirma} o que a outra \emph{rega}.}

%%>>> t69

\pg{59} \tese{69}{Critica das proposições opposías.}

a) A \emph{contradictorias} não podem ser ambas verdadeivas, nem ambas falsas; porque uma coisa não pode ser e deixar de ser ao mesmo tempo e sob o mesmo aspecto (\emph{principio de contradiego}) --- Por isso da falsidade de uma conclue-se a verdade da outra; e da verdade de uma, a falsidade da outra \footnote{As contradictorias : '\par a) Não podem ser \emph{ambas verdadeiras}, seja. qual for a sux materia. Porquanto, «\emph{firmando} uma o que z outra \emph{nega}, se ambas fossem verdadeiras; uma. coisa seria e deixaria de ser ao mesmo tempo e sob o mesmo aspecto. Assim, quando digo --- \emph{todos os homens são justos}, --- alfirmo que a justiça existe em \emph{todos} os homens; mas quando digo --- \emph{alguns homens mio são justos}, --- \emph{nego} que a justiça existe em \emph{todos} oS homens, Por isso, se lossem verdadeiras estas duas \emph{coniradictorigs}, 2. justiça existiria e não existiria em \emph{todos} os homens: o que é absurdo.\par 5) Não podem ser \emph{ambas falsas}. Porquanto, se é falsa a \emph{affirmatica}, é falsa, porque o que se affirm, \emph{não} é; e sc é lalsa \emph{i negativa}, é falsa porque o que se nega. é. Ora, como repugna que uma coisa, ao mesmo tempo e sob o mesmo aspecto, \emph{seja} e \emph{não seja}, \emph{ás contradictorias} não podem ser ambas \emph{falsas}. Assim, se é falsa esta proposição --- \emph{todos os humens} são \emph{justos}, é falsa porque a justiça não se \emph{encontra} em \emph{todos} os homens; e se é falsa esta outra --- alguns \emph{homens não são justos}, --- é falsa, porque a justiça se \emph{encontra} em \emph{todos} os homens: e por isso, se ambas são \emph{falsas}, a justiça \emph{não se encontra} e se encontra em \emph{todos} os homens: o que é absurdo.\par A razão é, porque entre uma e outra não ba \emph{meio termo}, como não ha meio Lermo entre a. altirmação e a negação. Urna coisa é, ou não é. --- Repetimos que as proposições \emph{singulares} São \emph{contradictorias}, quando differem entre si na \emph{qualidade}; porque, para se exeluirem reciprocamente, basta que uma seja affirmativa, e outra negativa. Assim são \emph{contradictorias} us seguintes --- \emph{Pedro é bom}, \emph{Pedro não} é bom.}.

b) As \emph{contrarias} --- não podem ser ambas verdadeiras, seja qual fôr a materia, de que tratam, --- não podem ser ambas falsas, \emph{quando a materia} é necessaria, --- podem ser ambas falsas, \emph{quando} a \emph{materia é contingente} \footnote{AS \emph{contrarias}:\par «) Não podem ser \emph{ambus verdadeiras}, seja. qual for a sua materia. Na verdade, uma \emph{affirma} 0 que a outr \emph{nega}, e por isso, se ambas fossem \emph{verdadeiras}, a mesma coisa seria e deixaria de ser a0 mesmo tempo e sob o mesmo aspecto.\par 5) Não podem ser ambas \emph{falsas}, quando a materia é \emph{necessaria}, Com efleito, neste caso, pelo facto de \emph{o} predicado convir \emph{essencialinernte} ao sujeito,\pgn{60}a \emph{contraria affinmativa} é verdadeira, pois o que convem \emph{essencialmente} a um sujeito, não pode deixar de convir lambem a todos os seres, que estão comprehendidos na extensão d'esse sujeito; e por isso a \emph{contraria negativa é falsa}; ex. \emph{todo o homem é racional} --- \emph{nenhum homem é racional}.\par '\emph{e}) Podem ser ambas \emph{falsas}, quando a materia é \emph{contingente}. Porquanto, neste caso, \emph{O} predicado, por não convir \emph{essencialmente} ao sujeito, não se encontra em \emph{todos} os individuos, comprehendidos no: sujeito, nem falta em \emph{todos}; e tanto pode ser falsa a proposição \emph{contraria}, que attribue a \emph{todos} o predicado, como a que o não attribue á \emph{ninguem}; como nestes exemplos: \emph{todo} o \emph{homem é bom} --- \emph{nenhum homem é bom}. --- Entre as proposições \emph{contradictorias e ás conírarias} dá-se esta diflerenca, que nas \emph{contradictorias} a negação de uma destroe a \emph{affirmação} da outra, e entre ellas não ha meio termo: mas. nas \emph{conérariías}, \emph{a negação} de uma não só destroe a \emph{affirmação} da outra, ma chera a outro extrémo, e por isso, existindo entre extremos oppostos um meio termo, entre as \emph{contrarias} ha um meio termo, que é a proposição \emph{particular}, \emph{ejffirnudiva} ou \emph{negativa}. (S. Thom. in I \emph{Períhe}., l. 11).}.

Portanto --- \emph{a}) da verdade de uma podemos deduzir a tal-. sidade da outra, seja qual for a materia, de que se trata; --- \emph{b}) da falsidade de uma deduzimos a verdade da outra, quando a materia é \emph{necessaria}; --- \emph{c}) da falsidade de uma não se pode deduzir a verdade da outra, quando a materia é \emph{contingente}.

€) As \emph{sub-contrarias} --- não podem ser ambas falsas, seja qual fôr a \emph{materia}, de que tratam, --- não podem ser ambas verdadeiras, quando a materia é \emph{necessaria}, --- podem ser ambas verdadeiras, quando a materia é \emph{contingente} \footnote{As \emph{sub-contrarias} :\par a) Não podem ser \emph{ambas falsas}, seja qual for a materia, de que tralam ; aliás as suas \emph{contradictorias}, que são as contrarias, seriam ambas verdadeiras.\par 5\emph{b}) Não podeni ser \emph{ambas verdadeiras}, quando a materia é \emph{necessaria}. . Porque, então, \emph{o} predicado convem a £\emph{odos} e a cada um dos individuos, e por isso a \emph{sub-contraria affirmativa} é verdadeira. e a \emph{negativa} € falsa ; como: \emph{alguns homens são racionaes}: \emph{alguns homens não} são \emph{racionaes}.\par \emph{c}) Podem ser ambas \emph{verdadeiras}, quando a matevia é \emph{contingente}. Porque, então, o predicado não existe em todos os \emph{individuos}, nem falta a \emph{todos}, e as duas \emph{sub-contrarias} não são propriamente \emph{oppostas}, pois não têm o mesmo sujeito, como: \emph{alguns homens são bons}, «\emph{Iguus homens} não \emph{são} bons. (Cf. S, Thomaz, In 1 \emph{Perih}., 1. 10).}.

Por isso --- \emph{a}) da falsidade de uma deduzimos a verdade da outra, seja qual fôr a materia, de que tratam; --- 5) da verdade de uma podemos deduzir a falsidade da outra, quando a materia é \emph{necessaria}; --- \emph{c}) da verdade de uma não podemos deduzir a falsidade da outra, quando a materia é \emph{contingente}.

%%>>> t70

\tese{70}{Conversão.}
4 \emph{conversito} é a \emph{transformação de} uma \emph{pro}- , \emph{posição} noutra. --- Faz-ze pela troca dos termos, passando o sujeito \pg{60} \pg{61}  para predicado e este para sujeito, sem alterar a \emph{qualidade} e a \emph{verdade} das proposições: assim a proposição « algum homem é sabio» converte-se nesta: «algum sabio é homem ».

%%>>> t71

\tese{71}{Divisão da conversão.}
A conversão pode ser \emph{simples}, \emph{accidental} ou \emph{contraposta}.

a) É \emph{simples}, quando permanece a mesma \emph{quantidade}, --- Convertem-se d'esta forma --- 4) as proposições \emph{tenfversaesnegalivas}; ex. «nenhum homem é planta» --- « nenhuma planta: é homem»; --- 5\emph{b}) e as \emph{particulares-affirmatievas} ; ex. « algum homem é sabio » --- «algum sabio é homem » ().

b) É \emph{accidental} (\emph{per accidens}) quando se altera a \emph{quantidade}. --- Convertem-se assim --- \emph{a}) as \emph{universaes-affirmativas} ; ex. «todo o homem é racional» --- «algum racional é homem »; --- b) as \emph{universaes-negalivas} ;- ex. » nenhuma planta é mineral » --- » algum mineral não é planta » \footnote{A' conversão \emph{simples} estão sujeitas --- as \emph{UNniversaes-negativas}, porque na universal \emph{negativa} o \emph{sujeito} e O \emph{predicado} tomam-se. em toda a sua \emph{extensão}, \& NA conversão conservam a sua universalidade, --- \& ás \emph{particulares-affirmativas}. porque, nestas proposições, o sujeito e o predicado são lermos \emph{particulares}, e ui conversão conservam a sua \emph{porticutaridade},} !

c) É \emph{coniraposia}, quando se conserva.a mesma. \emph{quantidade}, mas os termos, por meio da particula \emph{negativa}, se tornam \emph{indefinidos}. --- Convertem-se por esta forma --- \emph{a}) as \emph{UNtversaesaffirmativas}; ex. «todo o animal é ente sensitivo » --- « todo o ente não sensitivo é não animal », --- b) e as \emph{particulares-negativas}; ex. «algum corpo não é ente vivo» --- « algum ente não vivo não é não corpo» \footnote{4” conversão \emph{accidental} estão sujeitas --- as \emph{unseersaes-uffivinatívas}, pois nestas proposições o \emph{predicado é partientar}, \emph{e} por isso.” quando este, \emph{na conversão}, se torna \emph{sujeito}. deve ser \emph{particular}, --- e as universaes-\emph{megafivas}, porque vestas proposições o predicado é \emph{universal}. e por isso allri- Due-se a \emph{cada um} dos \emph{entes}, que estão comprehendidos sob a sua \emph{extensão}.}.

%%>>> t72

\tese{72}{Equivalencia.}
\emph{Equivalencia} é a reducgao de duas proposigóes \emph{oppostas} a mesma significação. --- Esta reducção faz-se por meio da particula \emph{negativa} «não» \footnote{A conversão \emph{contraposta} nunca ou rarissimas vezes se emprega.\par (A particula \emph{negativa} «\emph{não}» tem uma \emph{natureza maligna}, pois destroe tudo aquillo, a que é anteposla, isto é, \emph{da} a proposição um sentido opposto ào que linha. Assim, se 2 particula \emph{wegatioa} é anteposta a \emph{quatédade}, a. proposição de \emph{wnuicersal} se Lorna \emph{particular}, e vice-versa, e, Se é anteposia Á \emph{qualidade}, á proposição \emph{de affirmativa} se torna \emph{negativa}, e vice-versa.}.

%%>>> t73

\pg{62} \tese{73}{Regras da eqnivalencia.}

\textsc{ReGRA} 1º --- \textsc{Duas} \emph{contradictorias} tornam-se \emph{equivalentes}, antepondo a particula \emph{negativa} ao sujeito de uma d'ellas; porque, assim, a negação allera a \emph{quantidade} e a \emph{qualidade} da proposição. Às duas \emph{contradictorias} «todos os homens são sabios» e «alguns homens não são sabios» tornam-se \emph{equivalentes}, antepondo a particula \emph{negativa} ao sujeito da primeira, «nem todos os homens são sabios» que equivale a «alguns homens não são sabios».

\textsc{ReGRA} 2º --- \textsc{Duas} \emph{contrarias} tornam-se \emph{equivalentes}, pospondo a \emph{negação} ao sujeito d'uma d'ellas; assim a particula \emph{negativa} altera só a \emph{qualidade}. A proposição «todos os homens são eloquentes », pospondo «nio» ao sujeito «homens», torna-se equivalente a sua \emph{contraria}, que é «todos os homens -

\emph{Hn} não são eloquentes ».

. As sub- \emph{contrar its} não admittem \emph{equivalencia O}

%%>>> cap-4

\chapter*{CAPITULO TERCEIRO\\[0.3em]\large RACIOCINIO E ARGUMENTAÇÃO}
\addcontentsline{toc}{chapter}{CAPITULO TERCEIRO --- RACIOCINIO E ARGUMENTAÇÃO}

\textsc{SuUMMARIO}: --- Raciocinio, seus elementos e divisão. --- Argumentação e

syllogismo. --- Syllogismo simples e suas regras. --- Figuras e modos

do syllogismo simples; --- Syllogismo composto e sua divisão, --- 1n-

ducção e suas especies. --- Sophisma e suas especies.

%%>>> art-4

\section*{Artigo I.}
\addcontentsline{toc}{section}{Artigo I.}

\emph{Raciocinio}, \emph{seu elementos} e \emph{divisao}

%%>>> t74

\tese{74}{Raciocinio..}
\emph{Raciocinio} É a \emph{união} ou \emph{separação de duas ideas pelo facto de ambas convirem}, ou \emph{uma convir} e \emph{outra} 7\emph{ui}0, \emph{com} uma \emph{terceira idea}, com a \emph{qual foram comparadas}; ex.

* "Todo, o espirito é immortal ;

A alma. humana é espirito :

Logo a alma humana é immortal;

Comparando entre si duas ideas, muitas vezes não podemos perceber 4\emph{mmediatamente} a sua conveniencia ou desconveniencia\footnote{As \emph{sub-contrarias} não udwiltem \emph{eguivalencia}; porque. seja qual for o posto da partícula \emph{negativa}, uma proposição \emph{sub-contraria}. não pode \emph{reduzir-se a significação} da outra.} \pg{63} . Em tal caso recorremos a uma terceira idea, para com ella comparar as duas primeiras. D'esta comparação: resulta que --- ou ambas as ideas convêm com a terceira, --- ou ambas não convém, ------ ou uma convem e outra .não. Na \emph{primeira} hypothese, conclue-se a conveniencia das duas ideas; --- na M: , ; \emph{segunda}, não póde concluir-se a sua convenjencia nem a sua desconveniencia; --- na \emph{terceira}, conclue-se a sua desconveniencia. Assim, não conhecendo immediatamente a relação entre estas duas ideas «\emph{alma humana}» e «\emph{immortal} », comparamol-as com uma terceira «\emph{espirito}», e, pela sua conveniencia com esta, concluimos que ellas convém entre si \footnote{A perfeição das substancias intellectuaes consiste no conhecimento da verdade. Mas ha uma differença neste conhecimento, e a dilferença funda-se na propria natureza d essas substancias. As intelligencias superiores, que são os AÀnjos, conhecem sempre a verdade cum um só olhar, sem movimento ou discurso. As intelligencius inferiores, como é a do homem, não chegam sempre ao perfeito conhecimento da verdade com um só olhar, mas devem fazer um movimento, um discurso, passando do conhecido para o desconhecido, e assim descobrem e conhecem explicitimente muitas ver- 'dades, que estavam implicilamerte contidas em determinados principios. É, pois. necessario que a nossa intelligencia estabeleça os principios, e depois os analyse e resolva, para que se veja o que nelles se achava envolvido. D'esta analyse ou resolução derivam diversos juizos, dispostos entre si por modo, que um deriva dos outros. Esta deducgüo de um juizo de outros \emph{juizos} é propriamente o \emph{raciocinio} (\emph{Gf}. \emph{Sum}. \emph{TA}. p. l. q. 85. a. 5).\par Erram. portanto, --- \emph{Wundi}, \emph{quando} dá o nome de \emph{raciocinio} a ioda e qualquer forma de conbecimento (mesmo á simples sensação), --- e \emph{Locke}, \emph{quando} diz que o \emph{raciocinio} é uma simples acção \emph{intuitiva}. Pelo \emph{raciocinio}, a intelligencia não vê uma verdade ma outra, mas \emph{passa de uma verdade para outre}. Esta passagem realisa-se, quando a intelligencia não vê immediatamente a relação entre duas diversas ideas.}.

%%>>> t75

\tese{75}{Elementos do raciocinio.}
Os elementos do \emph{raciocinio} são dois, a saber: \emph{materia} e \emph{forma}.

aj A \emph{materia} é constituida pelas \emph{tres ideas}, que são a materia \emph{remota}, e pelo \emph{tres juizos}, que são a materia \emph{proxima},

Das tres ideas uma é \emph{sujeito} da conclusão, outra é \emph{predicado} da mesma conclusão, a terceira é o termo de \emph{comparação}. --- O sujeito e o predicado chamam-se \emph{extremos}; o sujeito é o \emph{extremo menor}, e o predicado é o \emph{extremo maior}; a idea, que é o termo da \emph{comparação}, diz-se \emph{medio} \footnote{O \emph{sujeito} diz-se extremo \emph{menor}, e o \emph{predicado} devomina-se extremo \emph{maior}; porque o predicado, referindo-se ao \emph{sujeito} como \emph{o principio formal} e \emph{determinante} ao \emph{principio material} e \emph{determinavel}, tem uma extensão \emph{maior}\pgn{64}que a do sujeito, --- O Lermo de comparação diz-se \emph{medio}; não só porque é o \emph{meio} pelo qual julgamos da \emph{conteniencia}, ou da \emph{desconveniencia}, do sujeito com o predicado, mas tambem porque a sua \emph{extensão} é. geralmente, \emph{media}" entre a extensão do sujeito e a do predicado.}.

As tres ideas comparadas entre si constituem \emph{tres jitizos}, --- um formado pela comparação do \emph{predicado} e da idea \emph{media}, e chama-se \emph{premissa maior}; --- outro, pela comparação do \emph{sujeito} e da idea \emph{media}, e chama-se \emph{premissa menor}; --- o terceiro, pela comparação do \emph{sujeito} e do \emph{predicado}, e diz-se \emph{conclusão}. --- Às \emph{premissas} formam o \emph{antecedente}; a \emph{conclusão} é o \emph{consequente}.

b) A \emph{forma} do \emph{raciocinio} é a disposição das \emph{ideas} e dos \emph{juizos}, pela qual, admittido o \emph{antecedente}, deve \emph{necessariamente} admittir-se o \emph{consequente}. --- O \emph{nexo}, que em virtude da \emph{forma} existe entre o \emph{antecedente} e o \emph{consequente}, chama-se \emph{consegquencia}. --- Não são, pois, uma e a. mesma coisa o \emph{consequente} e a \emph{consequencia} : e raciocinios ha, em que a conseguencia é verdadeira e o consequente é falso, --- ou em que a consequencia é falsa e o consequente é verdadeiro, como nestes dois exemplos:

' (\emph{consequencia verdadeira} --- \textasciicircum{} Todo o homem é racional;

\emph{e consequente falso}) Ora a planta é homem;

\textasciicircum{} UT ' Logo a planta é racional;

! XxX

(consequenciu. \emph{falsa} Todo o animal é sensitivo;

\emph{e consequente verdadeiro}) Ora o cão é animal;

Logo o homem é livre \footnote{O nexo entre o antecedente e O \emph{consequente} \& tão estreilo, «que, admittindo-se aquelle, deve admittir-se tambem este. Com etffeito, quando julsamos ser verdadeiro um principio universal, \emph{implicitamente} julyamos: ser lambem verdadeira uma conclusão, contida naquelle principio. Quando, pois, se deduz a conclusão, deve ella \emph{explicitamente} allicmar ou negar o que \emph{implicitamente} foi aflirmado ou .vegado no principio; aliás uma e a mesma coisa seria e deixaria de ser 20 mesmo tempo e sol: o mesmo aspecto. --- \emph{O juizo} ácerca da verdade da \emph{consequencia} pertence a Logica: mas o juizo \emph{acerca da} verdade do \emph{conseguente} pode pertencer a outras sciencias.}.

%%>>> t76

\pg{64} \tese{76}{Divisão do raciocinio.}
O raciocinio é \emph{affiinetivo}, ou \emph{negativo}. --- É \emph{affirmativo} se atfirma a \emph{conveniencia} do sujeito com o predicado, porque ambos convém com a \emph{idea media}. --- É \emph{negativo}, se affirma a \emph{desconveniencia} do sujeito com o predicado, porque um d'elles não convem com a \emph{idea media} \footnote{Já demos exemples do syllogismo \emph{afirmativo}. Um exemplo do syllogismo \emph{negativo} é o seguinte : \emph{todo o homem é racional}; \emph{o cão} não é \emph{konrem}; \emph{logo o} cão \emph{não é racional}. A conclusio diz que o predicado \emph{racional} não \emph{compete} au \emph{sujeito cão}, porque este vão convem com a \emph{ideu media}, que é \emph{homen}.}.

\emph{GAS} aims rn \emph{rnm} arm aims a

%%>>> cap-5

\chapter*{CAPITULO TERCEIRO\\[0.3em]\large RACIOCINIO E ARGUMENTAÇÃO --- 65}
\addcontentsline{toc}{chapter}{CAPITULO TERCEIRO --- RACIOCINIO E ARGUMENTAÇÃO --- 65}

%%>>> t77

\pg{65} \tese{77}{Principios do raciocinio affirmativo e negativo.}
O raciocinio «\emph{ffirmativo} lunda-se no principio de comveniencia, que é: \emph{Duas coisas}, \emph{cada uma das quaes é equal} a ema \emph{terceira}, \emph{são eguaes entre se}. --- O raciocinio \emph{negativo} funda-se no principio de \emph{desconveniencia}, quo é: \emph{Duas} coisas, \emph{das quaes} uma só é \emph{egual a} uma \emph{terceira}, \emph{nio são eguaes entre} sé \footnote{Advirta-se que duas coisas, egrüaes a um: Lerceira, são eguaes entre sb \emph{naquele sentido e naquella}. \emph{medida}, \emph{em que} são eguaes a essa \emph{terceira coisa}. Esta advertencia é da maxima. importancia, Porque não podemos dizer que são eguaes entre si duas coisas, uma das quaes é egual á \emph{terceira}, \emph{considerada}. \emph{sob} um \emph{aspecto}, e a outra é egual á mesm: Lerceira coisa, mas \emph{considerado}. \emph{sob} outro \emph{aspecto}; porque nesse euso0 o \emph{meio} termo não é um só, mas mulliplo, e o raciocinio é um sophisma. Assim é falso este raciocinio --- \emph{alguns homens são bons} --- \emph{alguns homenus são maus} --- \emph{logo alguns} bons \emph{São} naus, --- porque a terceira coisa --- \emph{alguns fiomens} --- não é considerada sempre sob um e o mesmo aspecto, mas na premissa \emph{maior} sob um aspecto e na menor sob outro e differente aspecto, pois os homens que são bons são diversos dos que são' maus, e assim a terceira coisa é duplice, e portanto não ha comparação das duas coisas com uma terceira, e falta a base ao raciocinio. ---- Advirta-se lambem que os dois principios de \emph{conveniencia}. (identidade) e de \emph{desconveniencia}. ('\emph{iscrepancia}) sê fundam e se resolvem no principio de \emph{contradicção} -- \emph{na}, \emph{coisa} não \emph{pode simultaneamente} ser e \emph{deixar de ser} (\emph{ident non potest simul} esse \emph{et non esse}). Aquelles dois \emph{principios} são verdadeiros, porque verdadeiro é o principio de \emph{contradicção}. Por isso este é principio \emph{remoto} do raciocinio, e os outros dois são principios \emph{pro.cimos}.}. .

%%>>> art-5

\section*{Artigo 11.}
\addcontentsline{toc}{section}{Artigo 11.}

\emph{Argumentação e syllogismo}

%%>>> t78

\tese{78}{Argumentação.}
\emph{Argumentação} é a \emph{expressão verbal do raciocinio}. --- Por isso, ludo quanto dissermos acerca da \emph{nalurefi} da \emph{argumentação}, das suas \emph{propriedades} e \emph{leis}, indo deve attribuir-se ao \emph{raciocinio}.

%%>>> t79

\tese{79}{Especies de argumentação.}
Às especies de \emph{argumentação} são duas: --- uma, em que de duas proposições \emph{descemos} para uma Lerceira, como do \emph{lodo} para a \emph{parie}: --- outra, em que de duas proposições \emph{subimos} para uma terceira, como das \emph{partes} para o \emph{todo}. --- A primeira especie é a argumentação \emph{deductiva}, chamada \emph{syllogismo} , a segunda é a \emph{zndauctiva}, e diz-se \emph{inducção}. --- Agora tratamos do \emph{syllogismo}; mais tarde, da \emph{tnduceão} (C).\footnote{O \emph{todo}, de que se falla aqui nà argumentação, --- não é o \emph{integral}, composto de partes \emph{integrantes}, como é, por ex., uma casa composta das \emph{paredes}, do \emph{tecto}, e do pavimento, --- nem é o \emph{potencial}, composto de partes\par Vol. 1. \emph{G} : \emph{E} A \textbar{} B\pgn{66}\emph{potenciaes}, como é, por ex., a nossa alma composta das faculdades de \emph{entender}. querer, etc., --- mas é o \emph{todo universal}, composto de partes \emph{subjectivas}, como é, por ex. o animal, cujas partes \emph{subjectivas} são \emph{o} animal \emph{racional} e o animal \emph{irracional}. Na verdade. atiribuindo-se, na argumentação, a uma \emph{parte} 0 predicado, que convem ao \emph{todo}, é necessario que essa \emph{parte} contenha em si o \emph{todo}. Ora isto não se verifica relativamente ao lodo \emph{integral}, porque uma parte \emph{integrante} não é o todo \emph{integral} (assim o \emph{fecto} não é a casa), --- nem ao todo \emph{potencial}, pois uma parte \emph{potencial} não equivale ao todo \emph{potencial} (assim o \emph{entendimento} não é toda a \emph{alma}) --- mas verifica-se relativamente ao todo \emph{universal}, pois que as partes \emph{subjectivas} contem o todo \emph{uni versal} (assim dizemos que o \emph{homem é animal}).}

3700

\emph{GA} --- SECCROFPRIMEIRÀ 66 \textbar{} OG

%%>>> t80

\pg{66} \tese{80}{Syllogismo.}
\emph{Syllogismo} é a argumentação, pela qual se conclue que o predicado de um sujeito \emph{universal} convem a um sujeito \emph{particular}. pelo facto d' este se achar contido naquelle sujeito «\emph{universal}, Assim no seguinte syllogismo :

Todo o ser creado é limitado ;

Ora \emph{o} homem é um ser creando ;

Loro o homem é limitado --- o predicado «\emph{limitado}» convem ao sujeito particular «\emph{homem} » por este se achar contido no sujeito universal « ser \emph{creado} ».

SL. Elementos do syllogismo. --- O syllogismo consta de dois elementos, que são a \emph{materia} e a \emph{forma}, --- A \emph{materia} é constituida por \emph{tres termos} (materia \emph{remota}) e por \emph{tres proposições} (materia \emph{proxima}). \emph{A forma} é a conveniente disposição dos termos e

Án das proposições em ordem á \emph{conclusão}. --- Os \emph{termos} e as \emph{proposições} tomam os nomes, que demos ás ires \emph{ideas} e aos tres \emph{juizos}, que constituem a \emph{materia} do \emph{raciocinio}. --- A proposição \emph{maior} indica que ao sujeito \emph{universal} compete \emph{necessariamente} um certo predicado; : menor, que o sujeito \emph{particular} está contido na. extensão do sujeito «:\emph{iversal}; a \emph{conclusão}, que ao sujeito \emph{particular} deve altribuir-se'o predicado do sujeito \emph{universal}, --- Por isso, a verdade \emph{particular}, expressa pela \emph{conclhusão}, recebe toda a sua força da verdade \emph{necessaria} e \emph{universal}, significada pela \emph{premissa maior} ().\footnote{A força, de que é dotada a \emph{conclusão}, não depende propriamente da \emph{universalidade} da premissa \emph{maior}, mas sim da sua \emph{necessidade}; á \emph{universalidade} é uma consequencia da \emph{necessidade}. --- A proposição necessaria, que é a base do syllogismo, pode ser necessaria \emph{absolutamente} ou \emph{hypoteticamente}, conforme a necessidade é independente, ou não, de toda e qualquer condição. No primeiro vaso, o \emph{predicado} da \emph{conclusão} exprime a essencia (\emph{Lotal ou parcial}) do termo \emph{nwedío}, ou uma \emph{propriedade essencial}, e então o nexo entre o \emph{predicado} da \emph{conclusão} \& 0 termo \emph{inedio é absolutamente} necessario\pgn{67}. No segundo caso, o \emph{predicado} da \emph{conclusão} exprime uma. propriedade, aMiribuida ao termo \emph{medio} pela \emph{experiencia}, e então \emph{o} nexo entre o \emph{predicado} da \emph{conclusão} € o Lermo \emph{medo} é necessario \emph{hypotheticamente}, isto €, suppostas certas condições. Mas, em ambos os casos, a propriedade, que \emph{necessariamente} convem ao termo \emph{medio}, attribue-se a todos os sujeitos \emph{particulares}, contidos na extensão do termo \emph{sedio}, que é o sujeito \emph{universal}.\par Dependendo do termo \emph{medio} o resultado da comparação entre o sujeito e o predicado, emquanto o termo \emph{medio} é 3. \emph{razão} por que um predicado convem ou repugoa a um sujeito, vé-se claramente i necessidade de encontrar esse \emph{medio} para se poder formar o syllogismo. Os antigos philosophos deram muitas regras para a invenção do \emph{medio}, Limitamo nos a apresentar duas d'essas regras; --- uma serve, quando a conclusão é uma proposição \emph{universal-affirmativa}, e por isso vale tambem quando a conclusão é uma \emph{particular-afirmativa} ; --- a outra serve, quando a conclusão é uma proposição \emph{universal-negativa}, e por isso vale tambem para a \emph{particular} negativa. --- 1.º) Quando a conclusão é uma proposição \emph{universal-affirmativa}, procure-se um \emph{medio}, que \emph{esteja}. \emph{contido no predicado}, \emph{mas contenha o sujeito}. Com effeito, no syllogismo, cuja conclusão é uma proposição \emph{unzversal-affirmativa}, o \emph{medio É sujeito} do \emph{predicado} na \emph{maior} € é \emph{predicado} do \emph{sujeito} na menor. Mas o \emph{medio}, --- se é \emph{sujeito} do \emph{predicado}, deve ser contido no proprio predicado, porque o predicado é mais universal que o sujeito, --- e, se é \emph{predicado do sujeito}, deve conter o proprio sujeito, pela mesua razão. Assim, querendo concluir o syllogismo por esta proposição \emph{universal-affirmativa} : \emph{logo toda} a. \emph{flor é vivente}, pergunto 2 mim mesmo: porque toda \& \emph{flor} deve ser vivente? E percebo que toda a flor deve ser vivente, porque é \emph{planta}. Mas É \emph{planta}, para que possa ser \emph{medio}, deve achar-se contida no predicado \emph{vivente}, e deve, pela sua vez, conter o sujeito \emph{toda} «a \emph{flor}. É realmente a planta está contida no \emph{vivente} c contem \emph{toda a flor}. Logo a planta € 0 \emph{medio}. Para maior clareza completamos o syllogismo : \emph{toda} a \emph{planta é vivente}; \emph{ora toda} a \emph{flor é planta}: \emph{logo tod}« a \emph{flor é vivente}. --- 2.º) Quando a conclusão € uma proposição \emph{universal-negativa}, procure-se \emph{um inedio} --- ou que \emph{esteja excluido da extensão do predicado}, mas \emph{inclua} 0 \emph{sujeito}, --- OU \emph{que esteja incluido} na \emph{extensão do predicado}, mas \emph{exclua o sujeito}. A razão d'esta regra é manifesta. No syllogismo, que tem por conclusão uma proposição \emph{universal-negativa}, uma das premissas deve ser e é \emph{negativa}; e é negativa, porque o \emph{merlio} --- ou repugna ao predicado, e convem ao sujeito, --- ou convem : 40 predicado, mas, repugna zo sujeito. --- De resto, na invenção do \emph{medio}, : Oo uso e a reflexão valem mais do que as regras.}

. SyHlogismo abreviado. --- A's vezes, por. brevidade, oa para maior força da argumentação, omilte- -se uma das premis-

sas, que facilmente se subentende. Neste caso o syllogismo diz-se \emph{abreviado}, ou, simplesmente, \emph{enthymeme},- ex.:

Todo o ser creado é limitado,

Logo o homem é limitado \footnote{Para se conhecer qual das duas premissas foi \emph{ométzda}, se a \emph{maior} OU a \emph{menor}, é necessario examinar a que ficou. Se a premissa «ue ficou contém o \emph{extremo menor}, então foi omitlida a \emph{premissa maior}, que devia\pgn{68}exprimir a relação entre o \emph{exiremo maior} e o0 termo \emph{medio}. Se a premissa que ficou contém o \emph{extremo maior} (\emph{como} no exemplo citado), então foi omittida a \emph{premissa menor}, que devia exprimir a relação entre o \emph{extremo} menor (sujeito) e o termo \emph{medio}. ]}.

%%>>> t89

\pg{67} \pg{68} \tese{89}{Principios do syllogismo.}
O \emph{syllogismo} funda-se em dois principios, que, são uma derivação dos principios de conveniencia e \emph{desconveniencia} e que se enunciam d'este modo: --- a) \emph{Tudo o que se} «\emph{ffirina de um sujeito universal}, \emph{deve affirmar-se de cada um dos particulares}, \emph{nelle contidos} (para o syllogismo \emph{affirmativo}): --- b) \emph{Tudo} o \emph{que se nega de um sujeito universal}, \emph{deve negar-se de cada um dos particulares}, \emph{nelle contidos} (para o syllogismo \emph{negativo}) \footnote{Evidente é o principio do syllogismo «\emph{afirmativo}. Por quanto, se não podesse attribuir-se a um sojeito particular um predicado, attribuido já a um sujeito universal. sob o qual o particular está contido, nm e o mesmo predicado, ao mesmo tempo e sob o mesmo aspecto, seria e não seria attribuido ao mesmo sujeito:particular; --- seria attribuido; porque. quando é attribuido ao sujeito universal, uttribue-se \emph{implicitamente} a todos os sujeitos particulares, contidos na extensão do sujeito universal, e por isso áquelle sujeito particular; --- não seria altribuido, por hypothese. --- Não menos evidente é o principio do syllogismo \emph{negativo}. Com effeito, quando negamos \& um sujeito universal um determinado predicado, \emph{implicitamente} negamos o mesmo predicado a lodos os particulares, que estão contidos nesse sujeito universal; e por isso não podemos, sem contradicção, deixar de negar \emph{explicitamente} tal predicado a algum d'esses particulares, quando nos occupamos \emph{directamente} do tal particular. --- O principio. que regula o syllogismo \emph{affirmativo}, foi chamado pelos escholasticos \emph{dicun de omni}, eo principio, em que se funda o syllogismo \emph{negativo}, foi denominado \emph{dictan de nullo}.\par Este dois principios. que regulam o syllogismo. «\emph{ffirmativo} e \emph{negativo}, e que, como dissemos, são uma derivação dos principios de \emph{conveniencia} e de. \emph{desconveniencia} (ou de \emph{identidade} e de \emph{discrepancia}), são chamados \emph{principios logicos}, porque dirigem as operações intellectuaes: emquanto os de \emph{conveniencia} e de \emph{desconveniencia}. são chamados \emph{principios metaphusicos}, porque dizem respeito á realidade. e exprimem o. que são as coisas em si mesmas, independentemente da. nossa inteligencia. --- Os principios \emph{logicos}, Como os \emph{metaphysicos}, faudam.sc, como vimos, no principio de \emph{con£radiccao}. D' este principio occupa-se a \emph{Ontologia}. \textbar{}\par Depois de termos explicado a natureza do \emph{syllogismo}, devemos agora examinar as objecções, adduzidas contra. o seu \emph{valor logico}. Limitamo-nos a Citar as de \emph{Stuart Mill}, de \emph{Spencer}, de \emph{Hacou} e \emph{Locke}, que são apenas umas reproducções das que foram apresentadas pelos escriptores antigos, e que são actualmente propugnadas pelos \emph{naterialistas}, \emph{positivistas} e \emph{modernistas}.\par «\emph{j John}. \emph{Stuart}. \emph{Mill} (\emph{Systema da}. \emph{Logica}) procura demonstrar que o syllo- \emph{E} gismo é um sophisma, um \emph{circulo vicioso}, visto que a verdade ger, que é : expressa pela premissa maior e de que a conclusão recebe toda a sua força,\pgn{69}depende da verdade já conhecida da mesma conclusão. Assim neste syllogismo --- \emph{todos os homens são} moriges; \emph{ora Pedro é homem}; \emph{logo Pedro} é \emph{mortal} --- conclue-se que este individuo é mortal, por serem mortaes todos os homens, quando é certo que se estabelece a premissa maior --- \emph{todos os homens são moríaes}, por se saber que é mortal cada um dos homens.\par Poucas palavras em resposta. A premissa maior do \emph{syllogismo} --- ou é um juizo \emph{analytico}, --- ou é um juizo \emph{synthetico}, --- Se é um juizo \emph{analifyco}, a sua verdade não se funda na verdade da conclusão, mas na propria essencia do . sujeito e do predicado: assim neste syllogismo --- \emph{todo o ser finito é contingente}; ora \emph{o mundo é ser finito}; \emph{logo o mundo é contingente} --- affirma-se que todo o ser finito é contingente, não por ser contingente o mundo, isto é, não pelo conhecimento da propria conclusão, mas pela analyse das duas ideas --- \emph{finito} e \emph{contingente}. --- Sc a premissa maior é un juizo \emph{synthketico}, baseado na observação e no \emph{raciocinio}, como acontece no seguinte syllogismo --- \emph{todos os homens são mortaes}; \emph{ora Pedro é homem}; \emph{logo Pedro} é mortal, --- tambem neste caso a verdade-d' aquella premissa não- depende da verdade da conclusão. Com effeito, a esta proposição \emph{universal} --- \emph{todos os homens são moríaes} --- não chega a nossa intelligencia pelo conhecimento do \emph{particular} ; vorque o \emph{perticalar}, por nós conhecido, refere-se ao \emph{passado}, e não ao futuro, \emph{ao} passo que aquella proposição \emph{universal} prescinde de todo o tempo, e abrange todos os homens, passados e futuros, e, onde se encontra o homem, ahi deve encontrar-se \emph{a} morte, se as leis da natureza não soffrerem uma suspensão. O equivoco dos adversarios consiste em verem na proposição --- \emph{foros os homens são mortaes} -- apenas uma referencia a \emph{collecção} de individuos, quando é certo que essi proposição exprime a natureza dos homens, pois reduz-se á esta: \emph{lodo o fiomem d mortal}. Logo a \emph{premissa maior} não recebe da \emph{conclusão} a sua força. --- Tambem é falso que a conclusão seja conhecida antes da \emph{premissa maior}; puis, embora aquella esteja contida nesta, todavia nós não o conhecemos \emph{émarediafaiente}, c por isso é necessario recorrer á umu outra verdade. D'onde se vê que o \emph{syllogismo} é um meio para \emph{descobrir} e conhecer \emph{distinctamente} novas verdades.\par Alem d' isso, é sabido que muitas vezes os sabios parlem dos mesmos principios universaes, admittidos por todos e chegam a; concfusoes inteiramente-oppostus. Ora isto não poderia acontecer: se antecedentemente se conhecessem todus as verdades particulares, contidas no principio universal, E se acontece, isto significa apenas que as conclusões oppostas têm um fundamento diverso. isto é, que os \emph{medios} empregados no raciocinio foram diversos. Portanto, a verdade do \emph{principio maéeersal} não deriva da \emph{verdade} das proposições \emph{particulares}, signilicadas pelas conclusões, mas a propria verdade da conclusão é o Iructo do \emph{meio fermo}. e por isso é uma coisa nova. --- Notamos que o argumento de \emph{Sfacrt Mill} se encontra em \emph{Bacon} e constitue uma objecção. que os velhos escholasticos faziam contra a legitimidade do syllogismo e a que respondiam Ltriumphantemente.\par O mesmo \emph{Start Mil} diz que a \emph{inducção} fornece as premissas \emph{maiores}\pgn{70}dos \emph{syllogismos}, pois não existe \emph{generalidade} sem \emph{induccio}. --- Hespondemos que a \emph{inducção} Sec toma em \emph{dois} sentidos. No sentido \emph{mais rigoroso}, exprime a \emph{argumentação} que vai do \emph{prerticuler} para o \emph{geral}. No sentido \emph{eis lato}, significa toda é qualquer \emph{generetisaczo}, embora feita \emph{inmediatamente} € sen \emph{raciocinio}. --- Tomada no primeiro sentido, a \emph{Zmduccao} fornece ao \emph{syllogismo} as premissas \emph{maiores}, quando ellas contém um juizo \emph{synthelico}, mas não quando contém um juizo «\emph{alytico}. Tomada no segundo sentido, á \emph{inducção} fornece as premissas \emph{maiores} do \emph{syllogismo}, pois é ecrlo que não existe \emph{fgeneralidade} sem \emph{indacco}, como diz \emph{Aristoteles}. (\emph{Post}. An. lib, 1, c. 1S).\par \emph{bd}) \emph{H}. \emph{Spencer} affirma que o \emph{silogismo} é inutil e arbítrario pelas seguintes. razões : --- 1.º) o syllogismo não pode ser compasto de proposições HUNiVErsaes, visto que o objecto-da. Logica, sendo constituido pelas coisas \emph{concretas}, \emph{singular}; --- 9 \textasciicircum{}) o nosso raciocinio natural nunca assume a forma de svliogismo; --- 3s) pelo \emph{syllogismo} não apprehendemos coisa alguma.\par Respondemne. que as razões sdduzidas por \emph{J}/. \emph{Spencer}: em confirmação\par da.sua opinião são falsas. --- É falsa a \emph{primeira} : porque. como provaremos, existem ideas \emph{universaes}, € \emph{O sensualismo}, que só admitte idéas \emph{particulares}, é antiscientifico. --- É fulsu a \emph{seguida}; porque, embora os nossos discursos não assumam \emph{ordinariamente} a forma de \emph{sy/Togisino}, nem por isso este deixa de ser util para as discussões scientificas e para a disciplina da intelligencia. --- É falsa a \emph{terceira}; porque muitas \emph{Vezes}” os sabios concordam nos mesmos principios geraes, mas discordam nas consequencias. E porque esta. divergencia? Porque não conhecem o.modo de.tirar.de determinados principios legilimas. consequencias. Ord é o \emph{syllogismo} que ensina a tirar de prin- €ipios universaes consequencias legitimas. Alem d'isso, o \emph{syllogismo}, como dissemos contra \emph{Mill}, serve para \emph{descobrir} e conhecer \emph{distinctamente} novas verdades, só implicitamente contidas nos principios geraes. Diz \emph{Sortais}: «Os progressos das sciencias, exactas attestam a fecundidade da \emph{deqducrao}. Se a objecção tivesse um fundamento, Lastaria conhecer OS axiomas, os postulados e as definições geometricas, para se conhecer a geometria inteira. E quando se diz que a premissa \emph{maior} contem a \emph{conclusão}, devemos entender que ella contem a razão pot; qual : conclusão é verdadeira, isto é, que a contem \emph{implicitamente}. Para que se torne \emph{explicite} tal inclusão, não basta 9 \emph{maior}; devemos recorrer á \emph{menor}, que mostra se a razão indicada na \emph{motor} conveuha \emph{ao sujeito} da \emph{conclusão}». (\emph{Traité de Phil}, v. M, n. 38).\par \emph{c}) \emph{Bacon} e \emph{Lock} rejeitam o \emph{syllogismo} e dizen que a \emph{inducção} é o unico meio para. adquirir x sciencia.\par Esta opinião é falsa. A \emph{imducção}, como veremos. é um dos nieios para a adquisicio da sciencia, mas não é o \emph{mico}. --- A legitimidade do \emph{xylTogismo}}.

S4. \textbf{Divisão do syllogismo}. ---. O syllogismo é é \emph{simples} ou \emph{composto}. --- Q \emph{simples} só consta. de proposições \emph{simples}. --- O \emph{composto} contém proposições \emph{compostas}, ou consta de syllogismos \emph{simples}. --- Agora tratamos do \emph{simples}; depois, do \emph{composto}.

\emph{Syllog sismo simples e suas regras}

85. Syllogismo \textbf{simples}. --- O syllogismo é \emph{simples}, quando

só consta de proposições \emph{simples}. Ex.:

\textbar{} Todo o espirito é incorruptivel ;

A alma humana é espirito ;

Logo a alma humana é incorruptivel.

86. \textbf{Regras} do \textbf{syllogismo simples}. .---: São oo \textbf{estas regras}: \emph{quatro} relativas aos \emph{termos}, e \emph{quatro ás proposições}.

I "Terminus este triplex, maior \emph{mediusque} minorque,

II. Latius hos, quam \emph{praemissae}, \emph{conclusio} non vult.

IH. Nequaquam \emph{aedium} capiat \emph{conclusio} oportet.

JV. Aut semel, aut iterum \emph{enedius generaliter} \textbf{este}.

V. \emph{Utraque} si \emph{praemissa} neget, nil inde sequitur.

VJ. Ambae \emph{affirmantes} nequeunt generare \emph{negantem}.

VII. \emph{Pejorem} sequitur semper \emph{concluso} partem.

VIIL Nil sequitur \emph{geminis} ex \emph{particularibus} unquam \footnote{A \emph{necessidade} da \emph{conclusão} no syllogismo depende exclusivamente do seu \emph{nexo} com as \emph{premissas}. Para que exista tal nexo, € necessario que uma das premissas contenha du conclusão, e a outra declare que tal continencia existe; porque então, postas as premissus, deve necessariamente deduzir-se a conclusão. Mas, para que as premissas contenham e Geclarem contida a conclusão, é necessario observar as citadas oito \emph{segras}.}.

HBecna L --- \emph{Os termos decem} ser \emph{tres}: maior, menor e \emph{medio}. --- Se forem mais ou menos que \emph{res}, não pode ter logar a comparação do \emph{sujeito} e do \emph{predicado} com o \emph{medio}. --- Falla lambem a esta regra quem emprega um termo equivoco, \textbf{Como}:

Todo o cão ladra:

Ora uma ceonstellação é cão;

Logo uma constellação ladra.

Rucra 11. --- \emph{Nenhum termo deve ler na conclusão} uma \emph{extensão maior que nas premissas}. --- Um termo, que é \emph{particular} nas \emph{premissas} e \emph{universal} na \emph{conclusão}, equivale a \emph{dois} termos eo \emph{syllogismo} tem \emph{quatro} termos; como no seguinte, em que o termo «\emph{vivo}» 6 \emph{particular} na \emph{maior} e \emph{universal} na \emph{conclusão} :

Todo o homem é vivo:

Ora o bruto não é bomem;

Logo o bruto não é vivo \footnote{Alem.d' isso, a conclusão, derivando das premissas, deve estar con- Lida nellas, como a. parte no todo. Ora, se um termo tem na conclusão umi extensão superior á que tem nas premissas, já não pode estar contido nellas, porque o \emph{waís} não està contido no menos. ---- Finalmente, na conclusão do syllogismo unimos ou separamos dois termos, porque foram comparados nas premissas com o medio. Ora se um termo fosse mais extenso na conclusão do que nas premissas. feria sido \emph{comparado} com o medio só numa \emph{parte} da sua extensão, e assim pereceria toda a estructura do syllogismo.}. é reconhecida por sabios nolaveis e insuspeitos. \emph{Leibnitz} escreve ( N. \emph{Ensaios}, LIV, e. 16): «Sustento que a descoberta do \emph{syHogismo} é uma das mais bellas e importantes do espirito humano; o \emph{syllogismo} é uma especie de mathematica universal, cuja importancia é desconbecida, más que \textbf{encerra} uma arte de infalli bilidade». --- \emph{Claudio Bernard} \textbar{} \emph{Introd}. ao \emph{estudo} da \emph{medicinm exp}): « O espirito humano funcciona sempre por meio do «\emph{yllogismo} .. Poderia demonstrar a minha proposição com argumentos physiologicos ».

72 “TT LoGICA ="SECÇÃO PRIMEIRA

\textsc{ReGRrAa} IIL --- \emph{O termo medio não deve} entrar \emph{na conclusão}, --- Recorremos ao termo.iedio, para com elle compararmos o \emph{sujeito} e o \emph{predicado}. Elfectuando-se a comparação nas \emph{premissas}, o termo \emph{medio} não pode entrar na \emph{conclusão}, que deve "declarar exclusivamente a conveniencia ou a desconveniencia entre o sujeito e o predicado. --- É falso, pois, este syllogismo :

Napoleão foi pequeno ;

Ora Napoleão foi capitão;

Logo Napoleüo foi pequeno capitão.

\textsc{RgcRA} IV. --- \emph{O fermo medio deve ser Universal}, «\emph{o menos} em uma \emph{das premissas}. --- Na verdade, se o \emph{medio} é particular em \emph{ambas} «\emph{s premissas}, equivale a \emph{dois} termos, pois exprime dois objectos ou individuos distinctos, e falta a comparação “dos extremos com um e o mesmo \emph{medio}; assim é falso o seguinte syllogismo: .

Alguns homens são ignorantes;

Ora alguns homens são sabios:

; Logo alguns sabios são ignorantes \footnote{É porém, verdadeiro o syllogismo, quando o termo \emph{medio é singular} e incommunicavel porque este, tomando-se em toda a sua \emph{extensão}, corresponde ao \emph{universal} e siguifics sempre um e \emph{o} mesmo objecto ou individuo, como neste exemplo: :\par Luiz de \textbf{Camões} foi pobre; 2\par Ora Luiz de Camões foi o auctor dos Lusíadas ;\par Logo o auctor dos Lusiadas. foi pobre.\par Dissemos --- quando Oo termo \emph{media é singular} e \emph{iucommnnica vel}; porque se o termo é \emph{comanuiicavel}, emlrora \emph{siuguler}, torna-se, de algum modo, \emph{particular}, e não pode ser \emph{medio}, Assim não podemos dizer: O Pae € \emph{Deus}; mas o Filho é \emph{Deus}; logo o Filho é o Pae: porque o termo \emph{Deus}. embora \emph{singular}, é commuwnicavel e não é proprio de uma só Pessoa divina.}.

\textsc{ReGRA} V. --- \emph{De duas premissas negativas nada} se \emph{pode concluir}. --- Sendo \emph{negativas} as. duas premissas, É signal de que \emph{o i1edio} não é medida nem do sujeito, nem do predicado, e por isso não podemos concluir a \emph{conveniencia} nem a \emph{desconveniencia}. do predicado com o sujeito. Assim d'estas proposições :

Toda z planta. não é immortal :

Ora a alma humana não é plan z

; Logo . . " - . . » não podemos concluir que a nossa alma é, ou não, immortal \footnote{Não pode deduzir-se nenhuma conclusão, --- nem \emph{affiraativa}, porque a conclusão \emph{affirmativa} suppõe a conveniencia de ambos os extremos com O medio, --- nem \emph{negativa}, porque esta exige que um dos extremos convenha\pgn{73}com a medio. --- Mas, para que esta. rera conserve loto. 0 seu vigor, não basta que ambas as premissas, consideradas em si, sejam \emph{negativas}, mas é necessario que sejam negalivaus com \emph{relação ao termo medio}, emquanto ambas \emph{negem}. X conveniencia dos dois extremos com o medio: aliás o \emph{syllogismo} Será legitimo. Assim é \emph{legitimo o} seguinte; \emph{quem}. \emph{não} ani a \emph{Deus não será}, \emph{Salvo} ; \emph{mas o impio não ama} « \emph{Deus}: \emph{Togo o} £\emph{mpio iiio será salvo}; \emph{porque a menor}, embora, considerada. em si, seja, \emph{regativa}. todavia, considerada \emph{em} relação a0 \emph{termo medio}, é \emph{affirmativa}, porque \emph{affrma} x \emph{conveniencia} do sujeito \emph{impio.com} 'o medio termo-zo «\emph{ema} c« \emph{Deus}, pois é equivalente a esta --- o \emph{nnmpío} é não \emph{amante de Deus}. que é certamente \emph{affirmativa}.}.

%%>>> cap-6

\chapter*{CAPITULO TERCEIRO\\[0.3em]\large RACIOCINIO E ARGUMENTAÇÃO --- 73}
\addcontentsline{toc}{chapter}{CAPITULO TERCEIRO --- RACIOCINIO E ARGUMENTAÇÃO --- 73}

\pg{69} \pg{70} \pg{71} \pg{72} \pg{73} RecRa VI. --- \emph{De duas premissas affirmalivas não se pode tirar} uma \emph{conclusio negativa}. --- Se as premissas são \emph{affirmativas}, o sujeito e o predicado convém com o termo \emph{medio}, e, por isso, entre si. É falso, pois, o seguinte syllogismo :

Todo o homem honesto é digno de louvor:

Ora Pedro é homem honesto;

Logo Pedro não é digno de louvor.

\textsc{REGuA} VII, --- 4 \emph{conclusão segue} sempre « \emph{premissa mais fraca}. --- A premissa \emph{mais fraca} é a \emph{negativa} relativamente á \emph{affirmative}, e a \emph{particular} relativamente á \emph{wnicersat}. Por isso esta regra resolve-se nas duas seguintes :

1*) \emph{Se uma premissa é negativa}, \emph{temubem}. \emph{deve ser negativa} a \emph{conclusão}. Porque, sendo \emph{negativa} uma das premissas, um dos \emph{extremos} não convem com o \emph{medio}, e por isso os dois extremos não podem convir entre si. É falso, pois, este syllogismo :

Todo o corpo é extenso :

Ora a alma humana não é corpo :

Logo a alma humana é extensa.

9?) \emph{Se uma das premissas é particular}, \emph{tumbem deve ser particular} «a \emph{conclusão}. Porquanto, se a conclusão fosse \emph{universal}, sendo uma das premissas \emph{particular}, um termo teria maior extensão na \emph{conclusão} do que nas \emph{premissas}, o que é contrario 4 \emph{Regra} 7\emph{I}. Por isso é falso o seguinte syllogismo :

Todas as acções bóas merecem louvor;

Ora alguma occupacao é bôa ;

Logo todas 1s oceupações merecem louvor.

Recta VIII. --- \emph{De duas premissas particulares nada pode concluir-se} --- As premissas particulares ---- \emph{a}) ou são «\emph{ambas affirmativas}, --- \emph{b}) ou \emph{ambas negativas}, --- \emph{c}) ou «\emph{ma affirmation} e \emph{outra negativa}.

\pg{74} a) Se as premissas particulares são \emph{ambas affirmativas}, nada se conclue. Porque o termo \emph{snedio} deve ser \emph{universal}, polo menos numa das premissas, como manda a \emph{Regra IV}; ora em duas premissas \emph{particulares-affirmativas} não ha algum termo \emph{universal}, pois os \emph{sujeitos} e os \emph{predicados} são lodos \emph{particulares}. Por isso é falso o seguinte syllogismo :

Alguns homens são sabios;

Alguns homens são ignorantes;

Logo alguns ignorantes são sabios.

b) Se as premissas \emph{particulares} são ambas \emph{negativas}, lambem não ha conclusão; porque o prohibe a \emph{Zegra} V.

c) Se de duas premissas \emph{particulares} - uma é \emph{affirinative} € outra é \emph{negativa}, não pode \textbf{tirar-se} nenhuma conclusão --- nem \emph{aflirmativa}, nem \emph{negativa}. Não pode tirar-se uma conclusão \emph{affirmativa}; porque o prohibe a \emph{Regra VII}. Nem pode tirar-se uma conclusão \emph{negativa}; porque, para tirarmos uma conclusão \emph{negativa}, deveria haver nas premissas dois termos \emph{universaes}, oceupados pelo \emph{predicado} da conclusão e pelo termo \emph{medio} (\emph{Regras lie IV}); ora nas premissas ha um só termo \emph{universal}, que é o \emph{predicado} da premissa \emph{negativa}. Assim é falso este syllogismo:

Alyuns homens são maus;

Alguns homens não são maus;

Logo alguns maus não são maus \footnote{\pgn{75}dos mesnios principios. --- A \emph{Logica} de \emph{Porto Feml} reduz á \emph{seis ás oito} regras, supprimindo a [ e a 111, porque evidentes. Mas, se este criterio fosse acceilavel, poderiam exceluic-se tambem outras, que lambem são muito evidentes, «como a Il: e. dizemos mais, poderiam todas reduzir-se 4 I. que é fundamental: pois é nella. que se busseam as outras \emph{sec}. --- A razão das \emph{oito} regras não é a sua evidencia, maior ou menor, mas é o fim pratico, para o qual tendem e que é evitar o erro, e descobril-0 fneilmente, quando proposto por outros. Sob este aspecto, as oito regras são todas ulilissimas, porque são varios moeios para chegarmos a conclusões verdadeiras.}.

(f) Esta regra soffre excepção, quando o \emph{sy/logismo} contem lermos \emph{singulares}, ou ternios \emph{particulares} equivalentes aos. \emph{sinyulares}; pois esses lermos, tendo uma e a mesma significação ou extensão, correspondem aos \emph{uriversaes}. Assim é legitimo o seguinte syllogismo: \emph{S}. \emph{Pedro é Galileu}; \emph{ora S}. \emph{Pedro é o Principe dos Apostolos} ; \emph{logo o Principe dos Apostolos é Caldeu}. Neste exemplo, os extremos são comparados com um e \emph{o} mesmo lermo medio, e assim salva-se o principio de \emph{conveniencic}.

Alguns escriptores affirmam que todas as leis do \emph{syllogismo} podem reduzir-se a uma só, que é a seguinte: \emph{O syllogismo simples é legitino}, \emph{quando} uma \emph{das premissas inche} «\emph{a conclusão}, \emph{e a outra esprime} essa \emph{inclusão}. --- Respondemos que essa Jei, a qual exige o nexo entre ás premissas e a com clusão, não só não exclue as \emph{oito regras} do syllogismo, mas exige-as: porque são ellas que ensinam o que se ha de fazer ou evitar para que exista O nexo entre o antecedente e o consequente.

Outros dizem que, para a legitimidade do sylozmismo, bastum os dois principios fundamentues de \emph{eomreniencia} ou de \emph{desconveniencia}. Mas elles não reflectem \emph{que} us oilo regras propostas servem para a. recta, applicação

%%>>> art-6

\section*{Artigo IV.}
\addcontentsline{toc}{section}{Artigo IV.}

\pg{75} \emph{Figuras e modos do syllogismo simples}

2

87. \textbf{Forma} do syllogismo. --- No \emph{syllogismo}, alem da \emph{meateria}, que é constituida pelos \emph{termos} e pelas \emph{proposições}, lemos a \emph{forma}, que é a \emph{conveniente disposição} dos \textbf{mesmos} \emph{lermos} e das \emph{premissas} em ordem á \emph{conclusio}. --- E, como a disposição dos \emph{termos} se diz \emph{figura}, e a das \emph{premissas} se.chama \emph{modo}, assim . neste artigo tratamos das \emph{figuras} e dos \emph{nodos} do \emph{syllogismo}.

S8. \textbf{Figura}. --- \emph{Figura} é a diversa disposição do \emph{termo medio} relativamente aos \emph{dois extremos}.

SO. Numero das \textbf{figuras}. --- As \emph{figuras} são \emph{quatro}. Na verdade, os diversos modos de dispor \emph{o fermo medio} com relação aos \emph{dois extremos} são \emph{quatro}: pois que o termo \emph{medio} pode ser; --- 12) \emph{sujeito} na \emph{premissa maior} e \emph{predicado} va \emph{menor}; --- 9,") \emph{predicado} na \emph{wuwtior} e \emph{sujeito} na \emph{menor}, --- 37) \emph{predicado} na \emph{maior} e na \emph{menor}, --- 4\&9?) \emph{sujeito} na \emph{maior} e na \emph{menor}.

Damos os exemplos das quatro figaras:

\emph{Fx}. \emph{de} 1º \emph{füjp} --- Todo o homem merece affeição:

Ora o teu inimigo é homem:

Logo o teu inimigo merece tileição.

\emph{Joe}. \emph{da} 2 \emph{fig}. --- Todo o homem é animal;

Ora todo o animal é substancia:

Logo alguma substancia é homem.

\emph{Fx}. \emph{da} 4\textasciicircum{} \emph{fig}. --- Toda a planta não é animal;

Ora todo o homem é animal:

Logo todo o homem não é planta.

\emph{Fx}. \emph{da} 4\textasciicircum{} \emph{fig}. --- 'Todo o homem é racional;

Ora algum homem é ignorante;

Logo algum ignorante é racional.

\pg{76} A \emph{primeira figura}. 6 a mais perfeita, porque de uma maneira evidente representa o organismo do syllogismo \footnote{Como a. \emph{segunda}. figura é apenas um modo \emph{indirecto} da \emph{primeira}; as figuras podem reduzir-se a \emph{tres}, conforme o termo \emph{medio} --- ou é uma \emph{vez} sujeito e ua \emph{vez} predicado, --- ou a \emph{duas} veces predicado, --- ou é \emph{duas vezes} sujeito. --- Todavia, a \emph{primeira} figura pode desdobrar-se e constituir uma \emph{duplice} figura, conforme o logar que o \emph{medio} occupa nas premissas, e assim temos a \emph{primeiro} figura \emph{directa} e \emph{indirecta}; --- \emph{directa}, se 0 \emph{medio é sujeito} da. premissa \emph{maior} e \emph{predicado} da \emph{menor}, --- \emph{indirecta}, se o \emph{medio é predicado} da premissa \emph{maior} e \emph{sujeito} da \emph{menor}. --- A \emph{primeira} figura é a mais perfeita; porque, nesta figura, o termo \emph{medio} é verdadeiramente \emph{medio}, emquanto, sendo \emph{sujeito} e \emph{predicado}, Participa da natureza dos dois extremos. que são o sujeito e o prediendo da conclusio. Mas a \emph{primeira} figura \emph{directa} é mais perfeita que a \emph{primeira indirecta}; não só porque, na \emph{dérecía}, os termos occupam na conclusão o logar de \emph{Sujeito} e de \emph{predicado}, que oceuparam nas premissas, mas tamhem porque a \emph{mdirecía} exprime, natural e conipletamente, o organismo e a indole do Syllogimmo, io passo que a \emph{indirecta} segue um processo contorto e.violento, nem mostra claramente o fundamento, em que o syllogismo se baseu. --- Menos perfeita é a figura, em que o \emph{medio é predicado} em ambas as premissas, porque então o \emph{medio} não é verdadeiramente \emph{medio} entre o sujeito e o predicado da conclusão. --- Todavia esta figura, em que o \emph{medio é predicado} em ambas as premissas é mais perfeita que a ultima, em que. \emph{o medio é sujeito} em ambas as premissas; pois é mais nobre o predicado -que o sujeito (S. Th.. \emph{Op}. 48, tr. S, c 4).\par As \emph{tres ultimas} figuras podem facilmente reduzir-se a \emph{primeira}. --- O se guinle syllogismo da \emph{primeira} figura \emph{indirecta}; --- \emph{todo o homem é animal}, \emph{oru todo o animal é substancia}; \emph{logo lguma substancia}. \emph{é homem} --- pode reduzir-se a \emph{primeira} figura \emph{directa}, dizendo: --- \emph{tado o cuinal.é substancia}; ora \emph{todo q homem é animal}; \emph{logo todo o homen é substancia}, --- Assim este syHogismo da \emph{segunda} (ou \emph{terceira}) figura: --- \emph{todo o bruto} não é \emph{racional} ; ora. \emph{todo o homem}. \emph{d racional}: \emph{loto todo e homem}. \emph{não é bruto} --- veduz-se a \emph{primeira} (\emph{directa}), dizendo: --- \emph{todo o racional não é bruto}: \emph{orc todo o homem é racional}; \emph{Togo todo o}. \emph{hoimem não é bruto}. --- Como tambem o seguinte syllogismo da \emph{terceira} (ou \emph{quarta}) lura: --- \emph{todo o homem é racional}; \emph{ora elgiin homem é impro}; \emph{logo algum impio a racional}, --- reduz-se á \emph{primeira} (\emph{directa}), dizendo: --- \emph{todo o homen}. \emph{é racional}: ora \emph{algum impio é homem}: \emph{logo algum Dupéo é racional}, --- Repetimos --- que, se a \emph{primeira} figura se desdobra, temos \emph{quatro} figuras, --- e que, se não se desdobra, só temos \emph{tres}.}.

%%>>> t90

\tese{90}{Principios das figuras.}
Cada uma d'estas figuras tem seus \emph{principios} proprios. --- A \emph{primeira} ea \emph{segunda} figura fundam-se nos dois principios, em que se funda o syllogismo e que, como dissemos, são uma derivação: dos principios de conveniencia e de \emph{desconveniencia}. --- A \emph{terceira} figura deve exprimir o seguinte principio: \emph{Duas coisas}, \emph{uma das quaes não está contida} no \emph{genero}, em \emph{que está contida a outra}, \emph{differem entre si}. ---- A \emph{quarta} figura basea-se neste axioma: \emph{Se alguns} \pg{77}  \emph{sujeitos}, \emph{subordinados a um genero}, \emph{pertencem} tambem \emph{a} outro \emph{genero}, \emph{estes dois generos tem em comunan alguns sujeitos} \footnote{Mas reconhece-se facilmente que os principios, ou axiomas, que regem a \emph{terceira} e a \emph{quarta} figura, são apenas diversos aspectos dos principios, em que se fundam a \emph{primeira} e a \emph{segunda} loura,}.

%%>>> t91

\tese{91}{Leis das figuras.}
Cada figura tem as suas leis proprias, que são dilferentes applicações das \emph{oito} regras do syllogismb \emph{simples} e que, por isso, não se podem violar, sem que ao mesmo tempo sejam violadas essas regras. --- As leis das \emph{figuras} são \emph{seis}: umu para a \emph{primeira} figura, \emph{tres} para a \emph{sequnde}, uma para a \emph{terceira} e uma para a \emph{quarta}. São signifiondas nos seguintes versos:

\emph{Para} a [2 \emph{fy} --- 1 Sit minor \emph{affirmans}, \emph{maior} vero \emph{generalis}.

\textbar{} ll. \emph{Maior} ubi \emph{affirmat}, \emph{generalis} tum \emph{minor} este.

HAE. Si \emph{iniuor affirmat}, \emph{conchisio} sit \emph{specialis}.

\emph{Para} o 2º \emph{fig}.

\textbar{} IV. Si \emph{wegat hace}, maior \emph{generalis} postulat esse,

\emph{Para} « 2* \emph{fig}. --- V. \emph{Una} negans este, nec \emph{maior} Sit \emph{specialis}.

\emph{Para} a á \emph{fig}. --- VL. Sit minor \emph{affirmans}, \emph{conclusio particularis} \footnote{Advertimos que, neslas leis, a »\emph{uttior e} a menor são tomadas em sentido rigoroso, e assim a premissa \emph{maior} é a que contém o extremo \emph{maior}, isto é, o \emph{predicado} (\emph{da conclusão}), e a premissa \emph{menor} é a que contem o extremo \emph{menor}, isto é, o \emph{sujeito} (\emph{da}. conclusão).}.

Lei I. --- A menor \emph{seja affirmativa}; \emph{a maior}, \emph{universal}.

a) A menor \emph{sejt affirmativa}. Por quanto, se a premissa \emph{menor é negativa}, lambem \emph{negativa} ha de ser a \emph{conclusão} (Regra VII do \emph{Syl}). Se a \emph{conclusão é negativa}, o seu \emph{predicado é universal}. Se o \emph{predicado} é \emph{universal} na \emph{conclusão}, deve ser \emph{universal} tambem na premissa \emph{maior} (Regra ]l \emph{do Syl}.). Sec \emph{predicado} da premissa \emph{maior é universal}, a propria premissa é \emph{negalioa}. Se a premissa maior é tambem \emph{negativa}, não pode haver conclusão (Regra V do \emph{SyL}). Logo a premissa \emph{menor} deve ser \emph{affirmativa}. Assim é falso o seguinte syllogismo : \textbar{} Todo o homem é vivente:

Ora nenhuma planta é homem;

Logo nenhuma planta é vivente.

b) A \emph{maior seja universal}, Na verdade, se a premissa \emph{maior} fosse \emph{particular}, o termo \emph{medio} seria \emph{particular} --- tanto na \emph{maior}, porque \emph{sujeito} de uma proposição \emph{particular}, --- como na menor, porque. \emph{predicado} de uma proposição \emph{affirmativa}; o. que se oppõe é Regra IV do \emph{Syllogismo}. Assim é falso o seguinte:

5g \emph{sec} --- --- Lodick- = SECÇÃO vRIEEÍRA--- .

Algum corpo é vivente; -

Ora a pedra é um corpo;

Logo a pedra é vivente.

Lei 1l. --- \emph{Se} « \emph{maior é affirmativa}, á \emph{menor deve ser universal}. Na ve rdade, se \emph{a'maior é affirmativa}, o termo \emph{medio}, que é \emph{predicado é particular}. Por isso, o \emph{medio} deve ser \emph{universal} na \emph{menor} (Regra IV 'do \emph{SyL}), e a mesma menor deve ser \emph{universal}, porque o \emph{medio} occupa ' o logar de \emph{sujeito}. Pecca contra esta regra o seguinte:

'l'odo o animal é vivente;

Ora algum vivente é planta;

logo alguma planta é animal.

Lei Hi. --- \emph{Se} a menor é \emph{affirmativa}, \emph{a conclusão deve} ser \emph{particular}. Com effeito, se a premissa \emph{menor é affir} mativa, o extremo \emph{menor}, que é \emph{predicado} d' essa premissa é \emph{particular}, e por isso, quando entra na conclusão como \emph{sujeito}, deve ser tambem. \emph{particular}, "e assim a conclusão é “necessariamente \emph{particular}. É, “pois, legitimo o seguinte

' ; Nenhum homem é bruto;

Ora todo o bruto é sensitivo:

Logo algum sensitivo não é homem,

Lei IV. --- \emph{Se} a \emph{conclusão é negativa}, \emph{a maior deve ser universal}. Por quanto, se a conclusão é \emph{negativa}, o seu \emph{predicado é universal}. Se o predicado é \emph{universal} na \emph{conclusão}, deve ser lambem \emph{universal} na premissa \emph{maior}, da qual é \emph{sujeito}; e por isso a premissa \emph{maior} deve ser \emph{universul}, porque \emph{universal} é a proposição, cojo sujeito é \emph{universal}. É, pois, legitimo o seguinte:

* Toda a planta é vivente;

Ora nenhum vivente é pedra;

Logo nenhuma pedra é planta.

Lei V. --- \emph{Uma das premissas seja negativa}, \emph{a maior seja universal},

a). \emph{Uma das premissas} seja \emph{negativa}, Na verdade, se \emph{ambas}

as \emph{premissas} são \emph{affir mativas}, o termo \emph{medio}, que é sempre \emph{predicado}, é sempre \emph{particular}, porque \emph{particular} é o \emph{predicado} do proposição \emph{affirinativa}, e [alla-se contra a EV Regra do \emph{Syl}. Por isso é falso o seguinte É

O homem é animal; :

Ora o cão é animal;

Logo o cão é homem.

b) A \emph{maior seja unicersal}. Um termo, que é, ao mesmo tempo. \emph{predicado} na \emph{conclusão} e \emph{sujeito} na \emph{maior}, não pode ser menos extenso na \emph{maior} do que na \emph{conclusão} (Regra 11 do \emph{Syl}.). Mas, na \emph{conclusão}, esse termo é \emph{universal}, porque \emph{predicado} de uma proposição \emph{negativa}. Logo deve ser tambem \emph{universal} na \emph{maior}. Ora o sujeito universal torna \emph{universal} a proposição. Logo a maior deve ser \emph{universal}. Por isso é falso « o seguinte:

Algum vivente é homem:

Ora nenhuma planta é homem;

. Lomo nenhuma planta é vivente.

Lei Vl. --- A menor \emph{seja affirinatioa}; \emph{a conclusão particular}.

a) A \emph{menor seje affirmativa}. Se a \emph{menor é negativa}, tambem a \emph{conclusão} ha de ser \emph{negativa}, e o \emph{predicado} da mesma \emph{conclusão é universal}. Sendo \emph{universal} na \emph{conclusão}, o \emph{predicado} ha de ser \emph{universal} tambem na premissa maior. Se o \emph{predicado é universal} na \emph{maior}, esta é \emph{negativa}, porque só na proposição \emph{negativa} o \emph{predicado} é tomado \emph{universalmente}. Mas de duas premissas \emph{negativas} não pode deduzir-se a conclusão. Logo a premissa \emph{menor} deve ser \emph{affirmativa}. É falso, pois, o seguinte:

Todo o animal é substancia ;

Ora nenhum animal é planta;

Logo nenhuma planta é substancia.

b) A \emph{conclusão seja particular}. Se a \emph{conclusão} fosse \emph{universal}, o seu \emph{sujeito} havia de ser \emph{universal}. Mas como o \emph{sujeito} da conclusão é tambem o \emph{predicado} da \emph{menor}, o \emph{predicado} da. menor devia ser universal. Ora o \emph{predicado} da menor não pode' ser \emph{universal}, porque é predicado de uma proposição \emph{a/ffirmaliva}, mas deve ser \emph{particular}. Portanto a \emph{conclusio} deve ser \emph{particular}. É, falso, pois, o seguinte:

Nenhum bruto é planta ;

Ora todo o bruto é vivente;

Logo todo o vivente não é planta.

\$2. Modo. --- \emph{Modo} é a diversa disposição das \emph{premissas} consideradas na sua \emph{quantidade} e \emph{qualidade}.

%%>>> t93

\tese{93}{Numero dos modos.}
Sendo o \emph{modo} a diversa dispo- .sição das \emph{premissas} consideradas na sua \emph{quantidade} e \emph{qualidade}, .6 claro que serão tantos os \emph{modos}, quantas forem as possiveis disposições \emph{das premissas}, assim consideradas. Ora \emph{as premissas}, consideradas na sua \emph{quantidade} e \emph{qualidade}, podem ter, em \emph{cada} \pg{78} \pg{79} \pg{80}  \emph{figura}, 16 diversas disposições ; e, como as figuras são \emph{quatro}, segue-se que todos os \emph{modos} do syllogismo se reduzem a 64.

9£, Numero dos modos legitimos. --- Nem todos os 64 modos s\&o legitimos, Por quanto, se excluirmos os que são abertamente contrarios ás Regras VI e VII do \emph{Syllogismo}, licam apenas 36, 9 em cada figura. E se estes 36 modos se considerarem em ordem á conclusão, que deve ser deduzida em conformidade não só das \emph{regras} communs do \emph{Syllog gismo}, mas tambem das \emph{leis} proprias de cada \emph{figura}, veremos que os modos \emph{legitimos} se reduzem a 19 \emph{OC}).\footnote{Apresentamos o schema. (a) das diversas disposições das premissas, «pnsideradas na sua \emph{quantidade} e \emph{qualidade}, recordando que \emph{za} letlra A exprime a premissa «\emph{niversal-affirmativa}, E a \emph{uneocersal-negatioa}, X a \emph{particular-affirmativa}, O a \emph{particular-negativa},\par (a) AA EA . lA OA\par t AE ER. IE \textbar{} OK\par Al \textasciicircum{} El Jl 'OI\par \textbar{} AO- EO 10. 00.\par ' Cada \emph{figura}, pois, tem 16 \emph{modos}. E, como as figuras são \emph{quatro} todos os \emph{modos} são 64. --- Mas nem todos estes modos são lerritimos ; pois alguns --- EE, EO, 17, 10. OE, OI, OO --- pecc am contra a regra V.* on'a V]IIL* E assim ficam 9 \textbf{modos} uteis em cada figura, que são os seguintes\par (b) AA EA [A OA\par AE \textbar{} --- E \textbar{} ---\par AL \textasciicircum{} EL --- ---\par AO --- --- --- '\par Mas estes 9 \emph{modos} ficam ainda.mais reduzidos, ge se avaliarem não só pelas regras communs a todos os syllogismos, mas tambem pelas \emph{leis proprias} de cada \emph{figura}. Assim ::\par Na 1.º Figura, \emph{os modos} --- AE, AO, IA, IE, OA --- não concluem legitimamente, porque peccam contra, a lei d'essu figura, que é: 4 \emph{menor seja effiriativa}; \emph{a maior}, \emph{universal}. Por i 1sso temos \emph{gratro} modos uteis, que são:\par (c) AA cuja conclusão é A\par AI + 1\par “EA » E\par ER. C » [5 * Na 2.º Figura, os \emph{modos} --- Al, AO. IE, OA --- s\&oillegitimos : porque OS dois primeiros peccam contra a lei: \emph{Se} a \emph{maior é affirmativa}, \emph{a} menor \emph{deve} ser \emph{universal}, e os outros dois peccam contra a lei : \emph{Se} a \emph{conclusão é negativa}, « \emph{maior deve} ser \emph{miversal}. Assim ficam cinco modos legitimos, que são:}\footnote{AA cuja conclusio é I\par : AE » E\par EA » o\par El » [9]\par IA » 1\par Na 3,º Figura, 08 \emph{modos} --- AA, AT, TA, IE, OA --- são legitimos, porque\pgn{81}peccam contra, a lei propria d' esta figura: \emph{Uma premissa seja negativa}; a \emph{maior}, universal. --- Assim, ficam \emph{quatro} modos legitimos que são:\par (e) AE cuja conclusão é E\par AO » o\par EA » E\par El » 0\par Na 4.º Figura, os \emph{modos} --- AE, AO, 1E --- são illegitimos, porque peccam contra a regra d'esta figura, que diz: A menor \emph{deve ser affirmativa}. Assim, os \emph{modos} legitimos são \emph{seis}, 1 saber:\par (£f) AA cuja conclusão é [\par AI » I\par EA » O\par EI » t\par LA. » I\par OA » o\par Sendo, pois, \emph{quatro} os \emph{modos} legitimos nu \emph{primeira} gura, \emph{cinco} na \emph{segunda}, \emph{quatro} na \emph{terceira}, \emph{seis} na \emph{quarta}, é claro que os \emph{modos} uteis são, ao todo, dezenove --- Os \emph{modas} mais perfeitos são Os da primeira figura, pois manifestam evidentemente a necessidade da conceiusão.\par Estes 19 modos legitimos são indicados pelos seguintes versos, que, embora barbaros na apparencia, são, comtudo, eonstruidus com grandissima arte, porque as \emph{ires} primeiras \emph{vogaes} de cada palavra significam \emph{a qualidade} e a \emph{quantidade} das premissas e da conclusão, € as \emph{consoantes} servem para \emph{o} fim, que mais adeante veremos. Lembramos que --- a significa a proposição \emph{universal-affirmaoliva}, --- e denota a universal-negativa, --- A significa a proposição \emph{particular-affirmativa}, --- o denota 2 particular-negativa. (1.º fig.) \emph{Barbara}, \emph{Celarent}, \emph{Darii}, \emph{Ferio}. --- (2.º Jig.) \emph{Baralip} (tou).\par Camentes, \emph{Dimatis}, \emph{Fresapno}, \emph{Fresisom} (orum). (3a fig.) \emph{Cesare}, Camestres, \emph{Festino}, \emph{Baroco}. --- (4.º fig.) \emph{Darapti},\par \emph{Felapton}, \emph{Disamis}, \emph{Datisi}, \emph{Bocardo}, \emph{F'erison}.\par Os-modos \emph{uteis} ou. \emph{legitintos}. dividem-se em. \emph{perfeitos}, e \emph{imperfeitos}, São \emph{perfeitos}, quando a conclusão \emph{clara} e \emph{directamente} deriva das premissas. São \emph{imperfeitos}, quando a conclusão só \emph{indirecta} e esciramente deriva das premissas. \emph{Perfeitos} são os \emph{quatro} modos da. 1.º figura, que de uma maneira immediata mostram a. propria força da conclusão. \emph{Dmperfeitos} são todos os outros \emph{quinze}, nos quaes a necessidade da conclusão não \emph{se} infere unicamente da posição das premissas, mas deve ser demonstrada por outro meio, Demonstra-se pela redneção e pela \emph{dedueção ao impossivel}.\par \emph{Jieducesro} é a revocação de um modo \emph{imperfeito} a um \emph{inodo perfeito}. pela \emph{conversão} das proposições, e ás vezes pela \emph{transposição} das premissas. --- Para que a reducção se faça legitimamente, é preciso considerar tres\par Vol. L 7\pgn{82}coisas : --- 1.5) \emph{qual} é o modo \emph{perfeito}, ao qual deve reduzir-se um modo imperfeito; --- 2\%) quaes as proposições, que se hão de sujeitar a conversão, e por que forma; --- 3.º) quando deve Ler logar a transposição das premissas, Ora estas Lres coisas são indicadus pelus \emph{consoantes} das palavras, que exprimem os \emph{modos}, --- Às lettras \emph{B}, \emph{C}, \emph{D}, F.'por que começam as palavras, significam que um determinado modo \emph{imperfeito} se deve reduzir áquelle dos \emph{quatro} modos \emph{perfeitos}. que comeca pela mesma letra, Assim \emph{JBaralip} reduzse a \emph{Barbara}, \emph{Cesare a Colarent}, \emph{Delepti} 4 \emph{Darii}, \emph{Festino a Kerio}. --- As leitras \emph{S}. e \emph{P}., que se acham no meio das palavras, indicam as proposições, que devem \emph{converter-se}, e n forma da \emph{conversão}, Assim, a lettra \emph{S} diz que a. proposição, significada. pela \emph{vogal antecedente}. deve estar sujeita a conversio \emph{simples}, isto é, madando o predicado em sujeito, mas conser-. vàndo a \emph{quantidade} A letra \emph{P} sisnilica que a proposição, indicada pela vogal anlecente, deve estar sujeita á conversão \emph{per accidens}, ou \emph{accidental}, mudando o predicado em sujeito, e alterando a quantidade. A lettra M, seja qual for o seu logar, indica a transposição .das premissas, de modo que a \emph{menor} se torne \emph{maior}, e vice-versa. O que levamos dito vê-se no seguinte syllogismo em \emph{Camestres} : --- \emph{todo o homem é racional}: \emph{ora}. \emph{nenhum cão é racional} ; \emph{logo nenTaun}. \emph{cão é homem}, A letira \emph{C}, que é inicial em \emph{Camestres}, indica que este modo imperfeito se deve reduzir ao modo períeito, que comeca. pela mesma letlra \emph{C} e que é \emph{Celarent}. A lettra. M. que na palavra \emph{Camestres} segue a primeira vogal, indica que as premissas se devem mudar. A lettra 5, que vem depois da segunda e terceira vogal, significa que a premissa \emph{menor} e a \emph{conclusão} se devem sujeitar á conversão \emph{simples}. Se tudo isto £c fizer, aquelle syllogismo imperfeito em \emph{Camestres} reduz-se ao seguinte perfeito em \emph{Celarent}: --- \emph{nenhum racional} é cão; ora \emph{todo o homem é racional}; \emph{logo nenhuns homen}. \emph{é cão}. --- Feita a redneção, vê-se claramente a necessidade da conclusão. --- Esta Teducção é possivel em \emph{todos} os modos imperfeitos, á excepção de dois --- de \emph{Baroco} e \emph{Bocardo}.\par A \emph{dedlucção} «\emph{o}. \emph{impossivel} consiste em obrigar o adversario, que nega a consequencia, depois de ter acceitado as premissas, a admittir o impossivel e o absurdo --- a verdade de duas contradictorias. Faz-se a deduccio do modo Seguinte. Toma-se a contradictoria da conclusão negada, compara-se com uma das premissas, e forma-se um syllogismo perfeito, cuja conclusão é contradictoria da outra premissa admitt'da. Assim o adversario deve admittir a verdade de duas contradictorias. Seja o seguinte syllogismo em \emph{Bocardo} --- \emph{algum animal não é racional}; ora \emph{todo o animal é sensitivo} ; \emph{logo algum sensitivo não é racional}. Se alguem concede as premissas mas nega a conclusão, admitte que é verdadeira a seguinte proposição, contradictoria da conclusão negada, --- \emph{todo o seusitico é racional}. Forma-se, então, o seguinte syllogismo: \emph{todo o sensitivo é racional}; \emph{ora todo} o \emph{animal é sensitivo}; \emph{logo todo o animal é racional}. Portanto. quem concede as premissas do primeiro syllogismo mas nega a conclusão, deve admittir a verdade d'estas duas contradictorias --- \emph{algun animal mio é racional}; \emph{iodo o animal é racional}. --- Esta prova de legitimidade vale para todos os modos imperfeitos, sem\pgn{83}excepção. --- Agora vê-se melhor a arte, com que foram compostas \textbf{essas} palavras, que exprimem os modos do syllogismo. E todavia, «não ha hoje um estudante que não saiba rir-se de um \emph{Barbara} ou de um \emph{Baroco}; mas não ha ninguem que não iucline a cabeça deante de um X ou deum X» (\emph{Grofry}).\par Como a proposição \emph{simples}, o syllogismo \emph{simples} divide-se em \emph{absoluto} e \emph{modal}. \emph{É absoluto}, quando as duas premissas são proposições \emph{absolutas}; \emph{é modal}, quando pelo menos uma das premissas é proposição modal. --- Digamos alguma coisa ácerca do syllogismo \emph{modal}; porque ao \emph{absoluto} deve. applicar-se tudo o que convem a9 \emph{simples}. --- O syllogismo \emph{moda}? é triplice ' --- \emph{modal necessario}, \emph{modal-impossivel}, € \emph{modal}. \emph{possivel}, . ou .\emph{contingente}, --- É \emph{modal-necessario}, quando uma ou ambas as premissas são proposições \emph{modacs-ecessarias}, como : --- é necessario que a creatura sja limitada ; ora é necessario que todo o homem seja ereutura logo é necessario que todo o homem seja limitado. --- É \emph{modal-impossivel}, quando uma só premissa é \emph{miodal-impossivel}, como: --- é impossivel que as almas humanas sejam materiaes; ora a alma de Pedro é humana: logo é impossivel que a alma de Pedro seja material. Dizemos --- uma só, porque, se ambas as premissas [orem proposições \emph{modaes-impossiveis}, 0 syllogismo não poderia ter conclusão, sendo \emph{negativa} a proposição \emph{modal-impossicel} (\emph{Kegra} VE \emph{do Syll}.). --- É \emph{snodal-possivel}, ou \emph{contingente}, quando uma ou ambas «as \textbf{premissas} são proposições \emph{modaes-contingentes}, OU \emph{possiveis}, como : --- É possivel que todo 0 homem seja branco; ora é possivel que todo o homem seja musico ; logo é possivel que o musico seja branco, --- Como um juizo \textbf{exacto acerca} do syllogismo \emph{modal} é muitas vezes diflicil, sobretudo porque a indole da conclus\&o depende geralmente da indole da materia exposta, por \textbf{isso}, para O processo se tornar mais simples e seguro, convem reduzir o syllogismo \emph{modal} ao \emph{absoluto}. Assim este syllogisino \emph{wiodal-necessario}: --- é necessario que toda a substancia extensa seja corpo; ora é necessario que toda a pedra seja substancia extensa: logo é necessario que toda a pedra seja corpo, --- reduz-se ao sesuinte \emph{absobdo} em \emph{Barbara}; -- toda a substancia extensa é necessariamente corpo; ora toda a pedra é necessariamente substancia extensa; \emph{logo} toda a pedra é necessariamente corpo.}

ÁrTIGO V

\emph{Syllogismo composto e sua divisao}

%%>>> t95

\pg{81} \tese{95}{Syllogismo composto.}
Syllogismo \emph{composto} é, como dissemos, o que contém nas premissas alguma proposição \emph{composta}, ou é formado de syllogismos \emph{simples}.

%%>>> t96

\tese{96}{Divisão do syllogismo composto.}
Se as premissas contém alguma proposição \emph{composta}, o syllogismo \emph{composto} pode ser \emph{copulativo}, \emph{causal}, \emph{condicional}, \emph{disjunctivo} e \emph{adversalivo}. -- \pg{82} \pg{83}  Se é formado de syllogismos \emph{simples}, o syllogismo \emph{composto} É ou \emph{dilemma}, ou \emph{soriles}, ou \emph{polysyllogismo}. .

%%>>> t97

\tese{97}{Syllogismo copulativo.}
Syllogismo \emph{copulaiico} é O que encerra uma ou ambas as premissas \emph{copulations}. --- Ex.

'l'odos os animaes são sensilivos; \emph{TOTO}

Ora o cão e a lebre e o coelho são animaes :.

Logo o cão e a lebre e \emph{o} coelho são sensitivos.

O syllogismo \emph{copulativo} resolve-se em syllogismos \emph{simples}, e depois estes avaliam-se pelo criterio commum \footnote{Nem sempre o syllogismo \emph{copalativo pode} resolver-se nos syHogzismos * \emph{sêmples}, como no seguinte exemplo: --- quem eumpre e ensina a lei divina é srande \emph{uo} reino dos ceus; ora Pedro ensina a lei divina; loro é grande no reino dos eeus. A rizão é porque o predicado \emph{grande no reino dos} ceus compete a quem \emph{cumpre} e \emph{ensina} a lei, e não a quem se limita a ensinal-q.}.

%%>>> t98

\pg{84} \tese{98}{Syllogismo causal.}
Syllogismo \emph{causal}, chamado tambem, \emph{epicherema}. é o que tem alguma proposição \emph{causal}. --- Ex.

Todo o espirito é immortal :

Ora. a alma humana é espirito, porque é independente da materia;

Logo a alma humana é immortal. .

Na apreciação d'este syllogismo, a razão adduzida avalia-se .Separadamente das proposições (').

%%>>> t99

\tese{99}{SyMogismo condicional.}
Syllogismo \emph{condicional} é aquelle em que a premissa \emph{inaior} é uma proposição \emph{condicional}. --- Conclue de dois modos: --- 1») \emph{affirmando} a \emph{condição} na \emph{menor} e o \emph{condicionado} na \emph{conclusão}; ex. «\emph{Se é homem}, \emph{é animal}; \emph{ora é homem}; \emph{logo é animal}»; --- 2") \emph{negando} o \emph{condicionado} na \emph{menor}, c á \emph{condição} na \emph{conclusão}; ex. « \emph{Se é homem}, \emph{é animal}; \emph{ora não é animal}: \emph{logo não é homem}» O.\footnote{\emph{A} parte \emph{causal} não se refere propriamente á \emph{forma} do syllogismo, mas apenas a proposição. E pode baver syllogismos legitimos na forma eg verdadeiros na materia, que tenham uma razão \emph{falsa}, como no seguinte exemplo: --- \emph{todo o expirito é inunmortal}: \emph{ora} a \emph{alma humana é espirito}, \emph{porque é simples} ; \emph{logo} « \emph{alma}. \emph{humana é inmortal}. A razão, por que a nossa alma é espirito, não é a sua \emph{simplicidade} (tambem a. alma dos brutos é \emph{simples}), \textbar{} mas é a sua independencia da materia.}\footnote{Dizemos que o syllogismo \emph{condicional} conclue de dois modos, Na verdade, \emph{a/firimando a condição} na \emph{menor}, podemos \emph{affirmar} 0 \emph{condicionado} na \emph{conclusão}; porque, estando \emph{o condicionado} contido na \emph{condição}, é claro que, admitida esta, deve sdmittir-se aquelle. Como tambem, \emph{negando} 0 \emph{condicionado} na menor, devemos \emph{negar} a. \emph{condição} na \emph{conclusão}; pois, estando o \emph{condicionado} contido nu, \emph{condição}, é evidente que, não se verificando aquelle, deve concluir-se que se não verificou esta,\par Mas o syllogismo seria illegitimo --- 1.º) se, \emph{affirmando} o \emph{condicionado}, \emph{affirmasse}: tambem a \emph{condição}, pois o \emph{condicionado} não contem a \emph{condição} ; assim seria falso o seguinte syllogismo: --- \emph{se é homem}, \emph{é animal}; mas é , \emph{animal}; \emph{logo é homem}; visto que O \emph{animal} não contem o \emph{hRonem}, e por isso um ente pode ser \emph{animal}, sem ser \emph{homem}: --- 2.º) se, \emph{negando} ux condicio Dà \emph{menor}, \emph{negasse} O \emph{condicionado} na. \emph{conclusão}; pois o \emph{condicionado} pode ser verdadeiro, independentemente da \emph{condição}; por isso € falso o seguinte Syllogismo*: \emph{se é homem}, \emph{é animal}; \emph{mas não é komen z logo não é animal}.\par Todavia, se prescindirmos d:t \emph{forwa logica} do syllogismo \emph{condicional}, e só considerarmos a \emph{mraferia}, podemos ás vezes deduzir da \emph{affirmação} do \emph{condicionado a affirmaczo} da \emph{condição}, e da \emph{negação} da \emph{condição a negação} da \emph{condicionado}, Isto acontece --- ou quando a \emph{condição} se identifica com o \emph{condicionado}, como neste exemplo: se é \emph{homem}, \emph{d animal racional}; mas é \emph{animal racional}; \emph{logo é homem}: --- ou quando a \emph{condição} é a razão unica e exclusiva do \emph{condicionado}, como reste exemplo: \emph{se é homem}, \emph{tem} a \emph{faculdade de rir}; ora não é \emph{homem}; \emph{foto não tem} a. \emph{faculdade de vir}.\par O syllogismo \emph{condicional} reduz-se facilmente ao \emph{simples}, prepondo a\pgn{85}proposição menor e a \emph{conclusão} do syllogismo \emph{condicional} uma. proposição \emph{maior}, que tenha por \emph{sujeito} o \emph{predicado} da \emph{condição}, tomado universalmente, e por \emph{predicado} o \emph{predicado} do \emph{condicionado}. Assim este svilogismo \emph{condicional}: --- se Pedro é homem, Pedro é racional ; ora Pedro é homem ; logo Pedro é racional, --- reduz-se ao seguinte \emph{simples}: --- todo o homem é racional: ora Pedro é homem ; logo Pedro é racional.}

%%>>> t100

\pg{85} \tese{100}{Syllogismo disjunetivo.}
Syllogismo \emph{disjunctivo} é aquelle, em que a \emph{premissa} maior é uma proposição \emph{disjuncliva}. --- Conclue de dois modos: --- 1." \emph{affirmando} uma \emph{parle} da disjuneção na menor e \emph{negando} todas os \emph{outras} na conclusão, ex.: « \emph{Este corpo} vivo é ou \emph{homem}, ou \emph{bruto}, ou \emph{planta}; ora a \emph{homem} ; \emph{logo não é bruto nem planta}: --- 2º) \emph{negando} na menor \emph{todas} as \emph{partes}, \emph{menos uma}, que se \emph{affirma} na \emph{conclusão}, ex.: « \emph{Este corpo vivo} é ou \emph{homem}, \emph{ou bruto}, ou \emph{planta}; \emph{ora não é bruto} nem \emph{planta}; \emph{logo é homem} \footnote{Para que a conclusão seja legitima, a premissa deve conter uma \emph{verdadeira disjuncção}, isto é, as partes da premissa devem excluir-se mutuamente, de modo que, se uma é verdadeira, a outra deve ser falsa, e vice-versa, Assim não seria legitimo este syllogismo: \emph{Lucas} ou é \emph{escriptor}, ou é \emph{pintor}; mas é \emph{eseriptor} ; \emph{loyo não é pintor}; porque a qualidade de escriptor e a de \textbar{}. pintor não se excluem mutuamente.}.

%%>>> t101

\tese{101}{Syllogismo adversativo.}
Syllogismo \emph{adversativo} é o que contém alguma proposição \emph{adversativa}. --- Ex.

Todo o bruto é sensitivo, mas não é racional;

Ora o cordeiro é bruto:

Logo o cordeiro é sensitivo, mas não é racional. Tambem este syllogismo se avalia pelas regras communs \footnote{Ao syllogismo \emph{composto} pertencem o syllogismo \emph{relativo} e O \emph{exponivel}. --- É \emph{relativo}, quando a premissa maior 6. uma proposição relativa; como: --- a nossa alma vale tanto, quanto vale Deus; ora Deus vale infinitamente: logo a nossa alma vale infinitamente. Os termos d' este syllogismo Bão sempre \emph{tres}, como se vê ua seguinte resolução: --- \emph{o} valor da alma humana é o mesmo valor de Deus; ora o valor de Deus é infinito; ; logo \emph{o} valor da alma humana é infinito. --- É \emph{exponivel}, quando uma ou ambas as premissas são proposições \emph{exponioeis}, como: --- \emph{o} bruto só possue a vida sensitiva; ora o homem possue tambem a. vida intellectual: lozo o homem não é bruto. Para se contiecer a legitimidade da conclusão, resolve-se a proposição \emph{expowicel}, e forma-se o seguinte syllogismo: --- o bruto possue \& vida sensitiva, mas não x intellectual: ora o homem possue tambem a vida intellectual ; \emph{logo} o homem não é bruto.}.

%%>>> t102

\tese{102}{Dilemma.}
\emph{Dilemma} é um syllogismo; cuja premissa . \emph{maior} é uma proposição \emph{disjunctiva}, composta: de \emph{duas} partes, ordenadas de forma que, admittida qualquer d'ellas. se conclue sempre contra o adversario. --- É celebre o dilémma de \emph{Tertul}-

\textasciicircum{}

\emph{liano}, demonstrando a injustiça da lei, \textbf{promulgada} por \emph{Trajano}, o qual ordenara que os christãos não \textbf{fossem inquiridos}, mas que fossem condemnados, se alguem os \textbf{denunciasse}:

Os chrislãos ou são reus, ou innocenles ;

Se são reus, porque probibes que sejum inquiridos?

Se são innocentes, porque os condemnas, quando denunciados *?

O \emph{dilemma} deve satisfazer ás tres condições \textbf{seguintes}: --- 1?) a \emph{disjuncção deve ser completa}; --- 2º) o \emph{nexo entre} a \emph{conclusão} e \emph{as parles da disiuncção deve ser necessario}; --- 3º) \emph{deve ser impossivel} a \emph{retorção}, de modo que não possa \textbf{tirar-se} uma segunda conclusão, inteiramente opposta \& primeira (").

%%>>> t103

\pg{86} \tese{103}{Sorites.}
\emph{Sorites} é um syllogismo formado de varias\footnote{Se os membros, da proposicio- \emph{disjwnctiva} são \emph{tres}, ou \emph{quatro}, ele. este syllogismo chama-se \emph{frilenmna}, \emph{quadrilemma}, etc. --- \textbf{Como} se vé, o \emph{dilemma}, --- 1.º) encerra, explicita ou \emph{implicitamente}, \emph{quatro} proposições: ii \emph{primeira} é \emph{disjunctiva}; \emph{a segunda} \& uma \emph{condicional-causal}, porque é constituida pela \emph{primeira} parle da \emph{düsfencecio}, e A ella conetuve contra o adversario, dando a razão da conclusão: x \emph{terceira} tambem € \emph{condicional-causal}, e com a segunda \emph{parie} da \emph{disjunceão} impugna o adversario; a \emph{quarta}. affirma \emph{categoricamente} do \emph{todo} o que, de um modo \emph{hypothetico}, foi affirmado das \emph{partes}; assim a conclusão do citado dilemma de "Tertulliano é a seguinte: --- logo a lei \emph{por} té promulgada é sempre c em todo o caso injusta; --- 8,0) contém em si dois \emph{entymemas}, como demonstra o exemplo dado; --- 3.º) differe do syllogismo \emph{disjfunctivo}, porque este, admittida uma parte, exclue a outra, \emph{o} vice-versa; ao passo que o \emph{diesona} deduz de cada uma das partes a mesma verdade contra o adversario.\par Mas, para que o \emph{dilenuma} couclua:efficazmente, devem observar-se as tres condições enumeradas no texto: --- 1.2) 4 \emph{disjuneção deve ser completa} ; iliás o adversario encontrará uma porta de evasão. Assim é Wlegitimo o \emph{dilemma} de Cicero (De \emph{Sem}, XIX), quando, querendo \textbf{demonstrar} que não devemos ter medo da morte, disse que a nossa alma --- ou \textbf{será} eternamente feliz, --- ou será mortal, e não se lembrou de -que a \textbf{nossa alma} pode ser tambem eternamente infeliz. --- 9.2) \emph{O nexo entre a conclusão} € us \emph{partes da disjuncção deve ser necessário}; aliás o adversario tirar-se-ha de todo o embaraço, negando a \emph{necessidade} da conclusão, Assim vicioso é \emph{o} argu-" mento de um philosopho, quando, querendo convencer um dos \textbf{seus} discipulos de que \emph{não revia acceitar} um emprero publico, disse-lhe: «Ou tu cumprirás com os teus deveres, ou não: no primeiro caso desagradarás aos homens ; no segundo, 2 \emph{Deus}». Como viciosa é a resposta do discipulo: «no primeiro caso, agradavel a Deus: no segundo, aos homens». Porquanto \emph{falta} O \emph{nexo} necessario entre a. boa administração e o desagrado dos homens, como não tem egual importancia o desagrado de Deus e o \textbf{desagrado} dos homens. --- 3.º) \emph{Deve ser impossivel a retorção}; pois é pouco honroso ser ferido pelas proprias armas. Falla tambhen a esta condição o \emph{dilemma} do \emph{tal philosopho}, \emph{que} queria levar o diseipulo a não acceitar eargos publicos.}

- CAPITULO TERCEIRO --- RACI!

proposições, dispostas de modo, que o \emph{predicado} da primeira é o \emph{sujeito} da \emph{segunda}, o \emph{predicado} da \emph{segunda} é o \emph{sujeito} da \emph{terceira}, e assim por deante até se chegar a \emph{conclusão}, formada do \emph{sujeito da primeira} e do \emph{predicado} da ultima. --- Ex.:

A alma humana é um ser capaz de formar ideas abstractas:

O ser capaz de formar ideas abstractas é espiritual :

O ser espiritual é incorruptivel :

O ser incorruptivel é naturalmente immortal ;

Logo a alma humana é naturalmente immortal.

O \emph{soriles} resolve-se em tantos syllogismos \emph{simples} da \emph{primeira} figura, quantas são as proposições, menos duas, ou quantos são os \emph{termos medios}. Assim o \emph{sorites} proposto resolve-se nos tres seguintes syllogismos :

O ser capaz de formar ideas abstractas é espiritual ;

Ora a alma humana é um ser capaz de formar ideas ahstractas :

Logo a alma humana é espiritual.

O ser espiritual é incortuptivel: .

Ora a alma humana é um ser espiritual :

Logo a alma humana é incorruptivel.

O ser incorruptivel é naturalmente immortal ;

Ora a alma humana é um ser incorruptivel:

Logo a alma humana é naturalmente immortal \footnote{\pgn{88}vem \emph{a}. espiritualidade, \& incorruüptibilidade, e a. immortalidade. Logo essus prerogativas devem tambem convir € convém i alma humana. --- Note-se que os syllogismos \emph{simples}, em que se .resolveu \emph{o} citado \emph{sorites}, são \emph{tres}, quantas são as proposições, menos duas, ou quantos são os \emph{termos medios} (que, no exemplo citado, são --- \emph{ser capaz de formar ideas abstractas}, --- \emph{ser espiritual}, --- \emph{ser incorruptivel}). .\par Mas para que o \emph{sorifes} seja uma. forma legitima de.argumenlacüo, deve satisfazer ás condições apontadas no texto. --- Antes de tudo, é necessario que se não introduza nenhum termo \emph{equivoco}, porque, sendo este termo equivalente a \emph{dois termos}, já não pode subsistir a estructura do syllogismo. --- Alem d' isso, todas as proposições devem ser \emph{absolutamente} verdadeiras e ligadas entre si: aliás sahirüo as conclusões mais desarrazoadas, que se podem imaginar, como neste exemplo: «quem bebe muito, derme bem; quem dorme Dem. não peeca; quem não pecca, é santo ;'logo quem bebe muito é santo». -- Mais: todas as proposiqóes sejam \emph{wnzcersaes}, exceptuada a \emph{primeira}, que pode ser \emph{particular}, A razão é porque. \emph{todas} as premissas do \emph{sorites}, á excepção da \emph{primeira}, representam as premissas \emph{maiores} dos syllogismos. da \emph{primeira} figura, nos quaes o proprio \emph{sorifes} se resolve; ova na \emph{primeira} futura a \emph{mcior} deve ser \emph{universal}. Mas se a \emph{primeira} premissa é \emph{particular}, não ha infracção de nenhuma regra, porque essa premissa, na resolução, occupa o logar da \emph{menor} e só \textbf{exige} que tambem \emph{particular} seja a conclusão. --- Finalmente, todas as proposições sejam a \emph{ffirmativas}. Na verdade, se \emph{todas} fossem \emph{negativas}, resolver-se-hia o \emph{Sorites} em syllogismos, compostos exclusivamente de proposições \emph{negativas}, e não haveria conclusão alguma. Se. fosse \emph{negativa a primeira}, como esta, na resolução, constitue a \emph{mener} do \emph{primeiro} syllogismo, teriamos um syllogismo da \emph{primeira} figura com a menor \emph{negativa}; o que é contrario \& lei d essa ligura, Se fosse \emph{negativa} qualquer outra: premissa, teriamos, na resolução, uma conclusão tambem \emph{negativa}; e, como esta seria a menor de um outro syllogismo, haveria tambem a infracção \emph{da} mesma lei da \emph{primeira} figura, Mas, Se a \emph{penulfime} proposição for \emph{negativa}, não deriva d'ahi inconveniente algum, comtanto que seja \emph{negativa} lambem a \emph{conclusão}. --- É muito perigosa esta especie de argumentação, (que, para \emph{Pracy}, todavia, é a mais simples e genuina); porque a \emph{rapida} successão das proposições não deixa a inteligencia o tempo necessario para reflectir, e para ver onde pode esconder-se \emph{o} sophisma. Por isso, se o adversario recorre ao \emph{sorites}, é necessario resolvel-o em syllogismos simples. conforme o conselho do \emph{Cicero}: « \emph{Vitiosi sunt soriles}: \emph{frangile igitur eos}, \emph{si potestis}, \emph{ne molesti}. \emph{siut}». (Acad., 1. IT).}.

Para a verdade do \emph{soriles} é necessario --- 1." que não haja Lermos \emph{eguivocos}, --- 2") que \emph{todas} as proposições sejam \emph{absolutamente} verdadeiras, e que uma derive necessariamente da outra, --- 3\textasciicircum{}) que \emph{todas} as proposições sejam universaes (só a \emph{primeira} pode ser \emph{particular}) e \emph{affirmativas} (só a \emph{penullima} pode ser \emph{negativa}) (). \textbar{}

(13 Portanto -- 1.º) no \emph{sorifes} 0 \emph{primeiro} logar é -occupado pela. premissa \emph{menor} do primeiro syllogismo, o \emph{segundo} pela premissa. \emph{maior} do mesmo syllogismo, o \emph{terceiro}, \emph{o quarto}, et, pela premissa \emph{seaior} dos syllogismos seguintes; --- \%º) no \emph{soritfes} encontram-se as premissas \emph{maiores} de todos os syllogismos, em que se resolve: as premissas menores omittem-sc, á excepção da menor do primeiro syllogismo, como tambem se omitlem us \emph{conclusões}, á excepção da conclusão do umo syllogismo ; --- 3.º) os «\emph{sujeitos} de \emph{todas} as premissas, á excepção do sujeito da \emph{prêmeiro}, fazem o olfieio de \emph{fermo medio} nos syllogismos, que o \emph{sorifes} contém.

(S3 O \emph{sorites} funda-se no seguinte \emph{principio}: \emph{guardo} um \emph{attribuito} se NE cont um \emph{Sujeito}, \emph{dudo o que consen} dO \emph{attributo}. convem \emph{tambem} ao \emph{sujeito}. Assim, no exemplo citado, affirmando-se que + ahna humana é um ser capaz de formar ideas abstractas, tudo Oo que convem a esse Ser, convem tambem A nossa alma, Ora a \emph{um} ser capaz de formar ideas abstractas con-

%%>>> t104

\pg{87} \pg{88} \tese{104}{Polysyllogismo.}
\emph{Polgsyllogismo} é uma argumentação formada de varios syllogismos, dispostos de modo, que a \emph{conclusão} do \emph{primeiro é premissa} do \emph{segundo}, e assim por deante. Ex.

Tudo \emph{o} que é independente da materia é espiritual ;

Ora o alma humana é independente da materia ;

Logo a alma humana é espiritual ;

Ora \emph{a}. que é espiritual é immortal; *

Logo a alma humana é immortal.

%%>>> art-7

\section*{Artigo VI.}
\addcontentsline{toc}{section}{Artigo VI.}

\pg{89} \emph{Inducgão e suas especies}

%%>>> t105

\tese{105}{Inducção.}
\emph{Inducção} É uma argumentação, em..que se altribue a um \emph{sujeito universal} ou \emph{geral} um predicado, que, pela \emph{experiencia}, conhecemos convir aos seres \emph{particulares}, contidos na extensão d'aquelle \emph{sujeito}. Ex.

O vão, o cordeiro, o le\&o... são sensilivos;

Ora o cão, o cordeiro, o leão... são brutos:

Logo Lodos os brutos são sensitivos,

%%>>> t106

\tese{106}{Divisão da inducção.}
A \emph{inducezxo} divide-se em \emph{completa} e \emph{incompleta}, --- E \emph{completa}, quando attribuimos ao \emph{todo} um predicado, que conhecemos convir a \emph{cada} uma das suas partes. Esta especie de \emph{inducção} não tem valor algum scientifico e por isso não nos oecupamos d'ella. --- É \emph{incompleta}, quando attribuimos ao \emph{todo} um predicado, que conliecemos convir a \emph{algumas} das suas parles. Esta segunda especie, que é propriamente \emph{inducção}, é um meio demonstrativo de novas verdades \footnote{Alguns escriptores allirmanm que a inducção \emph{incompleta} eva desconhecida pelos antigos, e que foi inventada. por \emph{Bacen} no sec, XVII. --- Esta\par são é falsa. É certo que \emph{Bacon} ensinou a maneira de fazer «so util e asse legitimo da \emph{inducção}, por meio de algumas regras, que elle coltigiu e que nós indicaremos mais tarde, Mas os philosophos antigos, principalmente \emph{Aristoteles} c os \emph{KEscholasticos} da edade media. sobretudo \emph{Santo Thomaz de} Aquino, já conheciam a inducção \emph{incompleta}; pois que fallam, com muita precisão e profundidade, da sua natureza e vantagens, e da experimentação, em que se basea a propria inducção. Se não empregaram mais frequentemente u inducção, a razão foi a falta dos instrumentos de observação, muito necessarios para evitar o erro neste zenero de artumentação, ou porque, sendo quasi todos ecclesiasticos os estudiosos d'esses tempos, preferiam \& investigação do sobresensivel a do sensivel.\par Alguns escriptores dizem \emph{perfeita} Qu incperfeita a. inducqiüo, que nós. chamamos \emph{completa} ou \emph{incompleta}. --- Mas essa divisão em \emph{perfeita} e \emph{imperfeita} convem á induec\&o \emph{incompleta}, conforme o numero dos particulares enumerados € \emph{sufficiente}, ou não, para a conclusão. A induec\&o \emph{incompleta} é \emph{perfeita}, quando leva á \emph{certeza}. de (que ao \emph{todo} convem um certo predicado, embora esse predicado tenha sido observado num só \emph{particular}, e leva á \emph{certeza}, quando se conhece que o lul predicado deriva \emph{da} propria valureza ou essencia. do mesmo particular. --- Q valor \emph{logico} e \emph{demonstrativo} da inducção \emph{incompleta} não depende, como veremos, do manero ou da \emph{quantidade} dos phenomenos, mas du sua \emph{qualidade}, \emph{emquanto} revelam a essencia, da qual derivam. (CL S. Thom, \emph{ia Hb}. \emph{IT Post}., \emph{lect}. XX).\pgn{90}\emph{al C}.},

107. Legitimidade da inducção incompleta, --- A passagem, \pg{90}  - que, na énduccüo \emph{incompleta}, se faz do \emph{particular} para o \emph{geral}, Li não é arbitraria. Por quanto, averiguando-se, mediante repetidas E experiencias, que uma \emph{propriedade} ou \emph{operação} deriva da essencia = de muitos individuos, rectamente .concluimos que a mesma e \emph{propriedade} ou \emph{operação} deve encontrar-se em \emph{todos} os outros \textasciicircum{} individuos, dotados da mesma essencia. Assim, tendo reconlie- \emph{MV} M Jj . H \emph{E} p cido que a \emph{extensão}, que se observa em \emph{muitos} corpos, \& uma --- propriedade que deriva da \emph{essencia} d'elles, com rectidão con- \emph{D} dui " I OYhOR SS soe (1 \textbar{} : \emph{Ll} cluimos que todos os corpos são \emph{extensos} (!).

%%>>> t108

\tese{108}{Processo da inducção incompleta.}
São duas, pois, UT as principaes operações, que constituem a inducção \emph{incompleta}: m a) Deve averiguar-se, pela experiência, se o nexo entre um \emph{OO} facto \emph{particular} e a sua causa é essencial ou \emph{accidental} (?). ---- . (UT (1) A \emph{inducção} é um verdadeiro raciocinio, como o é o \emph{syllogismo} ; pore que ambos esses processos --- passam do \emph{conhecido} para o \emph{desconhecido}, --- " --- haseam-se no mesmo principio de \emph{identidade}, --- contêm explicita ou \emph{imo} plicitamente tres proposições, --- e têm o mesmo fim, que é a \emph{certeza}. UM Segundo \emph{Bacon}, \emph{Wolff}, \emph{Tiberghien}. \emph{J}. \emph{Stuart Mill} e os \emph{positavistas}, as á proposições geraes, que são o resultado dà imducção, não são \emph{certas}, mas \emph{t} sómente \emph{provaveis}. --- Respondemos que, para chegarem a esta conclusão, --- era necessario que estes escriptores demonstrassem que os seguintes juizos \emph{D}\textasciicircum{} jnductivos: --- «\emph{todo o corpo é exiqnso}» --- «\emph{o homem}, \emph{cortada} a cabeça, \emph{nd não vice} » --- «\emph{ao dia succede sempre} a \emph{notte}» --- não são \emph{certos}, mas sómente \emph{So provaveis}. Parecem-nos pouco coherentes os \emph{positivistas}, que, depois de ter x elogiado tanto a \emph{inducção}, como o \emph{unico} meio para alcançar a certeza, acu- \emph{I}» . bam por dizer que os juizos \emph{énduciivos} são sómente \emph{provaveis}, Para elles, ES pois, a inducção é, e não é, um meio para a acquisição da \emph{certeza}! --- (2) O \emph{meio}, de que a inducção se serve para conhecer se um \emph{facto}, \emph{phe}- 27 \emph{nomeno}, ou \emph{propriedade}, não depende de circumstancias \emph{accidentaes}, mas = deriva da propria \emph{essencia} de um certo e determinado ente, é a \emph{experiencia}. \emph{D} --- Digamos alguma coisa ácerca da \emph{experiencia} e suas especies. \emph{T}7 A) Experiencia. --- A \emph{experiencia} é um conhecimento adquirido por meio \emph{o} dos \emph{sentidos}, \emph{internos} ou \emph{externos}, ou por meio da \emph{consciencia}, E SEN Divide-se em \emph{observação} e \emph{experimentação}.

a) Observação. ---- Observação é a percepção de um phenomeno, que sc :

Oflerece aos nossos sentidos \emph{espontaneamente} e Sem a nossa cooperação.

a) Requisitos da observação. --- Para a lezitimidade da \emph{observação} é neces-

sario --- Í.") que a percepção seja intensa, perseverante \emph{e} livre de tolo \emph{o} :

preconceito : --- 2.*, que o phenomeno se considere no set! conjuncto e em \emph{ES} 3 cada uma das suas partes, e este exame dê sempre o \emph{mesino} resultado; --- 2\emph{i} a 3.º) que se determinem todas as propriedades e cireumstancias especiaes de \emph{E}. um phenomeno \emph{e} as condições, em que elle se realiza. s po. 4 : " . » - . t b) Experimentação. --- \emph{Experimentação} É a operação, \emph{pela} qual reproduzi- NE mos um phenomeno com a nossa industria e arte, A natureza não apresenta n \textasciicircum{} * ---

b) Se o nexo é \emph{essencial}, o facto deve \textbf{elevar-se} a \textbf{expressão}. de uma \emph{lei universal}, de modo que possa \textbf{applicar-se} a \emph{todos} os factos \emph{particulares}, em todos os tempos e em todos os \textbf{logares} \footnote{A inducção \emph{scientifica}, pois, não consiste nà accumwulaczo das experiencias. Uma experiencia \emph{Dmminosa}. em que o espirito descubra claramente a relação entre um phenomeno e a Sua causa, vale mais do que mil experieneins ebscuras, nem estas accrescentam força áquella. A \emph{inducção} consiste em despir um phenomeno de todas as circumsianciàs. que o cercam, considera-lo na sua essencia ou natureza, e \emph{relaciona-lo} com uma causa proporcionada. Quando esta \emph{relação} apparece, não são necessarias ulteriores observações e experiencias, Podemos dizer que a inducção é uma \emph{iniuipzo} do espirito, que num tacto \emph{singular} e \emph{concreto} descobre e põe em relevo uma relação \emph{universal} --- \emph{Cl}. \emph{Bernard} diz: (\emph{Plienomenos da} vida). «Não é pela aceumulação dos factos que a sciencia progride, \emph{mas} sim \emph{pela Iuz da tdeu}\textgreater{}».\par Agora mais claramente se vé o erro de \emph{Stuart Mill}, que, para demonstrar que o \emph{sy/ogismo} é um \emph{circio}. \emph{vicioso}, diz que a \emph{premissa} maior do mesmo syllogismo recebe a sua força \emph{da conchisão}. Por quanto, essa \emph{premissa maior}, embora seja um \emph{juízo synéhetico} € por isso resulte da iínducção, comurdo não recebe da \emph{eeachtsaüo} a sua força: pois o \emph{geral}, a que chega a \emph{inducrão}, não representa a somma de todos os \emph{particulares} Por isso, O syllorismo é sempre um meio de \emph{descobrir} novas verdades, mesmo quando a sua \emph{premissa inaiíor} é O resultado da \emph{inducoro}.\par *}.

%%>>> t109

\tese{109}{Principio da inducção incompleta...}
\textbf{Para}. que. um \pg{91} \pg{92}  facto \emph{particular} possa elevar-se a \emph{lei universal}, é necessario applicar-lbe um principio \emph{analytico}, que lhe dé a propria universalidade; d'outro modo teriamos um effeito sem causa proporcionada. --- Esse principio \emph{analytico} é o seguinte: « \emph{Uma propriedade}, \emph{que} se \emph{encontra} num \emph{ente} e \emph{que deriva da} sua \emph{essencia deve encontrar-se} em \emph{todos os entes}, \emph{dotados da mesma essencia} ». Portanto, o principio da induegüo \emph{incompleta} é uma applicação do \emph{principio} de \emph{causalidade} \footnote{Como se vê, o principio, em que se funda a inducção, equivale ao principio, em que se funda o \emph{syllogismo}, nas na \emph{ordem inversa}. O prin cipio do \emph{sy/Togisito} pode exprimir-se do modo seguinte: \emph{deve} a \emph{ffirmar-se ou negar-se de todas} as \emph{especies de} um \emph{genero}, ou \emph{de todos os individuos de uma especie}, \emph{o} que se \emph{affirmox}, ox \emph{negou do} mesmo \emph{genero}, ou \emph{da mesnia especie}. O principio da \emph{inducção} pode tambem formularse assim: \emph{deve a ffirmar-se ou negar-se de um genero}, ou de uma \emph{especie}. \emph{o que se affirmow ou negoit} de \emph{todas} as \emph{especies d}' \emph{esse genero}, ou \emph{de todos} os \emph{individuos d}' \emph{essa especie}. --- O principio da inducção, como o do \emph{Syllogismo}, reduz-se ao principio de \emph{causalidade}, que é: a \emph{mesma} causa \emph{natural produz o mesmo effeito}, € portanto: o mesmo \emph{effeito natural exije} e \emph{suppõe} \& mesma causa.\par Dizem alguns philosophos, que o principio da \emph{inducção incompleta} 6 0 seguinte: \emph{as leis da natureza são constantes}. --- Mas erram. Por quanto, este juízo, sendo \emph{synthetico}, não é o \emph{principio} da \emph{inducção}, mas sim é uma \emph{consequencia} d' ella. Como sabemos nós que as leis da natureza são constantes? Pela. experiencia ? Não basta. A experiencia poderá attestar que um phenomeno se repete uma, duas, tres .. vezes do mesmo modo, mas não attesta que se repete sempre e em toda x parte e do mesmo modo, de forma (ue possa tirar-se esta conclusão geral : as \emph{leis da natureza são constantes}. É necessario, pois, basear a inducção \emph{incompleta} num juizo, não \emph{synthetico}, mas \emph{enalytico}, que. por ser \emph{absoluto}, tem a, maxima \emph{generalidade},\par \emph{J}. \emph{Stuart}. \emph{Mill} (\emph{Systema}. \emph{da Logica}). tendo dito que a marcha natural do espirito humano é inferir o \emph{particular} de outro \emph{particular}, allirma que + inducção É a «;reneralisacüo da experiencia »,.c que o axioma que lhe serve de base é \emph{o} seguinte: «o curso da satuwreza. \emph{E uniforme}»,\par Vejamos as contradicções do escriptor inglez. --- Se o espirito humano só póde inferir o \emph{particular} de. outro \emph{particular}, como póde dizer-se agora que elle passa do \emph{parficulor} para O \emph{geral}, e que applica a £\emph{ódos} os casos o que se encontrou só num? --- segundo \emph{Start AH}, passamos do \emph{partietear para} o \emph{geral} em virtude do axioma: \emph{o curso da natureza é uniforme}, Bem. Mas poderá um positivista, sem se contradizer, admittir a \emph{imiiformidade} do curso da natureza? \emph{Mi}! só reconhece como certo \emph{o} que Ihe refere a exporleneia. Ora a experiencia poderá attestar-lhe aquella uniformidade em algumas circumstancias \emph{particulares}, mas nunca em \emph{foras} as cireumstancias, em \emph{todos os} tempos, e em \emph{todas} as parles, Por isso tambem o axioma ---\pgn{93}\emph{o curso da natureza é uniforme} --- é já um \emph{resultado} da \emph{indução} incompleta, e suppõe um outro principio, de que recebeu a generalizução. Este outro principio, em ultima analyse, não póde ser por sua vez um juizo \emph{synthketico}, porque reappareceria a diificuldade, mas deve ser um juizo analytico. . Em summa, o principio da \emph{inducção} deve ser necessariamente um juizo \emph{analytico}. Os \emph{positivistas} ou admitlem os juizos \emph{analgticos}, ou não, Se os admittem, contradizem-se, porque o seu systema só reconhece juizos \emph{syatheticos}. Se os não admittem, então tiram toda a força e legitimidade 3 \emph{inducção}: pois a reduzem a um sophisma, qual é a passagem do \emph{particular} para o: \emph{geral}, sem um principio que a justifique.}. \emph{I}

%%>>> t110

\tese{110}{Differenca entre o Syllogismo e a indnecàáo incompleta.}
A differença entre estas duas especies de argumentação \pg{93}  está em que o \emph{syllogismo} conclue do \emph{universal} ou do \emph{geral} para o \emph{particular}. ou do \emph{todo} para a \emph{parte}, emquanto a \emph{inducção incompleta} conclne do \emph{particular} para o \emph{geral}, ou da \emph{parte} para o \emph{todo}. --- \textbar{} \textbar{} o. ol

a) No \emph{syllogisino}, pois, o \emph{sujeito} da \emph{conclusão} é uma \emph{parte} e o termo \emph{medio} é um \emph{todo}. --- Na \emph{inducção} o \emph{sujeito} da \emph{conclusão} é um \emph{todo} e o termo \emph{medio} é o complexo das \emph{partes},

Án que formam aquelle \emph{todo}.

b) No \emph{syllogismo}, a premissa maior affirma que o predicado

da \emph{conclusão} convem ao termo \emph{medio}; a premissa \emph{menor}, que o \emph{sujeito} está contido no termo \emph{medio}; a \emph{conclusão} que ao \emph{sujeito} convem o \emph{predicado}, que foi attribuido ao termo \emph{medio}. --- Na inducção, a premissa \emph{maior} aflirma que o \emph{predicado} da \emph{conclusão} convem ás \emph{parles}, que formam o termo \emph{medio}; a premissa menor, que o complexo das partes constitue o (\emph{odo}; a \emph{conclusão}, que ao \emph{todo} convem o \emph{predicado} attribuido. ás .\emph{partes}.. . o

D'ahi se deduz: --- \emph{a}) o termo, que. no. \emph{syllogismo é}. \emph{medio},

na \emph{inducção é extremo} menor; e vice-versa: --- \emph{b}) a \emph{proposição}, que no \emph{syllogismo é conclusão}, na \emph{inducção} é premissa \emph{maior};

\emph{A} e viceversa. \textbar{} \textbar{}

O que levamos dito, melhor se vê nos seguintes exemplos: (Fx. \emph{do syllogismo}). --- Todos os brutos são sensitivos 3

Ora o cão, o cordeiro, \emph{o} leão... são brutos:

Logo o cão, o cordeiro, o leão... são sensitivos. (\emph{Ex}. \emph{da induccio)y}. --- O cão, o cordeiro, o leão... são sensitivos:

Ora o cão, o cordeiro, o leão... são brulos;

Logo todos os brutos são sensilivos \footnote{I5 necessario accenLuar as differengas entre o \emph{syllogisino} e a \emph{énducgao}.\par O \emph{syllogismo} desce do universal para o \emph{particular}. do \emph{todo} para a \emph{parte}. das \emph{leis}, para os \emph{factos} e \emph{pluenomenos}, da causa para os \emph{effeitos}, do \emph{necessario} pira o \emph{contingente}, ele: A \emph{inducção}, pelo contrario, sote \emph{do particular} para\pgn{94}O \emph{universal} ou \emph{geral}, da \emph{parte} para o \emph{todo}, dos \emph{factos} e \emph{phenomenos} para as \emph{leis}, dos \emph{effeitos} para a causa, do \emph{contingente} para o \emph{necessario}. É por isso que o \emph{syllogismo} nos mostra o \emph{particular} contido no \emph{geral}, ao passo.que a \emph{inducção} nos manilesta. o \emph{geral} contido no \emph{particular}. --- Não ha illacüo ou consequencia, a que não desça o \emph{syllogismo}; como não ha altura ou generalidade, a que se não eleve a \emph{inducção}. --- No \emph{syllogismo}, em que deduzimos do \emph{universal}! o particular (pois 0 \emph{universal} extende-se ao \emph{particular}), considera-se principalmente a extensão das ideas; --- na \emph{inducção}, em que concluimos do \emph{particular} para o \emph{universal} (pois o \emph{geral} ou \emph{universal} tem a. mesma natureza e obedece ás mesmas leis que o \emph{particular}), considera-se sobre tudo a \emph{comprehkensao} das ideas. E\par Embora distinetlas no seu processo, estas duas formas do raciocinio vão sempre unidas, --- Não ha \emph{inducção}, propria e rigorosamente dicta, que não se funde no \emph{syllogismo}. Porquanto, a \emph{induecção} conclue que ao \emph{todo} convem o que se verificou na \emph{parte}. Ora esta passagem da \emph{parte} para o \emph{torto} baséa-se na força d'este syllogismo: s mesmas causas \emph{physicas} produzem \emph{os mesmos effeitos}; ora algumas causas \emph{desconhecidas tem} a \emph{mesma natureza das outras já conhecidas}; \emph{logo} essas causas \emph{desconhecidas hão de produzir os mesnmos effeitos}. --- Se a inducção suppõe o \emph{syllogismo}. este suppõe a \emph{inducção}, ao menos emquanto esta significa toda e qualquer \emph{generalização}, embora feita \emph{imediatamente} e \emph{sem raciocinio}; pois o nosso espirito eleva-se ao \emph{universal} ou \emph{geral} por meio do \emph{particular}. É neste sentido que Aristoteles diz: Não existe \emph{universal} Sem \emph{inducção}, nem \emph{inducrão} sem os \emph{sentidos} (An. \emph{post}., 1, 18). S. Thomas aecrescenta: «Si \emph{wnicersalia}, ex quibus procedit demonstratio, cognosci possent absque \emph{inductione}, sequeretur quod homo posset accipere scienliam eorum, quorum non habet sensum, \emph{Sed impossibile ext universalia speculari absque induclione} » (\emph{Post}. An., I. 30).\par Mas qual é o primeiro processo da nossa razão --- o \emph{syllogismo} ou a. \emph{inducpzo}? --- Respondemos que é a inducção, emquanto significa toda e qualquer \emph{generalisação}; porque, como dissemos, só percebemos o universal pelo \emph{particular}. --- Dor isso, o nosso espirito conhece primeiramente o \emph{particular}. Do \emph{particular}. cleva-se, imediatamente” e sem o auxilio do \emph{raciocinio}, ao \emph{universal} (inducção \emph{vulgar}). Em seguida, do \emph{universal} ou \emph{geral}, estabelecido como \emph{principio}, desce aos \emph{particulares} (syllogismo). Finalmente auxiliado pelo \emph{syllogismo}, forma uma argumentação, em que do \emph{particular} sobe ao \emph{universal} (inducção \emph{scientifica}).\par D'ahi se vê melhor a utilidade do \emph{syllogismo}. Porquanto, se a sciencia pela \emph{inducção} procura factos e descobre principios, por meio do \emph{ayllogismo} demonstra os mesmos principios, e organiza-se e completa-se pelo conhecimento das consequencias nelles contidas.}.

111. Corollario --- Sendo a \emph{inducção} tão distineta e diversa do \emph{syllogismo}, é claro que aquella não pode reduzir-se a este:

« pois que a conclusão do \emph{syllogismo} é essencialmente uma proposição \emph{particular}, ao passo que a conclusão da \emph{isducczo} é sempre uma proposição \emph{universal} \footnote{A \emph{inducção} reduzem-se a \emph{avalogia} e \emph{x hypothese}.\par 1) Analogia. --- \emph{Analogia} é a argumentação, peli qual se convede a\pgn{95}um \emph{sujeito particular} um predicado, concedido já a um. \emph{sujeito} par \emph{tieular}, semelhante ao \emph{primeiro}. --- Exprime-se do seguinte modo :\par Causas \emph{semelhantes produzem effeitos semelhantes} ;\par \emph{Mas} A e \emph{DB são} causas \emph{semelhantes} ;\par \emph{Logo} A e \emph{B producem effeitos semelhantes} ;\par \emph{Ora} a causa A \emph{produz o ejfeito C}:\par \emph{Logo} a causa B \emph{deve produzir o ejfeito} C. -\par \emph{a}) Relações, entre a inducção e a analogia.\par 11º) Ambas parlem do mesmo punto --- o \emph{parfécilar}; mas não chegam ao mesmo termo, porque a \emph{insdueção} chega ao \emph{geral}, 10 passo que a \emph{analogia}. sómente chega ao \emph{particular}.\par 2,0) Na \emph{inducção} attribuimos ao \emph{todo} um predicado. que necessariamente \emph{convem} a \emph{cada} uma das partes, e é impossivel o erro; mas nº \emph{analogia} attribuimos a ma coisa um predicado, pelo facto d'este convir a uma outra Coisa \emph{semelhante} á primeira. Por isso é facil \emph{o} erro na \emph{analogia}, porque “o predicado póde convir a uma coisa em razão do elemento \emph{differenciat} e \emph{proprio}, e não \emph{por causa} do elemento \emph{generico} e commum.\par 4) Vantagens da analogia, --- São importantissimas as vantagens, que da\par \$\pgn{96}\emph{analogia} tiram as sciencias naturaes e principalmente aquellas artes, que se fundam na imitação, como são a-eloquencia, a poesia, cto.\par b) Hypothese. --- \emph{IIypothese} é a argumentação, pela qual se considera como verdadeira uma proposição (\emph{huypolkese}), cuja verdade, é desconhecida mas que explica claramente alguns factos.\par Convem com a \emph{analogia}, em quanto ambas concluem do \emph{particular} para o \emph{particular}, e os seus resultados não excedem os limites da probabilidade, “--- \emph{A hypothese} é de grande utilidade, sobretudo nas sciencias experimentaes, e exprime-se com a formula seguinte:\par \emph{Se existisse} a causa A, \emph{dar-sehiam os factos b}, \emph{e}, \emph{d};\par «\emph{Ora os factos b}, c, \emph{d}, \emph{dão-se realmente};\par \emph{Logo} a causa A \emph{póde}. \emph{existir}.\par Esta conclusão é sómente pravavel. Se, porém, forem examinadas todas 2s hypotheses possiveis, e se provar que todas ellas, á excepção de uma, são inadmissiveis, então a excepiuada é certa.\par a) Divisio da hypothese. --- A Dypolhese é \emph{plarsica}, moral e \emph{logica}.\par 1.9) É \emph{physica}. quando se emprega para explicar os factos da aatureza, que dependem immediatamente das leis physicas.\par 2.º) É \emph{moral}, quando serve para explicar os factos, que dependem da liberdade dos seres intelligêntes; assim pela acções externas dos homens conhecemos os seus sentimentos internos.\par 3.º) É \emph{logica}, quando, na interpretação de palavras obscuras, propomos uma opinião, que nos parece conforme com o pensamento de quem laliou.\par \emph{b}) Requisitos da hypetese. --- A \emph{hugpothese} --- 1.9) deve ser \emph{possivel}; --- 2.º) nem em si, nem nas suas deducções deve \emph{repugnar a algum facto certo}; --- \$.") deve explicar todos os fuclos, para que se emprega,}.

\emph{Sophisma e suas especies}

%%>>> t112

\pg{94} \pg{95} \tese{112}{Sophisma.}
\emph{Sophisma} é uma argumentação, que parece verdadeira, mas na realidade é falsa. s

%%>>> t113

\tese{113}{Divisão do sophisma.}
O \emph{sophisma} divide-se em \emph{gramatical} e agito. --- O \emph{grammatical} consiste no abuso das

\emph{palavras}; o \emph{logico}, no abuso das \emph{ideas}.

%%>>> t114

\tese{114}{Especies do sophisma grammatical.}
As principaes especies do sophisma \emph{grammwuttical} são: a \emph{amphibologia}, a \emph{composição}, a \emph{divisão}, a \emph{extensão} e a \emph{restricgao}.

a) Amphibologia. --- A \emph{amphibologia} dá se, quando uma palavra se toma em dois sentidos diversos; assim o \emph{syllogismo} tem \emph{quatro} termos. --- Hefuta-se, distinguindo o sentido.

b) Composição. --- O sopbisma de \emph{composição} commette-se, quando a um \emph{sujeito} se attribuem. \emph{simultaneamente} dois .predi-. ..:cados, que só lhe podem. convir \emph{successivamente}. Se a proposição --- os \emph{cegos} vêm --- se toma no sentido que os homens; emquanto \emph{cegos}, vêm, é um sophisma de \emph{composição}.

€) Divisão. --- O sophisma de \emph{divisão} commette-se, quando a um sujeito se attribuem \emph{successivamente} dois predicados, que só lhe podem competir \emph{simultaneamente}. Se esta proposição --- o \emph{homem que é facinora nio pode} alcançar \emph{glorit} --- se toma no sentido que um homem, que é \emph{facinora}, nunca poderá \pg{96}  alcançar gloria, nem mesmo depois de deixar de ser \emph{facinora}, é um sophisma de \emph{divisio}. --- Este sophismattambem se refula com uma distincção.

d) Extensão. --- E sophisma de \emph{extensão}, quando de um sentido \emph{particular} se passa para um sentido \emph{universal}. Ex.

Todo o sabio merece louvor ;

Ora este impio é Babio:

Logo este impio merece louvor.

Refuta-se este sopbisma, \emph{distinguindo} a \emph{conclusdo:'o} impio merece louvor. em quanto é sabio, mas não \emph{em geral}.

e) Restrícção. --- O sophisma da \emph{resíriccio} dá-se, quando de um sentido \emph{a/firmativo} se passa para um sentido \emph{exclusivo}. Ex.

Tu deves ser sabio; \textbar{}

Logo não deves cuidar da sande.

Refuta-se, \emph{distinguindo} o antecedente: tu deves ser sabio: mas não deves ser \emph{sómente} sabio. \textbar{} .

%%>>> t115

\tese{115}{Especies do sophisma logico.}
As principaes especies do sopbisma \emph{logico} são: a \emph{ignorancia do elencho}, \emph{a petição de} \pg{97}  \emph{principio}, o \emph{circulo} vicioso, a \emph{não-causa} como causa, o \emph{consequente},

a \emph{inducção viciosa}, a enumeração \emph{imperfeita}, o \emph{accidente}.

a) Ignorancia do eleneho. --- A \emph{ignorancia do elencho} commette-se, quando por ignorancia ou malicia se desvia a questão do verdadeiro caminho. --- Neste caso avise-se o adversario.

b) Petição de principio. --- A \emph{pelição de principio} dá-se quando se suppõe como certo o que se ha de demonstrar. Ex.

O que não morre, é immortal;

Ora a alma humana não morre;

Logo a alma humana é immortal.

Este e os restantes sophismas refutam-se, denuociando-os.

€) Cirenlo vicioso. --- Commette-se o \emph{circulo} vicioso, quando duas coisas se demonstram uma pela outra. --- Assim é sophisma provar a \emph{immortalidade} da nossa alma pela sua \emph{espiritualidade}, e a \emph{espiritualidade} pela \emph{immortalidade}.

d) Não-cansa como causa. --- Commette-se este sophisma, quando se toma, como causa de um effeito, oqueéapenasuma --- condição, ou um antecedente ou consequente do mesmo effeito. Ex.

O imperio romano cahiu depois do Christianismo;

Logo o Christiaoismo foi a causa da queda do imperio romano.

e) Consequente. --- Este sophisma dá-se, quando, pelo facto de um consequente se deduzir de um antecedente, se julga que tambem este se possa deduzir d'aquelle. Assim posso dizer: « \emph{Pedro corre}; \emph{logo se} move» ; mas não: « \emph{Pedro} se \emph{move}; \emph{logo corre}».

f) Inducção vieiosa. --- A \emph{inducoüo viciosa} commette-se, quando se attribue a \emph{toda} a especie um predicado, que \emph{accidentalmente} se encontrou \emph{nalguns} individuos da mesma especie. Ex.

Este medico é ignorante;

"

Logo todos os medicos são ignorantes.

\&) Enumeração imperfeita. --- A enumeração \emph{imperfeita} dá-se, quando se attribue ao \emph{todo} um predicado, que é proprio sómente \emph{d'algumas partes} enumeradas. --- Ex.

Os europeus, os americanos e os asiatieos são brancos;

Logo todos os homens são brancos.

h) Accidente. --- Commette-se este sophisma, quando se attribue a um sujeito como \emph{essencial} um predicado, que é sómente \emph{accidental}. Ex.

Ha abusos no ensino da philosophia ;

Logo este ensino deve eliminar-se,

Vol, 1.

\textbf{LOGICA MATERIAL}

116. Divisio da Logica material --- Dividimos a \textbf{Logica} \emph{material} em \emph{tres} capitulos: no \emph{primeiro} iraiamos da \emph{verdade}; no \emph{segundo}, dos \emph{meios} pura a acquisicio da verdade: no \emph{terceiro}, do \emph{methodo}.

%%>>> cap-7

\chapter*{CAPITULO PRIMEIRO\\[0.3em]\large VERDADE}
\addcontentsline{toc}{chapter}{CAPITULO PRIMEIRO --- VERDADE}

\pg{98} Suwxaxio: --- Verdade, sua divisão e natureza. -- Certeza, sua \textbf{divisão},

existencia e propriedades, --- A existencia da certeza e o Sceplicismo. ---

Supremo e universal motivo da certeza, --- A certeza c á objectividade

das ideas. uüiversacs.

ArTIGO [\textbar{}

\emph{Verdade}, \emph{sus divisao e natureza}

%%>>> t117

\tese{117}{Verdade.}
\emph{Verdade} é a \emph{conformidade entre a intelligencia e o objecto}. --- O \emph{objecto} da verdade pode ser uma coisa \emph{real} ou \emph{ideal}, \emph{existente} ou \emph{possivel}, etc. \footnote{\emph{Veritas est adaequatio intellectus el vel}: Eis a definição \textbf{classica} da \emph{verrdade}, dada pelos eacholasticos, ou, melhor. pelo senso commum. Na opini\&o dos homens, uma pessoa possue a verdade, quando conhece um \textbf{objecto} como é em si mesmo, quando o seu pensamento corresponde á \textbf{realidade}, quando ha conformidade entre o que a intelligencia affirma do objecto-e o que o objecto é em si mesmo. Diz S. Thomaz: « Per \emph{conformitatem} \textbf{intellectus} et rei \emph{veritas} deiinitur ». (\emph{Swm}., \emph{Th}. p. 1,.q. 10, a. 2).\par D'esta definição resulta, que, para a constituição da verdade, \textbf{concorrem}, --- como \emph{materia remota}, dois termos, real ou logicamente distinctos, que são a \emph{intelligencia} e O \emph{objecto}, --- como \emph{materia proxima}, \emph{a especie} ou \emph{imagem intellectual}, \emph{representativa} da \emph{essencia} ou d' algum \emph{attriduto} do objecto, e a propria \emph{entidade} do objecto, percebida tambem pela intelligencia mas de um modo confuso e indeterminado, --- como \emph{forma}, a \emph{união intencional} d' aquella \emph{imagem} com a mesma \emph{entidade}. Esta união, para a inteligencia humana, consiste na affirmac\&o da conveniencia objectiva da \textbf{jmagem} ou forma com o objecto existente, e é significada pela proposição, em que \emph{unimos} o \emph{predicado}, que representa e exprime a forma ou natureza \textbf{percebida}, com o \emph{sujeito}, que fepresenla o objecto como é em si mesmo, mas\pgn{99}percebido de um modo confuso e indeterminado; como quando dizemos: Pedro é humem. Por isso o sujeito e o predicado, nas proposições, são conceitos \emph{objectivos}, emquanto denotam uma \emph{coisa}, e não meramente \emph{subjeclivos}, isto é, não são meras modificações da faculdade intellectual,\par \emph{A} adequação, ou conformidade, entre a intelligencia e o objecto --- não é \emph{physica,.-ou entilativa}, porque nem a intelligencia se torna o objecto, nem este se torna aquella, --- mas é \emph{representativa}, ou intencional, emquanto o pensamento corresponde á realidade, isto é, emquanto o ser \emph{ideal} do objecto está conforme com o ser \emph{real} d'elle. --- Para baver esta conformidade, não é necessario que a intelligeneia reproduza e represente \emph{todos} os elementos, que formam a \emph{comprehensão} do objecto conhecido ; busta que as \emph{propriedades} do objecto, que a intelligenciu representa, se encontrem realmente no proprio objecto. A conformidade não será perfeita, não será plena; mas será, sempre verdadeira. A \emph{conformidade completa}, \& \emph{adequação}, no sentido rigoroso da palavra, isto é, conbecer um objecto tanto, quanto é cognoscivel, é um dote só proprio da intelligencia infinita de Deus.\par Nem se opponha que, se a nossa intelligencia \emph{não} conhece o objecto tanto, quanto é cognoscivel, não podemos dizer que haja adequação entre a intelligencia e o objecto. Por quanto, aquella adequação, em que consiste a essencia da verdade, pode realizar-se, e realiza-se mesmo quando o pensamento representa \emph{ne só aliributo} do objecto, porque ha então uma \emph{equaldede} ou \emph{conformidade} entre a intelligencia e esse attributo conhecido. Repetimos : será uma \emph{conformidade}, ou egualdeade \emph{incompleta} com relação a \emph{todos} os attributos, a \emph{toda a comprehensdo} d' esse objecto, mas é uma \emph{conformidade}, ou egualdade completa com relação á \emph{existencia} d' esse attributo conhecido. Assim quando dizemos: o \emph{homem é racional}, ainda que ignoremos os outros attributos do homem, e a propria \emph{essencia} da racionalidade, comtudo ha uma \emph{egualdade} entre o nosso pensamento e o homem, porque alfirmamos que o homem é dotado de razão, e elfectivamente no homem \emph{existe} esta faculdade. --- Os escholasticos iuculeam esta doutrina, dizendo que o pensamento, para ser verdadeiro, deve conformar-se com o objecto \emph{formal}, e não com o \emph{material}, entendendo por objecto \emph{material} á coisa \emph{inteira}, com \emph{todos} os seus elementos, e por objecto \emph{format} a parte, o elemento, o aspecto, sob'o qual a coisa se apprehende.\par \emph{Kant} e, em geral, os \emph{swbjechivisías} não admittem que a \emph{verdade} consiste na \emph{conformidade} entre x intelligencia e o objecto, quando esse objecto é uma coisa real ou objectiva, porque, dizem, a realidade é \emph{émaccessivel} a nossa intelligencia. Por isso alguns escriptores rejeitam a classica definição da \emph{verdade}, proposta pelos escholasticos. \emph{Le Roy}. um dos chefes do \emph{modernismo}, din: « \emph{A} grande divergencia entre os escholasticos e nós funda-se uu propria noção de \emph{verdade} ». (\emph{Dorme et critique}, p. 355) Outros, entre os quaes \emph{Senirow}! e alguns \emph{meo-escholasticos} de \emph{Lowuvain}, embora conservem a formula escholistica da defivição da \emph{verdade}, alteram-Ihe comtudo o sentido,\pgn{100}dizendo que o \emph{objecto}, com que deve conformar-se o pensamento, nunca pode ser uma coisa real e \emph{objectiva} (porque não é possivel uma comparação entre o pensamento e a realidade), mas é apenas uma \emph{nota abstracta}, \emph{anteriormente conhecida}, um \emph{substituto mental}; dê modo que, nesta opinião, a \emph{verdade} se reduz a uma concordancia de ideas entre si.\par Critica. É certo que a definição de \emph{verdade}, dada pelos escholasticos, suppõe e exige o \emph{f«cfo} da conformidade do pensamento com a realidade, e por isso suppõe e exige que a nossa intelligencia \emph{posse} ter a intuição não só da sua imagem subjectiva, mas tambem do objecto rea! como é em si mesmo. Ora que a nossa intelligencia possa atlingir e attinja effectivamente o objecto real e externo, e de um modo \emph{immediato}, é coisa. certissima, Se não se admilüir esta intuição \emph{Zumnediata} da realidade, a. \emph{Metaphysica} reduzir-se-ba a uma sciencia puramente ideal, sem que haja a possibilidade de se encontrar um meio legitimo e racional, que nos melta em communicacio com o mundo exterior, Voltaremos ao assumpto.\par .A' \emph{verdade} oppóe-se o \emph{erro}, Diremos' da sua natureza, causas e \emph{remedios}.\par 'a) Natureza do erro. --- O \emph{erro} consiste na \emph{desconformidado} da intelligencia com o objecto. A \emph{desconformidade}, em que consiste o \emph{erro}, ou a \emph{falsidade}, por si importa uma contradicção, uma. \emph{falta de concordi} da intelligencia com o objecto, emquanto a intellizencia attribue a um objecto um predicado, que lhe não convem, ou nega-ihe um predicado, que Iherconvem. --- Aquella desconformidade, a qual consiste na falta de \emph{commprehensao} do objecto, emquanto a inlellizencia não conhece \emph{o} objecto tanto, quanto elle é cognoscivel, se é uma \emph{falta} de adequação \emph{absoliia}, não é uma falta de \emph{verdade}, não € um erro, ou uma \emph{falsidade}. --- O juizo, que traduz essa desconformidade ou discordia da intelligencia com o objecto. é e diz-se \emph{erroneo}, ou \emph{falso}. --- O \emph{erro}, pois. não é uma \emph{verdade incompleta}, como atfiemou \emph{Cousin}, mas é a privação da verdade. Porque o erro oppõe-se a verdade, como a privação da forma se ophõe a propria fortia; a;privacio não participa da forma, exclue-a. Assimr quem diz: a alma humana é \emph{mortal}, não diz uma verdade incompleta; diz ó contrario da verdade. Por isso Oo \emph{erro} é um \emph{juico}, que \emph{affirma} não ser o que é, \emph{ou ser o} que não é.\par b) Cansas do erro. --- Algumas são \emph{internas}, outras externas. --- As causas \emph{internas} são communs e \emph{particulares}: \emph{ÀS communs}. derivam da \emph{nelural propensão} e \emph{facilidade}, que os homens têm de julzar --- conforme o visivel, sem se- importar do invisivel, --- conforme a disposição da intelligencia, que nem sempre se regula com a devida prudencia, --- conforme os allectos do coração, que pode rejeitar o que lhe é contrario. As \emph{particulares} derivam da \emph{condição} e do \emph{fempercuinento} dos \emph{individuos}. --- As causas \emph{externas} são varias, como a ignorancia da Logica, a leitura de livros, onde se ensina o erro ou a verdade não é claramente exposta nem victoriosamente defendida.\par c) Remedios para o erro. --- São muitos. Os principaes são: --- o amor da verdade, --- \& pureza do coração, --- o estudo alurado da Pbilosophia escholastica, --- a duvida prudente, -- a oração assidua e fervorosiü.}.

%%>>> t118

\tese{118}{Especies de verdade.}
A \emph{verdade} \textbf{pode ser}: \emph{logica}, \emph{metaphysica} e moral. --- Verdade \emph{logica} é o \emph{conhecimento}, emquanto é conforme com o objecto conhecido. --- Verdade \emph{meta}- "\emph{physica} é o proprio \emph{objecto}, emquanto. corresponde ao conhecimento \pg{99} \pg{100} . --- Verdade \emph{moral} é -o0: \emph{discurso}, emquanto traduz a.

intenção de quem falla. --- Na verdade \emph{logica}, é a \emph{intelligencia} que deve conformar-se com o objecto; --- na \emph{metaphysica}, É o \pg{101}  \emph{objecto} que deve corresponder \& idea: --- na 2 moral, é o discurso

que deve exprimir o pensamento. s

Occupamo-nos agora da verdade \emph{logica}, reservando \textbf{para} mais tarde o estudo da verdade \emph{metaphysica} e da moral \footnote{A verdade \emph{logica} chama-se tambem \emph{accidental} ; a verdade \emph{ontologica} denomina-se tambem \emph{essencial}, \emph{metaphysica}, \emph{transcendental}. --- A verdade \emph{ontologica} é o proprio \emph{objecto}, emquanto corresponde ás ideas archetypas da intelligencia divina e é capaz de gerar um verdadeiro conhecimento de si mesmo na intelligencia creada. Neste sentido, todo o ente é verdadeiro, porque todo o ente é conforme com as ideas divinas e é por nós cognoscivel. --- A verdade \emph{logica} denota a. conformidade do nosso conhecimento com os objectos percebidos: e diz-se \emph{accidental}. porque tal conformidade não existe sempre no nosso pensamento. Neste sentido, o nosso conhecimento --- ---é \emph{verdadeiro}, quando percebemos existir nas coisas o que rentmente existe, ou não existir nellas o que realmente não existe, --- e é \emph{falso}, quando percebemos existir nas coisas o que realmente não existe, ou não existir nellas e que realmente existe. --- A verdade \emph{moral}, propriamente chamada \emph{veracidade}, consiste na conformidade do discurso com o pensamento, e existe quando se diz o que se pensa, quer este pensamento seja verdadeiro, quer seja falso. --- A verdade \emph{on£ologice} das coisas naturaes depende exclusivamente da intelligencia divina: porque as coisas dizem-se verdadeiras, emquanto correspondem ás ideas divinas. --- A verdade \emph{logica} da nossa intelligencia depende das coisas, e é causada por ellas: porque o nosso conhecimento é verdadeiro, quando e porque se conforma com ás coisas. --- A verdade moral depende dos nossos pensamentos; porque o nosso discurso é verdadeiro, quando se conforma com os nossos pensamentos. (Cf. S. Th., . lSent., Dist. 19, q. 5, a. 2 ad 2; \emph{De verit}, q. 1, a. 9). --- Portanto a verdade do \emph{disecwrso} é cmusada pelo pensamento; a verdade do pensamento, pelas coisas ;: a verdade das \emph{coisas}, pela intelligencia divina, D' este modo a intelligencia divina é a razão suprema, o fundamento ultimo de toda a intelligibilidade e de toda verdade.\par Á verdade, seja qual fôr \& sux especie. reside \emph{propria} e principalmente na \emph{Zn£elligencia}, e d'esta deriva para as \emph{coisas}, emquanto as coisas se referem isto é, pelo facto de se referirem a moesma intelligencia. Por quanto, a verdade consiste numa. conformidade e por isso numa \emph{relação de semelhança} da intelligzencia com a coisa. Oru a \emph{relação de semelhança} encontra-se exclusivamente nz intelligencia; porque nus coisas não pode encontrar-se senão O \emph{findamento} e i causa d' essa semelhavça ou adequação. Da \emph{intelligencia} \& verdade deriva para \emph{us coisas}; que, por jsso, se dizem \emph{verdadeiras}, emquanto se referem á mesma intellizencia, como os alimentos se chamam \emph{sãos}, emquanto produzem ou conservam a \emph{samidede} no \emph{animal}, em que a \emph{sanidade} propriamente se encontra (Cf. S. Th., \emph{doc}. \emph{cit}.).}.

%%>>> t119

\tese{119}{Divisão da verdade logica.}
A verdade \emph{logica} é --- \emph{a}) \emph{subjectiva} e \emph{objectiva}, --- \emph{b}) \emph{inmediata} e \emph{mediata}, ---- c) \emph{necessaria} e \emph{contingente}, --- \emph{d}) \emph{ideal} e \emph{real}, --- \emph{e}) \emph{metaphysica}, \emph{physica} e \emph{moral}, --- \emph{f}) \emph{directa} e \emph{reflexa}.

a) Verdade subjectiva e objectiva. --- É .\emph{subjectiva}, quando se considera como um \emph{producto} da \textbf{nossa} intelligencia, e por isso dotada de propriedades, que lhe convêm em virtude do principio, do qual deriva. --- É \emph{objectiva}, quando se considera em ordem ao \emph{objecto}, que representa, e por isso dotada de propriedades, que convêm ao proprio objecto.

b) Verdade immuediata e mediata. --- É \emph{immediata}, 1, quando

é percebida sem demonstração; ex. \emph{o todo excede} uma \emph{das} suas \emph{paries}. --- É \emph{mediata}, quando é deduzida de outras verdades pelo raciocinio; ex. \&' \emph{alma humana} é uma \emph{substancia espiritual}.

c) Verdade necessaria e, contingente. --- É necessaria, quando exprime uma coisa, que não póde deixar de ser o que é; ex \emph{o homem é animal} racional, porque o homem não pode conceber-se sem as notas de \emph{animalidade} e de \emph{racionalidade}. --- É \emph{contingente}, quando exprime uma coisa, que pode deixar de ser o que é; ex. \emph{o mundo existe}. \textbar{} U

.d) Verdade ideal e real. --- É \emph{ideal}, quando exprime um juizo \emph{abstracto}; ex. \emph{se} a \emph{pedra sentisse}, \emph{teria alma}. --- E \emph{real}, quando exprime uma coisa existente na realidade, independentemente do nosso pensamento; ex. o \emph{homem é livre}.

'e) Verdade metaphysica, physica e \textbf{moral}. --- É \emph{meta pharsica}. quando exprime um \emph{juizo} de ordem \emph{metaphysica}; ex. \emph{todo o effeito exige} uma causa. --- É \emph{physica}, quando exprime um juizo de ordem \emph{physica}; ex. \emph{o fogo queima}. --- É \emph{moral}, quando exprime um juizo de ordem «\emph{oral}; ex. \emph{o filho respeita seus paes}.

f) Verdade directa e reflexa. --- A \emph{directa} é a que contem um juizo \emph{directo}, e por isso exprime um \emph{objecto}, como é em si mesmo, isto é, affirma que um objecto é tal em si mesrao, qual é percebido pela intelligencia; ex. \emph{o home é capaz de sabedoria}. --- A \emph{reflexa} é a que contem um juizo \emph{reflexo}, e por isso exprime a \emph{conformidade} do conhecimento com o objecto, isto é, exprime a convicção de que o conbecimento é conforme com a realidade; ex. \emph{estou certo de que o homem é capaz} de \emph{sabedoria} \footnote{É evidente á \emph{differenca} entre \emph{a} verdade \emph{directa} e a \emph{reflexa}. \emph{A directa} tem por termo um \emph{objecto}, \emph{emquanto} este é conforme com a intelligencia ; ao passo que a \emph{reflexa} tem por termo a propria \emph{conformidade} da intelligencia. com o objecto conhecido. --- \emph{Que} a nossa intelligencia possa reflectir e reflicta efeclivamente \emph{sobre} os seus actos e conheça a propria \emph{conformidade} com os objectos, e \emph{assim possua} a verdade \emph{reflexa}, é coisa que lodos percebem\pgn{103} e admittem, --- A dificuldade consiste em determinar o \emph{modo} por que a intelligencia conhece tal conformidade. Entendem alguns que, para isso, seja preciso estabelecer um confronto entre o objecto conhecido e a imagem, em que e por que o objecto é conhecido; mas não advertem que, para estabelecer tal confronto, seria necessario conhecer o objecto por um outro modo, ou por um outro meio, distineto e diverso d'essa imagem; a não ser que se queira confrontar a imagem com um termo desconhecido.\par Não é possivel, nem é precisas tal solução. A coisa é muito simples. Para que a intelligencia possa conhecer a propria \emph{conformidade} com o objecto, deve apenas rellectir sobre \emph{o} seu acto cognoscitivo. Por quanto, a intelligencia, quando conhece um objecto, não é meramente \emph{passiva}, isto é, não se imita a receber em si a impressão \emph{e} a imagem do objecto, mas é eminentemente «\emph{activa}, porque, recebida a impressão e a imagem do objecto, veage no mesmo objecto, representado por aquella imagem, e, por este acto de reacção, recebe-o em si, apprehende-o, exprime-o intencionalmente, e, exprimindo- o. \emph{conforma-se} com elite. Esta \emph{conformidade} da intelligencia com o objecto é o termo \emph{natural} do acto cognoscilivo; porque a nutureza da intelligencia consiste propriamente \emph{em} se tornar \emph{conforme com o objecto}. E, se tal é a natureza do acto intellectual, isto é, se a verdade de um juizo resplandece na propria natureza do conhecimento \emph{directo}, é evidente que 2 intelligencia, \emph{reflectindo} sobre a. natureza do seu acto, percebe, de um modo immediato, a propria conformidade com o objecto, \emph{c} assim possue a verdade \emph{reflexa}, Tal é o ensino de S. Thomaz. (\emph{De} Ver., q. 1, a. 9).}.

%%>>> t120

\pg{102} \pg{103} \tese{120}{A verdade logica encontra-se perfeitamente no juizo.}

a) A verdade \emph{logica} encontra-se \emph{perfeitamente} \textbf{naquelle acto} da intelligencia, que encerra e exprime, \emph{de} um \emph{modo explicito}, a conformidade da mesma intelligencia com o objecto. Ora o \textbf{acto} da intelligencia, que encerra e exprime, \emph{de} um \emph{modo explicito}, a conformidade da propria intelligencia com o objecto, é o \emph{juizo}; porque pelo \emph{juizo} a intelligencia declara que um objecto é em si, como ella o percebe, isto é, que uma natureza, ou \textbf{essencia}, representada pela forma intencional, é inherente ao objecto,

£ existente na realidade, e por isso une os dois termos distinctos, que são o objecto e a sobredita forma intencional. Logo a verdade \emph{logica} encontra-se \emph{perfeitamente} no \emph{juizo} \footnote{\emph{A} nossa tellisencia não se conforma perfeitamente com o objecto senão no \emph{fuize}, e por isso a verdade não se encontra perfeitamente senão neste acto intellectual. Quando dizemos \emph{simplesmente}: \emph{Rouwm}. 2 NOSSA intelligencia exprime \emph{um} conceito. uma natureza, uma imagem: mas essa imagem não poda dizer-se conforme, \emph{nem} desconforme com o seu prototypo. Se a este realmienie não ae refira: assim como uma imagem, pinlada na tela, não pode dizer-se semelhantes \emph{ou} dessemelhante, ao seu exemplar, quando se considera ex \emph{ivaumente} em si mesma. sem relação ao dito exemplar, Ora a imagem \emph{de um objecto}, \emph{obtida pela} simples «oprebensão, não se refere \emph{explicitamente ao proprio objecto}, senão no \emph{fuizo}; assim a\pgn{104}\emph{humanidade} não se percebe claramente como natureza de um individuo, se não quando a refiro, de um modo \emph{explicito}, a. esse individuo, e uno essa forma intencional com o seu sujeito, dizendo: \emph{Pedro} é homem.\par Affirmando que a. verdade se encontra \emph{perfeitamente} no juizo, não queremos significar que com \emph{um só juizo} se possa conhecer tudo o que se refere a um objecto: mas êntendemos qué só \emph{por} meto \emph{do juizo} podemos conhecer \emph{completa} e \emph{perfeitamente} um objecto, porque só neste acto e por este acto á intelligencia se conforma \emph{plena} e \emph{perfeitamente} com o objecto real.\par Como tambem, quando affirmamos que a verdade se encontra perfeitamente no \emph{juizo}, referimo-nos sempre a intelligencia humana, e não a angelica, muito menos a divina. --- A idea da intelligencia humana é imperfeita, limitada, abstracta; e, embora represente um elemento, ou uma propriedade, ou a propria essencia, não representa, comtudo, o objecto \emph{infeiro}, como é em si mesmo; e, por isso, depois de termos considerado as partes, os elementos do objecto, representados pela idea. devemos recolher esses elementos, essas partes, e constituir o todo, como se encontra na realidade: e esta obra é do \emph{juizo}. --- A idea, orem, da inteligencia angelica, e, \emph{a-for}- \emph{Hori}, a idea da divina, exprime o L \textsc{OMM} como é em si mesmo, com a sua\par natureza, com o seu ser, e com todas as propriedades e notas individuantes; e por isso o Anjo e, \emph{a-fortiori}, Deus não precisam de \emph{compór} ou de \emph{dividir}, isto é, de formar juizos \emph{affirmativos} (em que o predicado \$e \emph{conwpüe} com o sujeito), nem juizos \emph{negativos} (em que \emph{o} predicado se \emph{divide} do sujeito). para alcançar a verdade. --- D' onde se vê que o \emph{juizo}, OU a composição e \emph{divisão}, É uma coisa \emph{accidental é verdade}, porque não constitue a sua \emph{essencia}, mas é apenas um \emph{meio} natural para a intelligencia \emph{humana} attingir a verdade. (\emph{Stan}, \emph{th}., p. 1, q. 11, a. 2; \emph{De} Ver, q. 1, a. 3).}.

b) A verdade \emph{logica}, sendo uma adequação da intelligencia com o objecto, deve encontrar-se \emph{perfeitamente} no acto, em que se encontra a perfeita adequação da intelligencia com o objecto. Ora essa perfeita adequação encontra-se no \emph{juizo}. Porquanto, o \emph{juizo}, exprimindo, de um modo explicito, a \emph{na£ureza} ' e O ser, ou \& \emph{existencia}, de um objecto, exprime o objecto como é em si mesmo, e por isso affirma a perfeita adequação da inteligencia com o mesmo objecto. Ássim quando digo: \emph{Pedro é homem}, não só denoto uma determinada natureza, que é a \emph{humanidade}, mas tambem affirmo que esta natureza é, ou \emph{existe}, num determinado individuo, que é \emph{Pedro}. Logo a verdade \emph{logica} encontra-se \emph{perfeitamente} no \emph{juizo} (t).\footnote{A verdade deve encontrar-se \emph{perfeitamente} naquelle acto intellectual, que adequa \emph{perfeitamente} o objecto. Ora todo o objecto (fallamos do objecto creado e \emph{finito}) consta, como veremos na \emph{Ontologia}, de dois elementos, reaes e distinctos, que são a \emph{essencia}, ou a \emph{natureza}, e OQ ser, ou a \emph{existencia}. \emph{A essencia}, para se tornar uma coisa real, deve receber o ser, ou \emph{a existencia}; e, só depois de reCeber o ser, ella é. ou \emph{existe}. \emph{A exsencia}, que recebe o \emph{ser}, \emph{É potencia}: O \emph{Bare} é recebido na \emph{essencia} e que a faz existir \emph{actnabnerte},\pgn{105}\emph{é acto}. Privada do ser, \& essencia é uma mera \emph{abstracção}, é um \emph{nada}. Por isso uma.coisa diz-se ente, não pela essencia, mas pelo \emph{ser}, que possue. D' onde se segue que o acto intellectual é perfeitamente adequado ao objecto, quando exprime, e de um modo \emph{explicito}, ambos esses elementos do objecto, isto é, a \emph{essencia} e o \emph{ser}. Ora tal acto é o \emph{juizo}; pois este, como sc vé nas proposições, contem sempre a essencia e O ser. Assim quando digo: Pedro é \emph{homem}, não só denolo uma \emph{essencia}, que é a hwmamnidace, mas affirmo que essa essencia é, ou \emph{existe}, e se encontra nesse individuo, que por Isso é um ente \emph{real}. --- Nem podia ser d'outro modo. A verdade \emph{ontologica} das coisas funda-se mais no \emph{ser} d'ellas, do que na sua essencia. Com effeito, uma coisa tem Lanto de \emph{verdade}, quanto tem de entidade; porque \emph{ente} e \emph{verdadeiro} são, nu. realidade, uma e a mesma coisa (ens \emph{et} verum converfwniwr) Sendo assim, é claro que a verdade \emph{Ingica} deve encontrar-se naquelle acto, que, de um modo \emph{explicito}, exprime não só a essencia das coisas, mas tambem o seu ser. Tal acto é o \emph{juizo}. (CC T. \emph{Sent}., Dist. 19,\par q a.d.\par Note-se que a intelligencia é determinada ao \emph{juizo}, logo que lhe se manifeste o \emph{ser} das coisas, Para isso, não é preciso que ella reflicta em si, --- nem que tenha consciencia da verdade do proprio juizo, ou da sua aptidão para conhecer, --- nem que esteja certa da \emph{suficiente} manifestação do mesmo ser. Se a manifestição é tal, qual a exige a faculdade perceptiva, á reacção d'esta, por meio de um juizo correspondente, -é um acto muito natural e necessario, por isso mesmo que é uma faculdade percepliva.\par Se a \emph{verdade} sc encontra perfeitamente no \emph{j«czo}, tambem no \emph{juizo} deve encontrar-se \emph{propriamente á falsidade}, 0u \emph{O erro}. Por quanto, a \emph{falsidade} só pode encontrar-se \emph{propriamente} naquelle acblo, em que a intelligencia affirma que uma coisa é o que não. é, ou nega. que é o que é. Ora este acto é 0 \emph{fuizo}. --- Alem disso, a \emph{falsidade} é contraria á \emph{verdade}. Ora, como os contrarios se referem ao mesmo sujeito (\emph{contraria sunt circa idem}), \emph{a}. \emph{falsidade} deve encontrar-se naquelle sujeito, em que se encontra a \emph{verdade}. Ora. a \emph{verdade} encontra-se na intelligencia emquanto \emph{compõe} ou \emph{divida}, isto é, emquanto \emph{julga}. Logo tambem a \emph{falsidade} deve encontrar-se na intelligencia, emquanto \emph{julga}. (\emph{Sum}. \emph{Th}., p. 1, q. 17, a. 4).\par Mas- o \emph{erro}, Mesmo BO \emph{juico}, € sempre. umi. coisa. \emph{accidental}. para. a. nossa intelligencia, Esta faculdade tende por si para a verdade; e, se erra, 2 culpa não se deve procurar nella, mas em causas extrànhas. --- Tacs causas dos erros reduzem-se a tres classes: ---- \emph{phujsiologicas}, \emph{psychologicas} e moraes. Reservando para mais tarde o estudo das causas \emph{phgsiologicas}, diremos alguma coisa acerca das outras duas especies. --- Às causas \emph{psy}- €\emph{hologicas} podem reduzir-se ao \emph{influxo da imaginação}, \emph{a falta de reflexão}, \emph{a incerteza} ou \emph{ambiguidade} dos \emph{vocabulos}, \emph{a ignorancia da arte do raciocinio}. --- Ás causa moraes dos erros, sobretudo em materia de Beligiio e de Moral, são as \emph{paixões desordenadas}. É um facto que o \emph{homem tem a} tendencia para julgar \emph{como} verdadeiro o que lhe agrado, \emph{e} que o amor proprio, mais\pgn{106}"Bo \emph{do} T \emph{j}- .\pgn{107}é redondo) o predicado se attribue ao sujeito, não se funda na experiencia, mas exclusivamente na propria natureza do sujeito, e a \textbf{verdade} de tal juizo é independente de toua a experiencia e subsistiriza. \textbf{sempre}, embora não existisse a experiencia. --- O escriptor allemão encontrou juizos de uma terceira especie, que não são nem completamente \emph{analyticos}, ou \emph{a-priori}, nem completamente \emph{syntheticos}, ou \emph{a-posteriori}, mas participam de uns e de outros, e por isso são \emph{symiheticos a-priori}. Exponhamos a sua theoria.\par O juizo, diz Kant, --- é \emph{analytico}, quando o predicado \textbf{está} contido, ainda que implicitamente, no conceito do sujeito, --- é \emph{synthetico}, quando o predicado está completamente fóra do conceito do sujeito, e por \textbf{isso} a união de um com o outro importa. uma \emph{synthese}. B conclue: «Por \textbf{isso}, são \emph{analyticos} os. juizos, em que não só ha união, mas tambem identidade entre o predicado e o sujeito, e são \emph{syntheticos} aquelles, em que ba união sem identidade.. Os \emph{aualyiicos} podem chamar-se \emph{explicaticos}; 0s \emph{synthketicos}. \emph{extensivos}: porque os primeiros nada «ecrescentam pelo predicado ao conceito do sujeito,. mas apenas dividem pela analyse o mesmo conceito nos seus conceitos parciaes, que se percebiam nelle (embora de um mudo confuso), ao passo que os segundos accrescentam ao conceito do sujeito um predicado, que de modo nenhum podia conceber-se contido no mesmo sujeito, nem podia , deduzirse pela analyse» (\emph{Critica da razão pura}, latrod., TV). --- Os juizos \emph{synilelicos}, segundo Rast, subdividem-se em syntheticos \emph{a-priori}, e em syn- Lheticos \emph{a-posteriori}, São synthelicos «\emph{priori}, quando têm algum elemento, que não é dado pot experiencia, mas deriva espontaneamente da intelligencia, e por isso exprime uma coisa \emph{wnioersal} € necessaria. São syntheticos \emph{a-posteriori}, qu \emph{empiricos}, quando se fundam exclusivamente na experiencia, e por isso exprimem uma coisa \emph{singular} e \emph{contingente}. --- Os \emph{syningticos a-priori} são verdadeiramente \emph{extensions}. € constituem propriamente a Sciencia e são objecto da Metaphysica; mas os \emph{syn£heticos a-posteriori}, por serem particulares e \emph{contingentes}, não são \emph{srientificos}, e, como os \emph{anal ticos}, não levam a novos conhecimentos.\par Mas como descobriu Manuel Kant esta terceira especie de juizos? A coisa foi simples. Disse elle que um juizo \emph{analytico}, \textbf{para} ser \emph{tal}, deve ter tres condições : 1.5) o seu predicado deve estar contido no \emph{conceito} do sujeito, --- 9,23) deve ser \emph{necessario}, --- 3."; deve ser \emph{universal}, Se falta uma d'estas ! condições, o juizo deixa de ser «\emph{nalyfico}. A's vezes pode "um juizo deixar de ter a primeira condição e possuir á segunda e a terceira, isto é, pode um juizo não ter o predicado contido no conceito do sujeito e comtudo ser necessario e universal. Neste caso o Juizo --- é \emph{synlhketico}, porque o predicado está fora do conceito do, sujeito e só pela experiencia se conhece a sua relação com o mesmo sujeito, --- mas é synthetico \emph{a-priori}, porque é \emph{universal} e \emph{necessario}, e por isso não se funda num dado da experiencia, que é por si \emph{singular} e \emph{contingente}, mas deve derivar exclusivamente da intelligencia, e deriva d' ella: espontaneamente, \emph{a-priori}, isto é, por uma especie\pgn{108}de instincto ou de força mechanica. Todas as intelligencias possuem naturalmente umas formas; dotadas de \emph{universalidade} e de \emph{necessidade}. que, applicadas successivamente aos objectos da experiencia, os toroam determinados e definidos. --- Kant apresenta varios exemplos d'estes juizos \emph{synthecos a-priori}. Citamos tres: --- \emph{todo} o \emph{effeito exige} uma causa: --- \emph{cinco mais sete são} eguaes a \emph{doce}; --- a \emph{linha recta é a mais} breve \emph{entre dois pontos}. Estes juizos --- são \emph{syntheticos}, porque o predicado não está contido no conceito do sujeito, --- mas são \emph{a-priori}, porque são dotados de \emph{universalidade} e de \emph{necessidade}, isto é, porque o predicado convem \emph{sempre} e \emph{necessariamente} ao sujeito. --- D' estes principios Kant tira duas consequencias: --- 1.º) os juizos \emph{analtyticos}. ou \emph{a-priori}, em que o sujeito e o predicado são termos universaes, por Serem \emph{necessarios} e \emph{universaes} e consequentemente por derivarem espontanea e exclusivamente da intelligencia, não têm valor algum \emph{objechivo} ; --- 2“) os juizos \emph{syntheticos} «\emph{priori}. que têm um sujeito de \emph{experiencia}, mas um predicado \emph{universal} e \emph{necessario}, por causa da completa \emph{subjectividade} d' esse predicado, são tambem destituidos de valor \emph{objectivo}.\par Agora. uma. hreve critica a estas affirmações do'eseripLor àallemão.\par ; Antes de tudo, as definições do juizo \emph{analytico} e \emph{synthelico}, embora tenham alguma coisa de verdade, podem todavia prestar-se a equivocos e abrir adito a erros, sobretudo se se tenha em vista. o systema de Kant. Este confunde sempre o \emph{sijeito} com o \emph{conceito}, passando do \emph{sujeito}, isto é, da \emph{coisa}, á qual se attribue um predicado, para o \emph{conceito}; come se podesse, indifferentemente e sem distineção alguma, dizer-se do \emph{conceito} o que se affirme. do \emph{sujeito}, quando o \emph{conceito} e o \emph{sujeito} são duas coisas distinctas, existindo o \emph{sujeito} na \emph{realidade} e \emph{o-conceito} na \emph{intelligencia}. 1, ainda que o conceito, tomado \emph{objectivamente}, se identifique com o sujeito, comtudo esta identidade refere-se é \emph{coisa representada}, e não \emph{ao} modo \emph{de representar}. --- Por isso, a affirmação que o juizo é \emph{analytico}, quando o predicado está contido no \emph{conceito} do sujeito, \textbf{pode} prestar-se a equivocos ; visto que. o \emph{conceito} é uma coisa \emph{subjectiva} e: por isso \emph{relativa} e \emph{variavel}. Se Oo \emph{conceito} exprimisse a natureza OU. \emph{essencia} do sujeito, vão haveria que dizer; mas, na theoria d' esse escriptor, não a exprime. --- Alem d'isso, é falso que o juizo \emph{analytico} não seja principio de novos. conhecimentos. Os predicados de um sujeito, no juizo. \emph{analytico}, estão todos contidos na \emph{essencia} e por isso no \emph{conceito} do mesmo sujeito; mas não estão todos contidos \emph{explicitamente} quanto a nós, e por isso esse juizo leva-nos ao conhecimento de coisas, que antes não se conheciam. --- Finalmente, ha juizos \emph{analyticos}, em que a'conveniencia do predicado com o sujeito não se conhece pela immediata intuição do sujeito. mas pelo \emph{raciocinio}. --- Não menos equivoca é a definição do juizo \emph{synthetico}, dada pelo eseriptor allemão. Diz elle que, neste juizo. o predivado está inteirnmente fóra. do conceito do sujeito. Ora, ainda que, antes do juizo, podesse dizer-se que o predicado não está contido no conceito, do sujeito, não se pode dizer o mesmo no acto do juizo, porque a intelligencia não poderia a. que um predicado convem a um\par sujeito, se não visse essa. conveniencia. Por isso é falso que, no juizo \emph{synthetico}, haja união sem \emph{identidade}. Quando digo: \emph{Pedro é bom}, não aflrmo\pgn{109} que \emph{Pedro é a bondade}, mas que é \emph{dotado} de bondade, e por isso affirmo que tudo o que o predicado exprime se encontra ou está contido no sujeito, isto €, affirmo a \emph{identidade} do predicado com o sujeito. Esta \emph{identidade} alfirma-se em todo e qualquer juizo aíffirmativo, quer \emph{analytico} quer \emph{synihetico}. (\emph{Sum}. \emph{Th}., p. I, q. 89, a. 5 ad 3).\par Ha mais. Os juízos \emph{ealyficos}, como os entendeu Kant, em que todos os elementos são \emph{subjectivos}, ou \emph{a-priori}, sem nenhuma dependencia da experiencia, sem nenhuma relação aos objectos, não existem e são verdadeiras chimeras. Por quanto, \emph{o} sujeito e o predicado, que constituem um juizo \emph{analytico}. embora sejam termos \emph{universaes}, comtudo estão \emph{intrinsecamente} contidos nos objectos da experiencia. Ássim quando digo: o \emph{homem é racional}, o homem não é uma coisa imaginaria, mas é um ente real, independente do nosso pensamento, e quando affirmo que o homem é racional, affirmo que esse atributo é uma coisa \emph{intrinseca} e essencial a. esse sujeito. Logo todos os juizos \emph{analyticos} têm o seu fundamento nas coisas ou nos objectos da experiencia, e por isso não são inteiramente \emph{subjectivos}, não dependem exclusivamente da inteligencia: --- O' sujeito \emph{e} o predicado do juizo podem ser termos. \emph{abstractos} \& universaes: mas, como. o \emph{abstracto} deriva do \emph{concreto} e O universal deriva do \emph{singular}, segue-se que o sujeito e o predicado são termos \emph{objectivos}.\par Não menos chimericos são os juizos \emph{syntheticos a-priori}, forjados por Kant. Primeiramente, esses juizos repugnam a indole da nossa. intelligencia. Por quanto, não é possivel 2 enunciação de um juizo, se não se conhece a conveniencia entre o predicado e o Sujeito, porque é essa mesma conveniencia que se deve affirmar ou negar ; aliás o juizo seria um acto irracional. A conveniencia entre o predicado e o sujeito só pode ser conhecida --- ou pela \emph{intuição}. --- ou pelo raciocinio, --- ou pela experiencia. Ora, nesses juizos \emph{syntheticos a-priori}, a conveniencia entre o predicado e o sujeito não é conhecida --- nem pela intuição, porque são \emph{syntheticos}, isto é, porque o predicado não está contido uo \emph{conceito} do sujeito, --- nem pelo \emph{raciocinio}, porque são \emph{espontaneos}, --- nem pela, experiencia, porque são \emph{a-priori}, e derivam do fundo da inteligencia sem correspondencia com o que está de fóra, --- Além d'isso, affirma Kant que os juizos \emph{syntheticos a-priori} São universaes e --- \emph{necessarios}, embora o predicado não esteja contido no \emph{conceito} do sujeito. Ora esta affirmação é absurda. Um juizo só pode ser universal e necessario, isto é. um predicado só pode convir \emph{sempre} e \emph{necessariamente} a um sujeito, emquanto ou porque o predicado està contido na propria natureza e no \emph{conceito} do sujeito, visto que só o que constitue ou acompanha a natureza, convem \emph{sempre} e \emph{necessariamente} ao sujeito, que nella subsiste.\par Com relação aos exemplos, que o escriptor aliem\&o adduz como prova da sua theoria, e que seriam juizos \emph{syn£heticos a-priori,no} sentido de Kant, não são elles senão juizos \emph{analyticos} ou \emph{synihelicos}, no sentido da \emph{Escholastica}. --- Assim o juizo: \emph{todo o effeito exige} uma causa, --- não é \emph{synihetico a-priori}, mas exclusivamente \emph{analgytico}. Por quanto, um juizo é \emph{amalytico}, Segundo o proprio Kant, quando o predicado esta contido no conceito do sujeito, \emph{ou} se encontra na resolução do sujeito. Ora a \emph{exigencia de} uma\pgn{110}causa está contida no conceito de \emph{effeito}. Na verdade, \emph{effeito} é uma coisa, que foi feita, que começou, não deriva de si mesma, mas de outro ente. Ora derivar de outro ente é depender. d' elle, como da causa do proprio ser. Logo a \emph{exigencia de} uma causa está contida no conceito de \emph{efeito}: --- Diga-se o mesmo acerca do juizo: \emph{cinco mais sete são eguaes} a \emph{doze}. Elle é exclusivamente \emph{analytico}: porque o predicudo se encontra na resolução do SA De facto, o conceito do predicado é o mesmo que o do sujeito, pois é unidade repetida \emph{doze} vezes, embora por differentes signaes. Nem vale a razão de Kant, quando diz que o predicado «\emph{dose}» não está contido no conceito do sujeito «\emph{cinco mais} sete», pelo facto de nos servirmos, ás vezes, de um signal externo, \emph{p}. ex., dos dedos, para juntar o numero \emph{cinco} ao numero “\emph{sete}. Com effeito, \emph{o} juizo consiste DA \emph{conveniencia} entre \emph{o} predicado e o sujeito ; e esta conveniencia é percebida exclusivamente pela inlelligencia, ' que analysa os éonceilos, ou examina as: partes. A \emph{experiencia} pode ser uma \emph{condição}, mesmo necessuria, para se averiguar se os conceitos convém,\par ão ou a causa da união. --- Como tambem " OU não, entre si, mas não é a ra \emph{onalytico} é o juizo: a \emph{linha recta} é a \emph{mais breve entre dois pontos}. O conceito da \emph{maxima brevidade} está contido” no conceito da \emph{Tinha} recta; porque, fixados dois pontos, percebe-se logo que a linha \emph{recía}, traçada. entre esses pontos, é mais breve do que qualquer linha curva. A figura sensivel, que exprime este \emph{juizo}, não influe na enunciação nem na verdade d'elle, mas Serve apenas para o fixarmos melhor na jnteltigencia. --- Todos os outros juizos, apresentados por Kant, (nudam-se --- ou na immediata relação dos conceitos, --- ou na dedueção do raciocinio, --- ou nos dados da experiencia; e por isso, ou são \emph{analyticos}, ou Sio \emph{syntheticos}, no sentido escholastico.\par D' onde se vê que, se não têm valor \emph{objectiro} os juízos \emph{analyticos} e os \emph{syntheticos a-priori}, como os entendeu Kant, porque insubsistentes e chimericos, têm, pelo contrario, valor \emph{ob/ectivo} os juizos \emph{analyticos} e os \emph{syntheticus}, como os entende a Escholastica. Nem se insisla com o escriptor allemão que um conceito «\emph{universal gpnecessario} não pode provir da experiencia, cujo objecto é \emph{singular} e \emph{conting} . Todas ás coisas \emph{singulares} e \emph{contingentes} têm em si um elemento \emph{nniversaN} e necessario, que é a essencia d' ellas, abslrabida de todas gescondigoes individuaes. (\emph{Sum}. \emph{Th}., p. 1, q. S6, a. 3).\par Vê-se tambem que, se os juizos \emph{analyticos} e \emph{synthelicos a-priori}, inventados por Kant, não concorrem para o progresso da sciencia (a qual tem por objecto as coisas, e não os sonhos), os juizos \emph{analyticos} e \emph{sypmihelicos}, tomados no seu legitimo Sentido, servem admiravelmente para o desenvolvimento de todas as sciencias, racionaes e experimentaes.\par D'aqui a pouco, fallaremos nas taes \emph{formas subjectivas}, que aos juizos \emph{syntheticos a-priori} communicariam a \emph{universalidade} e a \emph{necessidade}.}

%%>>> t121

\pg{104} \pg{105} \tese{121}{A verdade logica encontra-se imperfeitamente na idea.}

a) \emph{Encontra-se na idea}. --- A verdade consiste na \emph{conformidade} da intelligencia com o objecto. Ora a nossa intelligencia, quando apprebende simplesmente um objecto, forma uma \emph{idea}, \pg{106}  mM \textasciicircum{} \emph{e} H . E : H \emph{o} que é uma imagem ou semelhança do proprio objecto, e por i meio d' essa idea se torna \emph{conforme} com o objecto; porque \emph{T} uma faculdade perceptiva, informada pela especie ou imagem o \emph{v} que o amor da verdade, regula os seus juizos e não o deixa reconhecer os ES proprios erros. --- Não admira : que o homem \emph{erre}. O erro é compativel com V uma intelligencia finita e dependente, como é a nossa. O que, porem. repuo gna a razão é a pertinacia no erro, a qual querem dar o nome de \emph{caracter} = e que não é senão o privilegio dos \emph{insipientes}. Diz Cicero: « \emph{Cujusvis homim más est} errare : \emph{nullius nist insipientis in errore perseverare} », (Fili pp. XI, 2). ---

\emph{L}4 --- Que a verdade se encontre perfeitamente no \emph{juizo}, isto é, que a nossa ns intelligeocia, no juizo e pelo juizo, se ache em pleno accordo com o objecto, \emph{o}” é dictado do senso commum, mas nezado por todos os \emph{subjectivistas}. Dizem

\textbar{} estes que a intelligencia, pelo juizo, une duas coisas separadas, como são --- o predicado e o sujeito, e que por isso a união e a identificação do predi- . cado com o sujeito depende exclusivamente da intelligencia, e o juizo é ind inteiramente \emph{subjectivo} e não corresponde á realidade, isto é, não enuncia O o que o objecto-é em si- mesmo, mas o que a intelligencia pensou. --- Ora T isto é falsissimo. A inteligencia, quando julga, não é a causa por que um - predicado compete a um sujeito; s não faz Senão estabelecer e ratificar uM nm facto, isto é, que um predicado real está realmente contido num sujeito \emph{V} real; e por isso a união real do predicado com o sujeito não é effeito do : juizo, mas é anterior, é presupposto ao juízo e é a verdadeira causa d' elle. --- Assim, quando dizemos: \emph{Pedro} é bom, emtanto Se forma este juizo e se

affirma que a \emph{bondade} convem a. \emph{Pedro}. emquanto se averiguou que na rea-

lidade esse predicado está contido nesse sujeito, independentemente do : nosso pensamento. Logo o predicado, no juizo, não convem ao sujeito, " porque assim \emph{o} julga a inteligencia, mas a intelligencia assim o julga, \emph{o} porque o predicado convem ao sujeito. --- Quando, pois, se diz que a intel-

ligencia, no juizo, une o predicado com o sujeito, não se entende que a - intelligencia une duas coisas realmente separadas, uma da outra, mas deve

entender-se que ella julga estarem unidas na realidade duas coisas, que \emph{o} tinham sido consideradas separadamente. " O principal adversario da \emph{objectividade} dos juizos é \emph{Manuel Kant}. Como. = os erros d' este escriptor são seguidos por muitos racionalistas, livres peno7 sadores, protestantes, pelos sequazes do \emph{modernismo} (que o S. P. Pio X, ind na neci Pond S sei. 1907, chamou a \emph{s}, \emph{ynthese de tulas as heresias}), i A devem aqui ser expostos e e refutados. \emph{e}

O ponto de partida de todo o systema \emph{subjectivista ou idealista} de Kant : \emph{No} é constituido pelos celebres \emph{juízos syntheticos a-priori}. --- A Escholastica. " seguindo o bom senso, tinha ensinado que ha só duas especies de juizos : \emph{P} --- 0 \emph{analytico}, ou «-\emph{priori}, e o \emph{syntretico}, ou \emph{a-posteriori}, conforme a conve- i doc \emph{nieucia} do predicado com o sujeito é conhecida pela razão, ou pela expe- - \textbar{} \emph{riencia}: \emph{ou}, por outras palavras, conforme o principio da synthese ou da \emph{lo união é interior}, ou \emph{exterior}, a0 sujeito, Não ha meio termo. Embora os \textbar{}\textasciicircum{} \emph{elementos} do juizo \emph{enalytico} sejum dados pela experiencia, comtudo elle É \emph{e} sempre \emph{analytico}; visto que o \emph{motivo}, \emph{por} (que nesse juizo (por ex., \emph{o circulo} \pg{107}  de \textbf{uma coisa}, se torna \emph{intencionalmente} a propria \textbf{coisa}. Logo a \textbf{verdade} \emph{logica} encontra-se na \emph{idea}.

b) \emph{Mas encontra-se imperfeitamente}. --- O \textbf{conhecimento} de um objecto é perfeito, quando se conhecem, de um \textbf{modo} \emph{explicito}, \pg{108}  todos os elementos, de que consta e que são a \emph{natureza}, ou a essencia, e O ser, ou a \emph{existencia}. Ora a \emph{idea}, se representa, de \pg{109}  um modo \emph{explicito}, a \emph{natureza}, não exprime, do mesmo modo, a \emph{existencia}. Não representando a idea, de um modo \emph{explicito}, \pg{110}  \emph{existencia}, não pode realizar-se \emph{perfeitamente} a \textbf{união} dos dois termos distinctos, que são o objecto e a forma intencional e que constituem os elementos necessarios da verdade. Logo a verdade \emph{logica} encontra-se \emph{imperfeitamente} na \emph{idea} \footnote{A conclusão parece-nos certa. Se a verdade é uma adequação da. intelligencia com o objecto, e a iden é uma semelhança do proprio objecto, porque representa a \emph{essencia} d'elle, é claro que a intelligencia, informada\par pex\pgn{111}pela idea, se torna semelhante e adequada ao proprio objecto. Diz S. Thomaz: «Nos simples conceitos encontra-se a verdade, porque se conhece a coisa emquanto á sua essencia», (In L. IX, \emph{Met}. 1. 11). E noutro logar: « Ainda que o intelligivel incomplexo (a idea) não seja nem verdadeiro nem falso, comtudo a intelligencia, percebendo-o, é verdadeira \emph{emquanto se torna}, \emph{adequada a coisa entendida} ». (In 111 \emph{de Anima}, l. 11). --- Na idea, não ha apenas a verdade \emph{ontologica}, que é a \emph{essencia} do objecto, representada. pela propria idea; mas ba tambem a verdade \emph{logica}. porque a intelligencia conhece essa essencia, e, conhecendo-a, lhe se torna semelhante, ou adequada. O conhecimento será \emph{imperfeito}, porque tal essencia é representada de um modo confuso e indeterminado, mas é verdadeiro conhecimento. --- Tambem a conformidade ou adequação da intelligencia com o objecto, por meio da idea, é conhecida, de um modo \emph{directo}, pelo exercicio do proprio acto perceptivo, embora esse conhecimento não seja perfeito, como no juizo.\par Mas a verdade encontra-se \emph{imperfeitamente} na \emph{idea}, porque essa idea, embora represente, de um modo \emph{explicito}, a \emph{essencia} de um objecto, comtudo não representa, do mesmo modo, o \emph{ser}, ou a \emph{existencia}. Assim quando vejo um homem, e digo: \emph{homem}, esta idea de homem não diz \emph{explicitamente} se esse homem existe, uu não. --- Dissemos que a idea não representa o \emph{ser} das coisas de um modo \emph{explicito}, porque \emph{implicitamente} o representa. E isso é muito natural. Estando \emph{x essencia} de uma coisa unida intimamente a0 \emph{ser}, como a potencia está unida ao acto, é manifesto que não se pode conhecer x \emph{essencia}, Sem que, \&o mesmo tempo e de algum modo, seja conhecido o ser. ou a \emph{existencia}; porque as essencias não são conhecidas em \emph{abstracto}, mas em concreto, isto é, emquanto estão unidas com o \emph{ser}. Manifestando as suas \emph{qualidades} e a sua essencia, um objecto manifesta necessariunente, de algum modo, o seu \emph{ser}, ou a sua \emph{existencia}. Verifica-se isto não só com relação ás essencias das coisas \emph{concretas} e \emph{singulares}, mas tambem com relação ás essencias \emph{abstractas} e universaes, que têm um \emph{ser ideal}, fortnado pela razão. Percebendo-se, embora \emph{implicitamente}, 0 \emph{ser}, ou a \emph{existencia} das coisas pela idea e na \emph{idea}, e, dependendo a verdade de uma coisa mais do seu \emph{ser}, do que da sua. \emph{essencia}, á verdade encontra-se tambem na \emph{idea}, embora de um modo \emph{imperfeito}.\par Concluimos que tanto a \emph{idea} como o \emph{juizo} têm o mesmo \emph{objecto}; com a differenca que o \emph{juizo} representa e exprime, de um modo \emph{explicito}, a essencia € \emph{o ser} de um objecto, e por isso a elle compete \emph{perfeitamente} a verdade: € a \emph{idea} representa e exprime, de um modo \emph{explicito}, a essencia de um objecto, e, de um modo \emph{implicito}, o \emph{ser}, ou a \emph{existencia}, e por isso compele-Ihe a verdade \emph{imperfeitaomente}. --- Uma e a mesma impressão do objecto pode determinar e determina a intelligencia a formar uma \emph{idea} e a emittir um \emph{jo}. (Cf. in 1 \emph{Sent}., Dist. 19, q. 5, a. 1 ad 7).\par Mas a verdade, embor:. \emph{perfeitamente} se encontre no \emph{juízo}, e só \emph{inperfeitamente} na \emph{idea}, ou na \emph{simples apprehensão}, comtudo, pode separar-se\par \emph{E}\pgn{112}do \emph{juizo}, (pois ha muitos juizos falsos), e não pode separar-se da \emph{idea}, nem\par quando esta é \emph{incomplexa}, nem quando é \emph{complexa}. --- Não pode separar-se da \emph{idea}, quando esta é \emph{incomplexa}, isto é, quando denota a essencia das coisas, despida das notas individuantes. Por quanto, se o /\emph{wizo} pode discordar do objecto, emquanto pode affirmar que este é o que não é, ou que não é o que é, a \emph{idea}, quando é \emph{incomplexa}, não pode discordar do objecto; porque é impossivel que ella não represente. o que realmente representa; aliás representaria e deixaria de representar: o que é absurdo. Portanto, - assim como um objecto existente consta da sua entidade, assim tambem a idea consiste numa forma intencional, que é a expressão da coisa pensada e que não pode deixar de ser verdadeira. Assim a idea de \emph{Rkomem}, de \emph{virtude}, de \emph{sabedoria}, etc., pelo facto de representar simplesmente uma essencia, não pode deixar dé ter conformidade com a essencia, que representa. ---. -Nem pode - separar-se--da --idea,-quando esta é complexa. Por quanto, nós podemos conceber e concebemos um \emph{monte de ouro}; mas, se nos limitarmos a contemplar esse monte, e não affirmarmos, nem negarmos a sua existencia, a \emph{idea} nào poderá cbamar-se \emph{falsa}. --- Todavia, devemos admittir que a idea \emph{complexa} pode ser falsa, não por si, mas \emph{por} accidente, emquanto dá occasião a intelligencia de emitlir um juizo falso. Assim, quando nos representamos um \emph{monte d}' ouro, pode a intelligencia julgar que tal monte exista. realmente; Mas, então, \emph{o erro} está sempre no \emph{juízo}, e não na \emph{idea}.\par Pelo que temos dito pode avaliar-se a opinião de alguns \emph{meo-escholasticos}, quando dizem que na \emph{simples apprehensão}, ou nà \emph{idea}, não se encontra a \emph{verdade}, nem de um modo \emph{imperfeito}.\par (0A \emph{egualdade}, em que.propriamente consiste a razão ou a essencia da \emph{verdade}, deve existir --- não entre a intelligencia e \emph{o} objecto \emph{material}, isto é considerado: como é na \emph{reutidáde}. com \emph{todas} as suas motas, --- mas entre a intelligencia e o objecto \emph{formal}, isto é, considerado naquellas notas, que a intelligencia percebe e com as quaes se conforma. Ora a. verdade, emquanto se considera como uma \emph{egualdade} entre 'e intelligencia e o objecto \emph{formal}, não pode admittir varios graus, não pode ser mais ou menos verdade, por isso mesmo que duas coisas eguaes não podem ser mais ou menos eguaes, pois a egualdade està num ponto indivisivel; assim como duas linhas eguaes, emquanta são eguaes, não podem ser mais ou menos eguaes, e, se podem augmentar e iugmentam, não augmentam emquanto eguaes. --- Alem d'isso, a verdade, como dissemos, encontra-se perfeitamente no juizo, \emph{affirmativo} ou \emph{negativo}. \emph{No juizo affirmatiszo}, a. intelligencia \emph{afirma} a conveniencia entre o predicado sujeito; no \emph{negativo}, \emph{nega} aquella conveniencia. Ora a \emph{affirmacdo} e a ação são coisas indivisiveis, incapazes do mais e do\par menos. Logo a verdade \emph{logica é indivisivel}. --- Os que admittem graus na verdade, considerada sob este ponto de vista, confundem a \emph{egualdade} com \& \emph{semalhignga}. A semelhança aümiic graus, porque uma coisa pode ser}.

%%>>> t122

\tese{122}{A verdade logica é indivisivel.}
A verdade \emph{logica} é, como dissemos, a \emph{adequação}, ou a \emph{conformidade}, da intelligencia com o objecto; e por isso importa uma \emph{egualdade} entre o Juizo da intelligencia e o que se percebe do objecto. Ora a \emph{egualdade} \pg{111} \pg{112} , considerada como tal, é \emph{indivisivel}, não admitte graus: de , modo que duas coisas \emph{eguaes}, emquanto eguaes, não são mais ou menos eguaes, mas ou são eguaes, ou não o são. Logo a verdade \emph{logica}, emquanto. importa uma \emph{egualdade} entre o juizo da intelligencia e o que se percebe do objecto, é indivisivel (*).

P Dissemos --- emquanto importa uma \emph{egualdade}; \textbf{porque se} a verdade \emph{logica} se considera, não emquanto importa ama \emph{egualdade}, mas ou em relação á \emph{intfelligencia}, que . percebe, ou em

\textasciicircum{} . relação ao \emph{objecto}, que é percebido, então pode admittir varios graus, e um conhecimento pode ser mais ou menos verdadeiro :do que outro, --- conforme a maior ou menor perfeição da intelligencia, --- o maior ou menor numero das notas, que se percebem no objecto, --- a maior ou menor firmeza, com que a intelligencia, na percepção do objecto, pronuncia o juizo \footnote{\pgn{114}\emph{egualdade}. não 5 \emph{recebe o} mais e'o menos; --- mas que. se se considera o\par proprio \emph{ser} darcoisa; em-que-se-funda. a-verdade; então existe uma verdade \textbar{} (tanto \emph{ontologica}, como \emph{logica}) maior que outra, porque onde ha mais euti- .dade, ha mais verdade.(De \emph{Virtutibus}. q. 9, a. 9 ad 1). .\par Vejamos agora se-a \emph{falsidade} admitte, ou não, varios graus. A \emph{falsidade}, como dissemos, consiste na \emph{desconfor midade} da intelligencia com o objecto, é pode ser \emph{negativa ou positiva}. Ora, se a lalsidade \emph{negativa} não admitte graus, porque \emph{o não ser conforme} é uma coisa indivisivel, a falsidade \emph{positiva} admitte varios graus, não só com relação á \emph{intelligencia} e ao \emph{objecto}, mas tambem com relação á propria essencia ou razão de \emph{falsidade}. --- 1.º) Admitte vários graus na propria razão de \emph{falsidade}. Porquanto, a \emph{desconformidade} entre, a intelligencia e o objecto pode ser mais ou menos pronunciada. Diz S. Thomaz: «Quanto magis a veritate receditur, tanto major est fajsitas ; non enim est tam magnu falsitas aeslimantis \emph{fria} esse quinque, sicut ejus." qui aestimat \emph{tria}. esse \emph{cenfum}» (\emph{C}. Gentes, 1. 3, c. 139) A razão é manifesta. A falsidade não é apenas a falta de verdade, mas é . o \emph{contrario} da verdade; ora entre dois termos contrarios, ou oppostos, ba muitos pontos intermedios. Ha, porem, uma excepção, e dá-se quando a falsidade. do juizo consiste na mera negação: de alguma verdade; assim a falsidade d estes juizos: --- Deus não existe, o mundo não existe --- não admitte\par graus. --- 2.º) Admitte varios graus com relação a \emph{intelligencia}, Com efieito, uma intelligencia obtusa, ou totalmente dominada pela vontade, pode cahir e cahe em erros, mais grosseiros do que os erros, em que pode cahir e cahe uma intelligencia perspicaz, ou menos sujeita á influencia da vontade. --- 3.0) Admitte varios graus com relação ao \emph{objecto}, --- ou emquanto a um objecto podemos attribuir um maior ou menor numero de notas, que lhe vão convém, ou negar-lhe um numero maior ou menor das que lhe são proprias : ; --- ou emquanto um juizo falso em materia necessaria e \emph{inmubavel} é mais. falso que um juizo falso em materia \emph{contingente} e \emph{mutavel}; assim o juizo --- \emph{Deus não existe} --- é mais falso que o juizo --- \emph{o mundo} não \emph{existe}.}.

%%>>> t123

\pg{113} \pg{114} \tese{123}{A verdade logica € immnutavel.}
A verdade \emph{logica} pode exprimir --- ou a \emph{essencia} dos entes, --- ou um \emph{facto contingente}. Ora, tanto: num como noutro caso, ella é 4\emph{mmutavel}.

a) É immutavel, quando exprime a \emph{essencia} dos entes. --- Com effeito, a verdade \emph{logica}, consistindo na adequação da intelligencia com o objecto, deve corresponder á indole do mesmo, objecto; de modo que ella será immutavel ou mutavel, conforme immutavel ou mutavel fôr o objecto. O objecto da nossa intelligencia é a \emph{essencia} tanto do Ente infinito, como dos entes finitos. Ora a \emph{essencia} do Ente infinito é absolutamente immutavel ; e a \emph{essencia} dos entes finitos, embora seja mutavel quanto aos elementos \emph{reaes} (assim a essencia do homem é mutayel quanto ào \emph{corpo} e a \emph{alma}), comtudo é imuiufavel quanto aos elementos \emph{especificos} (assim 2 a essencia do homem é immutavel quanto

4 \pg{115}  a \emph{animalidade} e a \emph{racionalidade}). Logo a verdade \emph{logica}, quando exprime a \emph{essencia} dos entes, é \emph{imumulavel} \footnote{A verdade \emph{logica} da, nossa intelligencia. funda-se . . Da verdade .\emph{ontologica} das coisas percebidas. Por isso, como dizemos no texto, a verdade \emph{logica é immutavel} ou \emph{mutavel}, conforme immutavel ou mutavel é a verdade \emph{ontologica}. Ora a verdade \emph{ontologica} pode referir-se a Deus e ás creaturas. .\par Deus, e só Elle, é essencial e absolutamente immutavel; porque toda a mudança denota imperfeição, e Deus é a perfeição infinita. Se Deus é absolutamente immutavel, tambem immutavel é o conbecimento, que temos d'Elle. Por quanto, embora este conhecimento seja infinitamente inadeguádo.e deficiente, comtudo \emph{o} que atiribuimos a \textsc{Deis}, realmente se encontra nElle, embora de um modo infinitamente elevado e superior á nossa comprehensão. Assim, quando dizemos que Deus é bom, é justo, etc., estas perfeigóes absolutas existem realmente em Deus; e \emph{o} nosso conhecimento, embora seja inadequado e possa tornar-se mais perfeito, é todavia verdadeiro e immutavel. Concluunos que a verdade \emph{logica}, da nossa intelligencia, relativa a Deus, é immutavel, porque immutavel é a verdade \emph{ontologica} do proprio Deus.\par Se de Deus descemos ás creaturas, encontramos. em. todas ellas.a mudança. Os entes espirituaes, se não estão sujeitos á mudangas subsianmciaes, pois não são compostos de materia, estão sujeitos-a2- mudanças-accidentaes; emquanto \textbf{passam} da potencia para o acto operativo, porque não. s\&o'a propria actualidade. Os seres \emph{maleriaes} estão sujeitos a mudanças substanciaes e \emph{accidentaes}. --- Mas, embora todos \emph{os} sores creados estejam Sujeitos á mudanças, comtudo ha em todos uma coisa absolutamente immutavel, que é \emph{essencia}, considerada nos seus elementos \emph{especificos}. Se a essencia das creaturas, --- considerada nos seus elementos reaes, pode estar e está sujeita -- * mudanças, mesmo \emph{substanciaes}, --- considerada nos seus elementos \emph{especificos}, não está sujeita á mudança, mas é necessaria, immutavel ; assim 2 essencia do homem, se, considerada nos elementos reaes, que são a alma e O \emph{corpo}, é contingente e mutavel, --- considerada nos seus elementos \emph{especificos}, que São a \emph{unimalidade} e a racionalidade, € necessaria e immutavel; porque, embora o homem nunca tivesse existido, é e será sempre verdade que elle € \emph{animal racional}. D' ahi se segue que, se a verdade logica da nossa intelligencia, --- emquanto representa o que ba de mutavel nas creaturas, é mutavel, --- emquanto representa o que ba de immutavel nas mesmas creaturas, que é a \emph{essencia}, é immutavel. --- A immutabilidade das essencias creadas tem, como veremos, a sua razão suprema, o seu ultimo fundamento nà Essencia divina (Cf. in IV \emph{Ethic}, 1. 3).\par Immutavel quando exprime a \emph{essencia} das coisas, a verdade é immutavel tambem quando exprime 0s \emph{primeiros principios da razão}, que se formum immediatamente, \emph{e} por simplicissima abstracção,. depois da apprehensão da \emph{essencia} e que a \emph{ella} se referem e nella se fundam. Com eifeito, estes principios: --- uma coisa não pode simultaneamente ser e deixar de ser --- todo o effeito exige uma causa, ele., mostrum-se dotados de tal \emph{wnivcersali}- \textasciicircum{}. \emph{dade} e \emph{necessidade}, que não é possivel conceber o contrario, sem se renunciar 30 senso commum. Ora um principio \emph{universal} e \emph{necessario} é naturalmente \emph{imuuiavel}. --- Os \emph{prüneiros principios} são 0 lundamento de todas as sciencias, cojó objecto é \emph{constituido por} verdades \emph{tambem} :\emph{niversaes} c}.

\emph{H}

b) É \emph{immutavel}, quando exprime um \emph{facto contingente}. --- Por quanto, a verdade \emph{logica}, como dissemos, deve corresponder á indole do objecto; com o qual deve conformar-se a intelligencia. Ora um "facto \emph{contingente}, embora, por si, podesse deixar de dar-se, comtudo, depois que se deu, é absoluta e eternamente immutavel; porque é impossivel. que, tendo-se dado, não se tenha dado, ou não se tenha dado pelo modo, por que se deu. Sendo \emph{imnmutavel} o facto, é immutavel a relação dà nossa intelligencia com elle, e. por isso é \emph{immnutavel} a verdade. Logo a verdade \emph{logica}, quando exprime um facto \emph{contingente}, \emph{é âmmutavel U}). necessarias, que, por isso, se resolvem naquelles principios. . Todas as Sciencias têm por fim a descoberta, de taes verdades, e, depois de as terem descoberto, nem sequer duvidam de que ellas possam estar sujeitas a mudanças, ou possam tornar-se falsas. Nenhum eseriptor, que tenha ainda o juizo são, poderá sustentar que 9 mais 2 sejam eguaes a 5, ou que os tres angulos de um triangulo não sejam eguaes a. dois anguios rectos. Diz \emph{Cl}. \emph{Bernard} : « Nas mathematicás, quando se parle de um axioma, ou. principio, cuja verdade € absolutamente "hecessaria-e"consciente;:a liberdade deixa de existir, as verdades adquiridas são immutaveis. O geometra vão é livre de duvidar se .08 tres angulos de um triangulo São equaes-a dois rectos; e por isso não é livre de rejeitar as consequencias logicas; que se deduzem d' este principio... Sendo as verdades mathematicas immutaveis' e absolutas, a sciencia augmenta pela simples e successiva juxtaposicio de todas as verdades adquiridas» (\emph{Inirod}. \emph{é la medicine} experim. p. 65).\footnote{Assim esta verdade --- 'S. Jounna d'Arc derrotou os Inglezes em Orléans --- exprime um \emph{facto contingente}, que se deu mas podia não se ter dado, e todavia é \emph{immutavel}; porque. tendo-se dado. é e será eternamente verdade que S. Joanna d' Arc derrotou os Inglezes em Orléans. Se não fosse assim, se esse [acto contingente não fosse immutavel, deveriamos admittir ' que S. Joanna d'Arc derrotou e.não derrotou os lnglezes em Orléans: o que é absurdo (Sum. \emph{Th}.; p. I, q. '85 a. 1 ad 3).\par Agora uma advertencia. Quando Se diz que a verdade \emph{logica}, exprimindo um \emph{facto contingente}, se muda. pelo facto de se mudar o objecto material. não devemos entender que a verdade se transforma em falsidade, mas que a uma verdade succede, uma oultra/verdade, --- ambas immutaveis. Assim quando, vendo \emph{Pedro} de saude, digo: Pedro é são, e, depois, vendo-o: doente, digo: Pedro não é são, não podemos dizer que um juizo é verdadeiro e outro é falso, ou que a verdade, pela mudança do objecto material, se transformou em falsidade; mas devemos dizer que a uma verdade succedeu outra verdade, que ambos os juizos são verdadeiros, porque ambos '. correspondem a realidade. Dá-se a \emph{conéradicgao}, quando ao mesmo sujeito se attribue e se nega o mesmo predicado, sas \emph{ao} mesmo tempo e sob \emph{o} mesmo \emph{respeito}: \emph{o} que não se verifica no presente exemplo. E ambos esses juizos são \emph{tmmutlaveis}; porque é e será eternamente verdade que Pedro num certo tempo estava são e num outro tempo não estava são.\par Come se vé, consideramos a verdade \emph{logica}, emquanto \emph{objectiva}; porque;\par \textasciicircum{}\pgn{117}da immutabilidade do \emph{objecto} deduzimos a immutabilidade da \emph{verdade} --- Mas, se a verdade se considerar” emquanto- é \emph{subjectiva}. então é mutavel,como é mutavel a nossa intelligencia, embora o objecto permaneça immutavel. A nossa intelligencia, por ser limitada, pode passar, em relação á verdade, de um estado para outro diverso e até contrario, de um conhecimento verdadeiro para um outro falso, e viceversa, É uma \emph{fallibilidade}, não \emph{essencial} nem \emph{natural}, mas \emph{accidental} e eventual. (\emph{Sum}. \emph{Th}., p. 1, q. 16, a. S).\par A \emph{immwutabilidade} da verdade é negada por todos os sequazes do \emph{relativismo}. O nosso espirito, na opinião d'elles, não é perfeito quanto. á intelligencia, e não pode conhecer absolutamente a objectiva verdade das coisas. Tudo o que conhecemos é \emph{relativo} aos nossos sentidos, 20 nosso espirito, e a nossa experiencia (Comte), a qual, por sua vez, é relativa á nossa constituição (Kant). A verdade, portanto, não é uma \emph{coisa}, É uma vida, e, como tal, segue o destino da nossa existencia. Ora o espirito humano não :. é immutavel, mas evolve-se e muda-se constantemente. A estas evoluções e. mudanças do espirito. correspondem. as evoluções ..e.. mudanças: da verdade. - Portanto, nada é absoluto, isto é, nada é fixo, immutavel e do nosso pensamento; mas tudo é relativo, isto é, tudo € incerto, variavel e dependente do nosso espirito. D' onde se segue --- que os systemas contradictorios são apenas \emph{momentos diversos} da mesma \emph{verdade}, --- que uma coisa, hontem verdadeira, mesmo para todo o genero humano, pode ser falsa amanhã, --- que uma e a mesma coisa coisa pode ser, ao mesmo tempo e sob o mesmo aspecto, verdadeira e falsa, conforme a uns parece verdadeira e a outros falsa. Concluem que não devemos acceitar uma verdade senão \emph{condicionalmente}, isto é, senão com a condição que ella se acceita até que não se conhece verdadeira a proposição opposta.\par Esta opinião, tão antiga como Protagoras (que dizia ser a nossa intelligencia á \emph{medida} de todas as coisas), é sustentada per Comte, Bain, Spencer, Rénouvier, pelos Kantistas, positivistas e modernistas, e pelos sequazes da \emph{philosophia} nova (Bergson, Le Roy, Wilbois). --- Não sendo possivel citar... tudo o'que esses escriptores disseram, limitamo-nos a transcrever uns: bre-. ' vez trechos de Hamilton, de Schérer, e de Stuart-Mill.\par Hamilton diz: «O \emph{absoluto} é uma \emph{pseudo-idea}... O \emph{absoluto} é uma idea \emph{negativa}, é o \emph{não-relativo}. isto é, é a negação do \emph{relativo} ». - ---Schérer escreve: «Um principio apoderou-se do espirito moderno. e\par devemo-lo a Hegel. Entendo fallar do principio, em força do qual uma proposição não é mais verdadeira que a sua opposta.. Isto quer dizer que tudo é relativo e que os juizos absolutos são falsos. Esta descoberta do caracter \emph{reloficó} das verdades é o facto capital da historia do pensamento contemporaneo. Hoje, \emph{para nós}, uma coisa não é verdade, nem erro. É necessario inventar outras palavras. Em toda a parte só vemos \emph{gradações}; chegamos a admittir \emph{a identidade dos contrarios}. Conhecemos não a religião, mas as religioes ; não a moralidade, mas os costumes: não os principios, mas os factos, \emph{Tudo é relativo}; mais ainda, \emph{tudo} é apenas uma relação».\par Stuart Mill accrescenta : « Pode ser que noutros planetas, ou nas partes\pgn{118}- ainda desconhecidas do nosso globo, \emph{haja outra logica}. Quem sabe se \emph{amav} nhà, se d'aqui a pouco haverá outra seiencia, seja ella qual fôr, se haverá ."--- - . Quas coisas semelbantes, se haverá coisa \emph{alguma} ? ». Lu Examinemos os argumentos dos, \emph{relativistas}. --- : a) \emph{O absoluto}, dizem elles, \emph{não existe}; \emph{tudo é relativo}... DU Resposta. Mas o que é o \emph{absoluto}? o que é o \emph{relativo}? Se interrogarmos T o senso commum, este nos dirá --- qué o \emph{relativo} é o que se refere a um C outro, de modo que, sem este, não pode existir, nem pode conceber-se, --- e Rs que o \emph{absoluto} é o que é independente \&\& este vinculo de relação, de modo - que pode existir, e pode conceber-se, sem que se deva recorrer a outro ente. - E, como tal independencia admitte varios graus, assim ha varias formas \emph{I} de \emph{absoluto}, uma, mais perfeita que outra. Por isso pode chamar-se --- \emph{abso}- \emph{E luto} tudo o que exisle em si, e não depende do sujeito, que o conhece, --- --- e \emph{relativo} tudo o que existe exclusivamente no sujeito e pelo sujeito. Mais - \textasciicircum{} perfeito, infinitamente mais perfeito é o \emph{Absoluto}, que não depende de al- 77 guma hypotbese, mas até repelle toda a hypothese, toda a restricção, todo \emph{TUM} o limite; e, comparado com- Elle; todo: e- qualquer “outro \emph{ser}: é relativo, isto L é, \emph{contingente}, \emph{limitado}. --- Agora pecguntamos:. Existirá exclusivamente ó \emph{E relativo}. ou deverá existir tambem o absoluto? Se não quizermos cahir \emph{t} numa contradicção pueril, devemos responder que. se existe o \emph{relativo}, isto \emph{v}\textasciicircum{} é, se existe o ser, que se refere a um outro, deve forçosamente existir este \emph{s} outro, a que o. primeiro se refere \emph{e} que se diz e é \emph{absoluto}. O: \emph{relativo}, \emph{E} emquanto tal, suppõe o \emph{absoluto}; quem rejeita este, deve. rejeitar aquelle. = Logo, ou não existe coisa alguma, ou existe tambem o \emph{absoluto}. O proprio = H. Spencer, evolucionista, diz: «O conceito de \emph{relação} traz comsigo o \emph{cono} ceito dos seus dois termos. Pretender que concebamos a relação entre o . . relativo e o n\&o-relativo, sem termos consciencia de cada um dos dois termos, \emph{n}» é pretender que comparemos uma coisa, de que temos consciencia, com uma ox outra, de que nà\& a temos: a comparação é um acto de consciencia, que \emph{EU} não pode realizar-se, se não tivermos consciencia de ambos os termos. \emph{o} Sendo assim, que juizo" deve. fazer-se--da "affirmação --- que o absoluto só --- se concebe como uma mera negação da concebibilidade? Se o não relativo " ou o absoluto não é presente ao pensamento senão como uma mera nega- = ção, torna-se inintelligivel a relação entre elle e o relativo, porque um dos i termos da relação não se encontra na consciencia. Se a relação € \emph{ininiellim} givel, inintelligivel é o proprio relativo ; d' onde a impossibilidade de todo - o pensamento » (\emph{Premiers principes}. p. T, c. 4.\par \emph{b}) Os \emph{relativisías} replicam : O \emph{absoluto}, se existe, é \emph{éncognoscivel}. \emph{b}---- Resposta, E será elle £\emph{ncognoscivel}, porque é \emph{absoluto} ?... Porque, respon- \textasciicircum{} dem, a experiencia limita-se a observar os factos, os phenomenos, e não \emph{id} penetra até ao fundo das coisas, para ahi descobrir a \emph{essencia}... A resposta \textbar{}\emph{o} não satisfaz, Porquanto, ha duas especies de experiencia: uma \emph{eulgar}, outra d \emph{scientifica}. Se a \emph{vulgar}, que é commum a todos os homens, não reflecte, e \emph{NN} limita-se a attestar os factos, a \emph{scientifica}, que é "propria dos sabios, não \emph{jc} se limita a observar os factos, mas dos factos sobe ás causas, e das \emph{prouU} priedades a essencia.--- \emph{De Dominicis}, eseriptov insuspeilo, commentando as\par citadas\pgn{119} palavras de \emph{Spencer}, conclue: «O \emph{absoluto} existe, e a sciencia não pode prescindir d'elle. A sciencia: exige- imperiosamente- que.-além- da- existencia do phenomeno, se admitta qualquer outra coisa, e que, sen esta qualquer outra coisa, o conhecimento não é fundado, não é seguro, \emph{não é racional}» (\emph{La dolirina della evoluzione}, vol. 9, p. 56). --- No decurso d' este tratado teremos occasião de voltar ao assumpto.\par \emph{c}) Os adversarios accresceptam : Embora não seja \emph{relativo} tudo o que\par existe, todavia é \emph{relativo} todo o conhecimento bumano, é por isso é \emph{relativa} toda a \emph{verdade} da nossa intelligencia.\par Respondemos que \emph{o} conhecimento se pode chamar \emph{relativo} : á nossa intelligencia num sentido, e nio se pode chamar \emph{relativo} num outro.\par Pode chamar-se \emph{relativo} á nossa intelligencia, emquanto um acto de uma faculdade deve revestir e reveste a condicio, o caracter. \& natureza da propria faculdade. E nada mais razoavel. O acto é sempre proporcionado á faculdade, de que deriva. Por isso, o conhecimento não só é relativo á intelligencia dos diversos individuos bumanos, mas tambem a propria in- telligencia humana, considerada em. si mesma; de modo que: assim-como-a intelligencia humana é essencialmente. finita e limitada, assim tambem dequado e deficiente é o conhecimento que temos das coisas, mesmo seus veis e naturaes... Este conhecimento pode dizer-se tambem \emph{relativo}, isto é, imperfeito, limitado, contingente, em opposição ao conhecimento \emph{absoluto}, isto éperfeitissimo, infnito, necessario, que é exclusivamente proprio de Deus. --- Mas, que se seguirà logicamente d'esta \emph{relatividade}? Seguir-se-ha, por venlura, o que dizem os adversarios, isto é, que o8 nossos conhecimentos não correspondem aos objectos, que são mutaveis, que são falsos? Não; porque, como dissemos, o conhecimento inadequado não é um conhecimento falso; e uma verdade será sempre verdade, fixa e immutavel, mesmo quando a intelligencia, conformando-se cada vez mais com o objecto, chegará a descobrir novas notas, ou a conhecer de um modo claro e distincto uma nota, que antes tinba sido percebida de um modo obscuro e confuso. - Mas não é este o sentido dos adversarios, quando dizem que o conhecimento é \emph{relativo}. Entendem elles, por tal expressão, que o conhecimento é \emph{inteiramente subjectivo}, isto é, que é apenas uma forma da nossa intelligencia, \& qual não corresponde realidade alguma. --- A affirmac\&o não poderia ser mais absurda. De facto, o nosso conhecimento, --- se é \emph{subjectivo}, emquanto é um acto \emph{immanente}, que tem o principio e o termo no mesmo sujeito, --- é \emph{objectivo}, emquanto á \emph{origem} e a \emph{representação}, porque é causado pelo objecto e representa um objecto. E, como o objecto tem o ser independente do sujeito, segue-se que o nosso conhecimento, emquanto exprime o objecto, é fixo e immutavel, porque fixa e jmmutavel é a essencia. do proprio 'objecto. E, embora o objecto conhecido \emph{exista} na intellimeneia por um modo differente d'aquelle, por que existe \emph{na} sua realidade, comtudo não deixa de ser o que é, não perde nem muda a propria essencia. É, pois, absurdo dizer que. mudando-se o sujeito, se muda tambem o objecto, e que o que é verdade para \emph{um individuo}, pode ser falso para um outro, e viceversa. --- Os \emph{relativistas} esquecem uma coisa muito elementar --- que, no conhecimento,\pgn{120}fica e determina o sujeito; porque o sujeito é movido e passivo, ao passo que o objecto é movente e activo. As coisas são o que são em si mesmas, e não o que nos possam parecer; uma coisa, que é verdadeira para uma intelligencia, deve ser verdadeira para iodas as intelligencias. O principio de contradicção (uma coisa não pode simultaneamente ser e deixar de ser) é uma verdade do senso commum, e não pode ser contestada.\par OS. P. Pio X, no decreto « Lamentabili » de 3 de Julho de 1907, condemnou o \emph{relativismo}, contido nesta proposição (58º): « Veritas non est immutabilis, plusquam ipse homo; quippe quae cum ipso et per ipsum evolvitur».}

%%>>> t124

\pg{116} \pg{117} \tese{124}{A verdade logica é unica e multiplice.}
. A verdade

. - * . " . \emph{logica} pode considerar-se em Deus e na intelligencia \emph{creada}. --- \pg{118}  \textbar{} --- Considerada em Deus, a verdade \emph{logica} é unica, porque unica C é a Essencia divina, com a qual as propriedades e as operações \pg{119}  divinas se identificam. Esta verdade \emph{logica} divina é a causa \emph{exemplar}. de todas as verdades \emph{logicas} das intelligencias creadas.

--- Considerada na intelligencia \emph{creada}, a verdade \emph{logica} é, \emph{mulliplice}, porque são tantas as verdades, quantos são os conhecimentos conformes com a realidade,. que existem nas intelligencias das creaturas e que são outros tantos reflexos deficientes da unica verdade divina C )

*AnmTIGO IT

\emph{Certeza}, sua \emph{divisão}, \emph{existencia e propriedades}

125. Estados da intelligencia com relação á verdade: --- A nossa inteligencia alcança perfeitamente a \emph{verdade}, quando forma \emph{juizos}, \emph{correspondentes d realidade}. Ora essa faculdade, com relação á verdade, pode; encontrar-se em \emph{quatro estados}. ' Por quanto, --- ou não forma um juizo, porque não tem as ideas respectivas, e está na \emph{igmorancia}, --- ou possue a idea de um objecto, mas não se decidé a pronunciar um juizo, e está na duvida, --- ou.pronuncia um juizo, mas teme que se possa verificar .o. contrario, e. está na \emph{opinião}, --- ou julga sem recejo

\textbar{} algum de errar, e está na \emph{certeza}. --- É da \emph{certeza} que nos devemos occupar neste. e nos seguintes artigos \footnote{Os \emph{ontologistas},. depois de Malebranche. dizem que existe uma unica verdade --- Deus, e que por isso a nossa inteligencia não pode contemplar a verdade senão em Deus. --- Este erro relüta-se com o criterio exposto no texto. Embora a nossa: intelligencia seja uma participação da intelligencia divina, assim como os seres creados são uma. participação do ser divino, comtudo a nossa intelligencia é realmente distincta da divina, como .realmente distinctos do ser divino são os seres creados. Por isso possuimos uma verdade \emph{propria}, ainda que participada e accidental. e esta verdade é \emph{multiplice}, porque muitos e diversos são os conhecimentos, que temos dos muitos e diversos seres (\emph{Sum}. \emph{Th}.. p. l, q. 16, a. 6).}. não é o sujeito que forma ou modifica o. objecto, maséo objecto que modi-\footnote{São \emph{quatro} os estados da intelligencia com relação á verdade : \emph{ignorancia}, \emph{duvida}, \emph{opinião} e \emph{certeza}. Devendo tratar detidamente da \emph{certeza}, digamos alguma coisa ácerca dos tres primeiros estados.\par 2) ignorancia. --- \emph{Tgrnorancia} é o estado da inteligencia, quando não\pgn{121}profere um juizo; porque não tem as ideas necessarias. --- É \emph{universal} ou \emph{particular}, conforme a. ausencia dos conhecimentos € \emph{total} ou \emph{parcial}. --- A --- \emph{particular} pode ser mais ou. menos extensa. --- A ignorancia, como se vé, \emph{AL} é uma relação \emph{negativa}, isto é, é uma. \emph{falta} de relação da intelligencia para --- com a verdade. É effeito. da limitação. e da. fraqueza. da .nossa intelligencia,. “é da dilficuldade ou obscuridade dos objectos copuosciveis. - Muitas vezes, n \emph{ignorancia} pode sev-uma-culpa; como quando-é- a--ausencia-de conheci- mentos indispensaveis, sobretudo dos relativos á ordem moral e religiosa. --- É uma das mais deploraveis consequencias do peccado original... -\par b) Duvida. --- Duvida a o estado da inteligencia, quando não se decide \emph{a} emittir um juizo, affirmativo ou negativo. --- Divide-se em negativa e \emph{positiva}. \emph{É negativa}, quando não ha razões nem para 2 affirmação nem para a negação; é \emph{positiou}, quando ha razões para a affizmacüo e para a vegação, mas de \emph{egual} peso ou valor, de modo que a intelligencia esta suspensa. --- Toda a \emph{duvida}, quando \emph{deliberada}, suppõe um juizo ácerca do valor das razões \emph{pro} e \emph{contra}. Tal juizo, embora não attinja o proprio objecto, pode ser imprudente ou erroneo. como qualquer outro.\par €) Upinião. --- \emph{Opinião} é o estado da intelligenei», quando pronuncia um juizo, mas com \emph{perplexidade} e \emph{receio} de \emph{errar}. --- As razões, que neste caso determinam a.intelligencia, são mais ou menos provaveis. de modo que não tornam o juizo emillido.absolutamente verdadeiro, nem o contradictorio absolutamente falso. --- A \emph{probabilidade}, pois, é a somma.das razões serias e atteudiveis, que sustentam uma proposição. Admitte varios graus ; é maior . ou menor, conforme o numero e o peso das razões. Ha opiniões \emph{provaveis}, \emph{mais provaveis}, \emph{probabilissimas}. Mas tudo o que está na esphera da \emph{probabilidade} não é \emph{certo}. Por isso uma opinião \emph{probabilissima}. não é \emph{certamente verdadeira}, assim como uma opinião menos \emph{provavel} não é \emph{certamente falsa}, --- A' \emph{opinião} reduz-se a \emph{suspeita}; a qual, embora não seja um juizo, comtudo é uma inclivação para elle, por causa de alguns indicios ou motivos, que se descobrem no objecto. Emguanto não é juizo, a \emph{suspeita} ditfere da \emph{opinião}, que é um juizo, embora não firme: emquanto é uma \emph{inclinação para o} juizo, é um inicio de \emph{opinião} (\emph{Sum}, \emph{Th}., 11-11, q. 2. a. 1).}

196. Certeza. --- \emph{Certeza} é a \emph{firme adhesão da intelligencia}, = á \emph{verdade conhecida}, \emph{fundada}. \emph{num}. \emph{motivo capas} de \emph{excluir todo} o \emph{receio de errar}. --- A \emph{certeza} --- «) é adhesão \emph{firme}, porque a O intelligencia \emph{descança} na posse da verdade, sem ré receio de errar; --- = \emph{b}) é adhesão a \emph{verdade} conhecida, porque só a \emph{verdade}, \emph{conhe}- \emph{cida} como \emph{tal}, possue a força de attrahir irresistivelmente e \emph{m} socegar a intelligencia; --- c) é-adhesão \emph{fundada} mum motivo \emph{T capaz de excluir lodo o} receio \emph{de errar}, porque a \emph{firmeza} da adhesão é um effeito, que deve ter um motivo proporcionado. \emph{m} \textbar{} --- Como se vé, a \emph{certeza}, tomada no sentido rigoroso, é consti- = tuida por elementos \emph{subjectivos} e \emph{objectivos} \footnote{A \emph{adhesão firme} da intelligencia é o elemento \emph{subjectivo} da \emph{certeza}; a verdade \emph{conhecida}, e O \emph{motivo}, que determina a \emph{firmeza} da adhesão, são\pgn{122}os elementos \emph{objectivos}: --- "A firme \emph{adhesão} da intelligencia á verdade é o objecto \emph{material da} certeza; \emph{o motivo}, que produz essa. firme adheésão, é o objecto \emph{formal}. --- \&. Thomaz, diz que a \emph{certeza é}: \emph{firmitas cdhaesiouis Viríutis cognoscitivae} in suum \emph{cognoscibile} » (in TII \emph{Sent}., Dist. 26, q. 2, a. 4).\par A \emph{cerieca} é uma perfeição intrinseca do proprio acto cognoscitivo, mas emquabto este exprime uma \emph{verdade}, --- vão só porque, como dizemos no 'texto, só a verdade tem a força de prender e encadear a intelligencia, de modo que não oscille para o contrario, --- mas tambem porque uma faculdade não acha a perfeição e o descauco senão no objecto \emph{proprio}, e o objecto \emph{proprio} da intelligencia é a \emph{verdade}, ou o ente, emquanto \emph{verdadeiro}. --- Por isso --- 1.*) a adhesão ao \emph{erro}, o quai por falta da devida consideração pode ser recebido como verdade, não merece o nome de \emph{certeza}, porque não passa de uma \emph{opinido}, e, muitas vezes, nem é \emph{opinião}, mas pertinacia Sum. \emph{Th}., 1-11, q. 18, a. 4; --- \emph{De} Ver, q. 6, a. 3); --- 2.º) a firmeza d' esta adhesão deve fundar-se em motivos \emph{solidos} e \emph{seguros}, de modo que a proposic\&o, lida \emph{por-certãa}, seja. incompativel. com:.a2. sua contradictoria, e assim seja excluida a possibilidade do, erro; --- 3.9) assim. como para a perfeição da verdade \emph{logica} não basta que a intelligencia se conforme com o objecto, 'mas é necessario que conheça ainda essa conformidade, assim tambem para a. perfeição da. \emph{certeza} não basta que a intelligencia adhira sem hesitação a verdade conbecida. mas deve ainda conhecer. claramente, mesmo no proprio acto do juizo, que esse adhesão é incompativel com o.erro.\par Sendo a \emph{certeza} uma adhesão da intelligencia á verdade, é claro que deve encontrar-se onde e como se encontra a verdade, e por isso no \emph{juízo} de um modo \emph{perfeito}, e na \emph{idea}, ou simples apprehensio, de um modo \emph{imperfeito}, ou. \emph{inicial}. Quando um objecto se manifesta convenientemente a iotelligencia e imprime nella a propria imagem, despida das condições materiaes ou individuaes, a intelligencia, reagindo \textbf{nessa} imagem, forma a idea do objecto, que a impressionou, e assim alcança, embora de um modo deficiente, a verdade, e x esta adhere sem hesitação. --- Como é natural. os escriptores, qae não reconhecem. na \emph{idea} um vislumbre de \emph{verdade}, nem reconhecem nesse:acto um vislumbre de \emph{certeza}.\par Os philosophos apresentam outras definições da \emph{ceríezo}. Uns dizem que a \emph{certeza} é a consciencia \emph{da verdade}; --- outros, que é a \emph{impossibilidade da duvida}; --- outros, que é \emph{a verdade} cabalmente \emph{demonstrada} ; --- outros, que é a verdade \emph{conhecida nos} seus \emph{principios}.\par Estas definições são todas inaeceitaveis. --- É inacceitavel a \emph{primeira} ; não só porque não pode applicar-se a certeza \emph{directa}, que é verdadeira \emph{certeza}, mas tambem porque não exprime os elementos essenciaes da certeza. --- E' inaeceitavel a segunda: porque, alem de ser \emph{negativa}, não manifesta o que a certeza é em si mesma. --- É inacceitavel a. \emph{terceira}; porque não abrange as verdades conhecidas \emph{immediatamente} e que são as mais certas. --- É inacceitavel a quarta; porque, para estarmos certos de uma verdade, não é preciso conhecel-a sempre nos seus \emph{principios}, pois as vecdades de \emph{facto} são percebidas \emph{imunediatamenic} na sua realidade e são certissimas.}. \emph{p}

197. Especies de certeza. --- A \emph{certeza} pode considerar-se --- nos \emph{entes}, em que reside, ---- nos \emph{motivos}, em que se funda, --- \pg{123}  10 \emph{modo}, por que se conhecem os motivos, --- no \emph{principio}, de. jue deriva. --- Considerada nos \emph{entes}, em que reside, é \emph{subjetiva} e \emph{objectiva}. Considerada nos \emph{motivos}, em que se funda, é \emph{nelaphiysica}, \emph{physica} e moral. Considerada no \emph{modo}, por que se conhecem os motivos, é \emph{directa} e \emph{refleza}. Considerada no. .. principio, de que deriva, é \emph{natural} e \emph{sobrenatural}.

a) \textbf{Certeza} subjectiva e objectiva. --- A \emph{subjectiva} é a propria \emph{firmeza}, com que adherimos a verdade. Diz-se \emph{subjectiva}, porque reside no \emph{sujeito}. --- A \emph{objectiva} é a propria \emph{verdade}, que se manifesta «claramente á intelligencia, produzindo uma adhesão firme e excluindo todo o receio. de errar. Diz-se \emph{objectiva}, porque está no objecto. --- A \emph{objectiva} é causa da \emph{subjectiva} (9.

b) Certeza metaphysica, physica e moral.

\emph{a}) \emph{Certeza metaphysica}. --- Certeza \emph{metaphysica} é a que se baséa na \emph{essencia} das coisas; tal é a certeza d'esta proposic\&o: \emph{todo o effeito exige} uma: causa. --- "Esta. certeza é: \emph{absoluta} .--- não admitte excepções, porque.a essencia das. coisas immutavel. = --- Por isso a proposição, \emph{contradictoria} a que exprime esta certeza, é \emph{absolutamente} impossivel \footnote{Segundo \emph{Kant}, a certeza \emph{subjectiva} é a que se refere aos phenomenos \emph{internos} do \emph{sujeito.e} a \emph{objectiva} é a que tem por termo os objectos \emph{externos}. --- Mas a certeza é primaria e propriamevte \emph{subjectiva}, e só secundaria e impropriamente se atiribue ao \emph{objecto}, emquanto este é capaz dc produzil-a no, espirito humano (Sum. \emph{Th}, ilf. q. 18, a. 4). --- A certeza \emph{objectiva}, pode tambem defiuir-se : a \emph{infallibilidade} ou \emph{força}, de que é \emph{dotado o} motivo \emph{para} conseguir um \emph{assenso certo}. Todavia essa força pode deixar de produzir o seu effeito, por falta de consideração. ;}. \textbar{} :

\emph{b}) \emph{Certeza physica}, --- Certeza \emph{physica} é a que se funda na \emph{constancia} das leis \emph{physicas}; tal é a certeza da seguinte proposição: a \emph{terra move} se em \emph{volta do sol}. --- Esta certeza admitte excepções, porque as leis \emph{physicas} podem ser suspensas pela Omnipotencia de Deus. --- Por isso a proposição \emph{contradictoria} é só \emph{physicamente}. \emph{impossivel} \footnote{O objecto, pois, da certeza \emph{metaphysica} é constituido por verdades \emph{metaphysicamente} necessarias, isto é, por verdades, cuja necessidade depende do nexo \emph{essencial} entre o sujeito, e o predicado. E, como a \emph{essencia} das coisns é \emph{immutavel}, a certeza \emph{metaphysica} não admitte excepções.}. -\footnote{O objecto da certeza \emph{pluysica} é constituido por verdades \emph{plujsicamente} necessarias, isto é. por verdades, cuja necessidade deriva das leis \emph{physicas}, que regulam e mundo. E. como as leis \emph{physicas} estão sujeitas ao poder do Creador, que as estabeleceu, a certeza \emph{physica} pode admittir excepções.}

o) \emph{Certeza} moral. --- Certeza \emph{moral} é a que se funda na \emph{constancia} das leis moraes, as quaes regulam as inclinações e as acções dos homens: tal é a certeza d'esta' proposição: os \emph{paes} amam. seus. filhos. --- Esta certeza tambem admitte excepções, porque as leis moraes podem ser violadas pela liberdade humana. --- Por i isso a proposição \emph{contradictoria} é só \emph{moralmente} impossivel (* \emph{y}.

c) Certeza directa e reflexa. --- É \emph{directa}, quando a intelligencia adhere firmemente a verdade, porque conhece, de um modo \emph{dmplicilo} e.pelo \emph{facto mesmo da} sua \emph{adhesão}, os \emph{motivos}, em que se funda a sua determinagão; assim todos conhecem, pelo menos \emph{implicitamente}, os motivos que os levam a acceitar, como inconeusso, este principio: \emph{todo o effeito exige} uma causa. --- É \emph{reflexa}, quando a intelligencia adhere firmemente a verdade, porque conhece, de.um modo \emph{explicito}, e pela \emph{reflexão}, os \emph{motivos}, que a:levam a essa adhes\&o ; assim, reflectindo sobre os. conceitos- \emph{de ejfeito-\&'de-causa}; -Conhecemos, \emph{explicitamente}, as razões; -que- demonstram- a- verdade d* aquelle principio ().\footnote{O objecto da certeza \emph{moral} é constituido por verdades \emph{moralmente} necessarias, isto é, por verdades, cuja necessidade depende das leis moraes, que dirigem as acções dos' homens para umifim constante e universal, determinado por Deus. E, como as leis moraes podem 'ser violidas pela liberdade humana, a certeza \emph{morali} tambem pode admittir excepções.}\footnote{A differenca entre a \emph{certeza directa}, e a \emph{reflexa}. não é \emph{essencial}, mas \emph{accidental}; porque a \emph{reflexa} é apenas a \emph{directa}, mas \emph{aperfeiçoada}. É uma differênça de \emph{graus}, e não de \emph{especie}.\par Alguns escriptores sustentam que não ha certeza \emph{directa}, mas que toda a certeza é essencialmente \emph{reflexa}. --- Mas não é assim. Quando a intelligencia, deterininada .pelo objecto, que se lhe manifesta convenientemente; proouncia um \emph{juízo}, está certa.do que affirma, pem é preciso, para isso, que seja consciente da propria certeza. Nem toda a affirmação é acompanhada necessariamente pela reflexão; e, quando é acompanhada, a affirmação não presuppõe a consciencia de si\textasciicircum{} mesma, mas, por ser \emph{objecto} da consciencia, deve \emph{preceder} a propria consciencia, em razão de. natureza, se não em razão de \emph{tempo}. --- Alem a isso, a \emph{consciencia}, Ou a \emph{reflexão}, Dão ac- . crescenta nada a firmeza da adbesão, ou a firmeza \emph{directa}, que é determinada '. exclusivamente pela evidencia da verdade. Como lambem nada acerescenta a certeza, que temos de uma coisa e que 'se funda na. evidencia, a consciencia da verdade da propria coisa; de modo que, se, por impossivel, a intelligencia não tivesse meio algum- para verificar o \emph{valor} objectivo dos seus conhecimentos, nem por isso os juizos certos, pronunciados por ella, perderiam nada da sua firmeza e necessidade. --- ''odavia, devemos admittir que a \emph{consciencia} da \emph{certeza}, como tambem do \emph{valor objectivo} de um juizo, dá a ultima perfeição aos nossos conhecimentos; porque então possuimos\pgn{125}plenamente a \emph{certeza}, como tambem a \emph{verdade}, quando temos a. consciencia. \emph{PP} de que a possuimos. --- Como o acto, por que a intelligencia alcança a \emph{certeza}, é identico ao acto por que alcança a verdade, é evidente que a mesma intelligeocia; quando" reflecte” sobre: .a- \emph{natureza}. do. seu acto; tem--consciencia não só da \emph{verdade}, mas tambem da \emph{certeza} do juizo. \textasciicircum{} ---\par Erram, pois, os que alfirmam que \& certeza directa, chamada tambemn \textasciicircum{} \emph{vulgar} e espontanea, não é verdadeira. certeca, mas \emph{cega credulidade}. Na certeza \emph{directa} encontra-se tudo o que constitue a certeza, tomada no seu sentido rigoroso: a verdade êvidentemente conhecida, a firme adbes\&o da intelligencia, excluinão todo o receio de errar, e a sufficiente percepção dos motivos. Será uma cega credulidade, e não uma verdadeira certeza, a cunviceão com que todos admittem a existencia propria, a existencia dos seus actos, dos seus parentes, da sua casa, e de mil outras coisas ? N\&o. Diz bem o \emph{Halmes}: « Interrogai a criança, o homem maduro, o velho, se estão certos da existencia propria, da existencia dos seus actos internos e externos, dos seus parentes, dos seus amigos, das suas habitações, e não achareis em ninguem a menor sombra de hesitac\&o. Crianças, velhos, homens, todos responderão a mesma coisa; e O seu olhar de espanto vos dirá o seu pen- . samento. Mas é possivel que um homem serio pergunte “ácerca. de coisas tão claras \emph{t}». (\emph{Phil}. \emph{fund}., V. 1, c. 3). --- Mas, dizem, hão ha cerfeca, onde - não ha conhecimento \emph{explicito} dos motivos. Não é assim. Se \textbf{assim} fosse, não haveria verdadeira certeza se não de um minimo numero de verdades: o queré falso. «Os proprios philosophos, diz Fénelon, estão invencivelmente convencidos de um grande numero de verdades, embora.não as possam explicar claramente, nem confutar os argumentos contrarios, que difficultam a sua percepção» (\emph{Lettres}). . :}

d) Certeza natural e sobrenatural. --- É natural, quando a intelligencia, esclarecida pela propria luz \emph{natural} e levada por um motivo \emph{objectivo}, isto é. baseado. nos proprios \emph{objectos}, adhere \emph{O} firmemente a verdade..--- É \emph{sobrenatural}, quando a intelligencia,. uM \emph{EN} esclarecida pela luz \emph{sobrenatural} da Fé é levadà- pela \emph{axctori}- \emph{I dade de} Deus, adhere firmemente a verdade (\emph{O}. \emph{s}

128. A certeza existe. UU \textasciicircum{}

a) A \emph{certeza} é a firme adhesão da intelligencia. á verdade. --- conhecida. Ora, como attesta a consciencia, á nossa intelhgencia adhere com tal firmeza a muitas verdades de evidencia imme- \emph{L}. diata, como são o \emph{facto} da. \emph{existencia propria}, os \emph{principios da Ps}

\textbar{} =\footnote{Advirta-se a distincção entre o natura! e O sobrenatural. \emph{Natural} ---. é tudo o que convem a um ente em virtude da \emph{natureza}, de que é dotado, ou que deriva e depende dos principios intrinsecos, de que a mesma natureza é composta: assim é natural para O homem o conhecimento de Deus pela consideração das creaturas. --- \emph{Sobrenatural} é tudo o que excede as faculdades e as exigencias naturaes da creatura; assim é sobrenatural para o homem o conhecimento \emph{intuitivo} de Deus.} \pg{126}  \emph{razão}, ete., que nem pode sequer suspeitar da possibilidade do contrario. Logo a \emph{certeza} existe \footnote{\emph{Cicero} expõe o mesmo argumento com as seguintes palavras: «Ut enim necesse est lancem in libra ponderibus impositis deprimi, sic animum perspieuis cedere. Nam quomodo non potest animal ullum non appetere id quod: accommodatum ad naturam áppareat, sie non potest animus objectam rem perspicoam non.approbare» (4cad.,l. 11, e. 18). --- Accrescente-se que, mesmo por confissão dos adversarios, a \emph{sciencia} existe. Ora a sciencia não pode existir nem conceber-se, se não se;admitlem alguns principios certos “e por si evidentes ; porque toda a demonstração, que produz a sciencia, deve partir de principios, dotados'd'essas qualidades. (De Ver., q. 11, a. 1y..--- O que dizemos da sciencia, deve applicar-se ás \emph{artes}, \emph{ás virbwdes}, a sociedade, e a muitas outras coisas, que não poderiani subsistir, se não se admittissem, como \emph{certas}, muitas verdades, relativas á nossa vida individual e social, : (8) A \emph{existencia da certeza} não \emph{pode} ser rigorosamente demonstrada, Por quanto, quem, procura demonusirar uma coisa, suppõe como \emph{certo} que esta raciocinando e que o raciocinio tem a força de convencer. e por isso suppoe Já a \emph{existencia da certeza}, antes de a demonstrar. Por isso a nossa demonstração é apenas uma explicação reflexa. --- \emph{Nem precisa} de ser de- 3uonslrada. A \emph{existencia} da \emph{certeza é um} facto tão evidente, que se impõe a nossa intelligencia; o que se vé, não precisa de ser demonstrado. Nem}.

. b) Toda a faculdade, assim como tende necessariamente para o seu-objecto-proprio, assim tambem, depois de o ter alcançado, necessariamente nelle se repousa. Ora a nossa intelligencia tende naturalmente para a verdade, e, quando esta lhe apparece claramente, não pode deixar de adherir a ella por necessidade de natureza, é de achar nella o seu descanço, sem nenhum receio de errar; isto é, não pode deixar de alcançar a \emph{certeza}. Logo a \emph{certeza} existe \footnote{Quem pode duvidar do facto da \emph{existencia propria}? Diz S. Thomaz: «Nullus unquam erravit ii hoc, quod jon perciperet se vivere» (\emph{De} Ver, q. 10, a. 8 ad 2) Quem pode negar os \emph{principios da vazao}, como o principio de \emph{contradicção}, de \emph{causalidade}, etc., sem renunciar a propria razão ? O mesmo Santo, fallando: d' esses principios, diz: «Verissima esse constat, intantum ut falsa esse nec possibile sit cogitare» (\emph{Post}. 4\emph{naly}., l. 19).}. \textbar{}

c) Deve admittir-se necessariamente a existencia de uma coisa, quando essa existencia se affirma pelo facto mesmo de se negar. Ora a existencia da \emph{certeza} affirma-se pelo facto mesmo de se negar; porquanto, quem nega a existerreia da certeza, emianto a nega emquanto tem a \emph{certeza} de que a certeza não existe. Logo a \emph{certeza} existe \footnote{\pgn{127}pode haver receio de errar; porque, como diz S. Thomaz, «sieut circa ipsnm introitum domus, qui omnibus patet, et primo occurrit, nullus decipitur. ita etiam in consideratione: veritatis » (In 11. \emph{Met}, 1. 1). . o\par a) Dizem: A intelligencia não adquire a \emph{certeza}, se não for dotada de aptidão para isso. Ora nunca poderá demonstrar-se que a inteligencia tem essa aptidão; porque commetter-se-hia um \emph{circulo vicioso}, dando-se como provado \emph{o} que havia de se provar; pois a demonstração suppõe e exige a \emph{certeca} dos principios, de que parte.\par Resposta. \emph{A} objecção contém o sophisma da passagem da ordem \emph{real} para a ordem \emph{logica}. É certo que, na ordem. real, ou \emph{ontologica}, primeira- . mente.é a \emph{aptidão} e-depois o acto das: nossas faculdades; mas é. falso que na. ordem \emph{logica}, Ou \emph{cognuscitiva}, O \emph{conhecimento} d'essa aptidão deva preceder * o \emph{acto}. Assim, eu não posso ouvir uma bella melodia; "se' não tiver-aptidão ---------- para isso: mas não é preciso que essa aptidão seja por mim conhecida antes do \emph{acto}, pois conhece-se e verifica-se no proprio acto, visto que da/ \emph{existencia} do acto podemos deduzir a \emph{aptidão} para o produzir: Por isso, a \emph{aptidão}, de. ... que a nossa intelligencia é dotada para alcahgar a certeza, deve, na ordem real, preceder os proprios conhecimentos \emph{certos}, mas não é preciso ' que os preceda na ordem \emph{logica}, isto é, não é preciso que' antes conhecamos essa \emph{aptidão} e depois alcancemos a \emph{certeza}.\par 5) Continuam: É impossivel averiguar se entre a ordem real e a \emph{ideal} existe uma porte, um nexo. Ora sem a existencia d'esta ponte, d'este nexo, nunca poderemos estar certos de que uma ordem corresponde á outra. Logo devemos admittir não a \emph{certrsa}, mas a \emph{duvida} universal.\par Resposta. A. ponte, o Texo, que une a ordem \$\emph{deal} com a \emph{real}, é a luz do objecto, que illumina a nossa intelligencia. As duas ordens, como vimos e veremos, não estão separadas; a ordem \emph{ideal} é a propria ordem real, emquanto existe intencionalmente no nosso espirito.}. \textbar{} ; 199. A certeza é proporcionada á intelligencja e ao motivo. ------- Todo o effeito deve ser proporcionado á causa, de que de-

. t * \pg{127}  riva. Ora a \emph{certeza} é um effeito, que deriva, como de causa \emph{adequada}, da \emph{intelligencia} e do \emph{motivo}; porque é a intelligencia; que levada pelo motivo adhere a verdade, e por isso a sua fir- .meza é mais ou menos inteusa, conforme a maior ou menor perfeição da sua luz e a maior ou menor efficacia do motivo. . Logo a \emph{certeza} é proporcionada á \emph{intelligencia} e ao \emph{motivo} \footnote{Na certeza devemos considerar duas cojsas: O \emph{acto} da intelligencia, que é a adhesão a verdade conhecida, e a \emph{firmeza} d'esse acto. Ora tanto o \emph{acto} como a sua \emph{firmeza} são coisas proporcionadas a perfeição da intellivencia e á efficacia do motivo. Ássim uma intelligencia mais elevada, mais educadà pelo estudo e pela reflex\&o, comprebende mais distincta e intimamente a força de um motivo, do que uma intelligencia menos perfeita, e por isso a sua firmeza na adhesão a uma verdade será mais intensa. Como tambem, na :mesma inteligencia, a adbesão a verdade será mais ou menos \emph{firme}, conforme mais ou menos forte iór o motivo. Assim uma pessoa acredita mais 8o que ella viu, do que no que os outros ihe referem. Por isso diz S. Thomaz: «Tanto major est certitudo, quanto fortius est motivum, quod determinationem causat» (In HL \emph{Sent}, Dist. 23. q. 2. a. 2\$). ---\pgn{128}Todavia, a \emph{certeza}, se admitte varios, graus na sua parle \emph{positiva}, emquanto significa q \emph{firme adhesão} da \emph{intelligencia} \& \emph{verdade}, não os admitte na sua. parte \emph{negativa}, emquanto significa \emph{a exclusão de} 4\emph{odo} o \emph{receiu de errar}. Porquanto, essa exclusão é6.um \emph{todo indivisível}, porque é uma \emph{negação perfeita}, não admitte \emph{graus}, não podendo haver um meio termo entre o \emph{ser} e 0 \emph{nàao-ser}.}.

£

: 130. --- AC céiteza metaphysica, a physica e a moral \textbf{diffe}-

rem especificamente, e uma especie é superior á outra = a) \emph{Differem especificamente}. --- A \emph{certeza} é proporcionada

tambem- \emph{ao-motivo}; e-por isso-a-motivos \emph{especificamente} distinctos “correspondem certezas \emph{especificamente} distinctas. Ora os motivos da certeza \emph{metaphysica}, da, \emph{physica} e da \emph{moral} são \emph{especificamente} distinctos; porque o motivo da certeza \emph{metaphysica} é a . essencia \emph{das} coisas, o da \emph{physica} é a \emph{constancia das leis physicas}, e o da \emph{moral} é a \emph{constancia das leis moraes}. Logo a certeza \emph{metaphysica}, a \emph{physica} e a moral differem \emph{especificamente} \footnote{Sendo a \emph{certeza} a, adhesão lirme da intelligencia á verdade, isto é, \& uma das proposições contradictorias, sem receio de que se verifique a opposta, conforme a diversa necessidade da-conveniencia do predicado com o sujeito, deve. haver diversas especies de certeza, Ora um predicado pode convir ao-sujeito --- ou porque o exige a essencia das coisas, assim dizemos --- que o homem é animal \emph{racional}, --- ou porque o exigem as \emph{leis physicas}, assim dizemos que o \emph{fogo} queima, --- ou porque o exigem as \emph{leis moraes}, assim dizemos que as \emph{mães} amam seus \emph{filhos}. D'ahi as tres \emph{especies} de certeza : \emph{metaphysica}, \emph{physica}. e moral: --- São tres especies, e não tres \emph{graus} da mesma especie. Por quanto; os. actos. distinguem-se pelo seu objecto \emph{formal}; e por isso a objectos \emph{formaes especificamente} distinctos correspondem actos \emph{especificamente} distinctos. Ora, sendo os »\emph{nofivos}, que constituem os objectos \emph{formaes} da certeza \emph{metaphysica}, \emph{physica} e \emph{moral}, \emph{especificamente} distinctos entre si, tambem \emph{especificamente} distinctas entre si devem ser a certeza \emph{metaphysica}, \emph{á physica} e a moral.}.

b) \emph{Uma especie}; \emph{é superior} d \emph{outra}. --- Sendo a \emph{certeza}, proporcionada ao motivo, uma especie de certeza será superior á

outra, se o motivo d'aquella fôr superior ao d'esta. .Ora --- o motivo da certeza \emph{melaphasica}., que é a \emph{essencia das coisas}, é superior ao da certeza \emph{physica}, que é é a \emph{constancia das leis physicas}, eo motivo da certeza: \emph{physica}: é superior ao da certeza. \emph{moral}, que é a \emph{constancia das leis}: \emph{moraes}: Legn a certeza. \emph{etaphagsica} é superior á \emph{physica}, ea certeza \emph{physica} é superior á \emph{moral} (E).\footnote{As tres especies de certeza diflerem entre si, --- não na parte \emph{negativa}, porque todas tres excluem o receio do opposto, --- mas na parte \emph{positiva}, isto é, emquanto a \emph{firmeca da adhesão}. Esta firmeza é proporcionada ao \emph{motivo}; e por isso, como o motivo é diverso em cada uma d'ellas, a firmeza numa especie de certeza é mais intensa que a firmeza numa outra\pgn{129}especie, e assim uma-certeza é superior á outra, --- A firmeza na certeza. \emph{metaphysica} é mais intensa que a firmeza nas outras duas especies ; visto que, na certeza \emph{metaphysica}, O predicado convem \emph{necessariamente} 30 \emph{sujeito}, pois convem pela exigencia da propria essencia, e esta é \emph{necessaria} e \emph{immutavel}, de modo que não é possivel a excepção. --- A firmeza na certeza \emph{physica} é mais intensa que a na certeza \emph{moral}; porque, na certeza \emph{physica}, \& necessidade da conveniencia do predicado com o sujeito funda-se na \emph{constancia das leis physicas}, que não podem ser suspensas senão pela omnipotencia de Deus, ao passo que, na certeza \emph{moral}, a necessidade da conveniencia do predicado com o sujeito funda-se na \emph{constancia das leis} moraes, que podem ser violadas pelo livre arbitrio do homem.\par Mas, embora as leis \emph{physicas} possam ser suspensas pela omnipotencia de Deus, e as leis \emph{moraes} possam ser violadas pelo livre arbitrio do homem, todavia d'ahi não se segue que a certeza \emph{physica} e a \emph{moral} não mereçam . o nome e a dignidade de verdadeira \emph{cer(eca}, é que se reduzam a mera \emph{pro}-\emph{babilidade}. Com effeito, a \emph{certeza}. é.a. firme. adhesão-da intelligencia-a-uma. verdade, com a exclusão de todo o \emph{receio} de errar. É claro que, quando se falla do \emph{receio} de errar, entende-se fallar do receio \emph{prudente}, que é proprio das pessoas serias e que se funda em motivos graves, e não do receio \emph{imprudente}, que é proprio dos espiritos levianos e que se funda nos delirios de uma imaginação doente. Ora a certeza \emph{physica} e a \emph{moral} excluem todo o receio \emph{prudente} de errar. Por quanto, a certeza \emph{physica} funda-se nas leis \emph{physicas}, e a \emph{moral} nas leis moraes. Ora as leis \emph{physicas} São constantes, e só podem ser suspensas por Deus, que só as suspende por motivos graves e em circumstancias taes, que a suspensão pode ser facilmente averiguada. Logo uma proposição, que se funda mas leis \emph{physicas}, é certa, e merece toda a adhes\&o da intelligencia, emquanto não \emph{se provar claramente} que se verificou o contrario, isto é, que se deu a excepção; porque á excepção não se suppõe, mas deve provar-se, e, até que se não prove, não deve ser admittida. Assim, quando digo: este fogo queima, a proposição é certissima, e.só um louco. pode affirmar o contrario. --- Como tambem \emph{constantes} são as leis moraes, porque representam a natura! ordenação da natureza humana para o bem: e por isso, se num caso particalar são violadas pela livre vontade humana, tal violação se manifesta immediatamente pelas circumstaxicias. Logo uma proposição, que se funda nas leis \emph{moraes}, é certa e merece a firme adhes\&o da intelligencia, até que se não provar com evidencia que, num caso particular, essas leis foram violadas; pois ninguem deve ser tido por mau, senão depois de se provar a sua maldade. Assim, quando digo: \emph{esta} mãe ama seus \emph{filhos}, a proposição é certa, até que não se prova o contrario. --- Quando, pois, se trata da certeza \emph{plysica} e da \emph{moral}, embora não repugne o contrario, todavia, se não ha motivos verdadeiramente \emph{graves}, O contrario não deve admittir-se, e a adhesão deve ser e é firme, porque exclue todo o receio \emph{prudente} de errar. --- Concedemos que, não repugnando 9 contrario, seja \emph{possivel} o erro. Mas, nesse caso, O erro não se segue \emph{per}\par Yol. I. ' 10\pgn{130}se, mas. \emph{per} accidens, e.a certeza \emph{physica} e a moral não deixam de ter o nome e a realidade de \emph{certeza}; porque uma coisa deve ser avaliada pelo que é \emph{per se}, isto é, pela sua natureza, e não pelo que é \emph{per accidens}, isto é, pelo que lhe pode acontecer por circunstancias extranhas. . :\par Muitos escriptores não reconhecem senão a certeza \emph{metaphysica}, porque, dizem, a certeza, a verdadeira certeza deve.excluir todo o receio, tanto \emph{prudente} como \emph{imprudente}, e a possibilidade de todo o erro, tanto \emph{per} se como \emph{per accidens}. --- Não dividimos essa opinião. Todos os homens estão \emph{certos}, firmemente certos de muitas verdades da ordem \emph{physica} € \emph{moral}, não suspeitam do opposto, e, em conformidade com essa convicção, regulam a sua vida. Assim, todos estão convencidos --- de que o \emph{fogo}. queima, e não se lançam numa fornalha accesa, embora seja, por um milagre, possivel o contrario, --- e de que o abraço da mãe se pode acceitar, sem receio de que seja um attentado contra a propria. vida, embora isso não seja absolutamente impossivel, Se, pois, os homens descançam tatobem nas verdades da ordem: \emph{physica e.moral}, estas verdades são tidas como certas. Será uma . certeza menos perfeita que a \emph{metaphysica}; mas é sempre \emph{certeza}.\par Objectam: "A certeza-physica e-u-moral não existem. Porquanto, a certeza. exclue. iloda/a \emph{duvida}. Ora a certeza \emph{physica} e a \emph{moral} não excluem toda a \emph{duvida} ; porque, estando as leis \emph{physicas} e as moraes sujeitas a excepções, podemos sempre suspeitar que se tenha dado a excepção.\par \emph{A} resposta está dada; mas a repetição não será inutil. É falso que a certeza \emph{physica} e a \emph{moral}! não excluem toda a duvida, porque é falso que, em cada um dos casos particulares, podemos suspeitar que se tenha dado a excepção nas leis \emph{physicas} e moraes. Na verdade, uma excepção nas leis \emph{physicas} só pode provir de Deus; e, como Deus não opera caprichosamente, se não Houver uma razão seria para fundamentar a excepção a essas leis num determinado caso, toda suspeita acerca d' essa excepção é infundada. Como tambem, uma excepção ás leis \emph{morars} deriva do abuso da liberdade: por isso, se não houver uma razão que demonstre esse abuso, toda o suspeita é irracional. --- Embora, pois, uma excepção seja \emph{possivel} em abstracto, todavia podemos saber com a certeza que, num caso \emph{concreto}, não se realizou, e.isto basta para que o espirito adhira firmemente a verdade, sem receio de errar. --- Repetimos que a \emph{possibilidade}' da excepção, se prova-que a certeza \emph{physica} e a \emph{moral} não são tão perfeitas como a \emph{metaphysica}, não tira aquellas duas especies a \emph{essencia} da \emph{certeza}.\par Agora: uma advertencia. A certeza \emph{physica} e a \emph{moral} não podem exprimir-se pelas seguintes formulas: --- se Deus não suspende as leis physicas, este fogo queima: --- se as leis moraes não são violadas, esta mãe ama seus filhos. A razão é clara. Sendo intrinseco e essencial o nexo entre a condição e o condicionado, essas formulas não exprimem a certeza physica e a moral, mas a metapbysica. --- Essas duas especies de certeza, \emph{physica} e moral, devem exprimir-se pelas seguintes formulas: este fogo queima, esta máe ama seus filhos. Por quanto, suppõe-se nestes casos, e com razão, que as leis \emph{physica} € \emph{moral} seguem o seu curso ordinario, não haveudo razões graves para sc admittir o contrario.\par a}

me

i29

\emph{M}

131. A certeza sobrenatural é superior a \textbf{toda} a certeza natural --- A certeza é proporcionada á perfeição da intelli-

\emph{A}.

\emph{z}

. . . \emph{ME} gencia e a efficacia do motivo. Ora --- a \emph{perfeição}, que a intel- \pg{130}  : ligencia recebe da "luz sobrenatural da Fé, é immensamente superior á perfeição, que ella recebe da sua luz natural, ------ e: o \pg{131}  \emph{motivo}, em que se funda a certeza sobrenatural e que é a \emph{auctoridade} divina, sendo \emph{infallivel}, é immensamente mais firme que todo o motivo, em que se funda a certeza natural. Logo a certeza \emph{sobrenatural} é superior.a toda a. certeza. \emph{natural} \footnote{A proposição é apenas uma consequencia dos principios expostos nos numeros antecedentes. --- Devemos, porém, fazer uma importante distineção. Se por \emph{certeza} se entende a. \emph{firmeza}. da adhes\&o a verdade, então a intelligencia adhere muito mais firmemente ás:verdades- da Fé" (divina) do: : que ás verdades da sciencia ;. porque o motivo, em.que.se. funda a. Fé, sendo. infallivel, é mais forte que o motivo, em que se funda a sciencia. Mas, se por \emph{certeza} se entende o descanço \emph{da intellügencim} na verdade conhecida, então, embora a adhesão ás verdades da Fé seja \emph{mais firme} do que a adhe- . são ás verdades da sciencia, comtudo o descanço da intelligencia é \emph{maior} nas verdades da razão do que nas da Fé; porque a adhesão ás .verdades da razão é determinada pela verdade das coisas, que se manifesta evidentemente a propria intelligencia, ao passo que a adhesão ás verdades da Fé é determinada pela vontade, que move e inclina a intelligencia a adherir 2 uma verdade. que se não vé, isto é, em que se não vé o nexo entre o predicado e o sujeito. D' aqui se segue que, se nas verdades da sciencia, evidentemente conhecidas, a adbesão da intelligencia é necessaria e a duvida é impossivel, nas verdades da Fé a adhesão é livre, e um movimento \emph{indeliberado de duvida} pode levantar-se no espirito do crente (Cf. \emph{De} Ver., q. 10, a. 12 ad 6). --- Lembramos que a \emph{firmeza} da adhesão não se oppóe o movimento \emph{indeliberado de duvida}, mas só a \emph{duvida deliberada}.}....

ArTIGO [II

4 \emph{existencia da certeza e o Scepticismo}

%%>>> t132

\pg{120} \pg{121} \pg{122} \pg{124} \pg{125} \pg{128} \pg{129} \tese{132}{Scepticismo.}
\emph{Scepticismo} é o. systema, que, negando. a legitimidade dos \emph{meios} de adquirir a certeza, nega a \emph{existencia} da mesma certeza, e considera .a \emph{duvida} como o resultado definitivo da sciencia e como o estado normal do espirito humano \footnote{É muito antiga a divisão dos philosophos em \emph{dogmaticos} e \emph{scepticos}. Os \emph{dogmaticos} sustentavam que o homem pode alcançar um conhecimento \emph{certo} de \emph{todas} as coisas; os \emph{scepticos}, pelo contrario. ensinavam que o homem não pode alcançar a certeza de coisa alguma. Ambos erravam, como diz S. Thomaz (L. 11 \emph{Post}., l. 7. Mas, se o \emph{erro} dos \emph{dogmaticos} era grave, o dos \emph{scepéicos} era gravissimo. --- Fundador do \emph{Scepticismo} foi Pyrro de Elide (360-970 a. Cb.), da eschola de Aristoteles e de Epicuro, seguido por Timão de Philunte (325-935 a. \emph{Ch}.), por Sexto Empirico (200 a. Ch), por Enesidemo de Creta (\emph{fl}. SO-60 a. Ch.), e por outros, que foram chamados pyrronicos. Diziam elles --- que só vemos as apparencias das coisas, e não a realidade, --- que toda a proposição pode ser defendida e oppugnada por argumentos de egual peso e valor, --- que não podemos affirnar nem negar a possibilidade de \emph{alcançar} a certeza, --- e que, por isso, é preciso \emph{suspender}\pgn{132}todo o juizo e contentar-se com a duvida. Esta suspensão foi chamada \emph{o}5servação, e os pyrronicos foram denominados \emph{observadores}. --- Os \emph{pyrronicos} foram seguidos e excedidos. pelos \emph{academicos}, sequazes de Arcesilau de Eolia (315 a. Ch.). Estes negavam expressámente a possibilidade de alcançar a certeza e por isso ensinavam que, em todas as coisas, é necessario suspender o assenso \emph{certo} e contentar-se com o assenso provavel. Como se vé, a differenga entre os \emph{pyrronicos} e os academicos não é sensivel; porque, no fundo, todos negam a possibilidade de adquirir a certeza. A opinião de uns e de outros teve \emph{o} nome commum de \emph{Scepticismo}, significando-se com essa ? palavra «o systema dos que affirmam --- que nada pode ser conhecido com \emph{certeza} e que por isso \emph{se} deve \emph{duvidar de tudo} ». --- Tambem este erro, apezar de tão grosseiro, foi résuscitado, de um modo total ou parcial, nestes ultimos seculos, pelos que não acceilam a luz do Evangelho e desprezam a sapiencia da Escholastica. São conhecidos os nomes de Miguel de Montaigne (1533-1592), de Pedro Gharron (1041-1603), de Francisco Sanches (1562-1632), de Pedro Bayle (1647-1706). de David Hume (1711-1776), e do celebre Manuel Kant (1724-1804) com a turba toda dos subjectivistas, modernistas, etc.}.

%%>>> t133

\pg{132} \tese{133}{--- Divisão do Scepticismo.}
O \emph{Scepticismo} divide-se em \emph{universal} e \emph{particular}. --- \emph{Universal} é o que, negando a legitimidade de \emph{todos} os meios de adquirir a certeza, sustenta que : nada é \emph{certo} e que por isso \emph{se deve duvidar de tudo}. --- \emph{Particular} - --- é o que, negando a legitimidade de um ou \emph{outro} meio de adquirir a certeza, diz que se deve:duvidar de uma ou \emph{outra} especie de verdades, relativa a esse meio. --- Reservando para mais tarde a critica do Scepticismo \emph{particular}, aqui tratamos do universal \footnote{O scepticismo \emph{particular} tuma diversos nomes, conforme a diversidade dos \emph{meios}, cuja legitimidade se contesta. Chama-se --- a) \emph{positivismo} ou \emph{empirismo}, se nega a legitimidade da \emph{intfelligencia}, e” só acceita o que pode ser percebido pela \emph{experiencia} externa, --- b) idealismo, se não admitte o testemunho dos \emph{sentidos}, --- \emph{c}) \emph{racionalismo}, se nega a auctoridade da Revelação \emph{sobrenatural}, --- \emph{d}) \emph{fideismo}, se contesta o valor da \emph{ração} e só á admitte como \emph{certo} o que foi revelado, -- \emph{e}) \emph{scepticismo historico}, se não E acceita o testemunho dos homens.}.

%%>>> t134

\tese{134}{O Scepticismo universal é naturalmente absurdo.}

a) Uma opinião é \emph{naturalmente} absurda, 'quando repugna á natureza. Ora o \emph{Scepticismo} «\emph{universal} repugna á natureza; pois esta; pelo seu proprio peso, determina a intelligencia a adherir firmemente a verdades, tanto da ordem \emph{ideal} como da \emph{real}, dotadas de evidencia immediata, de modo que toda a du- "vida se torna absolutamente impossivel. Logo o \emph{Scepticismo universal é naturalmente} absurdo \footnote{Com razão diz \emph{Pascal}: «Duvidará o homem de tudo? Duvidará de que duvida? Duvidará de que existe É impossivel que se chegue a estes\par a \textasciicircum{} »}: m

b) Uma opinião é \emph{naturalmente} absurda, quando não pode 2 \emph{m} ser traduzida na pratica, embora se refira á pratica. Ora tal é " » o \emph{Scepticismo universal}; porque ninguem pode operar, se não (5 -está \emph{cerío} --- de que existe, --- . de que é. dotado.de faculdades... :": . operativas, --- de que existem, na realidade, algumas coisas, em \textbar{} = que possa exercer a sua actividade, etc. Logo o \emph{Scepticismo uni}- : \textbar{}

\emph{UM versal é naturalmente} absurdo \footnote{S. Agostinho deduz de existencia da \emph{duvida} a existencia da \emph{certeza} com as seguintes palavras: « Regulam ipsam, quam vides, concipe hoc modo: omnis, qui se dubitantem intelligit, verum \emph{intelligit}, et de hac re, quam intelligit, \emph{certus} est; de vero igitur \emph{ceriws} est. Omuis igitur, qui utrum sit veritas dubitat, in se ipso babet verum, unde non \emph{dubitet} ; nec ullum verum \emph{S} nisi ip verilate. verum est. Non itaque oportet eum de veritate dubitare, qui \textbar{} potuit uundecumqe dubitare» (\emph{De} vera \emph{religione}, \emph{c}. 30, n. 75). E \emph{La}\par PS\par do.}. \&

. 135. O Scepticismo universal é theoricamente absurdo. « \textasciicircum{}

a) Toda a opinião, que envolve contradicções, é \emph{íheorica}- \textasciicircum{} = \emph{mente} absurda. Ora o \emph{Scepticismo universal} envolve muitas con- \textasciicircum{} iradiecgóes. Com effeito, os scepticos, quando proclamam este' . principio: \emph{deve duvidar-se de tudo}, ou 0 reconhecem como ver- \textasciicircum{} \emph{dadeiro}, ou como \emph{falso}, ou como \emph{duvidoso}. Se o reconhecem como : \emph{verdadeiro}, então é \emph{certo} para elles que se deve duvidar de tudo; sec reconhecem como \emph{falso}, tambem é \emph{certo} para elles que não "se deve duvidar de tudo ;"se'o reconhecem como \emph{duvidoso}, tambem neste caso é \emph{certo} para-elles que: duvidam d'esse principio: --- Logo o \emph{Sceplicismo} universal é \emph{theoricamente} absurdo (2).

b) Uma opinião é \emph{fheoricamente} absurda, quando os seus \textbar{}. \emph{S}\textasciicircum{} defensores affirmam com os factos o que negam com.as palavras. - É o que acontece com os defensores do Scepticismo \emph{universal}. De facto, os scepticos, na defeza da sua opinião, empregam o

extremos; e penso que nunca houve effectivamente um pyrronico perfeito. --- = A propria natureza sustenta a razão impotente e não a deixa delirar até 7 este ponto» (\emph{Pensées}, ch. IV, par. 95). : \emph{ME}

(1) S. Thomaz nota --- esta impossibilidade de operar. em que se en- » U contra um sceptico, se quizer ser coherente com a sua opinião, --- e a contra- ; diec\&o, em que pode eahir e cahe realmente na vida pratica (\emph{Metaph}. L. XI, 1.6) --- É curioso o seguinte facto da vida de \emph{Pyrro}. Querendo conformar \emph{m} a Deatica \textbar{} com a theoria, prometteu que d'ahi por deante não evitaria o ; ---

-encontre: dos cavallos, dos carros, etc. Um dia, vendo ao longe um. cão. dam- 2 nado, que vinha ao seu encontro, \emph{Pyrro} deitou a fugir, e aos discipulos, \emph{m} que Ihe censuravam a incoherencia, respondeu : «É difficil despir-se da natureza --- Difficile est naturam exuere!». ----raciocinio \pg{134} , afirmam que os \emph{sentidos} e a \emph{razão} se enganam, disünguem a \emph{verdade} da \emph{falsidade}, a \emph{duvida} da \emph{certeza}, a \emph{affirmação} da \emph{negação}. Sendo assim, elles admittem, como: \emph{certo}, que são dotados de faculdades perceptivas, --- que existem e

duvidam, --- que alguma: coisa existe na realidade, --- que O principio de contradicção é verdadeiro, etc. Logo o \emph{Scepticismo} universal é \emph{theoricamente} absurdo (i) -

%%>>> t136

\pg{133} \tese{136}{Consequencias do Scepticismo universal.}
As conse-

queneias d'esta opinião não poderiam ser mais perniciosas. A negação da \emph{certeza} é a destruição da ordem”. \emph{religiosa}, \emph{moral} e \emph{social}: --- É a destruição da ordem \emph{religiosa}; porque não admitte ' a auctoridade divina, que é certissima; nem as verdades reveladas, que são infalliveis. --- É a destruição da ordem moral; porque não é possivel a honestidade da vida, quando se duvida dos primeiros principios dà moralidade e dos primeiras preceitos dà lei natural. ="É a destruição da ordem \emph{social}; porque, quando se duvida do principio da auctoridade, dos direitos e deveres, etc., falta u sociedade, o verdadeiro e essencial fundamento \footnote{E. para limitar a esphera da demonstração, diremos que a. pretendida existencia da \emph{duvida universal} é a destruição do \emph{Scepticismo universal}. Por quanto, como sabem os scepticos que se deve duvidar de tudo e que realmente duvidam de tudo? Pelo testemunho da \emph{consciencia}. Logo elles têm por coisa \emph{certa} o testemunho da consciencia; e assim renegam o seu systema. --- Se disserem que duvidam. tambem do testemunho da consciencia, renegam' egualmente Oo seu systema; porque, senão é certo o gue a consciencia.refere, nem é: certo qué existe a duvida.\par Objectam : Devemos evitar \emph{todos} os erros. Ora só quem duvida de tudo, evita \emph{todos} os erros. Logo deve seguir-se o Scepticismo universal.\par Resposta. É verdade que devemos evitar, com todo o empenho, todos os erros: mas é falso que evila todos os erros quem duvida de tudo. A duvida universal é o maximo dos erros, é a fonte de todos os erros. --- O argumento dos scepticos é semelbante a este: quem nada come, está Seguro de não ser envenenado ; logo devemos rejeitar todo o alimento! --- É prudencia não adherir a coisas duvidosas, que não têm solido fundamento ; mas é pertinacia duvidar de coisas evidentes. --- Diremos com S. Agostinho: « Miserabilias enim \emph{fallitur}, qui nunquam putat facultatibus cognoscitivis esse credendum, quam qui illis eredendo nonnumquau fallitur » (\emph{De civ}. \emph{Dei}, 1. 19, c. 1S).}.\footnote{Estas consequencias são evidentemente legitimas. --- A quem objectasse que o \emph{Scepticismo} só se refere aos conhecimentos da ordem natural e por isso não pode destruir a vossa Fé, responderiamos que a nossa Fé não é possivel, se não se admitlem, como certas, muitas verdades da ordem\pgn{135}\emph{natural}, como são a existencia de Deus, a sua auctoridade, o facto da Revelação divina, etc., e que por isso a destruição da ordem \emph{natural} traz comsigo a destruição da ordem \emph{sobrenatural}.}

\emph{Universal e supremo motivo da certeza}

%%>>> t137

\pg{135} \tese{137}{Motivo da certeza.}
\emph{Motivo} da \emph{certeza} € a razão, quê determina a \emph{intelligencia} a \emph{adherir firmemente a verdade}. --- Este \emph{motivo} deve existir; porque a \emph{certeza} é um effeito, e todo-o effeito exige uma causa proporcionada \footnote{O \emph{motivo} da certeza deve existir necessariamente. A nossa intelligencia, sendo uma faculdade cognoscitiva, não pode adherir firmemente a . uma verdade sem uma razão, sem um motivo. Uma adhesão firme, determinada pelo iastincto \emph{cego} da natureza, repugna á indole e a nobreza da inlelligencia, --- O \emph{motivo} deve poder explicar sufficientemente o effeito, que elle produz e que é a propria certeza.}.

%%>>> t138

\tese{138}{Especies do motivo da certeza.}
O \emph{motivo} da \emph{certeza}.

é --- \emph{a}) \emph{inirinseco} e \emph{extrinseco}, --- \emph{b}) \emph{universal} e. \emph{particular}, --- €) \emph{supremo} e \emph{proximo}.

a) Motivo intrinseco e extrinseco. --- E' \emph{intrinseco}, quando a razão, que determina a intelligencia a adherir a verdade, reside na \emph{propria verdade}, ou no \emph{objecto}. --- E' \emph{extrinseco}, quando essa razão reside, --- não na verdade, ou uo objecto, --- mas na testemunha, que nos refere uma. verdade \emph{C}

b) Motivo universal e particular. é' commum a todas as verdades, ás quaes adhere a intelligencia. --- E \emph{particular}, quando é \emph{proprio} de uma só especie ou de uma só ordem de verdades \footnote{Que o \emph{motivo da certeza} possa ser e seja \emph{intrinseco} e \emph{extrinseco} á verdade, deduz-se da existencia de uma duplice ordem de verdades, --- que são as verdades da \emph{scienciam} e as verdades da fé. Nas verdades da sciencia, a intelligencia não adhere senão por um motivo \emph{intrinseco} a propria verdade, por uma razão, que se encontra no proprio objecto e que, conhecida pela mesma intelligencia, determina esta faculdade a adherir. Nas verdades da fé, a intelligencia adhere a verdade por um motivo, que é \emph{exírinseco} \& propria verdade e que reside na pessoa, que nol-a refere \emph{ou} ensina. \textasciicircum{} (3 O motivo universal e O \emph{particular} não se excluem mutuamente, mas um é compativel com o outro, como o superior é compativel com o inferior. --- Cada \emph{especie} de certeza tem o seu motivo particular. Todavia, esses motivos \emph{particulares} devem ter uma nota \emph{commum} : e esta nota, como veremos, consiste na evidencia \emph{objectiva}.\par *}. Ut \emph{tiom em}

c) Motivo supremo e proximo. --- E? supremo, quando não depende de outro, e é a razão \emph{ultima}, que nos determina a adherir a verdade. --- E' \emph{proximo}, quando depende de um motivo superior, e nelle finalmente se resolve (\textasciicircum{}).

%%>>> t139

\pg{136} \tese{139}{O motivo da certeza é a evidencia do objecto .}
\footnote{Todos os entes, como dissemos, têm varias e diversas causas, de que constam ou de que dependem, e que, por isso, são \emph{intrinsecas} ou \emph{exirinsecas}. --- Umas e outras são \emph{proximas} ou supremas (\emph{remotas}). --- As causas. \emph{proximas} referem-se ás \emph{supremas}, em que acham a sua base e a sua explicação. --- Como as causas, conhecidas por nós, são \emph{motivos} da nossa certeza, assim esses motivos distinguem-se em. \emph{proximos} e \emph{supremos}.\par \emph{t} Expliquemos a \emph{matwreza} e as \emph{especies} da \emph{evidi}\par a) Evidencia. --- \emph{Evidência} é a \emph{clara manifestação do objecto}, ou \emph{da verdade}. Por quanto, todos dizem que um objecto é \emph{evidente}, quando sé manifesta claramente a faculdade perceptiva, de modo que actua nella e a determina a percebel-o como é em si mesmo. «llla \emph{véderi} dicuntur, quae per se ipsa movent intellectum nostrum vel sensum ad sui cognitionem » (\emph{Sum}. \emph{Th}.. I-II, q. 1, a. 4). --- Um objecto, para que se torne \emph{evidente}, ou visivel a intelligencia, deve ser cercado por uma luz, ou claridade conveniente, como deve ser illuminado um objecto \emph{material}, para que possa ser percebido pelos nossos olhos. Essa luz, ou claridade, é coisa propria da intelligencia, e da intelligencia irradia para os objectos e os torna visiveis; porque a intelligencia, dotada como é de uma energia espiritual, abstrabe dos pbantasmas as imagens ou especies immateriaes, que representam os objectos, e em que e por, que elles: são claramente perceptiveis e percebidos. Portanto, o phenomeno completo da. \emph{evidencia}, como o da visão externa, abrange tres coisas: O0 \emph{objecto}, a sua \emph{visibilidade} e a sua \emph{manifestação} a faculdade perceptiva. O \emph{objecto} e a sua \emph{visibilidade} constituem o elemento \emph{material} da \emph{evidencia} (0 \emph{objecto} é o elemento \emph{remoto}, a Sua \emph{visibilidade} é o elemento proximo), á. \emph{manifestação} é o, elemento \emph{formal}.\par b) Especies de evidencia, --- A \emph{evidencia} pode ser --- \emph{a}) \emph{objectiva} e \emph{subjectiva}, --- \emph{b}) \emph{intrinseca} e extrinseca, --- c) \emph{immediata}, e \emph{mediata}.\par \emph{a}) \emph{Evidencia objectiva} e \emph{subjectiva}. --- A \emph{objectiva} é o \emph{objecto}, emquanto, iluminado pela intelligencia, a: esta se manifesta claramente. A \emph{subjectiva} é a clara \emph{percepção} do objecto, ou da verdade, que se torna presente á intelligencia. --- A \emph{evidencia} é uma coisa propriamente \emph{objectiva}, pois é uma qualidade do objecto cognoscivel, embora deva tornar-se de algum modo \emph{subjectiva}, porque, se o objecto não se manifesta claramente a inlelligencia, esta não pode dar a propria adhegão. --- \emph{A} evidencia \emph{objectivo} é de varias \emph{especies}. Por quanto, sendo uma \emph{qualidade} do \emph{objecto}, ou da. \emph{verdade}, essa. \emph{evidencia}\pgn{137}, € \emph{metaphysica}, \emph{physica}, e \emph{moral}, conforme o \emph{objecto}, ou « verdade, se funda ' na ordem \emph{metaphysica}, na \emph{physico}, ou na moral. --- A evidencia \emph{subjectiva}, m admitte varios graus na mesma \emph{especie}, conforme a maior ou menor perspicacia e applicac\&o da propria intelligencia. DA\par \emph{b}) \emph{Evidencia intrinseca} e \emph{extrinseca}. --- São duas especies da evidencia --- \emph{objectiva}, --- A \emph{intrinseca} é propria do \emph{motivo intrinseco} da \emph{certeza}, e reveste \emph{Au} a propria \emph{verdade}, ou o \emph{objecto}; tal é x evidencia das verdades da \emph{sciencia}. --- A \emph{exirinseca} é propria do \emph{wofivo extrinseco} da \emph{certeza}, e reveste a \emph{auciori}- \emph{AE dade}, que attesta uma verdade; tal é a evidencia das verdades da fé, tanto \emph{sobrenatural} como \emph{natural}. --- Por isso, quando. a evidencia é intrinseca, vó- Se-claramente o nexo entre o sujeito e o predicado; quando é exírinseca, Dão apparece esse nexo. mas apparece a \emph{auctoridade} de quem attesta uma . coisa. A evidencia \emph{intrinseca} torna a verdade \emph{evidente}: \emph{a extrinseca} torna "a: verdade \emph{evidentemente}. \emph{crivel}.. ---. Advirta-se- qe-a--evidencia- "extrinseca. és \emph{tal} com relação á \emph{verdade}, a qual pc camente ; porque, com relação\par \emph{c}) \emph{Evidencia immediata e mediata}. --- São duas especies da evidencia \emph{intrinseca}. --- E' \emph{immediata}, quando a verdade se manifesta --- pela simples \emph{analyse} dos termos, como nos \emph{principios da razão}, p. ex. --- \emph{o todo é maior que} uma \emph{parte}, --- ou pelo \emph{testenunho dos sentidos}, como nos \emph{factos experiment«es}, p. ex. --- \emph{o sol resplandece}. - E' \emph{mediata}, quando a verdade se manifesta --- não por si mesma, --- mas pela luz dos principios \emph{immediatamente} evidentes, em que está contida e de que é deduzida pelo raciocinio, como nas conclusões, p. ex. ---- a \emph{alma}. \emph{humana é immortal}. --- Mas, tanto num como no outro caso, a \emph{evidencia} manifesta a propria \emph{entidade} da verdade, ou do objecto.}.

a) O \emph{motivo} da \emph{certeza} é

a razão, que determina a intelli-

verdade. Ora a razão, que determina a ingencia a adherir a teligencia a adherir a verdade, é a \emph{evidencia do objecto}, isto 6, a \emph{clara manifestação do objecto}; porque toda a faculdade é determinada ao acto pelo seu objecto, e por isso, quando este objecto \pg{137}  “lhe se manifesta \emph{claramente}, ella não pode ser indifferente, mas, --- levada pela sua propria natureza, deve adberir com firmeza ao L mesmo objecto e descancar nelle. Logo o \emph{motivo} da \emph{certeza} É a \emph{E evidencia} do \emph{objecto} \footnote{A nossa intelligencia, sendo por si indifferente para a percepção de . innumeraveis objectos, não pode perceber um certo objecto, se este não Ihe ; : Se manifesta convenientemente e não a determina ao acto. Esta manifestação e. determinação são absolutamente necessarias. --- Mas, depois de o objecto se manifestar e determinar a iotelligencia; esta não pode permanecer indifferente, mas \emph{necessariamente} deve perceber o objecto e adherir a elle com firmeza. Exige-o a propria natureza da faculdade cognoscitiva. Se assim não. fosse, se a faculdade podesse permanecer indifferente perante a acção do objecto, deveriamos dizer que a faculdadenão foi feita para o seu objecto, isto é, que a faculdade não é taculdade: o que repugna.\par Mais tarde explicaremos o modo, por que o objecto actua na intelligencja € a determina a percepção. Limitamo-nos agora a dizer que não é possivel o conhecimento, se \emph{o} objecto cognoscivel não se une com a faculdade cognitiva\textasciicircum{} 0 objecto \emph{externo} penetra. DO Orgão animado do sujeito, --- não pela Sua \emph{substancia}, porque o corpo é impenetravel, --- mas pela sua \emph{acczo phuysica}, porque a acção do agente se encontra no paciente. Ora esta acção é a. \emph{expressão}, mais ou menos adequada, da \emph{natureza} do objecto: porque o effeito é \emph{Semelhante} á propria causa. O sujeito adverte e percebe, por intuição ou de\par H\pgn{138}\emph{i}. \textasciicircum{} 4 \emph{c}---\pgn{139}e o vicio attribuir-se-hia ao proprio Auctor da natureza. --- Concluimos que, assim como a vontade não pode deixar de adberir a. um bem, que não tem . apparencia alguma de mal, assim tambem a intelligencia não pode deixar de adberir a uma verdade, quando não tem apparencia alguma de talsidade, isto é, quando se torna visivel ou evidente (S. Thom., \emph{De malo}, q. 6, a. un., ad 10).\par 1) O objecto, ou a verdade, quando se torna. evidente, \emph{determina} por tal modo a intelligencia a produzir o acto, que este não pode ser impedido \emph{directamente} pela vontade, embora o possa ser \emph{indirectamente}. --- Uma faculdade impede \emph{directamente} uma outra, quando a primeira não deixa operar a segunda, embora esta por si possua \emph{todos} os requisitos necessarios para o acto, Uma faculdade impede-indirectamente-uma-ountra; quando: esta-0á0 possue um requisito necessario para o acto, e tal \emph{falta} depende da influencia da. primeira faculdade. --- \emph{A} proposição pode desdobrar-se nas duas seguintes:\par \emph{a}) A producc\&o do acto da intelligencia não pode ser impedida \emph{directamente} pela vontade. Por quanto, a consciencia attesta que a intelligencia vão pode deixar de adherir a verdade, quando esta lhe se torna evidente, seja qual fôr \emph{o} esforço da vontade para obstar a essa adhesão. E não podia ser d'outro modo. Quando a intelligencia adhere a uma verdade evidente, \emph{nenhum} dos elementos, que concorrem para este acto (e que são a intelligencia, o objecto, a evidencia do objecto, a união do objecto com a faculdade), está sujeito ao dominio \emph{directo} da vontade. Portanto a vontade, assim como não pode obstar a que os olhos vejam um objecto, quando este se acha presente e aquelles estão bem dispostos e convenientemente applicados, assim tambem não pode obstar a que a inteligencia adhira a uma ' verdade, quando esta lhe se torna evidente. eom mM\par 5) A producção do acto da. intelligencia pode ser impedida \emph{indirecta}\par \emph{mente} pela vontade. Dá-se este phenomeno, quando á intelligencia falta, pela influencia da vontade, um ou outro requisito necessario para o acto. Assim, se a. intelligencia, pela imposição da vontade, distrahe de uma ver--- dade a sua attenção, não pode dar a propria adhes\&o, embora a verdade seja por si evidente: assim como, se a vontade impõe aos olhos que se féchem, estes não poderão perceber um obiecto, embora presente e lumi- \textbar{} noso; porque, então, falla um requisito essencial para o acto, qual é a applicação da faculdade ao seu objecto. --- Notamos que, se a verdade \emph{não} é OU. \emph{não Gpparece evidente}, á intelligencia, não sendo determinada necessarirmente pelo “objecto. pode ser impedida, no exercicio do seu acto, pela vontade, como pode ser determinada por ella; o que se verifica com relação --- ás verdades \emph{provareis}, que não são \emph{perfeitamente} evidentes, --- ás verdades da fé, quer bumana. quer divina, que não são dotadas de evidencia \emph{intrinseca}. (\emph{Sun}. (\emph{Th}., III, q. 17, a. 6: \emph{De} Ver., q. 14, \emph{a}. 1).}. \emph{BM Án m}

--- = \pg{138}  \emph{i} b) A verdade, a qual a intelligencia \textbf{adhere} ou pode adherir, \emph{z tL} --- ou é \emph{percebida por nós}, ou é \emph{attestada} por outrem, Ora, em [1 ambas as hypotheses, o \emph{motivo} da certeza é a evidencia \emph{do objecto}. \emph{E} : . E um modo immediato, essa \emph{acção} extranha, que actua nelle e o penetra, e, \emph{RC} E percebendo a acção, percebe o proprio agente, o proprio objecto externo. \emph{v} Assim, o objecto se torna manifesto e visivel, pela sua acção, a faculdade \emph{he} perceptiva ; e o sujeito, por tal evidencia objectiva, está tão certo da exis- \emph{E} - tencia do objecto externo, como está certo da existencia propria. --- O objec- NE to sensivel tem: dois aspectos. Um aspecto é physico e exterior, e este cale

H sob a intuição dos sentidos ; outro é immaterial e interior, e este cahe sob --- a intuição da intelligencia. Este aspect material e interior dos objectos

sensiveis é representado por uma imagem espiritual, e é nesta e por esta

imagem que a intelligencia vé o \emph{objecto} ea elle adbere € com firmeza. ed Dizemos que a intelligencia deve necessariamente adherir ao objecto, ou

a verdade, quando esta lhe se torna visivel ou evidente. --- Mas. ha uma

duplice necessidade; de \emph{especificação} e de \emph{exercicio}, A necessidade é de \emph{espe}-

\emph{cificação}, quando: uma faculdade é por tal modo determinada a produzir um

corto acto, que não pode produzir um acto diverso, ou contrario; assim a \emph{Fi} vontade dos bemaventurados ama necessariamente Deus, e é impossivel que \emph{Hc} o odeie. A necessidade é de \emph{exercicio}, quando uma faculdade é por tal modo t determinada a produzir o acto, que não pode deixar de o produzir; assim \emph{DUM} os nossos olhos, quando são bem dispostos e convenientemente applicados, Pu não podem deixar de perceber um objecto presente. --- A esta duplice neces- \emph{AD} 'sidade corresponde, em sentido opposto, uma duplice \emph{indifferenga}, ou \emph{liber}-

\emph{dade}, que tambem-se chama de \emph{especificação} e de \emph{exercicio}. \emph{e} Feita esta disiinecio, pergunta-se: A necessidade, em que se acha a

» .intelligencia, de adherir a um objecto, ou a uma verdade, que se torna \emph{E} evidente, é uma necessidade de \emph{especificação}, ou tambem de \emph{exercicio}? A ' esta pergunta damos a seguinte resposta. 3--- . 2) O objecto, ou a verdade, quando se torna evidente, determina por Dm tal modo'a'adlies\&o da intelligencia, que esta não pode adherir a \emph{proposimd} ção contraria. É a necessidade emquanto á \emph{especificação} do acto. A conclusão - é certissima. Toda a faculdade, como dissemos, deve necessariamente adberir = ao bem, para o qual tende e que lhe se manifesta convenientemente. Ora o - bem da intelligencia é a verdade, e por isso, quando esta se torna evidente, a \textbar{} : intelligencia não pode deixar de adherir a ella, e de rejeitar a proposição I contraria. --- Além d' isso, se \emph{a} intelligencia podesse deixar de adherir a uma NE proposição evidentemente. verdadeira. não haveria repugnancia alguma em M que ella podesse adherir a uma proposição evidentemente falsa; porque, na \textbar{} hypothese, admittic-se- hia a liberdade ou indifferenga \emph{de especificação}, a. qual ].---7 diz respeito a actos oppostos, e a adhesão a uma proposição evidentemente \emph{NE} falsa é um acto opposto a adhesão a uma proposição evidentemente verda- \textbar{} deira. Ora é impossivel que a intelligencia adhira a uma proposição eviden- \emph{PN} temente falsa. Por quanto a \emph{falsidade} é um mal da intelligencia, e, quando \emph{b}. € \emph{evidente}, é um mal total. sem alguma apparencia de bem. Ora nenhuma \textbar{} \emph{u} faculdade pode adherir ao seu mal \emph{total}; aliás seria naturalmente \emph{viciada}, \emph{t} \textbar{}\emph{o} ] ; --- \pg{139}  ---- É \emph{a evidencia} do \emph{objecto}, quando a verdade, percebida por nós, é \emph{necessaria}; porque a intelligencia não une um predicado “aum sujeito, senão quando vé, de um modo \emph{immediato} ou me- "diato, \emph{que} um \emph{convem-.ao--outro}.---- Ea. evidencia. \emph{do} ..\emph{objecto},. -. quando a verdade, percebida por nós, é \emph{contingente}; porque não admittimos um facto, ou um objecto, senão porque \textbf{esse} \pg{140}  facto, ou esse objecto, se revelou \emph{claramente} ás respectivas fa-. culdades perceptivas. --- E' a \emph{evidencia} do \emph{objecto}, quando a verdade é \emph{attestada} por outrem; porque não se admitte \textbf{essa} verdade, senão..emguanto..se torna evidente pela \emph{evidencia da auctoridade}. Logo o \emph{motivo} da \emph{certeza} é a \emph{evidencia do objecto} \footnote{A proposição menor pode ser mais desenvolvida, --- A verdade \emph{mecessaria} pode ser --- um \emph{principio de raso}, como: todo o effeito exige uma causa, --- ou uma conclusão, como: a alma humana é immortal Se é um \emph{principio de ração}; x-intelligencia adhere firmemente a essa verdade, porque a conveniencia do. predicado..com..o.sujeito.lhe.se mostra com. ums evidencia \emph{immediata} e \emph{propria}. Se é uma \emph{conclusão}, a intelligencia não adhere a ella, se não depois de lhe constar uma duplice \emph{evidencia}: a evidencia das \textbar{} \emph{premissas}, nas quaes a conclusão se resolve, e a evidencia do \emph{nexo logico} entre as premissas e a conclusio. Por isSo é' sempre a evidencia, embora \emph{mediato}, que leva a intelligencia a adherir a uma conclusão. Diz S. Thomaz: « Assentit intellectus alicui duplieifer : uno modo, quia ad boc movetnr ab ipso \emph{objecto}, quod est --- vel per se cognitum, sicut patet in principiis primis, --- vel per \emph{aliud} cognitum, sicut patet de conclusionibus » (\emph{Swm}. Th., II-II, q. 1, a. 4). --- Se a verdade é \emph{contingente}, é tambem a \emph{evidencia}, \emph{do objecto} que nos determina a adhes\&o. Porque estamos nós certos de que o mundo existe e o sol resplandece? Porque esses objectos são evidentes para os nossos olhos. Diz S. Thomaz: « Similiter determinatur judicium sensitivae partis ex hoc, quod sensibile subjacet sensui »(In I \emph{Sent}., Dist. 93, q. 2, a: 9. --- Sea verdade é \emph{affestada} por outrem, é Sempre a \emph{evidencia}, ainda que extrinseca, que deteceming,a adbesão da inteligedcia, porque não acreditamos numa coisa, referida “Hor outrem, se não é \emph{evidente} a \emph{auctoridade} da testemunha. --- Como, porém, só a evidencia intrinseco determina \emph{necessariamente} a. intelligencia a adhesão, a mesma inteligencia, quando se trata de verdades de fé. . ou dotadas apenas dé evidencia \emph{extrinseca}, pode soffrer a influencia da vontade, e esta pode obstar a que a intelligencia pratique o acto, ou pode determinal-a a adherir a uma das duas proposições contradictorias. Continua o S. Doutor: « \emph{Alio} modo intellectas assentit alicui, non quia sufficienter moveatur \emph{ab objecto} proprio. sed per quamdam electionem voluntatis, declinaus in unam partem magis quam aliam... EL si hoc sit cum certitudine absque formidine alterins partis, eril. \emph{fides} » (Sum. \emph{TÀh}., 11-11, q 1, a. 4). --- Concluimos, com S. Thomaz, que a \emph{evidencia objectiva} é o fundamento, o verdadeiro motivo da nossa certeza. « \emph{Certitudo}, quae est in scientia et intellectu, \emph{est} ex \emph{ipsa evidentia} eorwm, quae \emph{certa} esse dicuntur » (In [ \emph{Sent}., 1. c.).}.

140. A adhesão, causada. \textbf{pela} evidencia do objecto, é infallivel. 2M s

a) A adhesão, causada pela \emph{evidencia}, \emph{do objecto}, \textbf{será} \emph{infallivel}, se tudo o que é \emph{evidente}, fôr \emph{necesshiriamente verdadeiro} ;

porque, então, não é possivel o erro. Ora, effectivamente, tudo o que é evidente, é necessariamente verdadeiro ; porque a \emph{evidencia objectivrm} é a propria verdade, emquanto se manifesta

" \pg{141}  claramente a intelligencia, ou, por outras palavras, é a propria

\emph{e cônformidade} do pensamento com o objecto. Logo a adhesão, causada pela \emph{evidencia do objecto}, é infallivel \footnote{A evidencia, sendo um phenomeno objectivo e subjectivo, é o proprio objecto, é a propria verdade, emquanto se manifesta claramente a intelligencia. Ora desde que um objecto se manifesta claramente a \emph{intelligenejà}, esta percebe immediatamente o objecto como se manifesta, isto é, como é em si mesmo e por isso se acha necessariamente conforme com elle. Um erro na percepção do objecto evidente não é possivel, a não ser que se admitta um d' estes dois absurdos: ou que o objecto, que se manifesta claramente, não se manifesta tal, qual é em si mesmo, ou que a faculdade não é capaz de perceber o objecto proprio, quando este Ihe se manifesta conve- .nientemente. --- Alem disso, sendo um.dictame. de senso commum- que \emph{tudo o} que é \emph{evidente é verdade} (porque todos appellam para a evidencia, como --- .. para a prova infallivel da verdade), se a evidencia podesse estar unida com a falsidade, a inteligencia seria fatalmente arrastada para a mesma falsidade: o que importa a destruição da mais nobre das nossas faculdades --- Portanto, a \emph{evidencia} está por tal modo unida com a \emph{verdade}, que aquella não pode estar sem esía; embora a. \emph{verdade} possa estar sem a evidencia (\emph{intrinseca}), \emph{O} que se verifica nas verdades da fé.\par D' ahi se segue que este principio --- o que é evidente é verdade --- é analytico. Com effeito, um principio é \emph{analytico}, quando a conveniencia entre o predicado e o sujeito se percebe pela simples \emph{analyse} dos termos. Ora entre estes dois termos --- \emph{evidente} e verdade --- existe a maxima conveniencia; porque \emph{evidente} significa que o objecto apparece a intelligencia e é percebido como é em si mesmo, e \emph{verdade} significa que o objecto é na \textasciicircum{} intelligencia como é em si mesmo. Logo esse principio é \emph{amalytico}.\par Balmes, philosopho espanhol, diz (\emph{Phil}. \emph{fund}., 1. I, 9, 22) que esse principio, --- o que é evidente é verdade, --- embora certissimo, não é evidente, não é \emph{dnalytico}; porque a \emph{analyse} não pode, por si, descobrir se uma \emph{idea}, embora clara, corresponde, ou não, ao \emph{objecto}, sendo necessario, para isso, transpor um abysme, qual é a passagem do subjectivo para \emph{o} objectivo. --- Como se vê, esse escriptor parte de um principio falso. E o principio falso é que a \emph{evidencia} Seja uma \emph{idea}, completamente \emph{subjectiva}. Ora a evidencia € tambem, e sobretudo. \emph{objectiva}, porque é o proprio objecto, emquanto se manifesta claramente é jntelligencia, é a propria conformidade da idea com o obiecto, percebido \emph{directamente} por uma intuição clara e immediata. Por isso, o que é evidente é conforme com a realidade, ou é a propria realidade, emquanto é percebida pela intelligencia. --- Estes, e outros equivocos nascem de um grande erro, commum aos \emph{subjectivistas} de todas as raças, que negam a faculdade perceptiva a possibilidade de apprehender \emph{directo} e \emph{immediatamente} o proprio objecto, existente na realidade.}.

b) A adhesão, causada. pela. evidencia. do. objecto, pode .re----.. = ferir-se --- ou a uma verdade necessaria, --- ou a uma verdade \emph{contingente}, --- ou a uma verdade de fé. Ora, seja qual fôr a espetie de verdades, a que se refere, a adhesão é sempre \emph{infallivel} \pg{142} . --- E' \emph{infallivel}, quando a verdade é necessaria ;. porque, neste caso, a intelligencia vé, de um modo \emph{immediato} ou \emph{mediato}, que o predicado convem essencialmente ao sujeito, e porque vê essa conveniencia,-é-impossivel. que se. engane, --- E' \emph{infallivel}, quando a verdade é \emph{contingente}: porque, quando -o objecto se manifesta evidentemente ás faculdades perceptivas, estas etavam bem dispostas e convenientemente applicadas, e por isso não é possivel o erro. --- E' \emph{infallivel}, quando a verdade é de fé; porque quando a auctoridade da testemunha é evidente, é signal de que a mesma testemunha é dotada de' \emph{sciencia} e de \emph{veracidade} : o que é sufficiente para excluír todo o receio de errar. Logo a adhesão a verdade, causada pela \emph{evidencia do objecto}, é infallivel \footnote{Este argumento será mais desenvolvido e corroborado no capitulo seguinte, quando tratarmos: dos "\emph{etos} para. 3. consecugio da: verdade. dade. Não ha nem 1 pode baver \emph{u} uma evidencia \emph{verdadeira} e uma evidencia \emph{falsa}. O que pode haver, o que ha, é a evidencia \emph{real}, e a evidencia tida como \emph{tal}, quando o não é. --- Muitas vezes um escriptor proclama \emph{evidente} um juizo seu, que não é \emph{realmente} evidente. Deriva jaso da leviaudade, dos prejuizos, das paixões; e sempre do imperio dà vontade, a qual quer que haja evidencia onde não ha. --- A evidencia \emph{ndo} se \emph{demonstra}: porque éa visibilidade do que é, e o que se vé, ou \emph{o} que torna visivel uma coisa, não precisa de ser-demonstrado. ---- Mas: a evidencia, se não se demonstra, exa- »\textasciicircum{}\emph{ina-se}; porque a verdade não é visivel a. nossa intelligencia, como um objecto externo é visivel aos nossos olhos. --- O exame deve referir-se --- a) á \emph{existencia} da evidencia, --- \emph{b}) á sua \emph{localização}, --- c) á Sua \emph{extensão}.\par a) Quanto. á \emph{existencia}, devemos advertir --- que a evidencia, sendo uma luz, tem, como a luz do dia; varios. graus, é mais ou menos brilhante, --- e que, assim como não é possivel determinar exactamente o momento em que termina a noite e principia" o dia, assim tambem não é possivel determinar o grau de clareza, necessario para constituir a verdadeira evidencia. Portanto ha tambem evidencias incértas e \emph{duvidosas}. Nestes casos é preciso evitar a precipitação nos'juizos, para que não se diga \emph{certamente} eviden que o é apenas provavelmente. Quando, porém, a evidencia é plena, completit, clarissima, como.a luz do meio-dia, ella é o verdadeiro motivo, o. fundamento inabalavel da certeza e da infallibilidade dos nossos juizos.\par \emph{b}) Quanto á \emph{localisação}, devemos examinar se a evidencia é \emph{intrinseca} ou exirinseca, )sic é, se o objecto. que se diz \emph{evidente}, está \emph{dentro} ou \emph{fóra} de nós. A evidencia \emph{intrinseca} só mostra uma existencia \emph{intrinseca}, e a evidencia \emph{exirinseca} só mostra uma existencia \emph{extrinseco}. É, pois, erroneo concluir de uma evidencia \emph{in£rinseca} para uma existencia \emph{extrinseca}. O allucinado, que ainda está na posse de si mesmo, tem, ao mesmo tempo, --- á \emph{evidencia exírinseca da} sua visto normal, \emph{pela} qual \emph{percebe} um objecto\pgn{143}externo, --- e a evidencia \emph{inírinseca} da sua visão anormal, pela qual percebe os phantasmas da sua imaginação alterada. Ambas estas evidencias são verdadeiras. O essencial é que o allucinado não attribua uma existencia \emph{extrinseca}, á taes phantasmas, que só têm uma existencia \emph{in£rinseca}. --- Este exame ácerca da localisagzo da evidencia, se é possivel quando o sujeito se acha no seu estado normal, na. posse. de..si..mesmo,.. é. impossivel quando. elle.se acha num estado anormal, p. ex., no sonho ou na loucura; porque, então, o sujeito é victima --- não de uma \emph{falsa} evidencia exirinseca (porque a evidencia exírinseca, nestes casos, não existe) --- mas do falso habito, innato ou adquirido, de julgar como existente na realidade uma coisa, que só existe na propria imaginação.\par \emph{c}) Quanto á \emph{extensão}, é preciso examinar se o objecto, que se diz evidente, é \emph{simples} ou \emph{complexo}. Quando é \emph{complexo}, devemos averiguar se \emph{cada} uma das partes, de que se compõe, é dotada, ou não, de evidencia. No caso negativo, é um grande erro attribuir a um objecto complexo uma evidencia \emph{total}, quando é apenas \emph{parcial}. Um pequeno \emph{erro}, a minima \emph{probabilidade}, que se introduza na serie dos nossos raciocinios, basta para tornar falsa, ou obscura a conclusão. Mas, então, a falta não deve attribuir-se á evidencia, mas a nossa ignorancia e a limitação da intelligencia -bumana, que não pode abraçar, com um só relance, uma questão complexa. --- É para auxiliar a intelligencia nas investigações e leval-a é descoberta da evidencia, que a \emph{Logica} prescreve as suas regras.}.

%%>>> t141

\pg{143} \tese{141}{A evidencia do objecto é o universal e supremo motivo da certeza.}

a) É \emph{o motivo universal da certeza}. --- Motivo \emph{universal} da certeza é a razão commum a todas as verdades, isto é, é a vazão, pela qual adherimos a \emph{todas} as verdades, seja qual fôr a sua especie ou ordem. Ora essa razão commum é 2 \emph{evidencia do objecto}; porque, como attesta a consciencia, não podemos adberir, e não adherimos a nenhuma verdade, senão porque o seu objecto se manifesta \emph{evidentemente} a faculdade .perceptiva. Logo a \emph{evidencia do objecto} é o motivo universal da certeza \footnote{Se.\& \emph{evidencia do objecto} é o motivo \emph{universal} da certeza, segue-se que elle deve necessariamente applicar-se a todas as ordens de verdades. De facto, a evidencia \emph{exérinseca} é o motivo necessario e sufficiente para a intelligen.ia adherir ás verdades da fé, ou da ordem \emph{moral}; e a evidencia \emph{intrinseca é}. \emph{o} motivo necessario e sufficiente pata adherir ás verdades \emph{necessarias} e contingentes, ou da ordem \emph{metaphysica} e \emph{physica}.\par Advirta-se, porém, que, sendo \emph{diverso} o \emph{objecto}, que se diz \emph{evidente}, conforme a diversidade das verdades, \emph{diversa} tambem é a \emph{evidencia}, propria de cada ordem das mesmas verdades. \emph{A} evidencia, que é propria das verdades da ordem \emph{methaphysica}, é diversa da evidencia, que é propria das verdades da ordem \emph{physica} e da ordem \emph{moral}. D' ahi a Liriplice especie da evidencia: \emph{metaphysica}, \emph{physica} e moral. E isto é muito razoavel. A evidencia é uma qualidade, que se refere á propria entidade da verdade; e toda a qualidade,\pgn{144}que se refere á entidude de um sujeito, é diversa conforme a diversidade do proprio sujeito. --- D' este modo. \emph{cada especie} de \emph{evidencia} se torna motivo de certeza \emph{proprio} de \emph{cado ordem} de \emph{verdades}. Assim, 'a evidencia \emph{metaphysica}. é: o” motivo” de certeza proprio" das verdades da ordem \emph{metaphysica}; a evidencia \emph{physica}. \emph{é o} motivo. de certeza proprio das verdades da ' ordem \emph{physica}; e a evidencia \emph{moral} é o \emph{motivo} de certeza proprio das verdades da ordem \emph{morali}. --- D' onde se vê a differenca entre o motivo \emph{universal e} o motivo \emph{particular} da certeza, O motivo universal é a \emph{evidencia objectiva}, isto é, que reside no \emph{objecto em geral}. O motivo. particular é \emph{aevidencia}, que reside no objecto particular, isto é, na verdade de uma certa especie, Por isso, \emph{o} motivo \emph{particular} da certeza \emph{metaphysica} é a \emph{evidencia objectiva} da essencia das \emph{coisas}; o da certeza. \emph{pliysica} é a \emph{evidencia objectiva} da \emph{constancid} das \emph{leis physicas}; o da certeza \emph{moral} é a \emph{evidencia objectiva} da \emph{constancia} das \emph{leis moraes}, ou da \emph{auctoridade}. \emph{A evidencia objectiva} é o \emph{genero}; \emph{a evidencia objectiva} da \emph{essencia das} coisas, a da \emph{constancia} das \emph{leis physicas} e \emph{das leis} moraes' são \emph{especies A especie} deve conter o genero.}

b) \emph{E o motivo supremo da certeza}. --- Motivo supremo da certeza é o que não depende de outro, e é a razão \emph{ultima}, que nos determina a adherir a verdade. Ora tal é a \emph{evidencia do objecto} ;. porque. não..adherimos a verdade, senão. quando o seu objecto se torna \emph{evidente}, e, depois de se tornar evidente, não procuramos outro motivo para adherir, e para esta \emph{evidencia} appellamos; como para a razão \emph{ultima} da nossa adhesão. Logo a \emph{evidencia do objecto} é o motivo \emph{supremo} da certeza \footnote{A \emph{evidencia objectiva}, pois, é o motivo supremo da certeza, --- porque é a razão necessaria e \emph{sufficiente}. para a.intelligencia adherir a verdade, --- porque nenhuma outra razão é muis \emph{efficac} do que essa evidencia, --- porque todas as outras razões \emph{fundam-se} e \emph{resolvemw-se} na evidencia. :\par Dizem: Podendo a evidencia ser \emph{real} ou \emph{apparente}, deve admittir-se um outro motivo superior, que sirva para discriminar uma de” outra. Logo a evidencia não é motivo supremo da. certeza.\par Resposta. É certo que, para discriminar a certeza \emph{real} da apparente, é necessario um motivo. Este motivo, porém, --- não é uma coisa \emph{distinta} da evidencia, --- mas é a \emph{propria evidencia}, a qual manifesta por si a propria presença, quando é \emph{real}, ou a propria ausencia, quando é apenas \emph{apparente}. Com effeito, a evidencia é \emph{Tuc}; ora a luz manifesta-se por si mesma, quando é presente, e, quando não é presente, manifesta-se pela sua ausencia, pelas trevas. Por isso, a evidencia \emph{apparente} não é verdadeira \emph{evidencia}; e basta uma simples reflexão desapaixonada para averiguar que ella é apenas \emph{apparente}; ao passo que a evidencia \emph{real} se mostra tanto mais pura e brilhante, quanto mais intensa fôr a reflexão nos motivos da adhesaão. E}. \textbar{}

%%>>> t142

\tese{142}{A evidencia do objecto é percebida pela intelligencia por diversos meios.}
A \emph{evidencia do objecto} \& percebida exclusivamente pela intelligencia; mas esta percepção faz-se por diversos meios, conforme a diversa indole do objecto. --- Se o \pg{144} \pg{145}  objecto é constituido por verdades \emph{necessarias}, o meio, de que a intelligencia se serve para lhe perceber a evidencia, é a \emph{analyse}: d'onde a certeza \emph{metaphysica}. --- Se o objecto é consti- ;'tuido por verdades \emph{contingentes}, o. meio, de. que. a intelligencia . se Serve para lhe perceber a evidencia, é a \emph{experiencia}: d'onde a certeza \emph{physica}. --- Se o objecto é constituido por verdades da fé, o meio, de que a intelligencia se serve para lhe perceber a evidencia, é a \emph{auctoridade}: d'onde a certeza moral (').\footnote{A \emph{evidencia objectiva} é o supremo e universal motivo de certeza. Por isso, quando nos perguntam pelos motivos, por que adherimos a uma verdade, ou unimos um predicado a um sujeito, appellamos para a \emph{evidencia} ' \emph{do.objecío}, e dizemos: isto é assim, \emph{porque é evidente}. Depois d'esta resposta, não procuramos nem apresentamos outro motivo --- Mas pode ainda perguntar-se: Como é que se conhece a \emph{existencia} d' essa \emph{evidencia objectiva} ? Esta segunda pergunta, embora seja um complemento da primeira, comtudo .differe da outra; e a resposta, que lhe se deve dar, consiste em indicar os “\emph{meios}, que a intelligencia emprega para. -descobrir-essa. evidência mo dizemos no texto, são a \emph{analyse}, \emph{a experiencia} e "a auctoridade, Estes \emph{meios} devem ser \emph{especificamente} diversos, porque \emph{especificamente} diversas São as ordens de verdades, cuja \emph{evidencia objectiva} queremos descobrir, --- Da confusão d' estas. duas questões: \emph{qual é o motivo da certeza}? --- \emph{como se} co- -nhece.0 \emph{motivo da certeza} ? --- nasceram todas ás divergencias dos. philosophos, que, querendo resolver a primeira questão, resolveram a segunda, e -cada um a seu modo, isto é, querendo indicar o \emph{motivo} da certeza, indicaram o \emph{mein}, por que se descobre. --- Devendo, no capitulo seguinte, tratar d' esses \emph{meios}, que aqui apenas indicamos, vamos agora expôr e criticar as principaes opiniões dos philosophos ácerca do. \emph{motivo} da certeza.\par O \emph{motivo da} certeza costuma hoje chamar-se \emph{criterio da verdade}. Mas esta segunda denominação não tem sempre o mesmo sentido. Lendo os eseripios dos philosophos modernos, vemos que o \emph{criterio da verdade} --- para alguns significa a \emph{regra}. \emph{a medida}, pela.qual discriminamos .a.verdade da falsidade, --- para outros denota o \emph{meio}, o \emph{instrumento}, de que o homem se .Serve para alcançar a verdade, ou a certeza, --- para outros, finalmente, .indica o \emph{motivo}, a \emph{razão}, que determina a intelligencia a adherir com fir- -meza' a verdade. --- Se por \emph{criterio da verdade} se entende a norma, ou-a \emph{medida}, que sirva para distinguir a verdade da falsidade, a guestão pode ser prejudicada por um falso supposto. Podem suppór, como effectivamente Suppõem os-sequazes de Descartes, que a intelligencia deva, no principio das suas investigações, duvidar de \emph{tudo} e que, para depór. essa duvida Universal, seja necessario procurar uma \emph{norma}, ou uma \emph{regra}, que, appliesda. aos juizos, faça: distinguir o verdadeiro do falso. Ora, como vimos e veremos, tudo isto é impossivel. Ha verdades por si tão evidentes, .que se percebe:f immediatamente, sem necessidade de normas, ou \emph{regras}, e de cuja falsidade nem é possivel suspeitar. Estas verdades \emph{immediatamente} evidentes podem chamar-se \emph{regras}, ou \emph{criterios}, Ou \emph{principios} das conclusões, emquanto estas se acham virtualmente contidas e se resolvem naquellas, Mais tarde\par Vol. 1, u\pgn{146}voltaremos a este assumpto, --- Se por \emph{criterio da verdade} se entende o \emph{meto}, 0 \emph{instrumento}, de quê \emph{o} lómem se serve para alcançar a verdade, resolve-se a questão facilmente, indicando esses \emph{meios}, e mostrando a sua legitimidade. --- Se por \emph{criterio} da \emph{verdade} se entende 0 \emph{motivo}, que leva a adherir com firmeza a verdade, como hós \emph{o} entendemo\$, então esse motivo ' consiste na evidencia \emph{obiectiva}. --- Aqui o e a \emph{evidencia objectiva} € \emph{criterio} de verdade, não só quando por \emph{criterio} se entende o \emph{motivo}, que nos leva a adherir firmemente a verdade, --- mas tambem quando por \emph{criterio} se entende a \emph{medida} dos nossos juizos, e a \emph{regra}, pela qual discriininamos \emph{a} verdade da falsidade. Que a \emph{evidencia objectiva} seja, ao mesmo tempo, \emph{motivo} e \emph{regra}, apparece claramente no acto de fé, com que acreditamos \emph{hum} facto historico. Se nos perguntarem pelo motivo do nosso assepso, respondemos que é a auctoridade das testemunhas, E se nos perguntarem" ainda pela \emph{regra} d'este nosso assenso, respondemos sempre que é a auctoridade das testemunhas, Eis, portanto, que uma e a mesma coisa, qual é a \emph{auctoridade das testemunhas}, €, ao mesmo tempo, motivo e \emph{regra} da adhesão, E, se isto se verifica nas verdades da fé; porque não se havia. de verificar nas verdades da \emph{sciencia}? A razão é sempre a mesma. De facto, quando queremos conhecer o \emph{motivo} da nossa adhesão a uma verdade racional, recorremos a \emph{evidencia objectiva}. E. quando queremos convencer os outros de que os nossos juizos' são \emph{verdadeiros}, tambem recorremos a \emph{evidencia} ' \emph{objectiva}, como a \emph{uma regra} segura; porque resolvemos esses juizos nos \emph{primeiros principios}, se elles são mediatos, e nas \emph{ideas}, representativas da entidade das coisas, se são \emph{imwumediatos}: O que importa um recurso para a \emph{evidencia objectiva}. --- Todavia, a \emph{evidencia} não é \emph{motivo} pela mesma razão por que é \emph{regra}; porque \emph{motivo} e \emph{regra}, embora possam convir a uma e mesma realidade, comtudo differem nos seus \emph{conceitos}. ;\par : * :\par Rejeitando'a \emph{evidencia vbjectiva}, como universal e supremo \emph{motivo}, ou \emph{criterio} de certeza, os philosophos. apresentaram outros \emph{motivos}, OU \emph{criterios}, --- como elles -os:denominam. Pudem estes reduzir-se a duas classes: uns são \emph{inirinsecos}, Ou \emph{subjectivos}, outros são \emph{exírinsecos}, Ou \emph{objectivos}.\par 2) Üriterios intrinsecos. --- Os criterios \emph{intrínsecos}, OU. \emph{subjectivos}, São: \& intelligencia humana, a intuição directa e immediata da Verdade infinita, a idea clara e distincta, o. principio de contradicção, a incomprehensibilidade da proposição contradictoria, \emph{o} testemunho da consciencia, o instincto cego da natureza, o sentimento, a vontade livre, a verificação da experiencia sepsivel. --- De cada um d'.estes criterios faremos uma breve critica. [\textbar{} \emph{a}) A \emph{inteligencia humana}. --- Não faltaram phiülosophos, desde os\par tempos de Protagoras até aos nossos dias, que ensinassem ser a \emph{inteligencia humana a medida} das coisas, a unica \emph{regra} da \emph{verdade} e da \emph{falsidade}, do bem e do mal. --- A opinião não poderia ser mais absurda; porque é absurdo que uma e a mesma coisa seja, a0 mesmo tempo e sob o mesmo aspecto, verdadeira e falsa, boa e má, só porque a uns parece verdadeira e boa e a outros parece falsa e má. É o relativiemo, que leva ao \emph{scepticismo}.\par \emph{b}) A \emph{intuição directa e immediata} da Verdade \emph{infinita}. --- Éa opinião\pgn{147}dos \emph{onfologos}, --- Malebranche (1633-1715), Gioberti (1801-1859), Rosmini (1797- --- 1855), --- que só reconhecem como certas as verdades, correspondentes a Verdade infinita, da qual teriamos uma visão \emph{directa} e immediata. --- Mais tarde refutaremos o \emph{Ontologismo}. Limitamo-nos aqui a dizer que esse criterio, se é proprio dos bemaventurados, que contemplam a Deus \emph{intuitivamente}, não pode ser proprio dos mortaes, que só do conhecimento das coisas \textbf{creadas} se elevam ao conhecimento do Creador, e sempre .pela abstracção.\par \emph{c}). A \emph{idea clara} e \emph{distincta}. --- É a opinião de Descartes, para quem são \emph{certas} só as verdades, de que temos uma idea clara e \emph{distincta}. --- Este criterio não pode admittir-se. Por quanto, o \emph{criterio}, Sendo o \emph{motivo}, ou a \emph{regra}, que nos dirige no juizo \emph{relativo} a um \emph{objecto}, deve ser \emph{objectivo}: por- . que é o objecto que se impõe a faculdade, e não a faculdade ao objecto. Ora a \emph{idea clara} e \emph{distincta}, \& uma coisa \emph{subjectiva}: porque é a imagem \emph{intencional} do objecto, e a sua \emph{clareza} e \emph{distincção} dependem ordinariamente da capacidade do sujeito e do grau de attenção. --- Além d'isso, o criterio da verdade deve s ser immutovel, como o é E . propria verdade. Mas, a. \emph{idea}.\par homem pode perceber clara. e \emph{distinctamente} uma coisa, e. um. .outro-pode-. perceber, com a mesma clareza e distinecção, o contrario : d' onde o \emph{subjectsvisio}. --- Notamos que a formula de Descartes «o que se percebe clara e distinctamente é verdade» differe tanto da formula escholastica «o que é evidente é verdade» quanto o \emph{subjectivo} diflere do \emph{objectivo}. Por quanto, --- na formula cartesiana, exprime-se a \emph{percepção} Ou a \emph{intuição} da evidencia, isto é, do objecto evidente, ao passo que, na escholastica, exprime-se a propria \emph{evidencia} do objecto, ou o proprio \emph{objecto} emquanto evidente; --- na primeira, é a propria \emph{visão} que se indica como criterio; na outra, é a \emph{visi bilidade} do \emph{objecto}; --- na primeira, á \emph{evidencia} é uma qualidade do \emph{sujeito}; na outra, a \emph{evidencia} é por si e primeiramente .uma qualidade do \emph{objecto}.\par \emph{d}) \emph{O principio de} contradicção. --- E" a opinião de Wolf, que pretende deduzir d'esse principio \emph{todas} as verdades, como outras tantas conclusões dé um unico principio. --- Tambem esta opinião é falsa. O principio de contradicção, embora esteja contido em todas as verdades, que nelle se fundam e se resolvem, comLudo não é o \emph{motivo}, \emph{a razão}, pela qual adherimos a verdade. Porque estamos nós certos de que o sol resplandece? Estamos certos, --- não porque o \emph{sol} não pode, ao mesmo tempo, resplandecer e deixar de resplandecer, isto é, não pelo \emph{principio de} contradicção, --- mas porque Lal se manifesta aos nossos olhos, isto é, pela \emph{evidencia objectiva}. Porque estamos nós certos de que o effeito exige uma causa? Estamos certos, --- não porque o \emph{effeito} não pode, ao mesmo tempo, ser e deixar de Ser \emph{ejfeito}, --- mas porque o predicado apparece claramente contido no sujeito. --- Além d'isso, o principio de contradicção não pode ser o motivo supremo da certeza. Por quanto, se nos perguntarem pela razão, por que julgamos certo esse principio, respondemos: porque é \emph{evidente}. Logo tambem o principio de \emph{gontradiccZo} depende da \emph{evidencia}.\par €) A \emph{incomprehensibilidade da proposição contradictoria}. --- É o criterio de \emph{Spencer}, que affirmou ser \emph{certa só a} proposição, \emph{cuja contradictoria} É\pgn{148}\emph{incomprehensivel}; assim a. proposição --- o sol é luminoso --- é certa, porque é incomprehensivel a contradictoria --- o sol não é luminoso. --- Tambem "este \emph{criterio} é falso. A palavra \emph{incomprehensicel} pode significar que uma \emph{coisa não} se comprehende --- \emph{ot} porque o predicado \emph{repufmma} ao sujeito, --- oü porque não somos capazes de ver a relação entre os dois termos. Ora, tanto no primeiro como no segundo sentido, essa \emph{óncomprehensibilidade} não pode ser o \emph{motivo} ou \emph{criterio supremo} da certeza. --- Não pode sel-o no primeiro sentido ; porque aquella repugnancia suppõe já outro criterio, que a \emph{manifeste}. --- Não pode sel-o no segundo; porque aquella incapacidade é subjectiva, mutavel e relativa. lo facto de eu não comprehender esta proposição --- a terra gira em volta do sol --- não pode deduzir-se que seja certa a contradictoria --- a terra não gira em volta do sol.\par \emph{P}) O \emph{gastemunho da consciencia}. --- É a opinião de \emph{Galluppi} e d' alguns diseipulos'Me \emph{Kant}, --- Mas tambem esta é inacceitavel. O \emph{testemunho da} consciencia SÓ accusa a \emph{existencia} dos factos, que se realizam no espirito, mas-não n nifesta a \emph{essencia}. d' elles ; e por isso não manifesta a \emph{ração}, O \emph{motivo}, por que adberimos a verdade. --- Além d'isso, o acto da \emph{consciencia} ê: um: conhecimento \emph{reflexo} e-refere-se totalmente ao sujeito pensante, Ora “o \emph{motivo}, ou \emph{criterio} da verdade deve extender-se não só ao conhecimento \emph{reflexo}, mas tajnbem ao \emph{directo}, € referir-se não só ao \emph{sujeito}, mas tambem, e sobretudo, ào \emph{objecto}. --- Finalmente, a \emph{consciencia} é um \emph{meto} que nos manifesta « evidencia dos nossos conhecimentos. Ora o \emph{motivo}, que nos\par \emph{értos} de uma verdade, não é a \emph{manifestação} da evidencia, mas é a propria \emph{evidencia}. Logo o \emph{motivo}, Cu \emph{criterio} da verdade não é o testemunho da \emph{consciencia}.\par 4) \emph{O instincia cego da} natureca. --- E' a opinião de Reid (1710-1786), seguido por Dugald Stewart (1773-1835) --- Mas nem esta pode admittir-se. Porquanto, esse criterio repugna á natureza da nossa intelligencia, a qual, sendo uma faculdade cognoscitiva, não pode julgar com certeza, nem adherir firmenienteé, a uma verdade, se não lhe resplandece um motivo racional, como os olhos não podem ver sei] luz. Ora o \emph{cego instincto}. da natureza é destituido de todo e qualquer motivo objectivo de assenso, e por isso não pode levar-nos a verdade, como "um cego não pode indicar-nos o caminho. --- Mais, a nossa certeza, se se baseasse nó \emph{cego imstincto} da natureza, além de irracional, seria exclusivamente \emph{subjectiva}, porque não dependeria, de modo algum, da intuição do objecto. D' ahi o swbjeclivismo e o \emph{scepticismo}.\par \emph{h}) \emph{O sentimento}. --- É a opinião de lacobi (1743-1819), chefe da "eschola . \emph{Senfinenialisia} na Allemanha. --- Tambem este criterio, afim ao precedente, é inacceitavel, Na verdade, o criterio deve ser immulavel. Ora \emph{o sentimento} é naturalmente \emph{mutavel}. --- Além d'isso, o \emph{criterio da} verdade deve ser um e o mesmo para todos os homens. Ora o \emph{sentimento}, embora relativo á mesma coisa, varia conforme a diversidade do caracter de cada um, da sua constituição physica, da sua educação, dos seus costumes, prejuizos, e mais condições pessoaes. --- Finalmente, o \emph{sentimento} é ceo, deve ser guiado pela razão, e por isso não pode ser norma da verdade, nem \emph{motivo da} certeza.\par 1) A \emph{vontade livre}. --- É a opinião dos sequazes do \emph{pragmatismo} e do\pgn{149}do tempo, pelo bracharense Francisco Sanchez (1562-1632), que dizia: « Nada se sabe, só a Fé resolve toda. a duvida e ensina a verdade», ---- Mas este criterio da \emph{auctoridade divina}, \emph{ou da}. \emph{Revelação sobrenatural}, não é \emph{universal}, nem é o \emph{primeiro} na ordem do tempo. --- Não é \emph{universal}; porque, se é criterio das verdades da \emph{JF}, que nos são reveladas por \textsc{Deus}, não é criterio das verdades da \emph{Sciencia}, que descobrimos pela luz natural da razão. --- Nem é o \emph{primeiro} na ordem do tempo; porque não podemos admittir a auctoridade\pgn{150} divina, se não estivermos anteriormente certos da existencia de Deus, da sua infallibilidade, do facto e da autbenticidade da revelação, e de muitas outras verdades. Por isso o uso da razão precede a fé. CM \emph{b}) A \emph{auctoridade humana}. --- É a opinião de R. de Lamennais (1780-\par )\emph{e} " Bonnetty. (SS 879 que não reconhecem como \emph{certo} senão o que\par versa] ; \emph{porque} não se extende a muitas outras verdades, que descobrimos 8 luz natural da razão. Nem 6.0 \emph{primeiro} na ordem do tempo; porque\par não podemos acceitar do consenso do genero humano uma verdade, se antes. não estivermos certos de que esse consenso existe, de que não é o resultado dos prejuizos e se distingue das opiniões individuaes. Ora, para averiguar tudo isso, é necessario o uso de outros criterios. --- Alem d'isso, se o homem devesse admittir, como certo, só o que nos é ensinado pelo con- 9 Senso do genero humano. á maior parte dos homens achar-se-hia numa perpelua duvida, e o.estado normal do genero humano seria o \emph{Scepticismo}. Nem seria muito melhor a condição dos doutos, que, só depois de aturados estudos, poderiam chegar a conbecer \emph{o} consenso do genero ] humano e apenas com relação a um ou outro problema.}

%%>>> t143

\tese{143}{O motivo da certeza e o senso commun.}
Todos os homens, .doutos e ignorantes, concordam em que o \emph{motivo} da \pg{146} \pg{147}  \emph{ceriez} za, consiste na \emph{evidencia do objecto}. Com e(feito, quando elles desejam convencer alguem, procuram tornar-lhe \emph{evidente} o \emph{objecto}, \pg{148}  e para esta \emph{evidencia} appellam, como para a ultima e irrespon- .divel prova das suas convicções. D' ahi as expressões : isto é \pg{149}  evidente.» --- -« cedemos a evidencia » --- «não podemos con- \emph{PP} tradizer a evidencia ». --- E o senso commum affirma --- não Só L "immanentismo, que dizem ser verdade só o que pode ser zppetecido como N bom e pode ser util para “a vida" pratica; --- Tambem: este criterio é inad- ------------ : missivel. Porquanto, se só fosse verdade ou certo o que parece bom e «til, =. 0' homem seria a medida da verdade, conforme a opinião de Protagoras, e \textbar{} coisas contradictorias seriam simultaneamente verdadeiras ; visto que, muitas 2 vezes, para um é boa e util uma contradictoria, e para outro é boa e util outra contradictoria. Ora isto vão só não é criterio de verdade. mas é a = destruição de toda a verdade, porque não haverja a maneira de a distinguir do erro. --- Além d' isso, o acto \emph{da} vossa \emph{livre vontade} não pode ser o s«- \emph{O premo motivo} da nossa certeza. Porquanto, o que é que provocará esse acto ? --- Será um \emph{impulso cego}. ou um \emph{motivo intellectual}? Se é um impulso \emph{cego}, O E \textasciicircum{} acto da vontade não é \emph{livre}, mas é \emph{fatal} e \emph{instinctivo}, e teremos todos os \emph{Ps} inconvenientes, apontados na opinião de \emph{Reid}. Se é um \emph{motico intellectual}, \textasciicircum{} que dirige á nossa escolha entre opiniões diversas, acceitamos, com isso, ------ o criterio \emph{Zn£ellectualista}, --- a visibilidade do que é, como supremo motivo da \emph{ES} evidencia. --- Finalmente, o acto da vontade \emph{livre} não pode ser o principio \emph{num} adequado da nossa ndhesão a. verdade,. nem.ás. verdades da fé. Porquanto,: Ll. i adhesão ás verdades da fé

to \emph{intellectual}, deve fundarse e funda-se em razões certas, em motivos certos de \emph{credibilidade}. Ora estas razões, estes motivos são percebidos pela \emph{éntelligencia}.. Logo a \emph{f£ore vontade} não pode ser o criterio, o motivo de verdade, ou de certeza.

\emph{j}) \emph{A verificação da experiencia}, sensivel. --- É a opinião de Comte, Littré e de todos or \emph{positivisfes}, que reconhecem como certo só o que foi \emph{verificado} pela experiencia \emph{sensivel}. --- Não pode acceitar-se essa opinião. Antes de tudo, a experiencia \emph{sensivel}! pode ser e é um \emph{meio} para a. percepção, directa ou reflexa, de algumas verdades. mas.não é o proprio \emph{motivo}. que determina a intellizencia a adhes\&o. --- Além d'isso, esse criterio, se serve para algumas \emph{hypotheses}, não. serve para \emph{todas} as verdades, e por isso não é universal. A mathematica e as sciencias naturaes admittem, como certas, muitas coisas, que não foram, 'nem podem ser verificadas pela experiencia \emph{sensivel}. -Os astronomos predizem o phenomeno do eclypse, e estão certos , ... de que se ha de realizar, mesmo antes de verilicado pela experiencia. --- . Finalmente, por ser esta opinião uma consequencia dos principios do \emph{posifivismoy}; refüta-se com todos os argumentos, que serão; mais tarde, adduzidos contra esse systema.

b) Criterios extrinsocos. --- Os criterios \emph{exírinsecos}, ou \emph{objectivos}, são: a * auctoridade divina e a auctoridade humana.

a) A \emph{auctoridade divina}. --- É a opinião de N. Huet (1630-1721). G. de Bonald (1754-1840). L. Bautain (1796-1867), precedidos. pelo menos na ordem \pg{150}  que a \emph{evidencia do objecto} é o motivo da certeza, --- mas tambem que a intelligencia percebe essa evidencia pelos tres-meios, que são a \emph{analyse}, a \emph{experiencia} e a \emph{auctoridade}: porque os. bomens, quando querem mostrar que: uma coisa é \emph{evidente}, recorrem a um d'esses meios, conforme a indole da coisa \footnote{Com razão escreve o grande Bossuet: « A verdade é o ) que é; a falsidade é o que não é... Podemos entender o que é; mas não: podemos entender Oo que não é... A verdadeira regra para julgar bem é a de não julgar senão quando se vé claramente » (\emph{Connais}. \emph{de Dieu et de soi-méme}, c. Y.}.

%%>>> art-8

\section*{Artigo V.}
\addcontentsline{toc}{section}{Artigo V.}

A \emph{certeza} e a \emph{objectividade das ideas universaes}

%%>>> t144

\tese{144}{Nexo entre as duas questões.}
A questaóPrelativa a \emph{certeza}, está intimamente ligada com a questão relativa á \emph{objectividade das ideas universaes}. Por quanto, a certeza do \emph{raciocinio} suppõe a certeza do \emph{juiszo}. Ora o \emph{juizo} suppoe o valor \emph{objectivo} das ideas universaes. Na verdade, o \emph{juizo} é a affirmacao da conveniencia ou "da\textasciicircum{} nao-conveniencia entre duas. \emph{ideas}; e, como affirmando essa” conveniencia; ou \textasciicircum{} não-conveniencia das ideas, entendemos affirmar a conveniencia Ou a não-copveniencia \pg{151}  das \emph{coisas}, que as ideas representam, o nosso juizo será \emph{certo}, \emph{se:as"ideas} forem plenamente conformes com as coisas \footnote{A \emph{objectividade}, ou o valor \emph{objectivo} das idea \emph{universaes} é uma d'aquellas verdades, que não \emph{precisa} de ser demonstrada, nem o \emph{poe} ser. --- Náo \emph{precisa}; porque é o proprio senso commum que nos diz serem as nossas \emph{ideas} representativas das \emph{coisas}. Quando affirmo que «Pedro é homem», todos entendem que essa idea \emph{universal} «\emph{homem}» tem um valor \emph{objectivo}. entendem que nesse sujeito, que se chama Pedro, ha alguma \emph{coisa}, a qual corresponde á idea homem. --- Não \emph{pode} ser demonstrada. Porquanto, quem faz uma demonstração, deve servir-se e serve-se de \emph{ideas}. Estas --- ou são \emph{objectivas}, --- ou são \emph{ndo-objectivas}. Se são \emph{objectivas}. suppoe-se previamente 9 que se quer demonstrar, e commette-se uma \emph{petição de principio}. Se são \emph{ndo-objectivas}, toda a discussão é inintelligivel e-inutil; --- A nossa demons-: co tração, pois, reduz-se a uma explicação. --- "l'ratamos da \emph{objectividade} dus ideas «niversaes; porque é o valor \emph{objectivo} d'estas que se nexa. --- Se as ideas wüiversaes são \emph{objectivas}, \emph{a-fortiori} são \emph{objectivas} as singulares.}.

%%>>> t145

\tese{145}{Idea universal.}
Idea \emph{universal}, como dissemos, é a -que -representa uma. \emph{essencia}, mas. commum..a..muitos..entes, de modo que a todos e a cada um d'elles pode ser attribuida ; tal é a idea de \emph{humanidade}, que é commuwum a todos . os homens e que se attribue a todos e a cada um d'elles \footnote{\emph{A} idea \emph{universal} representa um objecto \emph{universal}; e esse \emph{objecto} é a \emph{essencia} das coisas, abstrabhida das propriedades individuaes. Por isso podemos dizer indifferentemente que o objecto \emph{proprio} da intelligencia --- é o \emph{Universal}, --- ou é a \emph{essencia} das coisas. S. Thomaz diz: «Intellectus est "niversalnun, non \emph{singularium} » (\emph{C}. \emph{Gent}., \textbar{}. Y, c. 44). « objectum intellectus est ipsa rei \emph{essentia} » (\emph{De} Ver. q. 10, a. 4 ad 1). --- Quando dizemos que a idea \emph{universal}, ou, simplesmente, o \emph{universal}, tem valor \emph{objectivo}, reterimo-nos ao seu coníeudo. que é a \emph{essencia}.}.

%%>>> t146

\tese{146}{Idea universal directa e reflexa.}
A idea \emph{siniversaldirecta} é a que representa. uma \emph{essencir}, considerada. nos seus. \emph{elementos constitutivos}, sem relação aos individuos, nos. quaes se, Encontra ou pode encontrar-se; tal é a idea de \emph{homem}, considerada exclusivamente nos seus elementos de \emph{animalidade} e de \emph{racionalidade}. --- A idea \emph{universal-reflexa} é a que representa a mesma essencia, mas considerada em \emph{relação} aos individuos ou. inferiores, nos quaes pode existir e aos quaes pode attribuir-se; tal é a idea de \emph{homem}; considerada--como- especie; que- abrange: todos os individuos humanos- \footnote{Quando a nossa intelligeneia percebe uma \emph{essencia}, p. e., a essencia do \emph{homem}, percebe-a emquanto é uma determinada natureza, constituida por certos e determinados elementos, sem se importar dos individuos, em que essa essencia se encontra ou pode encontrar-se. Neste caso, a consideração\pgn{152} da intelligencia limita-se exclusivamente a esses elementos. --- Mas a. intelligencia, depois de ter considerado a essencia nos seus elementos, pode tornar a consideral-a, não já em relação aos mesmus elementos, mas em relação aos individuos, em que se encontra ou' pode encontrar-se: Neste caso concebe:se a \emph{essencia} como uma coisa que se refere a. 'muitas, isto é, como uma \emph{forma}, que se encontra neste ou naquelle individuo, sem que nunca se. exgote quanto a extensão. --- No primeiro caso, temos o universal \emph{directo}, que se concebe por uma intuição \emph{directa}; no. segundo, temos o universal \emph{reflexo}, que é o producto de uma consideração \emph{reflexa}...\par Como \$e vé, para a formação do universal \emph{directo} não se requer, da parte da intelligencia, se não uma \emph{abstracção}, pela qual, no objecto proposto, considere exclusivamente a \emph{essencia}, prescindindo das notas ou propriedades individuaes. Assim, a \emph{humanidade}, que se encontra em Pedro, Paulo; André, etc.;. emquanto se considera “tal, qual- existe em cada um d' elles, é.uma.coisa. \emph{singular}. e \emph{individual}, porque a humanidade de Pedro não é a humanidade de Paulo. Mas, se essa humanidade, pela acção da inteligencia, é depurada das propriedades individuaes, de que se acha cereada no tempo e no espaço, isto é, na. sua subsistencia real, torna-se immediat\&mente \emph{universal}. Diz S. Thomaz: «Una et eadem natura, quae singularis erat et individuata per materiam in singularibus hominibus, efficitur: postea. \emph{universalis} per actionem intellectus depurantis ipsaib a conditionibus, quae sunt hic et nunc» (De \emph{Univ}. t. 1). --- Tal abstracção não' é real, mas \emph{ideal}; pois consiste em que a intelligencia considera exclusivamente a essencia, que é o seu objecto proprio, prescindindo dos caracteres, de quo se acha cercada nos individuos. Tal abstracção é legitima, e não pode ser causa ou occasião! de erro Quando abstrabimos, não julgamos que uma essencia existe em si mesma de um modo \emph{abstracto} e \emph{separada} dos\par caracteres individuaes; mas simplesmente consideramos-a essencia, e pres eindimos dos caracteres: individuaes. Acontece o mesmo com relação á percepção sensitiva: Podemos, por ex., perceber com os nossos olhos a cór de um fructo, sem.comtudo lhe percebermos o sabor (Cf. Sum. \emph{Th}., \emph{p}. I, q. 85, a..1 ad 1; --- \emph{De potentiis animae}, c. VI). \textasciicircum{}- Esta operação é proporcionada. \&.natureza. da. nossa intelligencia. Sendo esta uma facu'dade espiritual, istoé, independente e livre de toda a materia. á suu \emph{operação} deve ser tambem independente e livre da materia, e por isso livre e independente d' essas condições deve ser o seu \emph{objecto}. Ora um objecto é livre e independente das condições da materia, quando se concebe de um modo \emph{universal}. Por isso. o mesmo \emph{objecto}. que os sentidos percebem como \emph{singular}, pela intelligencia é percebido como \emph{universal}, conforme o celebre adagio de S. \emph{Sev}. \emph{Boecio}: « Singulare dum sentitur, universale dum intelligitur ».\par Se para a formação do universal \emph{directo} é necessaria e sulficiente, da parte da intelligencia, a abstracção, para a formação do universal \emph{reflexo}, é necessaria a abstracção e a \emph{comparação}. --- É necessaria a abstracção. Por\pgn{153}quanto, o universal \emph{reflexo} é, propria e formalmente, uma coisa commum: a muitas «unum versus alia». Ora essa coisa una, que é a essencia, subsiste de um modo material e concreto, e se acha: multiplicada, quanto ao ser, em muitos individuos. É, pois necessario que adquira uma subsistencia intencional, pela qual ella, que é multiplice e divisa.na realidade, se torne una e indivisa na propria intelligencia, e seja uma imagem, que represente o que esses individuos têm de commum, ou de semelhante --- Mas não basta este primeiro acto da jútelligencia É necessario um segundo, que é a' comparação. Na verdade, o universal \emph{reflexo} é, como dissemos, uma coisa. commum a muitas --- unum versus \emph{alia}. Ora, para baver. esse universal,. não basta que \emph{a} intelligencia forme, pela \emph{abstracção}, uma imagem intencional. que representa o que os individuos têm de commum e que por isso possue a capacidade, ou a não-repugnaüéia- para-existir: em muitos particulares mà; € necessario que, depois de a ter formado, \emph{reflicia} sobre essa essenci \emph{siracia}, que a compare com os seus inferiores ou particulares e veja se ella é positiva e eíleciivamente paxtieipnvel e participada. e tenha para. os mesmos. inferiores a razão de \emph{genero} ou de \emph{especie}. --- Só depois d' estas: duas operações é que temos o verdadeiro universal \emph{reflexo}, --- a essencia, que seja, ao mesmo tempo, «no e \emph{nmultiplice}: uma quanto ao ser, que ella tem na \emph{intelligencia} € que representa. o que muitos individuos têm de \emph{commem}, ou de \emph{semelhante}, --- \emph{multiplice} quanto ao \emph{ser}. que ella tem na \emph{realidade} e que se multiplica e divide conforme a multiplicação e a divisão dos entes singulares e individuos. A \emph{unidade} \textbar{} é \textbar{} effeito da \emph{abstracção} ; a \emph{multiplicidade} é effeito da comparação.\par Entre o universal \emph{directo} e 0-\emph{reflexo} ha. varias differenças. Apontamos.\par as principaes. --- 1.º) 4 denominação de universal convem mais ao universal \emph{reflexo}, do que ao \emph{directo}. Por quanto, universal significa uma \emph{corsa} -commum- .. a \emph{outras}; ora esta relição de uma coisa a outras realiza-se de um modo \emph{explicito} e \emph{positivo} no universal \emph{reflexo}, que actualmente se refere aos individuos, ao passo que no universal \emph{directo} não se verifica senão de um modo \emph{implicito} e negativo, emquanto a essencia, considerada em si, não é limitada a nenhum individuo e tem a capacidade de ser attribuida, pela reflexão mental, a todos os individuos. existentes e possiveis. --- 2.º) O universal \emph{directo}, ainda que no seu proprio conceito inclua a \emph{unidade} (porque significa uma só essencia), comtudo, quanto A sua individuação real, não importa nem \emph{unidade} nem pluralidade, mas \emph{prescinde} de uma e de outra, porque, se importasse \emph{unidade}, não poderia applicar-se a muitos, e. se imporlasse \emph{pluralidade}, não poderia applicar-se a um só (\emph{De Ente et-Essentia}. €. 4j. Pelo contrario, \emph{o}- universal \emph{reflexo}, importando expressamente uma relação a varios entes, importa. \emph{unidade} e \emph{pluralidade}, ---- \emph{unidade} quanto \& si mesmo, e \emph{pluralidade} quanto ao termo da sua relação; porque é uma só forma, ou uma só essencia, que pode encontrar-se em muitos individuos,\pgn{154}pelo . menos: possiveis 3.º) O- universal reflexo-encerra o \emph{directo}, ao qual accrescenta.a relação 20s individuos. Conforme a, variedade d' esta. relação, o \emph{reflexo} divide-se nos cinco. \emph{predicaveis}, que: são o \emph{genero}, \emph{a especie}, \emph{a differença}, .O \emph{wroprio} € \emph{O accidente}. O genero e a \emph{especie} exprimem. perfeitamente o universal \emph{reflexo}... O universal \emph{directo}. não é propriamente o \emph{genero}, nem a \emph{especie}, mas é o \emph{principio} do \emph{genero} e da \emph{especie}, emquanto representa o ] elemento: \emph{objectivo}, que se.acha contido na noção de genero ou de \emph{especie}. (S. Thom. \emph{De universalibus}, 1. \emph{I}). Note-se. que o animal emquanto animal, e o homem emquanto. homem, é o universal \emph{directo}, e.o animal emquanto . genero, e o homem emquanto- especie, é o universal \emph{reflexo}. --- 4.º) Não é o universal \emph{veflexo} que se attribue, por. \emph{identi"ade}. a este ou áquelle ente particular (assim não dizemos: \emph{Pedro é}. \emph{especie}), --- mas é o universal \emph{directo}, isto é; é o elemento \emph{objectivo}, que se acha contido na noção generica, ou . especifica, nias. sem relação aos entes particulares tassim dizemos; \emph{Pedro} é homem). Esse elemento \emph{objectivo}. exprime, como dissemos, a essencia, á qual, considerada em-:si mesma, é indiferente para existir em um ou em muitos particulares. (\emph{De Ente} et \emph{Essentia}. c. 4). --- É claro que este modo de formar e perceber o universal é \emph{proprio} da intelligencia bumana. Deus percebe perfeitamente o universal, qu: nós formamos pela abstracção; mas Elle não precisa de praticar. esse. acto para o perceber.\par (D Além d' isso, todos attribuem, por \emph{identidade}, a essencia, representada pela idea universal-directa, aos entes existentes na realidade: assim todos dizem : Pedro é homem. Ora ninguem poderia attribuir tal essencia \textbar{} \&os entes reges, se ella effectivamente não se encontrasse nelles; assim ninguem poderia dizer que Pedro é homem, sb em Pedro não existissem os elementos \emph{especificos}. que constituem o homem. Por quanto, o verbo é, exprime, numa proposição, a \emph{identidade} não a simples semelhança) dos termos; e por isso a proposição --- Pedro a homem --- seria falsa, se a idea de \emph{Lo}- "e", não denotasse o mesmo ser \emph{especifico}, que se encontra em Pedro. Logo a \emph{essencia}, representada pela idea universal-directa, existe nos entes reaes.}. \emph{MILL} E

*: 147. Objectividade da idea \textbf{universal-directa}. --- A idea \emph{universal-directa} terá um valor \emph{objectivo}, se a. essencia, que ella representa, existir nos entes \emph{reaes}, embora não exista nelles pelo mesmo modo, por. que. existe. na intelligencia.. Ofa. a essencia,

' 3 4 . . . L representada pela idea \emph{wnmiversal-direcía}, existe, effectivamente; nos: entes \emph{reaes}, embora não exista nelles pelo mesmo modo, por que existe na intelligencia. -. a) A essencia, \emph{representada}. pela \emph{idea niversal-directa}, \emph{existe} --- nos \emph{entes} reaes. --- A \emph{essencia}, representada pela idea. \emph{universal}- \textasciicircum{} , \pg{154}  \emph{directa}, é o conjuncto. dos elementos \emph{especificos}, que faz com que o ente seja o que é e seja distincto de todos os entes das outras --- especies. Ora. todo Oo ente real é intrinsecamente constituido por uns elementos \emph{especificos}, pelos. quaes é o que é e.se distingue de todos os entes das outras especies; assim o \emph{homem} é con- “stituido nela \emph{animalidade} e. pela \emph{racionalidade}, e por estes dois elementos é homem e não é bruto, e distingue-se dos brutos, “das Plantas, dos mineraes. Logo a \emph{essencia}, representada pela idea \emph{universal-direcía}, existe .nos entes \emph{reaes}. (1) .- -

b) A essencia \emph{não existe nos entes reaes pelo} “mesmo \emph{modo}, \emph{por que existe} na \emph{intelligencia}. A essencia, representada pelo \emph{universal-directo}, exisle. --- na \emph{intelligencia}, que a percebe, --- e nos entes \emph{reaes}, que por ella são o que são; mas na \emph{intelligencia} existe por um modo \emph{intencional} e \emph{abstracto}, ou \emph{universal}, porque \pg{155}  existe numa \emph{idea} e só exprime os elementos \emph{especificos}; --- ao passo que nos entes \emph{reges} existe por um modo \emph{real} e \emph{concreto}, ou \emph{singular}, porque .constitue o ente real na sua especie e se acha cercada e determinada. pelos. caracteres \emph{individuaes}. Logo a \emph{essencia}; representada pelo «\emph{nivcersal:direcío}, não. existe "nos entes \emph{redes} pelo mesmo modo; por que existe na intelligencia. \footnote{Diz S. Thomaz: «No universal-directo, isto é, na: essencia abstracta, das notas ou caracteres \emph{indiniduaes}, devemos distinguir. duas coisas: a \emph{essencia}, \emph{ou natureza}, e x \emph{abstracção}, OU \emph{universalidade}, A essencia, que se percebe, se abstrahe e se universalisa, não existe senão nos entes singulares; mas a propria percepção, a abstracção, a universalização encontra-se na \emph{intelligencia}. É o que se verifica na. percepção sensitiva. A vista percebe a côr de um pomo. mas não lhe sente o saber. Se, pois, se perguntar onde está a cór. que se percebe sem o saber, respondemos que a côr, que se vê, não se envontra senão no pomo. Mas. se a côr se percebe sem sabor; deve a razão d' isso procurar-se na propria vista, na 'qual existe a semelhança da côr, € não do sibor..Do .mesmo:.modo,: :a:humanidade:z'"que- €- -percebid : não se encontra senio oeste ou naquelle individuo humano. Mas, sc ell se \emph{percebe} sem propriedades individuaes, e, por isso, sé se ábstralie € se reveste da forma de universalidade, tudo isso acontece a humanidade emquanto é percebida pela intelligencia, na qual existe, a semelhança da natureza da especie, e não a semelhança dos principios” ou caracteres: ind viduaes». (\emph{Sum}. \emph{Th}., p. 1, q. 85. a. 2 ad 2).\par Como.o assumpto acerca da \emph{objectivifade} das: ideas universaes é de importancia capital, accrescentamos as seguintes considerações. Todo o individuo, existente na realidade, é determinado. a.uma certa especie, e por: isso tem alguns elementos, que Ihe são \emph{proprios}, que não pertencem a outros individuos da mesma especie, e tem outros elementos, que lhe são rommwns com os individuos da mesma especie e o constituem nessa especie, e não noutra. Tanto os elementos \emph{proprios}. como OS'communs são \emph{objectivos}, Ou reaes. São reaes os elementos \emph{proprios}; aliás um individuo \textbar{} não poderia distinguir-se de outro individuo da mesma especie. São reaes os elementos \emph{communs}; aliás os entes reges reduzir-se-hiam a entes ?\emph{deaes}. --- São reaes os elementos \emph{proprios} e oS communs, e são \emph{distinctos}. Por quanto; se não fossem \emph{distinctos}, os elementos communs não poderiam encontrar-se sem os elementos \emph{proprios}, isto é, a essencia não poderia existir senão num certo e determinado individuo, e assim haveria um só individuo subsistente em cada essencia ou especie: o que repugna á experiencia e a razão. --- Sendo \emph{distinctos}, os elementos \emph{proprios} € OS connus, se não podem separar-se \emph{realmente} uns dos outros, podem comtudo separar-se \emph{intencional} ou \emph{ideulmente}, emquanto podemos considerar os elementos \emph{communs}, sem considerarmos os \emph{proprios}. E, quando consideramos os elementos rommuns, e prescindimos dos \emph{proprios}, esses elementos communs não deixam de ser \emph{objectivos}, ou de existir na \emph{realidade}, como. não deixa de ser \emph{objectiva} e real a. \emph{cór} de uma pera, só porque a consideramos sem o \emph{saber}. \emph{A abstracção}, pela. qual consideramos uma coisa sem outra, é um acto da intelligencia, mas é\pgn{156}um: acto que não altera, não transforma, não destroe a objectiva realidade\par das coisas. Estas continuem a ser o que eram, --- Nem se diga que tudo o que é \emph{real}. \emph{é individuo}. Sé é \emph{individuo} tudo o que é real \emph{quanto} \& \emph{existencia}, porque só os individuos existem, não é \emph{individuo} tudo o que é real \emph{quanto} d essencia, porque esta é real, e, comtudo, n não é, \emph{por}. \emph{si}, \emph{individua}.\par Objectam: Nenhuma coisa pode tornar-se o que não é. Ora todo o objecto, existente na realidade, é \emph{individuo}. Logo. nenhum d'esses objectos pode tornar-se- universal. :\par Resposta. A affirmação --- nenhuma coisa pode tornar-se o que nio é ---\par não: pode aeceitar-se incondicionalmente ; porque ha muitas coisas que se tornam.o: que não são. “assim 0 marmore não é estatua, : e comtudo se torna estatua. Só. então uma coisa não se pode tornar o que não é, quando lhe falth a \emph{capacidade} para essa. mudança. Ora o objecto, existente na realidade, se é \emph{individuo}, ou particular, em \emph{acto}, \emph{é universal} em \emph{potencia}, porque tem "sa \emph{capacidade} de ser percebido de um modo \emph{abstracto}, na sua \emph{essencia}, livre das «condições proprias de uma, existencia \emph{concreta}. --- A objecção. suppõe que, para a formação do universal, seja necessario separar ou remover \emph{physica} oU. \emph{realmente} da essencia.as condições materiaes ou individuaes. Mas a supposic\&o é falsa. Por. quanto, como dissemos, essa separação, ou remoção, não é \emph{physica}; ou real, mas-é \emph{intencional}, ou \emph{ideal}. Escreve s. Thomaz: « Ista abstractio non est.intelligenda secundum. rem, sed secundum rationem » (De \emph{pot}, \emph{ànimae}, c. 6). --- Nem.se diga que o objecto, se não se muda \emph{physicamente}, se. muda \emph{idealmente}. Por quanto, se o objecto. quando é percebido- abstractamente, se muda.no seu modo \emph{de ser}, ou de \emph{existir}. não se muda: e» si\textasciicircum{} mesmo. ---- Muda-se- o- objecto no.seu \emph{modo} de \emph{existir}, porque começa a-existir na nossa. intelligencia de um modo \emph{abstracto} e immaterial, quando, na realidade, existe de üm modo \emph{concreto} e \emph{material}. E nada mais natural. Todo o'objecto; que é recebido num sujeito. segue o modo de existir, proprio: do mesmo sujeito. Ora, sendo a intelligencia o sujeito em que 0 objecto é recebido por meio' da \emph{idea}, é claro. que o mesmo objecto deve despir, por assim dizer. o medo de existir, que. tinha na realidade, para « revestir.o modo de existir. que é proprio da intelligencia. O modo de exis- 'fir da intelligencia é a plena isenção ou independencia da materia e das condições, que a acompanham. Logo tambem o objecto, recebido na intelligencia, deve existir nella sem a materia e sem as condições materiaes. O objecto sem \& materia e-sem.as condições materiaes é a \emph{essencia abstracta}, é O \emph{universal directo}. Diz S. Thomaz: « Similitudo rei recepitur in P n" q. 85, a. 5 ad 3). --- Mas o \emph{modo de ser}, ou de \emph{existir}, do objecto, nào " o secundum modum intellectus, et non secundum modum rei» (\emph{Sum}. \emph{Th}.,\par \emph{proprio objecto}. Se o objecto se muda no seu modo de ser, não se muda em\pgn{157}si mesmo: permanece sempre o que era.. A essencia, que existe na intelligencia por um. modo immaterial ou espiritual, é a mesma que existe na realidade por um modo material; porque os elementos, que a constituem, "são sempre os mesmos num e noutro modo de ser. ou de existir. . Ipsistem : Se a intelligencia formasse o \emph{universal}, ella communicaria ao objecto a \emph{necessidade} e a eternidade; visto que o universal é, por \$1, necessario e eterno. Ora isto é impossivel ; ; porque todo o objecto creado é \emph{contingente} e mutavel. Logo a intelligencia não forma o \emph{universal}.\par Resposta. A intelligencia, quando forma o universal (\emph{directo}), não com. munica coisa alguma. limita-se. a remover, a afastar do.objecto. os. caracteres, ' jndividuaes e concretos, para fazer sóbresahir e pôr em vista exclusivamente a \emph{essencia}. Se, pois; \& essencia, quando \emph{abstracta}, reveste os caracteres. de. .\emph{necessidade} e de eternidade, estas notas não derivam do acto abstractivo - para a essencia, porque esse acto.não as tem e ninguem dá o que não tem, --- mas brotam do proprio fundo da. essencia, e.esta não é [3 .producto,.. mas apenas o \emph{objecto} da iotelligencia. --- Dizemos que essas duas notas brotam do proprio fundo da essencia. Por quanto, se as coisas creadas são \emph{contingentes} e \emph{mutaveis} quanto á existencia, .porque .podiam e podem deixar de existir, são \emph{necessarias} e \emph{eternas} quanto á \emph{ia}. São necessarias; porque as \emph{essencias} não podem deixar de corresponder e de ser semelhantes aos archetypos da intelligencia divina; são eternas, emquanto as essencias não süv coarctadas entre um determinado espaço de tempo, mas, por si, prescindem do tempo, como tambem do logar. --- Todavia essa \emph{necessidade} e \emph{eternidade}, que se attribue ás creaturas, é infinitamente diversa da \emph{eternidade} e necessidade, que é propria de \textsc{Deus}. Por quanto, á \emph{necessidade}, que se attribue ás creaturas, é limitada a \emph{essencia} d'ellas, e não se extende a . \emph{existencia}. porque as creaturas não existem necessariamente, \&O0 passo que a necessidade, que é propria de Deus, abrange a \emph{essencia} e a existencia, porque Deus necessariamente existe, alem de -ser \emph{necessariamente o} que é. Como tambem, ao passo que a \emph{eternidade}, que se attribue \& essencia das coisas creadas, é apenas negativa, emquanto .\emph{abstrahe} do tempo, a eternidade, , que é propria de Deus, é \emph{posifiea}. emquanto, com \emph{a} sua indivisivel' presença, abrange todo o tempo (Sum. \emph{Th}., II, q. 93, a. 4 ad 4; \emph{Quodl}. IT, a. 1).}.

%%>>> t148

\pg{152} \pg{153} \pg{156} \tese{148}{Objectividade da idea universalreflexa.}
O objecto da idea: \emph{universal-reflexa}, embora \emph{exista unicamente} na \emph{intelligencia},. comtudo tem o seu \emph{fundamento} na realidade.

a) O.objecto--da-idea. \emph{universal-refleza} . \emph{existe unicamente} na \emph{intelligencia}. ---Por quanto, essa idea representa uma. essencia, que, sendo una, comtudo É comam. a muitos particulares, e \pg{157}  a cada um d'.elles pode ser attribuida. Ora tal essencia \emph{existe unicamente} na \emph{intelligencim}; visto que, na. realidade. não ba coisa nenhuma que seja, ao mesmo tempo, una e commum a --- varios seres, pois todo ente real é \emph{singular} e. \emph{individuo}, ea essencia de um individuo é realmente distincta da essencia dos outros judividuos. Logo o objecto da idea \emph{universal-reflexa}. \emph{existe unicamente} na \emph{intelligencia} \footnote{Além d'isso, se o universal \emph{reflexo}, isto é, o \emph{conteudo} da idea uni \emph{versal-reflexa}, existisse nos entes singulares ou individuos, poderia ser-lhes attribuido, como se lhes attribue o \emph{universal-directo}. Mas não é assim. Ao passo que podemos dizer : \emph{Pedro é homem}, não podemos dizer: \emph{Pedro é especie}. Qual é a razão d' esta differença? A razão é que, se essa essencia, denotada pelo termo homem, se encontra realmente em Pedro, não se encontra\pgn{158} nelle a essencia, denotada. pelo termo \emph{especie}. Devemos, pois, dizer\par que o \emph{universalreficxo} existe: \emph{unicamente} na intelligenicia,.- 500. 7\par "Para, maior clareza lembramos que a essencia pode considerar-se em \emph{si mesma}, nas coisas singulares, e na \emph{intelligencia}. --- Considerada em \emph{si} mesma, a essencia importa unicamente os \emph{elementos}, que a \emph{constituem} e que. fazem com que seja aquella. que é e não outra, e por isso \emph{prescinde} das condições concretas, de que se acha determinada va realidade, e tem os attributos, que lhe pertencem pela sua definição; tal é a humanidade; considerada nos seus elementos constitutivos de animalidade e de \emph{racionalidode}. Neste caso, a \emph{essencia} é universal-directo, ou \emph{metaphysico}, e pertence á ordem \emph{metaphysica}. --- Considerada nas coisas singulares, 2 \emph{essencia} està' cercada de propriedades concretas e individuaes (que, nas substancias compostas, são as condições da \emph{materia} individual), e por isso tem attributos, que lhe convém emquanto e porque existe nos seres. singulares; tal é a humanidade, que existe em. Pedro, Paulo, etc. Neste caso, a essencia pertence a ordem \emph{physica}. --- Considerada .na \emph{intelligencia}, \emph{a essência} importa uma coisa, que, sendo una, é commum a muitos particulares, e por isso tem um - ser ideal, e tem os attributos, que lhe convém por: causa d' esse ser, que recebeu da intelligencia; Lal é a humanidade, considerada como especie. Neste caso, a \emph{essencia} é uma coisa universal \emph{reflexa}, ou \emph{logica}, e. pertence á ordem \emph{logica}. ---- Como se vé, tanto a \emph{singularidade}, como a universalidade, ou \emph{communidade}, são coisas \emph{exírinsécas} e \emph{accidentaes} para a essencia, considerada em si mesma, isto é, nos seus elementos constitutivos. (Cf. \emph{De Ente et Essentia}, c. Vl). --- Ora' a essencia, que se chama universal, isto €. que se \emph{refere} actualmente aos particulares, --- não é a \emph{essencia} considerada em \emph{si} mesma, porque, se assim fosse, ella havia de ser sempre, e em toda a parle, universal, e, por isso, \emph{unica} e \emph{multiplice}, --- nem é u \emph{essencia}, considerada nos seres \emph{síngufares}, porque esta não é especie, nem é, ao mesmo tempo, unica e \emph{multiplice}, --- mas é a essencia, considerada como existe na \emph{intelligencia}, porque, só em relação a esse \emph{ser ideal}, a essencia pode ter \emph{ii} razão de especie, e ser, ao mesmo tempo, \emph{unica}. e \emph{nudtiplice} (Cfr. \emph{De Anima}, L, If, 1, 12),}.

s

b).O-objecto da, idea \emph{auniversal-reflexa} tem. o seu \emph{fundamento} na \emph{realidade}. --- O \emph{objecto} da idea \emph{universal-reflexa} terà o seu \emph{fundamento} na \emph{realidade}, se as coisas \emph{reaes} subministrarem a intelligencia. un. .elemento \emph{objectivo}, que, primeiramente depurado das condições materiaes e depois comparado com os indi- 'viduos; se conceba como ema \emph{forma una e commum}. a muitos \emph{Seres}. Ora as coisas reaes apresentam á intelligencia. esse \emph{ele}

" \emph{mente objectivo}, que é a \emph{essencia} das coisas; porque. a \emph{essencia} das coisas,. que é o. objecto do \emph{universal}- \emph{reflexo}, embora, \emph{quanto}

uo seu \emph{mnodo de'ser}, exista exclusivamente na intelligencia, comtudo, \emph{quanto aos seus elementos constitutivos}, nos quaes a propria intelligencia se funda para conceber a mesma essencia como una forma «na e comumwm, existe effectivamente na realidade \pg{159} ; assim a \emph{humanidade}, embora, considerada como \emph{especie}, só exista na intelligencia, comtudo, considerada nos seus: elementos de \emph{animulidade} e de \emph{racionalidade}, existe realmente \textasciicircum{} nos individuos bhumanos..Logo..o objecto. da.. idea universal. \emph{reflexa} tem o seu \emph{fundamento} na \emph{realidade} 0 UC\footnote{A conclusão parece-nos evidente, O universal \emph{reflexo} é O proprio o universal \emph{directo}, emquanto, pela reflexão, se considera como abstracto e como commum a muitos. seres. Ora o universal \emph{directo} exprime e denota uma coisa real e \emph{objectiva}, qual é a \emph{essencia} das coisas, considerada nos . seus \emph{elementos constitutivos}, Logo o universal \emph{reflexo} deve suppor e suppõe o universal \emph{directo}, e por isso tem o seu \emph{fundamento} na \emph{realidade}. --- É \emph{verdade} que o. universal \emph{directo} representa a essencia abstracta, não emquanto é abstracta, mas emquanto é essencia, ao passo que o universal \emph{refleco} representa a essencia \emph{abstracta}, emquanto \emph{abstracta} e comparada com os individuos; mas tambem é verdade que essa abstracção e \emph{comparação}, se dão á essencia um novo modo de ser na intelligencia, não lhe \textbf{tiram} o modo de ser, que tinha na renlidade, isto é; não transformam em \emph{ideaes} ou \emph{logicos} os elementos reaes ou \emph{objectivos}, que .a. constituem,-e que-permanecen:- senipre o que eram, e formam o \emph{substractum}, a base do universal \emph{reflexo}. --- Nem podia ser de outro modo. 4 nossa intelligencia, quando \textbf{forma} o universal \emph{reflexo}, não faz senão \emph{referir} uma essencia \emph{abstracta} aos individuos reaes, em que ella se encontra ou pode encontrar-se, e por isso a considera como uma coisa commum a muitos entes. Ora a intelligencia não poderia considerar uma coisa como conum a muitos entes, existentes na realidade, se nestes não existisse \emph{realmente} essa coisa, commum a todos, que não é Senão a \emph{essencia}. --- Portanto, --- se o universal \emph{directo} tem o seu fundamento \emph{immediato} ou \emph{proximo} na realidade, porque. a. essencia, que se considera em \emph{si mesma}, existe na realidade (embora na realidade não exista pelo mesmo modo por que existe na intelligencia), --- o universal \emph{reflexo} tem na realidade o seu fundamento mediato ou \emph{remoto}, porque a \textbf{essencia}, que se considera como \emph{abstracta} e commum a muitos entes, é uma coisa real e \emph{objectiva}. (Cfr. S. Thomaz in I \emph{Sent}., Dist. 19, q. 5, a. 1). --- Porisso,..-...--- quando dizemos que o universal € um ente \emph{logico}, ou \textbf{ente} de \emph{razão}, --- não nos referimos a essencia, que se abstrahe dos individuos pela consideração \emph{directa} da inteligencia e que é real Ou \emph{objectiva}, anterior a toda a consideração, --- nem nos referimos ao unzversal, considerado no seu \emph{fundamento} (pois este tambem, como vimos. é real), --- mas referimo-nos exclusivamente 4\emph{o} proprio caracter da \emph{universalidade}, emquanto se une \emph{accidentalmente} é essencia. Em poucas palavras: no universal devemos considerar duas coisas: a \emph{essencia} e a \emph{forma} de \emph{universalidade}, que a essencia reveste. A \emph{essencia} é ente \emph{real}; \emph{a forma de universalidade} é ente \emph{logico}, ou de \emph{razão}.\par Onde existem os \emph{universaes}?.. Eis uma das mais antigas e ruidosas questões, que se tem agitado no campo da sciencia. As muitas e diversas soluções, apresentadas pelos philosophos, podem reduzir-se a éres. A \emph{primeira} é a dos que affirmam que os \emph{universaes} não existem, de modo algum, Das coisas singulares e materiaes, existentes nu realidade. A segunda € á\pgn{160}dos que, pelo contrario, sustentam. que os universaes existem, emquanto taes, na realidade: A \emph{terceiro} é a. dos que, seguindo um- caminho medio, ensinam que os \emph{universaes} existem --- em \emph{potencia} ou \emph{fundamentalmente} na \emph{realidade}, --- em \emph{acto} ou \emph{formalmente} na \emph{intelligencia}. --- Tendo sulficientemente desenvolvido esta terceira. opinião, que é conhecida pelo nome de realismo \emph{moderudo} e que foi defendida por \emph{Aristoteles}. e depois por \emph{S}. \emph{Boecio}, pelo B. Alberto \emph{Magno}, por \emph{S}. Bonaventura, e, sobretudo, pelo Angelico "\emph{S}. \emph{Thomas} de Aquino, resta-nos examinar e criticar as outras duas opiniões.\par 3) Primeira opinião. --- A \emph{primeira} opinião, a qual, como dissemos, affirma - que xistem. de \emph{modo} algum, na realidade, [oi sustentada pelos sequa; és do nominalismo e do conceptualisio.\par \emph{i} a) O nominalismo: diz assim: Os chamados \emph{eniversaes} não eram, no .. principio, senão nomes \emph{singulares}, que denotavam \emph{individuos}. Esses nomes \emph{singulares} tiveram, em seguida, um sentido mais lato e extenso, e foram empregados para: denotar muitos outros individuos, semelhantes aos conhe- . cidos; d' ahi a \emph{universalidade} dos nomes. Portanto, os universaes existem exclusivamente nas palavras, não passam de simples \emph{nomes}, aos quaes não corresponde nenhuma realidade da natureza e nenhum conceito da intelligencia.: --- Defenderam o \emph{nominalismo} --- Epicuro (341-270 a. Ch.) e os \emph{sensualistas} na antiguidade, --- Rosceliino (j 1089) e Guilherme Occam (1280- 1343) na edade media, --- Thomaz Hobbes (1588-1619), David Hume, Smith, . Condillac, Dugald-Stewart, 'Taine, Ribot, Wundl, os modernos \emph{materialistas} *\emph{é.positivistas}, que só admittem imagens sensiveis e concretas.\par O \emph{nominalismo} € falso pelas seguintes razões. Os nomes 'são signaes externos das ideas. Logo, se existem nomes'universaes, devem existir tam- "bem ideas \emph{universaes}; aliás deveriarnos dizer que ha signaes, que não têm significação, e que ha palavras, que não dizem nada. --- Além d'isso, perguntamos aos nominalistas:.por que motivo, vendo um individuo da mesma especie, denotamo-lo com o mesmo nome?. Respondem: porque é semelhante aos outros. Optimamente. Mas como poderiamos nós conhecer aquella semelhança, se não possui\$semos já a idea de uma especie, ou essencia, que, \&bstrahindo dos caracteres individuaes, exprimisse exclusivamente o que é commum a varios individuos ? Logo a imposição do mesmo nome a individuos semelhantes presuppõe uma idea \emph{universal}. --- Finalmente, o \emph{noninaliso} destroe todas as \emph{sciencias}, porque estas têm por objecto as ideas universaes, e não \emph{as singulares}. t\par 5b) O \emph{concepiualismo} sustenta --- que os universaes existem nos conceitos, mas \emph{de modo algum} nx realidade, --- e que por isso os wniuersaes São meras \emph{ideas}, ás quaes nenhuma realidade \emph{objectiva} corresponde. --- Defen-}

' 149. Fundamento supremo da objectividade das ideas uni- . versaes. --- Às ideas \emph{universaes}, se têm o seu fundamento proximo mas coisas creadas, têm o seu fundamento \emph{supremo} na reali- : dade da. propria. Essencia infinita de Deus. Por quanto, todas as coisas creadas são outras tantas copias dos archelypos, que a intelligencia do Artifice eterno forma e que importam alguns dos infinitos modos, por que a Essencia divina pode ser representada \pg{161}  na creação. Assim, pode dizer-se que as coisas creadas - \$4 \emph{Ps} são impressões visiveis da invisivel sciencia divina. Medidas "EM pela sciencia divina, as coisas creadas medem, pela sua vez, a

deram esta opinião ---' os \emph{Sfoicos} na antiguidade; --- Pedro Abeilardo (1079- : 1149) na edade media, --- Descartes e Kant na edade moderna. M \emph{C}

O \emph{conceptualismo} é falso. Com effeito, assim como os nomes são signaes --- --- dos conceitos, nssim. tambem :os “conceitos - são essencialmente signaes das \emph{coisas}. Logo, se existem conceitos universaes, deve existir uma realidade objectiva, que de algum modo lhes corresponda. --- Além d' isso, todas as sciencias, á excepção da Logica, têm por objecto ás coisas reaes, c não as \emph{ideaes}. Mas, se fosse verdadeiro o \emph{conceptualismo}, ns sciencias occupar-se-hiam de phantasmas vãos ; d'ahi o \emph{scepticismo}. --- Finalmente, os \emph{conceptualistas} admitlem que nós percebemos a semelhança existente entre varios individuos .: e que, em vista d'essa semelhança, designamos com o mesmo conceito : AS muitos entes singulares. Ora não poderiamos afirmar uma semelhança entre 7 muitos entes singulares, se nestes não” existisse realmente um elemento ..: 1 \emph{Maus} commnum a. todos. Este elemento commzwm a muitos entes é o «\emph{universal}. Logo \textbar{}\emph{o} O universal existe de \emph{algum} modo na realidade \emph{objectiva}.

b) Segunda opinião:. ---.. A “\emph{segunda}. opinião affirma, como vimos, ue: os universaes existem, emquanto universaes, nas coisas reaes e objectivas; : A “de modo que fora da nossa intelligencia existe o \emph{genero}, a especie, etc. É O --- 3 \emph{NN} realismo \emph{exagerado}. --- Mas os sequazes d'essa opinião não explicam todos . do mesmo modo a pretendida existencia do universal, como \emph{tal}, na realidade, . \emph{MB} Alguns, como os \emph{Platonicos}, disseram que 0\$ \emph{wuntversaes}, como \emph{taes}, subsistem D em \emph{si mesmos}, tira da intelligencia. Outros, como Guilherme de Champeaux \textbar{} = (1070-1120) e David de Dinant, embora não admittissem 2a subsistencia dos universaes em. si mesmos, como formas separadas dos entes singulares, comtudo disseram que elles existem \emph{actualmente} na realidade, porque a essencia das coisas visiveis é \emph{realmente} unica e \emph{numericamente identica} em iodas, de modo que uma coisa não é. ou não parece distincta de outra senão pelos -. HH --- accidentes externos. Os transcendentaes modernos levaram \textbf{essa} opinião ao . mais crasso pantheismo, dizendo que existe uma unica substancia, da qual \emph{O}. ás coisas, que nos.cercam,. são. apenas. modificações. . Outros, sequazes de \emph{Scoto}, ensinaram que a essencia, considerada em si mesma, e fóra de toda 25.1 \emph{Man a} reflexão da intelligencia, tem x capacidade de existir, indifferentemente, : dii neste ou naquelle individuo, e é por si \emph{universal}. \emph{t}.

O \emph{realismo exaggerado} é ialso. Porquanto, se os \emph{universaes} existissem, como \emph{taes}, fóra da intelligencia, deverjamos dizer que, na realidade e inde- u pendentemenle de todo a acto intellectual], existem os \emph{generos} e aS \emph{especies}, --- 3 --- como \emph{íaes}; de modo que poderiamos dizer que um individuo humano, por ex., \emph{Pedro}, \emph{é especie}. Ora ninguem pode admittir tal absurdo. --- Alem d'isso, --- os universaes, que são o objecto da nossa intelligencia e de que se occupa \emph{bv} a sciencia, são immutaveis, incorruptiveis, necessarios, eternos, em relação \emph{B P}\textasciicircum{} \& essencia \emph{abstracta}, que elles exprimem, Ora tudo o que existe ou pode u existir na realidade é mutavel. corruptivel, contingente, temporal. Logo os \emph{Uuniversaes} não existem, como \emph{Íaes}, na realidade.

Pelo que deixamos dito podemos apreciar as diversas formas d' esta opinião. --- Quanto aos \emph{Plafonicos}, diremos com \emph{Aristoteles} que elles par-

Yol, X. 1a \pg{162}  Sciencia humana. Essas coisas, actuando na. nossa intelligencia, deixam ahi impressa uma imagem, que representa a \emph{essencia}, abstrahida das condições materiaes e por isso \emph{universal}: d'onde as ideas universaes. Logo as ideas universaes têm o seu \textbf{fundamento} supremo na realidade da Essencia infinita de \textbf{Deus} (1).

tem de \emph{um} falso principio, suppondo que, quando \emph{a} intelligenein \textbf{separa}, a essencia das notas individuaes, tal separação é real D'essa supposta \textbf{separação} real é claro que não pode derivar senão a subsistencia: dos universaes, como taes, fora da intellirencia. Ora o. principio estabelecido é falso; porque, como explicámos, essa separação não é \emph{real}, é apenas \emph{intencional} (Cfr. Arist, \emph{Post}. \emph{Analyt}., Db. ll, lect. 2; \emph{Metaph}., lib. VII, lect. 14). --- Quanto a \emph{Guilherme de Champeaux}, diremos com S. Thomaz que ás substancias singulares e individuas são diversas o disUncias entre si, não só porque as propriedades accidentaes de uma são diversas e distinctas das propriedades \&ccidentaes de outra, mas tambem: e sobretudo, porque a materia e a \textbf{forma} de uma não são a materia e a forma de outra. A forma, embora possa, por S), ser \emph{especificumente} una em varios seres, comtudo é numericamente multiplice, porque -se- torna" \emph{individua} em -cada ente. singular, e o principio d essa individuação é, como veremos, \emph{a} materia. --- Quanto aos \emph{Scotistas}, diremos que elles fundam . o seu raciocinio num equivoco. Não negamos que a essencia, considerada em \emph{si mesma}, e como existe na intelligencia, \emph{PE} seja dotada de uma certa \emph{communidade} ou. ?\emph{nultiplicidade} negativa. ou de uma indifferenca para existir neste ou naquelle individuo, e que, sob \textbf{este} aspecto, possa chamar-se \emph{waiversal}. Mas d'ahi não se segue que essa essencia exista na realidade. Por quanto, o-essencia --- se tem essa \emph{communidade} ou \emph{multiplicidade negativa}, quando se considera em si mesma e na intelli-. gencia, --- não a tem, quando se considera na realidade; pois a essencia não existê na realidade senão porque existe nos individuos singulares, e o que existe num individuo não pode pertencer a outro. Nem todas as propriedades, que convém a um ente, emquanto existe na intelligencia, convém ao mesmo, emquanto existe na realidade. .

(D A intelligencia humana possue. ideias \emph{universaes}. Essas ideas não se conlemplam na \emph{Essencia divina}, que nos é invisivel, mas são recolhidas das coisas creadas. Embora recolhidas das coisas creadas, as ideas universaes representam as essencias, não concretas \emph{e} singulares, como existem na realidade, mas abstractas e universaes. --- Transportadas numa ordem ideal, as essencias revestem os attributos de necessidade, immutabilidade, eternidade, etc, que, se.não se identificam com os correspondentes attributos de Deus, comtudo são \emph{reflexos} d' estes. Immutaveis, necessarias, eternas na intelligencia divina, as essencias das coisas creadas devem participar d'esses dotes, quando actuadas na cresção. Sendo immutavel, necessario, eterno o sello, tambem immutavel, eterna, necessaria é proporcionalmente a impres- S\&o. --- E, se a impressão foi feita na materia mutavel, contingente, temporal, basta uma simples abstracção da ipntelligencia para separar da materia á mesma impressão e fazer, d' este modo, respiandecer as nobilissimas propriedades, de que as essencias creadas são dotadas na iptelligencia divina, --- Por isso, \& sciencia humana participa, no sentido explicado, da

v D

%%>>> cap-8

\chapter*{CAPITULO SEGUNDO}
\addcontentsline{toc}{chapter}{CAPITULO SEGUNDO}

\pg{158} \pg{160} \pg{163} \textbf{MEIOS PARA} \emph{A} \textbf{ACQUISIÇÃO DA VERDADE} \textsc{SUMMARIO}; --- Intelligencia e razão, --- Consciencia e memoria intellectual. ---

Sentidos internos e externos..--- Testemunho, suas especies e transmissão:

AmnrrGO [.

\emph{Intelligencia e} razão

%%>>> t150

\tese{150}{Classificação d'esses meios.}
Os \emph{meios} para a acquisição da verdade são: a \emph{intelligencia}, a \emph{razão}, a consciencia, a memoria \emph{intellectual}, os \emph{sentidos internos} e \emph{externos}, e o \emph{testemunho}.--- Os \emph{primeiros} são \emph{inirinsecos}; o \emph{ultimo é estrinseco}. --- Neste artigo tratamos da \emph{intelligencia} e da \emph{razão}; nos seguintes, dos outros meios \footnote{A \emph{intelligencia} é propriamente á unica faculdade, que o homem possue para adquirir a \emph{verdade}; porque só ella, por ser espiritual, conhece . a \emph{relacito} entre \emph{o} predicado e o sujeito. Os sieios, pois, que enumeramos, --- ou são essa mesma faculdade intellectual. que toma diversos nomes (\emph{iniellijencia}, \emph{razão}, \emph{consciencia}, \emph{memoria}), conforme a diversidade dos actos, que exerce, ou das verdades, que adquire, --- ou são \emph{instrumentos} (\emph{sentidos} e testemunho), de que ella se serve para alcançar o seu fim. --- A classificação Jos \emph{smeios} para a acquisição da verdade é \emph{completa}. Cum effeito, todas as verdades, que podemos conhecer, --- ou são descobertas por nós --- ou nos são ensinadas por outros. --- Se são descobertas por nós --- ou representam : intidades ióminaferiaes, Ou exprimem entidades \emph{materiaes}. --- Se representam jntidades \emph{imunateriaes}, --- ou são ideas \emph{universaes} e juizos dotados de evilencia \emph{immediata}, e são percebidas pela \emph{inteligencia}; --- ou são juízos dotados le evidencia mediata, e são percebidas pela \emph{razão}; --- ou são as operações spirituaes, existentes \emph{actualmente} no espirito, e são percebidas pela conciencia; --- Ou, SãO as \emph{proprias} percepções \emph{intellectuaes passadas}, e são per- /ebidas pela memoria \emph{intellectual}. --- Se exprimem entidades materiaes, --- ou ie referem aos objectos \emph{externos}, e a iotelligeneia percebe-as por meio dos entidos \emph{externos}, --- ou se referem aos objectos \emph{infernos}, e a intelligencia jercebe-as por mejo dos sentidos \emph{internos}. --- Se as verdades nos são ensi- Xidas por outros, temos o \emph{testemunho}, que é divino ou humano. (Cf. Sum. \emph{Uh}., p. 1, q. 79, aa. 7, 8; \emph{De Verit}, q. 10, a. 8; c. \emph{Gent}, l. 1, c. 81).\par A \emph{Logica} e a \emph{Metaphysica}. occupam-se das nossas faculdades \emph{perceptivas}, nus \emph{sob diverso aspecto}. \emph{A Metaphysica} estuda-as emquanto têm uma certa Dablureza \emph{e} um correspondente \emph{modo} de operar: a \emph{Logica} considera-as mquanto são \emph{meios legitimos} para a consecução \emph{da verdade}.\par Às nossas faculdades \emph{perceptivas} serão meios legitimos para a consecução la verdade, se forem, pela sua propria natureza, verazes com relação ao seu\pgn{164}\emph{objecto proprio}. Que sejam ellas naturalmente dotadas d'essa \emph{veracidade}, é uma das verdades evidentes, que não precisam nem podem ser demonstradas \emph{sem petição de principio}. Por quanto, toda a demonstração, exigindo o o uso e.o exercicio das nossas faculdades perceptivas, suppõe como certa . e indiseutivel a sinceridade ou legitimidade d'essas faculdades. Esta ver- " dade, pois, pode ser mais \emph{explicada}, do que \emph{demonstrada}. A explicação é a seguinte. "Podas as faculdades perceptivas estão naturalmente ordenadas para conseguir a verdade, cada uma na-sua esphera e com relação ao proprio \emph{objecto}.. Ora não poderiam ellas estsr. naturalmente ordenadas para conseguir a verdade, se não fossem por si verazes; aliás o Auctor da natureza teria dotado o homem de instrumentos essencialmente, incapazes de conseguir o proprio fim. --- Portanto, se na percepção das nossas faculdades pode haver-e ha engano, este não depende-de-um- vicio-innate- ou- de-nmanatural ineptidão das mesmas, mas unicamente de circumstanocias accidentaes, como veremos (C. \emph{Gen}£, IIl, 107).}.

%%>>> t151

\pg{164} \tese{151}{Intelligencia.}
\emph{Jntelligencia}, considerada como meio especial para a consecução da verdade, é a faculdade, pela qual \textasciicircum{} apprebendemos a \emph{essencia} das coisas é os \emph{primeiros principios} “da demonstração. ----.Os. \emph{primeiros}. \emph{principios}. da demonstração ão \emph{verdades}, \emph{dotadas de epidenciá immediata} e \emph{percebidas pela} . \emph{Simples inspecção} ou \emph{analjjse dos termos}; como, por ex., o todo

é maior que uma sua parte ('). \textbar{}

189. Objecto da intelligencia. --- O \emph{objecto}, pois, da intelli-

gencia é constituido pela \emph{essencia} das coisas e pelos \emph{primeiros principios} da demonstração. --- Os principaes \emph{principios ord}.: I Pr. de \emph{identidade} --- um ente é identico a si mesmo. II. Pr. de \emph{contradicção} --- um ente não pode ser e deixar de

ser ao mesmo tempo e sob o mesmo aspecto. III Pr. de \emph{conveniencia} --- dois entes, cada um dos quaes é

egual a um terceiro, são eguaes entre si. AV. Pr. dé desconveniencia --- dois entes, dos quaes. um só é

\emph{i} egual x um-terceiro; não são eguaes entre si. V. Pr. de \emph{substancialidude} --- O accidente suppoe a substancia. VI. Pr. de \emph{causelidade} = o etfeito- \emph{exige} ou tem uma Causa.\footnote{Diz S. Thomaz:: « Nomen \emph{Z£níellectus} dupliciter accipi potest. \emph{Uno modo} secundum quod se habet ad hoc tantum a quo primo nomen impositum fuit (\emph{intelligere} enim est quasi \emph{íntus legere}); et sic dicimur proprie intelligere, cum apprehendimus \emph{qwidditatem} rerum, vel intelligimus illa quae statim nota sunt iatellectui, notis rerum quidditatibus, sicut sunt prima \emph{principia}, quae cognoscimus cum Lerminos cognoscimus... \emph{Alo} modo potest accipi intellectus communiter, secundum quod ad omues operationes se extendit, et Sic comprehendil \emph{opinionem} et \emph{ratiocinationem} » (\emph{De} Ver., q. 1, a. 12).}

- Ts ---

CAP. "n --- MEIOS PARA A ACQUISIÇÃO: :DA VERDADE 165 ---

VII. Pr. de \emph{razão sufficiente} --- todo o ente tem uma razão suf- \emph{Las}

' ficiente \footnote{\pgn{166}des \emph{contingentes}. Esta opinião não pode admittir-se. Porquanto, se é certo que em todas as verdades \emph{contingentes} se acha incluido o principio de \emph{razão sufficiente}, é certo tambem 'que este mesmo principio se funda no principio de \emph{rontradícção} : porque, se podesse dar-se uma coiea sem \emph{raczo sufficiente}, uma e a mesma coisa \emph{teria} e \emph{não feria} razão de ser: o que é ? contrario ao principio de contradicção.\par (B a sapientissima doutrina ,de S. Thomaz. Existe um nexo essencial entre as coisas; Uma está,-pela-propria natureza; ordenada para outra; e, em virtude. .d”.essa .ordem,.uma..opera-sempre do-mesmo- modo com relação á outra e alcança sempre o seu escopo. Este nexo, essencial e necessario, existe lambem entre a faculdade e o seu objecto \emph{proprio}. --- Objecto \emph{proprio} de uma faculdade é o que \emph{primeiramente} e por \emph{si mesmo} é percebido pela mesma faculdade. Diz S. Thomaz: «1d quod est \emph{primo} et \emph{per se} cognitum à virtute cognoscitiva est \emph{proprium} objectum ejus» (\emph{Sum}. \emph{Th}. p. Y. q. 85, a. S. --- A faculdade tem toda a sua razão de ser na existencia e na indole . do seu objecto \emph{proprio}, e, estando natural e necessariamente ordenada para elle, natural e necessariamente o alcança, ou não se engana na apprehensão do mesmo. Se assim não fosse, se a faculdade podesse enganar-se na apprehensão do objecto \emph{proprio}, devia dizer-se que ella está e não está ordenada para o seu objecto: --- está ordenada, porque tende para esse objecto. --- não está ordenada, porque o não alcança: o que é absurdo. --- Ora o objecto \emph{proprio} da intelligencia é a \emph{essencia} das coisas: porque é esta que \emph{pri meiramenie} e por si é percebida pela intelligencia. Ensina S. Thomaz: « Quidditas rei est proprie objectum intellectus » (\emph{De} ver., q. 1. a. 12). --- Portanto a intelligencia; não pode enganar-se na apprebensão da \emph{essencia} das coisas, Diz S. Thomaz: «Ad \emph{proprium}. objectum unaquaeque potentia, per se ordinatur, secundwm \emph{quod ipsa}. Quae autem sunt bujusmodi, semper eodem modo se habent. Unde, mauente potentia, non deficit ejus Judicium circa proprium objeetum. Objectum autem \emph{proprium} intellectus est \emph{quidditas} (\emph{essentia}) rei ; unde cirea \emph{quidditatem}. rei, per se loquendo, intellectus won \emph{fallitur} » (\emph{Sum}. \emph{Th}, p. L, q. 85. a. 6). --- Pode a \emph{apprehensao} da intelligencia ser \emph{deficiente}, \emph{imperfeita} e mesmo \emph{negativa}; mas nunca pode ser \emph{falsa}. E de que modo poderia ser falsa ? Talvez porque, em vez de representar a essencia de uma coisa, representa a essencia de outra? Mas, nesse caso, a apprehensão é verdadeira com relação a essencia, que é representada, e, assim, é sempre verdadeira. --- Se, pois, \emph{a intelligencia} se engana, alguma vez, na simples apprehensão \emph{de uina} essencia, \emph{não} se engana emquanto \emph{apprehende}\pgn{167}simplesmente uma essencia, mas emquanto nessa apprehensão enxerta, ex\par plicita ou implicitamente, um \emph{juizo}. E este juizo pode introduzir-se por dois modos: --- ou quando a essencia de uma coisa se attribue. a outra (p. ex. quando a definição do \emph{circulo} se attribue ao \emph{triangulo}]: --- ou quando a uma essencia se altribue um elemento extranho (p. ex. quando se define o homem : animal racional quadrupede)l Mas é claro que então o erro pertence pro- . priamente ao juizo. latente na apprehensão, e não a propria apprehens\&o. = Esta é sempre verdadeira (\emph{Sum}. \emph{Th}., p. I, q. 17, a. 31 q. 85, 1. c).}. : ---

(1). A \emph{intelligencia}, sendo faculdade \emph{immaterial}, tem por objecto o éwma- --- . \emph{terial}. --- O \emph{immaterial}. \emph{é iucomplexo} ou \emph{complexo}. O \emph{incomplexo} é a essência ven das coisas; O \emph{complexo} é uma \emph{verdade}. --- A \emph{verdade}, quando é dotada de evi- \emph{CS} . dencia \emph{immediata} e é percebida pela simples inspecção dos termos, torna-se N e é primeiro \emph{principio} da demonstração. (\emph{Posterior}., lib. I, lect. 7). --- "e

Expliquemos a natureza d'esses \emph{principios}. \emph{Principio}, em geral, é uma --- coisa, de que procede outra. Diz S. Thomaz: « Principium est omne id, à quo. aliquod procedit». (\emph{Sum}. \emph{Th}., p. I, q. 33, a. 1). Como a coisa, de que uma --- outra procede, pode ser uma \emph{realidade}, ou uma verdade, d' ahi um duplice e principio; --- o principio real, de que procedem os entes \emph{reaes}, --- e o prin- = cipio \emph{logico}, de que procede a \emph{demonstração}. --- O principio \emph{logico} pode ser \emph{primeiro} e segundo. O \emph{primeiro} é uma verdade, que --- se refere á nossa \emph{OP primeira} idea, que é a do \emph{ente}, emquanto \emph{ente}, --- é anterior a toda a these \textbar{} e a toda a hypothese, --- penetra, de um modo explicito ou implieito, em , C3 todas as demonstrações, --- e extende- -ge a todas as disciplinas e sciencias. ; --- o \emph{segundo} é uma verdade, que --- se refere a um \emph{particular modo}, do ente, e \emph{SA}: Lo --- tem ó sei fundamento e a sux resolução numa verdade anterior, --- e ás --- abrange uma especial sciencia-ou-disciplina. --- Os \emph{primeiros principios} co- -.. : stumam chamar-se tambem ---- \emph{axiomas}, ou \emph{dignidades}, porque oceupam um = logar proeminente na ordem dos conhecimentos bumanos, --- \emph{metaphysicos}, "porque são compostos de ideas \emph{universaes} e referem-se á essencia das coisas, \emph{O} --- \emph{analyticos}. porque, nelles, a conveniencia do predicado com o sujeito é \textbar{} percebida pela simples \emph{analyse} dos termos, --- \emph{immediatos}, porque são per- \textgreater{} cebidos de um modo \emph{immediato}, sem discurso. --- Estes \emph{primeiros principios P} distinguem-se dos \emph{principios de facto}, constituidos por verdades \emph{syntheticas}, \emph{NP} cujo predicado attribue-se ao sujeito, não pela analyse dos termos, mas uni- ---. camente pela experiencia, como : o sol resplandece, o mundo existe. !

Os \emph{primeiros principios}, ou axiomas, são verdades metapbysicas, e d'ellas ---

com se occupa, com o devido desenvolvimento, a \emph{Metaphysica}. Todavia, a \emph{Logica} \textasciicircum{}

") não pode deixar de indica-los, visto serem elles a base das verdades logicas e o fundamento de toda a demonstração, --- Reservando-nos, portanto, a provar, na \emph{Ontologia}. a verdade d'esses \emph{principios}, sobretudo dos principios de coniradicgzo e de \emph{causalidade}. limitamo-nos a um ou outro apontamento. --- Em todos os \emph{principios} esta contido o principio de \emph{contradicção} (porque em todas as ideas. está contida a idea de ente, da qual é contradictoria a idéa de \emph{nav-ente}); de modo que os outros- prtirctgios" \emph{se}" podem considerar como differentes formas d'esse principio. --- O principio de \emph{identidade} pode tambem formular-se assim: o que é, é. Applicações deste principio são os seguintes: --- a verdade deve ser absolutamente coherente a si mesma (principio de \emph{coherencia}), --- nenhum ente pode ter propriedades oppostas. --- O principio de \emph{ração sufficiente} não deve confundir-se com o de \emph{causalidade}, porque o \emph{ser} não pode confundir-se com o ser \emph{effeito}. Deus não pode ser \emph{effeito}, porque a. Causa primeira não pode ser causada, mas tem a \emph{razão} B \emph{sufficiente} da sua existencia na infinita perfeição da sua essencia. --- Leibnitz al. disse --- que o principio de \emph{contradicção} é o fundamento das verdades meces- : É \emph{sarias}, --- e que o principio de \emph{ração sufficiente} é o fundamento das verdu-

(153. A.intelligencia não \textbf{se engana na apprehensão da} essencia das coisas.

a) Nenhuma faculdade se engana na apprehensão do objecto \emph{proprio} ; - porque-existe-um-nexo-natural entre a faculdade e o seu objecto \emph{proprio}, e por isso, se a facaldade se enganasse com relação a. esse objecto, deveria dizer-se que ella foi feita e não foi feita para o seu objecto: o que é absurdo. Ora o objecto \emph{proprio} da inteligencia é a \emph{essencia} dás coisas. Logo a intelligencia não se engana na apprebensão da \emph{essencia} das coisas (O).

b) Se a intelligencia se'enganasse na apprebénsao da \emph{essencia} = das coisas, enganar-se-hia, ou porque attribue á essencia uma 3 propriedade, que não tem, ou porque lhe nega uma propriedade, que tem. Ora a intelligencia; quando-apprehende-aessencia-das.--- coisas, não lhe attribue, nem lhe nega propriedade alguma. Logo não se engana na apprehensão da. \emph{essencia} das coisas \footnote{A \emph{falsidade} encontra-se, propria e formalmente. no \emph{juízo}, quando a um sujeito attribuimos um predicado. que não lhe convem, ou quando lhe negamos um predicado, que lhe convem. Ora a nossa intelligencia, quando apprebende simplesmente uma \emph{essencia}, limita-se a considerar essa essencia, sem lhe attribuir ou negar alguma propriedade ou nota, de que esta ou pode estar cercada na realidade. e por isso a apprehensão não pode'deixar de ser. verdadeira. \emph{Abstrakentinm non est mendacium}; dizem os escholasticos. A . abstracção, pela qual consideramos uma coisa sem outra, não é erro, nem fonte de erros. O \emph{erro}, como acabamos de dizer, começa, quando se pronuncia um \emph{juizo}, explicito ou implicito (\emph{Sum.-Th}., p. T, q. \$5. a. 1 ad 1).\par (2j Chamamos a attenção sobre este argumento. A inteligencia, emquanto é um meio \emph{especial} para a consecução da verdade, pronuncia exclusivamente aquelles juizos. cujo predicado está contido na noção do sujeito. Ora a intelligencia ou tem a noção do sujeito, ou não a tem. Se não a tem. não pode pronunciar \emph{nenhum} juizo. Se a tem, conhece, ao mesmo tempo, que o predicado convem, ou não convem ao sujeito : e, quando affirma esta conveniencia ou desconveniencia por meio do juizo, é impossivel que se engane, como é impossivel que não conheça o sujeito, quando o conhece, Diz o Àngelico Doutor: «Jntellectus \emph{naturaliter} et ex \emph{necessitate} inhaeret primis principiis ». (\emph{Stowm}. \emph{Th}., p. I, q. 82, a. 9. --- Na doutrina do S. Doutor, a natureza\pgn{168}dos \emph{primeiros principios} é tal que não só é necessario que elles sejam por si verdadeiros, mas tambem é necessario que se veja essa. verdade. Diz assim: « Proprium horum principiorum (primorum) est quod non solum necesse est ea vera per se esse, sed etiam necesse est videri quod suut per se vera» (\emph{Post}. \emph{Anolyt}., I, \textbar{}. 19), E não pode deixar de ser assim. Estes \emph{principios} exprimem todos o principio de identidade, ou d' este recebem. toda a força. O sujeito e o predicado são € apparecem por tal modo unidos, que a affirmação 'ou a .negacio-de-um.importa a aífirmação: ou a negação do outro. \textasciicircum{}\par 'D'ahi se segue que, na apprebensão dos primeiros \emph{principios}, a intelligencja não se engana --- nem por \emph{si} mesina, ---.nem por uma \emph{accidentalidade}. Porque, estando \emph{o} predicado contido. na definição ou essencia do sujeito d'esses juizos, a intelligencia emtanto pode enganar-se na apprehensão dos principios, emguanto póde'enganar-se na apprehensão do \emph{sujeito}, Ora a intelligeneia não pode enganar-se na apprehens\&o do sujeito.d'esses principios --- nem \emph{por st} mesma, nem por uma \emph{aceidentalidade}. Não pode enganar-se \emph{por si mesma}; porque o \emph{sujeito} d' esses principios constitue o objecto \emph{proprio} da intelligencia, e esta não se engana nx apprehensão do objecto \emph{proprio}. Não pode enganar-se por uma \emph{accidentalidade}; porque o engano, se é possivel com relação aos conceitos \emph{mediates} ou \emph{complexos}, formados por precedentes juizos, não é possivel com relação aos conceitos \emph{zmmediatos} e \emph{Sinplicissimos}, que constituem o 'sujeito dos \emph{primeiros principios} e que são os de ente, \emph{todo}, \emph{parte}, causa, \emph{effeito}, etc. d}.

154. A intelligencia não se \textbf{engana} na apprehensão dos primeiros principios. \emph{VON}

a) Os \emph{primeiros principios} são juizos, cujo predicado está e se conhece, de uma maneira evidente e immediata, contido necessariamente na propria essencia do sujeito, de modo que quem percebe essa essencia, percebe nella e por ella o predicado; e assim a apprehensão dos primeiros \emph{principios} reduz-se á apprehensão da \emph{essencia} das coisas. Ora .a intelligencia não se engana na apprehensão da essencia das coisas. Logo nem se \textasciicircum{} engana na apprehensao dos \emph{primeiros} principios” E E

b) Se a intelligencia se enganasse na apprehensão dos \emph{primeiros principios}, \emph{o} engano podia provir --- ou da parte dos \emph{principios}, --- ou da parte da intelligencia. Ora o engano --- não pode provir da-parte-dos \emph{principios}, porque são pot si-certissimos e evidentissimos; --- nem pode provir da parie da \emph{intelligencia}, porque a intelligencia conhece os \emph{primeiros principios}, não por um \emph{termo medio}, mas imroediatamente e sem discurso, emquanto, por uma simplicissima e facillima analyse, conbece que o-predicado está contido va essencia do sujeito. Logo a intelligencia não se engana na apprehensão dos \emph{primeiros principios} ().

"\footnote{Notamos, para a efficacia do argumento, que o \emph{fermo medio} é necessario, quando a intelligencia ignora se. o predicado convem, ou não, ao sujeito ; porque, então, comparam-se o predicado e o sujeito com esse termo" medio, e o resultado da comparação é a affirmação da convenieneixz ou da desconveniencia do predicado com o sujeito. Mas, quando o predicado pertence á propria essencia do sujeito, pelo facto mesmo de conhecermos o sujeito, conhecemos que lhe convem esse predicado. Por isso, logo que sejam conhecidos os termos da proposição, nffirma-se \emph{immediatamente} á conveniencia ou a não-conveniencia do predicado com o sujeito. --- As proposições, que exprimem os \emph{primeiros principios}, dizem-se € são \emph{conhecidas} por \emph{si mesmas} e em \emph{relação} « nós (\emph{per} se \emph{notae et quoad} nos). São conhecidas --- \emph{por si} mesmas, porque \emph{o} predicado está incluido na noção do sujeito, --- e em relação a nós, porque vemos \emph{immediatamente} essa inclusão.(De Ver., q. 10, a. 12).\par O dito ácerca da legitimidade da \emph{intelligencia} é mais uma \emph{explicação} do} \pg{169}  . 155. \textbf{Regras para} a \textbf{intelligencia}. --- A \emph{inteligencia} será um \emph{OD} meio legitimo e infallivel para a acquisição da verdade, se ob- Ps servar as regras seguintes:

I. --- Não deve sahir-da sua-esphera; a qual;-como -disse--- 2 --------------------------------------- n B \textasciicircum{} u " ---. que uma demonstração, Por quanto, essa legitimidade é de uma evidencia ua tão iwmediata, que \emph{não pode} ser rigorosamente \emph{demonstrada}, --- nem seria- EN mente \emph{impugnada} sem \emph{petição de principio}. --- Não pode ser rigorosamente demonstrada ; porque, baseando-se a demonstração em principios, que, em. ultima analyse, hão de ser \emph{immeriatamente} evidentes, segue-se que. para \textasciicircum{} provar a legitimidade da \emph{intelligencia}, devemos já suppol-a como meio legitimo para o percepção do seu objecto. --- Mas do facto de a legitimidade da intelligencia não se poder demonstrar com todo o rigor. não se segue que ella seja menos evidente. Não pode demonstrar-se por isso mesmo que \& sua evidencia é \emph{immediata}. O sol não se demonstra, vê-se, --- Nem pode ser/seriameníe \emph{impugnada} sem \emph{peticzo} .\emph{de}. \emph{principio}. Quem impugnasse a legitimidade da \emph{intelligencia} na apprehensão dos \emph{primeiros principios}, deveria fazer uma. argumentação baseada, em ultima analyse, nos mesmos \emph{principios}. Uv e por isso admittiria a: Jegitimidade da \emph{inteligencia}. ao mesmo: temp Dque-a impugna. --- Portanto, quem nega os \emph{primeiros principios}, nega-os \emph{apparer} temente e por \emph{ostentação}, mas Não \emph{internamente} é por convicção, A Proposição \emph{contradictoria} de nm \emph{principio prineiro} é um absurdo. e o absurdo não se concebe, embora se exprima (ct. \emph{Post} 4\emph{nalyt}., . 19. :

Muitos escriptores tém desvi rtuado a \emph{natureza} dos \emph{primeiros principios}. ' ---. --- Façamos um breve exame critico d' algumas d' essas opiniões. ; \textbar{} o

a) Kant alfirma que, para explicar a \emph{necessidade evidente} dos \emph{primeiros principios} da. razão, basta admitir um \emph{habito cego} e \emph{innato}, que elle chama \emph{forma a-priori} da nossa mentalidade. --- Mas, sc assim fosse, todos os nossos juizos, que uma longa experiencia tornou habiluaes, --- p. e., o fogo queima, a corpo é pesado, --- teriam a mesma necessidade evidente. E todavia não é assim ; porque podemos conceber um fogo, que não queima, e'um corpo, que não é pesado. A \emph{necessidade evidente} dos principios da razão, longe de " ser o producto de um babito cego. é o resultado da clara percepção de uma... HB --- relação necessaria entre o sujeito e o preiticado. O habito poderá transformar E um juizo \emph{reflexo} em juizo \emph{espontaneo}, algo \emph{confuso}; mas um juizo espon- \textbar{} taneo, embora confuso. não é um juizo \emph{synthetico} a-priori, que se enuncie necessariamente, sem nenhuma razão, nem distincta nem confusa.

by O mesmo Kant allimma que os \emph{primeiros principios} c. em geral, os juizos \emph{analyticos} se reduzem a uma \emph{taufologia} (A = A). --- Contestamos a affirmação de Kant. O subjecto e o predicado, nos juizos \emph{analyticos}, não são \emph{formalmente} identicos, e o predicado exprime de um modo \emph{explicito} e \emph{distinco} o que o sujeilo exprimia de uin modo \emph{implicito} e \emph{confuso}. Além d isso, quando se compara \emph{o} predicado com 'o sujeito, apparecem novas Ap s. relações entre esses dois termos, novos respeitos, que antes não appareciam. ] Ui Finalmente, se assim fosse. tambem a \emph{Mathematica} reduzir-se-hia \& uma n ut \emph{tautologia} ; o que nem o proprio \emph{Kant} admitte, porque diz que as conclu- Us \emph{um} sões mathemalicas são \emph{Juros synthelicos a-privri}. ; Ut

e) Stuart \emph{Mill} (\emph{Log}, t. T, p. 26\%), seguido por todos os \emph{positivistas}, diz

\textasciicircum{} \pg{170}  di 2 Locicao SECÇÃO CEGUNDA

mos, é constituida pela essencia das coisas e pelos \emph{primeiros principios} da demonstração. \textbar{}

IP. --- Quando uma verdade, embora dotada de evidencia

que 08 \emph{primeiros principios}, OU \emph{axiomas}, S\&o \emph{as} mais universaes \emph{inducções da} experiencia \emph{sensivel}. --- Negamos que os \emph{primeiros principios} sejam as mais universaes indvucções da experiência: Por quanto, esses \emph{principios} exprimem verdades amnalyticas € \emph{metaphysicamente}: \emph{necessarias}, ao passo que as \emph{inducções da experiencia}, sensivel, embora geraes, são por si \emph{synthelicas}, e. \emph{metaphysicamente contingentes}. Ora 0 mais.não. pode identificar-se com o menos. Além d'isso, a intelligencia; para adherir aos \emph{primeiros principios}, deve perceber a existencia de um nexo \emph{necessario} e immutavel-entre o predicado e o sujeito. Ora a experiencia sensivel, só por si, é absolutamente incapaz de manifestar um nexo necessario e immutavel:; porque, se .attesta o \emph{facto} do nexo, não pode attestar a natureza d'elle, isto é, não pode dizer se o nexo é necessario ou contingente, immutavel ou mutavel. Logo os \emph{primeiros} principios não são as mais universaes indueções da experiencia sensivel,

*

, Pergunta-se :.Onde..têm-.o sen “fundamento. os \emph{primeiros principios} ?

i Respondemos que: os \emph{primeiros principios} têm o seu fundamento \emph{proximo}\textbar{}nos \emph{factos} da natureza, e O seu fundamen to ultimo na \emph{Inteligencia} de « Deus. --- Desenvolvamos esta resposta.

a) Os \emph{primeiros principios} têm o seu fundamento proximo nos \emph{Factos da natureca}. Este proposição tornar-se-ba evidente. se considerarmos a origem dos nossos conhecimentos intellectuaes.. O primeiro objecto, que a nossa intelligencia percebe no estado actual, é a \emph{essencia} das coisas materiaes, c esta essencia é o objecto \emph{proprio}. da. mesma faculdade. Diz S. Thomaz: « \emph{Primum}, quod intelligitur-a nobis secundum statum praesentis vitae, est \emph{quidditas} rei \emph{materialis}; quae est nostri intellectus objectum » (\emph{Sum}. \emph{Th}., p. 1, 5: S8, a. 3) --- A apprebensão . da \emph{essemcim} é o primeiro acto da intelligencia. Percebendo a essencia das.coisas, a. intelligencia percebe os elementos, de que a essencia é constituida. O primeiro d' estes elementos percebidos é o énie;.e por isso a primeira” 'e universalissima noção, que está con- ; tida em todas as noções e em que todas se resolvem, é a noção de \emph{ente}. --- Concebida a noção ou idea de ente. a intelligencia affirma immediatamiente a \emph{repugnancia do} ente com o nZo-ente, com o nada, e pronuncia o principio de contradiceção: uma coisa que é \emph{não pode deixar de} ser emquanto é. Como tambem, a formação de outras ideas (depois da, do \emph{ente}), por ex.. das ideas de \emph{substancia}, de causa, de ordem, etc.. dá origem a outros \emph{principios}, --- ao principio de \emph{causalidade}, de \emph{razão sufficiente}: etc. --- Esta promptidão e fa- 'eilidade, com que a nossa intellizencia, . depois da percepção da essencia, pronuncia os \emph{primeiros principios}. 6 chamada por S. Thomaz \emph{kabito} natural da \emph{intelligencia} humana, e os mesmos \emph{principios} são chamados \emph{innatos}. Mas, repetimos, os \emph{primeiros principios} podem chamar-se imnatos, não emquanto sejam infundidos por Deus na nossa intelligencia no momento da creação, --- mas emquantLo recebemos de Deus uma natural disposição para os pronunciarmos, e emquanto os pronunciamos depois de uma facillima e simplicissima \emph{abstracção} Diz S. Thomaz: «[ntellectas \emph{principiorum}, dicitur esse \emph{habitus naturalis} ; ex ipsa enim natura animae intellectaalis eonvenit homini

CAP.

H. --- MEIOS PARA A ACQUISIÇÃO DA VERRADE - 471

\emph{immediata}, \emph{é synthetica}, e não \emph{analytica}, então, \textbf{para} a intelli- gencia se não enganar, é indispensavel o exercicio de outra fa- \textasciicircum{}

quod statim, cognito quid. est. \emph{fofum} et. quid.est.pars, \textbf{cognoscat} quod. omae... \emph{totum}, \emph{est} maius sua \emph{parte}; et simile est in. ceteris. Sed quid sit fotu» ot quid sit \emph{pars}, coynoscere non potest nisi per species intelligibiles a \emph{pkantasmatibus} acceptas » (\emph{Sum}. \emph{Th}, LI, q. 51, a. 11. E o S. Doutor repete. a . cada passo que o conhecimento dos \emph{primeiros principios} se funda, no testemunho dos sentidos. e por isso nos factos particulares : «Tpsorum \emph{princi piorum} cognitio in nobis ex \emph{sensibilibus} causatur; nisi enim aliquod totum . Sensu percepissemus, non possemus intelligere quod totum esset maius parte» (\emph{C}. \emph{Gentes}, ]. 1L, c. 831. Concluimos que os \emph{primeiros principios} tém o seu fundamento \emph{proximo} nos tactos da natureza, .

P) Mas se os \emph{primeiros principios} têm o seu fundamento \emph{proximo} nos factos da natureza, não têm nesses factos o seu fundamento \emph{ultimo} (como pretendem os \emph{positivistas}). Os factos são \emph{mutaveis} e \emph{contingentes}. 20 passo que os \emph{primeiros principios} exprimem uma verdade immutavel, necessaria, \emph{eterna}, que não deixaria de ser o que é.. embora não existisse facto algum. --- Dizemos, pois, que esses \emph{principios}: t\&m-o seu. tundamento-seltimo na-7. \emph{telligencia de} Deus, Porquanto, o\& \emph{principios}, sendo verdades, devem residir na \emph{éntelligencia}, como em seu sujeito: Esta intelligência não pode sera do homem: porque os \emph{principios} impõem-se a nossa intelligencia, e devemos acceital-os; eram verdades antes de nós os percebermos, e não deixariam de o ser, ainda que a nossa intelligencia acabasse: Logo a intelligencia, \textasciicircum{} em que, como em sujeito proprio. residem os \emph{principios}, deve \emph{ser} a de Deus; porque verdades \emph{necessarias} só podem ter o fundamento \emph{ultimo} numa intelligencia necessaria. e esta é a de Deus. --- Bossuet escreve : « Áinda que não houvesse na natureza algum triangulo, seria sempre certo e inhubitavel que os tres angulos de um triangulo são eguaes a dois reclos. O que vemos na essencia do triangulo é certamente independente de todo o triangulo existente, Mais ainda: não é a intellizencia que dá \emph{o} ser a verdade; e por isso, ainda que a intelligencia fosse destruida, estas verdades subsistir am immutaveis» (\emph{Logique}, T, 361. E noutro logar: «Se pergunto a mim mesmo onde e em qual sujeito estas regras subsistem eternas e immutaveis, sou obrigado a admittir «um \emph{Ser}, \emph{onde a verdade é eternamente subsistente} e \emph{ondé ella é sempre percebida}; e este Ser deve ser a \emph{propria} Verdade e deve ser a Verdade toda» (\emph{Connais de Dieu} et \emph{de soi-máme}, IV, 5). --- Da Intelligencia de Deus os \emph{principios} refleclem-se na intelligencia humana, pelo facto d'esta ser effeito e semelhança d'aquella. Este reflexo, porém, não é \emph{directo} ou \emph{inmmediato}, pois faz-se, como dissemos, por uma facil abstracção das coisas particulares, que reproduzem em si os arelietypos da intelligencia divina.

Para entendermos mais claramente esta altissima doutrina, que é de S. Thomaz, tembremo-nos de que Deus nos instmue e nos torna participes das suas ideas por meio das coisas creadas, que são, por isso, outras tantas linguas, com que Elle nos \emph{falla na} orlem natural. Na verdade, dois elementos concorrem para i renlisacüo dos nossas conhecimentos: a \emph{faculdade intellectal}, \emph{c} 0 \emph{objecto intelligivel}. \emph{que} se une \emph{e} manifesta a nossa \emph{intelligeneia por una especie immalerizl} (Cf. \emph{Sum}. \emph{Th}. p. 1, q. 105, a. 3). Ora tanto \emph{a} nossa \emph{inteligencia}, \emph{quanto} o \emph{objecto intelligicel}, tem á sua razão \pg{172}  culdade. Assim, para pronunciar este juizo --- o hrundo \emph{existe}, --- a inteligencia deve servir-se do ministerio dos sentidos externos.

suprema, o seu fundamento ultimo em “Deus. --- Tem o fundamento \emph{ultimo} em Deus a \emph{nossa intelligencia}; porque esta é uma imagem ou semelhança da Jntelligeneia incereada. e a nossa, luz intellectual é um raio da luz divina, em que estão contidas as razões Nntelligiveis de todas as coisas. Diz S. 'Thomaz: «Ipsum lumen inlelleetuale, quod est in nobis, nihil est aliud quam quaedam participata similitudo Iumínis, increati, in qno continentur rationes (ideae).aeternae » (Sum. \emph{Th}., p. 1, q. \$\&.a. 5). --- Tem o fundamento \emph{ultimo} em Deus o \emph{objecto intelligivel}. Com -effeito, todas as coisas" creadas foram feitas em conformidade com os archetypos divinos, e. não se dizem e não são \emph{verdadeiras}, senão emquanto ' reproduzem fielmente esses archetypos. Diz S. Thomaz: « Res naturales dicuntur verae, 'secundum quod assequuntur similitudinem specierum. quae sunt in mente divinà » (\emph{Sum}. \emph{Th}., p. T. q. 16. a. 1 Sendo dotadas de verdade, as coisas creadas são tambem dotadas de \emph{intelligibilidade}, porque toda a verdade é por si intelligivel. Essa intelligibilidade é connatural e accessivel a nossa intelligencia; a qual. por ser uma: faculdade: espiritual; ias-existerité-nà materia, deve perceber o \emph{intelligivel}, escondido sob a concreção da materia. D'este modo, a intelligencia humana. --- conformando-se com as coisas creadas, conforma-se com a intelligencia divina, --- percebendo as essencias.das, coisas creadas. percebe os proprios archetypos divinos --- é. pronunciándo 0s \emph{primeiros principios}, exprime os principios da Inteligencia divina, qué os exprimiu e reproduziu nas creaturas. --- Estes \emph{principios}, sendo conhecidos por uma facillima e simplicissima. abstracção, não podem-ser ignorados, muito meros contestados seriamente por ninguem; e por isso constituem o patrimonio da intelligencia humana; emquanto tal. E, como os \emph{primeiros principios} são os criterios directores de todos os nossos.juizos e são um reflexo das razões ou ideas etdrnas dé \textsc{Deus}, podemos dizer que o homem vé tudo e julga de tudo nas razões ou ideas eternas de Deus. "Tal é a doutrina de S. Thomaz (\emph{C}. \emph{Gentes}, l IH, e. \&7: Sum. \emph{Th}., p. J. q '84. a. 5. --- Para maior clareza, adverlimos : que os archetypos divinos representar os varios e infinitos modos, por que \emph{n} a Essencia divina é e pode ser representada na creação. Por isso cada um d'esses modos representa a Essencia divina com limitação, e a nossa intelligencia. quando descobre nas coisas creadas esses archetypos, não descobre 'a propria Essenci» divina senão dc uma: maneira muito deficiente ou limitada, Em poucas palavras: as coisas creadas são reproducções ou imitações fidelissimas dos archetypos divinos; mas estes são imitações finitas da infinita Essencia divina. --- Concluimos que os \emph{primeiros principios}, pelo facto de se fundarem \emph{proximamente} .mas coisas creadas, as quaes reproduzem fielmente em si mesmas os archetypos divinos, têm a sua razito suprema, 0 Seu fundamento ultimo \emph{na} Inteligencia divini.

D'onde se segue que da existencia. dos \emph{primeiros principios}, que a NOSSA E inteligencia vé reproduzidos nas coisas sensiveis, pode lecitimamente dedu- : zir-se a existencia de Deus, --- Nem se diga que essa deducção importa uma passagem arbitraria da ordem ideal para a real. Porquanto esses principios, embora sejam \emph{percebidos} pela nossa intelligencia, todavia não são um mero producto logico, a que nenhuma realidade corresponda; mas são verdades \pg{173}  meio do \emph{raciocinio}, deduz de uma verdade conhecida uma ver- XD dade desconhecida. --- A verdade, que, pelo raciocinio, se \textbf{deduz} \emph{Rus} de outra verdade, é dotada de evidencia mediata \footnote{O. conhecimento de objectos \emph{izunediatamente} evidentes faz-se por um acto \emph{intuitivo} ada faculdade perceptiva intellectual. Esse conhecimento é 2 base de lodos os outros conhecimentos. Por isso, quando um objecto não é por si dotado de evidencia \emph{immediata}, devemos recorrer a um termo. medio, que nos mostra se ba uma relação necessaria, um nexo essencial entre esse objecto e os objectos \emph{immediatamente} evidentes. Existindo essa relação, esse nexo, é claro que, --- se são verdadeiros os objectos \emph{iuumediatamente} evidonies, devem ser tambem verdadeiros os objectos, que não têm essa evidencia immediata, mas que. estão relacionados. com os que. a têm, --- como tambem... se fossem falsos estes segundos objectos, deveriam tambem ser falsos os que são \emph{immediatamente} evidentes. Diz S. Thomaz: « Quaedam propositiones sunt necessarias, quae habent connexiouem necessariam cum \emph{primis} principiis, sicut conclusiones demonstrabiles, ad quarum remotionem sequitur \emph{remotio primorum principiorwm}. Et his intellectus ex necessitate assenüt, cognita \emph{connexione necessaria conclusionum} ad \emph{principio} per demonstrationis deductionem » (Sum. \emph{Th}.. p. I, q. 82, a. 2). --- Repetimos que a \emph{intelligencia} € á \emph{razão} não são duas faculdades distinctas, mas denotam duas distinctas funcções da mesma faculdade intellectual, a qual descobre a relação entre dois Lermos --- ás vezes de um modo \emph{immediato}, --- e outras vezes de um modo \emph{medieto} por um termo \emph{medio}. Por isso, alguns escriptores, dando o nome de \emph{razão} a faculdade intellectual pexecptiva, dividem-na em razão \emph{inheitiva} e em rasão \emph{discursiva}, em conformidade com essa duplice funeção. (Cf. \emph{Sum}. \emph{Th}., p. 1, q. 89, a. 8; \emph{De} Ver., q 15. a. 1).\par A \emph{razão} pode ser \emph{theoretica} € \emph{pratica}. \emph{A tleoretica}, ou \emph{especulativa}, limita-se\pgn{174} a percepção da verdade; a \emph{pratica} applica á verdade percebida a operação. Não são duas faculdades, mas- uma só; visto que O objecto da \emph{theoreretica é o} mesmo que o da \emph{pratica}, \emph{e} a applicação da verdade percebida a operação é uma coisa apenas accidental. --- Kant concede que a razão \emph{pratica seja} um meio legitimo para a consecução dà verdade; mas nega que tal seja a razão \emph{theoretica}. \emph{Mas} engana-se, Se a razão \emph{pratica} se identifica com a \emph{theoretica}, ou ambas são meios legitimos, ou ambas illegitimos.}. \emph{T} ------ \emph{UN objectivas}, que têm a sua expressio e realisação nas coisas concretas e visi- ens veis, e por isso representam uma entidade real, Sendo entidades \emph{objectivas P}. Oü reaes, OS \emph{primeiros principios} suppõem uma intelligencia objectiva e \textbf{real}; porque a causa. deve ser proporcionada ao elfeito. Esta intelligencia não pode : \emph{SO ser} a da creatura; não só porque a intelligencia da excatura, sendo contin- --- gente e temporanea, não poderia imprifhir nos primeiros principios um cunho \emph{um} de necessidade e de eternidade, mas tambem porque não teria 2 energia suf --- ficiente para os imprimir na concreção da materia. Logo deve ser a "intelli- i gencia de um Ente superior e increado, que tenha a força de conceber \textbf{essas} \emph{O} verdades e de as reproduzir nos entes finitos; visto que Loda a forma, que e existe no effeito, deve preexistir numa causa no estado de exemplar. Esse ente superior e inereado é Deus. --- Portanto o argamento, que deduz da existencia dos primeiros principios; conhecidos por nós, a existencia de \textbf{Deus}, é .\emph{arposteriori}, Porque. fundado num. \emph{effeito intelligivel}, emquanto d'essas ver:

te cleva-se a fonte de toda a verdade. Diz S. 1homaz: «. intellecta sit \emph{aeterna} quantum ad id quod intelligitur, non potest concludi quod anima sit aeterna, sed quod vcritates intellectae \emph{fundentur} in \emph{aliquo aeterno}. Fundantur enim in.jpsa prima veritate, sicut in eausa univorsali contentiva omnis veritalis » (\emph{C}. \emph{Gent}. \emph{II}, \$4). ”

157. Objecto da razão, --- O \emph{objecto}, pois, da \emph{razão} é constituido pelas verdades mediatas, isto é, pelas \textbf{verdades} deduzidas, pelo \emph{raciocinio}, de outras verdades, precedentemente conhecidas. O raciocinio, como dissemos, é \emph{deductivo} ou. \emph{znductivo}. \footnote{Para evitar repetições, remettemos o leilor para o que dissemos no cap. 111 da I parte da \emph{Logica} ácerca do \emph{raciocinio} e suas \emph{especies}.\par (\emph{B}) A razão não se engana \emph{por si}, isto é, pela sua propria natureza. Estando ella naturalmente, orientada ou ordenada para a verdade, não pode dizer-se, sem' uma contradicção flagrante; que ella está naturalmente orientada ou ordenada para o erro; e seria \emph{naturalmente} orientada ou ordenada para o erro, se se ênganasss \emph{por si}, pela sua propria natureza. --- Nem pode a \emph{razão} mudar essa orientação, essa ordem; porque não é uma faculdade \emph{livre}, mas \emph{necessaria}, e uma faculdade, \emph{necessariamente} ordenada \textbf{para} um certo tim, não pode ser ordenada para outro fim contrario. Essa orientação. ou: ordenação da razdo para a verdade refere-se tambem, e sobretudo, ás verdades \emph{mediaías}, isto .é, ás verdades, que não são por si \emph{immediatamente} evidentes, mas que são deduzidas de outras verdades por meio do \emph{raciocinio}. Por quanto, se a intelligencia angelica é naturalmente intuitiva, a humana é naturalmente discursiva. Percebida a \emph{essencia} de uma coisa, e pronunciados os \emph{primeiros principios}, \emph{a} nossa intelligencia não exbauriu, não comprehendeu tudo o que se roefere' a essa coisa, e precisa de completar o seu conhecimento, e completa-o por meio do \emph{juizo}. \emph{offirmativo} ou \emph{negativo}, passando de uma pruposição para outra, de uma verdade conhecida \emph{para} outra desconhecida. (\emph{De} Ver., q. 15, a 1).}.

1508. A razão não se engana ma percepção das \textbf{verdades} mediatas. "

a) Nenhuma faculdade .se engana com relação ao objecto \emph{proprio}; porque tende para esse objecto de um modo \emph{natural}, e por isso \emph{necessario}. Ora o objecto \emph{proprio} da \emph{razão} é a verdade \emph{mediata}, deduzida, pelo \emph{raciocinio}, de outra verdade. Logo a \emph{razão} não se engana na percepção das verdades mediatas (2).

b) Quando, pelo \emph{racivcinio}, deduzimos: de uma verdade conhecida uma outra verdade desconhecida, affirmamos \emph{explicitamente} na conclusão v que \emph{implicitamente} tinhamos affirmado nas premissas. Ora a intelligencia não se engana na percepção das p emássas, que. em ultima analyse, são verdades \emph{immediatas}.

] \pg{175}  Portanto, nem a razão se engana na percepção das conclusões, quando estas foram legitimamente deduzidas das \emph{premissas}. \emph{Às conclusões} são e dizem-se verdades \emph{mediatas}. Logo a \emph{razão} não se. engana na percepção das verdades mediatas \footnote{Nem pode ser de outro modo. Se a \emph{rac\&o} se enganasse na percepção das verdades mediatas, dar-se-hia o engano --- ou porque ella não pode estar certa da verdade das \emph{premissas}, --- ou porque não pode estar certa do nexo \emph{logico} entre o antecedente e o consequente. Ora a \emph{razão} pode estar certa --- tanto da verdade das \emph{premissas}, porque estas são \emph{immediatamente} evidentes, ou fundam-se e resolvem-se em verdades \emph{inonediatamente} evidentes, --- como da verdade do nexo \emph{logico} entre o antecedente e o consequente, porque podemos, pela reflexão, convencer-nos de que foram observadas as respectivas leis logicas e de que, por isso, a conclusão se funda e se resolve nos \emph{primeiros principios}, nas verdades \emph{inmmediai«mente} evidentes. Logo a \emph{razão} não se engana na percepção das verdades :\emph{sediatas}. --- Como se vê, 2 razão, no seu movimento, parte da intelligencia e volta para a intelligencia: esta é o começo e o termo d'aquella. Diz S. Thomaz: « Quia motus semper ab imruobili procedit et ad aliquid quietum terminatur, inde est quod rátiocinalio humana, secundum: viam-inquisitionis-vel- inventionis; \emph{procedit}. a quibusdam simpliciter intelleelis, quae sunt, prima prineipja, et rursus... .--- in via judicii, resolvendo \emph{redit} ad prima. principia, ad quae inventa examinat» (\emph{Sum}. \emph{Th}., p. 1, q. 79, a. 8). \textbar{} Mas, se a nossa \emph{razão} não pode enganar-se \emph{per} se, , pela sua propria \emph{natureza}, pode enganar-se \emph{per} accidens, por um motivo que obsla a reotidão do seu exercicio. O instrumento é bom; mas o modo, por que nos servimos d'elle. pode ser defeituoso e contrario ás regras da arte. \emph{E} .a arte do raciocinio é ume das mais complexas e delicadas. A falta” de estudo e reflexão, a influencia dos prejuizos e dus paixões, a difficuldade do assumpto, e outros motivos podem desviar a razão e fazer com que deixe de observar as leis da legitima deduccüo. D' onde o erro. --- \textbf{Para} evitar o erro (accidental), a razão percorra de novo o caminho feito. averigue se a conclusão desce legitimamente das premissas, e julgue a mesma conclusão i luz dos \emph{primeiros principios}; porque a certeza da conclusão depende inteiramente da certeza dos principios.\par No texto provamos em geral que a razão não se engana na percepção das verdades mediatas. Como essas verdades são adquiridas por meio da deducção ou da inducção, deveriamos provar que a \emph{razão} não se engana -nem quando, por meio do \emph{syllogismo}, \emph{desce do universal} para o particular, --- nem quando, por meio da \emph{inducção}, \emph{sobe} do \emph{particular} para o universal. Mas, como nos respectivos logares tratámos da \emph{legitimidade} do \emph{syllogismo} ? da énducgao, julgamos desnecessaria a insisteucia.\par Tambem a legitimidade da \emph{razão} é de evidencia immediata, e não pode "igorosamente demonstrar-se nem contestar-se sem \emph{petição de principio}. Na verdade, quem quizesse demonstrar ou contestar a legitimidade da rasão, levia fazer um raciocinio. Ora todo o raciocinio parte de principios, que, em ultima analyse, são verdades"dotadas de evidencia \emph{immediata}, e chega 1 conclusões, que são verdades dotadas de evidencia mediata. Logo quem defende ou impugna a legitimidade da \emph{razão}, já presuppõe a sua legitimilade no conhecimento das verdades de evidencia \emph{mediata}.}. .

%%>>> t159

\pg{165} \pg{166} \pg{167} \pg{168} \pg{171} \pg{174} \pg{176} \tese{159}{Regras para a razão.}
A \emph{razão}, para que seja meio legitimo para a acquisição da verdade, deve observar estas regras:

[*» --- Não deve sahir da sua esphera, que é o \emph{immaterial}. Para o conhecimento.do. material é indispensavel. o. exercicio de outra faculdade. A passagem do \emph{ideal} para o real é sophisma.

IP. --- Estabelecidas as \emph{premissas}, deve deduzir a \emph{conclusão} conforme as \emph{leis} do \emph{raciocinio}. Da inobservancia d'esta regra deriva a maxima parte dos'erros.

AnrIGO l1

\emph{Consciencia e memoria intellectual}

%%>>> t160

\tese{160}{Conscienceia.}
\emph{Consciencia} é a propria intelligencia,

j - o "que, pela \emph{reflexão}, percebe a existencia \emph{actual} das operações \emph{espirituaes}, que são as da \emph{inteligencia} e da \emph{vontade} \footnote{A nossa intelligencia vão \emph{so} “percebe os objectos, que estão fora de nós, mas, por ser uma faculdade espiribual e immaterial, pode dobrar- se sobre si mesma e perceber as operaqóes.espirituaes, que se realisam e estão presentes \emph{no intimo} do-nosso espirito e que derivam da \emph{intelligencia} e da \emph{vontade}. --- \emph{Emquanto} percebe essas operações \emph{espóriuaes}, \emph{actualmente} presentes no nosso espirito, a intelligencia toma o nome \emph{de consciencia}. Portanto o acto: proprio da \emph{consciencia} € á \emph{reflexão}, --- A \emph{reflexão}, propria. da \emph{consciencia}, não deve confundir-se com a acto, pelo qual a intelligencia torna a considerar um objecto já percebido, para ter d' elle um conhecimento mais distincto e completo; e que mais: propriamente podia chamar-se \emph{recogitação}. Tambem a \emph{recogitação} é: uma reflexão, mas \emph{objectiva} ou \emph{ontologica}, porque se refere ao objecto: percebido ; ao passo que o acto da \emph{coscienciu} é uma reflexão \emph{subjectiva} ou \emph{psychologica}, porque se refere ao \emph{sujeito} e ás suas \emph{operações}. --- Que a intelligeocia seja dotada de \emph{reflexão}, é coisa certissima. Não só percebemos os-objectos externos, mas percebemos que percebemos e queremos, isto é, percebemos as possas proprias operações expirituaes. Cada um de nós pode repelir com S. Agostinbo: « lntelligo me intelligere ; intelligo me velle» (\emph{De Z'rin}., 1. 10, c. 10 e 11; efr. \emph{Sum}. \emph{Th}.. p. 1, q, 87, 3. 2 e \&).}. .

161. Objecto. da consciencia --- O \emph{objecto} da consciencia é

; Jas \emph{Operações espidirecto} e \emph{indirecto}. \emph{riluaes}, \emph{actualmente} presentes: “O \emph{indirecto} é o proprio \emph{sujeito}, que exer ce aquellas operações Oo\footnote{A \emph{consciencia} percebe, pelo menos .\emph{confusamente}, duas coisas: a \emph{operação espiritual}, e \emph{o sujeito}, de que ella deriva. Por quanto a. operação percebe-se de um modo \emph{concreto}, como existe; ora existe não em si, mas no sujeito. Por isso, assim como os nossos olhos não percebem propriamente a côr, mas a coisa colorida, e as nossas mãos não apalpum propriamente a dureza, mas a coisa dura, assim tambem a consciencia não percebe propriamente o pensamento e a volição, mas o sujeito pensante e querente ; tanto\pgn{177}mais que ella, por ser uma e a mesma coisa com a iatelligencia, chega ao fundo das coisas. É verdade que a razão, pela qual a consciencia percebe o sujeito, é a operação, como a razão, pela qual os olhos percebem um objecto colorido, é a côr, e a razão, pela qual as mãos apalpam um objecto duro, é a dureza; mas tambem é verdade que a operação e o sujeito se percebem ao mesmo tempo, emquanto constituem um.só objecto, que. 6 \textbf{o.mesmo} sujeito modificado d'este ou d'Ajuelle modo. Por isso nós não percebemos que "somos e que temos: um-principio- de- vida;-senão-emquaoto-percebemos-queoperamos. (Cf. S. Thom. De \emph{mente}, a. Sj. --- Advirta-se --- que a operação \emph{espiritual} é objecto da consciencia, não emquanto representa uma coisa, mas \textbar{} :emquanto é um phenomeno \emph{interior}, Ou. uma .modificação da faculdade intellectual, --- e que a \emph{consciencia} attinge o sujeito não emquanto é dotado de uma ou de outra natureza, mas só emquanto é \emph{principio} e causo de operações \emph{espirituaes}, que, como dissemos, são as \emph{da-intelligencia} e da vontade.--- A \emph{consciencia} percebe lambem a existencia \emph{actual} das operações \emph{sensitivas}, mas só de um modo \emph{indirecto}, como explicaremos mais tarde.\par Dissemos que a consciencia percebe não só as operações da \emph{intelligencia}, mas tambem as da \emph{vontade}; porque estas tambem são immateriaes e por isso são objectos proporcionados de uma faculdade immaterial, como é a nossa \emph{consciencia}. --- Quando, pois, se affirma que o objecto \emph{proprio} da nossa intelligencia é o universal, 0 \emph{abstracto}, a àlfiriiação não é muito éxactà. Se o fosse, a intelligencia não poderia perceber as proprias operações, que não são coisas \emph{absiractos} Ou universaes,mas muito concretas e \emph{individuaes}. Portanto é mais exaeto dizer --- que o objecto proprio da intelligencia é o \emph{immaterial} --- e que, se esta faculdade percebe as coisas \emph{abstractas}, percebe-as emquanto são \emph{imunateriaes}. Fudo o que é \emph{abstracto é imunoteriol}; mas nem tudo o que é immaterial \& \emph{abstracto}. --- Mas não queremos por isso concluir que \emph{todos} os seres \emph{immateriaes} podem ser \emph{directamente} percebidos pela nossa inteligencia. Os seres, que excedem a nossa natureza, como os Anjos, e, Sobretudo, o Ser infinito, ainda que sejam immensamente intelligiveis em Si mesmos, não são, comtudo, objectos proporcionados da nossa fraca intelligencia, e não podem ser percebidos \emph{directamente}, mas só pelo raciocinio.\par D'onde se vé a excellencia e a utilidade da \emph{consciencia}. Pela \emph{intelligencia} € pela \emph{razão}, o nosso espirito eleva-se até ao mais alio grau de conhecimento, até á adoravel Divindade. Não pode subir mais. Mas, se não pode Subir mais, o nosso espirito pode dobrar-se, recolher-se em st mesmo, penetrar\par Yol. I. 13\par '\pgn{178}até ao fando da sua actividade. É o que faz pela \emph{consciencia}. É um novo campo de actividade, mais vasto, mais instruetivo, mais. delicioso, mais divino, como diz Aristoteles (\emph{Metaph}., \textbar{}. XI, c. 7). É, com effeito, neste campo que se encontram e se completam as duas ordens de conhecimentos, --- os principios e os factos, o abstracto e o concreto, o necessario e o contingente, a theoria € á pratica. --- Vê-se tambem a dilferenga entre a \emph{razão} e a \emph{consciencia}. As verdades, que constituem o \emph{objecto} da \emph{ração}, são \emph{necessarias} e dotadas. de evidencia mediata; mas as verdades, que constituem o objecto da \emph{vonsciencin} e que são as operações .espirituaes, são contingentes (porque essas operações são \emph{factos contingentes}) e dotadas de evidencia immediata.}

ans

162. A \textbf{consciencia} não se engana na \textbf{percepção das} operações espirituaes, nem na percepção do sujeito \textbf{agente}.

a) \emph{Não}. \emph{se} engana na \emph{percepção dos operações} espirituaes. --- Se a \emph{consciencia} se enganasse na-percepcio- das operações-espi--rituaes, o engano devia provir --- ou da alteração do \emph{objecto}, --- ou da alteração da \emph{faculdade}, --- ou do obstaculo do \emph{meio}. Ora o engano não pode provir --- nem da alteração do \emph{objecto}, porque este, sendo um phenomeno \emph{interior}, deve ser.attingido como .é em \emph{si} mesmo, e por isso com as alterações, que soffreu, se as soffreu, --- nem da alteração da \emph{faculdade}, porque uma faculdade \pg{178}  inorganica,. como é a consciencia, não pode estar sujeita á alteração, visto que toda a alteração de uma faculdade deriva da alteração do orgão, --- nem do obstaculo do meio, visto que os phenomenos. \emph{espirituaes}.. não. precisam de. \emph{meio} para se unirem com a faculdade, pois, estando presentes na alma, manifestam-se immediatamente a inteligencia, que é uma faculdade da mesma alma. Logo. a \emph{consciencia} não se engana na percepção das operações espirituaes \footnote{Devemos. ter sempre presente que o exclusivo munus da consciencia é o de attestar a presença das \emph{operações} espirituaes, que actualmente se encontram na alma,'sem se curar se essas operações são verdadeiras ou falsas, boas ou más; de modo. que, se uma outra faculdade se engana com relação ao seu \emph{objecto}, comtudo não se engana a consciencia, quando refere esse engano, como não mente quem refere historicamente as mentiras dos outros. --- Por isso, quando se trata da consciencia, o objecto real e o \emph{apparente} confundem-se, identificam-se; porque a \emph{modificação} do sujeito é sempre a mesma, e a consciencia, quando se limita a attestar a \emph{existencia} Y) essa modificação, é sempre verdadeira; como vcrdaderia é a afirmação de um doente de olhos, quando diz que vé todos os objectos de côr amarella, embora não sejam taes em si mesmos. Erra quem julga, não quem refere,}.

b) A \emph{consciencia} não \emph{se engana ma percepção do sujeito agente}. --- A \emph{consciencia} não se enganará na percepção do \emph{sujeito agente}, se este fôr \emph{intelligivel} e \emph{actualmente presente} á intelligencia; porque, neste caso, o engano refundir-se-hia no proprio Auctor da natureza. Ora o \emph{sujeito agente}, que é a propria alma, --- é \emph{intelligivel}; porque é \emph{immaterial}, --- e está \emph{actualmente presente}. á intelligencia nas sua \emph{operações}, de modo que é objecto de ümáa-simples- "\emph{apprehensao}: "Logo a consciencia não se engana-na.-percepção--do-sujeito -\emph{agente}- (\textasciicircum{}).\footnote{Para que a alma perceba a propria existencia e o que se passa nella, não é preciso mais do que a essencia da propria alma, que está intimamente unida com a inteligencia; pois é da essencia da alma que derivam as operações, nas quaes ella é acinalmente percebida, (S. \emph{Thom}., \emph{De} mente, a. S). ---\pgn{179}Mas, repetimos, a facilidade, com que a \emph{consciencia} percebe as \emph{operações} e a \emph{propria} alma, refere-se á \emph{existencia} das operações e da alma; porque a percepção da natureza das operações e da alma pertence á \emph{razão} e é coisa muito diflicil. (\emph{De} Ver., q. 10, a. S ad 8).\par Fizemos nós uma demonstragio ?"7Nào. Limitamo-nos a uma. simples --- explicação. A demonstração não é possivel, quando se.iratz dos fundamen- Los de toda. a sciencia. --- A \emph{veracidade} ou \emph{legitimidade} da consciencia. é de evidencia \emph{immediata} e não pode ser demonstrada nem contestada sem \emph{petição de principio}. Quem quizesse defendel-a ou conlestal-a, já a suppunha como \emph{cería}. Por quanto, quem aflirma ou nega que neste momento pensa, já admilie a \emph{existencia} do pensamento, porque alfirmar ou negar é pensar. Ora admittir a \emph{existencia} do pensamento é admittir a \emph{veracidade} da consciencia. Logo a \emph{veracidade} ou \emph{legitimidade} da consciencia não pode ser demonstrada nem contestada sem \emph{petição de principio}. \textbar{} ME\par Para 2 solução das. principaes difficuldades:cgütra: a: veracidade: da co \emph{Sciencia}, fazemos as seguintes advertencias. --- 1º.) Para que estejamos certos do testemunho, ou acto da \emph{consciencia}, não precisamos de outra consciencia. - Consistindo ella no poder de reflectir e tendo por objecto o que se passa no intimo do nosso espirito, pode ella dobrar-se sobre Si mesma e perceber O proprio acto. Por isso, para a certeza da consciencia, não é necessaria uma nova faculdade, mas é sufficiente um novo acto da mesma faculdade, --- 2º.) Muitas vezes duvidamos se um phenomeno interior deriva: d'esta ou d'aquella faculdade; mas a duvida não se refere á \emph{existencia}, mas sim á \emph{natureza} do mesmo phenomeno. Mas conhecer a \emph{natureza} de um pbenomeno, e por isso a causa d'elle, é munus da razão, não da \emph{consciencia}. --- 3º.) Quando os loucos e os sonhadores affirmam que vêm e ouvem o que realmente não existe, esses enganos devem attribuir-se --- não a consciencia, que se limita a referir o que se passa no espirito dos loucos e dos sonhadores, e nisso ella é veraz, --- mas sim á falta do uso legitimo do \emph{raciocinio}, porque o juizo relativo á correspondencia entre os phenomenos internos e os externos pertence á \emph{razão}, e nos loucos e nos sonhadores a razão está desarranjada ou não funcciona.\par Pergunta se: Perceberá a \emph{consciencia todas} e cada uma das nossas operações espirituaes, de modo que não possa haver uma só operação \emph{inconsciente} ? --- Uns respondem que \emph{fofa} a operação espiritual é necessariamente e sempre percebida pela \emph{consciencia}. Por isso dividem a consciencia em \emph{dzrecia} € \emph{reflexa}. \emph{A directa}, seria a \emph{presença} da alma nos seus actos; e, como \& alma está sempre presente nos seus actos, concluem que não pode haver acto, que escape. á consciencia. Escreve \emph{Balmes}: « Dizer que os phenomenos existem na nossa alma e que não os tem presentes, é uma contradiegão » (\emph{Phil}. \emph{fund}., 1. I, c, 23). --- Outros, porém, e com razão, admittem que pode haver e ha operações internus, de que não temos consciencia, Antes de tudo, repellem a distinccio da \emph{consciencia} em \emph{directa} o \emph{reflexa}. A consciencia \emph{directa} seria uma \emph{consciencia nio-consciencia}. O que constitue \emph{propriamente á consciencia} é a \emph{adeertencik}. Se não ha \emph{adoertencia}, não ha, nem \emph{pode}\pgn{180}haver \emph{consciencia}. A presença da alma nos seus actos não é consciencia, senão quando e porque \emph{reflecte} sobre elles. Ora, se é verdade que a alma pode \emph{habitualmente} perceber-sc a si mesma nos seus actos, emquanto tem sempre para isso uma grande e natural inéclinag\&o, não é verdade que, de\par \emph{facto}, ella se comtemple.sempre c sem descango. Para isso seria necessaria uma attenção- demasiadamente intensa. e. constante, de. que-a nossa alma não é capaz. Logo um facto \emph{psychico} não é é necessariamente um \emph{estado da consciencia}, ainda que que seja sempre um- \emph{estado da alma}, de que podemos ter consciencia por uma sulliciente reflexão. A idea e a consciencia da idea distinguem- se, como se distinguem a visão do objecto pensado c a visão do sujeito pensante. Podemos pensar, sêm pensar no pensamento,\par Todos admittem a \emph{legitimidade} Ou \emph{veracidade} da \emph{consciencia} quanto ao objecto \emph{directo}, mas alg: runs rejeitam-na ou desvirtuam-na quanto ao objecto\par imilamo-r i iões de \emph{Kant}, \emph{Descartes} e \emph{Hamilton}. \emph{Indirect} i\par \emph{pois--de Oi Ou}" posto” em “duvida; a a \emph{objectividade} . do \emph{pensamento}, admitte : a veracidade da c \emph{conscien} las sus , que esta. se - pode perceber os \emph{phenonwnos internos}, O. eu \emph{qhenomenico}, não pode attingir O \emph{sujeito} d esses phenomenos, o eu-nowmenico. --- A opinião de Kant é falsa em relação á esta ultima parie. Por quanto; pela consciencia, attingimos oS phenomenos internos como \emph{nossos}. Ora vão podemos attingil-os .como \emph{nossos}, sem attingirmos, ao mesmo tempo, o \emph{sujeito} d'elles, o eu-nowmenteo. Éo que a experiencia attesta. Alravez dos nossos pensamentos, affectos. desejos, ha sempre um que pensa, ama, deseja ; atravez dos variados phenomenos, que se succedem, uns 208 outros, \emph{no} nosso interior, descobrimos um sujeito, um noumeno, que, no meio de tantas vicissitudes, permanece senrpre o mesmo e que se chama: eu. --- L5 osse \emph{sujeito}, esse \emph{mowmeno}, esse eu, a consciencia não o percebe apenas como um facto: \emph{physiologico}, mas conhece-o tambem como um \emph{principio} de actividade, como uma causa das nossas operações, como uma verdadeira substancia individua, distincta e separada de todas as as outras substancias; embora, 1 repetimos, não perceba a natureza do eu. .\par \emph{Descartes}, na opinião de alguns criticos, admittiu que conhecemos o \emph{sujeito pensante}, mas que tal conhecimento é fructo, não de uma percepção \emph{tmimnediata}, mas de um \emph{raciocinio}. D'ahi o famoso enthymema: \emph{penso}, \emph{logo existo}. --- Esta opinião importa ema petição de principio. Quem diz: \emph{penso}, diz: eu sou \emph{pensante}, e por isso já admitte o que se diz na conclusão, isto é, a \emph{existencia} do sujeito pensante. E não pode deixar de ser assim. O eu é percebido por uma simples intuição da \emph{consciencia}, e anteriormente a todo o \emph{raciocinio}, que por isso, se é util para o conhecimento da natureza do sujeito pensante, é inutil para o conhecimento da \emph{existencia} do mesmo sujeito. Por quanto, sendo a acção absolutamente inseparavel do agente, não podemos perceber uma entidade sem outra, um modo de ser sem o proprio ser.\par \emph{Hamilton} diz que a consciencia, não podendo perceber o eu, se o não distingue do \emph{não-em}, percebe lambem \emph{directamente} O não-eu e por isso o proprio \emph{mundo}. --- Esta \emph{opinião} não pode acceitar-se, porque contraria \& propria \emph{consciencia}, alem de ser contraria á \emph{razão}, --- É verdade que a per-}

%%>>> t163

\pg{177} \pg{179} \tese{163}{Regras para a consciencia.}
As principaes são:

1º --- A \emph{consciencia} deve funccionar dentro dos limites do seu \emph{objecto}, que é constituido pelos phenomenos \emph{internos} do

\emph{z}

€------ \pg{180}  nosso espirito. Por isso não podemos dizer que temos \emph{consciencia} das leis physicae, do mundo, de \textbf{Deus}. \textbar{}

H*-A \emph{consciencia} deve referir exclusivamente a existencia

OD dos \emph{phenomenos} e do \emph{sujeito}, e não a natureza d'este ou d'aquelles; o conhecimento da natureza das coisas pertence á \emph{razão}. \emph{pu}

%%>>> t164

\pg{181} \tese{164}{Memoria intellectual.}
A memoria \emph{intellectual} é a --- mesma intelligencia,.emquanto \emph{conserva}... as \emph{proprias} operações. --- \emph{O passadas}, e as \emph{reconhece} como \emph{proprias} e como \emph{passadas} (*X. : =

%%>>> t165

\tese{165}{Objecto da memoria intellectual.}
O \emph{objecto} da me- = moria intellectual é triplice: \emph{directo}, \emph{indirecto} e \emph{implicito}. --- \emph{e} O \emph{directo} é constituido pelas operações \emph{intellecinaes passadas}, - \textasciicircum{} como taes; --- o \emph{indirecto} é formado pelos \emph{objectos} d' aquellas .. = Operações; --- O \emph{implicito} é O nexo existente entre as operações “---. \emph{intellectuaes} e O \emph{sujeito}, que as \emph{reconhece} como \emph{proprias} (\textasciicircum{}). "\emph{ah}

%%>>> t166

\tese{166}{A memoria intellectual não se engana na percepção Ps das proprias operações passadas.}
4 memoria \emph{intellectual} não Z se engana na percepção das \emph{proprias} operações \emph{passadas}, se O. uengano não é possivel --- nem quanto á conservação d essas : 7\textasciicircum{} operações, --- nem: quanto. ao” seu “\emph{reconhecimento}: -Ora.-o :engano-- = não é possivel ---nem- "quanto \&-conservação-d' essas operações; ---- porque tal. conservação, que é um dote de todas as faculdades, \emph{e} sobretudo espirituaes, por excluir todo o juizo, não está sujeita \textasciicircum{}: --- a erro, --- nem quanto ao \emph{reconhecimento} das operações como \emph{I proprias} e como \emph{passadas}, porque tal reconhecimento importa uma simples \emph{apprehensão comparativa}, pela qual percebemos a " \emph{identidade} entre a operação, que \emph{reproduzimos}, e a operação, \emph{Ps}

' cepc£o, do no-eu completi a percepção do \emph{eu}; mas é falso que ambas as ------ - --- percepções possam pertencer e pertençam a mesma faculdade. O não-eu é i percebido pelos sentidos \emph{externos}, c não pela consciencia. A alma é uma, mas os seus meios de informação são muitos e diversos, conforme a multi- \textasciicircum{} plicidade e diversidade dos objectos perceptiveis.

\footnote{E. pois, pela \emph{memoria}, que adquirimos o conhecimento da nossa \emph{existencia passada}, da \emph{identidade pessoal}, e \emph{da continuação da existencia}. Por quanto, não podemos lemhrar-nos \emph{de} que pensámos um \emph{dia}, sem conhecermos, ao mesmo \emph{tempo}, \emph{que} exisliamos nesse dia, que eramos então o que somos \emph{hoje}, e que temos \emph{continuado a existir} até ao presente. E M} A memoria \emph{intellectual} não é uma faculdade distincta da \emph{intelligencia}; porque duas faculdades são distinctas entre si, quando \emph{o objecto} de uma. : D € distincto do objecto de outra; ora o objecto da memoria \emph{intellectual} É --- identico 2o da. \emph{infelligenciam} ---2 verdade immaterial e intelligivel. (Cf. Sum. ! \emph{Th}.. p. E q. 79, a. 7). --- Como se vé, a \emph{memoria} e a \emph{consciencia} concordam \emph{O}. em que ambas têm por objecto as nossas modificações ou operações espirituaes: mas diferem em que a \emph{consciencia} refere aquellas operações como presentes, 20 passo que a \emph{memoria} refere-us como \emph{passadas}. Por isso pode dizer-se \emph{e} diz-se que i memoria é uma \emph{continuação da consciencia}.

que \emph{produzimos no passado}; e \textbf{esta} apprehens\&o com parativa 0\&0 pode ser. por si falsa. Logo a \emph{memoria intellectual} não se engana na percepção das \emph{proprias} operações \emph{passadas} \footnote{A \emph{memoria intellectual} vão só não se engana na percepção do seu objecto \emph{directo}, mas' nem na percepção do seu objecto \emph{indirecto} e do \emph{implicito}; porque tambem \emph{a} percepção d' estes objectos reduz-se a uma simples \emph{apprehensão}, em que não pode baver erro ou engano. --- Por isso, assim como não se engana \emph{a consciencia}, assim tambem não se engana'a memoria, que, como dissemos, é apenas uma \emph{continuação} da \emph{consciência}.\par \emph{A veracidade} da memoria é uma coisa absolutamente necessaria para a nossa vida. Se a memória não fosse veraz, nenhuma arte seria, possivel, nenhuma sciencia, nenhuma historia, nenhuma lingua, nenhuma lieção, etc. Ora o Auctor da natureza dá o que é necessario ás suas creaturas.\par Tambem. a. \emph{Zegitimidadoe..ou}. veracidade.da. \emph{memoria}, é. de uma. evidencia immediata, e não pode ser posta em duvida sem petição de \emph{principio}. Quem duvidasse d' essa \emph{legitimidade}, deveria formar um raciocinio: Ora todo o raciocinio suppõe a legitimidade da memoria; porque, quando de umas premissas deduzimos uma conclusão, essas premissas são. coisas passadas e não estão presentes no espirito senão pela memoria. --- É por causa d'esse nexo da memori, com o \emph{raciocinio} (cujas \emph{premissas} a memoria torna \emph{presentes}) que a. \emph{Logica} defende a. \emph{veracidade} ou \emph{legitimidade} d'essa faculdade.\par D'onde se segue que se, ás vezes, nos enganamos com relação ao \emph{passado}, o Engano não deve attribuir-se á memoria, mas unicamente á falta de criterio e de reflexão; porque ou se attribue á \emph{memoria} o que ella não refere, ou se considera como \emph{certo} o que a \emph{memoria} refere como \emph{duvidoso}. --- A memoria pode enfraquecer, pode mesmo esquecer as coisas; mas o esquecimento não'é \emph{erro}, assim como não é erro a ignorancia, e no que ella evidentemente refere como \emph{certo} não pode enganar-se. aliás deixaria de \emph{ser memoria}, isto é, deixaria de ser a faculdade de conservar e reconhecer o \emph{passado}. --- Aqui notamos que a \emph{memoria intellectual}, destinada para conservar as imagens immateriaes das coisas, é mais tenaz e menos movel do que a \emph{memoria sensitiva}; porque o orgão corporeo, a que está ligada a memoria \emph{sensitiva}, é ais movel, mais mutavel do que uma faculdade \emph{espiritual}, como é a memoria \emph{Zafellectual}, (\emph{De} Ver., q. 10, a. 2).\par Como a memoria pode esquecer e esquece o passado, é necessario robustecel-a, exercitando-a continuamente, e auxilial-a, applicando a \emph{atfencao} ás coisas que desejamos reter, \emph{e} ordenando as ideas de modo, que a recordação de uma traga naturalmente a recordação da outra. --- Os principaes \emph{dotes} de uma bóa \emph{memoria} são: \emph{facilidade}, \emph{tenacidade}, \emph{fidelidade}, \emph{pr j tidao}. \emph{É facil}, quando sem esforço recorda os conhecimentos passados; --- é £\emph{enac}, se os conserva por \emph{muito} tempo; --- é \emph{fiel}, se os reconhece com exauctidàao, --- é \emph{prompta}, se \emph{os reproduz} com celeridade.}.

%%>>> t167

\pg{182} \tese{167}{Regras. para. a. memoria intellectual.}
Para que a \emph{memoria intellectual} seja um meio seguro para a acquisição da verdade, devem observar-se, as seguintes regras:

1º --- Sendo uma continuação da \emph{consciencia}, a memoria \pg{183}  deve limitar-se a reproduzir as percepções ou operações \emph{passar das}, sem determinar as causas das coisas.

9» --- Devemos affirmar só o que a memoria reconhece, e “no modo, por que ella-o-reconhece... No- caso. contrario, 0. erro. . proviria --- não da \emph{memoria}, que não engana no que refere, --- mas da imprudencia ou precipitação nos juizos.

ArtTIGO HJ

\emph{Sentidos externos e internos}

%%>>> t168

\tese{168}{Sentidos externos.}
Seutidas \emph{exíernos} São os que existem na“ superficie \emph{externa} do nosso organismo e por cujo intermedio percebemos os seres \emph{saleriaes} e \emph{externos} (t).

.169. Numero dos sentidos externos. --- Os sentidos exfermos são \emph{cinco}: --- a \emph{vista}, pela qual percebemos as \emph{cór} ouvido, pelo qual percebemos" 68 sons, --- o \emph{olfacto}, pelo qu percebemos os \emph{cheiros}, --- "o \emph{gosto}, pelo-qual-percebemos- OS \emph{sabores}, --- o \emph{facto}, pelo qual percebemos os objectos externos, que nos impressionam pela sua resistencia \footnote{\emph{O} homem não só conhece as coisas invmateriaes, mas tambem as \emph{materiaes}. Para a percepção das coisas immateriaes, é dotado de faculdades \emph{immuteriaes} ; para, a percepção das coisas mnateriaes, é dotado de faculdades materiaes. A faculdade \emph{material} não reside exclusivamente na alma, nem exclusivamente no \emph{corpo}, mas no \emph{composto} de aima e de \emph{corpo}. --- ÀS faculdades perceptivas das coisas materiaes,emquanto \emph{materiaes}, SãO OS sentidos, que se distinguem em externos e internos.\par 9) Na \emph{Anthropologia} trataremos. com sufficiente desenvolvimento, da natureza e do numero dos sentidos \emph{externos}. e do objecto \emph{proprio} de cada um. Aqui, para o nosso escopo, bastam estes apontamentos. --- Tratamos das faculdades \emph{perceptivas}, não para investigarmos o modo, por que produzem as suas operações, mas para defendermos a sua veracidade ou \emph{legitimidade} nà percepção dos respectivos objectos.}. ao

%%>>> t170

\tese{170}{Objecto dos sentidos externos.}
O \emph{objecto} dos sentidos externos é \emph{triplice}: \emph{proprio}, commum e concomitante. --- É \emph{proprio}, quando é percebido por um só sentido; assim a côr é o objecto \emph{proprio} da vista, o som do ouvido, etc. --- É commum, quando é percebido por \emph{mais de} um sentido; assim a \emph{grandeza} é objecto commum, porque é percebida pela vista e pelo \emph{tacto}. --- É concomitante, quando não é percebido por si mesmo, mas, por meio de outro objecto, a que está naturalmente unido; assim \pg{184}  a \emph{substancia corporea} é objecto \emph{concomitante}, porque é percebida por meio das \emph{qualidades corporeas}, que a envolvem \footnote{Cora relação ao \emph{íriplice} objecto dos sentidos fazemos as seguintes advertencias. --- O'objecto"proprio-é- uma \emph{qualidade singular} e \emph{concreta}, que \emph{por si} e \emph{primeiramente} actua num só sentido, de modo que não pode actuar em nenhum dos outros, Os objectos \emph{proprios} são \emph{cinco}, tantos quantos ---\textbar{} os sentidos.---- O objecto commum é tambem uma \emph{qualidade singular} e \emph{concreta}, \textasciicircum{} que modifica, de algum modo, a acção dos sensiveis \emph{proprios}, é por isso, se actua nos respectivos sentidos por \emph{si}, não actua \emph{primeiramente}, porque não pode impressionar os sentidos senão. por meio dos sensiveis \emph{proprios}; assim a \emph{grandeza} é objecto \emph{commum} da \emph{vista} e do \emph{tacto}, porque refere-se á \emph{vista} por causa da \emph{cór}, \emph{c ao facto} por causa da \emph{resistencia}, Os objectos communs, como dissemos, são \emph{cinco}: O \emph{movimento}, \emph{a quielagao}, \emph{O numero}, \emph{a figura} e a \emph{grandeza}; mas, como diz S. Thomaz, todos reduzem-se a \emph{quantidade} (Cf. Sum. \emph{Th}. p. I. q. 78. a. 3) --- o objecto \emph{concomitante}, que os escholasticos chamam \emph{per} accidens, é uma coisa, que está naturalmente unida com os 'objectos \emph{proprios} ou communs, dé modo que não pode deixar de ser percebida na percepção d'esses objectos ; ; tal é, como dissemos, a \emph{substancia corporeu}, que, sendo o \emph{substracinm} dos sensiveis \emph{proprios} e communs, é perce- "bida por elles e com elles. A razão é que os sentidos percebem os sensiveis \emph{proprios} e comuns --- não de um modo \emph{abstracto}, --- mas de um modo \emph{concreto}; pois a vista não “percebe \& \emph{cor}, mas a \emph{objecto colorido}, o \emph{gosto} não percebe o \emph{sabor}; mas o \emph{objecto saboroso}, ete. É claro que os sentidos limitam-se a attestar o \emph{facto} de uma \emph{existencia} concreta, sem penetrar na essencia.}\footnote{A razão da natureza e do nmunero d' essas condições dedouz-se da propria natureza da percepção \emph{sensitiva}. Faz-se esta percepção, emquanto o objecto actua na faculdade e nella imprime uma sua especie, ou imagem, ou representação, a qual determina a faeuldade a percebelo como é em si mesmo. Portanto, tudo o que pode concorrer, de algum modo, para viciar ou alterar tal impressão, concorre tambem para viciar ou alterar a propria percepção, Para evitar todo e qualquer inconveniente, é indispensavel a verificação das condições \emph{enumeradas}. --- O orgão do \emph{sentido deve achar-se}\pgn{185}córes, nem um surdo ouve os sons), nem: estar: \emph{doente} (o gosto de-um- febri citante sente amarga-uma-coisa. doce), nem--deve---ser- \emph{alterado} (as mãos --- enregeladas são incapazes de apalpar e distinguir objectos delicados). --- O \emph{objecto} deve ser --- \emph{proprio}, isto é, deve referir-se a faculdade que é capaz de o perceber, --- e \emph{proporcionado} ao \emph{orgão}, isto é, não deve estar nem muito perto nem muito longe (a superficie do sol, da lua, por estar fora do alcance da vista, parece plana). --- O \emph{meto}, atravez do qual o objecto actua no sentido, não deve ser variado ou perturbado na sua condição natural (assim, pela interposição dos vapores, que refrangem os raios da luz, o sol. a lua e as estrellas parecem, ás vezes, quasi vermelhas; num pau, merjulhando nà agua, parece quebrado) --- A experiencia e a reflexão podem julgar Se as condições enumeradas se verificam, ou não. -}.

%%>>> t171

\tese{171}{Condições para a funcção dos sentidos.}
As condições, necessarias e- sufficientes, para a legitima funcção dos sentidos referem-se aos proprios \emph{sentidos}, ao \emph{objecto}, e ao meio. --- Da parte dos \emph{sentidos}, 6. necessario que se achem no seu \emph{estado} normal. e por isso estejam isentos de vicios, de doenças, e de qualquer affecção, que possa impedir ou perturbar o.exercicio da percepção. --- Da parte do \emph{objecto}, reguer-se que seja --- \emph{proprio}, de modo que se refira ao respectivo sentido, --- e \emph{proporcionado}, de modo que se encontre em condições convenientes para ser percebido. --- Da parte do meio (existente entre o sentido e o objecto), é preciso que tenha a capacidade de levar ao sentido, inalterada, a imagem do objecto (º).

a) Toda a faculdade alcança necessariamente o objecto, para o qual está \emph{naturalmente} ordenada, se não encontra obstaculos; porque, se não alcançasse..o.seu objecto, a. faculdade estaria e.. deixaria de estar ordenada para elle: o que é absurdo. Ora os: sentidos \emph{externos} estão \emph{naturalmente}. ordenados para perceber os objectos \emph{externos}, como o demonstra --- não só a sua situação na superficie \emph{exterma} do nosso corpo, --- mas tambem a sua \emph{adaptação} aos respectivos objectos \emph{externos}. Logo os sentidos \emph{externos} percebem os objectos \emph{externos} \footnote{A situação e a \emph{adaptação} (ou \emph{conformação physiologica}) dos orgãos dos sentidos \emph{externos} revelam-nos as intenções da natureza. --- PrimoiraméOnte a sifuacao. Se a natureza nos tivesse dotado de sentidos só pira conhecermos ás nossas disposições \emph{subjectivas}, deveria tel-os collocado na centro do organismo ou na cavidade do craneo. Mas não foi assim. Os orgãos dos sentidos \emph{externos} estão co!locados na superficie \emph{externa} do nosso corpo ou na extremidade dos membros. E para que? Para podermos relacionar-nos com \emph{o} mundo exterior. Se não tivesse sido este o escopo da natureza, teriamos um effeito sem razão sufficiente: o que é absurdo. Logo a disposição dos orgáaos'dos sentidos \emph{externos} manifesta a \emph{objectividade} das nossas percepções, isto é, a existencia dos \emph{objectos} percebidos. --- Secundariamente a. \emph{adaptação}. Todos os \emph{physiologistus}, examinando os orgãos dos sentidos \emph{externos}, reconhecem que elles estão maravilhosamente \emph{adaptados} para as respectivas funceções, isto é, para a percepção dos respectivos objectos \emph{exte} nos, --- a \emph{vista} para z percepção das \emph{córes}, 0 \emph{ouvido} para a percepção dos \emph{sons}, ete. Ora, se os sentidos \emph{externos} têm essa \emph{adaptação}, devemos concluir que \emph{elles percebem} os \emph{objectos externos}, se não houver algum \emph{impedimento}. \emph{Não pode negar-se a conclusão}, \emph{Sem} se negar o senso commum.}.

b) Todos os homens, de todos os tempos, de todas as classes e condições, á excepção dos loucos, estiveram e estão convencidos de que os sentidos \emph{externos} percebem --- não a \pg{186}  modificação do proprio órgão, --- nem uma coisa \textbf{imaginaria}, --- mas um objecto existente. na \emph{realidade}, fóra de nós. \textbf{Ora esta} convicção, sendo universal e constante,não pode \textbf{derivar senão} da natureza; e esta não pode enganar-se: Logo os \textbf{sentidos} \emph{externos} percebem os objectos \emph{externos} (i).

173. Os sentidos externos \textbf{percebem os objectos externos} por um \textbf{modo} immediato. o ,

a) Os sentidos \emph{externos} perceberüo os objectos \emph{externos}

(\emph{t}) Todos, mesmo .os adversarios, admittem a. existencia d'essa \textbf{convicção} universal c constante ácerca da \emph{objectividade} da percepção sensitiva \emph{externa}. Todavia, alguns dizem. com Taine, que a tal convicção é. uma \emph{sfusão verdadeira}, para. a distinguir das \emph{illusões falsas}. --- Mas, se toda a \textbf{illusão} é um engano, é uma falsidade, poderá haver uma \emph{illusão verdadeira}, \textbf{isto} é, uma \emph{falsidade verdadeira}? Ao bom senso a resposta. --- E não pode ser de outro modo. Quando um pobre homem: recebe uma bengalada na \textbf{cabeça}, \textbf{lanca-se} contra o adversario: e-faz todos-os esforços: para lhe arrancar das \textbf{mãos} a bengala. causa instrumental dá sua dor. O que mostra isso? Mostra que os sentidos \emph{externos}; isto é, os olhos do agredido não perceberam \textbf{uma coisa} imaginaria, ou uma modificação do proprio orgão, mas um objecto \emph{real}, um objecto \emph{externo}. --- O exemplo adduzido é vulgar; mas não faltam \textbf{outras razões} scientificas. É admittido por todos que-os nossos sentidos \emph{externos}, \textbf{como} todas ás nossas potencias ou faculdades, podem perceber um ou outro \textbf{objecto}, indifferentemente. Quando, pois, os sentidos percebem um determinado objecto, e não outro, é preciso dizer que elles foram determinados para. perceber \textbf{esse} objecto, e não outro. E d'onde, derivou esta determinação? A \textbf{determinação} podia derivar --- ou dos proprios sentidos --- ou da vontade, --- ou dos \textbf{objectos} \emph{externos}. Mas não derivou dos proprios sentidos; porque esses são por si indifferentes e indeterminados com relação aos seus objectos (se \textbf{fossem determinados} para perceber um objecto, e não outro, perceberiam \textbf{sempre} e só esse objecto). Nem derivou da nossa vontade; porque muitas \textbf{vezes os} sentidos percebem um objecto independentemente e contra o desejo \textbf{da nossa} vontade, Logo, se n\&o.quizermos admittir um effeito sem uma razão \textbf{sufficiente}, devemos dizer que essa determinação dos sentidos para que \textbf{percebessem} um certo objecto, e não outro, derivou do proprio objecto \emph{externo}. --- O \textbf{objecto} externo determina um sentido, emquanto, actuando no Proprio \textbf{sentido}, imprime nelle a sua. imagem, a qual. por isso, é uma qualidade \emph{accidental} do sujeito recipiente. O sentido, em seguida, por uma reacção vital, \textbf{appropria}-se d'essa imagem, e nella percebe o proprio objecto \emph{externo}, que a produziu. --- O que dizemos aqui, em resumo, será desenvolvido na \emph{Antkropologia}.

Objectam : A percepção dos sentidos \emph{externos} é um facto \emph{subjectivo}. c por isso não, pode referir o que está fóra do sujeito, o ndo-ex. --- Responemos que a percepção dos sentidos \emph{externos}, se é um [acto \emph{subjectivo} com relação ao seu \emph{principio}, que é \emph{a} faculdade, é \emph{objectivo} com \textbf{relac\&o} ao \emph{termo}, que é a coisa percebida... Por outras palavras: A percepção, \emph{considerada}. \emph{entitativamente}, \emph{é um facto subjectivo}; --- considerata. \emph{infencionat} ou \emph{representativamente}, \emph{é um facto objectivo}.

por um modo \emph{immediato}, se entre a acção, ou actuação, d'esses objectos nos respectivos sentidos e a paixão, ou impressão, recebida pelos mesmos sentidos, não existir meio algum.: Ora, effectivamente, entre..essa. acc\&0. e.essa. paixão não-.existe melio. algum; porque a acção do agente é recebida por si mesma no

\emph{A} paciente, e na acção e pela acção, que reflecte as propriedades - * do agente, este é percebido pelo paciente. Logo os sentidos \emph{externos} percebem os objectos \emph{externos} por um modo \emph{immediato}. .

b) Se os sentidos. \emph{externos} não percebessem os objectos -\emph{externos} por um modo. \emph{immediato}, deveria dizer-se que esses sentidos percebem --- \emph{primeira} e \emph{directamente} a propria: \emph{paixão}, isto é, a impressão subjectiva, recebida pela acção dos agentes externos, --- e \emph{segundaria} e \emph{indirectamente} a mesma \emph{acção} dos agentes externos. Ora a experiencia attesta --- que os sentidos B \emph{externos} percebem unica e \emph{directamente} a \emph{acção} d externos, --- e que a \emph{impressão} subjectiva dos m E não é percebida por-estes-de-modo-algum;-mas-exclusivamente pelos sentidos \emph{internos}. Logo os sentidos \emph{externos} percebem os objectos \emph{externos} por um modo \emph{immediato} (f). .\footnote{Na \emph{Anthropologic} provaremos que pelos nossos sentidos \emph{externos} per-. cebemos, de um modo \emph{immediato}, os objectos \emph{externos}. Limitamo-nos aqui a algumas reflexões necessarias e sufficientes para o nosso escopo. --- É principio inconcusso que \emph{o} acto do agente é recebido no paciente, --- e que o acto do agente (emquanto recebido no paciente) e o acto do proprio paciente não constituem senão um e o mesmo acto, dotado de um duplice aspecto, de um duplice caracter; porque esse acto --- relativamente ao agente, de que deriva, tem a natureza de \emph{principio}, --- e relativamente ao paciente, que o recebe, tem a natureza de \emph{fermo} e de \emph{paixão}. Dix S. Thomaz: « Unus et idem actüs ” sensibilis et sentientis, sed ralione dilferunt» (De Anima. L. YII, lect. 2). Recebendo em si mesmo o acto do agente externo, o paciente percebe o néo-eu no eu, e refere naturalmente a esse «d0-e« 0 acto, que o impressionou; visto que toda a \emph{paixzo} se refere necessariamente a uma acção. D'ahi a \emph{objectividade} da percepção. --- E não é Ludo. Se o aclo do agente é recebido e existe realmente no paciente, se a acção e a \emph{prizzo} formam um e o mesmo acto, com dois aspectos differentes, é claro que o paciente deve perceber \emph{directa} ou \emph{immediatamente}, sem intermediarios e antes da propria modificenção subjectiva, o acto do «\emph{do-eu}, \emph{do} agente externo. É, pois. inutil procurar uma \emph{ponte}, que nos sirva para passar do ex para o \emph{sdo-eu}; visto que entre o eu e o não-eu não ha distancia --- Tal é a opinião universal dos homens. Todos estão convencidos --- de que. \emph{percehendo} o mundo exterior, percebem \emph{o} que está fóra d'elles --- e de que o \emph{objecto directo} e \emph{inmediato} das suas percepções é \emph{constituido} pelos \emph{proprios objectos} externos.\par \emph{Objectam} :,\emph{O corpo} não pode \emph{actuar na} alma, nem esta \emph{naquelle}. \emph{Logo}\pgn{188}uma percepção \emph{immediata} dos corpos é impossivel. --- - Resposta. Concedemos\par que os corpos não podem operar na alma, quando separtda do organismo; mas negamos que não podem operar nella; quando unida ao organismo. o objecto “sensivel refere-se ao orgão \emph{ánimado}, como o agente ao paciente. E como o agente, \emph{pela} sua acção, existe no \emph{paciente}, aquelle pode e devo ser percebido por este, e de um modo \emph{immediato}, -\par *\par Na percepção do sentidos \emph{exteritos} poderá haver: engano ? Serão elles\par verases no que nos attestam ? --- Eis uma questão, que é a base de toda a theoria ácerca do conhecimento intellectual. Por quanto, dependendo o conhe- cimento intellectual do conhecimento sensitivo quanto ao \emph{objecto} (pois são os sentidos. que. subministram. a intelligencia o \emph{material} do pensamento), é claro; que; se se enganam -os sentidos, se engana tambem a inteligencia, --- 'Tratemos, "pois; da-veracidade-dos-sentidos \emph{externos}. --- -\par Antes de tudo, devemos advertir que nos sentidos não se encontra a verdade, tomada no seu conceito \emph{formal}, isto é;.no seu conceito proprio, emquanto importa uma “\emph{adequação} entre. a faculdade perceptiva e o objecto percebido. Com effeito, para que. nos sentidos, podesse encontrar-se a verdade; tomada no seu conceito proprio, deveriam elles perceber a \emph{relação} de identidade ou de discrepancia entre o sujeito e o predicado, e, para isso, deveriam: possuir varias ideas universaes, sobretudo a idea universal da forma, em que os dois termos convém. ou não convém. Ora as faculdades materiaes. como são os sentidos, não podem perceber a relação entre o sujeito e'o predicado, porque tal relação é uma coisa \emph{abstracta} ou uUniversal e essas faculdades não são capazes de ideas universaes: Logo nos sentidos: não se, encontra a verdade, . tomada no seu conceito proprio ou rigoroso. --- Se nos sentidos n\&o existe a verdade, tomada no seu conceito proprio ou rigoroso, nem se: encontra nelles o \emph{erro}, ou a \emph{falsidade}, tomada tambem no seu conceito proprio ou rigoroso, emquanto exprime uma \emph{desconformidade} entre a faculdade perceptiva \emph{e}.9 objecto percebido. Porquanto. dá-se essa \emph{desconformidade}, quando a faculdade perceptiva affirma --- que a um sujeito - convem um. predicado. que Ihe não convem, --- ou que a esse sujeito não convem um predicado, que lhe convem. Ora esta alfirmação, suppondo a percepção da relação entre o predicado e \emph{o} sujeito, não pode ser feita pela faculdade sensitiva Logo nos sentidôs não se encontra a falsidade, tomada \emph{no} seu conceito proprio ou rigoroso. --- Em poucas palavras: A verdade e a falsidade encontram-se propria e rigorosamente no \emph{juizo}.' Ora os sentidos não são capazes de pronunciar um verdadeiro juizo. Logo a verdade e a falsidade não se encontram propria e rigorosamente nos sentidos. Diz S. Thomaz: « Proprio foquendo, veritas est in intellectu componente \emph{et dividente} (\emph{h}. \emph{e}. \emph{judicante}), \emph{non} autem \emph{in} sensu, neque \emph{in} intellectu \emph{cognoscente quod} «\emph{quid est} (\emph{h}. e. \emph{simpliciter apprehendente essentiam}}

174. Os sentidos externos, \textbf{quando bem dispostos} e \textbf{convenientemente} applicados, não podem \textbf{enganar-se nà} percepção do objecto proprio.. \textbar{} U

a) Se-os-sentidos-ezlernos; quando bem dispostos e convenientemente applicados, podessem enganar-se na percepção do

\emph{T}

objecto \emph{proprio}, --- a \emph{vista} na percepção das \emph{cores}, O0 avido na 7 percepção dos \emph{sons}, etc., --- deveria dizer-se que elles são \textbf{naturalmente} incapazes de alcançar o objecto, para o qual \textbf{natural}- . mente tendem. Ora é absurdo dizer que uma-faculdade-é-natu- ----------- \emph{Ho}. ralmente incapaz de alcangar o objecto, para o qual \textbf{naturalmente} tende; pois toda a razão do ser e da natureza de uma : \emph{L}.

rei)» (\emph{Sum}. \emph{Th}., p. I, q. 16, a. 2). --- Se, algumas vezes, costumamos dizer P que os sentidos \emph{julgam}, emquanto distinguem \emph{um} sensivel de outro sensivel, que s\&o objectos do mesmo sentido ou de sentidos differentes, tomamos o ---. \emph{juico} em significação \emph{lata} ou menos \emph{rigorosa}. \emph{st}

Mas a verdade e a falsidade, se não se encontram \emph{formalmente} na facul- ---.--- ---dade sensitiva; podem encontrar-se nella \emph{materialmente}, enquanto a appre- ' Me hens\&o ou percepção dos objectos sensiveis corresponde, ou não, a realidade, \emph{AN} isto é, emquanto percebe, ou não, os entes objectivos, como são em si \textbf{mesmos}. cT D'abi a pergunta: Enganam-se, ou não, os sentidos \emph{externos} na percepção \emph{L}3 do seu objecto? --- E como esse objecto é triplice, tresdobra-se a pergunta: p .Enganam-se, ou não, os sentidos \emph{externos} na percepção do seu objecto \emph{pro}- " \emph{prio}, commum e \emph{concomitante}? --- A estas perguntas vamos \textbf{responder}:------------------------7 Lan

Repetimos que a nussa -não-é-uma \emph{demonstração}; muito menos-uma-- - z demonstração \emph{rigorosa}: é apenas uma declaração, uma \emph{resposta} aos \textbf{adver}- \emph{OT}. sarios. Por quanto á \emph{veracidade} dos sentidos externos é tão clara, tão imme- : vá dialamente evidente, que \emph{não precisa}, nem \emph{pode} ser demonstrada. --- \emph{Não pre}- . \emph{T cisa}. Na verdade, o senso commum appella para o testemunho dus sentidos \emph{Ll externos}, como para a prova irrefragavel da \emph{evidencia} de uma coisa. D'ahi \emph{us} as expressões : toca-se com as mãos, mette-se pelos olbos dentro, etc. --- --- \emph{Não pode}. Com effeito, toda a demonstração parle, em ultima analyse, de \emph{mnn} principios \emph{wnécersaes}; e estes são conhecidos pela inteligencia. Ora a ope- ---. ração da intelligencia é precedida pela operação dos sentidos \emph{externos}. Logo ur toda a demonstração presuppõe a \emph{veracidade} dos sentidos \emph{externos}. Ox

Esta convicção ácerca da veracidade dos sentidos (como tambem de. \emph{veracidade} de todos os outros \emph{meios}) não é elfeito de um \emph{cego instincto}. Na verdade, o \emph{instincto} é uma inclinação, que deriva da natureza e pela qual operamos sem um conhecimento sufliciente. Por isso, quando se trata de uma faculdade \emph{perceptiva}, O \emph{instincto} deve conceber-se come uma inclinação, proveniente da natureza, para apprehender um objecto, quando o proprio objecto não impressiona por si a faculdade. Por quanto, se esta impressão se realisa, a faculdade inclina-se para produzir \emph{o} seu acto em virtude da sua energia essencial; e tal inclinação não é, nem pode chamar-se \emph{instincto}. Assim, ninguem dirá que é o \emph{instincto} que nos leva a julgar que \emph{o todo é maior} que uma sua \emph{parte}; mas todos dirão que é a luz da razão, a energia intrinseca da nossa intelligencia. A necessidade, que, nestes casos, obriga 3 faculdade a affirmar a verdade, deriva do concurso de uma duplice causa de ordem exclusivamente cognitiva: da evidencia do objecto e da tendencia da faculdade para,o seu acto. A natural inclinação, pois, pela qual acreditamos no testemunho dos sentidos \emph{externos} (e de todas as outras faculdades perceptivas), é uma inclinação de ordem \emph{cognitiva}, e não uma inclinação \emph{cega} ou \emph{tnstinctiva}.

faculdade consiste na sua natural aptidão para conseguir o objecto \emph{proprio}. Logo os sentidos \emph{externos}, quando bem. dispostos e convenientemente applicados, Dão podem enganar-se na

» Se os sentitos externos, quando bem dispostos e con-

venientemente applicados, podessem enganar-se na percepção do objecto \emph{proprio}, dar-se-hia o engano --- ou porque elles não podem receber a impressão do objecto \emph{proprio}, --- ou porque esse objecto não .\emph{existe} na realidade. Ora nenhuma .d' estas duas hypotbeses pode admittir-se. Não pode admittir-se a primeira; porque a \emph{physiologia}, examinando o organismo. e.a disposição dos sentidos externos, reconhece-os aptissimos para receberem a impressão do respectivo objecto \emph{proprio}. Não pode admittir-se a segunda; porque os sentidos \emph{externos}, sendo por si inditferentes. para .receberem..a impressão de muitos objectos \emph{externos}, não poderiam aceusar a impressão de um certo objecto, se este não existisse na realidade. Logo os sentidos \emph{externos}, quando bem dispostos e convenientemente applicados, não podem enganar-se nu percepção do: objecto \emph{proprio} ()\footnote{Os sentidos \emph{externos}, quando bem dispostos e convenientemente applicados (no que se resumem todas as condições necessarias para o legitimo exercicio d'elles) não se enganam \emph{por si}, isto é, pela sua natureza, na percepção do objecto \emph{proprio}; porque, se se enganassem \emph{por si}, seriam feitos e deixariam de ser feitos para essa percepção: o que é contradictorio. --- Mas, se os sentidos \emph{externos} não podem enganar-se pela sua natureza, podem enganarse por uma causa \emph{accidental}, que lhes não deixa perceber o objecto, como é em si mesmo. Esta causa \emph{accidental} deve procurar-se na indisposição do \emph{organo}, ou na imperíeigüo do \emph{meio}. Por quanto, o sentido, para perceber o seu objecto, deve ser determinado por uma forma ou especie, que é a imagem intencional do mesmo objecto. -Mas, sé o sentido não estiver bem disposto. ou não fôr convenientemente applicado, não poderá receber, como devia, essa forma, ou .a forma não representará fielmente \emph{o} objecto. D'onde o engano ou a falsidade na percepção dos nossos sentidos. (O que dizemos dos sentidos, deve applicar-se a todas as faculdades \emph{passivas}, isto é, determinadas ao acto pela impressão de um agente). Diz S. Thomaz: « Sensus semper apprehendit rem \emph{ut} est, nisi sit impedimentum in \emph{organo} vel in \emph{medio}» (\emph{De} Ver. q. 1, a. 11; of. Sum. \emph{Th}, p. L, q. 17, a. 2). --- Como se vê, o engano ou o erro --- pode dar-se nos sentidos emquanto, por uma causa üceidental, não percebem uma coisa como é em si mesma, --- mas não pode dar-se emquanto accusam o estado proprio, a affeeção propria, embora possam occasionür um juizo falso da parte de uma faculdade superior. Assim, \emph{quando um} doente diz \emph{que o} assucar é amargo, o testemunho do gosto é\pgn{191}veraz, porque efectivamente o doente tem essa sensação, embora se engane por attribuir o amargo ao assucar.\par Se os sentidos podem enganar-se \emph{per} accidens, mesmo com relação ao objecto ou sensivel \emph{proprio}, Se podem não representar e não referir as coisas como são em si mesmas, é claro que podem tambem dar a intelligencia (que se baséa no testemunho d'elles) a \emph{occasizo} de emittir juizos \emph{falsos} ácerca dos objectos externos. --- Todavia a iotelligencia pode conhecer o engano de um sentido e a sua origem (se depende da indisposição do orgão ou da imperfeição do meio), e conhece-o por meio do mesmo ou de outro sentido, ou pela reflexão propria. Conhecido, o engano do sentido corrige-se facilmente. --- Por isso o juizo da intelligencia, que succede immediatamente á percepção dos sentidos, é sempre e necessariamente verdadeiro, quando estamos certos.da bóa disposição dos mesmos sentidos e da sua conveniente applicac\&o. --- Na \emph{Anthropologia} resolveremos algumas difficuldades contra a \emph{veracidade} dos sentidos externos (C£. \emph{De} Ver., q. /1, a. 11).}

--- ap \emph{I} noma \emph{emm s mm rm metn a ntm}

P i --- "MELOS PARA \emph{A} ACOUISIGÉO- DA \textbf{VERDADE} 191

175. Os sentidos externos, quaado bem dispostos e. \textbf{conve}- \textbar{}

nientemente applicados, não podem enganar-se na percepção do objecto commuun. - t.

a) A per cepcào. do. objecto. ou. sensivel commum. \&um.effeito, ---

para o qual estão por si e naturalmente ordenados os respectivos \textasciicircum{}7 sentidos, tomados collectivamente. Ora o effeito, para o qual estão por si e naturalmente ordenadas as faculdades, tomadas singular : ou collectivamente, não pode deixar de se alcançar, quando se verificam as devidas condições; aliás essas faculdades seriam na-turalmente insufficientes e deixariam de ser faculdades. Logo os sentidos \emph{externos}, quando bem dispostos e convenientemente applicados, não podem enganar-se na percepção do objecto commum.

b) Os sentidos \emph{externos}, quando bem dispostos e. convenientemente applieados, não podem enganar-se na percepção do objecto \emph{proprio}. Se não podem enganar-se na percepção do objecto "\emph{proprio}, nem podenret 1ganar-Sena-percêpção-do-objecto-connun 7 pois cada um -dos-diversos-.aspectos,-pelos-quaes--o--objectocommunm se refere a varios sentidos, constitue o objecto \emph{pro}- \emph{prio} de cada um d'esses sentidos. Logo os sentidos \emph{externos}, \textasciicircum{} \textasciicircum{} quando bem dispostos e convenientemente applicados, não podem enganar-se na percepção do objecto commum \footnote{Quando um objecto sensivel é commum a varios sentidos, a cada um dos quaes se refere sob um aspecto differente, ou por um attributo diverso, não pode ser adequada e completamente percebido se não por todos esses sentidos. Então, cada. sentido percebe do objecto \emph{commui} o silributo, que é objecto \emph{proprio} d'esse sentido. Ássim uma grandeza é percebida --- \emph{pela vista} quanto á \emph{côr}, --- e \emph{pelo tacto} quanto á \emph{resistencia}.\par -\pgn{192}D' \emph{este} modo, os.Yarios sentidos constituem uma faculdade \emph{completa} com relação ao \emph{objecto commum}, e não pode haver engano; pois este, se houvesse, deveria attribuir-se ao Auctor da natureza. --- Mas, se não pode haver, na percepção do sensivel comamum, um engano ou erro \emph{per se}, pode haver um engano \emph{per accidens}, quando, para essa percepção, não se empregam todos os "sentidos respectivos ; e assim a Jaculdade é insuflieiente, incompleta. inadequada comr relação ao seu objecto. --- Não admira, pois, que a vista, . não auxiliada \emph{pelo tacto}, se engane, ás vezes, na percepção da \emph{figura}, da ' grandeza, da \emph{distancia}, do \emph{muvimento}; porque estes são objectos communs. Mas, quando se dá o engano, este não pode attribuirse a um vicio ou defeito da \emph{vista}, mas unicamente á imprudeneia de quem julga, --- O engano de um só sentido com .relação a um objecto comtmwm pode corrigir-se pela applicação de um outro sentido, ou pela reflexão da intelligencia. ---\par O engano pode dar-se sobretudo na percepção da vista, emquanto dois objectos, dotados de «diversa grandeza, mas situados em distancias deseguaes, podem \emph{produzir} nesse sentido \emph{a} mesma impressão ; e esta impressão pode dar a intelligentia a occasião de concluir falsamente pela egualdade da grandeza dos dois objectos. Por isso, a percepção da \emph{vista} deve ser interpretada e contida nos seus verdadeiros limites por outro sentido, cuja percepção não varie segundo a variedade da distancia do objecto. Este sentido é o \emph{facto}, que conserva sempre a mesma distancia com relação ao seu objecto \emph{proprio}.}.

176. Os sentidos externos, quando bem dispostos e \textbf{convenientemente} applieados, não podem enganar-se na percepção do objecto concomitante.

a) Todos-os-homens-de-bom-senso estiveram sempre e estão convencidos --- de que os sentidos externos, quando bem dispostos e convenientemente applicados, percebem as qualidades sensiveis dos corpos \emph{esternos}, como são em si mesmas, adherentes. a um sujeito, e por isso percebem, ao mesmo tempo. e.com o mesmo acto, as qualidades e o sujeito, embora aquellas de um modo \emph{explicito} e este de um modo \emph{implicito}, --- e de que esta percepção não pode estar sujeita a engano. Ora tal convicção, universal e constante, não pode deixar de ser effeito da verdade; porque só a verdade é universal e. constante, e todo o effeito exige uma .causa proporcionada. Logo os sentidos \emph{externos}, quando bem dispostos e, convenientemente applicados, não podem enganar-se na" percepção do objecto \emph{concomitante}. -

1b) Se os sentidos \emph{exteríros}, quatido bem dispostos e convenientemente applicados, se enganam na percepção do objecto concomitante, dá-se o engano --- ou. porque. as.qualidades sensiveis, \emph{proprias} ou communs, actuam nos sentidos externos de um modo \pg{193}  \emph{abstracto}, isto é, como que separadas do sujeito, a que adherem e que as sustenta, --- ou porque essas qualidades são percebidas pelos sentidos de um modo \emph{abstracto}. Ora nenhuma destas duas hypotheses pode-admittir-se.- Não. pode admittir-se. a. primeira; porque as qualidades actuam nos sentidos, não \emph{separadas} do sujeito, mas \emph{existentes} no sujeito, pois é o sujeito que actua, embora pelas suas qualidades, nos sentidos. Nào pode admittir-se a segunda; porque as qualidades são percebidas pelos sentidos como \emph{são} em si mesmas, isto é, adherentes.a um. sujeito, pois os sentidos são incapazes de abstracção. Logo os sentidos externos, quando bem dispostos e convenientemente applicados, não podem enganar-se na percepção do objecto \emph{concomitante} \footnote{A percepção pode ser \emph{concreta} e \emph{abstracta}. É concreta, quando um objecto, composto de varios elementos, se percebe como é em si mesmo, como uma coisa só. É \emph{abstracta}, quando num objecto, composto de varios elementos, se "ercebe-0u-se-considera-uni-etemoento, v-1i0-se-consideram-o não se percebem. os outros... Assim, quando vemos um pedaço de \emph{marmo}: \emph{branco}, podemos considerar esse composto de marmore e de brancura, e podemos tambem considerar a brancura Sem o marmore, e O marmore Sem a. brancura. No primeiro caso, a percepção é concreta; no segundo, é \emph{abstracta},\par Posto isto, pergunta-se: A percepção dos sentidos \emph{externos} é concreta ou \emph{abstracta}? E como o objecto d' essa percepção é a substancia corporea emquanto modificada pelas qualidades sensiveis, a pergunta resolve-se na seguinte : Os sentidos \emph{externos} --- percebem a substancia corporea emquanto modificada pelas qualidades sensiveis, isto é, percebem o composto de substancia e de qualidades, --- ou percebem apenas a substancia sem as qualidades, ou as qualidades sem a substancia? A esta pergunta os philosophos respondem diversamente, conforme a diversidade dos seus principios. Uns, --- como \emph{Locke}, \emph{Berkeley}, \emph{Reid}, e sobretudo os \emph{Kantistas}, que só admittem phenomenos ou \emph{apparencias} c negam a existencia das \emph{substancias} ou dos nouwmenos, --- sustentam que a percepção: dos sentidos é \emph{abstracta}. Outros, como \emph{Aristoteles}, \emph{S}. \emph{Thomas}, seguidos por todos os escholasticos e philosophos de melhor criterio, ensinam que a percepção dos sentidos é concreta, emquanto os sentidos, percebendo as qualidades sensiveis como são (e são \emph{concretas} ou adherentes x uma coisa, a um sujeito), percebem o \emph{composto} de substancia e de qualidades, percebendo as qualidades come objecto \emph{proprio} ou communm, e a \emph{substancia} como objecto \emph{concomitante}. --- S. Thomaz diz que, para uma coisa ser objecto \emph{concomitante} dos sentidos, ou ser sensivel \emph{per} accidens, são necessarias tres condições : --- 4) que esteja unida a uma Outra coisa, que é sensivel por \emph{si}, --- \emph{b}) que seja percebida, de algum modo, pelo sujeito sensiente, aliás não seria sensivel, --- \emph{c}) que seja \emph{conhecida pelo} sensiente por meio de uma \emph{outra} faculdade percepliva, p. ex.. pela \emph{zm£eMigencia}, ou por uma outra faculdade sensitiva \emph{superior}. --- Adverte o mesmo S. Doutor que o sensivel \emph{per} «\emph{ccidens} não é tudo o que pode ser percebido num objecto sensivel, mas é o que a \emph{intelliz genclia immediatamente conhece}\par Vol. I 1a\pgn{194}após a impressão, que esse objecto exerceu nos sentidos \emph{externos}; assim, vendo um sujeito que falla ou caminha, percebemos, pela intelligencia, á vida d'elle, e dizemos: \emph{vejo} que esse \emph{sujeito} vive (De Anima, 1. 8, lect. 13).\par Que a percepção dos sentidos seja concreta, e não abstracta, é coisa manifesta. Na percepção \emph{abstracta}, opera-se uma \emph{analyse}, pela qual uma coisa se separa \emph{intencionalmente} de outra, isto \&, pela qual uma coisa se considera sem outra. Ora a \emph{analyse} é uma operação propria da intelligencia, e não dos sentidos. Os sentidos apresentam o composto como é em si mes- ' mo, na sua concreção, --- Nem podia ser de outro modo. O que alcançam os sentidos? Só a substancia? não; porque a substancia não pode actuar .directamente, por si mesma, nos sentidos, nem na intelligencia. A substancia está envolvida nas qualidades sensiveis, e não pode operar nem ser percebida.sem.ellas... Só as. qualidades. sensiveis? tambem não; porque as qualidades existem como: \emph{concretas}, isto €, adherentes a um sujeito, operam como concretas, e-por-isso-devem-ser-- percebidas. como concretas. Logo.a percepção dos sentidos é \emph{concreta}. :\par Mas não devemos exaggerar a parte dos sentidos nátpercepção do objecto\par \emph{concomitante}. Os sentidos. limitam-se a apprehender as qualidades sensiveis, como se encontram na realidade, e por isso concretas, isto é, existentes não em si mesmas mas num sujeito determinado, e essas qualidades concretas são por elles apresentadas a intelligencia. Eis tudo. Por quanto, se é certo que essas qualidades \emph{concretas}, percebidas pelos sentidos, constituem um composto de \emph{accidentes} e de \emph{substuncia}, é certo tambem que os sentidos não percebem que debaixo das: qualidades existe um sujeito, uma substanvia ; não sabém o que é \emph{substancia}, nem o que é accidente; não distinguem o elemento, que está em si, do elemento, que se apoia num sujeito, É a intelligencia, e só ella, que, recebendo dos sentidos as qualidades concretas, percebe immediatamente --- que estas, por isso mesmo que são qualidades concretas, devem residir num sujeito, numa substancia, e que esta substancia deve existir e deve ser dotada de uma natureza, proporcionada e correspondente a natureza das qualidades.\par Portanto os sentidos não se enganam \emph{per se}, pela sua natureza, na. percepção do objecto conconitante. Mas, se não se enganam \emph{por si}, os sentidos podem dar. a intelligencia a \emph{occasizo} de enganar-se, de emittir juizos falsos. Com elfeito, os sentidos não percebem \emph{directamente} a substancia, mas apenas as qualidades concretas. Ora pode acomiecer que as mesmas qualidades \emph{concretas} se encontrem em substancias \emph{diversas}, (porque não repugmna que substancias diversas tenham a mesma figura, a mesma côr, o mesmo peso, etc.), e pode acontecer tambem que as qualidades \emph{concretas} existam milagrosamente sem sujeito, sustentadas pela virtude omnipotente de Deus, Sendo assim, é claro que a intelligencia pode facilmente enganar-se, emquanto pode julgar ------ que certas qualidades concretas pertençam a uma certa substancia, quando \emph{realmente} pertencem a uma outra substancia diversa, --- ou que us qualidades\pgn{195}sensiveis adhiram actualmente a um sujeito, quando actualmente não adherem. \textasciicircum{} Mas, nesses casos, o engano --- não deve attribuir-se aos sentidos, que se limitam a apresentar as qualidades concretas, como são em si mesmas, sem: \&cerescentar se essas qualidades pertencem a esta-ou áquella substancia; e se existem, ou não, no sujeito proprio, --- mas deve attribuic-se a intelligencia, que não empregou todos so meios necessarios para evitar o erro e conseguir a verdade, --- É quaes são esses meios? Quando ha uma. seria duvida se as qualidades \emph{concretas} pertencem a esta ou áquella substancia, é necessario multiplicar as observações e as analyses, até se encontrarem outras qualidades mais proprias, \emph{mais} caracteristicas, de modo que possam convir a uma só substancia, e não a outra. Quando, por um milagre da omnipotencia divina, .as qualidades \emph{concretas} subsistem sem o sujeito proprio, é a Fé que se. encarrega de nos-instruir-e-preservar-de-urmr-vizo-erroueo-\par Se os sentidos externos percebem effectivamente. e por m modo \emph{imme dicato}, os corpos \emph{externos}, e se elles não se enganam \emph{por} si na percepção do seu triplice objecto --- \emph{proprio}. conmmum é concomitaate, é claro que a intelligencia, recebendo a relação dos sentidos e baseando-se exclusivamente no testemunho d'elles, pode e deve julgar immediatamente, sem nenhuma necessidade de raciocinio, que os corpos existem, que o mundo existe. O raciocinio é necessario, quando uma coisa não é \emph{evidente}. quando \emph{a} intelligencia não conhece \emph{ómunediatamente} a relação entre o predicado e o sujeito; mas quando conhece essa relação immediatamente, baseando-se na \emph{immediata} percepção dos sentidos. já não precisa de raciocinio. Repetimos: o que se vé, não precisa de ser demonstrado.\par *\par São tres as opiniões dos philosophos acerca da \emph{percepção} « e da \emph{veracidade} dos sentidos \emph{externos}.\par 1*.) Uns admittem --- que os sentidos \emph{externos} \textbar{} percebem, e de um modo\par immediato, os corpos \emph{externos}, --- e que a \emph{veracidade} dos mesmos sentidos é \emph{immediatamente} «\emph{vidente}, e não precisa nem pode ser demonstrada,\par 9\textasciicircum{}.) Outros sustentam que os sentidos \emph{externos} não percebem, de modo algum, os objectos externos; porque estes não existem, ou, se existem, não . podem ser percebidos pelos sentidos, --- e por isso negam. ou não discutem,\par a \emph{veracidade} dos mesmos sentidos.\par 3*.) Outros affirmam --- que a existencia dos Corpos \emph{externos} nos é manifesta, não por uma \emph{imediata} percepção dos sentidos \emph{externos}, mas por um \emph{outro} meio, --- e que a propria veracidade dos mesmos sentidos não € \emph{immediatamente evidente}, e precisa de ser demonstrada.\par Tendo exposto e defendido a \emph{primeira} opinião, que é a de todos os \emph{escholasticos}, antigos e modernos, e de todos os Padres e Doutores da Egreja, vamos expôr e refutar \emph{directamente a segunda} e a \emph{terceira}.\par A) A \emph{senda opinião}, como dissemos, é a que, negando a \emph{existencia} dos objectos \emph{externos}, nega a \emph{percepção} e a \emph{veracidade} dos sentidos externos.\par \textasciicircum{}\pgn{196}Para os sequazes d'esta opinião, uma percepção externa é uma contradicção; a percepção dos \emph{corpos} significa apenas a percepção das nossas ideas, 'das \emph{modificações} da nossa alma; Ó que vemos, ouvimos e, em geral, sentimos, é apenas uma serie de phenomenos internos do eu, uma serie de ideas ou de sensações meramente subjectivas, que, por uma illusão, projectamos para [ora de nós. Nada é real; só existem as nossas ideas. O mundo é apenas uma idea, uma apparencia sem realidade. --- É o \emph{idealismo}. Divide-se este em \emph{subjectivo} e \emph{objectivo}, conforme as imagens ou appurencias dos corpos, que a inteligencia possue, se fazem derivar unicamente --- da \emph{actividade} da alma, --- ou de uma acção do proprio Deus. O idealismo \emph{subjectivo} foi defendido sobretudo por Fichte; o \emph{objectico}, por Berkeley.\par O \emph{idealismo} é falso em \emph{si} e nas suas \emph{especies}.\par a) O \emph{idealismo}, em geral, é falso. :\par a) O \emph{idealismo} basea-se num falso fondamento. Por quanto, suppõe ellezzqueza-ahna-é-o-corpo-689-d\&as- substâncias \emph{completas}, uma estranha - a outra, --- que a alma está relegada num ponto do organismo, sem meio de cominunieação com a realidade externa, --- e que, por isso, não é possivel “uma impressão dos objetos, que, embora existissem, seriam sempre incognosciveis. Ora este fundamento, lançado por \emph{Descartes}, é inteiramente falso. A alma e o corpo, como veremos na \emph{Anthropologia}, são duas substancias incompletas, e substancialmente unidas, de modo que formam uma só \emph{pessoa} e uma só natureza ou \emph{essencia}, Não-é a alma só que sente, nem é o corpo só; mas é \emph{o} composto de alma e de corpo, é 0 \emph{orgão} animado, Logo O \emph{idealismo} basea-se num falso fundamento.\par \emph{b}). Se os objectos \emph{externos} não existem na realidade, mas exclusivamente nas nossas ideas, segue-se que os nossos sentidos \emph{externos} percebem o que não esta presente, o que não existe, Ora uma percepção sem a cois: percebida, -umaá visão sem a coisa vista, é simplesmente uma contradicção. Logo o \emph{idealismo}, em geral, é falso. " . \emph{eme fo} '\par b) O idealismo \emph{subjectivo} é falso. ! . . . \emph{a}) A causa das sensações não pode ser exclusivamente a \emph{actividade} da nossa alma. Por quanto, á nossa alma, quando se experimenta uma sensação, percebe que é passiva e que essa sensação é uma operação \emph{passiva}. Ora 0 \emph{principio} é á causa \emph{motriz} de uma operação \emph{passiva} não é o \emph{sujeito}, que é.por si indifferente, mas é o proprio \emph{objecto}. Diz S. Thomaz: « \emph{Objechum} comparatur ad actum potenliae \emph{passévae} sicul principium el causa movens » (\emph{Sum}. \emph{Th}., p. 1, q. 77, a. 3. Logo a causa das sensações não pode ser exclusivamente a actividade da alma.\par b) Devendo o effeito ser proporcionado \& propria causa, as sensações, se derivassem exclusivamente da actividade da nossa alma, não poderiam representar senão objectos \emph{simples} e \emph{espirituaes} ; porque simples e espiritual é a nossa alma, Ora as sensações representam coisas materiaes e corporeas, como os proprios :sdversarios confessam. Logo as sensações não podem derivar exclusivamente da actividade da alma, --- Além d'isso, se ás sensações derivassem uvicamente di alma, não haveria uma razão sulliciente\pgn{197}por que devessem representar, como effectivamente representam, os objectos corporeos e materiaes como existentes fóra de nós.\par e) O idealismo \emph{objectivo} é falso.\par a) Attesta-nos a consciencia que, quando experimentamos uma sensação, experimentamos, pela mesma impressão, a sensação e a causa d'ella; assim, quando nos approximamos do lume, sentimos o calor do lume, e, ao mesmo tempo. o lume quente. Ora, se a sensação fosse produzida por um agente extrinseco, como é Deus, as nossas mios, embora sentissem o calor, não poderiam sentir, como effectivamente sentem; o calor do Fame, nem o lume quente. Diz S. Thomaz: «Si illa species caloris in organo ab alio agente fieret, tactus, etsi sentiret calorem, non tamen sentiret calorem ignis, nec sentiret ignem esse calidum. cnm tamen hoc judicet sensus, cujus judicium in \emph{proprio} sensibili non errat» (De \emph{pot}. q. 2. a. 7). Logo as sensações não derivam immediatamente de \textsc{Deus}, mas dos [proprios objectos externos.\par \emph{Uc} 5) Se as- sensações: dos. objectos “\emph{externos} dexivassem: \textbf{exclusivamente}. da acção de Deus, e não da acção immediata dos proprios objects, a acção de Deus seria \emph{inutil} e má; --- il, porque Seria inutil produzirem nós am--- sensação de coisas materiaes, se estas não existissem realmente; --- má, porque Deus nos enganaria, levando-nos por um impulso irresistivel a julgar como existente o que realmente não existe. Ora é absurdo e- impio pensar que uma acção de Deus possa ser e seja \emph{inutil} e má. Logo as nossas sen-' sações não derivam exclusivamente da acção de Deus, mas da acção immediata dos proprios objectos. :\par Sc as nossas sensações têm por causa \emph{immediota} os objectos \emph{externos}, estes existem, e o testemunho dos sentidos \emph{externos}, que nos accusa a existencia dos objectos \emph{externos}, é legitimo ou veraz.\par E) \emph{A terceira} opinião, como dissemos, affirma --- que a existencia dos corpos ou dos objectos \emph{externos} não se manifesta por uma \emph{immediata} percepção dos sentidos \emph{externos}, mas por um \emph{outro meio}, --- e que, por isso, a --- veracidade dos mesmos sentidos externos não é \emph{immediatamente evidente}. --- \emph{oO} Este \emph{outro mein}, distincto da \emph{immediata} percepção dos sentidos, que nos lovaria a admittir com certeza a existencia dos objectos ou corpos externos, “é, --- para \emph{Descartes}, uma inclinação \emph{universal}, proveniente de Deus, --- para \emph{Malebranche}, a Revelação sobrenatural, contida na S. Escriptura, --- para \emph{Reid}, um \emph{instincto} natural, --- para \emph{Cousin}, \emph{Sentrou}! e outros non-escholastieps de \emph{L} in, o \emph{principio de causalidade}. --- Façamos uma breve critica das opiniões enumeradas, mostrando que o emprego d' esses meios, longe de substituir o \emph{exercicio} dos senUXdos \emph{externos}, presuppõe e exige necessariamente \emph{o exercicio} e a \emph{veracidade} dos mesmos sentidos.\par a) \emph{Descartes} ensina que uma \emph{inclinação universal}, proveniente de Deus, nos leva a admittir a existencia dos corpos, --- \emph{inclinação} que não pode deixar de ser \emph{veraz}, por isso mesmo que provem de Deus. Esta doutrina é falso pelas seguintes \textbf{razões}.\par \emph{a}) Na opinião de Descartes, o conhecimento do mundo seria \emph{posterior}\pgn{198}ao conhecimento de Deus. Ora dá-se exactamente o contrario. O conheci-. mento do mundo t-anterior-ao conhecimento de Deus, porque é pelo finito que nos elevamos ao infinito. Por isso não poderiamos conbecer a Deus, se não acreditassemos anteriormente na \emph{veracidade} dos sentidos \emph{externos}.\par 5) Alem d' isso, como: sabe “Descartes que essa énclinagzo- existe nos homens, e que é \emph{wnzversal;e} por isso impultavel.ao proprio Auctor da natu-. . reza? Elle não pode saber isso, senão pela percepção \emph{sensitiva esterna}, que \emph{o} mette em relação com os outros homens. Logo o juizo ácerca d' essa \emph{inclinagao universal presuppoe a veracidade} dos sentidos \emph{externos}.\par y) \emph{Malebranche} ensina que só a Revelação sobrenatural, contida na S. Escriptura, pode manifestar e manifesta ao homem a existencia dos corpos. --- Tambem esta opinião é inacceitavel. Para termos ou escutarmos o que diz á S. Escriptura, onde está contida a Revelação sobrenatural, precisamos do sentido da \emph{vista} ou do ouvido, e por isso já devemos acreditar na veracidade do testemunho dos sentidos \emph{externos}. :\par €) \emph{Reid} afirma que um \emph{instincto} natural, incapaz de enganar, leva irresistivelmente-os homens-a-admittir a-existencia \emph{dos-corpos}..--- Nem esta doutrina satisfaz. Esse \emph{insiincio}, seja qual fôr a sua. natureza, É necessariamente posterior á \emph{existencia} dos homens. Por isso, quem admitte a \emph{existencia}, dos homens, acredita na veracidade do testemunho dos. sentidos \emph{externos}, antes. de conhecer esse instincto ; porque é pelos sentidos externos que estamos certos da existencia dos hoiens. '\par d) \emph{Cousin}, \emph{Seniroul}, € outros \emph{neo-escholast de L} in, ensinam que admittimos a existencia dos corpos por uin \emph{raciocinio}, que da existencia das nossas \emph{sensações subjectivas}, cuja origem não pode attribuir-se ao eu, infere, em virtude do \emph{principio de causalidade}. a existencia de uma causa \emph{externa}. --- Esta \emph{opinião}, que se attribue tambem a Descartes (pois este \emph{infere} da existencia de Deus a existencia do mundo), não pode ser sustentada pelas seguintes razões.\par a) Esses: , escriptores admittem, como existentes na. realidade, umas coisas, ás quaes; dão.o nome. de. \emph{effeitos}, e outras, ás quaes; dàáo o nome de causas. Mas, de que modo poderão elles, sem o auxilio dos sentidos externos, averiguar que no mundo existem coisas, correspondentes e conformes ás noções abstractas de \emph{effeito} e de causa? E se não averiguaram, sempre pelo auxilio dos sentidos extersos, que entre um antecedente e um consequente existe um nexo de verdadeira \emph{efficientia}, como, poderão elles applicar o \emph{principio} de \emph{causalidade}? Portando o raciocinio, que da existencia das sensações pretende deduzir a existencia do mundo, presuppõe e exige o \emph{exercicio} e a \emph{veracidade} dos sentidos \emph{externos}.\par \emph{b}) Alem d'isso, a mesma experiencia, que nos attesta a existencia das modificações ou sensações \emph{subjectivas}, attesta-nos tambem a existencia das modificações ou sensações \emph{objectivas}, isto é, a existencia dos \emph{objectos externos}, que actuam nos nossos sentidos \emph{externos}. Portanto, se admittimos, como \emph{iminediatamente evidente}, á existencia das sensações \emph{subjechivas}, devemos tambem admittir, como \emph{émmediatemente evidente}, á existencia das sensações \emph{objectivas}. A mesma razão vale para ambos os casos. Logo não\pgn{199}pode acceitar-se essa opinião, que concede a evidencia \emph{émunediata} ás sensações \emph{subjechivas}, ao passo que a nega ás sensações \emph{objectivas}.\par c) Finalmente, o raciocinio, que da existencia das sensações subjectivas infere, em virtude do \emph{principio de} causalidade, a existencia do mundo exterior, tem uma conclusão mais ampla que as premissas. Por quanto, se as' nossas sensações são meramente \emph{subjectivas}, se são apenas modificações do eu, a unica conclusão logica é que essas sensações têm uma causa, mas uma causa \emph{inferna}, \emph{subjectiva}. Repugna por ventura uma cansa \emph{interna} de um phenomeno.- interno? Não. O \emph{subjectivo} não jnclue por-si-o \emph{objectivo}; e a passagem de uma para outra ordem é um sophisma. --- Mas, embora a existencia das sensações subjectioas nos levasse a admittir a existencia de uma causa externa, nunca poderia levar-nos ao conhecimento da natureza d' essa causa. Com effeito, na opinião de \emph{Cousin}, como na de \emph{Descartes}, os corpos \emph{externos} não produzem ás sensações \emph{subjectivas}, limitam-se a provo- -\emph{cal-as}; não são causas \emph{efficientes}, mas apenas \emph{occasionaes}. Ora um phenomeno, se nos leva a conhecer, de algum modo, a natureza da sua causa \emph{efficiente}, não nos leva a conhecer a natureza da causa occasional, que não se reflecte no effeito; um tiro de canhão não leva por si xo conbecimento do artilheiro, que o disparou. Por isso, o mundo, se devesse ser conhecido pelas sensações \emph{subjectivas}, seria incognoscivel na sua natureza, --- Se, porem, se admitte, com outros sequazes d'esta opinião, que os corpos externos \emph{produzem}. como causas \emph{eficientes}, as nossas sensações \emph{subjectivas}, em tal caso, como a acção da causa \emph{efficiente} é recebida e existe no sujeito paciente, este conhecerá a \emph{existencia} e a natureza da mesma causa, mas conhecel-as-ba por um modo \emph{immediato}, isto é, pela \emph{immediata}, percepção dos sentidos, sem ter a necessidade de recorrer ao raciocinio. Provar a existencia do mundo pelo raciocinio é explicar uma coisa clara por uma outra escura. --- Coneluimos que a existencia dos corpos \emph{externos é immediatamente} evidente, e por isso \emph{immediatamente evidente} é tambem a \emph{veracidade} dos sentidos \emph{externos}.}.

ibo \emph{o} Toei --- \textbf{Secção} FUNDA -

%%>>> t177

\pg{185} \pg{187} \pg{188} \pg{189} \pg{190} \pg{191} \pg{192} \pg{194} \tese{177}{Regras para os sentidos externos.}
Para que. os sentidos externos sejam meios seguros para a acquisição da verdade, devem observar-se as seguintes regras:

12) Quando-o--objecto-dos-sentidos é- \emph{proprio}, basta que \pg{195}  o \emph{orgão} esteja bem disposto e que o \emph{meio} não O impeça de \textbf{receber} a impressão do objecto. Pertence á razão julgar se \textbf{estas} duas condições se verificaram. :

II*) Quando o objecto é \emph{commum}, devem empregar-se ---- \pg{196}  \emph{todos} os sentidos relativos a esse objecto, e tambem a \emph{razão}; porque, embora o testemunho dos sentidos nunca seja \emph{por si} \pg{197}  falso, todavia pode levar a falsos juizos: o que se evita pela reflexão da razão. .

III*) Quando o objecto é \emph{concomitante}, --- \textbf{como} o erro é \pg{198}  mais facil, não só porque os sentidos. não são aplicados aos seus objectos \emph{proprios}, mas tambem porque substalicias diversas \pg{199}  podem ter \& mesma apparencia, --- é necessario o intervento da. razão, a qual, pela multiplicação das observações e analyses, chega a encontrar os traços proprios de cada substancia (\textasciicircum{}) : EC

%%>>> t178

\tese{178}{Sentidos internos.}
Sentidos-nlernos-s3o os-qué-exi----stem no: \emph{interior} do organismo, e pelos quaes o homem --- di- . \emph{recta} e \emph{immediatamente} attinge, por varios modos, as proprias \textasciicircum{} sensações, ou reproduz as \emph{imagens sensiveis} dos objectos \emph{externos}, --- \emph{indirecta} e \emph{mediatamente} percebe os mesmos objectos externos \footnote{Resumamos em poucas palavras o que deixamos dito. Os sentidos \emph{externos} são \emph{meios legitimos} para a acquisição da verdade. É por elles que entramos em communicação com o mundo exterior. Dizer com os \emph{subjectivistas} que.por meio dos sentidos externos só percebemos as \emph{nossas} sensações, e não os \emph{objectos}, é confundir a percepção \emph{externa} com a \emph{interna}, é contradizer : \emph{consciencia}, a qual attesta que, pelos sentidos externos, \emph{directa} c immediatamente percebemos os \emph{objectos}, \emph{e} que só, \emph{indirecta}. ou \emph{mediatamente}, por um sentido \emph{interno}, percebemos as nossas sensações.}.\footnote{Na \emph{Anthropologia} desenvolveremos a theoria ácerca dos sentidos internos. Aqui apontamos apenas a sua \emph{natureza} e O seu \emph{objecto}.}

%%>>> t179

\pg{200} \tese{179}{Objecto dos sentidos internos.}
O objecto, pois, dos sentidos \emph{internos} é duplice: \emph{directo} e \emph{indirecto}. --- O \emph{directo} é constituido pelas nossas \emph{sensações} e pelas \emph{imagens sensiveis} dos . objectos \emph{externos}.; --- o. \emph{indirecto}, - pelos \emph{objectos externos} \footnote{Donde 3 St \emph{VE} \textsc{QUE} O Exercicio dos sentidos \emph{tntermos} suppoe o exercicio dos sentidos. externos..e..que, se este não. se realisar ou for irregular, não poderá realisar-se ou não será regular aquelle. Un cego de nascimento nunca poderit formar a imagem sensivel da \emph{cór}.}. \textbar{}

%%>>> t180

\tese{180}{Numero dos sentidos internos.}
Os sentidos \emph{internos} são \emph{quatro}, a saber: o \emph{sentido} commum, a. \emph{imaginação}, a \emph{estimativa}, e a \emph{memoria sensitiva}.

a) Sentido commum. --- : \emph{Sentido} commum é a faculdade,

pela qual o homem percebe --- \emph{directamente} as sensações dos sentidos \emph{externos} e as suas \emph{differencas}, --- e \emph{indirectamente} o proprio. \emph{organismo} e o \emph{objecto} das mesmas sensações externas. --- É a continuação ou o complemento dos sentidos \emph{externos} \footnote{O \emph{sentido commun}. existe no homem. Não só \emph{vemos}, ouvimos, ete., mas \emph{percebemos} todas essas sensações externas, comparamol-as, \emph{distinguinmot-as} umas das outras .(o branco do \emph{amargo}, o som da cór, etc.). A faculdade, que tem por objecto as sensações, não pode deixar de ser \emph{sensitiva}. --- O homem, \emph{pelo sentido} commnum, não só percebe as proprias sensações externas, mas \emph{tambem o objecto} e 0 \emph{sujeito} d'ellas. Percebe o \emph{objecto} ; porque toda \emph{a} sensação termina'num objecto. Percebe o \emph{sujeito}; porque o homem percebe \emph{a} sensação não \emph{abstracta}, mas concreá; ora a sensação \emph{concreta} é o proprio sujeito modificado; e pór isso não se pode attingir a sensação, sem que, implicitamente, se attinja o sujeito. O \emph{sujeito sensiente} é o objecto conconwianie, OU per accidens, do : \emph{sentido}. commum. À percepção do \emph{sujeito} € \emph{direcía}, é inmediata, Mem requer 'um raciocinio, como pretende Reid.\par O \emph{sentido} commum chama-se tambem \emph{sentido intimo}. --- Sendo faculdade sensitiva, não deve elle confundir-se com a consciencia, que é a propria intelligencia, nem com \emph{o} chamado \emph{senso} commum, que é a disposição natural, de que o homem é dotado para perceber immediatamente e sem discurso os \emph{primeiros principios da razão}. --- Dá-se, ás vezes, ao \emph{sentido} commum o nome de consciencia; mas, então, para evitar confusões, é necessario accrescentar-lhe o adjectivo de \emph{sensitiva}. D' este modo distinguir-se-ha a consciencia \emph{sensitiva} da \emph{consciencia}, \emph{intellectual}. --- Como se vê, o \emph{sentido} commum faz na ordem \emph{sensitiva} O que a \emph{consciencia} faz na ordem \emph{intellectual}. O \emph{sentido} commim percehe os mesmos \emph{objectos}, que percebe a consciencia; com a diiferença que o \emph{sentido} commum percebe-os num estado real Ou \emph{concreto}, ao passo que a \emph{consciencia} percebe-os num estado \emph{ideal} ou \emph{abstracto}. É assim que os dados do \emph{sentido} comam se tornam, na consciencia, objectos de verdadeiros juízos, compostos de termos \emph{absíracios} ou universaes.}.

b) Imaginação. --- \emph{Imaginação} é a faculdade, pela. qual o \pg{201}  Homem conserva e reproduz as \emph{imagens} sensiveis dos s entes,

mesmo na ausencia d'elles \footnote{\emph{A imaginação} existe no homem. Vemos um animal, e, depois de o termos perdido de vista, conservamos ainda e reproduzimos a sua imagem sensivel. --- A imagem de um ente é sensivel, quando o representa cercado de propriedades concretas e \emph{materiaes}, como são a \emph{cór}, \emph{a figura}, \emph{a grandeza},\par cujas \emph{imagens sensiveis} a \emph{imaginação} conserva e reproduz, ' etc. --- Os entes, são os que foram percebidos \emph{immediatamente} pelos sentidos \emph{externos} e me "diatamente pelo “Sentido \emph{coin} Pot \emph{Siena} gi \emph{üppoe-o}-6 dos sentidos \emph{externos} e do sentido commum.\par A imaginação, no \emph{homem}. \emph{é reproductiva} € creadora. É Teproduciiva, quando reproduz as \emph{tinagens} dos objectos externos percebidos. É \emph{creadora}, quando une ou compõe varias imagens de objectos pe ebidos, pára formar \textasciicircum{} a imagem de um objecto, que nunca foi percebido e que talvez nem existe. --- Dissemos --- no \emph{Lkomem}; porque a \emph{imaginação} do animal irracional é apenas \emph{reproducliva}; visto que se esta faculdade é, em nós, tambem creadora, isso é devido ao facto de ella estar unida com a intellizencia.}.

c) Estimativa. --- \emph{Estimativa} é a faculdade, pela qual o “homem percebe nos objectos materiaes certas Propriedades de : “\emph{vantagem} ou de damno, que escapam- aos sentidos: externos \footnote{A \emph{estimativa} é laculdade \emph{sensitiva} e existe no homem e no animal irracional. Nos animaes irracionaes a existencia da \emph{estimativa} é mais manifesta. O cordeiro. vé pela primeira vez o lobo, e foge d'elle, como do inimigo. --- Se não é tão util e necessaria ao homem adulto, esta faculdade é muito util e necessaria na infancia, nas perturbações da razão, e em tantos movimentos, que: precedem a reflexão e que são regulados. por ella. --- A... \emph{estimativa} do homem é mais nobre que a do animal irracional. O animal percebe \emph{instinctivamente} as propriedades de \emph{vantagem} ou de damno, que existem nos objectos, mas não percebe a \emph{essencia} ou natureza d'esses objectos; ao passo que o homem, por ser tambem dotado de razão, conhece a natureza desses objectos, e assim os seus juizos são tambem comparativos, pois comparam o predicado com o sujeito. É por isso que a estimativa, no homem, toma o nome de razão \emph{particular}, porque compara termos \emph{particulares} (p. ex., \emph{esta herva é-me util}), Diz-se razão particular, para se distinguir da razão universal. ou \emph{infellectial}, que compara termos untiversaes.}.- -

d) Memoria sensitiva. --- \emph{Memoria sensitiva} é a faculdade, pela qual o homem conserva --- \emph{directamente} as \emph{proprias} per- \textasciicircum{}. cepções \emph{sensitivas passadas}, --- e indirectamente os \emph{objectos} d'aquellas percepções \footnote{No homem, pois, ha uma duplice \emph{memoria}: \emph{sensitiva e intellectual}. Pela memoria \emph{sensitiva}, lembramo-nos das \emph{nossas} percepções \emph{sensitivas pas}- * ' sadas; pela \emph{intellectual}, lembramo-nos das nossas percepções \emph{intellectuaes} passadas, --- Uma só faculdade não basta para perceber as percepções d' essa Quplice ordem --- \emph{intellecinal} e \emph{sensitiva}; porque deve haver proporção entre \textbar{} \& faculdade e o seu objecto directo.}. "mM -

%%>>> t181

\pg{202} \tese{181}{Os sentidos internos, operando nas condições indispensaveis, não podem enganar-se na percepção do seu objecto.}
Nenhuma faculdade, operando nas condições indispensaveis, pode enganar-se na percepção. do seu objecto; aliás a natureza. Arrastar-nos-hia fatalmente para o erro: o que é absurdo. Ora os sentidos \emph{infernos} são faculdades naturalmente destinadas a perceber as nossas \emph{sensações} ou a reproduzir ax \emph{imagens sensiveis} dos objectos \emph{externos}. Logo os sentidos \emph{internos}, operando nas condições indispensaveis, não podem enganar-se.na percepção do seu objecto \footnote{A razão é sempre a mesma. Uma faculdade, que por si, pela sua natureza, se engana na percepção do seu objecto, é e deixa de ser, ao mesmo tempo, uma faculdade; porque pode e não pode alcançar o seu objecto. Ora isto é absurdo.. --- É claro que, para a percepção sensitiva inferna ser verdadeira, devemos cumprir todas as condições indispensaveis, --- Mas, nesta-p crcepção, o-engano é menos frequente que---ba- percepão sensitiva externa. Por quanto, na percepção sensitiva \emph{externa}, 0 engano pode provir --- náoisó do \emph{orgão}, que pode não estar bem. disposto, --- como tambem do \emph{meio}, que pode alterar a acção do objecto. no respectivo sentido; ao passo que, na percepção Sensitiva interna, o engano pode provir unicamente do \emph{orgão}, que pode não estar bem disposto, pois, aqui, não ha \emph{meio} entre o objecto € a faculdade, mas o proprio objecto esta por si unido com a respectiva faculdade, ,determinando-a-immediatamente ao acto.\par Depois de termos provado a veracidade dos sentidos \emph{internos} em \emph{geral}, provaremos brevemente a veracidade de cada um d'elles em \emph{particular}.\par a) O \emph{sentido} communm não se engana na percepção do seu objecto.\par \emph{a}): 'Com oeffeito, esta. faculdade atlesiaà que vemos, ouvimos, gostamos, etc, que o nosso organismo se acha num estado normal ou anormal, que soffremos ou gozamos. Ora este. testemunho é tão directo, espontaneo e jrresistivel,. que:é necessario. ---..ou- admittil-o como. legitimo. --- ou cahir no mais completo e universal scepticismo. De facto, ninguem, fóra dos antigos scepticos, contesta a veracidade do \emph{sentido} comam.\par 5\emph{b}) Alem d' isso, o \emph{sentido} commum attinge a sensação dos sentidos externos emquanto é \emph{subjectiva} e adhere ao sujeito sensiente, Sendo assim, um engano: será possivel, se fôr possivel uma acção sem o objecto respectivo. Ora isto é absurdo.\par Para a solução das principaes difficuldades contra a veracidade do \emph{sentido commum}, apresentamos os seguintes principios; --- 1.º) Se, quando a nossa alma se acha absorvida, com excessiva vivacidade, num pensamento ouwnum sentimento. o \emph{sentido} commum não percebe as modificações dos sentidos externos, a razão é porque essas modificações foram apenas mate- Tiacs, \emph{e} não \emph{vitaes}. --- 2.º) O \emph{sentido} commun refere as sensações subjectivas, como ellas são; e, por isso, refere-as confusas, se são confusas, --- distinctas. se distinctas. A sua veracidade não exige mais. --- 3.9) Esta (aculdade limita-se a attestar \emph{of factos internos}; mas não emitte nem pode emittir um\pgn{203}juizo ácerca da natureza ou da causa d'esses factos. Portanto, \emph{o} erro rela- -livo a \emph{natureza} Ou a causa dos \emph{factos internos} deve imputar se a impru- dencia da razão e não ao testemunho do \emph{sentido l} ----4.?) Q-acto --- ------\par de perceber as sensações dos sentidos. "\emph{externos} é proprio do \emph{sentido} commum; o acto de \emph{reflectir} poluntariamente Sobre essas sensações « \emph{t} proprio da \emph{consciencia}. O primeiro acto é permanente e essencial ; o segundo é intermittente e accidental. --- 5.º) As modificações organicas e sensiveis, se \emph{passarem} o um \emph{estado} permanente, deixam de se manifestar ao: \emph{sentido} com- mum (porque a percepção de uma sensação exige uma mudança no estado habitual do orgão), e só tornam a manifestar-se, quando se dei uma uotavel modificação nesse estado habitual.\par b) A \emph{imaginação} não se engana na percepção do seu objecto. Com . efieito, o officio da \emph{imaginação} é o de conservar e reproduzir as \emph{imagens} ' \emph{sensiveis} dos objectos externos. Portanto, --- ou ella não conserva \emph{e} por isso não reproduz essas imagens, --- ou conserva-as e reprodul-as. Se' não as conserva, haverá nisso imperfeição, mas não engano. Se as conserva e reproduz, não pode deixar de as conservar e reproduzir, fielmente, --- -Todavia, a imaginação pode ser \emph{occasio} de erro: emquanto a intelligencia, habituada a dar um valor \emph{objectivo} ás proprias percepções, pode, sob \emph{a} influencia d' aquella faculdade (que reproduz as imagens dos objectos, mesmo quando estes não estão presentes), julgar como presente o que está realmente ausente. Mas, nesse caso, o erro deve attribuir-se a \emph{inlelligencia}..\par €) A \emph{estimativa} não se engana na percepção do seu objecto. Por quanto, os juizos da \emph{estimativa} ou são \emph{énstinctivos} ou \emph{comparativos}. Se são \emph{instinctivos}, são verdadeiros ; porque a natureza. de que derivam, não pode enganar. Se são \emph{comparativos}, são tambem verdadeiros; porque a \emph{razão}, dadas as condições necessarias, é meio legitimo para a acquisição da verdade. . d) A memoria \emph{sensitiva} não se engana na percepção do seu objecto. Na verdade, ou conserva as sensações \emph{passadas}, OU não ás conserva. Se não as conserva. mostra deficiencia, mas não se engana; se as conserva, cOonserva-as como são, e por isso não se engana.}.

%%>>> t182

\pg{203} \tese{182}{Regras para os sentidos internos.}
\textbf{Para} que os sentidos \emph{infernos} sejam meios seguros para a acquisição da, verdade, devem observar-se as regras seguintes: ---

E) Ós sentidos \emph{externos} devem estar bem \emph{dispostos} e convenientemente \emph{applicados} aos seus respectivos objectos; porque. o exercicio regular dos sentidos \emph{infernos} exige o exercicio regular dos sentidos \emph{externos}. \emph{vv}

IP.) Os orgãos dos sentidos \emph{internos} devem estar \emph{sãos}; aliás não percebem o seu objecto, ou não o percebem como elle é.

E) Os sentidos \emph{internos} devem limitar-se a a percepção do objecto directo, que é constituido pelas \emph{sensações} ou pelas \emph{imagens}. Para a percepção dos \emph{objectos} d' essas sensações ou imagens, é indispensavel o exercicio dos sentidos externos ().\footnote{Alem das regras, que são \emph{geraes} e communs a todos os sentidos\par E\pgn{204}internos, ha, outras-proprias-de cada-um: Mas. como---estas---são;-em-geral, applieações d'aquellas, limitamo-nos a apresentar as que dizemi respeito á imaginação, que, Tomo dissemos, pode induzir ao erro a inlelligencia. --- 1.º) Devemos, pela \emph{reflexão}, extremar, cuidadosamente, do elemento \emph{real} o elemento \emph{imaginario}. --- IL*) Não podemos aflirmar a actual existencia dos objectos apresentados pela imaginação, se essa existencia não fór atlestada pelos sentidos \emph{externos}. Por quanto, a \emph{imaginação}, pelo facto de representar um objecto, mesmo quando elle esta ausente, não pode ser testemunho fiel da actual e real existencia. das coisas. --- IlL.2) Devemos abster-nos, porquanto é ipossivél, de proferir.juizos, quando a \emph{imaginação} está excitada. A transgressão d'esta regra é causa de innumeraveis equivocos e desvarios.}

%%>>> art-9

\section*{Artigo IV.}
\addcontentsline{toc}{section}{Artigo IV.}

\pg{204} \emph{Testemunho}, \emph{suas especies e transmissão}

%%>>> t183

\tese{183}{Testemunho.}
\emph{Testemunho}, no sentido mais. rigoroso, é a \emph{manifestação de} uma \emph{verdade feita por} uma a \emph{outra pessoa}. --- Quem manifesta a verdade, diz-se \emph{testemunha}. --- A força moral, que a \emph{testemunha} possue para conciliar o assenso dos outros, denomina-se \emph{auctoridade}. --- Os requisitos necessarios, que coristituem a auctoridade, são. dois ;.. a \emph{sciencia} e a \emph{veracidade}. --- A adhesão, prestada pela intelligencia a uma verdade em força da auctoridade da lestemunba, chama-se fé. --- O \emph{testemunho}, embora pareça um meio menos nobre para a acquisição da verdade, é comtudo mais \emph{breve}, mais commuin, mais \emph{seguro} que os outros meios \emph{intrinsecos} \footnote{São dois, diz S. Agostinho, os caminhos por que chegamos ao conhecimento da verdade --- a \emph{auctoridade.e} a \emph{razão}. A auctoridade precede a\textasciicircum{} razão no tempo; á razão precede a. auctoridade na \emph{dignidade}. Eis as palavras .doS. Doutor: «Ad discendum dupBeiter.ducimur --- \emph{ewctoritate} atque \emph{rafione}. Tempore \emph{auctoritas}, re autem .\emph{ratio} prior est » (\emph{De Ordine}, 1; 9, e. 91. --- O homem conhece a verdade pela \emph{rázdo}. quando a descobre por si mesmo, ou claramente a percebe, se proposta por outrem, de modo que o motivo da adhesão é a evidente percepção da coisa. Neste caso, temos a \$ciencia, \& visdo,.. que é a \emph{adhesão da intelligencia}' a uma verdade ou immediatamente \emph{evidente} em \emph{si} mesma, ou \emph{evidente nos principios intrinsecos}. \emph{de que deriva}. --- O homem conhece a verdade pela \emph{auctoridade}, quando não a descobre, nem claramente a percebe, se proposta por outrem, mas a recebe \emph{unicamente} porque pro- Db9sta por outrem, de modo que o motivo da adhesão é \emph{exclusivamente} a Auetbridade. Neste caso, temos a fé, que é a \emph{adhesão} da \emph{intelligencia} a uma \emph{verdade}, \emph{evidente}, --- não \emph{immediatamente} em \emph{si} Mesa, nem mos \emph{principios intrínsecos}, --- mas \emph{num principio} extrinseco, que é a \emph{auctoridade}. No primeiro caso, o motivo da adhesão é a auctoridade da evidencia; no segundo, é a evidencia da \emph{auctoridade}. --- Dor isso a fé distingue-se da \emph{sciencia}, emquanto esta vé, e aquella não-vê. Diz S. Thomaz: «De ratione \emph{scientiae} ost quod\pgn{205}-habeat firmam inhaesionem cum \emph{visione intellectiva}; habet enim certitudinem procedentem ex inlellectu principiorum. \emph{Fides} aulem deficit a. scientia in eo \emph{quod non habei visionem} » (Sum. \emph{Th}., 1-2, q. 67, a. 3). E S. Agostinho jà disséra: «Quid est fides? credere quod non vides » (Tract. 40 in \emph{Joan}.).\par Por isso é preciso não confundir o \emph{magisterio} com a auctoridade do mesmo magisterio. A verdade, proposta pelo magisterio, pode o diseipulo vel-a, pela luz da propria intelligencia, intimamente ligada com os primeiros principios, e, neste caso, adhere a essa verdade, --- não pela \emph{auctoridade} do magisterio, --- mas pela \emph{evidencia} da coisa; e esta adhesão é \emph{sciencia}. A verdude, proposta pelo magisterio, pode o discipulo deixar de a ver em si mesma, e, neste caso, se adhere a ella, adhere unicamente pela auctoridade do magisterio ; e esta. adhesão é fé. Portanto o \emph{magisterio} não exclue a \emph{sciencia}. .:\emph{A.auctoridade} é constituida por dois dotes necessarios e suffücientes para == qué a testemunha merega credito --- a sciencia ea veracidade. Eur razão du \emph{sciencia}, á testemunha-não-se- engana ;. em. razão. da. veracidade, não en-. gana. --- Alguns auctores indicam tres requisitos para a auctoridade da testemunha : \emph{capacidade}, \emph{sciencia} e veracidade. Deve, porém, notar-se que a \emph{capacidade} está incluida na \emph{sciencia}; pois quem sabe, é \emph{capaz} de saber. .\par O \emph{testemunho} É mais. \emph{breve}, nais \emph{commum}, \emph{mais} seguro que os outros meios \emph{intrinsecos}. --- E mais \emph{breve}; porque, pela razão, chegamos a conhecer algumas verdades só depois de muito tempo e de muitos esforços, ao passo que, pelo \emph{testemunho}, conhecemos \emph{muitas} verdades em pouco tempo e com muita facilidade. --- E mais commum ; porque a maior parte dos homens, não podendo applicar-se aos estudos, só pode receber a verdade pelo \emph{testemunho}, \emph{humano} ou \emph{divino}. --- É mais seguro; porque a razão \emph{individual} està sujeita a frequentes erros, ao passo que o testemunho humano é menos sujeito ao engano, e o \emph{divino} é absolutamente infallivel.\par Os Kantístas, seguidos pelos \emph{moderiistas}, rejeitam toda e qualquer auctoridade, e admitlem, como dissemos, a \emph{immanencia} «omo lei e norma de pensar. Diz \emph{Blondel}: « O pensamento moderno, com uma susceptibilidade ciumenta, considera a noção de \emph{immanencia} como a condição propria da philosophia. E \emph{exactissimu} a idea que nada \emph{pode entrar} no homem sem \emph{que saia d'elle} e não \emph{corresponda}, \emph{a} uma n \emph{idade de expansão}, --- nada, nem \emph{como facto historico}, nem como \emph{doutrina tradicional}, nem \emph{como} . \emph{obrigação proveniente de fóra}; não ha para elle verdade, que se acceite, nem preceito, que imponha, sem que seja, de algum modo, \emph{autonomo} ou \emph{autochtoro}» (\emph{Lettres} sur \emph{les exigences de la}, \emph{pensée contemp}.).\par Resposta. Concedemos que nada pode entrar no nosso espirito, sem que Seja, de algum wodo, \emph{assimilado}, Uma verdade, para se tornar \emph{mossa}, deve ser percebida ; o inintelligivel nunca será \emph{objecto} da nossa intelligenciu. Mas negamos que a nossa intelligencia só deva admittir --- as verdades, que nasceram d' ella, ou que ella descobriu, --- ou as verdades, que nos são agradaveis... Como poderá o methodo da \emph{é»manencia} provar os \emph{factos historicos}\pgn{206}? Impossivel. Um facto.historico, para ser \emph{awfhentico}. deverá sahir do nosso espirito, ou corresponder: ás nossas \emph{necessidades d}' «\emph{xpansão} ? Não, com certeza. Um facto historico autonomo ou \emph{aufochtono} é uma contradicção. --- E não só recebemos \emph{de fóra} os factos \emph{historicos}, mas tambem tantos outros conhecimentos, relativos a todos os ramos da sciencia humana. O homem é essencialmepte \emph{tim animal} ensinado. A maior parte das verdades, que conhece, recebe-as dos outros. Os seus caleulos partem, geralmente, de verdades descobertas pelos outros. Diz bem \emph{Yonsegrive}: «O que seria da astronomia, se cada astronomo devesse descobrir, por si mesmo, a posição de todas as estrellas ? O que seria da \emph{chimica}, se cada chimico fosse obrigado a recomeçar. todas as analyses e todas as syntheses? O que seria da w \emph{gohysica} e, tambem, da \emph{mathematica}, 'se cada calculador devesse verificar as taboas dos logarilhmos e os instrumentos, de que se serve?» (\emph{Philosophie}, IL, p. 72). É, pois. impossivel que \emph{tudo possa sahir} do nosso espirito, a não ser que nesse espirito tudo ou muito tenha anteriormente entrado. Um exbrinseco, um \emph{objectivo} é absolutamente indispensavel, mesmo nas nossas operações mais simples e humildes. É o propro \emph{Blondel}, que, contradizendo-se, o affirma.:Diz assim: « O progresso dá nossa vontade obriga-nos a confessar a nossa \emph{insufficiencia}, leva-nos a admittir a necessidade de um \emph{supplemento}, e dá-nos a capacidade, não de o produzir ou de o definir, mas de o reconhecer e receber » (\emph{Lettres}, etc). --- A auctoridade, superior ao individuo, cuja necessidade é admittida por este eseriptor, não é só a humana, mas é tanibem. e sobretudo, a divina, porque \emph{Blondel} trata ahi do sobrenatural e da Re velacüo christá, \textbar{}} . !

%%>>> t184

\pg{205} \tese{184}{Especies do testemunho.}
O \emph{testemunho} pode ser --- \emph{divino} e humano, --- \emph{dogmatico} e \emph{historico}.

a) Testemunho divino e humano. --- O testemunho --- é \emph{divino}, quando é o proprio Deus quem manifesta uma verdade;

óbé \textasciicircum{}. LÓGICA --- Secção SEGUNDA -

--- é \emph{humano}, quando é o \emph{homem}. --- Para adherir convenientemente ao testmupho divino, a nossa intelligencia deve ser esclarecida e elevada pela luz sobrenatural da Fé. EE

b) Testemunho-dogmatico e histórico. --- O testémuünho --- é \emph{dogmatico}, quando a verdade manifestada é um \emph{dogma}, ou uma \emph{verdade scientifica} (o. ex., a immortalidade da alma humana), --- é \emph{historico}, quando a verdade manifestada é um \emph{facto historico}, p. ex, a victoria de Lépanto \footnote{A divisão do testemunho em \emph{divino} e humano funda-se na qualidade da \emph{testemunha}; e a divisão em \emph{dogmatico} e \emph{historico} funda-se na qualidade do \emph{objecto}. --- Deixamos outras divisões de menor importancia. :}. --- \emph{ete}

%%>>> t185

\pg{206} \tese{185}{O testemunho divino, tanto dogmatico como historico, é necessariamente infallivel.}
O testemunho divino será necessariamente \emph{infallivel}, tanto quando nos revela uma \emph{verdade Scientifica}, \emph{racional} ou \emph{sobreracional}, como quando nos manifesta um facto \emph{historico}, se Deus fór dotado de tão grande sciencia, que não possa enganar-se, e de tão grande \emph{veracidade}, que não possa enganar. Ora \textsc{Deus}, por ser a propria. Verdade, é.dotado de uma; \emph{sciencia}. \emph{infinita}, e, por ser a propria \emph{Santidade}, é dotado

:

de uma \emph{veracidade absoluta}. Logo o testemunho divino, tanto \emph{dogmatico} como \emph{historico}, é necessariamente \emph{infallivel} \footnote{- \emph{O testemunho divino} chama-se tamem \emph{palavra de} De«s ou Revelação. --- Como \emph{o} assumpto é de capital importancia, damos-lhe o conveniente desenvolvimento, sobretudo no que diz respeito aos \emph{mysterios}, reservando o resto para a Moral. --- Fallaremos, aqui, na \emph{essencia} da Revelação, no seu \emph{objecto}, na \emph{possibilidade} da Revelação dos \emph{musterios}, nas relações entre a Revelação e a \emph{sciencia}, ou entre a Fé e a \emph{razão}. \textasciicircum{} a) Revelação. --- Revelação, no sentido rigoroso da palavra, é a \emph{manifestagüo} de uma \emph{verdade}, racional ou \emph{sobreracionai}, \emph{feita por} Deus \emph{aos homens por um modo} sobrenatural. --- A Revelação --- \emph{a}] é a \emph{manifestução de uma verdade}; porque a Revelação, em geral, não se refere senão a. verdades, ensinadas por uma pessoa extranha; --- 5) é a manifestação de uma verdade racional ou \emph{sobreracional}; porque o \emph{objecto} da Revelação abrange não só as verdades, que não excedem a capacidade da nossa intelligencia, --- como a existencia de Deus, a immortalidade da alma humana, ete., --- mas tambem.as que a excedem e que se chamam \emph{muysterios}; --- \emph{c}) € á manife- .Stacüo feita por Deus; porque aqui a testemunha não é outro homem, mas . é o proprio Deus, que nos falla-directaiiente, ou por meio de-vutra pessoa;-------p. ex., por um propheta;------.é. a manifestação. feita..por.-um...modo- sob \emph{iural}; porque 3 Revelação é feita por um meio extranho \emph{e superior ás forças da} natureza. --- O testemunho divino, emquanto é \emph{manifestação de} uma \emph{verdade racional}, convem com .o. testemunho \emph{humano}, --- emquanto é manife- . stação \emph{de} uma \emph{verdade sobreracional}, \emph{feita por} Deus \emph{por} um modo \emph{sobrenatural}, differe do testemunho humano.\par D' onde se vê que a Revelação, propriamente dita, é sobrenatural, e esta differe da Revelação natural, que é a \emph{manifestação da verdade}, \emph{feita por} Deus aos homens por meio \emph{das obras} da creação. --- A Revelação sobrenatural e a Revelação \emph{natural} (impropriamente chamada \emph{Revelação}) differem quanto 20 \emph{principio}, a0 \emph{meio} e ao \emph{objecto}. --- Differem quanto ao \emph{principio}; porque, na Revelação natural, 0 \emph{principio} do conhecimento é a \emph{razão} humana, --- ao passo que, na \emph{sobrenatural}, é a Fé \emph{divina}. --- Ditferem quanto ao \emph{meio}; porque, ua Revelação \emph{natural}, o meio, por que Deus manifesta a verdade e que é constituido pelas \emph{creaturas}, é proporcionado á nossa natureza, --- ao passo que, na \emph{sobrenatural}, O meio por que Deus manifesta a verdade e que é constituido por uma communicação \emph{directa}, é superior a todas as exigencias da natureza humana. --- Differem quanto ao \emph{objecto}; porque, na Revelação \emph{natural}, nenhuma verdade excede a nossa capacidade, --- na \emph{sobrenatural}, \emph{ha} muitas verdades, que excedem a nossa comprehensão e que nunca teriam sido conhecidas, sem essa manifestação gratuita.\par Alem da Revelação \emph{natural} pelas \emph{creaturas} e da \emph{sobrenatural} pela \emph{Fi}, \emph{que} se referem á vida presente, ha uma Revelação \emph{sobrenatural} pela \emph{gloria}, \emph{que tem logar na} vida futura. A esta triplice Revelação corresponde uma \emph{triplice luz}, infundida por Deus na nossa alma ---- « luz da \emph{razão}, a luz da Fé pela \emph{graça}, e a luz da \emph{gloria}. --- Os tres graus de Revelação estão coordenados entre si; o \emph{segundo} suppõe o primeiro e prepara para o \emph{terceiro}.\par b) Objecto da Revelação. --- O \emph{objecto} da Revelação (sobrenatural) é cou- Sliluido, como dissemos, por verdades \emph{racionaes} e \emph{sobreracionaes}. --- Verdades\pgn{208} racionaes são as que não excedem as forças ou a comprehensão da razão creada; de modo que a mesma razão, entregue a sj mesma, pode descobril-as e perceber o \emph{inirinseco} nexo entre o sujeito e o predicado. --- Verdades \emph{sobreracionaes} são as que excedem as forças ou, a comprehens\&o da razão creada ; de modo que a razão, entregue a si mesma, não as pode descobrir sem a Revelação, nem, depois. da Revelação, pode perceber o \emph{inirinseco} nexo entre o predicado e \emph{o} sujeito. --- As verdades sobreracionaes cbamam-se, como indicamos, \emph{mysterios}. .\par Expliquemos a natureza do \emph{mysterio}. --- \emph{Mysterio}, em geral, é uma verdade occulta. Uma verdade pode ser \emph{occulta} por varias razões: --- ou porque, embora seja por si facilmente intelligivel, comtudo escapa á nossa experiencia, taes são os segredos dos corações; --- ou porque, ainda que seja conhecida quanto \& \emph{existencia}, cómiüdo não se percebe bem quanto á \emph{essencia}, assim dizemos que-o mundo-está-cheio-de mysterios;--- ou porque não pode ser conhecida quanto á \emph{existencia} sem a Revelação, embora, depois de revelada, facilmente se percebe quanto á essencia, taes são todas as instituições de caracter positivo, como, p. ex., a auctoridade do. Summo Pontifice ; --- ou porque excede por tal modo as. forças ou a capacidade da razão creada, que, sem a Revelação, não se pode conhecer a sua \emph{existencia}, e, depois de revelada, não se pode demonstrar nem perceber de modo algum a sua \emph{essencia}, ou o \emph{inirinseco} nexo entre o sujeito e o predicado. A verdade, que, por este ultimo modo, excede a capacidade da razão creada, é e diz-se, propria e rigorosamente, \emph{mysterio}, que por isso se define: \emph{Uma verdade}, \emph{que por'si} e \emph{absolutamente excede} .a \emph{capacidade da razão} creada.\par Advertimos que a distineção das verdades em racionaes e \emph{sobreraciomaes} não se funda em Deus, que é uma unica e simplicissima Verdade, mas exclusivamente na capacidade da razão creada, --- Advertimos tambem que o \emph{mysterio} se diz e é \emph{incompreensível}, --- não emquanto as palavras, que o exprimem, não tenham sentido algum (como pretendem os \emph{racionalistas}), --- mas*emquanto, mesmo depois de revelado, \& razão creada não chega a comprehender a \emph{essencia} d'elle, isto é, a intrinseca razão, pela qual o prediexdo convem ao sujeito. Assim, no mysterio da SS. Trindade, x nossa razão ercebe o que significa a \emph{unidade} de natureza e a \emph{trindade} de \emph{pessoas}, embora não comprehenda porque e como podem \emph{tres pessoas} ser um só Deus. Por isso, quando dizemos --- \textsc{Deus} é uno na \emph{essencia} e \emph{trino} nas \emph{pessoas}, a allirmação da conveniencia entre o predicado e o sujeito funda-se exclusivamente na \emph{auctoridade} de Deus. --- D' onde se vê a fatuidade dos incredulos € racionalistas, quando dizem que os mysterios são \emph{absurdos}. O \emph{mysterio}, como dissemos, é um juizo composto de dois termos, cuja relação é \emph{incomprehensicel} para a intelligencia creada. \emph{Se} a relação entre dois termos de um juizo é incomprehensivel, esse \emph{juizo} não pode chamar-se \emph{absurdo}. Por quanto, \emph{um juizo é absurdo}, \emph{quando comprehendemos} que o predicado\pgn{209}repugna ao sejeito. Mas, se não \emph{se comprehende} a relação entre os dois . termos, não pode dizer-se que o juizo é absurdo, nem que não é absurdo. : ---. Logo os \emph{mysterios}, por isso mesmo que são \emph{mysterios}, não podem cbamar-se \textasciicircum{} Vul absurdos. O que é \emph{sobre} a razão, não pode dizer-se que é contra a razão,\par c) A Revelação dos mysterios não repugna. --- Se repugnasse a Revelação dos \emph{misterios}, esta repugnancia dar-se-hia --- ou da parte de Deus, --- ou da. parte do homem, --- ou da parte dos proprios \emph{mysterios}. Ora essa repugnancia não se dà --- nem da parte de Deus, --- nem da parte do homem, --- nem da parte dos proprios \emph{mysterios}. --- Não se dá da parte de Deus; porque, com a manifestação dos \emph{mysterios}, Deus não só não deprime, nem humilha a .nossa intelligencia, mas, pelo contrario, eleva e nobilita esta faculdade, de- .. + ---: scobrindo-lhe novos horisontes, que ella nunca teria podido descobrir. --- Nem... ! E dá da parte do homem; porque o homem, se não pode comprefender- 08 --- desperta em nós um elevado sentimento da grandeza de \textsc{Deus} cum grande\par amor pela sua infinita bondade, que se revelou, de um modo ineffavel, so- .: bretudo nos beneficios da Encarnação, da Redempção e da Eucharistia. ---.... .Nem se dà da parte dos proprios \emph{mysterios}; porque as verdades \emph{sobreracionaes} não são contrarias ás verdades racionaes, e se os \emph{mysterios} não têm a certeza, gue nasce da \emph{evidencia} do nexo entre o predicado e o sujeito, têm a certeza, que nasce da \emph{evidencia} da \&uctoridade divina,\par d) Relações entre 2 Revelação e a sciencia, ou entre a Fó e a razão.\par \emph{a}) \emph{Entre} a Revelação e « \emph{sciencia}, \emph{ou} entre a Fê e a ração, não \emph{pode haver verdadeira oppostição}. Por quanto, a luz da Fé e a luz da razão são dois ratos, que partem do mesmo Deus, que é a Luz increada e infinita de toda a verdade. Ora, se entre dois raios, que derivam do mesmo principio, pode haver e ha distincção, não pode haver e não ha opposição : aliás haveria opposição no proprio principio: o que é absurdo. --- Se é verdade uma conclusão da razão, é absoluta verdade um artigo da Fé. Ora a verdade não pode ser opposta a verdade, mas só a falsidade. Uma contradieção certa\par entre uma verdade da Fé e uma verdade da sciencia é impossivel.\par 5) A Fé \emph{e} a \emph{racio presiam-se muto auxilio}, \emph{cada} uma \emph{en confirmidade com} a \emph{propria indole}. Diz o Concilio Vaticano « Neque solum Fides et ratio inter se dissidere nunquam possunt, sed opem quoque sibi mutuam ferunt » (Const. \emph{Dei Filius}, c. lV) --- \emph{A} razão auxilia a Fé, ou a Sciencia auxilia a Revelação, ou a Philosophia auxilia a Theologia revelada; e vice-versa.\par €) A \emph{razão auxilia} « Fé. (S. Thom. \emph{in} Boet., \emph{de Trinitate}, q. lI. a. 3). : Em \emph{primeiro} logar, a razão demonstra certas verdades, que são pream- : . bulos ou \emph{presuppostos} com relação á Fé, e a Theologia revelada; taes são a existencia e a personalidade de Deus, á possibilidade e \emph{o} facto da Revelação sobrenatural, etc. Ora a Fé \emph{presuppoe} estas verdades. Porquanto a Fé não é uma experiencia do divino nz nossa alma, come pretendem os \emph{chama}-\par Vol. 1 15\pgn{210}dos \emph{modernistas}, mas é uma firme adhesão-da \textbf{nossa} intelligencia a verdades, que não são por si evidentes e que'nos são propostas \textbf{pelo} magisterio divino. Mas, para que a nossa intelligencia possa adherir a verdades, que não são por si evidentes, é necessario que conheça, anteriormente, que Deus fallou e por isso que Deus existe. Sem este previo conhecimento, o acto de Fé não só seria irracional, mas até seria: impossivel. \emph{E} como estas verdades --- a existencia de \textsc{Deus}, o facto da Revelação, etc.,--- não são \emph{immediatamente} evidentes, devem ser \emph{demonstradas}. Esta demonstração é o officio da razão, ou da Philosophia. Por isso a Fé presuppõe a \textbf{razão}, ou a sciencia.\par Em \emph{segundo} logar, a, razão, ou a Philosophia, \textbf{dá noções} exactas de alguns termos, que se encontram nos artigos da Fé, ou se empregam na Tbheologia, como. os termos de \emph{essencia}; de pessoa, de. substancia, de \emph{accidente}, etc. --- Tocamos um ponto muito melindroso, mas muito importante.\par A Theologia revelada recebe os seus prineipios, que são os artigos da Fé, não de uma sciencia natura! ou creada, mas exclusivamente da sciencia divina, pela Revelção sobrenatural. Estes artigos são uma participação, uma impressão das verdades divinas na intelligencia bumanx As verdades divinas não podem penetrar e imprimir-se na intelligencia humana senão enuuciadas e concretizadas num \emph{juizo}; porque é no juizo, que, para nós, se encontra propriamente a verdade. Ora, quando Deus revela sobrenaturalmente uma verdade, não infunde na intelligencia bumana as ideas, que constituem o \emph{juizo}, em que se concretisa a propria verdade (tal infusão é possivel, mas não é o meio ordinario da Providencia), mas serve-se das ideas, que já existiam na nossa inlelligencia, estabelece entre ellas uma. união, que teria escapado a toda a mente creada, e d" este modo \textbf{manifesta as} suas verdades. Às nossas ideas; de que Deus se serve e que une num juizo, foram abstrahidas das coisas creadas, e importam uma analogia ou semelhança, ainda que remotissima, com a bondade divina; porque em todas as coisas creadas existe um vestigio, mais ou menos perfeito, da essencia infinita,\par Portanto, em toda e qualquer verdade revelada, ou artigo de Fé, devemos distinguir dois elementos --- o Amano e o \emph{divino}. O elemento humano € constituido pelas \emph{idéas}, que existem naturalmente na nossa intelligencia e que Deus une num juizo; 9 \emph{divino} consiste ne: propria \emph{unico} d' essas ideas, feita por Deus num \emph{juizo}, que exprime a verdade revelada. --- O elemento hwmalo \emph{presuppoe-se} ao \emph{divino}; porque, se não existissem na nossa intelligenci\& essas ideas, a verdade revelada seria incomprehensivel para o homem, como é incomprehensivel esta verdade natural --- o \emph{todo é maior} que « \emph{parte} --- para quem ignora o que é o \emph{todo} e o que é a \emph{parte}. --- Todavia, apezar d'essa \emph{presupposição} do elemento Àwmano, a verdade revelada deve chamar-se e é simplesmente \emph{divina} e nova. É \emph{divino}; pois é Deus, e só Elle, \emph{quem se e} associa aS \emph{ideas} preexistentes na mente humana, e por isso forma o \emph{juizo}, que \emph{exprime} a verdade revelada, visto que o juizo não consiste\pgn{211} nas ideas, consideradas isoladamente, mas na \emph{união} d' essas ideas; assim altribue-se \emph{a} um homem um juizo, embora elle não tenha inventado os termos, que o compõem. É «ova; porque novo é O nexo, nova a \emph{união} entre as ideas, --- nexo e \emph{união}, que, como dissemos, teria escapado a intelligencia humana.\par Adverlimos que a verdade divina, considerada na nossa intelligencia, é inadequada em relação á mesma verdade, considerada na propria intelligencia de Deus, mas não é \emph{falsa}. É inadequada ; porque as ideas humanas, abstrabidas das coisas creadas, não podem exprimir a grandeza infinita. Mas não é \emph{falsa}; porque, como dissemos, as coisas creadas são um reflexo, uma semelhança, um vestigio das perfeições do Creador. :\par As ideas, de que se compõe uma verdade revelada, devem exprimir-se por meio de Zermos-oraes-em-formulas-dogmalicas.-Estesctermos-nada-acere--- scentam a verdade; limitam-se a manifeslal-ba. A Egreja, a qual pertence formular os dogmas, escolhe esses termos, e fixa o sentido d'elles, Podia\par Egreja empregar um Lermo, que é proprio de uma eschola; mas, então, seria absurdo dizer que a dita eschola influiu na concepção do. proprio dogma. . --- Se uma eschola emprega um termo no sentido, em que o emprega a Egreja, 1», conclusão é que essa eschola está no caminho da verdade.\par Dadas estas explicações, é facil assignar a razão, ou a Philosophia, o seu officio com relação á Fé, ou a Theologia revelada. Este officio consiste --- em acceitar a verdade revelada, como Deus a revelou e a Egreja, depositaria e mestra da Revelação, a propõe, --- em analysar e illustrar as ideas, de que são compostas as formulas dogmaticas, seguindo sempre o magisterio da Egreja, --- em definir e explicar os \emph{termos}, que exprimem essas ideas, tomando-os sempre no sentido, em que os toma a Egreja, --- em procurar "nas coisas creadas umas analogias ou semelhanças, para que a nossa-intelligencia possa, de algum modo, elevar-se ás coisas divinas, que excedem a capacidade nossa. (Sum. \emph{Th}., p. 1, q. I, a. 9 ad 2).\par Em \emph{terceiro} logar, sendo necessariamente verdadeiros os \emph{principios}, em que se funda a Theologia e que são as verdades reveladas, e necessariamente verdadeiras as conclusões, que d' esses principios derivam por legitima dedueção, são necessariamente falsas as \emph{objeccóes}, que contra esses principios e essas conclusões levantam os adversarios. E essas objecções são falsas --- ou porque se baseam em falsos principios, --- ou porque não derivam legitimamente de principios verdadeiros, O officio da Pbhilosophia consiste em descobrir a falsidade dos principios, ou a illegitimidade das consequencias.\par . A estas tres especies de auxilio, que a razão presta a Fé, ou a Philosophia presta a Theologia revelada, devemos accrescentar que a razão --- \emph{prova}, com argumentos \emph{proprios} eG \emph{inírinsecos} a essencia das' coisas, as verdades reveladas, que não excedem 2 capacidade da nossa razão, --- \emph{deduz} das verdades reveladas uma multidão de conseguencias e de \emph{applicações}, que estavam implieita \emph{e} confusamente contidas naquellas verdades e que só se tornam\pgn{212}explicitas e distinctas pela analyse e pelo raciocinio, --- une \emph{entre si e} ordena «s \emph{verdades reveladas}, de modo que possam constituir um perfeito corpo de doutrina, e assim" forma a mais bella entre as sciencias, --- a Theologia.\par d)A Fé \emph{auxilia} a \emph{razão}, --- O Concilio Vaticano enumera os beneficios, que derivam da Fé para a razão, e por isso da Theologia revelada para a Philosophia, com as seguintes palavras : « \emph{Fides rationem} ab \emph{erroribus liberat} et \emph{tuebur}, eamque \emph{multiplici cognitione instruit} » (Const. \emph{Dei Filius}, \textbar{}. c.). Na verdade, a Fé adverte a razão para que não siga opiniões falsas, que ' apresentam uma especie de verdade, --- descobre-lhe muitas verd:des, que ella, a razão, entregue a si mesma, nunca teria descoberto, --- confirma muitas outras de ordem natural. A historia attesta que, antes do Christia- ,nismo;-os-philosophos;-ainda-os-ma\par : celebres; cahiram.em. muitos e graves eiro\$ cerca --- da Divindade, --- da origem das coisas, --- da natureza, origem” e destino do homem, --- e dé outras questões, cuja solução verdadeira é absolutamente necessaria para a humanidade. --- Estes beneficios, que derivam da Fé, ou da Theologia revelada, deve a razão, ou a Philosophia;.acceital-os com espirito de gratidão e de-reverencia. --- Quando se trata de verdades reveladas, que excedem a capacidade da nossa intelligencia, o Philosopho, como dissemos, não deve negal-as, nem interpretal-as com o seu criterio, mas deve acceital.as com submissão, e rejeitar, como certiamenle falsas e contrarias á razão, todas as conclusões contrarias abs prin- 'eipios da Fé, porque esta não pode errar. --- Quando se trata de verdades reveladas, que não excedem a capacidade dà nossa intelligencia, o Philosopho deve tel-as presentes, isto é, deve pensar que são reveladas, para que se não afaste do 'reclo caminho, quando quer demonstral-as com argumentos \emph{inírinsecos}, \emph{a isto} é, baseados na propria essencia das coisas; e, se [he acontece de- chegar a uma conclusão contraria a uma verdade revelada, deve; \emph{no} interesse \emph{da} .\emph{Fé e} da propria verdade scientifica, repetir as investigações, até que a conclusão scientifica se harmonise com a verdade revelada,\par Nem esta submissão da Philosophia a. Theologia, ou da razão àá Fé, pode chamar-se uma \emph{degradação} para a razão humana, --- não só porque ninguem se degrada, quando acceita o auxilio de mão benefica, que o salva do precipicio, --- mas tambem porque a Philosophia, empregando as proprias forças e os proprios principios, demonstra com argumentos \emph{intrínsecos}, baseados não \emph{na auctoridade:mas} na \emph{essencia} das coisas, a verdade racional, que a Theologia lhe aponta, A aftirmação dos racionalistas --- que a \emph{liberdade do erro} é, preferivel a \emph{escravidão da verdade} --- é um absurdo tão \emph{paradoxal}, que só é possivel nos loucos. A nossa razão é tanto escrava da Fé, quanto o navegante é escravo do pharol!...}.

- - in

186. A adhesão ás \textbf{verdades} reveladas por \textbf{Deus} tem \textbf{certeza} metaphysica, superior a \textbf{toda} outra certeza metaphysica.

a) Tem \emph{certeza melaphisica}. --- A nossa adhesão a uma verdade é sempre proporcionada: ao wwolívo, que a determina; de modo que ella é dotada de \emph{certeza setaphaysica}, \emph{pharsica}, ou \emph{moral}, \pg{209}  conforme a qualidade da \emph{ceríeza} do proprio motivo: Ora a \emph{auctoridade} divina, que é.o motivo da adhesão da nossa intelligencia ás verdades reveladas, é dotada de certeza sneta \emph{physica}; pois a. sciencia e a veracidade,-as quaes-constituem-a \emph{auctoridade} divina, competem por tal modo a Deus, que não podem deixar \pg{210}  de Jhé/competir, sendo Elle a Verdade e a \textbf{Santidade} subsistente. Logo a nossa adhesão ás verdades \textbf{reveladas} por \textbf{Deus} tem “certeza \emph{metaphysica}.

b) \emph{Esta} certeza é superior \emph{a toda} outra \emph{certeza metaphysica}. --- A certeza é, como dissemos, proporcionada á \emph{intelligencia} e \pg{211}  ao \emph{motivo}. Ora o acto praticado pela nossa intelligencia; \textbf{esclarecida} pela luz \emph{sobrenatural} da Fé, é por si mais \textbf{firme} que todo e qualquer acto \emph{natural} ; como tambem \emph{a} auctoridade divina, em que se funda a certeza da nossa adhesão ás verdades reve- . ladas, é um motivo mais firme que o proprio motivo, em que \pg{212}  se funda toda outra certeza \emph{melaphjsico} e que é a \emph{essencia} das

coisas; porque o engano e a falsidade repugnam mais na Verdade primeira do que na maxima evidencia da natureza. Logo a certêza dà nossa adhesão ás verdades reveladas por Deus é superior a toda outra certeza \emph{metaphysica} \footnote{Não insistimos no assumpto, já tratado no capitulo precedente,\par Objectam: A certeza, que temos das coisas, que \emph{se vêem}, é mais firme do que a certeza das coisas, que \emph{se} não \emph{vêem}, Ora a certeza da Fé \emph{divina}\pgn{213}refere-se a coisas, que \emph{se} não vêem. Logo a certeza da Fé \emph{divina} não é mais ' firme do que a \emph{natural} certeza \emph{metaphysica}. \textasciicircum{}. 7\par Resposta. Devemos, antes de tudo, \emph{distinguir} a proposição maior. A certeza, que temos das coisas, que se \emph{vêem}, é mais firme do que a certeza das coisas, que se \emph{não vêem}, --- emquanto se considera' no \emph{sujeito}, ---- mas ' não emquanto se considera no \emph{objecto}. Por quanto, a certeza do \emph{sujeito} não está sempre na razão directa da certeza do \emph{objecto}. O \emph{objecto} pode ser visivel em si mesmo, e todavia o \emph{sujeito} pode não velo, ou não velo tanto, quanto elle é visivel; assim o sol é visibilissimo em si mesmo, mais visivel do due todas as coisas, que elle illumina e torna visiveis, e comtudo os nossos olhos véem melhor as coisas illuminadas pelo sol, do que o proprio sol. Appliquemos esse principio ao nosso caso. As verdades reveladas por Deus, que são o objecto da nossa Fé, possuem a virtude de produzir no nosso espirito uma certeza mais.firme do que a certeza, que produzem, as verdades da sciencia; mas, como, nas verdades reveladas, não percebemos o nexo \emph{inirinseco} entre o predicado e o sujeito, pode a nossa inteligencia adherir menos firmemente, e pode mesmo não adherir. Mas, repetimos, a nossa incerteza ou o nosso erro não depende da incerteza do \emph{objecto}, mas da indisposição do \emph{sujeito}. --- Devemos tambem distinguir a proposição menor. As verdades reveladas --- não são vistas em \emph{si}, 10 nexo inírinseco, que une o predicado e o sujeito, --- mas são \emph{vistas} na auctoridade, que as manifesta. A nossa intelligencia, para que possa adherir a uma verdade, deve necessariamente \emph{vel-a} de \emph{algum modo}, --- OU em \emph{si mesma}, --- Ou na auctoridade, que a propõe. Ora as verdades reveladas não são vistas em \emph{si} mesma. Logo devem ser vistas na \emph{auctoridade divina}. Repetimos que, nesse caso, se não temos a auctoridade da \emph{evidencia}, temos a \emph{evidencia}, da \emph{auctoridade}. --- Em conformidade com estas distinccoes, negamos a conclusão.}.

%%>>> t187

\pg{207} \pg{208} \pg{213} \tese{187}{O testemunho humano-dogmatico é infallivel em relação aos juizos de senso commum.}
Os \emph{juizos} de \emph{senso} commum, por serem \emph{constantes}, universaes e communs, SÓ podem derivar do intimo da \emph{natureza racional}; porque só esta, existindo \emph{sempre}, em \emph{toda}, \emph{a parle} e em \emph{todos os homens}, tem a capacidade de produzir effeitos, dotados d'aquelles caracteres. Ora o que deriva do intimo da natureza \emph{racional} não pode ser falso, --- não só porque Deus, Áuctor da natureza e Verdade subsistente, não pode imprimir no nosso espirito uma natural inclinação para o erro, --- mas tambem porque a adhesão de todos os homens a estes juizos não pode ter outra causa a não ser a evidencia da verdade, porque só a verdade, evidentemente conhecida, pode produzir a unanimidade e universalidade de consenso. Logo o testemunho \emph{humano-dogmatico} é infallivel em relação aos juizos de senso commum. (\textasciicircum{}).\footnote{Existem no seio da humanidade alguns juizos, que são \emph{constantes}, «\emph{niversaes}, \emph{communs}, isto é, cuja verdade foi e é reconhecida \emph{sempre}, \emph{por}\pgn{214}\emph{toda} a \emph{parte e por todos}. São os juizos de senso \emph{commum}, relativos á nossa vida physica, intellectual c moral. Taes são: --- existem os corpos, --- as leis da natureza são constantes, --- Deus existe e merece o nosso eulio, o acto . honesto distingue-se do deshonesto, --- aos paes é divido o respeito dos filhos, --- a alma bumana é immortal, ete. --- Para que, pois, im juizo possa chamar-se de senso commum, deve ter os seguintes caracteres: --- 1.º) deve ser \emph{universal}, \emph{diwtwrno}, \emph{constante}, não só entre os povos selvagens, mas tambem entre es civilisados, de modo que não seja. esquecido com o andar do tempo e com o augmento da cultura; --- 2.º) deve. convir com os principios da razão, de modo que, reduzido a esses principios, manifeste mais claramente a sua verdade; --- 3.º) não deve provir da ignorancia, ou dos prejuizos, ou da corrupção dos costumes, ou da superstição, ou do interesse, ou de outras semelhantes causas. --- É claro que a universalidade, de que deve ser dotado um juizo de senso commum, é a moral. Se um ou outro individuo-cotitesta---uni-juizo-de-senso-comntm.-nem-por-isso---enfraqueve: \emph{o} testemunho da humanidade. Ha monstros tambem na ordem moral.\par \textbar{} Os juizos de senso commum constituem o patrimonio, de que está dotado \emph{o} genero humano e que a todos e a cada um dos. homens: serve immensamente para o uso da vida. Nessa distribuição, a natureza Dão foi exclusiva, nem parcial. --- Todos os homens, mesmo os selvagens, por isso mesmo que têm o bem da intelligencia, deduzem dos communs principios da razão, ou das naturaes inclinações, esses juízos, e os deduzem com um raciocinio espontaneo, embora vulgar e imperfeito por falta de disciplina. De facto, não é nada difficil deduzir --- da vista do mundo \emph{a} existencia de Deus, --- ou do desejo da felicidade a immortalidade da alma, etc. --- O que deixamos dito é confirmado pela historia de todos os povos.\par Estes juizos não podem derivar senão da natureza racional. Por quanto, devendo a causa ser proporcionada ao effeito, um effeito \emph{constante}, Universal e \emph{comnmum},. como é todo o juizo de senso commum,; só pode derivar de uma causa \emph{constante}, universal é commum. Ora a causa, dotada d' estes caracteres, não pode ser senão a \emph{natureza} racional, porque só esta, no meio de tantas divergencias de opiniões, de estudos, de habitos, de caracteres, etc., existe \emph{sempre}, \emph{por toda} a parte. em \emph{todos os homens}, Com razão escreve Cicero: «Omni in re, consensio omnium gentium \emph{lex} naturae putanda est». (\emph{Tuscul}, 1. 1, 1:13). --- Derivando da natureza racional, os juizos de senso commum não podem ser falsos, ou erroneos; porque a natureza racional, feita ou creada unicamente para alcançar a verdade, não está \emph{por si} sujeita a engano. O erro deriva sempre de uma causa \emph{contingente}, \emph{partic«lar}, como é a ignorancia, o prejuizo, a paixão, etc. A propria natureza pode, por uma d' estas causas, desfallecer num ou noutro individuo, mas não é possivel! que desfalieça em todos. (\emph{C}. \emph{Gentes}, 1. IT, c. 34).\par D' onde se vé --- que a causa, ou 0 \emph{principio}, á' onde derivam os juízos de senso \emph{commum}, é a, \emph{razão humana}, --- e que O \emph{senso} commum é a propria \emph{razão natural}, emquanto, sempre e em toda a parte, leva os homens\pgn{215}a admittir verdades, immediata ou mediatamente evidentes, sobretudo as que nos são necessarias para a vida physica, intellectual e moral.\par Os \emph{tradicionalistas} encontram essa causa \emph{exclusivamente} na revelação, feita por Deus ao primeiro homem e transmittida ás seguintes gerações. --- O erro d' esses escriptores está todo nesse adverbio exclusivamente; porque concedemos que a causa \emph{adequada} de muitas. verdades, tidas “como juizos do senso comun, seja o \emph{recto} uso da \emph{razão} e a revelação \emph{primitiva}.\par \emph{Reid} opina que O senso commum é um \emph{cego}, \emph{instincto} da natureza, pelo qual somos determinados a admittir as verdades fundamentaes, não porque as vemos claramente, mas porque uma força irresistivel não nos deixa duvidar ácerca d'ellas. \emph{Balmes} accrescenta que 0 .senso commum é um meio \emph{especial}, que é distincto dos outros e que pode chamar-se \emph{instincto inel}-, lectual. --- Mas estas opiniões são falsas. \emph{A} adhesão da humanidade a certos juizos fundamentaes não é determinada por um instincto cego, mas unica- "mente pela intuição-da- verdades porque esses juizos ou-são immediatamente --- evidentes, ou percebem-se por um Íacillimo raciocinio. Alem d'isso, 2 ad- missão de um cego insüincto intellectual déprinié a natureza racional, tira" a sciencia a dignidade de \emph{racional}, e favorece O \emph{sceplicismto}.\par Alguns dos modernos \emph{aanti-intellectwalistas}, como \emph{Bergson}, admittem, em theoria, a força do \emph{senso} commniwm, mas, na pratica, contestam-na. Outros, como \emph{Le Roy}, negam que os juizos do senso commum (como tambem as verdades reveladas) tenham um valor intellectual, e affirmam que esses juizos têm apenas a razão de \emph{preceitos} Ou \emph{regras} para a vida \emph{pratica}. Assim, O axioma, p. ex., 2 -- 2---- \&náío contem nenhuma verdade racional, mas nos insinua que devemos \emph{operar}, como se 2 + 8 fossem == 4. --- Mas esta theoria, longe de ser uma explicação, é uma mutilação do senso commum. Os axiomas tiram a força, não da utilidade, mas unicamente da verdade.\par Identifüicando-se o senso comunum com a inlelligencia e a \emph{razão}, é claro --- que aquelle não constitue um meto, 'distincto. d' estes dois, --- e que os seus juizos têm uma certeza \emph{metaphysica} Ou \emph{moral}, conforme a natureza dos objectos. --- Todavia o senso \emph{comum}, pelo facto de ser citado como \emph{auctoridade}, colloca-se tambem entre os meios \emph{exírinsecos}. é uma coisa moral e os seus juízos tem uma certeza moral, e não \emph{metaphysica}.\par *\par Sc o testemunho Jnunano-dogmatico é infallivel quanto aos juízos de senso comumum, não o é quanto a outros juizos. A intelligencia humana pode errar e erra. Ora um testemunho, que está sujeito a erros, não pode., produzir certeza, mas apenas uma opinião, mais ou menos provavel, segundo. a maior ou menor auctoridade das testemunhas, Por isso o testemunho humano, quando ensina um \emph{dogma}, ou uma verdade scientifica, não é por \emph{si}, emquanto \emph{testemunho}, um meio infallivel para a acquisição da verdade. --- “Todavia esse testemunho é muito util, porque leva a intelligencia á descoberta da verdade; --- e não só é util para os ignorantes, que não têm a capacidade de descobrir a verdade, d' onde o adagio: \emph{perítis} in \emph{arte} credendum\pgn{216}, --- mas tambem é üül para os sabios, que não podem abranger a esphera toda dos conhecimentos humanos, e por isso devem aoeceitar dos outros sabios as verdades, relativas a um ou outro ramo de sciencia. Acceitando esse testemunho, pode 0 homem enganar se; mas, como diz S. Agostinho, «etsi auctoritate decipi miserüm sit, longe tamen est miserius ab eadem non commoveri » (\emph{De util}, \emph{credendi}, c. 16). tol\par Fallamos \emph{na auctoridade} do testemunho humano-dogmatico, emquanto € \emph{auctoridade}. É neste sentido que S. Thomaz diz: «Locus ab \emph{auctoritate}, quae fundatur super rationem humanam, est \emph{infirmissimus} » (\emph{Sum}. \emph{Th}., p. I, q. 1, a, 8 ad 2). Por quanto, se um homem acceita uma verdade proposta por outro, e acceita essa verdade, não porque proposta por outro, mas porque fundada em bóas razões, nesse caso não é a \emph{auctoridade} o criterio da don- i 'trina, mas a \emph{doutrina} é o criterio da auctoridade, e, por isso, o meio, que leva o discipulo a consecução da verdade, não é o \emph{testemunho}, mas é a propria \emph{ração}, a qual como vimos, é por si infallivel. --- Para maior clareza, \textasciicircum{} vejamos-a-inffuencia-do-mostre-no discipulo: Adquirese'a sciencia, quando se deduzem dos primeiros principios as conclusões, que nelles estão contidas. Ora os primeiros principios não são infundidos pelo mestre na intelligencia do discipulo; são conhecidos por este de um modo \emph{immediato}, por uma simpliecissima abstracção, e presuppóem-se- ás explicações do mestre. Em que consiste, pois, o officio do mestre ? O officio do mestre consiste em \emph{guiar} 1 a inteligencia do discipulo no discurso, que parte dos primeiros principios e chega á conclusão, D' este modo, o diseipulo adhere firmemente ás conclusões, apoiando-se, não na auctoridade do mestre, mas nas hóas razões, que \emph{ouviu} do mestre e fez suas, e que lhe mostram o nexo intimo e essencial, existente entre os principios e a conclusão. (Cf. \emph{De} Ver., q. 11, a. 1).}

%%>>> t188

\pg{214} \tese{188}{O testemunho humano - historico pode ser um meio .}

infallivel para a acquisição \textbf{da verdade}. --- O testemunho \emph{humano-historico} poderá ser um meio infallivel para a acquisição \pg{215}  da verdade, se a sciencia e a \emph{veracidade} das \textbf{pessoas}, que attestani um facto, se tornarem evidentemente \emph{certas}; \textbf{porque}, em tal caso, toda a duvida é irracional. Ora a sciencia e a veracidade \pg{216}  das pessoas, que attestam um facto, podem tornar-se \emph{evidentemente certas}. Logo o testemunho \emph{humano-historico} pode ser um meio infallivel para a acquisição da verdade \footnote{A ;consideracio das leis, que regulam os costumes humanos, levam-nos a prestar fé a \emph{auctoridade}. Estas leis podem reduzir-se ás duas seguintes: 1.º) \emph{O} homem \emph{tende} naturalmente \emph{para} a \emph{verdade}. Por isso, quanto mais prudente e diligente for a testemunha, tanto menos: possivel será o receio do erro e do engano. --- I1:º) \emph{O hontem} ndo \emph{opera sento pelo} amor \emph{de algum} bem. Por isso a mentira, que é por si vergonhosa e indigua do homem, não pode attrahir \emph{o} homem, senão sob u apparencia de \emph{utilidade} ; se não existe esta. utilidade, nem deve existir a suspeita de mentira, pois ninguem é man gratuitamente (nome \emph{gratis est} molus). Sobretudo não ha receio de engano, quando a testemunha é um homem de bem, que se deixa guiar pelas leis da honestidade, e não pelas. da utilidade propria.\par Dizemos que o testemunho humano \emph{historico} é meio Seguro para a acquisição da verdade, se constar evidentemente da \emph{sciencia} e da \emph{veracidade} das testemunhas. É, pois, da maxima importancia estabelecer algumas regras, para se everiguar ee o testemunho é dotado das duas qualidades essenciaes. --- Estas regras formam \emph{a critica historica geral}.\par Critica historica geral. --- Aponlamos as principaes regras. Algumas servem para avaliar a \emph{sciencia}, e outras a \emph{veracidade} das testemunhas.\par Ll. \emph{Regras relativas} \& \emph{sciencia das testemunhas}. É certa a \emph{sciencia}. das\pgn{217}testemunhas, se os seus sentidos, no acto da percepção, --- a) estavam sãos e bem dispostos. --- b) e se foram attenta e diligentemente applicados.\par a) Deve admittirse que os sentidos das testemunhas, no acto da percepção, estavam sãos c bem dispostos, --- 1.º) se não constar o contrario; porque a integridade e a bóa disposição dos sentidos é natural, a0 passo que o defeito é accidental e conhece-se; --- 2.º) se as testemunhas forem muitas e concordes, porque então toda a duvida é infundada; --- 3,9) se o facto era de tal natureza, que podia ser facilmente percebido mesmo por sentidos defeituosos, como seria a existencia de uma cidade.\par b) Deve idmittirse que os sentidos das testemunhas foram \emph{attenta}- : diligentemente applicados, .---. 1.\textasciicircum{}).s6.0. facto. foi tão. sensivel, que \emph{podia} se percebido sem muita diligencia e sem medo de errar; --- 9.º) se o facto foi i\&o publico e notavel, que devia naturalmente attrahir as attonções..de todos; --- 3.º) se entre as testemunhas havia algumas que eram dotadas de --- prudencia e perspicacia, ou que tinham interesse e motivos particulares para bem aveciguarem o facto,\par ll. \emph{Regras relativas} \& \emph{veracidade das testemunhas}. --- E certa a \emph{veracidade} das testemunhas nos seguintes casos:\par \emph{a}) Quando as testemunhas são muitas e diversas entre si na edude e nos caracteres, na. patria e na religião, nos estudos e nos talentos, na educação e nos costumes, e todavia affirmam todas a mesma coisa: porque, em cireumstancias tão diversas, não só toda a falsidade é moralmente impossivel, mas tambem essa uniformidade não pode derivar da má fé. --- b) Quando as testemunhas são poucas, mas narram um facto na. presença de pessoas, ás quaes o facto é nocivo e que podem descobrir o engano; porque ninguem, ainda que mau, em taes circumstancias, falta a verdade. ! \emph{c}) Quando as testemunhas são poucas, mas narram um facto que é contrario ás opiniões, que ellas professam, e ao partido, a que pertencem ; porque, neste caso, só a verdade tem a. força de prevalecer e manifestar-se.\par d) Quando as testemunhas são poucas, mas a narração de um facto, em vez de procurar-lhes riquezas, glorias e prazeres, lhes pode causar ignominia, pobreza, tormentos e morte; pois é necessario acreditar em testemunhas, que sabem morrer pelo que attestam.}. E

189. O testemunho humano-historico e os \textbf{factos}: sobrena-

turaes. --- O testemunho \emph{hwumano-historico} pode ser meio infallivel para a acquisição da verdade, não só quando refere factos de ordem \emph{naífural}, mas tambem. quando. refere factos de ordem. . sobrenatural, como são os \emph{milagres}; porque, tambem na relação dos factos sobrenaturaes, a \emph{sciencia} e a veracidade das testemunhas podem tornar-se evidentemente certas \footnote{Os incredulos, como \emph{Dupuis}, \emph{Volney}, \emph{Hegel}, \emph{Strauss}. \emph{Renan}, --- applicam aos factos sobrenaturaes o systema do \emph{symbolismo} OU \emph{mythismo}, que foi inventado por Vice, \emph{Herder}, \emph{Nieburh}, e que não reconhece nos factos mais \emph{antigos} senão \emph{symbolos}, \emph{allegorias}, \emph{personificações} de \emph{ideas}, --- e aflir- . mam que o facto, chamada \emph{sobrenatural}, não passa de um facto \emph{natural}, ao qual a imaginação dos povos deu forma e caracter de sobrenatural. --- Como\pgn{218}se vê, os incredulos consideram os factos, não como estes existem na realidade, mas como existem na cabeça d' elles, e não conformam as suas ideas com a historia, mas conformam a historia com as suas ideas..\par Tendo provado.a veracidade dos \emph{meios intrinsecos} c \emph{extrinsecos} para a acquisição da verdade, devemos fallar da \emph{critica}, que Manuel \emph{Kant} fu d'esses meios, sobretudo da \emph{razão}. Quem não conhece o \emph{eriticisimo} de \emph{Kant}?\par O escopo de \emph{Kant} é o de examinar e verificar o valor cognitivo da nossa razão, Diz assim: «Devemos fazer a critica, \emph{não} dos livros nem dos systemas, mas do poder da razão em geral, em relação a todos os conhecimentos, - que ella procura attingir fóra da experiencia» (\emph{Critica da vacio pura}. Inrod.) O problema, portanto, que elle se propõe resolver, é o seguinte: Os actos da razão. que. pela sua natureza, parecem representar um objecto e\par rrósistirelménte-a-adunittir-fóra-de-nós. um-Lermo- distin. ele, . que- nDos-levaàn\par direito de j i ülgar as coisas ts objectivamente verdadeiras e » dispostas em si, como\par as julgamos ? --- Em menos. palavras :. Os objectos são em si mesmos, como parecem a nossa mtelligencia? --- Eis o problema. Vejamos a resposta.\par Começa \emph{Kant} por analysar os vossos .conhecimebtos e diz que todos elles resultam compostos de dois elementos: um particular, variavel, contingente ; outro universal, immutavel, necessario. O primeiro deriva da experiencia, e pode chamar-se \emph{materia} ; o segundo emana espontaneamente da propria constituição do sujeito coxnoscente, e diz-se \emph{forma}. --- O duplice elemento encontra-se em todos os actos das nossas faculdades coguilivas. --- Estas sio tres: a \emph{sensibilidade}, a. intelligencia e a \emph{razão}. .\par \emph{a}] A \emph{sensibilidade} (cuja critica chama-se \emph{Esthetica tra lental}) Pre. cede a intelligencia; emquanto é uma condição essencial. para o acto e desenvolvimento: d* esta faculdade..--- A \emph{materia} da \emph{sensibilidade} é consti.\par tuida pelas \emph{impressões}, ou sensações, que hão de ser elaboradas; e a \emph{forma}, á qual a mesma sensibilidade sujeita essas impressões, ou sensações, para que sejam percebidas, é duplice: a do \emph{espaço} e do \emph{tempo}, --- á primeira para as impresgões externas, a segunda para as/infernas. Na verdade, a analyse do conhecimento sensitivo mostra que nelle intervem um elemento \emph{a-priori}, isto é, que não depende da experiencia, mas deriva espontaneamente do sujeito. As nossas representações dos \emph{corpos} não são possiveis senão emquanto estes nos apparecem no espaço. O espaço. pois. é como uma \emph{forma}, em que os sentidos “apresentam os objectos e que por isso deve \emph{preexistir} ao conhecimento, deve ser \emph{a-priori}. Como tambem as impressões \emph{internas} não podem perceber-se senão \emph{simultanea} OU successivamente, isto é, senão no \emph{tempo}: e por isso tambem a forma do \emph{fempo} deve ser \emph{a-priori}. --- D' onde se vé que o \emph{tempo} c O espaço não são coisas \emph{objectivas}, independentes do sujeito cognoscente ; mas são formas \emph{subjectivas}, São \emph{moros de} ser da sensibilidade, que tornam possivel \emph{o} conhecimento das impressões externas e internas.\pgn{219}Estas formas imprimem-se na materia, e revestem-na naturalmente das qualidades do proprio sujeito cognoscente. Toda o qualquer impressão ou sen----- sação vem coar-se neste duplice molde. --- A \emph{forma} tanto do espaço como. do \emph{tempo} é intuição \emph{pura}. É intwigao. emquanto é uma representação no espaço e no \emph{fempo}; mas. não é conceito, porque não é universal; --- e é pura, emquanto é \emph{a-priori}, emquanto: representa as leis \emph{necessarias} das nossas sensações, que, se são externas, só se realisam no \emph{espaço}, C, Se são internas, só podem realisar-se no \emph{fempo}. --- D' onde se vê tambem a natureza do \emph{phenomeno}, Este é a propria. \emph{impressão}, Ou sensação, emquanto se torna objecto de \emph{intuição}, isto é, emquanto é \emph{situada} no \emph{espaço} e no \emph{tempo}. Portanto, para a formação do \emph{phenomeno}, é necessaria a união da impressão com a \emph{ficio}. --- O conhecimento da \emph{sensibilidade} aperfeiçõa-se na \emph{intelligencia}. .\par \emph{b}) \emph{A intelligencia} (cuja critica chama-se \emph{Analytica} transcendental) é superior á sensibilidade-e-torna----perfeita---o---scientifica----s-experiencia--- DO nxeto. proprio e especifico da inteligencia . o \emph{juizo}. \emph{A materia} do juizo é constituida pelos \emph{phenomenos}, e é transmittida pela sensibilidade; a \emph{forma} é o modo por que o predicado convem ao suiejto, e é subministrada pela propria intelligencia. \emph{O juizo}, pois, resulta da: applicação da \emph{forma} aos phenomenos. --- \emph{A forma é a-priori}; independente de toda a experiencia. De facto, ba\textasciicircum{} juizos, que são o resultado de uma \emph{synfLese} (porque o predicado não está contido no conceito do sujeito) e que, ao mesmo tempo, têm umas notas de necessidade, São os juizos \emph{syntheticos a-priori}. Não podendo essas notas provir de um elemento \emph{empirico}, ou experimental (pois. o objecto da experiencia é naturalmente \emph{contingente} e \emph{singular}), devem provir de um elemento \emph{a-priori}, subjectivo e anterior a toda a experiencia. --- Se a \emph{forma} é o modo por que o predicado convem ao sujeito, é claro que basta determinar todas .as especies do juizo, isto é, definir, abstrahindo do contendo ou da materia, as normas, que a intelligencia segue nos seus juizos, para se conhecer o... numero das \emph{formas} puras, latentes no fundo do nosso espirito. Ora as di- 'versas especies do juizo, a que se reduzem todos os juizos possiveis e cada. uma das quaes dá origem a um conceito, são \emph{doze}. --- De facto, os juizos podem variamente formar-se quanto á \emph{quantidade}, \emph{a qualidade}, á relação, á \emph{modalidade}, --- Quanto á \emph{quantidade}, o juizo pode ser \emph{singular}. particular, \emph{universal}; d' onde O \emph{conceito} de unidade, de \emph{pluralidade} ou \emph{multidão}, de \emph{tutalidade}. Quanto á qualidade, o Juizo pode ser \emph{affirmativo}. \emph{negativo}, \emph{limiialivo} ou \emph{infinito} (indeterminado): d' onde O \emph{conceito} de realidade ou \emph{affirmação}, de \emph{negação}, de limitação. Quanto a relação, o juizo pode ser \emph{categorico} (se affirma ou nega absolutamente), \emph{hypothetico} (se inclue uma relação .de causalidade e de \emph{dependencia}), \emph{disjunctivo} (se importa uma relação de\par mutua dependencia); d' onde o conceito de substancia € \emph{modo} Ou \emph{accidente} (0 modo está para a substancia, como, no juizo \emph{categorico}, \emph{o} predicado está ' para o sujeito), de causa e \emph{effeito} (o \emph{effeito} depende da causa, como o condicionado depende da condição), de acção e \emph{paixão} (ha \emph{reciprocidade} entre o\pgn{220}agente e o paciente); Quanto á \emph{modalidade} (isto é, a \emph{qualidade} da relação entre os termos), o juizo pode ser. \emph{problematico} (se afirma uma relação \emph{possivel}). \emph{asserlorio} (se affirma uma relação \emph{real}), \emph{necessario} (se affirma uma relação \emph{necessaria}); d' onde O \emph{conceito} de \emph{possibilidade} e \emph{impossibilidade}, de \emph{existencia} e \emph{não-existencia}, de \emph{idade} e \emph{contingencia}. Sendo \emph{doze} as varias especies do juizo; \emph{doze} são tambem as \emph{formas puras} e OS" \emph{conceitos} da. intelligencia. --- Xstes \emph{doze conceitos}, producto do acto subjectivo do juizo, são as \emph{classes supremas}, ás quaes podem e devem reduzir-se todos os nossos conceitos, e por isso chamam-se \emph{categorias}. --- Portanto, as \emph{categorias} não representam os \emph{medos da existencia}. mas os \emph{modos do.juico}; são conceitos \emph{puros a-priori}, necessariamente logieos; são leis subjectivas da actividade intellectual: s\&o formas vazias, moldes ou typos sem conteudo, mas capazes de receber a materia. fornecida pelos sentidos. e de a modelar de um modo e não de outro. As categorias appliceam-se aos objectos, percebidos pela sensibilidade e: torüaim possiveró pensamento d\textasciicircum{} "elles: Nuda pode ser objecto do nosso pensamento.ou juizo.seuio..por..nreio. de uma ou outra d' estas formas, e em conformidade com ellas; de modo que a razão, por que um objecto se torna intelligivel e se representa: de uma modo determinado, depende exclusivamente dos \emph{conceitos} ou das \emph{categorias}, Assim, vão é a intelligencia que recebe a forma do objecto, mas '6 o objecto que recebe a forma da intelligencia; isto é, não é a intelligencia que depende do objecto, mas é o objecto que depende da intelligencia. --- As \emph{categorias} estão como que em germen ou em potencia em todas as intelligencias. Quando se apresenta um phenomeno, ellas brotam espontaneas do espirito e applicam-se inslinctivameute ao proprio phenomeno, adaptando-o.2 si mesmas \emph{e} tornando-o cognoscivel e conhecido. Assim, quando percebemos que \emph{o} sol aquece a. pedra, salta do fundo do espirito a categoria de causa e de \emph{effeito}, que se imprime sobre essa percepção e a torna intelligivel, formando-se este juizo geral: \emph{desde que o'sol} appareca, a pedra \emph{deve} aquecer, Tal é o conhecimento iutellectua]: uma synthese, em que\textasciicircum{}/o espirito applica ao phenomeno uma determinada categoria. --- Por isso, o officio das categorias não é o de representar os objectos, mas o-'de encadear c coordenar as representações, A \emph{sensibilidade} subministra.,os phenomenos: as \emph{categorias} dão a forma, que os relaciona e coordenà por uma necessidade subjectiva..--- Mas o conheci- « mento intellectual só se torna perfeito pela \emph{razão}.\par \emph{c}] \& \emph{razão} (cuja critica chama-se \emph{Dialectica transcendental}] é uma faculdade superior á intelligencixn. A sua funcção é-a de estabelecer os \emph{principios universaes} e \emph{absolutos}, que devem verificar-se tambem nos entes reaes fóra de nós e em todos os entes possiveis. Para isso, ella sujeita a algumas \emph{formas innatás}, ou '\emph{a-priori}, OS \emph{juízos} e OS \emph{conceitos} da jntelliyencia, para os coordenar e unificar. D' ahi umas syntheses, compostas d' essas \emph{formas} e A esses \emph{juízos} Ou \emph{conceitos}, --- AS formas \emph{innatas}, ou a \emph{priori}, são tres \emph{ideas}; --- \& idea do eu, isto é, da totalidade dos phenomenos \emph{internos} (idea \emph{psychologica}); --- a idea do mundo, isto é, da totalidade dos phenomenos \emph{exter}-}.

, 190 \textbf{Transmissão} do \textbf{testemunho humano} - \textbf{historico}. --- O testemunho \emph{humano-historico} transmitte-se por \textbf{tres} meios: pela - \emph{tradição}, pelos monumentos, e \textbf{pela} \emph{historia}.

a) \emph{Tradição} é a. transmissão. \emph{oral} de um facto por uma serie \pg{219}  não interrompida de testemunhas. --- Será evidentemente manifesta a \emph{sciencia} e a \emph{veracidade} das testemunhas tradicionaes, se a \emph{tradição} fôr --- a) \emph{perpetua}, de modo que \textbf{seja} como uma .. cadeia formada de tantos anneis, quantas são as gerações das \pg{220}  testemunhas, e extendida entre nós e a epocha em que o facto se -, deu, sem nenhuma interrupção; --- 5) \emph{uniforme}, isto é, invariavel e constante quanto á substancia do facto; '--- \emph{c}) ampla, por forma \pg{221}  que cada linha da serie seja representada por um grande nu- \emph{HT} mero de testemunhas; --- \emph{d}) \emph{importante}, isto é, que se refira a = factos publicos e notaveis* --- A ausencia ou o silencio dos do- '

nos (idea \emph{cosmologica}); a idea de \textsc{Deus}, isto é, da totalidade de todas as ] realidades e perfeições (idea \emph{theologica}). : ] v

Da analyse das tres faculdades cognitivas segue-se que nós conhecemos os \emph{phenomenos} e a \emph{ordem}, com que se succedem, --- mas não conhecemos o que está debaixo d'estas apparencias, d' estes phenomenos, não conhecemos a coisa em si, O nowmenó \emph{A} sensibilidade attesta que \emph{Ra} alguma coisa' É atraz do phenomeno, porque o phenomeno é o nosso modo de representar 2 a coisa em si, mas não diz o que essa coisa é. Nem o dizem a intelligencia .-. \emph{Ps} e a razão, ambas limitadas pela experiencia \emph{sensitiva}; pois as \emph{categorias} e o. os \emph{principios} só se referem aos \emph{phenomenos}. --- A razão, por si, sente-se ES inclinada a \emph{transcender}, a elevar-se fóra e acima dos phenomenos, procurando um principio primeiro objectivo, e assim conclue para a existencia da \emph{alma}, do mundo e de Deus. Mas, se ella chega a taes conclusões, nem por isso o é o seu principio. Por quanto, se se pretendesse. que a realidade fOBBe- c \textbar{} conforme com: os discursus-da-razão,-seriam-- inaumeravels ---\emph{as-an£énomáias}, Hu ou contradicções, em que ella havia de cabir fatalmente. . .

As principaes \emph{antinomias} são as quatro seguintes : " : d

1.3) O mundo £\emph{eve}. começo no espaço e no tempo, ou não \emph{teve} ;.

2.3) O mundo resulta de elementos simples, ou \emph{compostos} ;

3.3) A liberdade \emph{existe} no imundo, ou não \emph{existe};

4.º) \emph{Existe} uma causa necessaria do universo, ou \emph{não existe}.

Em cada uma d' estas proposições \emph{disjunctivas}, a these está em contradicção com a antithese; e todavia os argumentos, que se adduzem em - \textbar{} --- favor de uma, têm o mesmo valor dos argumentos, adduzidos em favor de \textbar{} * outra, Como poderá a razão sahir d' esta obscuridade e livrar-se d'esta ter- ------ --- tura? O remedio é facil. Essas \emph{proposições} são contradictorias na ordem dos ' noumenos, isto é, das realidades objectivas, mas não são taes na ordem dos phenomenos, isto é, das apparencias. Se se supprimir a \emph{realidade objectiva}, 7 que se oppõe a conciliação e a verdade de duas proposições contradictorias, cessará toda a contradicção; já não haverá motivo para controversias; a these e a antíthese serão ambas verdadeiras; haverá a verdade subjectiva, que n consiste na \emph{concordia das ideas}; e isto basta. « Então as duas partes ficarão E convencidas de que, se podem confutar-se nrutgamrente, é porque \emph{discutent} , \emph{por} um nada, e uma certa apparenceiu transcendental lhes representou (no : no tempo e no espaço) uma realidade, que não exislia.. A questão não existe, \textasciicircum{} porque não existe um objecto determinado; a questio portanto é essencialment inutil, destituida de sentido.... O vosso \emph{objecto} só existe na vossa cabeça é não pode existir fora d' ella» (\emph{Critica da razão pura}).

Todavia as verdades, que se referem á \emph{alma}, ao mundo, a Deus, e que a razão especulativa Ou \emph{iheorelica} não pode demonstrar, são \emph{promulgadas} " pela \emph{razão pratica}. Por quanto, á \emph{razão pratica}, inculeando que se \emph{deve evi}- . \emph{Ps tar o} mal e \emph{fazer o bem}, deve \emph{implicita} e necessariamente admittir a \emph{Zberdade} e a \emph{inunortalidade da alina}, \emph{a realidade do mundo é au existencia}, \emph{de} dH.

cumentos é uma razão exclusivamente \emph{negativa}, que \textbf{nada} prova contra a \emph{tradição}.

b) \emph{Monumentos} . são obras de arte mais duradoiras, que ser-

\textsc{Deus} ; porque, se não se admittirem estas tres verdades, n\&o é é- possivel evitar o mal e fazer o bem. --- A \emph{razão pratica} não \emph{demonstra} estas verdades, attinge-as pela fé, isto é, por umz'certeza \emph{subjectiva} e destituida de \emph{evidencia}; de modo que a alma, o mundo, e Deus se apresentam conio \emph{postulados}, necessarios e inexlicaveis. Assim a fé toma o logar da \emph{sciencia}. .

Tal é a theori\&/de \emph{Kant} ácerca do conhecimento humano. A coisa é \textbar{} abstrusa, mesnio porque o escriptor allemão teve o gosto --- de emprega algumas palavras extranhas para exprimir as linhas fondámentaes do \emph{systema}, --- de mudar a significação usual de muitas outras, ---.e, sobretudo, de.ser obscuro para se mostrar profundo. A sua vaidade. levou-o a comparar-se com Coperuico; porque, assim como Copernico fez uma revolução na . Astronomia, mostrando que não é o Sol que gira em volta da terra, mas é x terra que gira em volta do Sol, assim tambem elle, Kant, fez uma teve- Iuc\&o na Philosophia, mostrando que não é a intelligencia que se accommoda. ás coisas, para as conhecer.como são em. si mesmas, mas são as coisas que se: accommodam' a Inteligencia, revestindo as formas, que ella lhes applica.----- Todaávia, nem-o--merito-da: invenção é devido a Manuel Kant. Ha muitos seculos, Protagoras tinha dito: Q espirito .é a medida (\emph{forma}) das coisas, e não são as coisas a medida (\emph{forma}) do espirito (Cf. Aristot., XT \emph{Metaph}., c. 5); e por este seu scepticismo foi exilado de Athenas. Nos tempos mais recentes, e anteriormente a Kant, tinham ensinado, com menos exaggero, embora com maior incoherencia, esta theoria, entre outros. Descartes, Leibnitz, Malebeanche, Reid, Cousin. sobretudo Hume, tanto que escreveu \emph{Compayré}: «Ce n'est pas Kant, \emph{c} est Hume qui a fondé la pbilosophie critique» (\emph{La philosophie} de \emph{David Hume}). --- Não é nossa intenção de refutar aqui todas e cada uma das allirmações de Kant: algumas já foram refatadas, e outras o serão no decurso d' este livro.

\emph{M}

O \emph{criticismo} de Kant é talso --- nos. seus \emph{fundamentos}, --- na sua \emph{essencia}, --- nas suas \emph{consequencias}. "

A) É falso nos seus \emph{fundamentos}.

a) O systema parte de um principio, que é a \emph{duvida} do \emph{valor} da razão: e, para averiguar se essa duvida é, ou não, attendivel, serve-se da mesma razão, e por isso já suppõe come certo o que dá por duvidoso. A contradicção é manifesta, --- Se é Balural e legitimo que o nossa razão. supposita sempre a propria aptidão para alcançar a certeza, reílieta sobre os seus actos e reconheça os processos, que seguiu na investivação da. verdade; é ilegituno e absurdo que se pretenda formar raciocinios e -Oorganisar systemas, quando não só se duvida da legitimidade da razão, mas se sustenta que tudo é uma mera illusão. .

5) Um dos fundamentos do \emph{criticismo} é a lei da immanencia. Diz o eseriptor allemão que o conhecimento, sendo um \emph{facto da} consciencia, não pode attingir o que está fora d'ella, e por isso é impossivel a percepção di \emph{coisu em si mesma}, do nowmeno. --- Mas essa base é insubsisteute. O conhecinento é um \emph{facto da consciencia}, é um \emph{lacto subjectivo} quanto ao seu \pg{223}  vem para perpetuar a memoria de uma” pessoa ou de um acontecimento notavel; taes são as estatuas, as columnas, os templos, as pinturas, as medalhas, etc. --- Os \emph{monumentos}, \textbf{para} que te-

\emph{principio}, que é a faculdade cognitiva: mas não é \emph{subjectivo} quanto ao \emph{fer}- . \emph{mo}, que é o objecto conhecido. Uma coisa conhecida \emph{em si} mesma não se altera pelo facto de ser conhecida, mas permanece sempre no seu estado de absoluta; porque não é ella que se transforma na intelligencia, mas é a intelligencia, que, actuada pela forma da coisa; lhe se torna semelhante, attingindo-a como é em si mesma, de modo que ha adequação entre o conceito da intelligencia e a realidade da coisa.

c) Kant admitte tambem que o nosso espirito, extrahindo do proprio fundo a verdade toda, é \emph{autonomo}, isto é, independente de todo o objecto e de toda a. regra extrinseca da verdade, e por isso deve rejeitar todo o elemento, que não é \emph{immanente} e \emph{autochtone}, quer nas construceções das sciencias physicas, quer na elaboração das sciencias logicas e moraes. --- Mas, se assim fosse, como opina Kant, para que haviam os sabios de suar tanto para interrogar a natureza e arrancar-lhe \emph{os} segredos? E se elles julgam .. indispensavel melter-se em communicação com os objectos externos, á razão - € porque essa cominunicag\&o -fornece-dados,-quezpodemos-conbecer. --- Se. o nosso espirito não é independente dos objectos” externos, nem é independente das leis logicas e moraes. Estas leis dominam o homem, impõem-se jmperiosamente a sua intelligencia e a sua vontade. Se são \emph{inferiores}, emquanto se tem consciencia d'ellas, são \emph{exteriores}, emquanto se impõem a cada um de nós. O homem percebe-as como. verdadeiras, mas não é elle que as torna taes. A ordem das coisas, que essas leis exprimem, é essencial e immutavel, e ficaria sempre a mesma, embora a razão do homem se escurecesse pa loucura, ou a sua vontade se corrompesse no vicio. '

\emph{d}) A \emph{autonomia} de nossa intelligencia, se é uma gloria, é tambem, segundo Kant, uma necessidade, visto que, na percepção, o sujeito, que percebe, e o objecto, que seria percebido, são separados por um abysmo insuperavel. --- Ora é [also que o objecto seja sempre inaccessivel a faculdade percepliva, Quando se trata da percepção do \emph{sentido} commum e da \emph{conscien}- \textasciicircum{} cia, o sujeito e o objecto não somente não são distinctos, mas idendificam-se, É impossivel que o sujeito se separe do objecto e que o não attinja. O sujeito consciente percebe-se a si mesmo Immediatamente, como umu existencia e uma realidade certissima, e descobre em si mesmo uma multidão de actos vitaes, sensitivos, intellectuaes, de ideas objectivas. Dizer com Kant que'nos pereebemios a nós mesmos atravez de formas \emph{a-priori} da nossa mentalidade; que alteram todos os objectos, mesmo os objectos internos da consciencia. é emittir uma hypothes \emph{gratuita} e \emph{contradictoria}: --- \emph{gratuita}, porque não só não se baséa em nenhuma experiencia, mas é contraria a toda a experiencia, --- con£radictoria, porque uma percepção sem um objecto percebido é uma conlradiegão... Quando se trata da percepção \emph{sensitiva externa}, tambem neste caso o sujeito e o objecto não estão distantes. O \emph{acto} do \emph{agente} encontra-se no \emph{paciente}, do modo que o agente e o paciente se . acham reunidos num todo \emph{completo}. Sendo assim, é claro que \emph{o} paciente, ou O sujeito, não precisa de suhir de si mesmo, \emph{fóra da} sua \emph{pelle} (\emph{como dizem} \pg{224}  - nham auctoridade, devem ser, --- \emph{a}) \emph{aulhenticos}, isto é, perten-

" centes ao tempo, ao logar, ás pessoas, a que são attribuidos; --- \emph{b}) \emph{contemporaneos} dos factos, que elles representam, ou das

os adversarios) para perceber o agente, ou o objecto. É, pois, falso que .entre o objecto e o sujeito corra um abysmo, que não é poseivel transpór.

B) É falso na sua essencia,

3) A essencia do \emph{criticismo} consiste, principalmente na pypotliese das \emph{formas} innatas, \emph{a-priori}, que existiriam em germen no fundo do espirito e . de que derivariam a \emph{determinação}, \emph{a necessidade}, a universalidade, para os ] \emph{phenomenos} da \emph{sensihilidade}, para os \emph{juizos} da \emph{intelligencia} e para os \emph{principios} da \emph{razão}. Ora essa hypothese é \emph{gratuita}, \emph{inutil}, \emph{irracional}.

a) É \emph{gratuite}. Ninguem tem consciencia da existencia de taes formas: e Kant não adduz nenhum argumento serio. . Existem certamente na intelligencia humana conceitos universaes, que se applicam ás coisas singulares: : mas estes não são meras illusões, não derivam \emph{a-priori}, espontaneamente, do fundo do espirito, mas têm o seu fundamento num principio extrinseco á nossa intelligencia, nas proprias coisas singulares, e. por isso não são formas.impostas ás. coisas pela inteligencia. mas. são formas impostas a intelligencia pelas coisas. : \emph{b}) \emph{É inutil}. Kant-admitte- as formas -\emph{a-priori}, para explicar as notas de \emph{universalidade} e de \emph{necessidade}, que se encontram nos nossos pensamentos. Ora, para isso, não era necessario recorrer a. essas formas \emph{a-priori}. Desde o momento que se reconheça, como é necessario .reconhecer, que a nossa intelligencia tem a força de penetrar no amago das coisas e descobrir nellas umas notas de \emph{universalidade} e de \emph{necessidade}, percebe-se immediatamente que o conceito \emph{universal} e 0 necessario tem o seu fundamento nos objectos \emph{singulares} e \emph{coniinjgentés}. --- O que se diz dos \emph{conceitos}, deve applicar-se aos : \emph{juicos} e aos \emph{principios}, que, sendo constituidos pelos \emph{conceitos}, devem ter e E têm o'seu fundamento, não numa necessidade \emph{subjectiva}, mas numa realidade \emph{objectiva}, extrinseca a0 acto cognitivo.

\emph{c}] \emph{É irracional}. Por quanto. todos os homens estão convencidos de que alguns conceitos, como o da \emph{quantidade}, da \emph{qualidade}, etc., não são coisas \emph{subjectivas}, como aifirma Kant, mas representam coisas \emph{objectivas}, independentes do nosso pensamento, --- Alem d' isso, se os conceitos \emph{universaes} e \emph{necessarios} não podessem derivar da experiencia, por ser esta singular e . \emph{contingente}, nem poderiam derivar do nos\&o espirito, que tambem € \emph{singular} d e contingente; porque \emph{omne aens}. \emph{a git sébi simile}. --- Mais: os \emph{conceitos} de \textbar{} Kant são \emph{puros}. isto é, sem vontettdo; sem materia; ora \emph{o} conceito é-necessariamente o conceito de alguma coisa, e por isso “deve ter um conteudo: \textbar{} alias deixará de ser conceito, e será uma especie de \emph{circulo quadrado}. D p) Kant admitte duas formas \emph{énnatas} e \emph{absiracias} na \emph{sensibilidade}: a. do espaço e a do \emph{tempo}. Antes de tudo, não se concebe como é que formas \emph{abstractas} possam ter por principio e por sujeito uma faculdade material e concreta, como é \emph{a sensibilidade}. --- Além d'isso, todos os homens estão convencidos de que a experiencia sensivel percebe certas determinações, ou formas, e, entre outras, a do \emph{espaço} e do fempo, como \emph{inkerentes} aos objectos externos, e independentes do nosso neto perceptivo. Separar, como faz Kant, da forma \emph{pura} do phenomeno, isto é, da forma subjectiva do «spaço e do \pg{225}  pessoas, a que se referem; --- c) \emph{conformes} ás tradições \textbf{legitimas}

e aos documentos certos.

c) \emph{Historia} é a narração, methodicamente \emph{escrzpta}, de \textbf{factos} notaveis para instrucção. da. humanidade. --- Um. livro \emph{historico}... será meio seguro para a acquisig\&o da verdade, se fôr \emph{geruino},

\emph{tempo}, á impressão sensivel, que procede do objecto; é uma arbitrariedade sem fundamento. --- Mais: na theoria de Kant, a materia das nossas sensações, indeterminada por si mesma, é determinada pelas formas. \textbf{innatas} da sensibilidade; e assim a variedade, que se encontra na figura e na duração dos objectos, depende exclusivamente do nosso conhecimento, Ora esta conclusão é contraria á experiencia, a qual attesta --- que a variedade na figura e na duração dás coisas é independente do nosso conhecimento, --- e que, o nosso espirito não pode introduzir nessas representações do \emph{espaço} € do \emph{tempo} á mais leve alteração ou mudança. --- Finalmente, Kaot admitte a existencia da realidade \emph{extensa}, embora d' ella só possa saber-se que é \emph{ex tensa}, e. acerescenta que temos a representação da \emph{extensão} pela applicação da forma subjectiva do \emph{espaço}. Ora isto tambem é contrario á experiencia Não é a idea: da-extensão-yue-derivadaides-do- espaço; muasé a-idenad --- espaço, que, pela abstracção e reflexão, de

\textbar{} \emph{extensão}.

igencia, a ém das observações

e) Com relação ás formas innatas da inte feitas, apresentamos as seguintes.

a) Kant, fazendo a \emph{critica da razão pura}, começa pelo \emph{juizo}, porque \textbar{} diz que \emph{julgar é pensar}, descobre os celebres \emph{juízos syntheticos a-priori}, e, para explicar as notas de \emph{necessidade} e \emph{universalidade}, que se encontram nesses juizos, recorre ás formas \emph{innatas a-priori}, latentes no fundo do nosso espirito. --- Uma arbitrariedade a apoiar outra arbitrariedade \textbar{}... Mas devemos notar o seguinte. Kant começa pelo \emph{juizo} e omitte a questão importantisstma da noção, ou da \emph{idea}. Antes de investigar o modo por que um predicado convem a um sujeito, é necessario conhecer a origem e o valor objectivo d'esses termos. --- Kant occupou-se do \emph{juizo}, do comwosto, e não das partes componentes; e como o juizo não representa a realidade senão pelos seus Lermos, é claro que, prescindindo da realidade das partes, devia tambem prescindir da realidade do composto. É verdade que elle admitte os \emph{concezios}: mas esses conceitos não são \emph{anteriores} ao juizo, são \emph{posteriores}; --- não são noções simples e objectivas, mas compostas e subjectivas, --- No \emph{juizo}, Kant só vé o \emph{subjectivo}. O verbo é não exprime, para elle, o \emph{acto de ser} dos estes reaes, mas. apenas a união ideal de termos puramente \emph{idoces}. \textbf{D'ahi} o. \emph{subjectivisino}. --- É falso que pensar e \emph{julgar} sejam uma e a mesma coisa; i simples apprehensão do objecto não é um juizo, é um pensamento.

5) Na theoria de Kant, a nossa intelligencia applicaria instinctivamente as formas innatas aos phenomenos e sentic-se-hia constrangida por uma fatal necessidade a emittir o juizo. --- Ora. isto repugna a factos averiguados. Áltesla-nos \& consciencia não só que nos sentimos logicamente necessidades a pensar e a julgar de um modo, e não de outro, mas tambem que essa necessidade deriva exclusivamente do objecto extrinseco ao sujeito pensante, Pensamos e julgamos \emph{assim}, porque o \emph{objecto} é é \emph{assim}. Portanto não é uma ' necessidade cega que nos leva x julgar, mas é o conhecimento da relação,

Vol, 1 16

*\emph{inleiro}. e \emph{veridico}. --- \emph{a}) Um livro diz-se \emph{genuino}, se realmente

foi escripto pelo auctor, a quem o attribuem, ou pertence ao

tempo, a que o referem. - No "caso contrario, diz-se \emph{apocripho}. \emph{E} O..livro-.será. genuino-----1.9).se. fôr attribuido.a um certo auctor

; existente entre o sujeito e o predicado, e independente do nosso trabalho \emph{D} intellectual. Uma necessidade mental não pode ser e não é senão a conse- , quencia da visão da verdade.

c) Segundo Kant, as formas énnatas, que são moldes vazios, conceitos

puros, brotam, na oceasião de-uma sensação, do fundo da alma e appli-

cam-se por si aos phenomenos, para communicar a estes a \emph{universalidade}.

--- Mas, perguntamos, d' onde tiram essas formas o poder de tornar univer-

Saes o8 phenomenos ? Como é que, quando vemos um objecto, aquecido pelo

Sol, surge em nôs espontaneo o conceito de causa, e quando vemos uma

arvore, desponta immediatamente o conceito de \emph{substancia}? Dir-se-ha que

a nossa intelligencia, ao apparecer de uma imagem sensivel na phantasia,

produz uma imagem espiritual, correspondente ao objecto externo. Isto é

certo; mas emtanto é certo, emquanto se admitte que o objecto externo,

“---produzind 0a sua imagem sensivel mr-phuntasia, « \emph{concorre} para 3 producção \emph{BU} da imagem espiritual na inteligencia. Por isso.a imagem sensivel é.0 \emph{prinm} cipio e a medida da imagem espiritual. Mas esta explicação da boa philo- \emph{E} Sophia, que é tambem a “do genero humano, é rejeitada por Kant, que não \emph{e} admitte nenhuma influencia do objecto no sujeito e faz das imagens das \emph{LS} - coisas um producto interior e total do espirito.

\emph{d}) No systema de Kant, a \emph{universalidade} do juizo é apen: tS \emph{logica} ou

\emph{subjectiva}, porque deriva da forma \emph{subjectiva}, applicada ao phenomeno. ---

Mas basta uma pequena reflexio: para se concluir que a applicagio das

formas innatas aos phenomenos nem a \emph{universalidade subjectiva} pode com- \emph{DEN} municar ao juizo. De facto, essa \emph{universalidade}, segundo os principios do \emph{po} escriptor ullemão, não pode deixar de ser limitada a tm só \emph{individuo}. Como \emph{ii} € possivel que, quando digo: \emph{o mel é doce}, todos os homens convenham \emph{üT} commigo, e assim o juizo “tenha um caracter de \emph{universal}? Para isso, é

necessario --- que se conheça a existência dos outros homens, isto é, das \emph{iT coisas em si}, dos \emph{nowmenos}, --- que se admitla que todos 16m a mesma orga-

nisacio e tambem: que o mel produz em todas as linguas a mesma impres-

são. Só depois de ter conhecido e averiguado tudo isto é que se pode for-

mular um juizo \emph{subjectivamente universal}. Ora Kant é um homem \emph{a-priori}, 2\emph{i} não admitte nowmenos, não reconhece qualidades determinadas nos objectos; \emph{M} e portanto não pode, logicamente, admittir a \emph{universalidade} subjectiva do \textbar{}: juizo. E então para que servem as formas innatas?

e) Finalmente, a disposição dos conceitos, ou das \emph{categorias}, é arbitraria,

torta, desordenada, irracional, (já o vimos), em que a paixão pela symetria

se sobrepõe a seriedade e ao senso commum. 'Diz \emph{X}. \emph{Fischer}, admirador de

Kant: «Na disposição das categorias apparece claramente a phantasia ar- \emph{PTT} chitectonica de Kant, o seu gosto pela apparencia; é preciso não dar uma \emph{NE} demasiada importancia a symetria, que ahi se ostenta complacentemente. \emph{h D} Como estas categorias são abstrahidas dos juizos, falta a este dodecalogo \emph{INE} de conceitos puros do entendimento a forma de systema, que nio pode ser

---- \pg{227}  por uma tradição nunca interrompida, que vem \textbf{desde} o tempo, em que o livro appareceu, até nós; --- 2?) se o estylo e as circumstancias da narração convierem a pessoa do auctor; --- 3º.) se concordar - com-. os - codices-..mais. antigos ; --- 4.).se-não--- ---

substituida pelos jogos do architectonico ». [\emph{Histoire de la}. \emph{Phsl}. \emph{Mod}., li, 363). --- É um jogo infantil o de Kant, mas vão é um jogo inoffensivo.

d) As considerações, que fizemos ácerca das formas \emph{innatas} em geral, podem e devem applicar-se ás formas \emph{innatas} da \emph{ração}.--- As ideas da alma, do \emph{mundo}, de Deus, são, como todas as outras, «\emph{adquiridas} e \emph{objectivas}; como se demonstirará nos logares proprios d'estes Elementos --- Só notaremos o erro de Kant, que faz da \emph{razão} uma faculdade \emph{distância} da \emph{intel ligencia}, e \emph{superior} á esta, quando esses dois termos significam apenas, como vimos, diversas funcções de uma e mesma faculdade.

Feita a analyse das tres faculdades cognitivas, o escriptor allemão conclue que só conhecemos os ph enos, AS \emph{apparencias}, mas nos é impossivel conhecer as coisas em si mesmas, os nowmenos. Diz assim: «As \textbf{coisas} apresentam-se como objecto existente fóra de nós, mas nós nio pod ' saber o que são -exc-si-mesmas;-S0--conbecemos---os-p 2 ; isto-é representações, que ellas operam em nós, quando impressionam os nossos --- sentidos ». --- Antes de tudo, notamos uma contradiecto de Kant. Affirma que os phenomenos São Os \emph{primeiros} e unicos objectos do nosso conhecimento. Mas, sendo assim, elle não podia determinar os elementos \emph{a-priorz} que constituiriam os mesmos phenomenos; pois não € possivel o conhecimento de uma coisa, \emph{anterior} a todo é conhecimento. --- Alem d' isso, \textbf{este} escriptor admitte, e por palavras muito claras, a existencia do mowmeno, da \emph{coisa} em \emph{si} mesmo; aliás, diz elle, haveria uma \emph{apparencia} sem a coisa, que apparece. Optimamenle:; e já se reconlece, implicitamente, o principio de --causalidade. Mas, se \emph{existe} a coisa, que \emph{apparece}, ella deve apparecer na propria apparicio ; porque não se concebe uma apparição abstracta, isto é, uma apparição separada da coisa, que apparece. Portanto, devemos dizer-ou que com O phenomeno apparece tambem o nouwmeno, ou que não apparece e não existe coisa nenhuma.- - \emph{vM}

A este \emph{subjeciivismo total}, a este \emph{nihilismo absoluto} levam logicamente us principios do auctor da \emph{Critica da razão}. No verdade, o phenomeno, no systema de Kant, sendo constituido pela \emph{i»:wressao}, ou sensação, e pela \emph{intuição}, que são coisas subjectivas, nada pode ter e nada tem de \emph{objectivo}. E não pode dizer-se que a impressão, ou sensação, contenha algum elemento \emph{objectivo}, que refira e apresente o \emph{objecto}, que impressionou os sentidos. Na verdade, para se concluir para a existencia do \emph{objecto}, teria de admittir-se o principio de \emph{causalidade}; visto que, só em virtude d' este principio, podemos passar da impressão recebida no sentido para o objecto, que a produziu. Mas, para Kant, a causa é uma forma \emph{innata}, \emph{a-priori}, sem, valor \emph{objectivo}, e a \emph{causalidade} não existe fóra de nós. Logo, no systema de Kant, não ha nada de \emph{objectivo}, ou, se houver, não pode ser de modo algum conhecido. --- Alem d'isso, se é pela forma do espaço que percebemos o plirenomeno, e se ! a forma do \emph{espaço} é toda subjectiva, porque não havia de ser inteiramente \emph{M} Subjectivo o proprio phenomeno? De facto, este escriptor diz e repete Ó \pg{228}  houver argumentos positivos, que neguem a sua genuinidade: --- b) Um livro é inteiro, quando se conserva tal qual foi escripto pelo auctor. Deixa de ser \emph{inteiro}, quando --- ou por mutilação se lhe tirou alguma. coisa, que o..auctor escreveu (\emph{truncado}), ---

continuamente --- que os objectos da experiencia não passam de illusões ou de representações subiectivas, --- que o phenomeno é um producto exclusivo do sujeito, --- que a impressão é coisa completamente subjectiva, --- que a materia é um phenomeno subjectivo, --- que a propria realidade é um. conceito puro da intelligencia..Cilemos alguns trechos, extrahidos da \emph{Critica}: « Os objectos externos, os corpos são apenas phenomenos, são modos das nossas representações, cujos objectos só existem em virtude d' essas mesmas representações, e, tomados separadamente, não são nada... Às coisas externas, isto é, a materia com todas as suas formas e mudanças, não passam de phenomenos, isto é, de representações em nós».- Como se vê, só temos phenomenos e representações, --- phenomenos, ou apparencias, que nada fazem apparecer, --- \emph{representações}, que nada representam. --- Pelo menos, d' este naufragio universal se salvasse o \emph{eu}! Mas nem elle se salva. O eu, para Kant;,-é-tambem-p 1 1;---e---.exute-eu-ph \emph{al} -não-é-um- eu \emph{real}». \emph{Penso}, \emph{logo} é. Tal Lé o. desvario fundamenta! de Kant.

Apraz-nos citar o juizo que do criticismo de Kant fazem \emph{Hartmann} e \emph{Schopenhauer}. --- Hartmann, mostrando que \emph{o} systema de Kant é uma hypothese \emph{a-priori,sem} nenhum fundamento, escreve : « A simples verdade que tudo o que posso representarzne, não é senão á \emph{minhu} representação, e que tudo o que posso conceber não pode ser outra coisa senão a \emph{minha} concepção, --- esta simples verdade é a origem de todo o idealismo subjectivo: todos os seus fundamentos são outras tantas frivolidades, em que se encontra sempre este argumento, mais. ou menos escondido, e é d'elle que os outros tiram toda a sua força» (Examen \emph{crit}. du \emph{realisme trascendental}, p. 27) --- Schopenhauer denuncia as contradiegbes de Kant com as seguintes palavras: « Kant baséa a existencia da coisa em si, embora alravez de mil subterfugios, sobre.um raciocinio, fundado no principio de causalidade; porque a impressio recebida nos nossos orgãos deve ter uma causa extrinseca. Ora, conforme a sua propria descoberta, o principio de causalidade é'por nós conhecido \emph{a-priori}, e por isso é uma funecüo da nossa intelligencia, e tem uma origem subjectiva; além d'isso, a propria impressão sensivel, a qual applicamos o principio de causalidade, é indubitavelmente subjectiva; e finalmente \emph{o} proprio espaço, em que, em virtude d'essa applicação, coltocamos, como objectos, as causas da dita impressão, é uma forma mental \emph{a-priori}, € por isso subjectiva. Portanto toda a intuição empirica existe inteiramente no mundo subjectivo ». (Le monde en \emph{tant} que volonté, 1, 490). --- Segundo o mesmo \emph{Schopenhauer}, quem sustenta que o mundo é apenas um estado do proprio organismo, deveria ser fechado numa casa de saude!

O \emph{eriticismo}, pois, é \emph{um} systema, destituido de base e de seriedade. O nosso conhecimento não deriva de uma actividade meramente subjectiva e productiva do objecto, mas é o resultado da intelligencia e do proprio objecto conhecido. --- Nem, admittindo a perfeita correspondencia entre o objecto e o conhecimento, se verificam essas obacuridades e torturas da razão, deseriptas

E ou por interpolação se lhe accrescentou uma coisa, que o auctor não escreveu (\emph{inlerpolado}) O livro será \emph{inteiro}, --- 1º.) se o estylo, a forma e as doutrinas forem sempre coherentes; --- 2\textasciicircum{}) se com córes aterradoras por Kant. Às taes antinomias são tambem formas a \emph{priori}, existentes unicamente na cabeça do escriptor allemão, e só elle teve a habilidade de encontrar argumentos, dotados do mesmo valor, em favor de cada uma das partes contradictorias d'essas proposições disjunctivas, No decurso d'este livro, teremos occasião e necessidade de mostrar que os argumentos favoraveis a these não o são a antilhese, e por isso que a intelligencia deve acceitar uma parte das antinomias, e excluir outra. --- Só acerescentaremos que o remedio, que foi proposto por esse escriptor para alliviar as torturas da razão e que, como vimos, consiste na suppressão da renlidade objectiva, não é nada eficaz para o escopo. Porquanto, essa suppressão, se faz dasapparecer o conflicto do espirito com o objecto, não faz desapparecer, antes intensifica o conflicto do espirito comsigo mesmo, emquanto elle é condemnado a uma perpetua. illus\&o, --- recebendo ou formando representações, que nada representam, --- julgando como objectivas as pro- . priedades das coisas externas, que a final são apenas modos da precepção e affecqóes do sujeito, --- e admittindo-a identidade e a verdade simultanea de duas proposições contradictorias. E será esta. uma concordia de-ideas?!... Não admira portanto, se a solução de Kant não satisfez ninguem. Alguns dos seus sequazes, mais coherentes que o mostre, despenharam-se no fundo de um absoluto scepticismo ; outros, levados pelo instincto da reacção, abracaram o mais intempoerante dogmatismo. --- Mas Kant confiava tanto na sua descoberta, que desejou applical-a tambem ao regime das nações. Dizia que a metaphysica divide os povos, porque admitte a realidade objectiva, e que, supprimida essa realidade, já não haveria motivo para luctas entre os diversos paizes. Para esse systema. de \emph{pacificação}, chegou a pedir o favor dos Governos d' então; ninguem attendeu ao pedido: só a Convenção nacional de França lhe responden, honrando-o com o titulo de \emph{cidadão frances} !.,.

Se” Kant não trata bem a razão especulativa, declarando-a incapaz de perceber a realidade e sujeita a perpetuas illusões, não trata melhor a razão pratica. Esta tambem é cega; e, se admitte algumas verdades, como a \emph{liberdade} e a \emph{immortalidade da} alma, a realidade \emph{do mundo}, a \emph{existencia de} Deus, admitte-as, não como conclusões de uma demonstração, mas unicamente como \emph{postulados}, attingidos pela fé, isto é, por uma certeza \emph{subjectiva}, e destituida de toda evidencia. A que estado reduziram estes philosophos Éranscendentaes x pobre razão humana !... Mas escutemos o que diz ácercz d' estes \emph{postulados} de Kant um celebre escriptor, que por fim percebeu toda a vaidade e o absurdo d'este systema. Fallamos de Ausonio \emph{Franchi}, que na. \emph{Ultima critica} escreve as linhas seguintes: « Absurda e impossivel torna-se toda a \emph{sahida} da. consciencia (\emph{razão pratica}) fóra \emph{de si} mesma e toda a. percepção de verdades transcendentes, apezar de todo o recurso a qualquer \emph{postulado}, E por onde havia de ir a consciencia, sahindo fóra de si, se fóra d'ella não ha nada? E como poderia ella conhecer as verdades tránscendenies, se além ou a cima d'ella não ha nada? Os \emph{postulados}, que haviam de fazer o grande milagre. referem-se á existencia de Deus, á liberdade e a immortalidade da alma. Pois bem, postular seriamente a existencia de Deus, \pg{230}  alguns trechos, citados por outros auctores, estiverem conformes com os respectivos do livro, ou só differirem ern coisas de pequena importancia. --- c) Um livro é \emph{veridico}, quando refere os quando se tem dito que \textsc{Deus} com a sua existencia é um abysmo de contradicções, postular a liberdade e a immortalidade da alma, quando se tem dito que \& alma é um ente imeginario, não é porventura \emph{postular} que a philosophia se transforme em comedia, e que a razão se reduza por si, de proposito, ao estado de loucura \emph{raciocinante}? É certo que se alguem raciccinasse d'este modo em qualquer outra materia, nas occorrencias da vida privada, domestica. civil, seria tido por todos, sobretudo pelos crilicistas, como um demente digno de ser recommendado a psychiatria e ao manicomio; e deverá depois admirar-se como eminentemente scientifico e superlalivamente philosophico este modo de raciocinar sobre materias de metaphysica e de psycbologia ? (Vol. Jf, p. 430). .

€) O \emph{criticismo}, falso no seus \emph{fundamentos} c na sua essencia, não é menos falso nas suas consegquencias; por quanto, é a destruição da verdade, da sciencia \emph{experimental}, da sciencia especulativa, da sciencia \emph{pratica}.

a) É a. destruição da. \emph{verdade}. --- A \emph{verdade}, como dissemos, consiste na \emph{conformidade do} nosso pensamento cou, as \emph{coisas}: « \emph{Veritas consistit} in \emph{adaoequatione intellectus} et rei» (Sum. Th.;p-1,q- 21, a. 91 Ella é \emph{immutavel}; porque fixo ou immutavel é o elemento objectivo, com que o pensamento se conforma. Mas, como no.systema de \emph{Kant} não ha uma verdade \emph{ontologica} fixa e immutavel, ou, se ha, é para nós \emph{completamente} desconhecida, nem, por consequencia, pode haver uma verdade logica. --- No entretanto, não querendo supprimir o vocabulo de \emph{verdade}, desvictuaram-lhe o sentido, dizendo que a verdade consiste na \emph{concordia das ideas}, e que por isso é a visão \emph{do que nos parece}, \emph{e} não \emph{do} que é. Esta seria a verdade \emph{subjectiva}, e por isso mutavel e \emph{subjectiva}, em opposição a verdade \emph{objectiva}, e por isso \emph{fixa} c absoluta.: D' ahi a' negação de todo o \emph{dogma}, \emph{philosophico} e \emph{theologico}.

5) É'a destruição de toda a sciencia experimental. --- Por quanto, sc o homem pode perceber unicamente as proprias ideas, e não pode, de modo algum, atlingir os objectos como são em si mesmos, toda a sciencia \emph{experimental} será apenas a sciencia das nossas ideas, será uma. pPsychologia. Diz S. Thomaz: «Si ea quae intelligimus, essent solum species, quae sunt in anima, sequeretur quod scientiae omnes non essent de rebus, quae sunt exira animam, sed solum de speciebus intelligibilibus, quae sunt in anima » (\emph{Sum}. \emph{Th}. p. 1, q. \$5. a. 2). --- É verdade que \emph{Kant} quiz limitar o seu scepticismo, ou agnosticismo, ao dominio do sobresensivel, ou da metaphysica, c pretender: ser dogmatico no dominio da experiencia, quer nas sciencias physicas, quer nas sciencias mathemativas e gcometricas, que tambem são, de algum modo, experimentaes ; mas não é com isso que se salva o valor das sciencias \emph{experimeniaes}. Por quanto, segundo \emph{Kant}, a propria percepção experimental, quer no seu conjuncto quer nos seus elementos, não passa, como vimos, da uma construccio ou creação do espirito, que é determinada pela matureza do sujeito pensante, e que seria differente, se fosse differente essa natureza. Por isso, a sciencia experimental não apprehende o que é na \emph{realidade}, mas apenas os effeitos psycbologicos do nosso espirito, proporcionados á sua extructura e ás suas disposições presentes.

CAP. II

factos, come estes se deram. O livro \textbf{será} \emph{veridico}, se o escriptor possue sciencia e \emph{probidade}. O escriptor possue sciencia, quando os factos, que narra, são publicos, notaveis e \textbf{observados} por

Tudo é, pois, \emph{subjecíivo}. Os sabios, quando dizem que apresentam os" phenomenos astronomicos, physicos ou biologicos, não apresentam realmente senão phenomenmos psychologicos, que não têm relação alguma com os primeiros. Será admissivel tão grande desvario?.. Não, com \textbf{certeza}. Não 6.0 accordo de uma idea com outra idea, que os sabios pretendem atlingir, --- mas é o accordo, cada vez mais exacto e perfeito, das suas \textbf{ideas} com a realidade, com os factos, com as leis da. natureza.

€) É a destruição de toda a sciencia \emph{especulativa}. --- Por quanto, é pos- . sivel por ventura uma sciencia exclusivamente \emph{subjectiva} ? Não é possivel. \emph{P} O accordo de varias ideas \emph{falsas}. embora logicamente ordenadas entre si, não pode constituir a \emph{verdade}. --- Mas nem \textbf{esse} accordo é possivel. Com elfeito, todo o accordo deve ser regulado por principios. Mas, se, para os' \emph{Kantístas}, não existe nenhum principio \emph{absolutamente} certo, a não ser o principio de contradicção, --- e se este mesmo principio é, como os outros principios, um habito cego e innato de alguns espiritos, de que outros podem ser destitiidos, --- Tão temos: a\textasciicircum{}certeza-absoluta-de-que-nao-podemosatiribuir e negar, ao,mesmo.tempo..a. mesma coisa. á mesma idea. E, não . tendo esta certeza, nem podemos estar certos de ter alcançado o accordo entre as nossas ideas, e o proprio accordo é impossivel. D'este modo, passa-se. fatalmente do scepticismo objectivo e parcial para o scepticismo subjectivo e total, para o \emph{nihilismo} do pensamento. A destruição da sciencia \emph{objectiva} .traz comsigo a destruição de toda a sciencia, mesmo da \emph{subjectiva}.

\emph{d}) É a destruição de toda a sciencia \emph{pratica}. --- \emph{Kant} admitte que a lei do dever proclama, no fundo das nossas consciencias, a distincção entre o bem e o mal. Mas perguntamos: As ideas \emph{moraes} de dever, de bem, de mal, etc., são \emph{objectivas}, qu meramente subjectivas? Se são \emph{objectivas}, porque rejeita elle a \emph{objectividade} das ideas \emph{theoreticas}? --- Se são meramente \emph{subjectivas} e \emph{illusorias}, com qual direito poderá elle impol-as a si mesmo e aos outros? \emph{Kant} responde que não pode impor \textbf{essas} ideas a ninguem por meio de uma demonstração; limita-se a impol-as a si \textbf{mesmo} por um acto de fé crê no dever, na \emph{liberdade}, na sancção \emph{divina}, na immortalidade da alma, e crê porque sente a necessidade \emph{de crer}. Por isso esse acto de fé é \emph{cego}, não é racional nem humano. E assim a \emph{Moral}, a mais imperativa das sciencias, funda-se numa crença \emph{cega}, que não pode impor-se a ninguem! --- Foram, portanto, logicos os discipulos de \emph{Kant}, transformando a Moral \emph{obrigatoria em} Moral \emph{facultativa}, e desvirtuando as ideas de \emph{dever}, de \emph{Hberdade}, etc. Para elles, \emph{a} Moral não passa de uma \emph{technica}, de uma \emph{erfe}, sem um ideal certo, sem preceitos certos, sem obrigações, sem sancção. Se uma doutrina não é creadora de forças e de felicidade senão pela certeza, que tem e que dá, couclufmos que o \emph{críticismo}, tendo destruido pela base a certeza, deve aniquilar todas as forças de acção e toda a felicidade da vida humana.

Não podemos deixar de advertir que este sceptcismo de \emph{Kant} não só destroe a verdadeira Philosopbia, mas tambem a "Theologia revelada. Sc o . homem nio pode nem deve acceitar nenhuma regra exirinseca de verdade absoluta e superior a elle; se elle deve rejeitar tudo o que não é o producto \pg{232}  muitas testemunhas; quando nenhum escriptor contemporaheo se levantou para contradizer o facto narrado, principalmente se assim o exigia o interesse de muitas pessoas; quando a tradição

immanente da propria personalidade, toda a auctoridade --- divina e humana --- é posta de parte, como coisa absurda ou inutil. Nada de testemunho externo, nada de Revelação sobrenatural O que se chama \emph{dogma} deve sahir, como toda a crenca e toda a regra de vida, do fundo intimo do homem, deve ser o fructo do esforço vital da alma. A religião, que se diz \emph{revelada}, é o producto espontaneo da consciencia humana, que vai procurando um ser superior. Os dogmas são concepções provisorias, sujeitas a mudanças e \& progressivas evoluções, \emph{tendo} por norma unicamente á lei do progresso \textasciicircum{}" humano. (Cf. \emph{Farges}, La crise de la certitudo; \emph{T}. \emph{Pesch}; le Kanlisme).

Os \emph{neo-kantistas}, para salvar de algum modo a certeza \emph{objectiva} dos nossos conhecimentos, modificaram o systema do mestre, c assim nasceram o voluntarismo, O \emph{pragmatismo}, o \emph{dogmatismo moral}. e o \emph{anti-intellectualismo}. --- Poucas palavras para cada uma d'estas opiniões.

a) Voluntarismo. --- Foi ensinado por \emph{Xenouwvier},---- 0 verdadeiro fundador do meocériticismo, € Seguido por Secretan, -ABirma que a certeza \emph{objectiva} dos nossos.conhecimentos. dependo.--- não-da razão. naturalmente impolenle para alcançar a verdade, --- mas da \emph{livre vontade}. E a \emph{vontade} que faz a \emph{verdade} c o \emph{erro}; de modo que a \emph{liberdade} individual é a \emph{creadora} da verdade scientifici e moral. \emph{m} o .

O \emph{voluntarismo} é falso. Por quanto, esse systema, se explica a existencia de alguns erros (porque o erro depende, ás vezes, do influxo da vontade). não explica a existencia de nenhuma verdade; porque « percepção du verdade é acto exclusivamente proprio da \emph{razão}. Depende, por ventura, do nosso livre arbitrio que os tres angulos de um triangulo sejam eguaes a dois angulos rectos? Não, com certeza, Com muita graça dizia \emph{Fowillóe} ao mesmo \emph{AEenowvier}: «O dia, em que bastasse a um asironomo um acto de livre axbitrio para ver na extremidade do lelesconio uma estrella, seria o dia da ruina da sciencia astronomica.».(L« \emph{liberté} et le \emph{deferminisine}, p. 12). A verdade é'c não pode deixar de ser impessoal.

\emph{b}) \emph{Pragmatismo}. --- Bo systema de \emph{Blondel}, W. \emph{James}, seguidos por \emph{De Saily}, \emph{Tyrrel}, \emph{Denis}, \emph{Loisy}, e outros. Rejeita o valor do conhecimento intellectual; porque todo o movimento \emph{transeunte} do objecto para o sujeito, como é o movimento da percepção objectiva, é impossivel, --- e sustenta que o unico meio para attingir a verdade \emph{objectiva} é-a-praxe; \& \emph{acção}, porque só o movimento do sujeito para o objecto, como é o movimento da vontade, é capaz de attingir o objecto. Escreve \emph{Blondel}: «\emph{A verdade} e O \emph{bem} só se demonstram pela acção e pelas, consequencias moraes» (\emph{L'action} ete.). O \emph{progmatismo} é falso. Por quanto, este systema, ao passo que rejeita o valor do conhecimento intellectual, porque afirma ser impossivel o movimento ou a acção \emph{transeunte} do ubjecto no sujeito, admitte o valor da acção voluntaria, a qual não é senão a acção transeunte do sujeito no objecto. Se a acção \emph{transeunte} é por si impossivel, como dizem os neo-Kantistas, deve ser impossivel sempre, em todos os casos, --- Além d'isso, o \emph{pragmatismo} admitte que o objecto não é percebido no conhecimento \emph{especulativo}, \pg{233}  eis monumentos confirmam a narração. [o eseriptor possue . a \emph{probidade}, quando é homem de bem e' bonesto ; ou escreveu L. por auctoridade publica; ou contou factos contrarios au seu interesse ou a sua gloria-().. .. \emph{cc}: "\emph{e} - \emph{C} --- € ] ] \emph{OT}. mas na \emph{acção livre}. Ora nenhuma acção \emph{livre} pode querer e atlingir uma coisa, que não tenha sido anteriormente percebida pela intelligencia. É \emph{CS} velho mas sempre novo o aphorismo: \emph{mihel volitwm nisi praecognibum}. Sea --- intelligeneia nada mostra a vontade, esta; nada quer. - . :

\emph{c}) \emph{Dogmatismo moral}. --- É «systema de \emph{Laberthonniére}; Allirma ----que ------ não é a intelligencia que alcança a certeza objectiva, mas é a \emph{boa vontade}, - xS --- e que a \emph{utilidade} é o criterio da verdade. E \emph{ES}

O \emph{dogmatisino moral} tambem é falso. Por quanto, se copcedemos que = a \emph{bôa vontade} possa determinar e applicar a intelligencia a um certo objecto, \emph{O} c que por isso é devido a essa \emph{bóa vontade} sc um conhecimento é mais UU claro, se uma convicção é mais forte, se uma opinião escolhida é mais pro- \emph{te} vavol, --- negamos, com o bom senso, qae a \emph{vontade} perceba a verdade. A vontade. é uma feculdade \emph{cega} ; quer, mas não vé o que quer. --- Alem d'isso. --- sc a \emph{bón vontade} fosse a faculdade percepliva da- Verdade, - "os homens demó \emph{vontade} não conheceriam -nenhuma...verdade, ----sobretado...de. ordem... «religiosa e moral: o que é falso. --- Finalmente, nem tudo o que é \emph{ut é O} verdade. Se a \emph{utilidade} pode satisfazer as exigencias do coração, só a. ver- " \emph{dade} satisfaz as exigencias da intelligeneia. . \emph{E}

\emph{d}) \emph{Anti-intellectuatissmo}. --- É a opinião de \emph{Bergson}. Sustenta. que tudo --- o que existe está sujeilo a um movimento \emph{vital} c \emph{continuo} ; e, como a. nossa intelligencia é incapaz de perceber a \emph{vida}, 0 \emph{monimento}, \emph{a continuidade}, por isso as nossas ideas relativas a objectos fixos e immutaveis, a verdades necessarias e eternas, são vãs e illusorias. Apresenta um remedio para tornar «\textasciicircum{}

OS possivel a metaphysica \emph{objectiva}, e o remedio é a intuição. e

Não menos inacceitavel é x opinião de \emph{Bergson}. A assercüo que a nossa ---- intelligencia não percebe o \emph{movimento}, \emph{a vide}, a continuidade, é gratuita e = contraria á experiencia. --- Alem dº isso, é falso que as nossas ideas, rela- . "\emph{m} tivas ás essencias e aos primeiros principios, sejam vãs e illusorias. Ellas 7 \emph{u} são tão fixas e immutaveis, como o são os objectos representados. --- Nem pode ser admiltido o remedio, apresentado por esse escriptor. A énfuipzo seria, para elle, uma nova faculdade, substituida a intelligencia, uma especie de \emph{instincto consciente}. Ora repugna que possa existir um \emph{tnstincio} con-

: --- \emph{sciente}, numa faeatdade dotada; ao. mesmo. tempo, de instincto. c.de reflexio. . .. .: EE Uma intuição sem ideas é uma coniradiecüo ; porque é uma visão que não CG \emph{Ts} é visão, --- Notamos que. para \emph{Hergson}, o instincto, perdido pelo homem, in quando na evolução deixou de ser puro animal, só pode recuperar-se e ter- C nar-se \emph{consciente} por um esforço violento e doloroso do espirito. Como se vé, \emph{L} i inventiva de Bergson não é menos genial que a do seu Kant! ;

(1) Não basta demonstrar que um livro é \emph{gemino}, \emph{inteiro} e \emph{verídico}. 6 --- necessario sahel-o \emph{inferpretar}. Para isso serve a \emph{Rhermeneutica}. E

2) llermenenutiea. --- \emph{Hermeneutíca} é a arte que nos ensina a interpretar --- o sentido das palavras de um eseriptor. \emph{Ax}

b) Elementos para a interpretacüo de um livro. --- Para se conhecer o ver- \emph{O}

Es,

CAPITULO TERCERIO

METHODO

\textsc{SUMNARIO}; --- Ponto de partida: Verdades immediatamente evidentes. --- Pro-

cesso: Methodo analytico e synthetico. --- Termo do movimento: Sciencia.

ARTIGO [

\emph{Ponto de partida} --- \emph{Verdades immediatamente evidentes}.

%%>>> t191

\pg{217} \pg{218} \pg{222} \pg{226} \pg{229} \pg{231} \pg{234} \tese{191}{Methodo.}
\emph{Methodo} é o \emph{caminho}, \emph{que a intelligencia}

\emph{deve seguir para a}. \emph{acquisição da} sciencia; ou, mais propriamente,

é a \emph{ordenada disposição das operações inlelleciuaes}. --- A nossa

intelligencia chega ao perfeito conhecimento da \emph{verdade}, ou a

acquisição da sciencia, por varias e e diversas operações, que de-

dadeiro sentido. das. palavras de umeseriptor, deve attender-so aos seguintes

elementos : \emph{grammatical}, \emph{logico}, \emph{systematico} e \emph{historico}.

\emph{a}) O elemento grammatical é lormado pelas palavras e phrases, que o

escriptor empregou para manifestar o seu pensamento.

.5) O elemento \emph{logico} consiste nas relações existentes entre as diversas

ideas ou pensamentos, que do livro constarem.

\emph{c}) O elemento \emph{systematico} é o laço intimo, que liga e subordina todas

as partes a um centro commum de unidade. --- Muitos escriptores reduzem

- o elemento \emph{systematico} ao \emph{logico}, e ensinam que o elemento \emph{Zogico} consiste

na harmonia de uma parte com outra e de todas com um \textbf{centro} commum. A : \emph{d}) O elemento \emph{historico} refere-se ás circumstanocias biographicas, que B possam revelar o pensamento do proprio auctor. --- t) Regras da hermeneutica, --- A hermeneutica, dà as regras para a exacla : \textasciicircum{} interpretação 'd' esses quatro elementos. . \emph{TU a}) Para. a interpretação do elemento grammatical, é necessario --- 1.º)

conbecer perfeitamente não só a lingua, em que o livro está escriplo, mas

tambem o sentido, que as palavras e as phrases tinham ne epocha, em que . \& viveu o auctor; --- 9*.) tomar as palavras no sentido proprio c natural,.

salvo, sc assim tomadas, exprimirem um absurdo; --- 3.º) consultar os

codices manuscriptos, ou as edições mais fieis e emendadas pelo auctor. : 5) Para a interpretacüo do elemoeuto \emph{logico}, deve attender-se aos lozares d \emph{parallelos}, aos principios, «que o escriptor estabelece, ás consequencias, que \emph{do} Lira, ás definições, que formula, c ás semelhanças, que emprega, etc. i \emph{c}) Para a interpretação do elemento \emph{systematico}, deve considerar-se todo \emph{ji} o conlexlo do obra e o escopo, que o auctor se propoz. .

\emph{d}) Para a interpretação do elemento \emph{historico}, deve attender-se ao tempo, ; em que o auctor viveu, is cireümslaneius, que o levaram a escrever, a

doutrina, que professava, etc.; porque tudo isso concorre para manifestar \emph{E} € esclarecer o seu pensamento.

vem ser convenientemente dispostas para se alcançar esse fim. Tal conveniente ou ordenada disposição chama-se \emph{methodo} \footnote{A utilidade e 2 necessidade do \emph{methodo} são manifestas. A sciencia não é um aggregado qualquer de noções ; mas é uma serie ordenada de principios e conclusões. relativas \emph{a} um objecto determinado. Ora é o \emph{methodo} que nos ensina o modo de coordenar os nossos conhecimentos.}.

%%>>> t192

\pg{235} \tese{192}{Elementos do methodo.}
No caminho intellectual, como no material;-devemos distinguir: tres-coisas:- o-ponto--de --- \emph{parlida}, o \emph{processo} e o \emph{termo}..--- O \emph{ponto}, de que parte a intelligencia, é constituido pelos \emph{principios} do raciocinio; O \emph{processo} é o modo, por que a intelligencia vai successivamente deduzindo das premissas as conclusões; o \emph{fermo}, a que se chega, é a \emph{sciencia}. --- O \emph{methodo}, porém, consiste propriamente no \emph{processo} ou movimento da intelligencia. \textbar{}

Neste artigo tratamos \emph{do ponto de partida}; no seguinte, do \emph{processo}; no terceiro, do \emph{termo}.

193. O ponto da partida das demonstrações 'philosophicas

não pode ser a duvida universal. --- Às conclusões do raciocinio estão contidas nos principios, e d'estes recebem a sua. luz \textbar{} e a sua forga. Ora, se a duvida \emph{universal} fosse o principio dos raciocinios, ou o \emph{ponto de partida} das dewonstrágoes plillosóphicas, nenhuma conclusão seria certa e evidente, porque a duvida é a negação de toda a certeza e evidencia, e assim teriamos . . o \emph{sceplicismo universal}. Logo o \emph{ponto de partida} das demonstrações philosophicas não pode ser a \emph{duvida universal} (\textasciicircum{}).\footnote{Estabelecamos. com exactidão, o estado d' esta questão importantissima. O inicio do methodo. o ponto de partida para a descoberta da verdade, para a acquisição da sciencia, poderá ser a \emph{duvida}: \emph{universal}, \emph{seria} c real? Poderemos e deveremos começar, duvidando, \emph{seria} e \emph{realmente}, de tudo, mesmo das coisas mais evidentes e da propria evidencia ? Não; não é possiyel duvidar, real e \emph{seriamente}, das coisas mais evidentes, nem da propria evidencia; não só porque essa disposição de espirito não depende da nossa vontade, mas tambem porque não haveria nenhum meio efficaz para procurar e descobrir a verdade. O ponto de partida seria um ponto de immobilidade, seria o sepulero de todo o processo scientifico.\par A esta questão antiracional pode substituir-se esta outra: Começando a procurar a verdade.-deveremos começar como \emph{se duvidassemos} de \emph{tudo}, isto é, deveremos partir de uma \emph{ficticia duvida} de \emph{tudo}, mesmo da evidencia,--- ou deveremos admittir, como certos, alguns principios \emph{primeiros}? A esta. segunda pergunta damos uma duplice e diversa resposta, conforme se trata . de \emph{descobrir}, ou de.demonstrar a verdade.\par \emph{a}) Se se \emph{trata} de \emph{descobrir} a verdade, essa duvida \emph{ficticia e} universal €\pgn{236}\emph{possivel}. \emph{O} homem, que adquiriu muitos e diversos conhecimentos, pode, \emph{A} levado por uma legitima curiosidade, reflectir sobre si mesmo e sobre as \emph{l Hi} suas certezas espontaneas, ou adquiridas, para as justificar perante OS seus olhos. Não é que duvide real e \emph{seriamente} d' essas certezas ; mas' quer exa- 5 minaleas scientificamente. Este exame pode ser \emph{universal}, sem exceptuar \emph{a P propria evidencia}. --- Ora isso não é impossivel. Sendo o estado primitivo\par \emph{HE} do nosso espirito uma £\emph{abwua rasa}, na qual nada está escripto mas tudo se \emph{BE} pode escrever e se vai escrevendo gradualmente, é claro que, pela reflexão, di podemos regressar a esse estado inicial, seguir a gradual evolução das \emph{noso} ' sas cerlezas, espontaneas ou adquiridas. e determinar a razão e o modo por \emph{HER que} cada uma d' ellas se gravou tão profundamente na possa intelligencia, \emph{ix} que nenhuma duvida real tem podido abalar, O resultado d' esse examoe psy- .chol Rico seria --- não. só a eliminação de muitas opiniões civelmente erroneas, productos de prejuizos, de paixões, de leviandades, --- mas tambem a \emph{convicção} de que-muitas: cerlezas; --- .como-a- existencia do mundo e de nós mesmos, a aptidão natural das nossas faculdades perceptivas, etc., foram adquiridas, não pela cega propensão da nossa mentalidade, mas pela evidencia \emph{immediata} das coisas. D'este modo, a \emph{evidencia objectiva} apparecer-nos-hia como o \emph{motivo} da todas as nossas certezas. --- Nem ha aqui uma contradicção, Contradieção haveria. se admittissemos a \emph{evidencia}, depois de sc ter negado ou \emph{duvidado} d'ella de um modo \emph{positivo}. Mas não fizemos isso. Ad- ! mittimos \emph{elfectivamente a} evidencia. antes do exame reflexo ; só, no decurso d'êsse exame, fizemos abstracção de todas as nossas certezas \emph{evidentes}, ficamos \emph{indifferentes} com relação a ellas, para adquirirmos a conviecio de que essas certezas tinham effectivamente por motivo a \emph{eu:dencia objectiva}. Uma \emph{b} coisa é a duvida \emph{positiva}, proveniente de uma \emph{razão} positiva de duvidar; outra coisa é.a duvida negativa. proveniente de uma \emph{ignorancia ficticia} de todas as certezas espontaneas, mesmo, da \emph{evidencia}. É esta duvida \emph{ficticia} c \emph{universal} o \emph{ponto} de partida na \emph{descoberta} dà verdade. S. Thomaz diz da \emph{Metaphysica}: «I8ta. scientia sicut habel universalem considerationem de vorilate, ita etiam ad eam pertinet \emph{universalis} dubitatio de veritate; et ideo non \emph{particulariter} sed simul \emph{wniversalem dwbitationem} prosequitur» (In DI \emph{Mefaph}.,lec. Y). Este, exame, pela verificação das nossas certezas, faz com que --- a nossa confiança \emph{espontanea} na aptidão das nossas faculdades, e na evidencia objectiva, se transforme em confiança \emph{reflexa}, --- as vossas certezas naturaes se tornem certezas \emph{seienfificas}, --- e O accordo entre a actividade \emph{espontancça} c á actividade \emph{reflexa} da nossa natureza, posto agora em evidencia, nos convença de que a machina humana está maravilhosamente organisuda.\par 5) Se. porém, se trata, não de \emph{descobrir}, mas de \emph{demonstrar} aos outros \emph{uma verdade} já descoberta, então é evidente que o ponto de partida, se pode ser a duvida \emph{particular}, pelo menos \emph{ficticia}, da verdade demonstravel, não pode ser a \emph{duvidal} universal, nem real nem \emph{ficticia}. Uma discussio com um adversario, \emph{ou} com um discipulo, deve partir --- de verdades, anteriormente\pgn{237}admittidas por ambas as partes, --- de principios. jà demonstrados ou ad- . mittidos e capazes de servirem de premissas para a demonstração, Sem isso, toda a demonstração é absolutamhnte impossivel. Assim como nada poderia mover-se, se tudo precisasse de ser movido; assim tambem nada poderia provar-se, se tudo precisasse de ser provado. O movimento do pensamento, como. todo e qualquar movimento, precisa de um motor nao-movido, ou immovel. Concluimos que o ponto de partida, na \emph{demonstração}, não pode Ser a \emph{duvida universal}, nem real nem \emph{ficticia}.\par Dijssemos que o ponto de partida pode ser a duvida particular, pelo\par menos \emph{ficticia}, da verdade demonstravel. Qual é a indole d' \textbf{essa} verdade, de que podemos legitimamente duvidar? Por quanto, ba verdades \emph{primarias} e \emph{inrüemonstraveis}; ha verdades \emph{certas} e \emph{demonstraveis}; e ha verdades \emph{duvzdoses} e \emph{cognosciveis}. --- Quanto ás verdades primarias e indemonstraveis, O homem não pode ter duvida, nem \emph{ficticia}; porque repugna a \emph{ficção} de uma coisa intrinsecamente impossivel, e tal é a duvida das verdades \emph{primarias}, . como, p.e., é a nossa existencia. --- Quanto ás verdades certas é \emph{demonstra}- \emph{veis}, o homem, que. já as conhece.como cortas,.--- nio.pode.ter.uma. duvida.\emph{real}, que seria uma mentira e uma affronta feita á propria verdade, --- mas 'pode ter uma duvida \emph{ficticio}, para descobrir, por meto do raciocinio; a. \textbf{razão} explicita dos seus conhecimentos e a força dos motivos da sua certeza. --- Quanto ás verdades \emph{duvidosas} e \emph{cognosciveis}, \emph{o} homem pode ter uma duvida \emph{real}, -que pode ser tirada por fneio da investigação e do raciocinio, --- Esta duvida \emph{particular} ditffere da duvida \emph{universal} dos scepticos, não só em razão da \emph{extensão} (porque uma é particular e outra é \emph{universal}), mas tambem em razão do \emph{fim} (porque a \emph{universal} tem por fim afastar os homens da investigação da verdade, ao passo que a \emph{particular} os estimula para essa investigação) e em razão do \emph{processo} (porque a \emph{universal} é O \emph{fermo} da inquisição, e a \emph{particular} é apenas o \emph{ínicio}).\par * 1\par Agora algumas palywras ácerca do systema de \emph{Descartes}, relativo \& presente questão. Este escriptor, no livro « \emph{Discours de la méthode} », estabelece i \emph{duvida universal} como o ponto de partida nas invesligações e demonstrações scientifieas, e como meio seguro para distinguir a verdade da falsidade. Damos o resumo do systema. --- O homem, para levanlar em bases firmes o- edificio. scientifico, deve collocar-se no estado de \emph{duvida universal}. Mas, principiando a duvidar de tudo, logo percebe, pelo testemunho da \emph{consciencia}, que não pode duvidar da existencia do proprio pensamento; porque duvidar é pensar. Admittida a existencia do pensamento, deve admittir a existencia do sujeito pensante. D'ahi a senlenca: Penso, \emph{logo existo}. Esta sentença é o primeiro principio, o fundamento de toda a philosophia, e d'ella deriva o grande criterio: \&\emph{udo o} que é \emph{conhecido clara} e distinctamente, é verdade, Em seguida, analysando os objectos do pensamento, encontra na sua intelligencia a idea do infinito. Nesta idea percebe-se claramente a existencia de Deus. Logo Deus existe. Se Deus existe, \emph{o} homem não pode eunganar.se --- nem na percepção das verdades immediatamente evidentes, ---\pgn{238}\emph{E} subir de um principio para outro, sem nunca chegarmos a um\par \emph{dio} primeiro, cuja verdade seja manifesta pela simples inspecção É dos termos. Esta segunda hypothese não pode admittir-se, por-\par nem na percepção dos phenomenos externos; aliás o erro refundir-se-hia no proprio Deus, Auctor das faculdades humanas,\par Do livro citado não apparece claramente se a duvida \emph{universal} é, para \emph{Descartes}, real, OU apenas \emph{ficticia}. É certo, porém, que esse escriptor não\par é seeptico; porque a duvida \emph{universal} não é \textbar{} para \emph{Descartes}, como é para os scepticos, o fermo das investigações, mas é é apenas « o \emph{inicio}, de que. parte para chegar a certeza. '\par Todavia, o systema não pode, admittir-se pelas. seguintes razões:\par \emph{a}) \emph{A duvida universal} €, para \emph{Descartes}, \emph{real} ou \emph{ficticia}. --- Se é \emph{real}, ella é impossivel como \emph{estado normal} do nosso espirito. Ha verdades tão evidentes, que se impõem ao nosso espirito, sem nos deixarem a possibilidade da duvida, --- Se é \emph{ficticia}, é impossivel como \emph{fundamento} das demonstrações. Duvidando de \emph{tudo}, mesmo dos primeiros principios da razão e da nossa capacidade de alcançar a verdade, não é possivel resolver a duvida por meio do raciocinio, que suppõe, como coisas certas, aquelles principios e aquella capacidade. Diz toxito ven \emph{Ratier}: « \emph{Descartes} metteu=se' no estado de um -bomem, .que,..pelo. receio. de se apartar. do .verdadeiro caminho. certa as proprias pernas e assim não pode dar mais um passo» (\emph{Logique}).\par \emph{b}) \emph{Descartes} começa por admittir como certo um só \emph{facto}: a \emph{existencia do} pensamento. Mus um facto, embora incontestavel, embora verdade real, se não é esclarecido por um principio racional, por uma verdade \emph{ideal}, permanece sempre esleril, infecundo. Se \emph{penso}, que consequencia poderá d'ahi tirar a sciencia? --- Se penso, responde \emph{Descartes}, \emph{existo}: eis a consequencia. Mas não vé elle que essa consequencia não se pode tirar e não se tirou senão pela applicação de um principio \emph{idea}? --- o principio de \emph{causalidade}, de que \emph{elle} duvida e que todavia admitte como certo? Sendo assim, a \emph{duvida} universal de \emph{Descartes} não é \emph{universal}. Mas, se a sua büa fé está salva, a sua logica está muito conipromettida;\par c) A celebre sentença --- \emph{penso}, \emph{logo éxisio} --- é, na opinião de \emph{Descartes}, \emph{O} primeiro principio, o fundamento de toda. 2 philosopbiz, e d'ella deduz \emph{ires} coisas: --- o principio'«f£udo \emph{o que é concebido} clara e \emph{distinctamente}, \emph{é verdade}», --- a. existencia de Deus, percebida claramente va idea de \textsc{Deus}, --- á veracidade das nossas faculdades, fundada na propria veracidade de Bleus, Ora essa sentença --- penso, \emph{logo existo}, --- quer seja um raciocinio (como diz Gassendi), quer seja a simples analyse de um facto interno (como pensa Cousin), não pode ser o \emph{primeiro principio} de toda a demonstração philosophica; não só porque o \emph{primeiro principio} não pode ser um \emph{facto}, como é a existencia do pensamento, mas tambem porque o primeiro principio é, como vimos, o de contradicção. --- Alem d'isso, aquelle principio --- tudo o \emph{que é concebido} de \emph{f} modo claro e distincto, é \emph{verdade} --- \& uma E consequencia de um \emph{facto} c \emph{e} um facto \emph{particular}, como é a. existencia do pensamento de \emph{Descartes}. Ora, se este escriptor duvida de todos os primeiros principios, não tem o direito de \emph{generalisar} um [facto particular, nem q pode deduzir d'elle uma regra \emph{geral}, applicavel a todos os factos. Toa \emph{d}) \emph{Descartes} demonstra --- a veracidade da nossa razão pela existencia\pgn{239}e veracidade de Deus, --- e a existencia e a veracidade de Deus pela veracidade da nossa razão; e assim demonstra o primeiro pelo segundo, e o segundo pelo primeiro: o que é um \emph{cireulo vicioso}. --- Do facto de Deus ser o principio da ordem \emph{ontologica} e de toda a existencia, não se segue que a idea de Deus seja o principio da ordem \emph{logica} e de todos os nossos conhecimentos.\par \textbar{} \emph{e}) Este systema leva ao \emph{idealismo}. \emph{Descartes} somente tem como certo o testemunho da consciencia. Mas a consciencia é uma faculdade interna, cuja acção está limitada a ordem \emph{subjectiva}. Para que, pois, possa o homem passar da ordem \emph{subjectiva} para a \emph{objectiva}, é necessaria uma outra faculdade, distincta da \emph{consciencia}. --- Como lambem, a \emph{consciencia}, se attesta a \emph{existencia} e a operação das outras faculdades, não attesta a \emph{veracidade} d'ellas, e não julga dos objectos, proprios das outras faculdades. --- Além d' isso, é um contrasenso deduzir da existencia da idea do \emph{infinito} a existencia real de Deus. De uma idea só segue uma idea; e a idea do \emph{infinito} pode ser formada pela nossa intelligencia:\par \emph{D}) O systema. da \emph{duvida universal} não só leva ao \emph{idealismo}, mas tam- \textbar{} bem ao \emph{scepticismo} universal. Se \emph{Descartes} duvida da propria evidencia, não é possivel admittir, sem contradiccio, como certa a existencia do pensamento e do sujeito pensante. Se nem essa existencia é \emph{certa}, temos o \emph{scepticismo} universal, Se não temos a contradicção. --- Se esse escriptor duvida da legitimidade da \emph{ração}, dos \emph{sentidos externos}. etc., porque não deveria duvidar tambem da legitimidade da \emph{consciencia}? Se nio duvida da legitimidade da \emph{consciencia}, porque deveria duvidar da legitimidade \emph{da razão}, dos \emph{sentidos externos}, etc. ? O argumento vale para: todas as faculdades. Logo, ou temos a contradicção, OU O \emph{scepticismo} universal.}

%%>>> t194

\pg{236} \tese{194}{O ponto de partida das demonstrações philosophicas deve ser constüfuido por verdades immediatamente evidentes.}
A demonstração consiste em deduzir de principios certos conclusões certas. Se- esses- principios não sao immediatamente evi-

7

”

dentes, devem ser demonstrados por outros, e estes por outros, e assim por deante. D'bhai resulta que --- ou devemos \textbf{parar} num principio, que seja immediatamente evidente, --- ou devemos

ToGlca --- SEcção SEGUNDA -

855 que destroe toda a certeza. Logo o \emph{ponto de partida} das demon-

strações philosophicas deve ser constituido por verdades \emph{immediatamente evidentes} \footnote{O inicio das demonstrações deve ser, pois, constituido por verdades \emph{imnmediatamente} evidentes. --- D'estas verdades --- umas são \emph{fundamentaes} e \emph{féraes}, emquanto devem ser presuppostas necessariamente a toda a demonstração ou discussão, seja qual for o seu objecto, --- outras são \emph{fundamentaes} e \emph{especiaes}, emquanto constituem o principio de uma \emph{sciencia especial}, ou da. \emph{parte especial} de uma sciencia. --- As verdades, que são fundamentaes e \emph{especiaes}, reduzem-se, de algum modo, ás verdades, que são fundamentaes e \emph{geraes}.\par As verdades, gue. se vonsideram como o fundamento \emph{geral} de toda a\par ir-se a \emph{tres}, e são, --- 19) o \emph{facto} da nossa \emph{existencia}, emquanto Somos seres racionaes e capazes de attingir a verdade, isto é, de attingir as coisas como são em si mesmas; --- 2º.) \emph{a existencia} dos \emph{objectos}, emquanto têm a capacidade de ser attingidos pelas nossas facul dados; 3.9) a \emph{evidencia objectiva}, que é o signal que indica quando o objecto é attingido por nós e que é o motivo da nossa certeza. --- É claro que, admitindo a \emph{evidencia objectiva} como signal de verdadeira percepção e motivo da nossa certeza, devemos admittir, implicitamente, todas as verdades evidentes e indemonstraveis, legitimadas pela sua propria evidencia, quer pertençam a ordem \emph{racional}, que pertençam a ordem \emph{experimental}. --- Se * não admitlirmos estes tres postulados, se duvidarmos da capacidade per-}.

195. O ponto de partida das demonstrações philosophicas, quanto á parte racional, \textbf{deve} ser constitaido: por verdades immediatamente evidentes e necessarias. --- O \emph{ponto} de \emph{partida} das demonstrações philosophicas, quanto á parte \emph{racional}, deve ser constituido por aquellas verdades, em que a razão resolve, por ultimo, as mesmas demonstrações e baséa toda a certeza das conclusões deduzidas. Ora a razão resolve, por ultimo, as suas demonstrações é baséa toda a certeza das conclusões deduzidas nas verdades \emph{immediatamenlie evidentes} e \emph{necessarias}. Logo o \emph{ponto de partida} das demonstrações philosophicas, quanto á parle \emph{racional}, deve ser constituido Dor verdades \emph{immediata}- . \emph{mente evidentes} e \emph{necessarias} \footnote{Diz S. Thomas :' « Certiludo scientiae tota oritur ex certitudine principiorum. Tune enim illationes per certitudinem sciuntar, quando resolvuntur in principia» (\emph{De} Ver., q.:1, 3.1). --- Na Philosophia devemos distinguir duas partes: uma \emph{racional}, outra, experimental, conforme trata da \emph{essewmcia} ou da \emph{existencia} das coisas. --- Na parte \emph{racional} da Philosopbia, a força .das conclusões deriva dà força das premissas. Essas premissas, na mesma parte da Philosophia, como em todas as sciencias \emph{racionaes}, são verdades secessarias, porque denotam uma relação necessaria ou essencial entre o predicado e o sujeito, e devem ser, em \emph{ultima analyse}, \emph{imediatamente evidentes}: porque, aliás, não é possivel a certeza scientifica, --- Dizemos --- em \emph{wine} :}. ceptiva da intelligencia, ou da capacidade que o objecto tem de se manifestar, ou da efficacia da evidencia objectiva, a demonstração não é possivel.

Alguns. philosopbos admittem, como postulado, n\&o a \emph{evidencia objectiva}, mas o \emph{principio} de contradicção. --- Não mos parece "acceitavel esta substituição. -Por- quanto, .0 \emph{principio} .de contradicção, reconhecendo-se verdadeiro pela sua \emph{evidencia}, não é \emph{anterior}, más \emph{posterior á evidencia}... Alem d'isso, o \emph{postulado} deve conter, de algum modo, a explicação de todas as verdades. Ora o \emph{principio de contradicção} --- não explica as verdades \emph{comtingentes}, (come a. existencia nossa, a do mundo, ete.), que são admittidas pela sua \emph{evidencia}, e não por esse principio, --- c, se mostra a \emph{possibilidade} das coisas, não pode mostrar a \emph{realidade} d' ellas.

Se um só dos tres postulados indicados não é admiliido pelo adversario, nem pode encetar-se a discussão; pois as armas devem ser eguaes. Todavia, se, nesse caso, não é possivel o methodo da demonstração, é possivel o methodo' da descoberta, OU \emph{verificação}. Deve levar-se o adversario a reflectir sobre si mesmo, a interrogar a sua consciencia (psychologica), fazendo-lhe tocar com as mãos como sobre essa laboa rasa, que é, no principio, o nosso espirito, se tem gravado --- a \emph{existencia do} eu pelo exercicio do pensamento, --- a \emph{existencia} dos \emph{objectos externos}, que uctuaram em nós e nos transmittiram a propria imagem, --- e o \emph{motivo} da nossa adhesão ás verdades conhecidas, constituido, não por um cego impulso da natureza, mis pela clara intuição dos. objectos. Feito este exame, o adversario, se estiver na hia fé, não pode deixar: de idmiltir or tres \emph{postulados}..

196. O ponto de partida das demonstrações philosophicas, quanto a parte.experimental, deve ser constituido por verdades immediatamente evidentes e contingentes. --- O \emph{ponto de partida} das demonstrações philosophicas, quanto á parte \emph{experimental}, deve ser constituido por aquellas verdades, que não .precisem de demonstração e que, ao mesmo tempo, se refiram directamente a \emph{existencia real} das coisas; porque só estas verdades podem ser o fundamento solido de raciocinios, pelos quaes chegamos ao conhecimento da existencia real de outras coisas. Ora as verdades, que não precisam de demonstração e que, ao mesmo tempo, se referem directamente a \emph{existencia real} das coisas, são as verdades \emph{immediatamente evidentes} e \emph{contingentes}. Logo o \emph{ponto de partida} das demonstrações philosophicas, quanto á \emph{parie} experimental, deve ser constituido por verdades immediatamente \emph{evidentes}. e \emph{contingentes} (). \textbar{} analyse. Por quanto, não é necessario que toda a sciencia \emph{racional} parts B de verdades \emph{immediatamente} evidentes; pode partir de verdades; recebidas de uma sciencia racional \emph{superior}. na qual ellas possuem essa evidencia immediata. --- As verdades \emph{immediatamente evidentes} constitaem o \emph{ponto de partida} das demonstrações na: mais nobre e elevada entre todas as Sciencias racionaes, que é a \emph{Metaphysica}. --- Repetimos que admittir uma verdade sem demonstração não é admittil-a sem motivo. O motivo existe, e é a evidencia immediata da propria verdade.

Objectam : A intelligencia, na resolução das suas demonstrações, pára, por ultimo, na veracidade das nossas faculdades. Logo é d' esta veracidade que deve partir a propria demonstração.

Resposta. A intelligencia, quando resolve a demonstração nos principios, de que partiu, serve-se da reflexão. Ora a reflexão é \emph{ontologica} e \emph{psychologica}. Esta considera \emph{o} elemento \emph{subjectivo} do acto; aquella, o elemento \emph{objectivo}. Se o acto se considera e se resolve no elemento \emph{subjectivo}, então o termo ultimo é a veracidade das nossas faculdades (e isto mostra que não se. pode racionar, sem que anteriormente se admitta essa veracidade). Mas, se o acto se considera e se resolve no elemento \emph{objectivo}, o termo ultimo, que é o motivo ultimo e a fonte de certeza, é constituido, não pela veracidade das f\&culdades, mas pelas verdades \emph{vmmediatamente evidentes} e necessarias, que.se manifestam pela sua propria luz e nos levam ao descobrimento de outras verdades, com que estão relacionadas.\footnote{Na parte \emph{experimental} da Philosophia, e, em geral, nas sciencias experimentaes, a força das conclusões deriva, como sempre, da força das premissas. Estas premissas referem-se á \emph{existencia} real das coisas, e por isso são verdades \emph{contingentes}, porque \emph{contingente} é essa existencia, Assim quando digo: o mundo \emph{existe}, enuncio uma verdade \emph{contingente}, porque \emph{a} existencia para o mundo não é absolutamente necessaria, mas podia ser e podia deixar de ser. --- As premissas dus demonstrações das sciencias racionaes não podem ser premissas das demonstrações das sciencias expe-\par Vel, I Wo\pgn{242}\emph{rimeníaes}; porque de uma premissa de ordem ideal, como é a premissa das sciencjas racionaes,não podemos rasoavelmente passar para uma conclusão de ordem \emph{real}, como é a conclusão das sciencias sxperimentaes. A conclusão deve-ser-sempre-proporcionada-ás-premissas,- em- que-está-contida.\par As verdades \emph{immediatamente evidentes} e \emph{contingentes} denotam factos \emph{externos} e factos \emph{internos}. Se denotam factos \emph{externos}, foram adquiridas pela intelligencia, baseada no testemunho dos \emph{sentidos externos}. Se denotam factos \emph{internos}, foram adquiridas pela intelligencia, baseada no testemunho da \emph{consciencia}. --- E, como o acto dos \emph{sentidos externos}, que é \emph{directo}, precede : o acto da consciencia, que é \emph{reflexo}, segue-se que o processo philosophico, devendo repetir o processo natural, deve começar, quanto á parte \emph{experimental}, pela consideração do \emph{mundo}, percebido pelos sentidos \emph{externos}.\par Quando dizemos que as verdades \emph{immediatamente evidentes} e \emph{contingentes} são o fundamento da Philosophia quanto á parle experimantal, e. em geral, das sciencias \emph{experimentaes}, referimo-nos ao fondamento \emph{proximo} e \emph{proprio} da demonstração, ou da sciencia; porque o fundamento \emph{ultimo} e ' commum a todas as sciencias ó constituido pelas verdades necessarias e dotadas de evidencia \emph{immediata}.}

%%>>> art-10

\section*{Artigo II.\\ \textbar{}}
\addcontentsline{toc}{section}{Artigo II. --- \textbar{}}

\pg{237} \pg{238} \pg{239} \pg{240} \pg{241} \pg{242} \emph{Processo} --- \emph{Methodo analytico} e \emph{synihetico},

%%>>> t197

\tese{197}{Methodo analytico e synthetico.}
O \emph{methodo} divide-se em \emph{analytico} e \emph{synthelico}. --- \emph{Analytico} é o que, baseado na \emph{experiencia}, vai dos objectos que têm mais para os que têm menos \emph{comprehensão}. Por este methodo, a razão \emph{sobe} dos \emph{effeitos} para as causas, dos \emph{factos} para as'leis, do \emph{todo} para a \emph{parie}, etc. --- \emph{Synihetico} é o que, fundado na legitimidade da razão, vai dos objectos que têm menos para os que têm mais \emph{comprehensão}. Por este methodo, a razão \emph{desce} das causas para os \emph{effeitos}, das \emph{leis} para. os \emph{factos}, da \emph{parie} para o \emph{todo}, etc. \footnote{São, pois, funcções do methodo analytico as seguintes operações : --- a inducção, que de factos particulares chega a conclusões geraes (um facto \emph{particular} tem maior \emph{comprehensão} que uma conclusão \emph{geral}, porque aquelle exprime e contêm esta); --- a abstracção, que de uma idéa particular extrahe uma idéa geral; --- toda a \emph{demonstração}, que sobe do eifeito para a causa, dos phenomenos para os principios, do composto para o simples, do concreto para o abstracto, do particular para o geral, do todo para a parte. --- Parece contradictorio o dizer que o methodo \emph{amalytico} sobe do \emph{particular} para o universal, e do \emph{todo} para a \emph{parte}. Mas não o é. O \emph{particular} está para o \emph{universal}, como o \emph{todo} para a \emph{parte}; pois o \emph{particular} contém em si todas as notas do \emph{universal}, \emph{e} mais algumas outras notas proprias.\par São, pois, funcções do methodo \emph{synihetico} as seguintes operações; --- o \emph{syllogismo}, que applica os principios geraes a objectos particulares; --- a \emph{definição}, pela qual consideramos o \emph{todo} formado de varias partes; --- \& \emph{divisão} e a \emph{classificação}, quando vão do genero para a especie, e d' esta para os individuos ; --- toda a demonstração, que \emph{desce} da causa para \emph{o}\pgn{243}efeito, dos principios para as consequencias, do simples para o composto, do abstracto para o concreto, do universal para o particular, da parte para o todo. --- Note-se que o \emph{effeito}, como tal, tem maior \emph{comprehensão} que a causa; pois a procedencia da causa é apenas uma das notas do effeito.\par O methodo \emph{analytico} chama-se tambem --- de invenção, porque geralmente serve para descobrir a verdade, --- de resolução, porque resolve um todo nos seus mais simples elementos, --- \emph{empirico}, porque se baséa principalmente na experiencia, --- \emph{inductivo}, porque faz muito uso da inducção, para subir dos phenomenos particulares para as leis universaes.\par O methodo \emph{synihetico} denomina-se tambem --- de doutrina, porque é mais opportuno para ensinar a verdade descoberta, --- de composição, porque lira dos principios universaes as conclusões particulares, --- deductivo, porque emprega especialmente a argumentação \emph{deductiva}, que € O \emph{syllogismo}.\par Erra, pois, \emph{Condillac}, quando reduz a \emph{analyse} a uma \emph{divisão}, e a \emph{synthese} a uma composição. Um e outro methodo, manifestando a relação dos effeitos com as suas causas, e d'estas com aquelles, explicam a coisa pelas seus principios, e por isso são processos \emph{scienfificos}, e as sciencias, que investigam as causas dos entes, devem empregar um ou outro.\par É nacessario não confundir o metbodo \emph{analytico} com o juizo \emph{analytico}. nem o methodo \emph{synéketico} com o juizo \emph{symthetico}. --- O methodo \emph{analytico} indica o processo da razão, que, baseada na experiencia, sobe do particular para o universal. O juizo \emph{analytico} exprime uma verdade necessaria, em que a conveniencia entre o sujeito e o predicado se conhece pela simples \emph{analyse} ou resolução d' estes dois termos. --- O methodo \emph{synthetico} é o caminho da inteligencia, que do universal desce para o particular. O juizo \emph{synthetico} é uma verdade \emph{contingente}, em que a relação entre o sujeito e o predicado se conhece, não pela \emph{analyse} dos dois termos, mas pela experiencia.}.

: 198. Differenca entre os dois methodos. --- O methodo \emph{analytico} segue a ordem \emph{logica}, ou do \emph{conhecimento}, porque o \emph{effeito} conhece-se antes da causa, os \emph{factos} antes da \emph{lei}, --- O methodo \emph{synthetico} segue a ordem \emph{real}, ou da \emph{natureza}, em que a causa precede o \emph{effeito}, as leis são anteriores aos \emph{factos} \footnote{Estas duas palavras: « \emph{analyse}» e «\emph{synthese}» têm uma significação - \textasciicircum{} \textasciicircum{} muito clara e determinada nas sciencias naturaes. Assim, todos entendem que a \emph{chimica}, por meio da \emph{analyse}, \emph{resolve} a agua DOS seus elementos, para a tornar a \emph{compor} por meio da \emph{synthese}.\par Mas não acontece o mesmo na Philosophia. Uns consideram como \emph{analyse} O processo que para outros é \emph{synthese}. Assim, dividindo O \emph{genero} nas suas \emph{especies} e concluindo d'aquelle para estas, alguns julgam que é uma «nalyse, quando, na realidade, é uma \emph{synthese}.\par A origem d'esta divergencia nasce da confusão, que fazem entre a \emph{extensão} e a \emph{comprehensão} das ideas, entre o \emph{todo logico} e o \emph{todo} real. --- Consideremos estas duas ideas --- \emph{animal} e \emph{homem}. A idea de animal é \emph{mais extensa} do que a de \emph{homem}, mas é menos \emph{comprehensiva} do que esta. Logo, na ordem da \emph{extensio}, o \emph{animal} a O \emph{todo}, \emph{o homem} é a \emph{parte}, \emph{e} quem conclue do \emph{animal} para o homem realisa uma \emph{analyse}; --- mas, na ordem d\& \emph{comprehensao}, o \emph{animal} é a \emph{parte}, \emph{e O homem} € o \emph{todo}, e quem conclue do \emph{animal} para o \emph{homem}, realisa uma \emph{synéhese}. Por isso, uma \emph{analyse} na ordem da\pgn{244}\emph{extensão} é uma \emph{synthese} na ordem da \emph{comprehensão}, e vice-versa; pois o que é \emph{composto} relativamente á extensão é \emph{simples} relativamente a comprehensão, e vice-versa. --- Assim se explicam as divergencias que existem entre os escriptores a respeito d'estes dois methodos --- \emph{analytico} e \emph{synthetico}.\par Foi para evitar esses equivocos que, no texto, apresentamos as definições dos dois methodos muito explicitas, dizendo que o methodo \emph{analytico} vai dos objectos que têm mais pata os que têm menos comprehensao, e que o \emph{synthetico} vai dos objectos que têm menos para os que têm mais \emph{comprehensao}. Por isso estas palavras --- \emph{particular}, \emph{geral}, \emph{todo}, parte, --- consideram-se na sua \emph{comprehensão}, e não nà sua extensão; assim. quem do conceito de \emph{animal} desce para o conceito de komem, realisa um processo \emph{synthetico}.}.

%%>>> t199

\pg{243} \pg{244} \tese{199}{Regras communs a ambos os methodos.}
Damos as. regras mais geraes, communs a ambos os methodos.

l1. Devemos principiar pelo facil e claro, e subir successivamente para o difficil e transcendente. A razão é manifesta.

\emph{IT}. O -processo deve ser \emph{ordenado}, de modo que entre as muitas e diversas conclusões não exista lacuna, mas haja um nexo intimo e gradual.

IIl. Não póde pretender-se a mesma especie de certeza em todas e cada uma das sciencias; mas devemos procurar aquella certeza, que corresponde á natureza da materia tratada. --- Erram, pois, os que exigem para todas as sciencias a certeza \emph{malthematica}, que é a \emph{metaphysica} \footnote{(Cf. S. Th., in II \emph{Mefaph}., lect, 5). --- Estando \emph{certos} de uma verdade, devemos pôr em pratica a seguinte regra de \emph{Bosswei}: « Não deve abandonar-se uma verdade demonstrada, seja qual for a difficuldade de a conciliar com outra verdade ».}.

%%>>> t200

\tese{200}{Regras proprias do methodo analytico.}

I. Devemos dividir e subdividir o \emph{todo} nas suas mais simples partes, observando rigorosamente as leis proprias da \emph{divisão}.

HH; Cada "uma -das-partes do todo-deve ser examinada e avaliada pelos principios da razão, ou da experiencia, conforme o \emph{todo} pertence á ordem \emph{ideal}, ou a \emph{real}. :

III. Finalmente; deve attribuir- -se ao \emph{lodo} só o que se reconheceu convir ás partes \footnote{Como o methodo \emph{analytico} faz muito uso da inducção, devem observur-se cuidadosamente todas as regras, que demos, para que esta especie de: argumentação seja legitima e nos leve ao conhecimento da verdade,}.

%%>>> t201

\tese{201}{Regras proprias do methodo synthetico.}

I. Na disposição das questões devemos principiar por aquellas que são mais \emph{genericas}, e ir suecessivamente descendo para as que:o são menos. \textbar{} :

11. Na. resolução dos problemas, devemos primeiramente estabelecer, como ultimo fundamento dos \emph{syllogismos}, as verdades dotadas de evidencia immediata, e tirar as conclusões particulares, segundo as regras do \emph{syllogismo}.

ITI. Não acabe o tratado sem que esteja exgoltado o assumpto,

expondo ordenadamente a existencia, a natureza, as especies, as propriedades e os effeitos do objecto \footnote{Baseando-se o methodo \emph{synihetico} no \emph{syllogismo}, Lodas as leis, que estabelecemos para a legitimidade d'esta forma de argumentação, hão de ser escrupolosamenie observadas pelos que empregam esse methodo.}.

. Indole do methodo scientifico. --- O methodo \emph{scientifico}, devendo levar ao conhecimento dos phenomenos pelo conhecimento das suas causas, não póde ser exclusivamente \emph{analytico}, , nem exclusivamente \emph{synthetico}. --- Não póde ser exclusivamente \emph{analytico}; porque pela \emph{experiencia}, em que este methodo se funda, conhecemos o phenomeno em si mesmo, mas não na sua causa. --- Não póde ser exclusivamente \emph{synihetico}; porque, sem o auxilio da experiencia, não podemos perceber os factos ou phenomenos, e por isso nos faltará o fundamento das invesligagóes e as nossas ideas não corresponderão a realidade das coisas.: Logo o methodo \emph{scientifico} deve ser \emph{misto}, isto é, deve ser \emph{analstico-synthetico} (\textasciicircum{}).\footnote{A sciencia, por ser o conhecimento de um objecto pelas suas causas, adquire-se no fim da \emph{syn£hese}, e não da \emph{analyse}: porque no fim da \emph{analyse} conhecemos a causa pelo phenomeno, ao passo que no fim da \emph{syníhese} conhecemos o phenomeno pela sua causa.\par X\par Os philosophos não estão de accordo com relação á indole do methodo \emph{scientifico}, \emph{e} propõem diversos methodos, conforme a diversidade do meio, . de que, na opinião d'elles, o homem se serve para adquirir a sciencia. --- Esse meio é \emph{inírinseco} Ou \emph{exírinseco}. --- Se é \emph{extrinseco}, --- ou é a \emph{tradição}. e temos o methodo \emph{tradicionalista}, --- ou é o \emph{conjuncto} dos diversos systemas, e temos o methodo eclectico. --- Se é \emph{inírinseco}, --- ou é a visão de Deus, conhecido por um modo \emph{directo} e \emph{intuitivo}, e temos o methodo ontologico, --- ou consiste exclusivamente nos \emph{sentidos}, e temos o methodo empirico, --- ou é exclusivamente a \emph{ração}, e temos o methodo \emph{idealista}, --- ou É a \emph{razão} auxiliada pelos \emph{sentidos}, e temos o methodo \emph{analytico-synthetico}. --- Logo os methodos são \emph{seis}, a saher: \emph{tradicionalista}, \emph{eclectico}, \emph{oníologico}, \emph{empirico}, \emph{idealista}, e \emph{analytico-synthetico}.\par D'estes methodo faremos uma breve exposição com a critica respectiva.\par 2) Methodo tradicionalista. --- Este metbodo, seguido por \emph{Bwet}, \emph{Bawutain}, \textbar{} \emph{Lamennais}, Ventura, faz consistir o \emph{meio} dos nossos conhecimentos na \emph{tradição} oU Revelação \emph{sobrenatural}; pois os seus sequazes affirmam que a nossa razão é tão fraca que não pode levar-nos a acquisição de nenhuma verdade.}

%%>>> t203

\pg{245} \pg{246} \tese{203}{Indole do methodo das diversas sciencias.}
Todas as sciencias, não Só os \emph{racionaes}, que lem por objecto as verdades \emph{necessarias}, mas lambem as \emph{experimeniaes}, que tém por

Este methodo, embora se neceite na Theologia revelada, que tira os \emph{principios} da. Revelação \emph{sobrenatural}, não pode acceitar.se nas sciencias bumanas, cujos principios se baséam na luz natural da razão, ou na experiencia sensivel, --- Além d" isso, é falso que a nossa razão não possa alcançar nenhuma verdade scientifica sem o auxilio da Revelação. --- Finalmente, se a razão fosse tão fraca que não podesse chegar ao conhecimento de nenhuma, verdade, não \emph{poderia demonstrar} o \emph{facto} e a \emph{legitimidade} da Revelação \emph{sobre}- “\emph{natural}, e assim nem poderiamos acceitar essa Revelação: d' ahi o \emph{scepticisnto}.

b) Methodo eclectico. --- O methodo \emph{eclectivo}, chamado tambem syncretismo, foi proposto por V. \emph{Cousin}. Partindo do principio que o erro é uma verdade incompleta, \emph{Cousin} ensina que o verdadeiro metbodo consiste em reunir os diversos systemas, que tratam do mesmo assumpto e que se completam. mutuamente, para assim formar um só systema.

É evidente a falsidade do \emph{eclec/ismo}. Primeir. erro não é uma \emph{verdade incompleta}, mas é a negação da “é impossivel reunir as diversas doutrinas pbilosophicas para formar um só systema, porque este unico systema seria um cumulo de contradicções. Assim os materialistas dizem que a nossa, alma é material; os espiritualistas sustentam . que ella é espiritual; o \emph{eclecfismo} reuniria os dois systemas e diria que a nossa alma é espiritual e material. Poderia conceber-se coisa mais absur. da? --- É verdade que alguns \emph{eclecióicos} mais moderados fazem consistir o seu methodo em tirar dos diversos sytemas só as verdades, para com ellas formar um corpo de doutrina. Mas, ainda assim, este methodo é incompleto, exclusivamente \emph{historico}, e \emph{perigoso}. --- É incompleto: porque, para fazer aquella escolha, é necessario um outro criterio, baseado em principios vertos, e \emph{o'ecleciismo} suppõe esses principios, mas não os ensina. --- É exclusivamente \emph{historico}, quando é certo que as sciencias hão de progredir, não pela narração das verdudes adquiridas, mas pela descoberta de outras ainda não conhecidas. --- É \emph{perigoso}, porque, sem um criterio precedente, seguro e firme, é muito facil cahir no erro ou no \emph{scepticismo}.

€) Methodo ontologico. --- O methodo \emph{ontologico}, seguido por \emph{Malebranche}, Gioberti, \emph{Rosmini}, \emph{Tiberghien} (sequaz do \emph{pamnentfheixmo}) consiste em collocar na visto \emph{intuitiva} e \emph{directa} de \textsc{Deus} e do seu conteudo, que é o mundo, \emph{o} meio do conhecimento das coisas creadas.

Tal methodo tambem deve rejeitar-se. Esse modo de perceber as creaturas no Creador, além de ser contrario á natureza da nossa razão, que se eleva ao conhecimento de Deus pelas creaturas, e não vice-versa, é desmentido pela nossa consciencia.

4) Methodo empirico. --- O methodo \emph{empirico}, seguido pelos materialistas e positivistas --- \emph{Comte}, \emph{Littré}, \emph{Stuart Mill}, \emph{Bain}, Taie, Wwndt, \emph{Spencer}, etc., consiste em admittir, como certo, só o que é attestado e verificado pela experiencia \emph{sensivel}, rejeitando os principios da razão, o testemundo da consciencia, o valor do syllogismo, etc.

Este methodo é inacceitavel pelas seguintes razões:

1.0) É falso --- que \emph{a unica} [ente de certeza é a. experiencia \emph{sensivel}, --- \pg{247}  objecto as verdades \emph{contingentes}, devem empregar o methodo \emph{analytico-synthelico}; pois este, como vimos, é o methodo \emph{scientifico}. Quando, pois, se affirma que nas sceiencias \emph{experimeniaes}

que os primeiros principios da razão são falsos, --- que o syllogismo não tem valor demonstrativo, --- que o testemunho da consciencia engana.

2.» O methodo \emph{empirico} limita-se a accumular \emph{factos}, sent se importar de os analysar severamente e de procurar as suas mais elevadas causas; por isso foi chamado «o methodo da \emph{formiga}». --- Ora este processo é \emph{antiscientifico}; pois não só não favorece, mas até contraria o progresso das proprias sciencias \emph{mafwraes}, porque estas sciencias, como observa \emph{Claudio Bernard}, progridem, não tanto pela aceumulação dos factos averiguados, como pela luz \emph{da idea}, que encerra a explicação dos factos conhecidos e a «\emph{anticipacso} dos desconhecidos » (\emph{Phenomenos} da \emph{vide}). Diz \emph{Helmholtz} : «O conhecimento dos factos é nada; a sciencia só começa, quando começam a descobrir-se as causas e as \emph{leis} d' esses factos » (Revue \emph{scient}.). .

3.0) Este methodo rejeita o sylogismo, e basea-se exclusivamente na \emph{inducção}. --- Isto chama-se mutilar a razão e a sciencia. O \emph{syllogismo}, como dissemos, longe de obstar ao progresso da \emph{inducção}, favorece-o. Por quanto, ... a inducção descobre as leis dos phenomenos sensiveis, --- leis, que são \emph{certas} ------ ou \emph{hypotheticas}. Oro o \emph{syllogismo explica} as leis \emph{certas}, \emph{confirma} as \emph{hypotheticas}, e, pelo methodo da \emph{analogia} e da dependencia, descobre leis novas. --- Além d' isso, a \emph{inducção}, como já notámos, nas mãos dos materialistas e dos positivistas, é uma arma inutil e sem força. Na verdade, esses escriptores, não baseando a inducção no principio de \emph{causalidade}, reduzem-na a um sophisma, qual é a passagem do particular para o geral, sem razão sufficiente.

Apraz-nos aqui citar o juizo, que do methodo \emph{positivista} deu \emph{M}. \emph{Pasteur}. \emph{No} discurso de recepção na Academia Franceza, em 97 de Abril de 1882, \emph{Pasteur} censurou o \emph{positivismo} por este ter confundido a observação vulgar com a \emph{experimentação}, e por ter feito tão pouco uso d' esta ultima. Eis as suas palavras. « O experimentador, que interroga e explora a natureza, encontra-se incessantemente deante de factos, que não são ainda manifestos e não existem, na maior parte, senão na potencia das leis naturaes. O desconhecido no possivel e não no realisado :'eis o seu dominio, e para o explorar recorre a esse maravilhoso methodo experimental, do qual se póde dizer com verdade que, se não \emph{basta para tudo}, engana raras vezes e sómente os que d'elle se servem mal. Este methodo elimina certos factos, provoca outros, interroga a natureza, obriga-a a responder, e só se detém quando o espirito se dá por plenamente satisfeito. O attractivo dos nossos estudos, o encanto da sciencia, se assim nos podemos exprimir, consiste em que, em toda a parte e sempre, podemos apresentar a justificação dos nossos principios e a prova das nossas descobertas. O erro de Augusto Comte e de M. Litiré consiste em confundir este methodo com o methodo restricto da observação. \emph{Alkeios} ambos d \emph{experimentação}, dão á palavra «experiencia » o sentido, que se lhe dá na conversação vulgar, inteiramente diverso do \$entido, cm que se toma na linguagem. scientifica. No primeiro caso, a experiencia é \emph{a} observação das coisas, é a \emph{inducção} que mais ou menos legitimamente conclue do passado para o possivel. O \emph{verdadeiro methodo experimental vai até á prova sem} replica. Às condições e O resultado quoti- \pg{248}  \$e deve seguir o methodo \emph{enalylico}, e nas sciencias racionaes o methodo \emph{synihetico}, deve entender-se que naquellas sciencias \emph{predomina a analyse} com a \emph{inducção}, e nestas predomina a \emph{synthese} com o \emph{syllogismo} (\textasciicircum{}). diano do trabalho do homem de, sciencia dispõem o seu espirito a não altribuir uma idea de \emph{progresso} senão x uma idea de invenção. Para julgar, pois, do valor do positivismo, o meu pensâmento foi procurar melle a invenção. \emph{Todavia não a encontrei}... M. Littré e Augusto Comte acreditavam e fizeram acreditar aos \emph{espiritos superficiaes} que o seu systema se baseava hos mesmos principios de methodo scientifico, que teve por fundadores \emph{Arcbimedes}, Galileu, Pascal, Newlon. Lavoisier. \emph{Foi d'ahi que veio} a \emph{illusão dos espiritos}. A que erros conduz esta pretendida identidade dos dois methodos! \emph{Arago} dissera de \emph{Comte}: Elle não possue titulos mathematicos, nem grandes nem. pequenos», U

e) Methodo idealista. --- O methodo \emph{idealista} segue um caminho contrario ao \emph{empirico}, e, desprezando a \emph{anolyse} e a. inducção, emprega. exclusivamente o methodo \emph{synéhetico}, que procede só pela deducção. \textbar{}

Este-systema;que-foi-chamado-o"systema: \emph{da} “\emph{aranha} (pois este animal” tira a teia do interior do seu organismo), é falso. A realidade \emph{objectiva} das coisab não pode adivinhar-se com a imaginação, nem fingir-se \emph{a-priori}, mas \emph{deve} ser conhecida pela experiencia.

f) Methodo analytico-synthetico. --- O methodo \emph{analytico-synthetico}, seguido pelos escholasticos, é o termo medio entre o empirismo e o \emph{idealismo}, empregando a \emph{analyse} e a \emph{synthese} para descobrir a verdade. Pela analyse, examinamos os factos ou phenomenos, e, por meio da inmducção, subimos a causa d' esses \emph{factos}. Pela \emph{synthese}, \emph{d da descoberta}, e, por meio do \emph{syllogismo}, applicamol-a aos varios effeitos, que ella produz ou é capaz de produzir. \emph{Newton} diz: «O methodo \emph{analytico} consiste em fazer experiencias, examinar os phenomenos e d' ahi. pela inducção, tirar conclusões geraes. O '\emph{synthetico} consiste em estabelecer, como principios, as causas conhecidas e averiguadas, para assim explicar e confirmar os phenomenos que d'ellas derivam » (\emph{Optica}, lib. II. --- O methodo \emph{analytico-syntethico} diz-se tambem \emph{psychologico-racional}; pois auxilia-se da \emph{experiencia} e da \emph{ração}.

Este systema é:o verdadeiro, porque corresponde á natureza do homem e a da \emph{sciencia}. --- Corresponde á natureza do komem. Na verdade, \emph{o} homem é dotado de faculdades \emph{sensitivas} e intellectuaes. Logo, no conhecimento da verdade, \emph{o} homem deve auxiliar-se de umas e de outras, sendo \emph{o} conhecimento \emph{sensitivo} anterior ao \emph{intellectual}. Ora o methodo \emph{analytico-synthetico} segue esta ordem. --- Corresponde á natureza da \emph{sciencia}. A sciencia é uma Serie ordenada de \emph{principios} e conclusões. Os \emph{principios} conhecem-se pela \emph{analyse} ou pela \emph{inducção}, e as conclusões deduzem-se dos \emph{principios} pela \emph{synthese} ou pelo \emph{sy/logismo}. Logo o metbodo \emph{scientifico} deve ser \emph{analytico-synthetico}.\footnote{Os methodos \emph{particulares} das diversas sciencias (exceptuada a \emph{Philosophia}) podem reduzir-se a tres grupos, conforme predimina em cada um d'elles o emprego de um d' estes tres \emph{meios} para a acquisig\&o da verdade, que são \emph{O raciocinio}, \& \emph{experiencia} ivel e a \emph{toridade}. O \emph{racioci} predomina na \emph{matfhematfhica}; a experiencia \emph{sensivel}, nas \emph{sciencias physicas} ; a \emph{auctoridade}, nas \emph{sciencias} moraes, D\textasciicircum{} ahi o methodo racional, ou \emph{a-priori},\pgn{249}--- o methodo experimental, ou. \emph{a-posteriori}, --- o methodo da auctoridade. O methodo \emph{racional} leva-nos a certeza \emph{metaphysica} e \emph{absoluta}; O experimental, \& certeza \emph{physica} e \emph{condicional}; o da \emph{auctoridade}, á certeza \emph{moral} (tratando-se da \emph{auctoridade} humana). --- Yindicaremos brevemente o methodo das sciencias \emph{mathematicas}, \emph{physicas} c \emph{moraes}, reservanüo-nos tratar do methodo da \emph{Philosophia} no seguinte numero. !\par A) Methodo da mathematica. --- Na \emph{ocathematica}, predomina a \emph{synthese}. Na verdade, a \emph{mathematica} parte dos axiomas e das definições, que se formam por meio da absíracozo é da \emph{analyse}, e depois, pela \emph{symfhese} ou dedwecao, passa de uma quantidade para outra, demonstrando, pela comparação de ambas com uma terceira, que uma é egual a outra. --- Por isso a operação, a que se dá \emph{o} nome de \emph{analyse} mathematica, propriamente € uma \emph{synthese}.\par B) Methodo das sciencias physicas. --- Se na \emph{matematica} predomina a \emph{synthese}, nas sciencias-physicas-predom-na-a-analyse.---------A5-sciencias-plusicas são \emph{naturaes} ou experimeníaes. .\par a) \emph{Methodo} das \emph{sciencias} naturaes. --- ÀS sciencias waiwraes, como a \emph{Zoologia}, a \emph{Botanica}, \emph{a Geologia} etc., têm por fim apresentar a \emph{definição} dos seres, mas uma definição explicativa, scientifica, que penetre até é essencia, pondo em relevo a lei, que preside a geração e a constituição dos proprios seres. --- Para alcançar esse fim, as sciencias naturaes empregam a \emph{classificação}, visto que esta operação, se não for inteiramente artificial, exprime o que muitas coisas têm de commum e o que têm de proprio, e assim denota. a essencia, que é composta de dois elementos --- um \emph{generico} e \emph{conumwm}, outro \emph{differencial} e \emph{proprio}. (Veja se o que dissemos, ua Secção \emph{Primeira}, ácerca da \emph{classificação}).\par b) \emph{Methodo das sciencias experimeniaes}. --- Às Sciencias \emph{experimentaes} desenvolvem-se \emph{principalmente} pela inducção. --- Esta começa pela observação dos phenomenos, progride pela \emph{supposição}, ou \emph{hypothese}, e completa-se pela \emph{verificação} On experimentação.\par a) A observação recolhe as circumstancias, que estiveram \emph{presentes}, us que \emph{faltaram}, e as que variaram ; e com ellas forma tres tabuas, chamadas por Bacon --- tabua da presença. tabuu da ausencia, tabua dos \emph{graus}. Estas fabuas podem levar-nos ao conhecimento do nexo entre os phenomenos e a sua causa. Mas tal nexo não passa de uma \emph{Rypothese}.\par b) A \emph{Rypothese} é indispensavel na sciencia; porque serve para limitar os pontos da discussão, e para formar o objecto da \emph{verificação}.\par \emph{c}) \emph{A verificação}, OU experimentação, pode realisar-se por cinco methodos differentes: --- pelo methodo da concordancia (quando todas as cireumstancias, que costumam acompanhar um phenomeno, variam constantemente, a Í excepção de uma, esta é a causa do phenomeno), --- pelo da \emph{differença} (se "\emph{dl} um pbenomeno se realisa num só dos varios casos, que convém entre si em : b todas as circumstancias, exceptuada uma, esta é a causa do phenomeno), --- \emph{UR} Z3 pelo da concordancia e da \emph{differenga} (\emph{quando} se começa pelo methodo da 1 \emph{EN} concordancia, e o da \emph{differença} serve de contraprova), --- pelo das variações \emph{LS}\pgn{250}\emph{concomitantes} (um phenomeno, que augmenta ou diminue; augmentando on diminuindo uma causa, é effeito d' esta causa), --- dos \emph{restos} (se- de um phenomeno complexo se subtrahir \emph{a} parte que, por precedentes indueções, se averiguou ser effeito de uma determinada Circumstancia, o resto do phenomeno-é effeito-da- circumstancia;-que- ficou). :\par C) Methodo das sciencias moraes. --- Sciencia moraes são as que determinam as leis reguladoras dos costumes humanos. --- As sciencias moraes são muito complexas. Digamos alguma coisa ácerca das \emph{sciencias sociologica} e \emph{ltistorica}.\par 2) \emph{Methodo da sciencia sociologica}. --- O methodo d'esta sciencia é muito complexo e delicado, pois exige um espirito observador, dotado de senso pratico, de vistas elevadas e largas, de grande conhecimento do coração humano, como tambem.do paiz e da epocha. --- Não se deve facilmente deixar o antigo, quando é bom; mas tambem não se deve rejeitar uma, novidade, quando esta é exigida pela mudança das circumstancias sociaes. --- Em todo o caso dará prova de altissima subedoria e de generosa dedicação quem se inspirar num ideal francamente christão, que, entre todos, é o mais racional e o mais adequado ás necessidades do bomem.\par V) \emph{Methodo da historia}. --- \emph{A historia} é a base das sciencias moraes. O "melbodo da \emph{Astoria} deve evitar os dois systemas oppostos, que são o \emph{ypercriticismo} e a \emph{credulidade} cega..--- O \emph{hypercriticismo}, ao passo que acceita muitas bypotheses.como certas, rejeita, sem razão plausível, antiquissimas tradições, baseadas em solidos fundamentos.--- A \emph{credwulidade cega} acceita, sem exame, todas as tradições, reconhecendo-as como verdadeiras.}

%%>>> t204

\tese{204}{Indole do methodo da Philosophia.}
Applicando os principios expostos, é facil concluir que o metbodo, que deve seguir-se na Philosophia, é tambem o \emph{analytico-synihelico}. Por quanto, nas discussões philosophicas deve -seguir-se- aquelle \pg{249} \pg{250}  methodo, que é mais conforme com a indole da nossa razão.

Ora é proprio do homem elevar-se do \emph{effeito} para a causa, das \emph{propriedades} para a \emph{essencia} das coisas, dos \emph{phenomenos particulares} pora as \emph{leis geraes} da natureza (methodo \emph{analytico}). --- Mas não basta. Devendo a ordem \emph{logica} dos conhecimentos corresponder á ordem \emph{real} das coisas conhecidas, é necessario fazer em seguida o caminho inverso, \emph{descendo} da causa para o \emph{effeito}, . da. essencia para as \emph{propriedades}, das leis para os \emph{phenomenos} (methodo \emph{syntelico}). --- Todavia, dominando nesta sciencia a \emph{synthese} com o \emph{syllogismo}, diz-se vulgarmente que o methodo philosophico é \emph{synihetico}, ainda que em rigor seja \emph{analytico-synthelico} \footnote{As opiniões dos escriptores acerca de methodo \emph{philosophico} são tantas. quantas são as que enumerámos ácerca do methodo \emph{scientifico}; pois a Philosophia é a mais nobre entre as sciencias, que procedem pela luz natural da razão. Por isso ao methodo \emph{philosophico} deve applicar-se o que dissemos ácerca do methodo \emph{scientifico} em geral.\par Até aqui fallamos do methodo de invenção, isto é, d' aquelle methodo que nos leva ao descobrimento da verdade.\par Resta-nos dizer alguma coisa ácerca de methodo do \emph{ensino}, pelo qual o mestre por tal forma dirige as operações intellectuaes do discipulo, que este, pela força da sua razão \emph{natural}, chega a alcançar a sciencia. --- Ha tres opiniões a este respeito: --- a) uns dizem que o \emph{methodo} do ensino deve Ser exclusivamente \emph{synthetico}; ---- b) outros, que deve ser «\emph{nualytico-synthetico},\pgn{251}mas a \emph{syníhese} deve sempre preceder a \emph{analyse}; --- €) outros, que o methodo do \emph{ensino} deve corresponder ao methodo da \emph{invenção}, isto é, que as sciencias devem ser ensinadas pelo mesmo modo por que se formaram.\par Façamos uma breve critica d' estas opiniões.\par 2) A \emph{primeira} opinião, sustentando que o methodo do \emph{ensino} deve ser exclusivamente \emph{synthetico}, é inacceitavel. Por quanto, ninguem terá \emph{pleno} conhecimento de uma sciencia, se, embora conbeca o modo de inferir dos principios as conclusões, ignora a força dos primeiros principios, de que as conclusões dependem, Ora é por meio da analyse que se descobrem e conhecem os primeiros principios. Assim, para a perfeita acquisição das sciencias naturaes, é preciso conhecer as observações e experiencias, que se fizeram para a descoberta das varias leis pbysicas.\par b) A \emph{segunda} opinião, admittindo que o methodo \emph{synthketico} deve sempre preceder o analgtico, é tambem inacceitavel. Na verdade, quando o mestre produz a sciencia no discipulo, produl-a emquanto lhe propõe aquelle mesmo processo que elle (mestre) empregou para chegar ao descobrimento da verdade, e assim o discipulo repete a invenção feita pelo mestre. Ora este methodo é \emph{analytico}, e por isso é falso que o methodo \emph{suynthetico} deve sempre preceder o \emph{analytico}.\par €) A \emph{terceira}. opinião, sustentando que uma sciencia deve ser ensinada pelo mesmo methodo ou processo, por que foi \emph{descobería}, é verdadeira. Por quanto. o methodo de \emph{invenção} dilfere do methodo de ensino, como a. natureza differe da \emph{arte}. Ora a \emph{arte} é a imitação da natureza. Logo o methodo, por que se ensina uma sciencia, não deve differir \emph{do} methodo, por que ella foi \emph{descoberta} ou \emph{formada}. Assim a \emph{Physica} deve enzinar-xe por experiencias feitas sob os olhos dos diseipulos: a \emph{Mathematica}, \emph{pela}. \emph{dedueção} ou \emph{synthese}; \emph{a Philosophia}, pela \emph{analyse} e \emph{synthese}, \emph{etc}.}

ArTIGO [II

\emph{Termo do movimento} --- \emph{Sciencia}.

%%>>> t205

\pg{251} \tese{205}{Sciencia.}
\emph{Sciencia} é um \emph{systema de conhecimentos certos}, \emph{relativos} ds causas ou a \emph{razão sufficiente dos} entes, \emph{e deduzidos pela demonstração}. --- lóxpliquemos esta definição.

\emph{a}) \emph{Sciencia} é um \emph{systema}. A sciencia deve ser um complexo de verdades dispostas de modo, que uma esteja subordinada á outra e todas a uma suprema, de que recebam luz e força. --- Concdusóes desconnexas não formam uma \emph{sciencia}.

b) E' um systema de \emph{conhecimentos certos}. Nào ha sciencia sem \emph{certeza}, --- Todavia não é sempre necessaria, nem sempre é possivel, a certeza "\emph{metaphysica}; basta a \emph{physica} ou a \emph{moral}.

c) E" um systema de conhecimentos \emph{relativos} ds causas, ou d \emph{razão sufficiente dos entes}. Todos concordam em que o conhecimento \emph{scientifico} se distingue do \emph{vulgar}, porque este observa sómente os phenomenos, quando aquelle procura as'causas dos \pg{252}  phenomenos e explica a \emph{razão}, \emph{por que} uma coisa é o que \&, e não pode ser ou não é d'outro modo. ---- As causas ou as razões dos entes são \emph{proximas} e \emph{supremas} \footnote{De proposito: dissemos que a \emph{sciencia} é uma systema de conhecimentos certos, relativos ás \emph{casas}, \emph{ou a razão sufficiente} dos entes. Por quanto, se um ente não tem sobre si uma causa, de que dependa, poderá ser conhecido de um modo \emph{scientifico}, se for conhecido pela \emph{razão}, pela. qual é o que é, Assim Deus, sendo o Ser subsistente pela sua \emph{essencia}, será conhecido \emph{Scientificamente} pela mesma \emph{essencia}.\par Que 2 \emph{sciencia} deva occupar-se das causas dos entes, é coisa manifesta. Verdadeira sciencie é a que subministra um conhecimento perfeito. O conhecimento perfeito exige que a coisa seja presente á intelligencia, como é em ' si mesma, de modo que a ordem ideal corresponda plenamente \& ordem real. Ora a coisa é o que é pelas suas causas, ou pela sua razão sufficiente. Logo, se \& coisa não é conhecida por estas causas, on por esta razão, não se encontra na intelligeneia como é em si mesma, e por isso não é conhecida perfeitamente. Diz S. Thomaz: «Scire aliquid est perfecte cognoscere ipsum ; hoc autem est perfecte apprehendere eius veritetem. Eadem enim sunt principia esse rei, et veritatis ipsius. Oportet ergo scientem, si est perfecte co- . gnoscens, quod cognoscat causam rei scitae». (In I \emph{Post} «nal., \textbar{}. 4).}.

d) E' um systema de conhecimentos \emph{deduzidos pela demonstração}. Conhecimentos \emph{certos} são o resultado d'aquella especie de raciocinio, que parte de principios \emph{certos}. Essa especie de raciocinio é a \emph{demonstração} (\textasciicircum{}).\footnote{O termo \emph{sciencia} toma-se em tres sentidos: Em sentido mais lato, significa todo e qualquer \emph{conhecimento}. Em sentido mais restricto, denota o conhecimento de uma coisa \emph{pelas} suas causas. Em sentido proprio, indica um \emph{systema} de conhecimentos. ó\par Os positivistas pretendem que a. sciencia deva occupar-se só do \emph{positivo}, isto é; só do que-foi percebido pela experiencia sensivel, e não das causas. --- Esta pretensão é absurda. A sciencia socega e tranquillisa o nosso espirito. Mas quando poderá o espirito descançar? Só quando tiver conhecido OS porquês, as razões dos phenomenos, que observa. A sciencia, pois, não só não deve prescindir do conhecimento das causas, mas sem este conhecimento não é, nem pode chamar-se sciencia. \emph{Felix qui poluit} rerum \emph{cognoscere} causas ! repeté com o poela \emph{O} senso commum.\par Dizem os positivistas: «A experiencia só mostra o phenomeno, e não a causa d'elle». --- A esta dificuldade responde Tyndal, que é insuspeito: « A experiencia \emph{guia} sempre para uma coisa, que está além do phenomeno experimental. Ella produz sempre um effeito superior a si mesma, e toda a differença entre \emph{um} sabio eminente e um sabio mediocre consiste principal mente na \emph{faculdade} da \emph{extensão ideal}» (5º \emph{Conf}. \emph{sobre} a \emph{luz}).\par Declaremos a natureza. da. \emph{demonstração}, seus \emph{elementos} e \emph{especies}.\par a) Demonstração. --- \emph{Demonstração} é o raciocinio, que parte de premissas \emph{certas}. --- A \emph{demonstração}, pois, differre das outras especies de raciocinio,\pgn{253}nio quanto á \emph{forma}, mas quanto á \emph{materia}; porque as outras especies prescindem d'aquella certeza nas suas premissas. ' Advertimos que \emph{demonstração}, em geral, é o raciocinio pelo qual o espirito: progride de uma verdade conhecida para outra desconhecida. --- Divide-se em \emph{apodictica}, \emph{dialectica} € \emph{sophistica}, conforme parte de premissas \emph{certas}, \emph{provaveis}, ou \emph{falsas}. A \emph{apodictica} produz a \emph{ceríeza}: a dialectica, a \emph{opinião}; a \emph{sophistica}, O \emph{erro}. --- Só un \emph{apodictica} é strictamente demonstração.\par v) Elementos da demonstração. --- Os elementos da demonstração São \emph{quatro}, \& Saber: \emph{these}, \emph{principio}. \emph{meio} e \emph{conclusão}. --- A \emph{these} é a verdade, que deve ser demonstrada; --- o \emph{principio} é a proposição \emph{maior}, em que se funda a conclusão; --- o meio é \& proposição menor, exprimindo que u conchesão ; está contida no \emph{principio}; --- a conclusão é a propria \emph{these}, emquanto deduzida das premissas, em que estava contida. Ex.\par (These) A alma humana é immortal\par (Principio) O espirito é immortal ;\par (Meio) . Ora a alma humana é espirito;\par (Conclusão) Logo a alma humana é immortal.\par Digamos alguma coisa ácerca dos quatro elementos da \emph{demonstração}, -\par \emph{a}) A \emph{these} deve ser uma verdade, que \emph{deva} e \emph{possa} ser demonstrada; aliás a demonstração será inutil ou \emph{impossivel}. !\par \emph{b}) O \emph{principio}, em que se baséa a conclusão, deve ser \emph{immediatamente} evidente, ou deve fundar-se nos \emph{primeiros principios}.\par \emph{c}) O \emph{meio} deve indicar a causa proxima e “a, por que o prediado convém ao sujeito. !\par \emph{d}) A \emph{conclusão} não é sempre da mesma natureza. Quando a proposição maior é um \emph{axioma} e a menor é uma verdade de \emph{fucto} ou de experiencia, 3 conclusão é necessaria e absolutamente \emph{verdadeira}; como neste exemplo; \emph{Todo o effeito tem} a sua causa; \emph{ora} o mundo é \emph{effeito}; \emph{logo o} mundo \emph{tem} a sua causa. --- Quando a \emph{maior} é uma verdade de \emph{facto} e a menor É uma proposição \emph{singular}, a conclusão não é necessariamente verdadeira, como: \emph{Os mortos} não resuscitam ; \emph{ora Lazaro é morto}; \emph{logo Lazaro não reswscita}.\par c) Terminologia da demonstração. --- Na demonstração, sobretudo mathematica, empregam-se palavras, que precisam de explicação, --- São as seguintes :\par \emph{a}) \emph{Definição}. D' esta já tratámos na Logica \emph{formal}.\par \emph{b}) \emph{Proposição theoretica} e \emph{pratica}. --- \emph{A theoretica} allirma o que uma coisa é; ex. «o triangulo é uma figura». A \emph{pratica} declara \emph{o} que se póde ou deve fazer; ex. « devemos amar a patria».\par c) Axioma é uma proposição theoretica e indemonstravel, porque a sua verdade se torna evidente pela simples inspecção dos lermos. \textbar{}\par d) \emph{Postulado} é uma proposição pratica e immediatamente evidente pela simples explicação dos termos; ex. «de um ponto podemos tirar \emph{uma linha} i até ao outro ponto ». ;\par €) \emph{Theorema} é uma proposição theoretica, que se deve demonstrar, \emph{i P} Problema é uma proposição pratica, que se deve \emph{ou} póde resolver. i 9\emph{g}) Lemma é uma proposição, que uma sciencia presi a \emph{outra para lhe} servir de base nas demonstrações,\pgn{254}SN Li}

%%>>> t206

\pg{253} \tese{206}{Objecto da sciencia.}
\emph{Objecto} da sciencia é tudo aquillo, de que a sciencia se occupa. --- Divide-se em \emph{material} e \emph{formal}.

O \emph{material} é a propria \emph{coisa}, que a sciencia attinge: assim os \emph{m corpos} são o objecto \emph{material} da \emph{chimica}. O \emph{formal} é o \emph{aspecto}, \emph{M} sob o qual a sciencia considera o objecto \emph{material}; assim as

mudanças-.subsiamctoes,.que..se realisam nos. corpos, formam o" IS . objecto \emph{formal} da \emph{chimica} (*). e 907. Especies da sciencia. --- A sciencia póde ser --- \emph{a}) \emph{espe}-

\emph{mud A}) \emph{Corollario} é uma. proposição; que emana evidentemente de outra pro- \emph{EE} posição já demonstrada e não precisa de nova demonstração, \textbar{} == \emph{i}) \emph{Schkolio} é a nota que, accresc entada a uma proposição, ou a esclarece, \emph{oo} ou a defende das objecções dos adversarios. \emph{o} 4) Especies de demonstração, --- A demonstração pode-ser --- a) \emph{direciía} e \emph{o}» \emph{indirecta}. --- \emph{b}) \emph{a-priori} e \emph{a-posteriori}, --- \emph{c}) absoluta e \emph{relativa}, --- \emph{d}) \emph{pura}, \emph{bo} empirica, e mista. E ' \emph{a}) \emph{Demonstração directa} e \emph{indirecta}. --- É \emph{directa}, quando se prova que L- o predicado convem ao sujeito por causa de algumas propriedades, que este A possue; assim concluimos que a nossa alma é espiritual, porque é indepen- M---Á" \emph{AED} dente da materia, --- É \emph{indirecta}, quando se conclue que uma. proposição 3 : é verdadeira, porque, admittindo se a contraria, resultaria o absurdo; assim

provamos que os bons devem receber o" premio, pois, no caso contrario,

dar-se-hia o absurdo que Deus é injusto. "m \emph{b}) \emph{Demonstração a-priori} e \emph{a-posteriori}. --- É \emph{a-priori}, quando da causa \emph{X} descemos ao seu \emph{effeito}: assim pelo facto de a alma humana ser espiritual ---- (causa) dem onstramos que ella é immortal (\emph{effeito}). --- É \emph{a-posteriori}, quando \emph{DUM} do \emph{effeito} subimos a sua \emph{causa}; assim da existencia do mundo (\emph{effeito}) su- Lo bimos ao conhecimento da existencia de Deus (causa).

DL dq Estas duas especies de demonstração, reunidas, formam uma outra espe- T cie, que se chama circular ou \emph{regressiva}, e que se dá quando, depois de "0 termos \emph{subido} dos effeitos para a causa, descemos d' esta para aquelles. Este O processo é de grandissima vantagem e leva ao perfeito conhecimento das \emph{mE} coisas e das suas relações. \emph{Do c}) \emph{Demonstração} absoluta e relativa. ---Ê \emph{absoluta}, quando a conclusão

\emph{p} se infere. de premissas \emph{por si} certas ; assim a \emph{existencia} de Deus \emph{demousirazd} se pela \emph{contingencia} das coisas visiveis. --- É relativa, quando a conclusão --- se tira de premissas admittidas pelos adversarios, embora falsas; assim

refutamos os \emph{scepticos}, Servindo-nos dos seus principios. \emph{i}» \emph{d}) Demonstração \emph{pura}, \emph{empirica} e \emph{mista}. --- É pura. quando ambas as - premissas são proposições \emph{analyticas}. --- E \emph{empirica}, quando ambas as pre- . missas são proposições \emph{syntheticas}. --- 1 \emph{mista}, quando uma premissa é e analytica e outra é \emph{synthetica}. i (1) Por isso um objecto, que é \emph{materialmente} uno, pode ser \emph{formal}- \textbar{}\emph{j mente} multiplice; como tambem um objecto, que é materialmente multiplice, \textbar{} Um pode ser \emph{formalmente} uno. Assim o \emph{organismo animal}, embora seja matei--- \emph{rialmente} uno, é \emph{formalmente} multiplice, porque pode ser considerado sob "\emph{M} diversos aspectos --- na sua constituição, nos seus elementos, funecóoes, \textbar{}--- mudanças, etc. ; --- como, pelo contrario, o \emph{sol}, a lua, as \emph{plantas}, 08 \emph{homens}, \emph{u} embora sejam muitos objectos \emph{materiaes}, todavia constituem um só objecto . \emph{formal}, emquanto, pela sua \emph{côr}, se referem a um só sentido, --- 4 \emph{vista}, a \pg{255}  \emph{culativa} ou \emph{pratica}, --- \emph{b}) \emph{racional} ou \emph{experimental}, --- c) \emph{subordinada} ou \emph{subordinante}.

a) Sciencia especulativa ou pratica. --- É \emph{especulativa}, quando o seu fim unico e exclusivo é a contemplação da verdade. E' \emph{pratica}, quando considera a verdade, mas em ordem á operação. --- A sciencia \emph{pratica} distingue-se da \emph{arte}, pois esta limita-se a ensinar as regras para uma obra, e aquella expoe as \emph{razões} das mesmas regras \footnote{\emph{Arte} é O conhecimento do modo de fazer alguma coisa. Desenvolve-se como um \emph{habito}, que aperífeiqóa a natureza. --- A arte \emph{humana} 'é immensamente superior a dos \emph{animaes}; pois aquella é racional, consciente, progressiva, ao passo que esta é instinctiva, inconsciente, eslacionaria, --- Ha muitas e diversas \emph{artes}; mas podem reduzir-se a duas especies principaes --- \emph{Hberaes} e \emph{smechanicas}. As primeiras têm por objecto alguns actos da razão, como são o raciocinio, a, \emph{palavra}, etc. ; d' ahi a \emph{grammatica}, \emph{a eloguencia}, etc. As segundas applicam-se ás obras corporaes. --- Entre as artes \emph{Hberaes} e as \emph{mechanicas} collocam-se as bellas \emph{artes} (pintura, esculptura, musica, etc.), que procuram o bello, como as \emph{liberaes}, e tendem para o \emph{util}, como as \emph{mechanicas}.}.

b) Sciencia racional ou experimental. --- É racional, se tem por objecto as verdades conhecidas pela \emph{razão}; tal é a \emph{melaphysicu}, a \emph{moral}, a malthemalica, etc. --- X experimental, ou \emph{empiricz}, se tem por objecto os factos percebidos pela \emph{experiencia}; tal é a \emph{physiologia}, a \emph{biologia}, a \emph{chimica} etc. Q).

€) Sciencia subordinada ou subordinante. --- É \emph{subordinada} a que depende de outra. --- É \emph{subordinante} aquella, de que uma outra depende: =-Gmasciencia-póde-depender-de-outra-por duas. razões: ou porque o \emph{fim} de uma. é meio com. relação. ao \emph{fim} da outra, assim a \emph{pharmacia} depende da \emph{medicina}; --- ou porque uma aceeita, como \emph{principios}, as conclusões demonstradas na superior, d' este modo a \emph{astronomia} depende da \emph{mechanica} \footnote{As sciencias \emph{experimentaes} têm por \emph{objecto} os \emph{factos}, percebidos pela \emph{experiencia}. Mas estes factos não devem ser isolados ou simplesmente accumulados, como pretendem os \emph{materialistas} e Os \emph{positivistas}, mas hão de ser explicados, coordenados, interpretados pela razão, como pensam e ensinam os melhores naturalistas, por. ex., C. \emph{Bernard}, \emph{Pasteur}, ete. Esta explicação, coordenação e interpretação não pode realisar-se sem os primeiros principios da razão. Dá-se, pois, a sciencia do \emph{contingente}, mas só emquanto este se considera sob a luz ou reflexo do \emph{necessario} e do \emph{absoluto}.}.\footnote{\emph{A}' \emph{Metaphysica}, que é a sciencia dos \emph{primeiros principios}, estão e devem estar subordinadas \emph{todas} as sciencias, \emph{racionaes} e \emph{experimentaes}, SE quizerem guiar-se pelo senso commen. --- Não podemos conceber como possa ter havido e possa ainda haver espiritos tão desvairados, que expulsassem e pretendam \emph{expulsar} a \emph{Metaphysica} do numero das sciencias! --- Mas d'isto occupar-nos-bemos em breve,}

908. As sciencias distinguem-se \textbf{pelo} seu objecto \textbf{formal}.

a) Às sciencias distinguem-se pelo seu \emph{objecto}. Tudo o que se refere a uma coisa, distingue-se por ella; assim a \emph{vista} é disstineta do ouvido, porque a \emph{côr} é distincta do som. Ora as sciencias referem-se ao seu \emph{objecto}; pois este é o \emph{termo} da sciencia. Logo as sciencias distinguem-se pelo seu \emph{objecto}.

.b) As sciencias distinguem- -se pelo seu objecto \emph{formal}. Às sciencias devem distinguir-se' pelo objecto, a que ellas immediata e \emph{principalmente} se referem. Ora o objecto, a que as sciencias \emph{inimediata} e \emph{principalmente} se referem, é o \emph{formal}. Logo as Sciencias distinguem-se pelo seu objecto \emph{formal} \footnote{A nossa \emph{sciencia}, \emph{depende} realmente do \emph{objecto}. Uma sciencia sem um objecto, que se conheca, é um absurdo; porque, se conhecemos, conhecemos alguma coisa, que é o objecto e o termo do conhecimento. Portanto, a sciencia depende \emph{realmente} do objecto. --- E não é o contrario, isto é, o objecto não depende realmente da nossa sciencia. Por quanto, o conhecimento que temos de um objecto pode acabar, ou pode mesmo não ter existido, sem que por isso acabe:o objecto. --- Sendo o objecto o \emph{termo} da sciencia, é claro que a sciencia se distingue e se especifica pelo seu \emph{objecto}, como todo o caminho se distingue e se especifica pelo seu termo. --- Todavia, como provamos no texto, a sciencia distingue-se pelo objecto, emquanto este se considera \emph{formalmente}. D' onde se segue que a unidade ou a multiplicidade das sciencias deriva, --- não da unidade ou multiplicidade do objecto \emph{material}, pois muitos objectos »\emph{aferiaes} podem ser considerados sob um só aspecto e assim constituir uma só sciencia, --- mas da unidade ou multiplicidade do objecto \emph{formal}, pois um só objecto material pode ser considerado sob dilferentes \emph{aspectos} e constituir \emph{varias} sciencias. Assim, a philosophia é uma \emph{unica} sciencia, porque o seu objecto, ainda que \emph{matevialnenie} seja multiplice, todavia \emph{formalmente} € unico, porque o:mundo, as plantas, os animaes, o homem com as suas operações intellectuaes e moraes, são considerados sob um unico aspecto, --- nas suas causas supremas. A anatomia e a physiologia são duas sciencias distinctas, porque o objecto d' ambas --- o organismo animal, --- ainda que materialmente seja unico, todavia \emph{formalmente} é duplice, visto que a anatomia considera-o na sua constituição, e a physiologia nas suas funcções,\pgn{257}------------------}.

909. Classes supremas das sciencas. --- As classes supremas das sciencias são cinco, a saber: \emph{Physica}, \emph{Mathematica}, \emph{Metaphysica}, \emph{Logica} e \emph{Moral}.

Na verdade, as seiencias distinguem-se pelo seu objecto \emph{formal}. O: objecto-formal-das-sciencias-é-o-abstracto; pois tudo o que a nossa intelligencia - conhece, conhece-o emquanto é \emph{abs}-

\emph{tracto}. Logo as classes das sciencias serão tantas, quantos forem os differentes graus de \emph{abstracção}. \textbar{}

Ora os differentes graus de \emph{abstracção} São cinco. --- É infimo o grau de abstracção, quando prescindimos da materia indivi \emph{dual}, e só consideramos a materia: \emph{commu}, sujeita ás mudanças sensiveis : d'abia \emph{Physica}. --- É superior a este, quando prescindimos tambem da ujateria commum-e-sensivel, e. só. con-. sideramos a \emph{quantidade} : d'ahi a \emph{Mathematica}. --- É mais alto ainda, quando prescindimos da materia commum e da quantidade, e consideramos O ente immaterial na sua realidade \emph{objectiva}, isto é, como existente em si e sem dependencia do nosso pensamento : d'abi a \emph{Metaphysica}. --- É altissimo o grau de abstracção, quando prescindimos de toda a realidade \emph{objectiva}, e consideramos o ente \emph{immaterial} --- ou na sua entitade ideal, isto é, na entidade com que existe na \emph{intelligencia}: A ahi a \emph{Logica}: --- ou na Sua entidade \emph{moral}, isto é, na entidade com que existe na vontade : d' ahi a \emph{Moral}. :

Logo as classes supremas das sciencias são cinco (\emph{Q}).\footnote{Esta classificação, que é devida aos scholastico 2 das muitas que se têm apresentado. Na verdade, ella é racional no seu \emph{fundamento} € completa na SUA divisão. --- É \emph{racional} no seu \emph{fiundanento}. Por quanto, o aspecto geral, sob o qual a sciencia considera qualquer objecto, é a su \emph{abstracção}, pois vão ha sciencia do \emph{concreto}; € por isso os diversos graus da abstracção devem distinguir ás diversas classes das sciencias. --- É \emph{completa} na Sua divisão ; pois DÃo hà sciencia, que não se reduza a alguma das cinco classes. Na verdade, toda a sciencia ou é especulativa, ou € \emph{pratica}. Se é especulativa, está subordinada á Pbysica, ou \& Mathematica, ou \& Metaphysica. Se é pratica, está subordinada é Logica ou a Moral, conforme aperfeicgoa os actos da intelligencia ou 0S da vontade.\par É immenso o empenho, que OS modernos philosophos mostram para resolver o problema da \emph{classificação das sciencias}, --- problema que tem apenas uma importancia secundaria. .\par As principaes classificações, apresentadas \emph{pelos} escriptores, são as de Bacon, de Ampére, de A. Comte, e de 4. \emph{Spencer}.\par a) Bacon tomou por base da sua classificação as tres faculdades da alma, que são a \emph{Imaginação}, a memoria ea \emph{rezão}, € assim divide todas as sciencias em ires classes. A imaginação corresponde a \emph{poesia}; a memoria,\par a \emph{historia}, que é \emph{civil} ou natural; a \emph{ração}, \emph{a philosophia} e a \emph{theologia}.\par Esta classificação não póde acceitar-se. A sciencia é um systema de\par conhecimentos. Ora os conhecimentos são diversos, emquanto se referem a\par diversos objectos \emph{formaes}. Às sciencias, pois, distinguem-se pelos objectos\par \emph{formaes}, e não pelas faculdades cognitivas. --- Além d' isso, \& \emph{poesia} não é\par \emph{sciencia} ; á \emph{razão} € 9 memoria São indispensaveis em todas as sciencias.\par 1) Ampére Loma. por base da sua classificação O \emph{objecto material}, e\par assim divide as sciencias em \emph{cosmologicas} € \emph{noologicas}, segundo consideram\par Ya. 1. 18\pgn{258}objectos materiaes OU espiritunes. --- AS sciencias \emph{cosmologicas} dividem-se em \emph{cosmologicas} propriamente \emph{ditas} e em \emph{phys:ologicas}. \emph{Às cosmologicas propriamente ditas} são a \emph{maíhematica} e a \emph{physica}. \emph{Às physiologicas} são as naturaes € ás \emph{medicas}. --- As sciencias \emph{noologicas} dividem-se em \emph{noologicas propriamente ditas} e em sociaes. ÀS mnoologicas propriamente \emph{ditas} sio as \emph{Dphilosophicas} e as \emph{nootecknicas}. \emph{ÀS sociaes} são as \emph{historicas} e as \emph{politicas}.\par Esta classificação tambem é inadmissivel. As sciencias, como dissemos, distinguem-se, --- não pelos seus objectos swaferiaes, visto que diversas sciencias podem ter um e o mesmo objecto material, --- mas sim pelos seus objectos \emph{formaes}; assim o \emph{komem} é o objecto da anthropologia, da medicina e da anatomia, porque pode ser considerado sob diversos aspectos.\par , €) \emph{Comte} divide as seieneids ein \emph{aostractas} € concretas. --- ÀS \emph{concretas}, como a \emph{mineralogia}, \emph{a botanica}, \emph{a coologia}, segundo \emph{Comte}, não estão ainda constituidas; e por isso não as classifica. ---- As sciencias \emph{abstractas}, que são fundamentaes, formam \emph{seis} grupos e, \emph{dispostas} na sua \emph{ordem hierarchica}, São : a mathematica, a \emph{astronomia}, a \emph{physica}, a chimica, a \emph{biologia} e a sociologia. --- Esta classificação baséa-se na ordem \emph{logica} e \emph{chronologica} das sciencias. --- Baséa-se na ordem \emph{logica}. Por quanto, é proprio da intelligencia passar do mais simples e abstracto \textbf{para} o mais complexo e concreto. Ora, nesta classificação, a primeira scieneja, que é a \emph{mathematica}, é mais simples e abstracta que a segunda --- a \emph{Gastronomia} ; esta o é mais que a terceira --- a \emph{physica}, etc. --- Baséa-se na ordem \emph{chronologica}. Com effeito, a 'primeira sciencia, que se constituiu, foi a \emph{mathematica}; depois a \emph{astronomia}; em seguida a \emph{physica}, \emph{a chimica}, \emph{a biologia} e, por ultimo, a \emph{sociologia}, que é o termo de todas. . .\par A \emph{classificação}, feita por Comte, é falsa pelas seguintes razões:\par 1.º) \emph{Comte}, na sua \emph{classificação}; parte das sciencias mais abstractas para chegar ás mais complexas; porque opina que ás leis abstractas encerram a mais completa explicação dos f\&ctos concretos. Ora isto é falso. Na verdade, embora um [acto concreto se resolva, \emph{pela analyse}, numa lei abstracta, comtudo uma lei abstracta não póde dar uma explicação \emph{completa} dos factos concretos; pois ficaria sempre para explicar \emph{Como}, \emph{por} quem e \emph{porque} leis abstractas se uniram de um certo modo \textbf{para dar} em resultado um determinado facto concreto. Por outras palavras: quando \emph{Comte} quer explicar todos os phenomenos pelas propriedades elementares da materia, não só confunde a \emph{condição} com a cuusa, mas cahe no absurdo de querer encontrar a Tazão do mais no menos, do superior no inferior, do ente no nada.\par 2.2) \emph{Comte} ensina que as sciencias \emph{positivas} se constituiram \emph{Successiva mente}, seguindo a ordem \emph{chronologica}, indicada por elle. --- Ora isto é falso. Por quanto, como observa \emph{H}, \emph{Spencer}, as sciencias não se desenvolvem \emph{successivamente}, mas \emph{simultaneamente}, e nenhuma das que oecupam os primeiros\pgn{259} graus, tanto na ordem logica como na \emph{chronologica}, podia progredir sem o auxilio das que estão collocadas nos ultimos graus, --- De mais, \emph{Comte} pretende que a \emph{astronomia} precedeu, na sua constituição, a \emph{physica}. Ora isto não é certo: porque, para determinar a posição de uma estrelia, é necessario ter em conta as condições, que podem modificar o phenomeno da visão, --- condições, que formam o objecto da \emph{physica}. --- Finalmente, é impossivel que as sciencias \emph{abstractas} precedessem, na sua constituição, as sciencias concretas; pois que, na realidade, o \emph{concreto} precede sempre o abstracto, porque este forma-se pela generalisação d' aquelle. Nem. se diga. que o nosso espirito passa do \emph{simples} para o \emph{composto} e do abstracto pura o concreto; porque este processo faz-se na ordem racional, Mas não Da ore \textgreater{} dem \emph{experimental}; e todos sabem que os \emph{positivistas} só acceitam o methodo experimental (Cf. Ferraz, \emph{Socialisme}, \emph{naturalisme et positivisme}, ch. Vi).\par 3.3) \emph{Comte} e us positivistas, ensinando que 8Ó pode ser objecto da. seiencia o \emph{positivo}, isto é, o que pode perceber-se pela experiencia sensivel, exeluem do quadro das sciencias a \emph{Metaphysica}. --- Mais tarde provaremos que a Metaphysica é sciencia. Agora dizemos que os \emph{positivistas}. se contradizem ; porquanto, ao passo que só querem reconhecer como certo o que foi verificado pela experiencia sensivel, depois adimittem, como certas, muitas coisas, que escapam ao alcance dos sentidos. Invesligam a natureza do eifeilo e da causa, da força e da materia, do meio e do fim: lentam estabelecer leis immutaveis; procuram a certeza ácerca das coisas sensiveis c das leis universaes : empregam o syllogismo e a inducção ; analysam, gene- - ralisam, applicam todos os processos do methodo discursivo. Os positivistas, pois, excluem a \emph{Metaphysica} do quadro das sciencias, ao passo que aeceitam muitas coisas, que formam o objecto da \emph{Metaphysica}. Com razão. disse \emph{P}. \emph{Janet} que O \emph{positivismo} não tem de \emph{positivo} senão o nome (\emph{Phil}. \emph{fran}). Nem podia ser d' outro modo; pois, como diz o proprio \emph{Fowillee}, positivista : «o homem é um \emph{animal metaphysico}, e a Metaphysica ha de durar em quanto existirem cerebros humanos » (\emph{I}/ \emph{Avenir de} la, \emph{Métaphysique}).\par d) Spencer divide as sciencias em «\emph{bstracias}. \emph{abstracto-concretas}, e \emph{toncrelas}. --- A's \emph{abstractas} reduzem-se a \emph{logica} eu \emph{authematico}: ás \emph{abstracto-concretas}, \emph{a physica}, \emph{a chimica}, etc. ; ás \emph{concretas}, a astronomia, a \emph{geologia}, \emph{a biologia}, \emph{á psychologia}, a \emph{sociologia}, ete.\par Esta classificação tambetn não pode acceitar-se. Por quanto, a divisão da sciencia em \emph{abstracta} e concreta não é racional: pois toda a sciencia, sendo fructo da \emph{generalisação}, é mais ou menos \emph{abstracta}. Esta observação vale tambem contra a classilicação de \emph{Comte}. --- Além d'isso, o objecto \emph{Tnaterial} não pode tonfar-se como base de \emph{classificução},.}

910. --- Reducção das sciencias a \textbf{estas} classes, --- Nào ha Sciencia, que se não possa reduzir a \textbf{alguma} d'estas classes.

a) A sciencia \emph{physica} abrange --- a \emph{physica}, que estuda as mudanças \emph{accidentaes} dos corpos, --- a \emph{chimica}, que se occupa das mudanças \emph{substanciaes} dos mesmos Corpos, --- a. \emph{astronomia}, que estuda os astros, --- a \emph{geologia}, que considera a terra na \pg{259}  Sua constituição intima, --- a \emph{physiologia}, que tem por objecto as funegóes regulares do organismo animal, --- a biologia, que tem por objecto .a vida em geral, etc.

b) A sciencia \emph{mathematica} comprehende --- a \emph{arithmetico}, que é a parte elementar da sciencia dos numeros, --- a \emph{algebra}, que simplifica e generalisa as questões relativas aos numeros, --- a \emph{geometria}, que estuda as formas, as, propriedades e as medidas da extensão, etc.

€) A sciencia \emph{metaphysica} abrange --- a \emph{ontologia}, que trata do ente em geral, --- a \emph{cosmologia}, que se occupa do mundo, --- a \emph{anthropologia}, . que estuda o homem, ---.a £\emph{heodiceo}, que procura conhecer..a existencia . e..a á perfeição do Ente Infinito com a luz natural da razão.

d) A sciencia \emph{logica} comprehende --- a \emph{glotlologia}, que considera a origem e as vicissitudes da linguagem, --- \emph{a logica}, que estuda o raciocinio para a consecução da verdade, ete.

e) A sciencia moral comprehende --- a \emph{ethica}, que trata das operações da nossa vontade, emquanto se referem ao fim, --- o \emph{direito natural}, que se occupa. das leis, impostas aos homens por Deus, Auctor da natureza, --- o \emph{direito positivo}, qué se occupa das leis impostas pela auctoridade legitima, --- a \emph{pedagogia}, que tem por objecto a educação, etc. --- A \emph{sociologia} reduz-se ao \emph{direito}, de que é parte. --- O direito \emph{positivo} subdivide-se em \emph{civil}, penal, \emph{constitucional}, inter nacional, ete \emph{C}"

! \textasciicircum{}

2 As seguintes observações são importantes.

1.º) Não pretendemos fazer a enumeração completa de todas as sciencias constituidas, mas apenas indicar o modo por que ellas poderiam ser reduzidas a alguma das cinco classes.

2») As classes supremas das sciencias, embora sejam distinctus e irreductiveis (pois os seus objectos \emph{formaes} são irreductiveis), todavia, sob differentes aspectos, estão sobordinadas umas ás outras, Assim, é principal --- a \emph{Moral}, emquanto descobre o \emph{fim ultimo}, que é a regra de todos os actos da vida; --- a \emph{Logica}, emquanto é o \emph{instrumento} de todas as sciencias --- a \emph{Metaphysica}, emquanto trata do \emph{absoluto}, do \emph{immutavel} e, sobretudo, do \emph{Ente supremo}; --- a \emph{Mathematica} é a \emph{Physica}, emquanto são muito uteis : para as artes, pata a industria, e para a vida material da. sociedade. --- A todas estas sciencias humanas, que: se guiam pela luz da razão, está absolutamente superior a Theologia revelada, que se guia pela luz da Fé, e que trata das mais altas e consoladoras verdades. Entre esta e aquellas existe e não pode deixar de existir a mais perfeita harmonia. Deus é um centro de luz, de que partem, como. outros tantos raios, todas as verdades e todas as sciencias. Por isso Deus é chamado o \emph{Senhor das sciencias}.

\pg{254} \pg{256} \pg{257} \pg{258} \pg{260}
